\documentclass[11pt]{amsart}
\usepackage[margin=1.1in]{geometry}
\usepackage{amsmath,amssymb,amsthm,mathrsfs,mathtools}
\usepackage{enumitem}
\usepackage{microtype}
\usepackage{hyperref}
\usepackage[nameinlink,capitalize]{cleveref}
\hypersetup{colorlinks=true,linkcolor=blue,citecolor=blue,urlcolor=blue}
\allowdisplaybreaks
\pdfstringdefDisableCommands{%
  \def\({}%
  \def\){}%
  \def\tilde#1{#1}%
  \def\mathfrak#1{#1}%
  \def\mathrm#1{#1}%
  \def\mathbb#1{#1}%
  \def\tau{tau}%
  \def\lambda{lambda}%
  \def\infty{infty}%
}

\theoremstyle{plain}
\newtheorem{theorem}{Theorem}[section]
\newtheorem{lemma}[theorem]{Lemma}
\newtheorem{proposition}[theorem]{Proposition}
\newtheorem{corollary}[theorem]{Corollary}
\theoremstyle{definition}
\newtheorem{definition}[theorem]{Definition}
\newtheorem{assumption}[theorem]{Assumption}
\theoremstyle{remark}
\newtheorem{remark}[theorem]{Remark}

\newcommand{\dd}{\mathrm d}
\newcommand{\RR}{\mathbb R}
\newcommand{\C}{\mathbb C}
\newcommand{\supp}{\mathrm{supp}}
\newcommand{\nn}{\nabla}
\newcommand{\tr}{\mathrm{tr}}
\newcommand{\fint}{\mathop{\ooalign{\(\int\)\cr\hidewidth\(-\)\hidewidth\cr}}}

\title[Type II Regularity]{Type II Regularity for Navier--Stokes}
\author{Guillem Duran Ballester}
\email{guillem@fragile.tech}
\date{}

\begin{document}

\begin{abstract}
We establish a local Type II exclusion criterion for the three-dimensional
Navier--Stokes equations near a possible singular point.  From a positive
Caffarelli--Kohn--Nirenberg concentration sequence we select renormalized local
windows around a retained core and analyze the resulting alternatives by
compactness, pressure reconstruction, and local energy estimates.  The argument
combines a compact single-core exclusion criterion, a repaired-gauge
representation for nondegenerate concentration cores, local Calderon--Zygmund
pressure control, a Caccioppoli estimate on compact cylinders, multibubble and
cascade reductions, and scale-collapse cost estimates.  The conclusion is
conditional on the analytic inputs and rigidity hypotheses stated in the paper;
no additional stationary or self-similar classification theorem is used
implicitly.
\end{abstract}

\maketitle

\setcounter{tocdepth}{2}
\tableofcontents


\clearpage
\section{Introduction}

The three-dimensional Navier--Stokes regularity problem is governed by the
interaction between local energy control, scale-critical concentration, and the
possible formation of singular structures at arbitrarily small parabolic
scales.  Since Leray's construction of global weak solutions
\cite{Leray1934}, the suitable weak-solution framework and the partial
regularity theory of Scheffer and Caffarelli--Kohn--Nirenberg have provided the
basic local language for this interaction
\cite{Scheffer1976,CaffarelliKohnNirenberg1982,Lin1998}.  In that theory,
singularity formation is detected by scale-invariant velocity-pressure
quantities, and the critical \(L^3\) scale plays a central role alongside the
local energy inequality, pressure reconstruction, and compactness.  The
analytic background also includes the classical local and critical-space
theory for Navier--Stokes
\cite{Ladyzhenskaya1969,Temam1977,ConstantinFoias1988,Kato1984,LemarieRieusset2002},
the \(L_{3,\infty}\) regularity theorem of Escauriaza--Seregin--Šverák
\cite{EscauriazaSereginSverak2003}, and the rigidity theorem of
Nečas--Růžička--Šverák for Leray self-similar profiles
\cite{NecasRuzickaSverak1996}.

This paper studies the local Type II alternative near a possible singular point.
Starting from a positive Caffarelli--Kohn--Nirenberg concentration sequence, we
select renormalized local windows around a retained core and analyze all
branches that can arise from the scale-center dynamics.  The result is a local
Type II exclusion theorem in the class of alternatives covered by the analytic
inputs stated in the body of the paper.  It is deliberately formulated as a
conditional PDE theorem: pressure reconstruction, repaired-gauge compactness,
terminal profile decoupling, cost comparison, and the relevant rigidity
hypotheses are assumed or proved explicitly where they are used.  No additional
stationary, ancient, or self-similar classification theorem is invoked
implicitly.

The proof begins with local concentration and the Type I/Type II separation.
For a nondegenerate retained core, the solution is rewritten in local
scale-translation variables.  The repaired-gauge representation produces a
renormalized Navier--Stokes equation with controlled modulation terms, local
pressure reconstruction modulo functions of time, and compact-window suitable
estimates.  Within this representation, a compact single-core branch with
positive CKN density and vanishing localized dissipation is impossible: local
compactness gives a spatially constant limit, the zero constant contradicts
pressure-stable persistence of the positive core, and a nonzero constant gives a
bounded selected-window limit incompatible with the Type II condition.

The remaining sections remove the ways in which a Type II branch could avoid
that compact single-core mechanism.  Caccioppoli estimates convert finite outer
velocity-pressure control into inner windowed \(H^1\) control, so rough-core
loss is returned to an explicit compact-cylinder alternative.  Multibubble and
cascade arguments separate active profiles, same-point cascades, and
gauge-degenerate branches from a single retained core.  Scale-collapse branches
are treated by localized cost identities: finite scale-collapsing cost, and the
corresponding finite absolute cost, are incompatible with a nonvanishing
retained core, while thick autonomous windows lead only to the stated
modulation, compactness, or rigidity-covered stationary alternatives.  The
terminal profile and corrected monotonicity sections then turn these local
reductions into a branchwise classification, and the final assembly eliminates
each covered Type II alternative.

The conclusion should therefore be read in its local and conditional sense.  In
the covered class, a positive local Type II concentration sequence cannot
survive the repaired-gauge, compactness, multibubble, rough-core, and
scale-collapse reductions.  What remains outside the theorem is not hidden in
the proof: it is precisely the failure of one of the stated analytic inputs, or
one of the explicitly named Liouville or rigidity problems for terminal
autonomous states.

\clearpage
\section{Further Local Reductions for Rough Cores, Scale Collapse, and Cascades}
\label{sec:s2-s5-pde-reductions}

This section records the remaining local reductions used to separate the
rough-core, scale-collapse, multibubble, and same-point cascade alternatives.
All statements are conditional analytic statements for represented Type II
branches.

\subsection{Vorticity-Controlled Rough-Core Alternative}

Let \((V,P,a,b)\) be a repaired-gauge renormalized Navier--Stokes branch,
\[
  \partial_\tau V+(V\cdot\nabla)V+\nabla P
  =
  \nu\Delta V+a(\tau)(V+y\cdot\nabla V)+b(\tau)\cdot\nabla V,
  \qquad \nabla\cdot V=0,
\]
and set \(\Omega=\nabla\times V\).  Then, in distributions,
\[
  \partial_\tau\Omega+(V\cdot\nabla)\Omega-(\Omega\cdot\nabla)V
  =
  \nu\Delta\Omega+a(\tau)(2\Omega+y\cdot\nabla\Omega)
  +b(\tau)\cdot\nabla\Omega .
\]
At the smooth approximation level, a localized enstrophy estimate contains
\[
  \int(\Omega\cdot\nabla V)\cdot\Omega\,\chi_R^2 .
\]
The standard bound is
\[
  \left|
  \int(\Omega\cdot\nabla V)\cdot\Omega\,\chi_R^2
  \right|
  \le
  C\|\Omega\chi_R\|_{L^2}^{1/2}
  \|\nabla(\Omega\chi_R)\|_{L^2}^{3/2}
  \|\nabla V\|_{L^2(B_{2R})}.
\]
After Young's inequality the differential inequality has a coefficient
\[
  C_\nu\|\nabla V\|_{L^2(B_{2R})}^4
\]
in front of the localized enstrophy.  Thus the vorticity equation does not by
itself close a uniform enstrophy-window estimate without an independent
gradient-window estimate.  The vorticity argument below is therefore a conditional
rough-core exclusion, while the primary rough-core exclusion remains the
Caccioppoli estimate in \cref{paper7:prop:uniform-windowed-gradient}.

\begin{lemma}[Local div-curl estimate]
\label{lem:s2-local-div-curl}
For each \(R>0\) there is \(C_R<\infty\) such that every
\(U\in L^2(B_{2R};\mathbb R^3)\) satisfying
\[
  \nabla\cdot U=0\quad\text{in }\mathcal D'(B_{2R})
\]
and
\[
  \nabla\times U\in L^2(B_{2R};\mathbb R^3)
\]
belongs to \(H^1(B_R;\mathbb R^3)\), with
\[
  \|\nabla U\|_{L^2(B_R)}
  \le
  C_R\left(
  \|\nabla\times U\|_{L^2(B_{2R})}
  +\|U\|_{L^2(B_{2R})}
  \right).
\]
\end{lemma}

\begin{proof}
Choose \(\chi\in C_c^\infty(B_{2R})\) with \(\chi\equiv1\) on \(B_R\), and set
\[
  W=\chi U.
\]
After extension by zero, \(W\in L^2(\mathbb R^3)\) is compactly supported.
Moreover,
\[
  \nabla\times W
  =
  \chi\nabla\times U+\nabla\chi\times U\in L^2(\mathbb R^3),
\]
and, using \(\nabla\cdot U=0\),
\[
  \nabla\cdot W=\nabla\chi\cdot U\in L^2(\mathbb R^3).
\]
For compactly supported distributions with \(W\), \(\nabla\times W\), and
\(\nabla\cdot W\) in \(L^2\), the Fourier div-curl identity gives
\[
  \|\nabla W\|_{L^2(\mathbb R^3)}^2
  =
  \|\nabla\times W\|_{L^2(\mathbb R^3)}^2
  +\|\nabla\cdot W\|_{L^2(\mathbb R^3)}^2 .
\]
Indeed, Plancherel and
\[
  |\xi|^2|\widehat W(\xi)|^2
  =
  |\xi\times\widehat W(\xi)|^2
  +|\xi\cdot\widehat W(\xi)|^2
\]
yield the identity after integration in \(\xi\).  Hence \(W\in H^1(\mathbb R^3)\).
Since \(W=U\) on \(B_R\),
\[
  \|\nabla U\|_{L^2(B_R)}
  \le
  \|\nabla W\|_{L^2(\mathbb R^3)}
  \le
  C_R\left(
  \|\nabla\times U\|_{L^2(B_{2R})}
  +\|U\|_{L^2(B_{2R})}
  \right).
\]
\end{proof}

\begin{lemma}[Critical \(L^3\) control gives local \(L^2\) control]
\label{lem:s2-l3-local-l2}
If
\[
  \sup_{\tau\ge\tau_0}\|V(\tau)\|_{L^3(\mathbb R^3)}\le M,
\]
then, for every \(R>0\),
\[
  \sup_{\tau\ge\tau_0}\|V(\tau)\|_{L^2(B_R)}^2
  \le C R M^2 .
\]
\end{lemma}

\begin{proof}
Hölder's inequality gives
\[
  \int_{B_R}|V(y,\tau)|^2\,dy
  \le
  |B_R|^{1/3}
  \left(\int_{B_R}|V(y,\tau)|^3\,dy\right)^{2/3}
  \le C R M^2 .
\]
Taking the supremum in \(\tau\) proves the estimate.
\end{proof}

\begin{proposition}[Vorticity-window control implies local windowed \(H^1\)]
\label{prop:s2-vorticity-window-h1}
Assume \(V\) is a represented divergence-free branch,
\[
  \sup_{\tau\ge\tau_0}\|V(\tau)\|_{L^3(\mathbb R^3)}\le M,
\]
\(\Omega=\nabla\times V\) in distributions, and
\[
  \forall R>0:\quad
  \sup_{T\ge\tau_0+1}
  \int_T^{T+1}\|\Omega(\tau)\|_{L^2(B_{2R})}^2\,d\tau<\infty .
\]
Then
\[
  \forall R>0:\qquad
  \sup_{T\ge\tau_0+1}
  \int_T^{T+1}\|V(\tau)\|_{H^1(B_R)}^2\,d\tau<\infty .
\]
\end{proposition}

\begin{proof}
Fix \(R>0\).  By \cref{lem:s2-local-div-curl}, for almost every \(\tau\),
\[
  \|\nabla V(\tau)\|_{L^2(B_R)}^2
  \le
  C_R\left(
  \|\Omega(\tau)\|_{L^2(B_{2R})}^2
  +\|V(\tau)\|_{L^2(B_{2R})}^2
  \right).
\]
Integrating over \([T,T+1]\), the vorticity term is uniformly bounded by the
assumed vorticity-window control, while the local \(L^2\) term is uniformly
bounded by \cref{lem:s2-l3-local-l2}:
\[
  \int_T^{T+1}\|V(\tau)\|_{L^2(B_{2R})}^2\,d\tau
  \le C_R M^2 .
\]
Therefore
\[
  \sup_{T\ge\tau_0+1}
  \int_T^{T+1}\|\nabla V(\tau)\|_{L^2(B_R)}^2\,d\tau<\infty .
\]
The same local \(L^2\) bound gives
\[
  \sup_{T\ge\tau_0+1}
  \int_T^{T+1}\|V(\tau)\|_{L^2(B_R)}^2\,d\tau<\infty ,
\]
and the asserted windowed \(H^1\) bound follows.
\end{proof}

\begin{corollary}[Vorticity-controlled rough-core exclusion]
\label{cor:s2-vorticity-rough-core}
Let a represented Type II branch have positive finite critical mass and local
vorticity-window control as in \cref{prop:s2-vorticity-window-h1}.  Then the
rough-core alternative, namely failure of local windowed \(H^1\) control, cannot
hold.
\end{corollary}

\begin{proof}
The hypotheses give the local windowed \(H^1\) bound by
\cref{prop:s2-vorticity-window-h1}.  This is the negation of the rough-core
alternative.
\end{proof}

Consequently, a represented bounded-critical-norm rough-core branch must fail
at least one of the following analytic hypotheses: local vorticity-window control,
the representation or critical-mass hypotheses used before the vorticity
argument, bounded modulation, Caccioppoli regularity, or one of their
integrated refinements.

\subsection{Scale-Collapse Autonomous Reduction}

Let a represented repaired-gauge Type II branch be in the remaining
scale-collapse class when it satisfies the representation hypotheses, positive
finite critical mass, a nonzero localized \(L^2\) core floor, bounded
modulation, local windowed \(H^1\) control, \(\lambda(\tau)\to0\), and the
scale-collapse drift obstruction
\[
  \int_{\tau_0}^{\infty}a_-(\tau)M_R(\tau)\,d\tau=\infty,
  \qquad
  M_R(\tau)=\int |V(y,\tau)|^2\phi_{R(\tau)}(y)\,dy .
\]
The core floor supplies constants \(m_0>0\), \(\tau_1\), and moving cutoffs such
that
\[
  M_R(\tau)\ge m_0
  \qquad\text{for a.e. }\tau\ge\tau_1 .
\]
This is the remaining scale-drift class; it is distinct from the radiative and
rough-core alternatives.

\begin{definition}[Thick negative-drift windows]
\label{def:s3-thick-negative-windows}
The branch has thick negative-drift windows if there exist \(T>0\), \(c>0\),
and \(\tau_n\to\infty\) such that
\[
  \frac1T\int_{\tau_n}^{\tau_n+T}a_-(\tau)M_R(\tau)\,d\tau\ge c
  \qquad\text{for every }n.
\]
If \(M_R\le M_1\) on these windows, then
\[
  \frac1T\int_{\tau_n}^{\tau_n+T}a_-(\tau)\,d\tau\ge \frac{c}{M_1}.
\]
\end{definition}

The divergence of \(a_-M_R\) alone does not imply the existence of
\cref{def:s3-thick-negative-windows}; a bounded nonnegative function can have
infinite integral while its unit-window averages tend to zero.  Without thick
windows, the branch remains in the thin-drift alternative rather than entering
a self-similar regime.

\begin{definition}[Autonomous modulation limit]
\label{def:s3-autonomous-modulation-limit}
On fixed windows \([\tau_n,\tau_n+T]\), the branch has an autonomous modulation
limit \((a_\infty,b_\infty)\) if, after translation,
\[
  a(\tau_n+\cdot)\to a_\infty,
  \qquad
  b(\tau_n+\cdot)\to b_\infty
\]
in a coefficient topology strong enough to pass the products
\[
  a(V+y\cdot\nabla V),
  \qquad
  b\cdot\nabla V
\]
to the distributional limit.  Strong \(L^1(0,T)\) convergence is sufficient.
The parameter \(a_\infty\) may be negative, zero, or any value allowed by the
bounded modulation range.  The case \(a_\infty<0\) is a generalized
self-similar regime, while \(a_\infty=0\) gives a stationary Navier--Stokes
regime, possibly with constant translation drift \(b_\infty\).
\end{definition}

\begin{lemma}[Thick autonomous drift gives a negative scale parameter]
\label{lem:s3-thick-negative-parameter}
Assume \cref{def:s3-thick-negative-windows} holds on windows
\([\tau_n,\tau_n+T]\), \(M_R\le M_1\) on those windows, and
\[
  a(\tau_n+\cdot)\to a_\infty
  \quad\text{in }L^1(0,T).
\]
Then
\[
  a_\infty\le -\frac{c}{M_1}<0 .
\]
\end{lemma}

\begin{proof}
The upper mass bound and thick weighted drift give
\[
  \frac1T\int_{\tau_n}^{\tau_n+T}a_-(\tau)\,d\tau\ge \frac{c}{M_1}.
\]
After translation and \(L^1\) convergence of \(a\) to the constant
\(a_\infty\), the left-hand side converges to \((-a_\infty)_+\).  Hence
\[
  (-a_\infty)_+\ge \frac{c}{M_1},
\]
and \(a_\infty\le -c/M_1<0\).
\end{proof}

\begin{definition}[Compactness on autonomous windows]
\label{def:s3-compactness-windows}
The branch has compactness on the windows \((T,\tau_n)\) if, for every ball
\(B_R\), the translated sequence
\[
  V_n(y,s)=V(y,\tau_n+s),
  \qquad 0<s<T,
\]
is precompact strongly in \(L^2(0,T;L^q(B_R))\) for every \(q<6\), weakly in
\(L^2(0,T;H^1(B_R))\), and weak-* in
\[
  L^\infty(0,T;L^3(\mathbb R^3)).
\]
The pressures are normalized modulo constants and converge in a topology
sufficient to pass
\[
  -\Delta P=\partial_i\partial_j(V_iV_j)
\]
to the local limit.  The compactness, modulation convergence, and any thick
drift condition are always imposed on the same window sequence \(\tau_n\).
\end{definition}

\begin{theorem}[Autonomous reduced-limit theorem]
\label{thm:s3-autonomous-reduced-limit}
Assume the represented scale-collapse hypotheses above, compactness on
autonomous windows, and an autonomous modulation limit
\((a_\infty,b_\infty)\).  Then, after passing to a subsequence, the translated
branch converges to a nonzero weak solution \((U,\Pi)\) on
\(\mathbb R^3\times(0,T)\) of
\[
  \partial_s U+(U\cdot\nabla)U+\nabla\Pi
  =
  \nu\Delta U+a_\infty(U+y\cdot\nabla U)+b_\infty\cdot\nabla U,
  \qquad
  \nabla\cdot U=0.
\]
Moreover,
\[
  U\in L^\infty(0,T;L^3(\mathbb R^3)),
\]
the pressure satisfies the Navier--Stokes pressure reconstruction, and the
localized core floor passes to the limit on a smaller compact cutoff.  In
particular \(U\not\equiv0\).
\end{theorem}

\begin{proof}
The weak-* \(L^\infty L^3\) bound follows from the positive finite critical
annulus.  The local \(L^2H^1\) bound and time-derivative compactness are part
of the compactness hypothesis, so Aubin--Lions gives strong local convergence
in \(L^2L^q\) for \(q<6\).  This strong convergence, together with weak
\(L^2H^1\) convergence, passes the nonlinear term to
\((U\cdot\nabla)U\) in distributions.  The pressure convergence passes the
Poisson reconstruction to \(\Pi\).

The modulation-limit hypothesis passes the coefficient terms to
\[
  a_\infty(U+y\cdot\nabla U),
  \qquad
  b_\infty\cdot\nabla U .
\]
Therefore the limit solves the autonomous reduced equation.  Finally, the
localized \(L^2\) core floor gives positive localized mass on the selected
branch.  The compactness hypothesis includes convergence strong enough for
this localized pairing, or for a compact subcutoff below it, to pass to the
limit.  Hence \(U\not\equiv0\).
\end{proof}

\begin{definition}[Stationary omega-limit hypothesis]
\label{def:s3-stationary-omega}
For a parameter pair \((a_\infty,b_\infty)\), the autonomous limit of
\cref{thm:s3-autonomous-reduced-limit} has a stationary omega-limit if there is
a nonzero stationary profile \((W,Q)\) satisfying
\[
  (W\cdot\nabla)W+\nabla Q
  =
  \nu\Delta W+a_\infty(W+y\cdot\nabla W)+b_\infty\cdot\nabla W,
  \qquad
  \nabla\cdot W=0,
\]
with
\[
  W\in L^3(\mathbb R^3),
\]
and with the admissibility hypotheses required by the rigidity theorem to be
invoked.
\end{definition}

The stationary omega-limit hypothesis is an additional analytic hypothesis: an
autonomous bounded trajectory need not be stationary merely by time
translation.  A Lyapunov or asymptotic-compactness argument is required to
produce \(W\).

\begin{definition}[Rigidity-covered stationary class]
\label{def:s3-rigidity-covered}
For \((a_\infty,b_\infty)\), the stationary class is rigidity-covered if every
admissible stationary solution of
\[
  (W\cdot\nabla)W+\nabla Q
  =
  \nu\Delta W+a_\infty(W+y\cdot\nabla W)+b_\infty\cdot\nabla W,
  \qquad
  W\in L^3(\mathbb R^3),
\]
is zero, or at least is not a terminal Type II blow-up branch in the class
under consideration.  In the classical backward self-similar normalization,
after rescaling and constant translations, this is the role of the
Nečas--Růžička--Šverák rigidity theorem.  For forward self-similar classes,
the required hypothesis is not nonexistence of profiles, but the narrower statement
that the admissible forward profile is not a terminal Type II blow-up branch.
\end{definition}

\begin{theorem}[Conditional scale-collapse exclusion]
\label{thm:s3-conditional-scale-collapse}
Assume the represented scale-collapse hypotheses, compactness on autonomous
windows, an autonomous modulation limit, a stationary omega-limit, and
rigidity coverage for the corresponding parameter pair.  Then the
scale-collapse alternative is excluded for that branch.
\end{theorem}

\begin{proof}
\Cref{thm:s3-autonomous-reduced-limit} gives a nonzero autonomous limit.  The
stationary omega-limit hypothesis yields a nonzero stationary
profile \((W,Q)\) in the admissible class for
\((a_\infty,b_\infty)\).  Rigidity coverage says that no such nonzero profile
can represent a terminal Type II branch in the class under consideration.  This contradicts the
assumption that the original branch remains in the remaining scale-collapse
class.
\end{proof}

If, in addition, the branch is in the supercritical Type II scale alternative
and the analytic implication from scale-collapse exclusion to that Type II
alternative is assumed, then \cref{thm:s3-conditional-scale-collapse} excludes
the corresponding Type II scale alternative.

The scale-collapse drift obstruction implies
\[
  \int_{\tau_0}^{\infty}a_-(\tau)\,d\tau=\infty
\]
only when \(M_R\) is bounded above on the same moving cutoffs, because then
\[
  a_-(\tau)M_R(\tau)\le M_1 a_-(\tau).
\]
The local mass floor is used for nontriviality of limits, not for this
unweighted-drift implication.  Thus the drift obstruction proves sustained
weighted contraction, but by itself does not prove convergence of \(a(\tau)\)
to a negative constant, nor uniformly negative fixed-window averages.

The scale-collapse alternatives are therefore:
\begin{enumerate}[label=\textup{(\roman*)}]
\item thick autonomous drift, compactness, modulation convergence, and a
stationary omega-limit reduce the branch to a generalized self-similar
stationary profile, which is excluded precisely in parameter ranges covered by
rigidity;
\item thin or oscillatory drift prevents the self-similar rigidity reduction;
\item an autonomous limit with \(a_\infty=0\) gives the stationary
Navier--Stokes equation with possible constant advection, hence a stationary
\(L^3(\mathbb R^3)\) Liouville problem unless an appropriate theorem is
available.
\end{enumerate}
The remaining analytic alternatives are thin drift, compactness failure,
failure of modulation convergence, failure to obtain a stationary omega-limit,
and absence of the relevant rigidity theorem.

\subsection{\texorpdfstring{Local \(L^3\)-Compactness from \(L^2\)-Compactness}{Local L3-Compactness from L2-Compactness}}
\label{subsec:local-l3-compactness}

We record a standard interpolation argument that converts local compactness in
\(L^2\), under uniform local \(L^6\) control, into local compactness in \(L^3\)
for translated renormalized Navier--Stokes trajectories.  Let \(R>0\),
\(I=(0,1)\), \(\overline I=[0,1]\), and let the shifted trajectories be defined
on \(\overline I\times B_{2R}\) by
\[
  W_n(s,y):=V(\tau_n+s,y),
  \qquad (s,y)\in \overline I\times B_{2R},
\]
where \(\tau_n\to\infty\).  All spacetime norms below are taken over \(I\);
the notation \(W_n(0,\cdot)\) denotes the endpoint trace whenever that trace is
used.

\begin{lemma}[Spacetime compactness in \(L^2_tL^2_x\)]
\label{lem:local-l3-compactness-spacetime-l2}
Assume
\[
  \sup_n\|W_n\|_{L^2(I;H^1(B_{2R}))}<\infty
\]
and
\[
  \sup_n\|\partial_s W_n\|_{L^2(I;H^{-1}(B_R))}<\infty .
\]
Then there are a subsequence, not relabeled, and a limit
\[
  W\in L^2(I;L^2(B_R))
\]
such that
\[
  W_n\to W
  \qquad\text{strongly in }L^2(I;L^2(B_R)).
\]
\end{lemma}

\begin{proof}
Restrict \(W_n\) to \(I\times B_R\).  The first hypothesis gives
\((W_n)\) bounded in \(L^2(I;H^1(B_R))\), and the second gives
\((\partial_sW_n)\) bounded in \(L^2(I;H^{-1}(B_R))\).  The embeddings
\[
  H^1(B_R)\Subset L^2(B_R)\hookrightarrow H^{-1}(B_R)
\]
are, respectively, compact and continuous.  The Aubin--Lions--Simon compactness
theorem therefore implies that \((W_n)\) is relatively compact in
\[
  L^2(I;L^2(B_R)).
\]
After passing to a subsequence, there is
\[
  W\in L^2(I;L^2(B_R))
\]
such that \(W_n\to W\) strongly in this space.
\end{proof}

\begin{lemma}[Spacetime compactness in \(L^2_tL^3_x\)]
\label{lem:local-l3-compactness-spacetime-l3}
Assume the hypotheses of Lemma~\ref{lem:local-l3-compactness-spacetime-l2}.  Assume also
\[
  \sup_n\|W_n\|_{L^2(I;L^6(B_{2R}))}<\infty .
\]
Then, after passing to a subsequence if necessary,
\[
  W_n\to W
  \qquad\text{strongly in }L^2(I;L^3(B_R)).
\]
\end{lemma}

\begin{proof}
By the \(L^2(I;L^6(B_{2R}))\) bound, the restricted sequence is bounded in the
reflexive space \(L^2(I;L^6(B_R))\).  Hence, after passing to a subsequence,
\[
  W_n\rightharpoonup \widetilde W
  \qquad\text{weakly in }L^2(I;L^6(B_R)).
\]
The same subsequence converges strongly to \(W\) in \(L^2(I;L^2(B_R))\) by
Lemma~\ref{lem:local-l3-compactness-spacetime-l2}.  To identify the weak
\(L^2_tL^6_x\) limit, let
\[
  \varphi\in C_c^\infty(I\times B_R;\mathbb R^3).
\]
Strong \(L^2_tL^2_x\) convergence gives
\[
  \int_I\int_{B_R} W_n\cdot\varphi\,dy\,ds
  \to
  \int_I\int_{B_R} W\cdot\varphi\,dy\,ds,
\]
while weak convergence in \(L^2_tL^6_x\) gives the same limit with
\(\widetilde W\) in place of \(W\).  Hence \(W=\widetilde W\) in
\(\mathcal D'(I\times B_R)\), and therefore
\[
  W\in L^2(I;L^6(B_R)),
\]
with
\[
  \|W_n-W\|_{L^2(I;L^6(B_R))}
\]
uniformly bounded in \(n\).  Indeed,
\[
  \|W_n-W\|_{L^2(I;L^6(B_R))}
  \le
  \sup_k\|W_k\|_{L^2(I;L^6(B_R))}
  +\|W\|_{L^2(I;L^6(B_R))}.
\]
The Gagliardo--Nirenberg interpolation inequality on the bounded domain \(B_R\)
gives, for almost every \(s\in I\),
\[
  \|W_n(s)-W(s)\|_{L^3(B_R)}
  \le
  C_R\|W_n(s)-W(s)\|_{L^2(B_R)}^{1/2}
  \|W_n(s)-W(s)\|_{L^6(B_R)}^{1/2}
\]
because \(W_n(s)-W(s)\in L^2(B_R)\cap L^6(B_R)\) for a.e. \(s\).  Squaring and
integrating in \(s\) gives
\[
  \int_0^1\|W_n(s)-W(s)\|_{L^3(B_R)}^2\,ds
  \le
  C_R\int_0^1
  \|W_n(s)-W(s)\|_{L^2(B_R)}
  \|W_n(s)-W(s)\|_{L^6(B_R)}\,ds .
\]
Hölder's inequality in time, applied to the product of the
\(L^2(B_R)\)- and \(L^6(B_R)\)-norms, yields
\[
  \|W_n-W\|_{L^2(I;L^3(B_R))}^2
  \le
  C_R
  \|W_n-W\|_{L^2(I;L^2(B_R))}
  \|W_n-W\|_{L^2(I;L^6(B_R))}.
\]
The first factor tends to zero by
Lemma~\ref{lem:local-l3-compactness-spacetime-l2}, while the second factor is uniformly
bounded.  Therefore the left-hand side tends to zero.
\end{proof}

\begin{lemma}[Fixed-time \(L^3\) convergence from fixed-time \(L^2\) convergence]
\label{lem:local-l3-compactness-fixed-time-l3}
Assume that, after passing to a subsequence, the time slices satisfy
\[
  W_n(0,\cdot)\to f
  \qquad\text{strongly in }L^2(B_R)
\]
and
\[
  \sup_n\|W_n(0,\cdot)\|_{L^6(B_R)}<\infty .
\]
Then
\[
  W_n(0,\cdot)\to f
  \qquad\text{strongly in }L^3(B_R).
\]
\end{lemma}

\begin{proof}
First we show that the strong \(L^2\) limit belongs to \(L^6(B_R)\).  Since
\(L^6(B_R)\) is reflexive, the uniform \(L^6(B_R)\) bound gives a subsequence
and a function \(g\in L^6(B_R)\) such that
\[
  W_n(0,\cdot)\rightharpoonup g
  \qquad\text{weakly in }L^6(B_R).
\]
For every \(\phi\in C_c^\infty(B_R)\), the strong \(L^2(B_R)\) convergence
also gives
\[
  \int_{B_R} W_n(0,y)\cdot\phi(y)\,dy
  \to
  \int_{B_R} f(y)\cdot\phi(y)\,dy .
\]
The weak \(L^6\) limit gives the same distributional limit with \(g\) in place
of \(f\), so \(g=f\).  Thus \(f\in L^6(B_R)\).  Returning to the original
sequence, the triangle inequality gives
\[
  \sup_n\|W_n(0,\cdot)-f\|_{L^6(B_R)}
  \le
  \sup_n\|W_n(0,\cdot)\|_{L^6(B_R)}+\|f\|_{L^6(B_R)}<\infty .
\]
Interpolating between \(L^2(B_R)\) and \(L^6(B_R)\) at the fixed time \(s=0\),
\[
  \|W_n(0,\cdot)-f\|_{L^3(B_R)}
  \le
  C_R
  \|W_n(0,\cdot)-f\|_{L^2(B_R)}^{1/2}
  \|W_n(0,\cdot)-f\|_{L^6(B_R)}^{1/2}.
\]
The \(L^6\) factor is uniformly bounded and the \(L^2\) factor tends to zero by
hypothesis.  Hence the \(L^3(B_R)\) norm tends to zero.
\end{proof}

\begin{remark}[Limitation at a fixed time]
\label{rem:local-l3-compactness-fixed-time-limitation}
The preceding two spacetime compactness lemmas do not by themselves imply
convergence at the fixed time \(s=0\).  The fixed-time conclusion in
Lemma~\ref{lem:local-l3-compactness-fixed-time-l3} requires a separate \(L^2\)
compactness hypothesis for the traces.  One sufficient condition is
\[
  \sup_n\|W_n\|_{L^\infty(I;H^1(B_R))}<\infty
\]
and, for some \(q>1\),
\[
  \sup_n\|\partial_s W_n\|_{L^q(I;H^{-1}(B_R))}<\infty .
\]
Simon's compactness theorem then gives precompactness of \((W_n)\) in
\[
  C([0,1];L^2(B_R)),
\]
and hence the fixed-time \(L^2\) convergence after subsequence extraction.
Without such an additional condition, fixed-time convergence at \(s=0\) is not
a consequence of Aubin--Lions compactness in \(L^2(I;L^2(B_R))\).
\end{remark}

\begin{lemma}[Time selection from averaged localized dissipation and local \(H^1\) control]
\label{lem:local-l3-compactness-time-selection}
Let \(J_n=[s_n,s_n+\ell]\) be time intervals with fixed length \(\ell>0\), and
let \(\delta_n\downarrow0\).  Assume that
\(\widetilde{\mathfrak D}_{R_0}\) is a nonnegative measurable function and, for
a fixed radius \(R\), that
\[
  \int_{J_n}\widetilde{\mathfrak D}_{R_0}(\tau)\,d\tau\le\delta_n
\]
and
\[
  \int_{J_n}\|V(\tau)\|_{H^1(B_R)}^2\,d\tau\le C_R .
\]
Then there exist times \(\sigma_n\in J_n\) such that
\[
  \widetilde{\mathfrak D}_{R_0}(\sigma_n)\to0,
  \qquad
  \|V(\sigma_n)\|_{H^1(B_R)}
  \le
  \left(\frac{2C_R}{\ell}\right)^{1/2}.
\]
Consequently, after passing to a subsequence,
\[
  V(\sigma_n)\to V_*
  \qquad\text{strongly in }L^2(B_R)
\]
and
\[
  V(\sigma_n)\to V_*
  \qquad\text{strongly in }L^3(B_R).
\]
\end{lemma}

\begin{proof}
Define
\[
  E_n
  :=
  \left\{\tau\in J_n:
  \widetilde{\mathfrak D}_{R_0}(\tau)\le\delta_n^{1/2}
  \right\}.
\]
By Chebyshev's inequality,
\[
  |J_n\setminus E_n|\le\delta_n^{1/2}.
\]
Indeed,
\[
  \delta_n^{1/2}|J_n\setminus E_n|
  \le
  \int_{J_n\setminus E_n}\widetilde{\mathfrak D}_{R_0}(\tau)\,d\tau
  \le
  \delta_n .
\]
For all sufficiently large \(n\), \(\delta_n^{1/2}\le\ell/2\), and hence
\(|E_n|\ge\ell/2\).  Since
\[
  \int_{E_n}\|V(\tau)\|_{H^1(B_R)}^2\,d\tau\le C_R,
\]
there exists \(\sigma_n\in E_n\) such that
\[
  \|V(\sigma_n)\|_{H^1(B_R)}^2\le \frac{2C_R}{\ell}.
\]
For otherwise the integral over \(E_n\) would be strictly larger than
\((2C_R/\ell)|E_n|\ge C_R\), contradicting the displayed bound.  Because
\(\sigma_n\in E_n\),
\[
  \widetilde{\mathfrak D}_{R_0}(\sigma_n)
  \le
  \delta_n^{1/2}\to0.
\]
The estimate on \(\|V(\sigma_n)\|_{H^1(B_R)}\) gives a bounded sequence in
\(H^1(B_R)\).  By weak compactness in \(H^1(B_R)\), after passing to a
subsequence there is \(V_*\in H^1(B_R)\) such that
\[
  V(\sigma_n)\rightharpoonup V_*
  \qquad\text{weakly in }H^1(B_R).
\]
The Rellich compactness theorem, after passing to a further subsequence if
necessary, gives the corresponding strong convergence
\[
  V(\sigma_n)\to V_*
  \qquad\text{strongly in }L^2(B_R).
\]
The same \(H^1(B_R)\) bound and the Sobolev embedding
\[
  H^1(B_R)\hookrightarrow L^6(B_R)
\]
give a uniform \(L^6(B_R)\) bound for \(V(\sigma_n)\).  Since
\(V_*\in H^1(B_R)\), the same embedding gives \(V_*\in L^6(B_R)\), and
\(\|V(\sigma_n)-V_*\|_{L^6(B_R)}\) is uniformly bounded by the triangle
inequality.  Interpolation gives
\[
  \|V(\sigma_n)-V_*\|_{L^3(B_R)}
  \le
  C_R
  \|V(\sigma_n)-V_*\|_{L^2(B_R)}^{1/2}
  \|V(\sigma_n)-V_*\|_{L^6(B_R)}^{1/2}.
\]
The \(L^6\) factor is uniformly bounded and the \(L^2\) factor tends to zero.
This proves strong convergence in \(L^3(B_R)\).
\end{proof}

\begin{remark}[Role of the time selection]
\label{rem:local-l3-compactness-time-selection-role}
Lemma~\ref{lem:local-l3-compactness-time-selection} does not require endpoint trace compactness.  Instead
of first fixing times \(\tau_n\) and then proving compactness of
\(V(\tau_n)\), one selects times of small localized dissipation inside
intervals where the same intervals also carry local \(H^1\) control.  This
supplies the fixed-time \(L^2\)-compactness hypothesis in
Lemma~\ref{lem:local-l3-compactness-fixed-time-l3}
whenever the argument provides common intervals with these two properties.
\end{remark}

\begin{corollary}[Compactness at translated renormalized times]
\label{cor:local-l3-compactness-renormalized-times}
Assume that, for each fixed radius \(R\), the shifted sequence satisfies the
hypotheses of Lemma~\ref{lem:local-l3-compactness-spacetime-l2} and the time slices
satisfy the two hypotheses of
Lemma~\ref{lem:local-l3-compactness-fixed-time-l3}.  Then, after passing to a
subsequence,
\[
  V(\tau_n)\to V_*
  \qquad\text{strongly in }L^3_{\mathrm{loc}}(\mathbb R^3),
\]
where \(V_*\) is the local \(L^2\) limit of the time slices \(V(\tau_n)\).
\end{corollary}

\begin{proof}
Apply Lemma~\ref{lem:local-l3-compactness-fixed-time-l3} first with \(R=1\) and
pass to a subsequence along which \(V(\tau_n)\) converges strongly in
\(L^3(B_1)\).  Apply the same argument to that subsequence with \(R=2\), and
continue inductively.  The diagonal subsequence converges strongly in
\(L^3(B_m)\) for every integer \(m\ge1\).  The limits on nested balls agree
because the convergence in \(L^3(B_m)\) restricts to convergence in
\(L^3(B_k)\) whenever \(k<m\).  Hence these local limits define a single
function \(V_*\in L^3_{\mathrm{loc}}(\mathbb R^3)\), and
\[
  V(\tau_n)\to V_*
  \qquad\text{strongly in }L^3_{\mathrm{loc}}(\mathbb R^3)
\]
along the diagonal subsequence.
\end{proof}

\begin{remark}[Role in the compactness argument]
\label{rem:local-l3-compactness-role}
The compactness argument uses the preceding statements as follows:
\begin{enumerate}[label=\textup{(\roman*)}]
\item a Caccioppoli-type estimate gives the \(L^2(I;H^1)\) bound in
      Lemma~\ref{lem:local-l3-compactness-spacetime-l2};
\item a time-regularity estimate gives the \(\partial_s W_n\) bound;
\item a local \(L^6\)-type estimate gives the hypothesis of
      Lemma~\ref{lem:local-l3-compactness-spacetime-l3} and fixed-time \(L^6\) boundedness;
\item Lemmas~\ref{lem:local-l3-compactness-spacetime-l2}
      and~\ref{lem:local-l3-compactness-spacetime-l3} give
      spacetime \(L^3\) compactness;
\item an additional fixed-time \(L^2\)-compactness hypothesis gives
      Lemma~\ref{lem:local-l3-compactness-fixed-time-l3};
\item Lemma~\ref{lem:local-l3-compactness-fixed-time-l3} gives the fixed-time strong
      \(L^3\)-compactness used in vanishing-dissipation contradiction arguments.
\end{enumerate}
This is the local compactness step needed to pass from renormalized
\(L^2_{\rm loc}\) compactness to the critical local \(L^3\) topology at the
selected times.
\end{remark}

\begin{corollary}[Local compactness from a common sequence of selected times]
\label{cor:local-l3-compactness-selected-times}
Assume there is a single sequence of selected times \(\sigma_n\) such that, for
every fixed radius \(R\),
\[
  \sup_n\|V(\sigma_n)\|_{H^1(B_R)}<\infty
\]
and
\[
  \widetilde{\mathfrak D}_{R_0}(\sigma_n)\to0 .
\]
Then, after diagonal subsequence extraction,
\[
  V(\sigma_n)\to V_*
  \qquad\text{strongly in }L^3_{\mathrm{loc}}(\mathbb R^3).
\]
For each fixed radius, Lemma~\ref{lem:local-l3-compactness-time-selection} supplies a
radius-dependent sequence from intervals with local \(H^1\) control.  If the
time selection is made on common intervals for an exhaustion of balls, then
Lemma~\ref{lem:local-l3-compactness-time-selection} replaces the separate fixed-time
compactness hypothesis in Lemma~\ref{lem:local-l3-compactness-fixed-time-l3}.
\end{corollary}

\begin{proof}
Fix an integer \(m\ge1\).  The uniform \(H^1(B_m)\) bound gives a subsequence
and a limit \(V_*^{(m)}\in H^1(B_m)\) such that
\[
  V(\sigma_n)\rightharpoonup V_*^{(m)}
  \qquad\text{weakly in }H^1(B_m).
\]
Rellich compactness gives, after passing to a further subsequence,
\[
  V(\sigma_n)\to V_*^{(m)}
  \qquad\text{strongly in }L^2(B_m),
\]
after relabelling the subsequence.  The Sobolev embedding
\(H^1(B_m)\hookrightarrow L^6(B_m)\) gives a uniform \(L^6(B_m)\) bound, and
the limit \(V_*^{(m)}\) also belongs to \(L^6(B_m)\).  Therefore
\[
  \|V(\sigma_n)-V_*^{(m)}\|_{L^6(B_m)}
\]
is uniformly bounded by the triangle inequality.  The interpolation inequality
\[
  \|V(\sigma_n)-V_*^{(m)}\|_{L^3(B_m)}
  \le
  C_m
  \|V(\sigma_n)-V_*^{(m)}\|_{L^2(B_m)}^{1/2}
  \|V(\sigma_n)-V_*^{(m)}\|_{L^6(B_m)}^{1/2}
\]
then gives strong convergence in \(L^3(B_m)\).  Repeating this construction
for \(m=1,2,\ldots\) and taking a diagonal subsequence gives strong convergence
on every \(B_m\).  The nested limits agree by restriction, so they define a
single \(V_*\in L^3_{\mathrm{loc}}(\mathbb R^3)\), and the diagonal subsequence
converges to \(V_*\) strongly in \(L^3_{\mathrm{loc}}(\mathbb R^3)\).
\end{proof}

\subsection{Reduction to Multibubble Alternatives}

Let the structural Type II hypotheses consist of exhaustion of the local Type II
alternatives, repaired-gauge representation, positive finite critical mass,
pressure reconstruction, bounded repaired-gauge modulation, Caccioppoli
regularity, and the cost comparison used for the compact Type II criterion.
If one of these hypotheses fails, the branch is classified by the corresponding
analytic failure rather than by a multibubble alternative.  If, in addition, the
Navier--Stokes implication from compact scale exclusion to the relevant Type II
scale alternative is included, the compact conclusion excludes that alternative.

Under these structural hypotheses, the following non-multibubble strata are
excluded under separate analytic hypotheses.
\begin{enumerate}[label=\textup{(\roman*)}]
\item Far radiation consists of radiative critical mass whose physical radius
stays away from the blow-up point:
\[
  \lambda(\tau_n)R_n\ge r_*>0
\]
along a radiative annular sequence \(R_n\to\infty\).  Under a single-point
blow-up localization hypothesis, this mass lies in a physically regular
exterior region.  An exterior-regularity hypothesis asserts that such
mass can be separated from the concentrating Type II core and does not count
as a core-radiation alternative.
\item The rough-core stratum is excluded by the Caccioppoli implication
\[
\begin{gathered}
  \text{representation},\quad \text{bounded critical norm},\quad
  \text{pressure reconstruction},\\
  \text{bounded modulation},\quad \text{Caccioppoli regularity}
\end{gathered}
\]
which yields local windowed \(H^1\) control.
\item The scale-collapse drift stratum is excluded whenever the scale-collapse
argument extracts a nonzero stationary \(L^3\) profile in a Nečas--Růžička--Šverák
rigidity-covered self-similar class.  The nonzero conclusion uses the same
positive critical-mass and tightness mechanism as the compact Type II
criterion.
\end{enumerate}

\begin{definition}[Same-point secondary concentration]
\label{def:s4-same-point-secondary}
Same-point secondary concentration means that at least two concentration
scales occur at the same physical blow-up point:
\[
  \lambda_1(t)\ll\lambda_2(t)\to0 .
\]
In the primary renormalized coordinates, the secondary bubble appears as
radiation at radius
\[
  R(t)=\frac{\lambda_2(t)}{\lambda_1(t)}\to\infty .
\]
Regular strict cascades of this form are excluded by
\cref{thm:s5-regular-cascade-exclusion}, and all same-point cascades are
excluded under the nonlinear decoupling hypothesis of
\cref{def:s5-same-point-nonlinear-decoupling}.  Before that hypothesis is imposed,
the remaining same-point alternative is precisely the failure of nonlinear
no-splitting/profile compatibility for scale-separated profiles.
\end{definition}

\begin{definition}[Multi-point concentration]
\label{def:s4-multipoint-concentration}
Multi-point concentration means that at least two physical concentration
centers are active.  In a renormalization centered at one bubble, the other
bubble appears as an escaping profile.  This case is excluded by a
separated-point decoupling hypothesis reducing each active physical point to
its own single-center branch.  Before that hypothesis is imposed, the remaining
multi-point alternative is failure of separated-point decoupling.
\end{definition}

\begin{theorem}[Reduction to multibubble alternatives]
\label{thm:s4-reduction-multibubble-residue}
Assume an admissible Navier--Stokes Type II branch satisfies the structural
Type II hypotheses above, the single-point blow-up localization and
exterior-regular removal hypotheses for far radiation, and the scale-collapse
rigidity hypothesis for scale-collapse drift.  If the branch is not a multibubble branch
in either the same-point scale-separated form of
\cref{def:s4-same-point-secondary} or the multi-point form of
\cref{def:s4-multipoint-concentration}, then the compact Type II
scale-exclusion conclusion holds.  If the Navier--Stokes implication from this
scale exclusion to the relevant Type II scale alternative is also included,
then that Type II scale alternative is excluded.
\end{theorem}

\begin{proof}
The structural hypotheses supply representation, critical-mass control,
pressure reconstruction, bounded modulation, Caccioppoli regularity, and cost
comparison.  If the branch is globally \(L^3\)-tight, the Caccioppoli estimate
gives local windowed \(H^1\) control.  Together with positive finite critical
mass and tightness, the compact Type II criterion gives infinite localized
renormalization cost, and the cost comparison gives the Type II scale-exclusion
condition.

If uniform compact-cylinder tightness fails, the radiative branch is present.
Far radiation is excluded by the single-point blow-up localization and
exterior-regular removal hypotheses, because it lies in a physically smooth
exterior region and is separated from the concentrating Type II core.  Regular
strict same-point cascades are excluded by
\cref{thm:s5-regular-cascade-exclusion}, and all same-point cascades are
excluded under the nonlinear same-point decoupling hypothesis.  Separated-point
multibubbles are excluded under the separated-point decoupling hypothesis.  If
the branch is assumed not to be in either multibubble form, no
radiative or noncompact alternative remains.

If scale-collapse drift occurs, the scale-collapse rigidity hypothesis extracts a nonzero
\(L^3\) stationary profile in a rigidity-covered self-similar class.  Rigidity
forces this profile to vanish, contradicting nonzero extraction.  The
scale-collapse drift alternative is therefore excluded.

The remaining rough-core branch is excluded by the Caccioppoli implication,
because the structural hypotheses include the pressure, modulation, and
regularity hypotheses required for local windowed \(H^1\) control.  Thus all
non-multibubble alternatives are exhausted, and the compact Type II
scale-exclusion conclusion holds.  The final Type II scale conclusion follows
from the corresponding Navier--Stokes implication when that hypothesis is
included.
\end{proof}

Thus, after the structural hypotheses, same-point regular cascade estimates, and
separated-point reductions are supplied, the only remaining multibubble
alternatives are failure of same-point nonlinear decoupling and failure of
separated-point decoupling.  Once nonlinear profile evolution and terminal
active-frame compactness provide those decouplings, a multibubble failure is an
earlier failure of profile completeness, representation, Caccioppoli
regularity, or scale-collapse rigidity, not an additional singularity class.

\subsection{Same-Point Multiscale Cascade Exclusion}

For \(\lambda>0\), \(x_0\in\mathbb R^3\), and
\(\phi\in L^3(\mathbb R^3)\), define
\[
  (\Lambda_{\lambda,x_0}\phi)(x)
  =
  \lambda^{-1}\phi\left(\frac{x-x_0}{\lambda}\right).
\]
Then
\[
  \|\Lambda_{\lambda,x_0}\phi\|_{L^3(\mathbb R^3)}
  =
  \|\phi\|_{L^3(\mathbb R^3)} .
\]
Two parameter sequences \((\lambda_i^n,x_i^n)\) and
\((\lambda_j^n,x_j^n)\) are asymptotically orthogonal if
\[
  \frac{\lambda_i^n}{\lambda_j^n}
  +
  \frac{\lambda_j^n}{\lambda_i^n}
  +
  \frac{|x_i^n-x_j^n|}{\max(\lambda_i^n,\lambda_j^n)}
  \longrightarrow\infty .
\]

\begin{lemma}[Finite \(L^3\) decoupling]
\label{lem:s5-finite-l3-decoupling}
Let \(K<\infty\), let \(\phi^1,\ldots,\phi^K\in L^3(\mathbb R^3)\), and set
\[
  g_n^j=\Lambda_{\lambda_j^n,x_j^n}\phi^j .
\]
If the parameter sequences are pairwise asymptotically orthogonal, then
\[
  \left\|\sum_{j=1}^K g_n^j\right\|_{L^3}^3
  =
  \sum_{j=1}^K\|\phi^j\|_{L^3}^3+o_n(1).
\]
\end{lemma}

\begin{proof}
It suffices first to prove the statement for
\(\phi^j\in C_c^\infty(\mathbb R^3)\).  The general \(L^3\) case follows by
density.  For each \(\varepsilon>0\), choose
\(\psi^j\in C_c^\infty\) with
\[
  \|\phi^j-\psi^j\|_3<\varepsilon .
\]
Since each \(\Lambda_{\lambda_j^n,x_j^n}\) is an \(L^3\) isometry,
\[
  \left\|
  \sum_{j=1}^K\Lambda_{\lambda_j^n,x_j^n}(\phi^j-\psi^j)
  \right\|_3
  \le K\varepsilon .
\]
On bounded \(L^3\) sets, \(f\mapsto\|f\|_3^3\) is locally Lipschitz:
\[
  \bigl|\|f\|_3^3-\|h\|_3^3\bigr|
  \le
  C(\|f\|_3+\|h\|_3)^2\|f-h\|_3 .
\]
Thus the smooth compactly supported case implies the general case after
\(\varepsilon\downarrow0\).

Assume now \(\phi^j\in C_c^\infty\).  The proof is by induction on \(K\).  The
case \(K=1\) is the scaling invariance of the \(L^3\) norm.  Assume the result
for \(K-1\), and write
\[
  G_n=\sum_{j=1}^{K-1}\Lambda_{\lambda_j^n,x_j^n}\phi^j,
  \qquad
  g_n^K=\Lambda_{\lambda_K^n,x_K^n}\phi^K .
\]
Applying the inverse scaling of the \(K\)-th profile gives
\[
  (\Lambda_{\lambda_K^n,x_K^n}^{-1}G_n)(y)
  =
  \sum_{j=1}^{K-1}
  \frac{\lambda_K^n}{\lambda_j^n}
  \phi^j\left(
  \frac{\lambda_K^n y+x_K^n-x_j^n}{\lambda_j^n}
  \right).
\]
For each \(j<K\), asymptotic orthogonality gives, along a subsequence, one of
the following alternatives:
\[
  \lambda_K^n/\lambda_j^n\to0,\qquad
  \lambda_j^n/\lambda_K^n\to0,
\]
or the scale ratio stays bounded above and below while
\[
  |x_K^n-x_j^n|/\lambda_j^n\to\infty .
\]
In the first case the displayed term is bounded pointwise by
\[
  \frac{\lambda_K^n}{\lambda_j^n}\|\phi^j\|_\infty\to0 .
\]
In the second case, if \(\operatorname{supp}\phi^j\subset B_A\), the set on
which the term is nonzero is
\[
  E_n^j
  =
  \left\{
  y:
  \left|
  \frac{\lambda_K^n y+x_K^n-x_j^n}{\lambda_j^n}
  \right|\le A
  \right\}.
\]
Since \(\lambda_j^n/\lambda_K^n\to0\), the diameter of \(E_n^j\) is
\[
  O(\lambda_j^n/\lambda_K^n)\to0 .
\]
Thus \(\mathbf 1_{E_n^j}\to0\) pointwise outside at most one limiting point
after passing to a subsequence, and the term converges to zero a.e.  In the
third case, the spatial translation tends to infinity in the rescaled
variables; hence for each fixed compact \(B_R\), the argument of \(\phi^j\)
lies outside \(\operatorname{supp}\phi^j\) for all large \(n\).  Therefore
\[
  \Lambda_{\lambda_K^n,x_K^n}^{-1}G_n\to0
  \quad\text{a.e. on }\mathbb R^3.
\]
The sequence is bounded in \(L^3\), because the inverse scaling is an
\(L^3\) isometry and the sum is finite.  The Brézis--Lieb lemma, applied to
\[
  \phi^K+\Lambda_{\lambda_K^n,x_K^n}^{-1}G_n,
\]
gives
\[
  \left\|
  \phi^K+\Lambda_{\lambda_K^n,x_K^n}^{-1}G_n
  \right\|_3^3
  =
  \|\phi^K\|_3^3
  +
  \left\|
  \Lambda_{\lambda_K^n,x_K^n}^{-1}G_n
  \right\|_3^3
  +o_n(1).
\]
Using the \(L^3\) isometry again,
\[
  \|G_n+g_n^K\|_3^3
  =
  \|\phi^K\|_3^3+\|G_n\|_3^3+o_n(1).
\]
The induction hypothesis gives
\[
  \|G_n\|_3^3
  =
  \sum_{j=1}^{K-1}\|\phi^j\|_3^3+o_n(1),
\]
and the result follows for \(K\).
\end{proof}

\begin{definition}[Admissible same-point critical cascade]
\label{def:s5-admissible-same-point-cascade}
Let \(u_n=u(\cdot,t_n)\), \(t_n\uparrow T^*\).  A same-point critical cascade
decomposition at \(x_*\) consists of profiles
\(\{\phi^j\}_{j\in\mathcal J}\subset L^3(\mathbb R^3)\), scales
\(\lambda_j^n\downarrow0\), centers \(x_j^n\to x_*\), and remainders \(r_n^J\)
such that, for every finite \(J\subset\mathcal J\),
\[
  u_n
  =
  \sum_{j\in J}\Lambda_{\lambda_j^n,x_j^n}\phi^j+r_n^J .
\]
The decomposition is admissible if:
\begin{enumerate}[label=\textup{(\roman*)}]
\item the parameters are pairwise asymptotically orthogonal:
\[
  \frac{\lambda_i^n}{\lambda_j^n}
  +
  \frac{\lambda_j^n}{\lambda_i^n}
  +
  \frac{|x_i^n-x_j^n|}{\max(\lambda_i^n,\lambda_j^n)}
  \to\infty
  \qquad (i\ne j);
\]
\item critical \(L^3\) decoupling with remainder holds:
\[
  \|u_n\|_3^3
  =
  \sum_{j\in J}\|\phi^j\|_3^3
  +
  \|r_n^J\|_3^3
  +o_n(1);
\]
\item there is \(\eta>0\) such that
\[
  \|\phi^j\|_3\ge\eta
  \qquad\text{for every }j\in\mathcal J .
\]
\end{enumerate}
\end{definition}

\begin{theorem}[Uniform bound on cascade length]
\label{thm:s5-uniform-cascade-length}
Assume
\[
  \sup_n\|u_n\|_{L^3(\mathbb R^3)}\le M<\infty .
\]
If \(u_n\) admits an admissible same-point critical cascade decomposition and
every active profile satisfies
\[
  \|\phi^j\|_3\ge\eta>0,
\]
then the number of active profiles is finite and
\[
  \#\mathcal J
  \le
  \left\lfloor\frac{M^3}{\eta^3}\right\rfloor .
\]
\end{theorem}

\begin{proof}
Let \(J\subset\mathcal J\) be finite.  By \(L^3\) decoupling with remainder,
\[
  \|u_n\|_3^3
  =
  \sum_{j\in J}\|\phi^j\|_3^3+\|r_n^J\|_3^3+o_n(1).
\]
The remainder contribution is nonnegative, hence
\[
  \limsup_{n\to\infty}\|u_n\|_3^3
  \ge
  \sum_{j\in J}\|\phi^j\|_3^3 .
\]
Since \(\|u_n\|_3\le M\),
\[
  M^3
  \ge
  \sum_{j\in J}\|\phi^j\|_3^3
  \ge
  (\#J)\eta^3 .
\]
Thus
\[
  \#J\le \frac{M^3}{\eta^3}
\]
for every finite \(J\subset\mathcal J\).  If \(\mathcal J\) had more than
\(\lfloor M^3/\eta^3\rfloor\) elements, choosing
\[
  \#J=\left\lfloor\frac{M^3}{\eta^3}\right\rfloor+1
\]
would contradict the preceding inequality.
\end{proof}

\begin{definition}[Regular strict same-point cascade]
\label{def:s5-regular-same-point-cascade}
Assume there are \(J\ge2\) active profiles ordered by scale:
\[
  \lambda_1^n\ll\lambda_2^n\ll\cdots\ll\lambda_J^n,
  \qquad
  \lambda_j^n\to0,
  \qquad
  x_j^n\to x_* .
\]
Set
\[
  \mu_j^n=\frac{\lambda_1^n}{\lambda_j^n}\to0
  \qquad (j\ge2).
\]
The cascade is regular in the innermost coordinates if, after passing to a
subsequence:
\begin{enumerate}[label=\textup{(\roman*)}]
\item
\[
  \xi_j^n=\frac{x_*-x_j^n}{\lambda_j^n}
  \longrightarrow \xi_j\in\mathbb R^3
  \qquad (j\ge2);
\]
\item each outer profile is \(C^2\) near \(\xi_j\);
\item for every \(R<\infty\),
\[
  \sup_{|z-\xi_j|\le R\mu_j^n}
  \bigl(|\phi^j(z)|+|\nabla\phi^j(z)|+|\nabla^2\phi^j(z)|\bigr)<\infty
\]
uniformly for all large \(n\);
\item the innermost repaired-gauge representation is available:
\[
  V_1^n(y,\tau)
  =
  \lambda_1(t)u(x_*(t)+\lambda_1(t)y,t)
\]
and solves the repaired-gauge Navier--Stokes equation on compact
\((y,\tau)\)-cylinders, modulo the outer-bubble perturbation.
\end{enumerate}
If \(|\xi_j^n|\to\infty\), the corresponding outer bubble is locally
invisible in the innermost variables by \cref{lem:s5-escaping-offset}, provided
the pressure interaction satisfies the decoupling hypothesis stated below.
\end{definition}

\begin{lemma}[Inner expansion of regular outer bubbles]
\label{lem:s5-inner-expansion}
For \(j\ge2\), define
\[
  W_j^n(y)
  =
  \lambda_1^n(\lambda_j^n)^{-1}
  \phi^j\left(\frac{\lambda_1^n y+x_*-x_j^n}{\lambda_j^n}\right)
  =
  \mu_j^n\phi^j(\xi_j^n+\mu_j^n y).
\]
For every fixed \(R<\infty\),
\[
  W_j^n(y)
  =
  \mu_j^n\phi^j(\xi_j)
  +
  (\mu_j^n)^2D\phi^j(\xi_j)y
  +
  \mathcal R_j^n(y)
  \qquad\text{on }B_R,
\]
where
\[
  \|\mathcal R_j^n\|_{L^\infty(B_R)}
  \le
  C_{j,R}\left((\mu_j^n)^3+\mu_j^n|\xi_j^n-\xi_j|^2
  +(\mu_j^n)^2|\xi_j^n-\xi_j|\right)
  =
  o(\mu_j^n)+O((\mu_j^n)^3).
\]
In particular,
\[
  W_j^n=c_j^n+A_j^n y+o_{L^\infty(B_R)}(\mu_j^n),
\]
with
\[
  c_j^n=\mu_j^n\phi^j(\xi_j),
  \qquad
  A_j^n=(\mu_j^n)^2D\phi^j(\xi_j),
  \qquad
  \|A_j^n\|=O((\mu_j^n)^2).
\]
\end{lemma}

\begin{proof}
Fix \(R\).  For \(y\in B_R\), Taylor's theorem at \(\xi_j\) gives
\[
\begin{aligned}
  \phi^j(\xi_j^n+\mu_j^n y)
  &=
  \phi^j(\xi_j)
  +D\phi^j(\xi_j)(\xi_j^n-\xi_j+\mu_j^n y)\\
  &\quad
  +\frac12(\xi_j^n-\xi_j+\mu_j^n y)^T
  D^2\phi^j(\zeta_{j,n,y})
  (\xi_j^n-\xi_j+\mu_j^n y),
\end{aligned}
\]
where \(\zeta_{j,n,y}\) lies on the segment between \(\xi_j\) and
\(\xi_j^n+\mu_j^n y\).  Multiplying by \(\mu_j^n\) yields
\[
\begin{aligned}
  W_j^n(y)
  &=
  \mu_j^n\phi^j(\xi_j)
  +\mu_j^nD\phi^j(\xi_j)(\xi_j^n-\xi_j)
  +(\mu_j^n)^2D\phi^j(\xi_j)y\\
  &\quad
  +O\left(\mu_j^n|\xi_j^n-\xi_j+\mu_j^n y|^2\right).
\end{aligned}
\]
The term \(\mu_j^nD\phi^j(\xi_j)(\xi_j^n-\xi_j)\) is independent of \(y\) and
may be absorbed into the constant part if one uses
\[
  c_j^n=\mu_j^n\phi^j(\xi_j^n).
\]
With the displayed choice \(c_j^n=\mu_j^n\phi^j(\xi_j)\), it is part of the
remainder.  Since \(\xi_j^n\to\xi_j\) and \(\mu_j^n\to0\), the stated estimate
follows.
\end{proof}

For the dynamical reduction it is useful to use the exact constant
\[
  c_j^n=\mu_j^n\phi^j(\xi_j^n),
\]
so that
\[
  W_j^n(y)
  =
  c_j^n
  +(\mu_j^n)^2D\phi^j(\xi_j^n)y
  +O_{L^\infty(B_R)}((\mu_j^n)^3R^2).
\]
Since \cref{thm:s5-uniform-cascade-length} gives \(J<\infty\),
\[
  c_{\rm out}^n=\sum_{j=2}^J c_j^n
  =
  O\left(\sum_{j=2}^J\mu_j^n\right)\to0
\]
on every fixed cascade sequence, provided the profiles are locally bounded at
their offsets.

\begin{lemma}[Absorption of constant ambient velocity]
\label{lem:s5-constant-ambient-absorption}
Suppose on a compact cylinder
\[
  V=V_{\rm in}+c(\tau)+S,
\]
where \(c(\tau)\) is independent of \(y\) and \(S\) is a perturbation.  Ignoring
\(S\), the equation for \(V_{\rm in}\) is
\[
  \partial_\tau V_{\rm in}
  +(V_{\rm in}\cdot\nabla)V_{\rm in}
  +\nabla P_{\rm in}
  =
  \nu\Delta V_{\rm in}
  +a(V_{\rm in}+y\cdot\nabla V_{\rm in})
  +(b-c)\cdot\nabla V_{\rm in},
\]
modulo spatially constant forcing terms, which are absorbed by the moving
translation gauge.  Thus \(b_{\rm eff}=b-c\).
\end{lemma}

\begin{proof}
Since \(c\) is independent of \(y\),
\[
  (V_{\rm in}+c)\cdot\nabla(V_{\rm in}+c)
  =
  (V_{\rm in}\cdot\nabla)V_{\rm in}
  +c\cdot\nabla V_{\rm in}.
\]
Terms depending only on \(\tau\), such as \(\partial_\tau c-ac\), are spatially
constant forces and are absorbed in the equivalent translation-gauge
description.  With the stated sign convention, the effective translation
parameter is \(b-c\).
\end{proof}

The leading constant part of every outer bubble in the inner coordinates is
absorbed into \(b_{\rm eff}\).  Since
\[
  |c_{\rm out}^n|
  \le
  \sum_{j=2}^J\mu_j^n|\phi^j(\xi_j^n)|,
\]
finite cascade length and local boundedness of the profiles imply
\[
  \sup_n|c_{\rm out}^n|<\infty,
  \qquad
  c_{\rm out}^n\to0 .
\]
Thus bounded modulation is preserved whenever the unperturbed repaired-gauge
\(b\)-parameter is bounded.

\begin{definition}[Perturbative scale-collapse robustness]
\label{def:s5-perturbative-s3}
Let \(E_n\) be the remaining inner-coordinate error after subtracting the
innermost bubble and absorbing \(c_{\rm out}^n\) into \(b_{\rm eff}\).  It
contains the linear shear terms
\[
  \sum_{j=2}^J(\mu_j^n)^2D\phi^j(\xi_j^n)y,
\]
Taylor remainders of size \(O((\mu_j^n)^3)\) on fixed balls, the rescaled
critical remainder, and pressure cross-terms generated by the interaction of
the innermost profile with the outer perturbation.  Perturbative scale-collapse robustness
means that on every compact cylinder \(B_R\times[\tau_0,\tau_0+L]\):
\begin{enumerate}[label=\textup{(\roman*)}]
\item \(E_n\to0\) in \(L^1_\tau H^{-1}_y(B_R)\), or in a stronger topology
sufficient for the scale-collapse compactness passage;
\item the pressure generated by cross-terms converges to zero in the pressure
topology used by the repaired-gauge and scale-collapse arguments;
\item local compactness, autonomous modulation, and stationary-limit
extraction in the scale-collapse limit are stable under adding \(E_n\to0\);
\item the limiting stationary equation is the
Nečas--Růžička--Šverák-covered self-similar profile equation, not a perturbed
equation.
\end{enumerate}
\end{definition}

Under \cref{def:s5-perturbative-s3}, the innermost bubble gives a
single-bubble scale-collapse branch with bounded modulation and positive
critical mass.

\begin{theorem}[Regular same-point cascade exclusion]
\label{thm:s5-regular-cascade-exclusion}
Let \(u\) be a Leray--Hopf solution on \(\mathbb R^3\times[0,T^*)\) with finite
initial energy and
\[
  \sup_{t<T^*}\|u(t)\|_{L^3(\mathbb R^3)}\le M<\infty .
\]
Let \(t_n\uparrow T^*\) be a concentrating sequence at \(x_*\).  Assume:
\begin{enumerate}[label=\textup{(\roman*)}]
\item \(u_n=u(\cdot,t_n)\) admits an admissible same-point critical cascade
decomposition;
\item every active bubble satisfies \(\|\phi^j\|_3\ge\eta>0\);
\item the active bubbles form a regular strict same-point cascade;
\item the innermost scale admits the repaired-gauge representation with bounded
modulation, pressure reconstruction, Caccioppoli regularity, and positive
finite critical mass;
\item perturbative scale-collapse robustness holds;
\item the scale-collapse extraction of a nonzero limiting profile and
Nečas--Růžička--Šverák rigidity hold for the innermost branch.
\end{enumerate}
Then the cascade cannot occur.
\end{theorem}

\begin{proof}
By \cref{thm:s5-uniform-cascade-length}, the number of active bubbles is finite:
\[
  J\le \left\lfloor\frac{M^3}{\eta^3}\right\rfloor .
\]
Let \(\lambda_1^n\) be the smallest scale, and for \(j\ge2\) set
\[
  \mu_j^n=\lambda_1^n/\lambda_j^n\to0 .
\]
\Cref{lem:s5-inner-expansion} gives, on every fixed \(B_R\),
\[
  \lambda_1^n(\lambda_j^n)^{-1}
  \phi^j\left(\frac{\lambda_1^n y+x_*-x_j^n}{\lambda_j^n}\right)
  =
  c_j^n+(\mu_j^n)^2D\phi^j(\xi_j^n)y
  +O_{L^\infty(B_R)}((\mu_j^n)^3R^2).
\]
Summing over \(2\le j\le J\), the constant part
\[
  c_{\rm out}^n=\sum_{j=2}^J c_j^n
\]
is bounded and tends to zero.  By \cref{lem:s5-constant-ambient-absorption}, it
is absorbed into the repaired-gauge translation parameter:
\[
  b_{\rm eff}^n=b^n-c_{\rm out}^n .
\]
The bounded-modulation hypothesis for \(b^n\), together with boundedness of
\(c_{\rm out}^n\), gives
\[
  \sup_n|b_{\rm eff}^n|<\infty .
\]
The nonconstant outer contribution is
\[
  \sum_{j=2}^J(\mu_j^n)^2D\phi^j(\xi_j^n)y
  +
  O_{L^\infty(B_R)}
  \left(\sum_{j=2}^J(\mu_j^n)^3R^2\right),
\]
which converges to zero uniformly on each fixed \(B_R\).  The rescaled
profile-decomposition remainder and pressure cross-terms are included in
\(E_n\), and perturbative scale-collapse robustness gives \(E_n\to0\) in the
topology required by the scale-collapse compactness and rigidity passage.

Consequently, the innermost branch satisfies the repaired-gauge
Navier--Stokes equation with bounded effective modulation, positive finite
critical mass, pressure reconstruction, Caccioppoli regularity, and a
perturbation that vanishes in the scale-collapse limit.  The scale
\(\lambda_1(t)\to0\) gives the scale-collapse hypothesis.  The scale-collapse
rigidity hypotheses therefore produce a nonzero stationary
\(L^3(\mathbb R^3)\) profile
in the Nečas--Růžička--Šverák-covered self-similar class.

The Nečas--Růžička--Šverák rigidity theorem rules out every nonzero
\(L^3(\mathbb R^3)\) stationary profile in that class.  This contradicts the
nonzero extraction, which ultimately comes from \(\|\phi^1\|_3\ge\eta>0\).
Hence the regular strict same-point cascade cannot occur.
\end{proof}

The preceding theorem uses the following analytic hypotheses: strong \(L^3\)
profile decomposition with Brézis--Lieb decoupling and remainder decoupling;
a uniform active-bubble mass floor \(\|\phi^j\|_3\ge\eta\); finite-offset
\(C^2\) nested-cascade regularity; repaired-gauge representation of the
innermost branch with pressure, Caccioppoli regularity, bounded modulation, and
positive finite critical mass; perturbative scale-collapse robustness; and the
scale-collapse extraction of a nonzero limiting profile together with
Nečas--Růžička--Šverák rigidity.  Bounded critical mass alone proves only uniformly bounded cascade
length.  Full exclusion of the regular same-point cascade requires the
inner-decoupling and perturbative scale-collapse hypotheses.

\begin{lemma}[Escaping offsets vanish locally]
\label{lem:s5-escaping-offset}
Let
\[
  W_j^n(y)=\mu_j^n\phi^j(\xi_j^n+\mu_j^n y),
  \qquad
  \mu_j^n=\frac{\lambda_1^n}{\lambda_j^n}\to0.
\]
If \(\phi^j\in L^3(\mathbb R^3)\) and \(|\xi_j^n|\to\infty\), then, for every
fixed \(R<\infty\),
\[
  \|W_j^n\|_{L^3(B_R)}\to0.
\]
If additionally
\[
  \|\nabla W_j^n\|_{L^3(B_R)}
  +
  \|D^2W_j^n\|_{L^3(B_R)}
  \to0
\]
for every fixed \(R\), then \(W_j^n\to0\) in \(W^{2,3}_{\rm loc}\).
\end{lemma}

\begin{proof}
With \(z=\xi_j^n+\mu_j^n y\) and \(dy=(\mu_j^n)^{-3}\,dz\),
\[
  \|W_j^n\|_{L^3(B_R)}^3
  =
  \int_{B_R}(\mu_j^n)^3
  |\phi^j(\xi_j^n+\mu_j^n y)|^3\,dy
  =
  \int_{B(\xi_j^n,\mu_j^n R)}|\phi^j(z)|^3\,dz .
\]
Since \(|\xi_j^n|\to\infty\) and \(\mu_j^nR\to0\), the balls
\[
  B(\xi_j^n,\mu_j^nR)
\]
eventually lie in \(\{|z|>A\}\) for every fixed \(A<\infty\).  Since
\(\phi^j\in L^3(\mathbb R^3)\),
\[
  \int_{|z|>A}|\phi^j(z)|^3\,dz\to0
  \quad\text{as }A\to\infty .
\]
Thus, for all sufficiently large \(n\),
\[
  \|W_j^n\|_{L^3(B_R)}^3
  \le
  \int_{|z|>A}|\phi^j(z)|^3\,dz,
\]
and the right-hand side is arbitrarily small for large \(A\).  This proves
local \(L^3\) vanishing.  The \(W^{2,3}_{\rm loc}\) conclusion is precisely the
additional derivative-tail hypothesis.
\end{proof}

The pressure contribution of an escaping-offset bubble is not automatic from
local velocity vanishing, because pressure is nonlocal.  The needed
escaping-offset pressure-decoupling hypothesis is that, for every compact
\(B_R\) and compact time window, the pressure generated by the cross source
\[
  U_n\odot W_j^n+W_j^n\odot U_n+W_j^n\otimes W_j^n
\]
converges to zero modulo constants in the pressure topology used for the
repaired-gauge and scale-collapse passages.  Under this hypothesis,
escaping-offset outer profiles are perturbative for the innermost
scale-collapse branch.

\begin{proposition}[Perturbative estimates for the same-point cascade]
\label{prop:s5-perturbative-estimates}
Let \(U_n\) be the innermost renormalized branch after absorbing
\(c_{\rm out}^n\) into \(b^n\), and let \(S_n\) be the remaining outer
perturbation on \(B_R\times I\).  Assume
\[
  \|S_n\|_{L^\infty_\tau W^{2,\infty}_y(B_R)}
  +
  \|\partial_\tau S_n\|_{L^1_\tau H^{-1}_y(B_R)}
  \longrightarrow0,
\]
\[
  \|U_n\|_{L^\infty_\tau L^3_y(B_{2R})}
  +
  \|U_n\|_{L^2_\tau H^1_y(B_{2R})}
  \le C_R,
\]
and assume the pressure reconstruction satisfies the local
Calderon--Zygmund estimate for cross terms.  Then every perturbative term
produced by replacing \(U_n+S_n\) with \(U_n\) vanishes in
\[
  L^1_\tau H^{-1}_y(B_R).
\]
\end{proposition}

\begin{proof}
The time and Laplacian perturbations satisfy
\[
  \|\partial_\tau S_n\|_{L^1H^{-1}}\to0,
  \qquad
  \|\Delta S_n\|_{L^1H^{-1}}
  \le C_R\|S_n\|_{L^1W^{2,\infty}}\to0.
\]
The linear transport terms obey
\[
  \|(U_n\cdot\nabla)S_n\|_{L^1H^{-1}(B_R)}
  \le
  C_R\|U_n\|_{L^2L^2(B_R)}
  \|\nabla S_n\|_{L^2L^\infty(B_R)}
  \to0,
\]
and
\[
  \|(S_n\cdot\nabla)U_n\|_{L^1H^{-1}(B_R)}
  \le
  C_R\|S_n\|_{L^\infty L^\infty(B_R)}
  \|\nabla U_n\|_{L^1L^2(B_R)}
  \to0.
\]
The quadratic perturbation satisfies
\[
  \|(S_n\cdot\nabla)S_n\|_{L^1H^{-1}(B_R)}
  \le
  C_R\|S_n\|_{L^\infty L^\infty}
  \|\nabla S_n\|_{L^1L^\infty}
  \to0.
\]
The modulation terms satisfy
\[
  \|a_n(S_n+y\cdot\nabla S_n)+b_n\cdot\nabla S_n\|_{L^1H^{-1}(B_R)}
  \le
  C_R(M_{ab})\|S_n\|_{L^1W^{1,\infty}(B_R)}
  \to0.
\]
For the pressure, the cross source is
\[
  2U_n\odot S_n+S_n\otimes S_n.
\]
On \(B_{2R}\),
\[
  \|U_n\odot S_n\|_{L^1_\tau L^{3/2}_y}
  \le
  \|S_n\|_{L^\infty_{\tau,y}}
  \|U_n\|_{L^1_\tau L^3_y}
  \to0,
\]
and
\[
  \|S_n\otimes S_n\|_{L^1_\tau L^{3/2}_y}\to0.
\]
The local pressure estimate then gives convergence of the cross pressure
modulo constants in the required pressure topology, hence its gradient
vanishes in \(L^1_\tau H^{-1}_y(B_R)\).
\end{proof}

For a regular finite strict cascade, \cref{lem:s5-inner-expansion} and bounded
relative modulation give the required \(S_n\to0\) estimates.  Therefore
regular nested-cascade structure, bounded modulation, pressure reconstruction,
and local windowed \(H^1\) control imply perturbative scale-collapse
robustness.

\begin{definition}[Same-point nonlinear decoupling]
\label{def:s5-same-point-nonlinear-decoupling}
Same-point nonlinear decoupling means that every finite same-point strict
scale-separated profile family satisfies:
\begin{enumerate}[label=\textup{(\roman*)}]
\item strong critical decomposition,
\[
  u_n
  =
  \sum_{j=1}^J\Lambda_{\lambda_j^n,x_j^n}\phi^j+r_n,
  \qquad
  \|u_n\|_3^3
  =
  \sum_{j=1}^J\|\phi^j\|_3^3+\|r_n\|_3^3+o_n(1);
\]
\item nonlinear profile stability: profiles below the critical small-data
threshold are perturbative, and every active singular profile has \(L^3\)-mass
at least \(\varepsilon_{\rm sd}\);
\item innermost velocity decoupling: after rescaling by the smallest scale and
absorbing all constant ambient velocities into the translation gauge,
\[
  V_n=U_n+S_n,
  \qquad
  S_n\to0
\]
in the local velocity topology required by perturbative scale-collapse
robustness;
\item pressure decoupling: all pressure cross-terms between the innermost
profile, outer profiles, and the remainder vanish modulo constants in the
repaired-gauge scale-collapse pressure topology;
\item modulation compatibility: the effective repaired-gauge parameters of the
innermost branch remain bounded after ambient constants are absorbed;
\item scale-collapse topology compatibility: the full interaction error
vanishes in the compactness and stationarity passage.
\end{enumerate}
\end{definition}

\begin{theorem}[Conditional no same-point multibubble theorem]
\label{thm:s5-no-same-point-multibubble}
Assume an admissible Navier--Stokes Type II branch has the structural hypotheses
needed to enter the repaired-gauge Type II class, and assume same-point
nonlinear decoupling, the innermost repaired-gauge representation, and the
scale-collapse extraction of a nonzero limiting profile together with
Nečas--Růžička--Šverák rigidity.  Then no
same-point multibubble cascade of active singular profiles can occur.
\end{theorem}

\begin{proof}
Let a same-point strict scale-separated cascade be given.  By
\cref{def:s5-same-point-nonlinear-decoupling}, the decomposition is strongly
decoupled in critical \(L^3\), and every active singular bubble carries at
least \(\varepsilon_{\rm sd}\) critical mass.  Since the ambient Type II branch
has bounded critical mass,
\[
  \sup_n\|u_n\|_3\le M,
\]
\cref{thm:s5-uniform-cascade-length} gives
\[
  J\le \left\lfloor\frac{M^3}{\varepsilon_{\rm sd}^3}\right\rfloor .
\]
Choose the innermost scale \(\lambda_1^n\).  The nonlinear decoupling
hypothesis says that, in innermost variables and after translation-gauge
absorption of constant ambient velocities,
\[
  V_n=U_n+S_n,
  \qquad
  S_n\to0
\]
in the perturbative topology required for scale-collapse compactness, including pressure and modulation
compatibility.  Hence the innermost branch \(U_n\) satisfies the repaired-gauge
Navier--Stokes equation with bounded effective modulation, positive finite
critical mass, pressure reconstruction, and Caccioppoli regularity, up to an
error that vanishes in the scale-collapse limiting passage.

The scale \(\lambda_1^n\to0\) gives the scale-collapse hypothesis.  Therefore the
scale-collapse rigidity hypotheses produce a nonzero stationary
\(L^3(\mathbb R^3)\) profile
in the Nečas--Růžička--Šverák-covered self-similar class.  Nonzero extraction
follows from the active bubble mass floor:
\[
  \|\phi^1\|_3\ge\varepsilon_{\rm sd}>0.
\]
The Nečas--Růžička--Šverák rigidity theorem rules out such a nonzero
\(L^3\) stationary self-similar profile, contradicting the extraction.
\end{proof}

Thus same-point cascade exclusion leaves only the failure of same-point
nonlinear decoupling, namely the no-splitting/profile-compatibility
multibubble alternative.  If a regular strict same-point cascade of active singular profiles is
not excluded, then one of the following has failed: profile decomposition or
remainder decoupling, active-bubble small-data separation, escaping-offset
pressure decoupling, local profile regularity or innermost representation, or
the scale-collapse extraction of a nonzero limiting profile and rigidity
hypotheses.  These are analytic profile,
representation, interaction, or rigidity failures; a below-threshold profile is
regular and is not a singular cascade bubble.

\clearpage
\section{Local Single-Core Criterion}\label{sec:single-core-criterion}

\subsection{Setup and definitions}

For \(R>0\), write
\[
  Q_R=B_R\times(-R^2,0).
\]
For a suitable pair \((v,q)\) on \(Q_R\), define
\[
  A_v(R)=R^{-1}\operatorname*{ess\,sup}_{-R^2<s<0}
  \int_{B_R}|v(y,s)|^2\,dy,
\]
\[
  E_v(R)=R^{-1}\int_{Q_R}|\nabla v|^2\,dy\,ds,
\]
\[
  C_v(R)=R^{-2}\int_{Q_R}|v|^3\,dy\,ds,
\]
\[
  D_q(R)=R^{-2}\int_{-R^2}^0\int_{B_R}
  |q(y,s)-(q)_{B_R}(s)|^{3/2}\,dy\,ds.
\]
All pressure quantities are understood modulo functions of time.

\begin{definition}[Bounded selected-window limit]\label{paper1:def:bounded-window-limit}
Let \((V_n,P_n)\) be a sequence of represented local rescalings on a compact
cylinder \(Q_r\).  A selected-window subsequence has a bounded selected-window limit
on \(Q_r\) if, after passing to a subsequence,
\[
  V_n\to V_*
  \qquad\text{strongly in }L^3(Q_r),
\]
where \(V_*\in L^\infty(Q_r)\).  The limit is nonzero if
\(V_*\not\equiv0\) on \(Q_r\).
\end{definition}

\begin{definition}[Local Type II sequence]\label{paper1:def:typeII}
A represented sequence is called a local Type II sequence if it has positive
local Caffarelli--Kohn--Nirenberg concentration on some fixed compact cylinder
and no selected-window subsequence has a nonzero bounded selected-window limit on
any compact cylinder in the sense of \cref{paper1:def:bounded-window-limit}.
\end{definition}

\begin{definition}[Compact single-core vanishing-dissipation branch]\label{paper1:def:single-core-branch}
Fix radii \(0<r_*<R_*\), constants \(M<\infty\), \(\eta_0>0\), and a sequence
of represented local rescalings \((V_n,P_n,a_n,b_n)\) on \(Q_{R_*}\).  The
sequence is a compact single-core vanishing-dissipation branch if:
\begin{enumerate}[label=\textup{(\roman*)}]
\item each \((V_n,P_n)\) is suitable on \(Q_{R_*}\) and solves the represented
      Navier--Stokes equation
      \begin{equation}\label{paper1:eq:represented-ns}
      \partial_s V_n+(V_n\cdot\nabla)V_n+\nabla P_n
      =\nu\Delta V_n+a_n(s)(V_n+y\cdot\nabla V_n)+b_n(s)\cdot\nabla V_n,
      \qquad \nabla\cdot V_n=0;
      \end{equation}
\item the compact-cylinder suitable estimates are uniformly bounded:
      \[
      A_{V_n}(R_*)+E_{V_n}(R_*)+C_{V_n}(R_*)+D_{P_n}(R_*)\le M;
      \]
\item the selected core has positive CKN concentration on \(Q_{r_*}\):
      \[
      C_{V_n}(r_*)+D_{P_n}(r_*)\ge\eta_0
      \qquad\text{for all }n;
      \]
\item the localized scale and center are nondegenerate for the selected core,
      so the representation theorem gives the variables in
      \eqref{paper1:eq:represented-ns};
\item the modulation coefficients are uniformly bounded on the selected cylinder:
      \[
      \|a_n\|_{L^\infty(-R_*^2,0)}+\|b_n\|_{L^\infty(-R_*^2,0)}\le M;
      \]
\item the branch has vanishing localized dissipation in the sense of
      \cref{paper1:def:localized-dissipation} below.
\end{enumerate}
\end{definition}

\begin{definition}[Vanishing localized dissipation]\label{paper1:def:localized-dissipation}
Let \(r_*<\rho<R_*\).  Selected windows are first translated and rescaled to the
fixed cylinders \(Q_{r_*}\Subset Q_\rho\Subset Q_{R_*}\).  The localized
dissipation on \(Q_\rho\) is
\[
  \mathcal D_n(\rho)=\int_{Q_\rho}|\nabla V_n|^2\,dy\,ds.
\]
The branch has vanishing localized dissipation if, after passing to a
selected-window subsequence, there is \(\rho\in(r_*,R_*)\) such that
\[
  \mathcal D_n(\rho)\to0.
\]
\end{definition}

\subsection{Local analytic inputs}

The compact single-core argument is based on the following local analytic statements.

\begin{theorem}[Local concentration dichotomy]
\label{paper1:thm:paper0-dichotomy}
Let \((u,p)\) be a suitable weak solution on
\(\mathbb R^3\times(T-\delta,T)\), and suppose \(\Sigma(T)\ne\emptyset\).
Then, using the pressure-normalized quantities \(C,D\), there exist
\(x_0\in\Sigma(T)\), \(\eta>0\), and \(r_n\downarrow0\) such that
\[
  C((x_0,T),r_n)+D((x_0,T),r_n)\ge\eta
  \qquad\text{for all }n.
\]
Moreover, at the concentration point either a local Type I scale-invariant
bound holds, or the point lies in the local Type II alternative, meaning that
no such local Type I bound holds.
\end{theorem}

\begin{proof}
This is the local concentration dichotomy used throughout the paper.
\end{proof}

\begin{proposition}[Passage to local Type II sequences]
\label{paper1:prop:paper0-typeII-reduction}
Let \((x_0,T)\) be a concentration point in the local Type II alternative of
\cref{paper1:thm:paper0-dichotomy}.  After choosing the positive concentration scales
given there and applying the local scale-translation selection used in the
present section to a nondegenerate compact single core, the represented selected
windows form a sequence with positive local CKN concentration on a fixed compact
cylinder.  In the selected-window Type II branch, no subsequence has a nonzero
bounded selected-window limit on any compact cylinder.  Hence the represented
selected windows are a local Type II sequence in the sense of \cref{paper1:def:typeII}.
\end{proposition}

\begin{proof}
The local concentration dichotomy supplies the positive CKN concentration sequence and the exclusion of
the local Type I alternative at the chosen concentration point.  The
scale-translation selection rewrites those shrinking cylinders as fixed
selected windows.  The last condition is the selected-window
formulation of the Type II branch considered here: any branch with a nonzero
bounded selected-window limit belongs to the bounded-window alternative,
not to \cref{paper1:def:typeII}.
\end{proof}

\begin{proposition}[Local representation input]\label{paper1:prop:representation-input}
For a nondegenerate selected single core, the local scale-translation
construction gives represented variables \((V_n,P_n,a_n,b_n)\) satisfying
\eqref{paper1:eq:represented-ns} on compact cylinders.  Failure of the localized
nondegeneracy condition is a multibubble, cascade, or gauge-degenerate
alternative, not a compact single-core branch.
\end{proposition}

\begin{proposition}[Local pressure decomposition and stability]
\label{paper1:prop:pressure-reconstruction}
For every \(B_r\Subset B_R\), the represented pressure decomposes as
\[
  P_n=P_{{\rm loc},n}+H_n,
  \qquad
  -\Delta P_{{\rm loc},n}=\partial_i\partial_j(\zeta V_{n,i}V_{n,j}),
\]
where \(\zeta\in C_c^\infty(B_R)\) equals one on \(B_r\), and \(H_n\) is
harmonic in \(B_r\).  The local part satisfies the Calderon--Zygmund bound
\cite{CalderonZygmund1952,Stein1970}
\[
  \|P_{{\rm loc},n}\|_{L^{3/2}(B_R)}
  \le C\|V_n\|_{L^3(B_R)}^2
\]
for a.e. time, and the harmonic part is controlled modulo spatial means by the
local pressure norm.

Moreover, if \(V_n\to0\) strongly in \(L^3(Q_R)\), then for every \(r<R\),
\[
  P_n-(P_n)_{B_r}\to0
  \quad\text{strongly in }L^{3/2}(Q_r).
\]
\end{proposition}

\begin{lemma}[Localization of pressure oscillation bounds]\label{paper1:lem:D-localization}
Let \(0<r<R\), and let \(q\in L^{3/2}(Q_R)\).  Then
\[
  D_q(r)\le C\left(\frac{R}{r}\right)^2D_q(R),
\]
where \(C\) is universal.
\end{lemma}

\begin{proof}
For a.e. \(s\in(-r^2,0)\), Jensen's inequality gives
\[
  |(q)_{B_r}(s)-(q)_{B_R}(s)|^{3/2}
  \le \frac1{|B_r|}\int_{B_r}|q(y,s)-(q)_{B_R}(s)|^{3/2}\,dy .
\]
Hence
\[
  \int_{B_r}|q-(q)_{B_r}(s)|^{3/2}\,dy
  \le
  C\int_{B_r}|q-(q)_{B_R}(s)|^{3/2}\,dy .
\]
Integrating over \((-r^2,0)\) and using \(Q_r\subset Q_R\),
\[
\begin{aligned}
  D_q(r)
  &\le
  Cr^{-2}\int_{-r^2}^0\int_{B_r}
  |q-(q)_{B_R}(s)|^{3/2}\,dy\,ds \\
  &\le
  Cr^{-2}\int_{-R^2}^0\int_{B_R}
  |q-(q)_{B_R}(s)|^{3/2}\,dy\,ds
  =
  C\left(\frac{R}{r}\right)^2D_q(R).
\end{aligned}
\]
\end{proof}

\subsection{Compactness and pressure stability}

\begin{proposition}[Compactness on selected windows]\label{paper1:prop:compactness}
Let \((V_n,P_n)\) be suitable on \(Q_{R_*}\), solve \eqref{paper1:eq:represented-ns}, and satisfy
\[
  A_{V_n}(R_*)+E_{V_n}(R_*)+C_{V_n}(R_*)+D_{P_n}(R_*)\le M,
\]
with
\[
  \|a_n\|_{L^\infty(-R_*^2,0)}+\|b_n\|_{L^\infty(-R_*^2,0)}\le M.
\]
Then for every \(r<R_*\), after passing to a subsequence there are \((V,P)\)
such that
\[
  V_n\stackrel{*}{\rightharpoonup}V
  \quad\text{in }L^\infty(-r^2,0;L^2(B_r)),
\]
\[
  V_n\rightharpoonup V
  \quad\text{in }L^2(-r^2,0;H^1(B_r)),
\]
\[
  V_n\to V
  \quad\text{strongly in }L^3(Q_r),
\]
and
\[
  P_n-(P_n)_{B_r}\rightharpoonup P
  \quad\text{in }L^{3/2}(Q_r).
\]
After passing to a further subsequence, \(a_n\) and \(b_n\) converge weak-* in
\(L^\infty\) to bounded coefficients, and the limit satisfies the corresponding
represented equation distributionally on smaller cylinders.
\end{proposition}

\begin{proof}
The bounds on \(A\) and \(E\) give weak-* compactness in \(L^\infty L^2\) and
weak compactness in \(L^2H^1\) on smaller cylinders.  Fix
\(r<r_1<R_*\), and choose \(m\) so large that
\[
  H^m_0(B_{r_1})\hookrightarrow W^{1,\infty}_0(B_{r_1}).
\]
We claim that \(\partial_sV_n\) is uniformly bounded in
\[
  L^1(-r_1^2,0;H^{-m}(B_{r_1})).
\]
Let \(\varphi\in H^m_0(B_{r_1})\).  Pairing \eqref{paper1:eq:represented-ns} with
\(\varphi\), integrating by parts in space where needed, and using the
embedding above gives
\[
  |\langle (V_n\cdot\nabla)V_n,\varphi\rangle|
  =
  \left|\int_{B_{r_1}}V_{n,i}V_{n,j}\partial_j\varphi_i\,dy\right|
  \le
  C\|V_n\|_{L^3(B_{r_1})}^2\|\varphi\|_{H^m},
\]
\[
  |\langle \nu\Delta V_n,\varphi\rangle|
  \le
  C\|\nabla V_n\|_{L^2(B_{r_1})}\|\varphi\|_{H^m},
\]
\[
  |\langle \nabla P_n,\varphi\rangle|
  =
  \left|\int_{B_{r_1}}(P_n-(P_n)_{B_{r_1}})\nabla\cdot\varphi\,dy\right|
  \le
  C\|P_n-(P_n)_{B_{r_1}}\|_{L^{3/2}(B_{r_1})}\|\varphi\|_{H^m}.
\]
The modulation terms obey
\[
  |\langle a_nV_n,\varphi\rangle|
  \le
  C\|V_n\|_{L^2(B_{r_1})}\|\varphi\|_{H^m},
\]
\[
  |\langle a_n\,y\cdot\nabla V_n,\varphi\rangle|
  =
  \left|\int_{B_{r_1}}a_n V_{n,i}\partial_j(y_j\varphi_i)\,dy\right|
  \le
  C\|V_n\|_{L^2(B_{r_1})}\|\varphi\|_{H^m},
\]
and
\[
  |\langle b_n\cdot\nabla V_n,\varphi\rangle|
  =
  \left|\int_{B_{r_1}}V_{n,i}b_{n,j}\partial_j\varphi_i\,dy\right|
  \le
  C\|V_n\|_{L^2(B_{r_1})}\|\varphi\|_{H^m},
\]
where the constant uses \(r_1\) and the uniform \(L^\infty\) bounds for
\(a_n,b_n\).  Integrating these estimates in time gives a uniform bound because
\[
  \int_{-r_1^2}^0\|V_n(s)\|_{L^3(B_{r_1})}^2\,ds
  \le
  r_1^{2/3}\|V_n\|_{L^3(Q_{r_1})}^2,
\]
\[
  \int_{-r_1^2}^0\|\nabla V_n(s)\|_{L^2(B_{r_1})}\,ds
  \le
  r_1\|\nabla V_n\|_{L^2(Q_{r_1})},
\]
\[
  \int_{-r_1^2}^0
  \|P_n(s)-(P_n)_{B_{r_1}}(s)\|_{L^{3/2}(B_{r_1})}\,ds
  \le
  r_1^{2/3}
  \|P_n-(P_n)_{B_{r_1}}\|_{L^{3/2}(Q_{r_1})},
\]
and
\[
  \int_{-r_1^2}^0\|V_n(s)\|_{L^2(B_{r_1})}\,ds
  \le
  r_1^2
  \|V_n\|_{L^\infty(-r_1^2,0;L^2(B_{r_1}))}.
\]
The uniform \(C\), \(E\), \(D\), and \(A\) bounds therefore prove the claimed
time-derivative bound; the \(D\)-bound on \(Q_{r_1}\) follows from the bound on
\(Q_{R_*}\) by \cref{paper1:lem:D-localization}.

The compact embedding \(H^1(B_{r_1})\Subset L^2(B_r)\), together with
Aubin--Lions compactness \cite{Aubin1963,LionsMagenes1972,Simon1987}, gives
\[
  V_n\to V\qquad\text{strongly in }L^2(Q_r)
\]
after passing to a subsequence.  The same \(L^\infty L^2\cap L^2H^1\) bounds
give the parabolic interpolation estimate
\[
  \|V_n\|_{L^{10/3}(Q_r)}\le C(r,M).
\]
Indeed, for a.e. \(s\), the Sobolev and Gagliardo--Nirenberg inequalities give
\[
  \|V_n(s)\|_{L^{10/3}(B_r)}
  \le
  C\|V_n(s)\|_{L^2(B_r)}^{2/5}
  \|\nabla V_n(s)\|_{L^2(B_r)}^{3/5}.
\]
Raising to the power \(10/3\), integrating in time, and using the \(A\) and
\(E\) bounds yields
\[
  \int_{-r^2}^0\|V_n(s)\|_{L^{10/3}(B_r)}^{10/3}\,ds
  \le
  C
  \left(\operatorname*{ess\,sup}_{-r^2<s<0}
  \|V_n(s)\|_{L^2(B_r)}^2\right)^{2/3}
  \int_{Q_r}|\nabla V_n|^2\,dy\,ds
  \le C(r,M).
\]
The same estimate applies to \(V\) by weak lower semicontinuity, and therefore
to \(V_n-V\) by the triangle inequality.  Interpolating between \(L^2\) and
\(L^{10/3}\), one obtains
\[
  \|V_n-V\|_{L^3(Q_r)}
  \le
  \|V_n-V\|_{L^2(Q_r)}^{1/6}
  \|V_n-V\|_{L^{10/3}(Q_r)}^{5/6}.
\]
The second factor is uniformly bounded and the first tends to zero, so
\[
  V_n\to V\qquad\text{strongly in }L^3(Q_r).
\]

The pressure compactness follows from the uniform \(D\)-bound:
after subtracting spatial means,
\[
  P_n-(P_n)_{B_r}
  \rightharpoonup P
  \qquad\text{in }L^{3/2}(Q_r).
\]
We pass to the limit distributionally using the preceding convergences.  It
remains to justify the modulation terms.  For every compactly supported
test field \(\varphi\),
\[
  \int a_n(s)\, y\cdot\nabla V_n\cdot\varphi
  =
  -\int a_n(s)\,V_{n,i}\,\partial_j(y_j\varphi_i),
\]
where summation over repeated spatial indices is understood.  Equivalently,
\[
  \partial_j(y_j\varphi_i)=3\varphi_i+y_j\partial_j\varphi_i .
\]
Hence this term passes to the limit by strong local \(L^2\) convergence of
\(V_n\) and weak-* convergence of \(a_n\): if \(F_n\to F\) in \(L^1\) and
\(a_n\rightharpoonup^* a\) in \(L^\infty\), then
\[
  \int a_nF_n-\int aF
  =
  \int a_n(F_n-F)+\int (a_n-a)F\to0.
\]
The term \(a_nV_n\) is immediate from the same argument, while
\[
  \int b_n(s)\cdot\nabla V_n\cdot\varphi
  =
  -\int V_{n,i}\,b_{n,j}(s)\,\partial_j\varphi_i
\]
passes to the limit by strong local \(L^2\) convergence of \(V_n\) and weak-*
convergence of \(b_n\).  These passages give the limiting represented equation
in the sense of distributions.
\end{proof}

\begin{remark}[Use of bounded modulation]\label{paper1:rem:bounded-modulation}
The uniform \(L^\infty\) bounds on \(a_n\) and \(b_n\) enter the compactness
argument in exactly two places.  First, they give the negative Sobolev
time-derivative bound for the three lower-order modulation terms in
\eqref{paper1:eq:represented-ns}.  Second, after subsequential weak-* convergence of
\(a_n,b_n\), they allow the modulation terms to pass to the distributional
limit once \(V_n\to V\) strongly in \(L^2_{\rm loc}\).
\end{remark}

\begin{lemma}[Stability of local pressure decomposition]\label{paper1:lem:pressure-stability}
Let \(V_n\to V\) strongly in \(L^3(Q_R)\), and use the pressure decomposition
of \cref{paper1:prop:pressure-reconstruction}.  Then for every \(r<R\),
\[
  P_{{\rm loc},n}\to P_{\rm loc}
  \quad\text{strongly in }L^{3/2}(Q_r),
\]
where \(-\Delta P_{\rm loc}=\partial_i\partial_j(\zeta V_iV_j)\).  If
\(V\equiv0\) on \(Q_R\), then
\[
  P_n-(P_n)_{B_r}\to0
  \quad\text{strongly in }L^{3/2}(Q_r).
\]
More locally in time, if \(V_n\to0\) strongly in
\[
  L^3(B_R\times(-r^2,0)),
\]
then
\[
  P_n-(P_n)_{B_r}\to0
  \quad\text{strongly in }L^{3/2}(Q_r).
\]
\end{lemma}

\begin{proof}
Strong convergence in \(L^3\) gives
\[
  V_{n,i}V_{n,j}\to V_iV_j
  \quad\text{strongly in }L^{3/2}(Q_R).
\]
Indeed,
\[
  \|V_{n,i}V_{n,j}-V_iV_j\|_{L^{3/2}(Q_R)}
  \le
  \bigl(\|V_n\|_{L^3(Q_R)}+\|V\|_{L^3(Q_R)}\bigr)
  \|V_n-V\|_{L^3(Q_R)}.
\]
Calderon--Zygmund estimates for the compactly supported localized source yield
strong convergence of the localized pressures in \(L^{3/2}\) on smaller balls:
\[
  \|P_{{\rm loc},n}-P_{\rm loc}\|_{L^{3/2}(Q_r)}
  \le
  C_{r,R}
  \|V_{n,i}V_{n,j}-V_iV_j\|_{L^{3/2}(Q_R)}.
\]

Assume now that \(V\equiv0\) on \(Q_R\).  Then \(P_{\rm loc}=0\).  By the
zero-limit stability part of \cref{paper1:prop:pressure-reconstruction},
\[
  P_n-(P_n)_{B_r}\to0
  \quad\text{strongly in }L^{3/2}(Q_r).
\]
The same proof is local in time.  If \(V_n\to0\) strongly in
\(L^3(B_R\times(-r^2,0))\), the Calderon--Zygmund estimate and harmonic
pressure control are applied only for times \(s\in(-r^2,0)\), giving the
strong convergence on \(Q_r\).
\end{proof}

\begin{lemma}[Zero-dissipation limit]\label{paper1:lem:zero-dissipation}
If
\[
  V_n\rightharpoonup V\quad\text{in }L^2(I;H^1(B_\rho))
\]
and
\[
  \int_I\int_{B_\rho}|\nabla V_n|^2\,dy\,ds\to0,
\]
then \(\nabla V=0\) a.e. on \(B_\rho\times I\).  Hence for a.e. time,
\(V(\cdot,s)\) is spatially constant on \(B_\rho\).
\end{lemma}

\begin{proof}
Weak lower semicontinuity gives
\[
  \int_I\int_{B_\rho}|\nabla V|^2\,dy\,ds
  \le \liminf_{n\to\infty}\int_I\int_{B_\rho}|\nabla V_n|^2\,dy\,ds=0.
\]
Thus \(\nabla V=0\) a.e. on \(B_\rho\times I\).  By Fubini, for a.e.
\(s\in I\) one has \(\nabla V(\cdot,s)=0\) in \(L^2(B_\rho)\).  Since
\(B_\rho\) is connected, every \(H^1(B_\rho)\) function with zero weak
gradient is a.e. equal to a constant.  Hence \(V(\cdot,s)\) is spatially
constant for a.e. \(s\).
\end{proof}

\begin{lemma}[Spatially constant limits]\label{paper1:lem:constant-limit}
Let \(V\) be a limit on \(Q_\rho\) and suppose \(\nabla V=0\) a.e. there.
Then \(V(y,s)=c(s)\) on \(Q_\rho\), with
\(c\in L^\infty(-\rho^2,0)\).  If \(c\not\equiv0\) on \(Q_\rho\), then
the approximating selected-window subsequence has a nonzero bounded
selected-window limit.
\end{lemma}

\begin{proof}
The spatial constancy follows from \cref{paper1:lem:zero-dissipation}.  Since
\(V(\cdot,s)=c(s)\) on \(B_\rho\), the local \(L^\infty_sL^2_y\) bound gives
\[
  |c(s)|^2|B_\rho|\le \operatorname*{ess\,sup}_{-\rho^2<s<0}
  \int_{B_\rho}|V(y,s)|^2\,dy,
\]
so \(c\in L^\infty(-\rho^2,0)\).  If \(c\not\equiv0\) on \(Q_\rho\), then the
strong \(L^3\) convergence of \cref{paper1:prop:compactness} gives a subsequential
limit to a nonzero bounded function on \(Q_\rho\), which is precisely a nonzero
bounded selected-window limit in the sense of \cref{paper1:def:bounded-window-limit}.
\end{proof}

\subsection{Contradiction lemmas}

\begin{lemma}[Persistence of positive core concentration]\label{paper1:lem:persistence}
For a compact single-core branch in the sense of \cref{paper1:def:single-core-branch},
there are fixed \(r_*>0\) and \(\eta_0>0\) such that along the selected-window
subsequence,
\[
  C_{V_n}(r_*)+D_{P_n}(r_*)\ge\eta_0
  \qquad\text{for all }n.
\]
\end{lemma}

\begin{proof}
The lower bound is part of the definition of the selected compact core and is preserved by
passing to subsequences of the selected windows.
\end{proof}

\begin{lemma}[Zero constant contradicts positive CKN concentration]\label{paper1:lem:zero-contradiction}
Let \(r_*<\rho\).  Suppose the compactness limit in \cref{paper1:prop:compactness} is
spatially constant on \(Q_\rho\) and is identically zero on \(Q_{r_*}\).  Then
\[
  C_{V_n}(r_*)+D_{P_n}(r_*)\to0.
\]
Consequently this case contradicts \cref{paper1:lem:persistence}.
\end{lemma}

\begin{proof}
Strong \(L^3(Q_{r_*})\) convergence to zero gives
\[
  C_{V_n}(r_*)
  =
  r_*^{-2}\int_{Q_{r_*}}|V_n|^3\,dy\,ds
  =
  r_*^{-2}\|V_n\|_{L^3(Q_{r_*})}^3
  \to0.
\]
It remains to control the pressure term.  Since \(V\) is spatially constant on
\(Q_\rho\), the identity \(V=0\) on \(Q_{r_*}\) implies
\[
  V=0
  \qquad\text{on }B_\rho\times(-r_*^2,0).
\]
Indeed, for a.e. \(s\in(-r_*^2,0)\), the spatially constant value on
\(B_\rho\) vanishes on the nonempty subball \(B_{r_*}\), hence vanishes on all
of \(B_\rho\).  The strong \(L^3_{\rm loc}\) convergence therefore gives
\[
  V_n\to0
  \qquad\text{strongly in }L^3(B_\rho\times(-r_*^2,0)).
\]
Applying \cref{paper1:lem:pressure-stability} with \(R=\rho\) and \(r=r_*\), in its
local-in-time form, gives
\[
  P_n-(P_n)_{B_{r_*}}\to0
  \qquad\text{strongly in }L^{3/2}(Q_{r_*}).
\]
Therefore
\[
  D_{P_n}(r_*)
  =
  r_*^{-2}
  \|P_n-(P_n)_{B_{r_*}}\|_{L^{3/2}(Q_{r_*})}^{3/2}
  \to0.
\]
Thus
\[
  C_{V_n}(r_*)+D_{P_n}(r_*)\to0,
\]
contradicting the uniform lower bound
\[
  C_{V_n}(r_*)+D_{P_n}(r_*)\ge\eta_0.
\]
\end{proof}

\begin{lemma}[Nonzero constant contradicts the Type II regime]\label{paper1:lem:nonzero-contradiction}
If the compactness limit is spatially constant and nonzero on a compact
subcylinder, then the original sequence is not a local Type II sequence.
\end{lemma}

\begin{proof}
By \cref{paper1:lem:constant-limit}, the selected-window subsequence has a nonzero
bounded selected-window limit.  This contradicts \cref{paper1:def:typeII}.
\end{proof}

\subsection{Main theorem}

\begin{theorem}[Local single-core Type II criterion]\label{paper1:thm:single-core-criterion}
Under the analytic inputs
\cref{paper1:prop:paper0-typeII-reduction,paper1:prop:representation-input,paper1:prop:pressure-reconstruction},
a compact single-core vanishing-dissipation branch in the sense of
\cref{paper1:def:single-core-branch} cannot be a local Type II sequence.
\end{theorem}

\begin{proof}
By the definition of vanishing localized dissipation, after passing to a
selected-window subsequence there is a fixed radius \(\rho\in(r_*,R_*)\) such
that
\[
  \int_{Q_\rho}|\nabla V_n|^2\,dy\,ds\to0
\]
and the positive core bound on \(Q_{r_*}\) remains valid along this subsequence.
Apply compactness on the fixed compact
cylinder \(Q_\rho\) and pass to a subsequence.  The compactness proposition
gives
\[
  V_n\rightharpoonup V\quad\text{in }L^2(-\rho^2,0;H^1(B_\rho))
\]
and strong convergence in \(L^3\) on every smaller cylinder.  Since the
localized dissipation tends to zero, the zero-dissipation lemma gives
\[
  \nabla V=0\quad\text{a.e. on }Q_\rho.
\]
Thus \(V(y,s)=c(s)\) on \(Q_\rho\), and in particular on the selected core
cylinder \(Q_{r_*}\).

If this spatially constant limit is zero on the core cylinder, the zero-limit
contradiction lemma gives
\[
  C_{V_n}(r_*)+D_{P_n}(r_*)\to0,
\]
contradicting persistence of positive core concentration.

If the spatially constant limit is not identically zero on \(Q_{r_*}\), the
nonzero-limit lemma produces a nonzero bounded selected-window limit, contradicting
the definition of a local Type II sequence.  The two cases exhaust all
spatially constant limits, so the branch cannot occur.
\end{proof}

\begin{remark}[Minimal local inputs]
The proof uses two local analytic inputs from the Type II argument: the
localized representation for a nondegenerate selected core, the compact-ball
pressure reconstruction with stability on smaller balls.  The
compact-cylinder \(A,E,C,D\) bounds and vanishing localized dissipation are
part of the branch hypothesis, not separate analytic inputs.  Multibubble,
cascade, gauge-degenerate, and compactness-failure alternatives are outside
this single-core theorem; they are treated by the local multibubble and
cascade exclusion theorem.
\end{remark}

\clearpage
\section{Repaired-Gauge Construction}\label{sec:repaired-gauge-construction}
\subsection{Assumptions and notation}\label{paper2:sec:assumptions}

Let \(u:\RR^{3}\times (T-\delta,T)\to \RR^{3}\), \(p:\RR^{3}\times (T-\delta,T)\to\RR\)
be a suitable weak Navier--Stokes solution with viscosity \(\nu>0\).
For \(z=(x_0,t_0)\) and \(R>0\),
\[
Q_R(z):=B_R(x_0)\times (t_0-R^2,t_0),\qquad
Q_R^*:=B_R\times I_*.
\]
Fix a reference physical center-time \((x_\bullet,t_\star)\) around which \(B_{R_*}(x_\bullet)\) is taken.
For a fixed spatial center \(x_0\) we also use
\[
Q_R(t):=Q_R(x_0,t).
\]
Fix a compact window \(I_*:=[\tau_0-\ell,\tau_0+\ell]\), \(\ell>0\), and write
\(I_{\rm ph}:=[t_\star-\Lambda_{\max}^2\ell,\; t_\star+\Lambda_{\max}^2\ell]\).

For a cylinder \(Q_R(x_0,t_0)=B_R(x_0)\times (t_0-R^2,t_0)\), set
\[
\begin{aligned}
A(u;Q_R(x_0,t_0))&:=R^{-1}\sup_{s\in(t_0-R^2,t_0)}
\int_{B_R(x_0)}|u(x,s)|^2\,dx,\\
E(u;Q_R(x_0,t_0))&:=R^{-1}\int_{Q_R(x_0,t_0)}|\nabla u|^2\,dx\,ds,\\
C(u;Q_R(x_0,t_0))&:=R^{-2}\int_{Q_R(x_0,t_0)}|u|^3\,dx\,ds,\\
D(p;Q_R(x_0,t_0))&:=R^{-2}\int_{t_0-R^2}^{t_0}\int_{B_R(x_0)}
\left|p(x,s)-\langle p(s)\rangle_{B_R(x_0)}\right|^{3/2}\,dx\,ds.
\end{aligned}
\]
Here \(\langle p(s)\rangle_{B_R(x_0)}:=|B_R|^{-1}\int_{B_R(x_0)}p(x,s)\,dx\);
pressure is always understood modulo time-dependent constants.

\begin{definition}[Single-core hypotheses]\label{paper2:def:SC}
Let \((R_*,r_*,M,\eta,\kappa,\Lambda_{\min},\Lambda_{\max})\) satisfy
\(\Lambda_{\min}\le \Lambda_{\max}<\infty\).
We say \((u,p)\) satisfies \(\mathrm{SC}(R_*,r_*,M,\eta,\kappa)\) on \(Q_{R_*}(x_\bullet,t_\star)\)
if the following hold:
\begin{enumerate}
\item[(\textup{SC1})] \(A(u;Q_{R_*}(x_\bullet,t)) + E(u;Q_{R_*}(x_\bullet,t))
  + C(u;Q_{R_*}(x_\bullet,t)) + D(p;Q_{R_*}(x_\bullet,t))\le M\) for every
  \(t\in[t_\star-R_*^2,t_\star)\).
\item[(\textup{SC2})] There exists a physically selected core map
  \((x_{\rm c}(t),\rho_{\rm c}(t))\) on \(I_{\rm ph}\) such that
  \(\rho_{\rm c}(t)\in(0,R_*]\), \(x_{\rm c}(t)\in B_{R_*/2}(x_\bullet)\), and for some
  \(\Theta_0>0\),
  \[
  \int_{B_{r_*}(0)}|\rho_{\rm c}(t)^{-1}u(x_{\rm c}(t)+\rho_{\rm c}(t)y,t)|^3\dd y\ge \eta
  \quad \text{for a.e. }t\in I_{\rm ph}.
  \]
\item[(\textup{SC3})] There exists \(0<r_{\rm sep}\le r_*/2\) such that for every
  \(0<r\le r_{\rm sep}\), every competitor pair
  \[
  \mathcal C_{\rm core}(t):=\left\{(\bar x,\bar\rho):|\bar x-x_\bullet|\le \frac{R_*}{2},\
  0<\bar\rho\le R_*\right\},
  \]
  \((\bar x,\bar\rho)\in \mathcal C_{\rm core}(t)\setminus\{(x_{\rm c}(t),\rho_{\rm c}(t))\}\)
  with
  \[
  \rho_{\rm c}(t)^{-1}|\bar x-x_{\rm c}(t)|
  +\left|\frac{\bar\rho}{\rho_{\rm c}(t)}-1\right|\ge r_{\rm sep}
  \]
  has separated local weighted mass in the same core-scale variables:
  \[
  \left|\int_{B_r}|\rho_{\rm c}(t)^{-1}u(x_{\rm c}(t)+\rho_{\rm c}(t)y,t)|^3\,\dd y
  - \int_{B_r}|\bar\rho^{-1}u(\bar x+\bar\rho y,t)|^3\,\dd y\right|
  \ge \eta.
  \]
  This quantitative non-competing core condition is imposed on the
  selected physical core directly; it is not a consequence of the repaired-gauge
  construction.
\item[(\textup{SC4})] For the selected core,
  \[
  \sup_{t\in I_{\rm ph}}\int_{B_{2R_*}}|y|^{-p}|\rho_{\rm c}(t)^{-1}u(x_{\rm c}(t)+\rho_{\rm c}(t) y,t)|^3\,dy\le M,
  \]
  for some \(0<p<3\).
\item[(\textup{SC5})] The preliminary frame associated with the selected core is
  the absolutely continuous representative
  \[
  0<\Lambda_{\min}\le \Lambda_{\rm pre}(\tau)\le \Lambda_{\max},\quad
  \Lambda_{\rm pre}\in AC(I_*),\quad
  X_{\rm pre}\in AC(I_*;\RR^3).
  \]
  It is tied to the selected physical core by the time change
  \(dt_{\rm c}/d\tau=\Lambda_{\rm pre}(\tau)^2\), \(t_{\rm c}(\tau_0)=t_\star\):
  one has
  \[
  X_{\rm pre}(\tau)=x_{\rm c}(t_{\rm c}(\tau)),\qquad
  \Lambda_{\rm pre}(\tau)=\rho_{\rm c}(t_{\rm c}(\tau))
  \]
  for a.e. \(\tau\in I_*\).
\end{enumerate}
\end{definition}
\noindent
The parameter \(\kappa\) denotes the quantitative constants obtained in the local gauge estimates and is not used in items \textup{(SC1)--(SC5)}.

\begin{remark}
The selected core \((x_{\rm c},\rho_{\rm c})\) and its separation/nondegeneracy
conditions are specified in physical variables; no repaired-gauge constraint is used in
their definition.  This avoids the circularity of defining the core via the chart
to be constructed.
\end{remark}

\subsection{Preliminary frame map from the selected core}\label{paper2:sec:preliminary-frame}

The transport map is constructed from the selected core; the core selection
itself is part of the hypotheses.

Define
\[
\frac{dt_{\rm c}}{d\tau}=\Lambda_{\rm pre}(\tau)^2,\qquad t_{\rm c}(\tau_0)=t_\star,
\]
and use the compatibility in \textup{(SC5)} to identify
\(\Lambda_{\rm pre}(\tau)=\rho_{\rm c}(t_{\rm c}(\tau))\) and
\(X_{\rm pre}(\tau)=x_{\rm c}(t_{\rm c}(\tau))\) a.e.

\begin{proposition}[Preliminary frame map regularity]\label{paper2:prop:prelim-frame}
Under \(\mathrm{SC}(R_*,r_*,M,\eta,\kappa)\), the map
\[
\Phi_{\rm pre}(\tau,Y):=(x,t)=(X_{\rm pre}(\tau)+\Lambda_{\rm pre}(\tau)Y,\ t_{\rm c}(\tau))
\]
is well-defined and absolutely continuous in \(\tau\).
Moreover, \(1/\Lambda_{\rm pre}\), \(\Lambda_{\rm pre}\in L^\infty(I_*)\),
\(X_{\rm pre}\in W^{1,1}_{\rm loc}(I_*;\RR^3)\), and for every \(R\le R_*\),
\[
\Phi_{\rm pre}:Q_R^*\to \Phi_{\rm pre}(Q_R^*)\subset Q_{2R_*}(x_\bullet,t_\star)
\]
is bi-Lipschitz with \(\det \nabla_Y\Phi_{\rm pre}=\Lambda_{\rm pre}^3\), hence
Jacobian and inverse-Jacobian bounds depend only on \(\Lambda_{\min},\Lambda_{\max}\).
\end{proposition}
\begin{proof}
From \textup{(SC5)}, \(\Lambda_{\rm pre}\in AC(I_*)\),
\(0<\Lambda_{\min}\le \Lambda_{\rm pre}\le\Lambda_{\max}<\infty\), and
\(X_{\rm pre}\in AC(I_*;\mathbb R^3)\). The scalar ODE
\[
\frac{dt_{\rm c}}{d\tau}=\Lambda_{\rm pre}(\tau)^2,\qquad t_{\rm c}(\tau_0)=t_\star
\]
therefore has an absolutely continuous strictly increasing solution on \(I_*\), and
\[
\frac{dt_{\rm c}}{d\tau}=\Lambda_{\rm pre}(\tau)^2
\quad\text{a.e. on }I_*.
\]
The compatibility condition \textup{(SC5)} identifies this absolutely continuous frame
with the selected physical core a.e.; no differentiability of \(x_{\rm c}\) or
\(\rho_{\rm c}\) beyond the frame regularity in \textup{(SC5)} is used.
For fixed \(\tau\), \(\nabla_Y\Phi_{\rm pre}(\tau,\cdot)=\Lambda_{\rm pre}(\tau)I\), so
\[
\det\nabla_Y\Phi_{\rm pre}(\tau,\cdot)=\Lambda_{\rm pre}(\tau)^3,
\qquad
\Lambda_{\min}^3\le \det\nabla_Y\Phi_{\rm pre}\le \Lambda_{\max}^3.
\]
Hence, if \(Y_1,Y_2\in B_R\),
\[
\Lambda_{\min}|Y_1-Y_2|
\le |\Phi_{\rm pre}(\tau,Y_1)-\Phi_{\rm pre}(\tau,Y_2)|
\le \Lambda_{\max}|Y_1-Y_2|,
\]
and therefore \(\Phi_{\rm pre}(\tau,\cdot)\) is bi-Lipschitz on \(B_R\). Its inverse has
Lipschitz constant \(\Lambda_{\min}^{-1}\), so Jacobian and inverse-Jacobian bounds
are controlled by \(\Lambda_{\min},\Lambda_{\max}\).

For image inclusions, if \(x=X_{\rm pre}(\tau)+\Lambda_{\rm pre}(\tau)Y\) with
\(|Y|\le R\), then
\[
|x-X_{\rm pre}(\tau)|\le \Lambda_{\max}R,\qquad
X_{\rm pre}(\tau)\in B_{R_*/2}(x_\bullet),
\]
so \(x\in B_{R_*/2+\Lambda_{\max}R}(x_\bullet)\). As \(\tau\in I_*\), we also have
\(t=t_{\rm c}(\tau)\in I_{\rm ph}\). Therefore
\[
\Phi_{\rm pre}(Q_R^*)\subset B_{R_*/2+\Lambda_{\max}R}(x_\bullet)\times I_{\rm ph}.
\]
\end{proof}
\begin{definition}[Preliminary profile]\label{paper2:def:prelim-profile}
For \((x,t)=\Phi_{\rm pre}(\tau,Y)\), define
\[
V_{\rm pre}(Y,\tau):=\Lambda_{\rm pre}(\tau)\,u(x,t),\qquad
P_{\rm pre}(Y,\tau):=\Lambda_{\rm pre}(\tau)^2\,p(x,t).
\]
Hence \(u(x,t)=\Lambda_{\rm pre}(\tau)^{-1}V_{\rm pre}(Y,\tau)\) and
\(p(x,t)=\Lambda_{\rm pre}(\tau)^{-2}P_{\rm pre}(Y,\tau)\) under \(\Phi_{\rm pre}\).
\end{definition}

\subsection{Exact renormalized change of variables}\label{paper2:sec:change}

Define renormalized fields on \(Q_{2R}^*\) by
\[
Y=\frac{x-X(\tau)}{\Lambda(\tau)},\qquad
u(x,t)=\Lambda(\tau)^{-1}V(Y,\tau),\qquad
p(x,t)=\Lambda(\tau)^{-2}P(Y,\tau),
\]
where \(dt/d\tau=\Lambda(\tau)^2\), \(X:\,I_*\to\RR^3\), \(\Lambda:\,I_*\to(0,\infty)\),
and \(\Lambda,X\in AC(I_*)\).
Write
\[
a(\tau):=-\Lambda'(\tau),\qquad
b(\tau):=-X'(\tau).
\]

\begin{lemma}[Exact differential identities]\label{paper2:lem:change-id}
For smooth data and fixed \((x,t)\), the chain rule gives
\[
\partial_t u
=\Lambda^{-3}\left(\partial_\tau V-(X'+\Lambda' Y)\cdot\nabla_Y V-\Lambda'V\right),
\]
\[
\nn_x u=\Lambda^{-2}\nn_Y V,\qquad \Delta_x u=\Lambda^{-3}\Delta_Y V,\qquad
\nn_x p=\Lambda^{-3}\nn_Y P.
\]
In particular
\[
\nn_x\cdot u=\Lambda^{-2}\nn_Y\cdot V.
\]
\end{lemma}
\begin{proof}
For fixed \((x,t)\), \(\tau=\tau(t)\) and \(Y=(x-X(\tau))/\Lambda(\tau)\).
Differentiate
\[
u(x,t)=\Lambda(\tau)^{-1}V\bigl(Y,\tau\bigr)
\]
in \(t\):
\[
\partial_t u
=-\frac{\Lambda'}{\Lambda^2}V
+\frac1\Lambda\partial_\tau V\,\frac{d\tau}{dt}
+\frac1\Lambda\nabla_YV\cdot\frac{dY}{dt}.
\]
Since \(d\tau/dt=\Lambda^{-2}\) and
\[
Y_\tau=-\Lambda^{-1}X'-\frac{\Lambda'}{\Lambda}Y,\qquad
\frac{dY}{dt}=Y_\tau\,\frac{d\tau}{dt},
\]
therefore
\[
\partial_tu
=\Lambda^{-3}\left(\partial_\tau V-(X'+\Lambda'Y)\cdot\nabla_YV-\Lambda'V\right).
\]
For spatial derivatives, \(Y=(x-X)/\Lambda\) gives \(\nabla_xY=\Lambda^{-1}I\), hence
\[
\nabla_xu=\Lambda^{-2}\nabla_YV,\qquad
\Delta_xu=\Lambda^{-3}\Delta_YV.
\]
Similarly \(p(x,t)=\Lambda^{-2}P(Y,\tau)\) yields
\[
\nabla_x p=\Lambda^{-3}\nabla_YP.
\]
Taking divergence of \(u=\Lambda^{-1}V\) gives
\[
\nabla_x\cdot u=\Lambda^{-2}\nabla_Y\cdot V.
\]
\end{proof}

\begin{theorem}[Renormalized Navier--Stokes equation]\label{paper2:thm:ren-NS}
Assume \(u,p\) solve
\[
\partial_t u+(u\cdot\nn_x)u+\nn_x p-\nu\Delta_xu=0,\qquad \nn_x\cdot u=0
\]
in distributions. Under the transformation above, \(V,P\) satisfy, on the image
renormalized cylinder, the exact equation
\[
\partial_\tau V+(V\cdot\nn)V+\nn P-\nu\Delta V
+a(\tau)\bigl(V+Y\cdot\nn V\bigr)+b(\tau)\cdot\nn V=0,\qquad \nn\cdot V=0,
\]
with the sign convention fixed by \(a=-\Lambda'\), \(b=-X'\).
Conversely, any \(V,P\) satisfying this equation pull back to a distributional
solution \(u,p\) of Navier--Stokes on the physical image.
\end{theorem}
\begin{proof}
To avoid hidden regularity loss, first write the physical equation tested against
\(\Phi\in C_c^\infty\) in the physical cylinder and then use the smooth pull-back formula
(an approximation argument gives distributions). By Lemma~\ref{paper2:lem:change-id},
\[
\partial_tu+\bigl(u\cdot\nabla_x\bigr)u+\nabla_xp-\nu\Delta_xu=0
\]
becomes
\[
\Lambda^{-3}\!\left(\partial_\tau V-(X'+\Lambda'Y)\cdot\nabla V-\Lambda'V\right)
+\Lambda^{-3}(V\cdot\nabla)V
+\Lambda^{-3}\nabla P
-\Lambda^{-3}\nu\Delta V=0.
\]
Multiplying by \(\Lambda^3\) and using \(a=-\Lambda'\), \(b=-X'\) gives
\[
\partial_\tau V+(V\cdot\nabla)V+\nabla P-\nu\Delta V
+a\bigl(V+Y\cdot\nabla V\bigr)+b\cdot\nabla V=0.
\]
The divergence condition is
\[
\nabla_x\cdot u=\Lambda^{-2}\nabla\cdot V=0
\quad\Longrightarrow\quad
\nabla\cdot V=0.
\]
Conversely, assume \((V,P)\) solve this equation on the renormalized cylinder.
Set \(u,p\) by the same change-of-variables map.
Then the preceding identities run in reverse and give
 \[
\partial_t u+(u\cdot\nabla_x)u+\nabla_xp-\nu\Delta_xu=0,\qquad
\nabla_x\cdot u=0
\]
in \(\mathcal D'\) on the physical image.
\end{proof}

\subsection{Domain localization and support}\label{paper2:sec:domains}

Let \(R>0\) be fixed and \(0<\ell\le \ell_0\). Set \(Q_R^*=B_R\times I_*\).
With \(\Phi\) as above define physical support maps
\[
\mathcal Q^{\rm ph}_{R}:=\Phi(Q_R^*),\qquad
\mathcal Q^{\rm ph}_{2R}:=\Phi(Q_{2R}^*).
\]

\begin{proposition}[Compact support transport]\label{paper2:prop:domain}
Assume \(\Lambda_{\min}\le \Lambda\le \Lambda_{\max}\) on \(I_*\).
Then
\[
\begin{gathered}
X(\tau)\in B_{R_*/2}(x_\bullet),\qquad |x-X(\tau)|\le \Lambda_{\max}R,\\
\mathcal Q^{\rm ph}_{R}\subset B_{R_*/2+ \Lambda_{\max}R}(x_\bullet)\times I_{\rm ph},\\
\mathcal Q^{\rm ph}_{2R}\subset B_{R_*/2+2\Lambda_{\max}R}(x_\bullet)\times I_{\rm ph},
\end{gathered}
\]
for every \(\tau\in I_*\) and \(x=X(\tau)+\Lambda(\tau)Y\) with \(Y\in B_{2R}\).
In particular, if \(|c(\tau)|\le \delta\) from Proposition~\ref{paper2:prop:preliminary-correction},
then \(X(\tau)\in B_{R_*/2+\Lambda_{\max}\delta}(x_\bullet)\) on \(I_*\).
Hence, \(\chi_R\in C_c^\infty(B_{2R})\) pulls back to
\(\chi_R^\tau:=\chi_R\circ \Phi^{-1}\in C_c^\infty(\mathcal Q^{\rm ph}_{2R})\)
with support depending only on \(R,\Lambda_{\min},\Lambda_{\max}\).
\end{proposition}
\begin{proof}
Fix \(\tau\in I_*\) and \(Y\in B_R\). Since \(x=X(\tau)+\Lambda(\tau)Y\),
\[
|x-X(\tau)|\le |Y|\,\Lambda_{\max}\le \Lambda_{\max}R,\qquad
\Lambda_{\max}^{-1}|x-X(\tau)|\le |Y|.
\]
Because \(X(\tau)\in B_{R_*/2}(x_\bullet)\), this gives
\[
x\in B_{R_*/2+\Lambda_{\max}R}(x_\bullet).
\]
Similarly, if \(Y\in B_{2R}\), then \(x\in B_{R_*/2+2\Lambda_{\max}R}(x_\bullet)\).
Together with \(t=t_{\rm c}(\tau)\in I_{\rm ph}\), these give the claimed
inclusions for \(\mathcal Q_{R}^{\rm ph}\) and \(\mathcal Q_{2R}^{\rm ph}\).

For pull-backs of cutoffs, \(\Phi\) is bi-Lipschitz at every \(\tau\), so composition
\[
\chi_R^\tau(x,t):=\chi_R\!\left(\frac{x-X(\tau)}{\Lambda(\tau)}\right)
\]
is smooth on \(\mathcal Q^{\rm ph}_{2R}\), supported where
\(Y\in B_{2R}\), hence supported where \((x,t)\in\mathcal Q^{\rm ph}_{2R}\). The support size depends
only on \(R,\Lambda_{\min},\Lambda_{\max}\) and the fixed \(C_c^\infty\) norm of
\(\chi_R\), not on the particular branch.
\end{proof}

\subsection{Transport of local quantities under the chart}\label{paper2:sec:local-transfer}

Let \(t=t(\tau)=t_{\rm c}(\tau)\) and, for \(R>0\), define
\[
J_\tau^{-}(R):=\bigl(t(\tau)-\Lambda_{\max}^{2}R^2,\ t(\tau)\bigr].
\]
Define renormalized local quantities on \(Q_R^*(\tau)\) by
\[
\begin{aligned}
A_V(\tau,R)&:=\sup_{s\in(\tau-R^2,\tau)}\int_{B_R}|V(\cdot,s)|^2\,dY,\qquad
E_V(\tau,R):=\int_{\tau-R^2}^\tau\!\!\int_{B_R}|\nabla V|^2\,dY\,d s,\\
C_V(\tau,R)&:=\int_{\tau-R^2}^\tau\!\!\int_{B_R}|V|^3\,dY\,d s,\qquad
D_V(\tau,R):=\int_{\tau-R^2}^\tau\!\!\int_{B_R}\Bigl|P-\langle P\rangle_{B_R}\Bigr|^{3/2}\,dY\,d s.
\end{aligned}
\]
For interval notation and a center \(X\), define
\[
A_u(J,R;X):=\sup_{s\in J}\int_{B_R(X)}|u(\cdot,s)|^2\,dx,\quad
E_u(J,R;X):=\int_J\int_{B_R(X)}|\nabla u|^2\,dx\,ds,
\]
\[
C_u(J,R;X):=\int_J\int_{B_R(X)}|u|^3\,dx\,ds,\quad
D_u(J,R;X):=\int_J\int_{B_R(X)}|p-\langle p\rangle_{B_R(X)}|^{3/2}\,dx\,ds.
\]
In the inequalities below we abbreviate
\[
\begin{aligned}
A_u(J,R)&:=A_u(J,R;X(\tau)),& E_u(J,R)&:=E_u(J,R;X(\tau)),\\
C_u(J,R)&:=C_u(J,R;X(\tau)),& D_u(J,R)&:=D_u(J,R;X(\tau)).
\end{aligned}
\]

\begin{proposition}[Cylinder comparison on compact windows]\label{paper2:prop:local-transfer}
Assume \(\Lambda_{\min}\le \Lambda\le \Lambda_{\max}\) and formulas of
Lemma~\ref{paper2:lem:change-id}.
Then for every \(\tau\in I_*\), \(R>0\),
\[
A_V(\tau,R)\le \Lambda_{\min}^{-1}A_u(J_\tau^{-}(R),\Lambda_{\max}R),
\]
\[
E_V(\tau,R)\le \Lambda_{\min}^{-3}E_u(J_\tau^{-}(R),\Lambda_{\max}R),
\]
\[
C_V(\tau,R)\le \Lambda_{\min}^{-2}C_u(J_\tau^{-}(R),\Lambda_{\max}R),
\]
\[
D_V(\tau,R)\le \Lambda_{\min}^{-2}D_u(J_\tau^{-}(R),\Lambda_{\max}R).
\]
In particular, all four renormalized quantities are bounded whenever the physical
quantities are bounded on the corresponding compact physical cylinders.
\begin{proof}
These are direct changes of variables under \(\Phi_{\rm pre}\):
\[
V(\cdot,\tau)=\Lambda(\tau)u(\cdot,t(\tau)),\qquad
P(\cdot,\tau)=\Lambda(\tau)^2 p(\cdot,t(\tau)),
\]
\[
\nabla_YV=\Lambda\nabla_xu,\qquad dY\,d\tau=\Lambda^{-5}\,dx\,dt,\qquad
d\tau=\Lambda^{2}\,dt.
\]
Insert these formulas into each definition and use
\(\Lambda_{\min}\le\Lambda\le\Lambda_{\max}\), plus
\[
\tau-s\in(0,R^2)\implies t(\tau)-\Lambda_{\max}^{2}R^2\le t(s)\le t(\tau)-\Lambda_{\min}^{2}R^2.
\]
\end{proof}
\end{proposition}

\subsection{Local repaired gauge map}\label{paper2:sec:gauge-map}

Fix \(R>0\), \(0<p<3\), and \(\chi_R\in C_c^\infty(B_{2R})\), \(\chi_R\equiv1\) on
\(B_R\). Let the local energy profile space be
\[
\mathcal V_R:=\left\{V\in L^3(B_{2R})\cap W^{1,3/2}(B_{2R};\RR^3):
|Y|^{-p/3}V\in L^3(B_{2R})\right\},
\]
with norm
\[
\|V\|_{\mathcal V_R}:=\|V\|_{L^3(B_{2R})}
+\||Y|^{-p/3}V\|_{L^3(B_{2R})}
+\|\nn V\|_{L^{3/2}(B_{2R})}.
\]
For differentiating the repaired scale and centering gauges we use the
admissible gauge class
\[
\mathcal A_{p,R}:=\left\{V:
\begin{array}{l}
V\in L^3(B_{4R})\cap W^{1,3}_{\rm loc}(B_{4R};\RR^3),\\[2mm]
\displaystyle \int_{B_{4R}}|Y|^{-p}|V|^3\,dY<\infty,\quad
\displaystyle \int_{B_{4R}}|Y|^{-p}|V|\,|Y\cdot\nabla V|\,dY<\infty,\\[2mm]
\displaystyle \int_{B_{4R}}|Y|^{-p-1}|V|^3\,dY<\infty
\end{array}\right\}.
\]
The larger ball is only a domain buffer: all gauge integrals are supported in
\(B_{2R}\), while \((\lambda,c)\) is restricted near \((1,0)\) so that
\(\lambda^{-1}(B_{2R}-c)\subset B_{4R}\).

\begin{definition}[Gauge actions]\label{paper2:def:gauge-action}
For \((\lambda,c)\in (0,\infty)\times\RR^3\), define
\[
V_{\lambda,c}(Y):=\lambda^{-1}V(\lambda^{-1}(Y-c)).
\]
For pressure profiles use
\[
P_{\lambda,c}(Y):=\lambda^{-2}P(\lambda^{-1}(Y-c)).
\]
We also write \(V^{\lambda,c}:=V_{\lambda,c}\) and \(P^{\lambda,c}:=P_{\lambda,c}\).
\end{definition}

\begin{definition}[Finite-dimensional gauge map]\label{paper2:def:gauge-map}
For \(V\in\mathcal V_R\) and \( (\lambda,c) \in (0,\infty)\times\RR^3\), set
\[
\mathcal F(\lambda,c;V)
:=\bigl(\mathcal F_0,\mathcal F_1,\mathcal F_2,\mathcal F_3\bigr),
\]
with
\[
\mathcal F_0(\lambda,c;V):=
\int_{\RR^3}\chi_R(Y)\,|Y|^{-p}|V_{\lambda,c}(Y)|^3\dd Y-\Theta_0,\qquad
\mathcal F_j(\lambda,c;V):=\int_{\RR^3}Y_j\,\chi_R(Y)\,|V_{\lambda,c}(Y)|^2\dd Y
\]
for \(j=1,2,3\). The gauge conditions are \(\mathcal F(\lambda,c;V)=0\).
Codomain is \(\RR^4\), domain \((0,\infty)\times\RR^3\times\mathcal V_R\).
\end{definition}

\begin{definition}[Localized modulation matrix]\label{paper2:def:gauge-nondeg}
For \(J\subset I_*\) and a profile path \(W:J\to\mathcal V_R\), define
\[
J_W(\tau):=
D_{(\lambda,c)}\mathcal F\bigl(1,0;W(\tau)\bigr)
\in\mathbb R^{4\times4}.
\]
The rows are ordered as the scale constraint followed by the three centering
constraints.
\end{definition}

\begin{lemma}[Weighted integrability and differentiation under integrals]\label{paper2:lem:weighted-finite}
Assume \(V,W\in\mathcal A_{p,R}\). Then \(G_0(V):=\int\chi_R|Y|^{-p}|V|^3\)
is finite and
\[
\frac{\dd}{\dd\epsilon}\Big|_{\epsilon=0}
\int \chi_R|Y|^{-p}|V+\epsilon W|^3\,\dd Y
=\int 3\chi_R|Y|^{-p}|V|\,V\cdot W\,\dd Y.
\]
\end{lemma}
\begin{proof}
For \(W\in\mathcal A_{p,R}\), define
\[
F_\varepsilon(Y):=\chi_R(Y)|Y|^{-p}|V(Y)+\varepsilon W(Y)|^3.
\]
For each \(Y\), \(\varepsilon\mapsto F_\varepsilon(Y)\) is \(C^1\) and
\[
\partial_\varepsilon F_\varepsilon
=3\chi_R|Y|^{-p}|V+\varepsilon W|(V+\varepsilon W)\cdot W.
\]
For \(|\varepsilon|\le1\),
\[
|F_\varepsilon(Y)|+|\partial_\varepsilon F_\varepsilon|
\le C\chi_R|Y|^{-p}\bigl(|V|^3+|W|^3\bigr)\in L^1(B_{2R})
\]
by the defining weighted \(L^3\) condition of \(\mathcal A_{p,R}\).
Hence dominated convergence allows differentiation under the integral, and
\[
\frac{\dd}{\dd\varepsilon}\Big|_{\varepsilon=0}\int_{B_{2R}}\chi_R|Y|^{-p}|V+\epsilon W|^3\,\dd Y
=\int_{B_{2R}}3\chi_R|Y|^{-p}|V|\,V\cdot W\,\dd Y.
\]
\end{proof}

\begin{lemma}[Gauge covariance]\label{paper2:lem:gauge-covariance}
For \(V\in\mathcal V_R\), \((\lambda,c)\in(0,\infty)\times\RR^3\), and
\(\alpha=0,\dots,3\),
\[
\mathcal F_\alpha(1,0;V_{\lambda,c})=\mathcal F_\alpha(\lambda,c;V).
\]
\begin{proof}
For each component this is a direct substitution:
\[
\mathcal F_0(1,0;V_{\lambda,c})
=\int\chi_R(Y)|Y|^{-p}|V_{\lambda,c}(Y)|^3-\Theta_0\,\dd Y
=\mathcal F_0(\lambda,c;V),
\]
\[
\mathcal F_j(1,0;V_{\lambda,c})
=\int Y_j\chi_R(Y)|V_{\lambda,c}(Y)|^2\,\dd Y
=\mathcal F_j(\lambda,c;V),\quad j=1,2,3.
\]
\end{proof}
\end{lemma}

\begin{proposition}[Differentiability of \(\mathcal F\)]\label{paper2:prop:gauge-C1}
For each fixed \(V\in\mathcal A_{p,R}\), the map
\((\lambda,c)\mapsto \mathcal F(\lambda,c;V)\) is \(C^1\) in a neighborhood of
\((1,0)\). In each component \( \mathcal F_\alpha\) all derivatives under the
integral sign are legitimate.
\begin{proof}
Fix \((\lambda,c)\) near \((1,0)\), and abbreviate
\[
V_{\lambda,c}(Y):=\lambda^{-1}V\!\left(\lambda^{-1}(Y-c)\right).
\]
By direct differentiation,
\[
\partial_\lambda V_{\lambda,c}
=-\lambda^{-1}V_{\lambda,c}-\lambda^{-1}(Y-c)\cdot\nabla V_{\lambda,c},\qquad
\partial_{c_k}V_{\lambda,c}=-\lambda^{-2}\partial_kV(\lambda^{-1}(Y-c))=-\partial_kV_{\lambda,c}.
\]
The neighborhood of \((1,0)\) is chosen so that the arguments of \(V\) remain
inside \(B_{4R}\) for \(Y\in\supp\chi_R\). The scale derivative is controlled by
the two weighted integrability conditions
\[
\int |Y|^{-p}|V|^3<\infty,\qquad
\int |Y|^{-p}|V|\,|Y\cdot\nabla V|<\infty,
\]
transported by the smooth change of variables. The translation derivative in
the scale component is treated by the source identity
\[
3\int |Y|^{-p}|V|\,V\cdot\partial_kV
=p\int Y_k|Y|^{-p-2}|V|^3,
\]
obtained first for smooth compactly supported profiles and then by approximation
in \(\mathcal A_{p,R}\); the right-hand side is finite because
\(\int |Y|^{-p-1}|V|^3<\infty\). The translation components are supported in \(B_{2R}\)
and are controlled by \(V\in L^2_{\rm loc}\) and \(\nabla V\in L^3_{\rm loc}\).
These bounds give an integrable majorant for the difference quotients, so
dominated convergence applies.
\[
\mathcal F_0(\lambda,c;V)
=\int_{\RR^3}\chi_R(Y)\,|Y|^{-p}|V_{\lambda,c}|^3\,\dd Y-\Theta_0,
\]
for which dominated convergence gives
\[
\partial_{\lambda}\mathcal F_0
=\int_{\RR^3}\chi_R|Y|^{-p}\,3|V_{\lambda,c}|\,V_{\lambda,c}\cdot\partial_\lambda V_{\lambda,c}\,dY,
\]
\[
\partial_{c_k}\mathcal F_0
=\int_{\RR^3}\chi_R|Y|^{-p}\,3|V_{\lambda,c}|\,V_{\lambda,c}\cdot\partial_{c_k}V_{\lambda,c}\,dY.
\]
Similarly
\[
\mathcal F_j(\lambda,c;V)=\int_{\RR^3}Y_j\chi_R(Y)\,|V_{\lambda,c}|^2\,\dd Y,\quad j=1,2,3,
\]
hence
\[
\partial_\lambda \mathcal F_j
\ =\int_{\RR^3}2Y_j\chi_R\,V_{\lambda,c}\cdot\partial_\lambda V_{\lambda,c}\,dY,
\]
\[
\partial_{c_k}\mathcal F_j
=\int_{\RR^3}2Y_j\chi_R\,V_{\lambda,c}\cdot\partial_{c_k}V_{\lambda,c}\,dY.
\]
Thus each component admits pointwise \(C^1\) formulas and these are continuous in
\((\lambda,c)\) near \((1,0)\). Therefore \((\lambda,c)\mapsto\mathcal F(\lambda,c;V)\in C^1\)
near \((1,0)\).
\end{proof}
\end{proposition}

\begin{proposition}[Full local Jacobian at \((1,0)\)]\label{paper2:prop:full-jacobian}
Let \(V\in\mathcal A_{p,R}\).
Set \(V_0(Y)=V(Y)\).
At \((\lambda,c)=(1,0)\), the matrix
\[
J(V):=D_{(\lambda,c)}\mathcal F(1,0;V)\in \RR^{4\times4}
\]
has rows
\[
J_\alpha=\bigl(\partial_\lambda \mathcal F_\alpha,\ \partial_{c_1}\mathcal F_\alpha,\
\partial_{c_2}\mathcal F_\alpha,\ \partial_{c_3}\mathcal F_\alpha\bigr),\quad \alpha=0,\dots,3.
\]
Writing \(W_{\rm scl}=V_0+Y\cdot\nn V_0\), one has
\begin{align}
\partial_\lambda \mathcal F_0(1,0;V)
&=
-3\int\chi_R|Y|^{-p}|V_0|\,V_0\cdot W_{\rm scl}\,\dd Y,\label{paper2:eq:jac00}\\
\partial_{c_k}\mathcal F_0(1,0;V)
&=-3\int_{\RR^3}\chi_R |Y|^{-p}|V_0|\,V_0\cdot\partial_kV_0\,\dd Y,\label{paper2:eq:jac0k}\\
\partial_\lambda \mathcal F_j(1,0;V)
&=-2\int_{\RR^3}Y_j\chi_R\,V_0\cdot W_{\rm scl}\,\dd Y,\label{paper2:eq:jacjlambda}\\
\partial_{c_k}\mathcal F_j(1,0;V)
&=-2\int_{\RR^3}\chi_R Y_j V_0\cdot\partial_kV_0\,\dd Y,\label{paper2:eq:jacjk}
\end{align}
for \(j,k=1,2,3\).
\begin{proof}
Differentiate the formulas in Definition~\ref{paper2:def:gauge-map} at \((\lambda,c)=(1,0)\).

Since \(V_{\lambda,c}(Y)=\lambda^{-1}V(\lambda^{-1}(Y-c))\),
\[
\partial_\lambda V_{\lambda,c}
 =-\lambda^{-1}V_{\lambda,c}-\lambda^{-1}Y\cdot\nabla V_{\lambda,c},\qquad
\partial_{c_k}V_{\lambda,c}=-\lambda^{-2}\partial_kV(\lambda^{-1}(Y-c)).
\]
At \((1,0)\), \(\partial_\lambda V_{\lambda,c}|_{(1,0)}=-(V+Y\cdot\nabla V)\),
\(\partial_{c_k}V_{\lambda,c}|_{(1,0)}=-\partial_kV\).

For \(\mathcal F_0\), write
\[
\mathcal F_0(\lambda,c;V)=\int \chi_R(Y)|Y|^{-p}|V_{\lambda,c}(Y)|^3\,dY-\Theta_0.
\]
Derivative in \(\lambda\):
\[
\partial_\lambda \mathcal F_0
\ =\int\chi_R|Y|^{-p}3|V_{\lambda,c}|\,V_{\lambda,c}\cdot\partial_\lambda V_{\lambda,c}\,dY
\]
Evaluating at \((1,0)\) yields \eqref{paper2:eq:jac00}.

For \(k=1,2,3\), \(c_k\)-derivative of \(\mathcal F_0\) has no cutoff term:
\[
\partial_{c_k}\mathcal F_0(1,0;V)
=\int \chi_R|Y|^{-p}\,3|V|\,V\cdot(-\partial_kV)\,dY,
\]
which is \eqref{paper2:eq:jac0k}.

For \(\mathcal F_j\),
\[
\mathcal F_j(\lambda,c;V)=\int Y_j\chi_R(Y)\,|V_{\lambda,c}|^2\,dY.
\]
Hence
\[
\partial_\lambda \mathcal F_j
=\int 2Y_j\chi_R\,V_{\lambda,c}\cdot\partial_\lambda V_{\lambda,c}\,dY,
\]
and at \((1,0)\):
\[
\partial_\lambda \mathcal F_j(1,0;V)
=-2\int Y_j\chi_R\,V\cdot W_{\rm scl}\,dY,
\]
which is \eqref{paper2:eq:jacjlambda}. Similarly,
\[
\partial_{c_k}\mathcal F_j(1,0;V)
=\int 2Y_j\chi_R\,V\cdot\partial_{c_k}V_{\lambda,c}\big|_{(1,0)}\,dY
=-2\int \chi_RY_jV\cdot\partial_kV\,dY,
\]
giving \eqref{paper2:eq:jacjk}.
\end{proof}
\end{proposition}

\begin{proposition}[Translation block and Schur-complement invertibility]\label{paper2:thm:quant-jac}
Let \(V:J\to\mathcal A_{p,R}\) satisfy the repaired gauge constraints
\[
\mathcal F(1,0;V(\tau))=0\qquad\text{for a.e. }\tau\in J.
\]
Write the modulation matrix in block form
\[
J(V(\tau))=
\begin{pmatrix}
\alpha(\tau)&u(\tau)^T\\
v(\tau)&A(\tau)
\end{pmatrix},
\]
where \(\alpha=\partial_\lambda\mathcal F_0\), \(u_\ell=\partial_{c_\ell}\mathcal F_0\),
\(v_j=\partial_\lambda\mathcal F_j\), and \(A_{j\ell}=\partial_{c_\ell}\mathcal F_j\).
Assume that for a.e. \(\tau\in J\)
\[
m_\chi(\tau):=\int |V(Y,\tau)|^2\chi_R(Y)\,dY\ge m_0>0,
\]
\[
\int_{A_R}|V(Y,\tau)|^2\,dY\le\varepsilon_0,\qquad
3C_\chi\varepsilon_0\le \frac{m_0}{2},
\]
where \(A_R=\{R\le |Y|\le2R\}\) and
\[
C_\chi:=\max_{j,\ell}\|Y_j\partial_\ell\chi_R\|_{L^\infty}.
\]
Assume also
\[
C_u:=\operatorname*{ess\,sup}_{\tau\in J}\int |Y|^{-p-1}|V(Y,\tau)|^3\,dY<\infty
\]
and, with \(C_\nabla:=\||Y|\nabla\chi_R\|_{L^\infty}\),
\[
6p\left(\frac{2}{m_0}\right)C_u\,R C_\nabla\varepsilon_0
\le \frac{p\Theta_0}{2}.
\]
Then \(J(V(\tau))\) is invertible for a.e. \(\tau\in J\), and
\[
\operatorname*{ess\,sup}_{\tau\in J}\|J(V(\tau))^{-1}\|_{\infty\to\infty}
\le K(R,p,\Theta_0,m_0,C_u,\varepsilon_0,\chi_R).
\]
\end{proposition}
\begin{proof}
The entries are finite by \(V(\tau)\in\mathcal A_{p,R}\) and the compact support
of \(\chi_R\). With the present parameter convention,
\(\partial_\lambda V_{\lambda,c}|_{(1,0)}=-(V+Y\cdot\nabla V)\). Hence the
scale entry satisfies the signed transversality identity
\[
\alpha(\tau)=-p\int |Y|^{-p}|V(Y,\tau)|^3\,dY=-p\Theta_0
\]
on the gauge surface. Thus \(|\alpha|\ge p\Theta_0\).

For the centering block, integration by parts gives
\[
A_{j\ell}
=-\delta_{j\ell}\int |V|^2\chi_R
-\int |V|^2Y_j\partial_\ell\chi_R.
\]
Writing \(A=-m_\chi I_3+E\), the annular support of \(\nabla\chi_R\) gives
\[
|E_{j\ell}|\le C_\chi\int_{A_R}|V|^2\le C_\chi\varepsilon_0,\qquad
\|E\|_{\infty\to\infty}\le3C_\chi\varepsilon_0\le m_0/2.
\]
Thus \(A=-m_\chi(I_3-m_\chi^{-1}E)\) and
\(\|m_\chi^{-1}E\|_{\infty\to\infty}\le1/2\), so \(A^{-1}\) exists and
\(\|A^{-1}\|_{\infty\to\infty}\le2/m_0\).

For the scale-to-translation block, the scale-constraint formula gives
\[
\left|u_\ell\right|
\le p\int |Y|^{-p-1}|V|^3\,dY
\le p\int |Y|^{-p-1}|V|^3\,dY,
\]
so \(\|u\|_\infty\le pC_u\). For the translation-to-scale block,
\[
v_j=-\int |V|^2Y_j(Y\cdot\nabla\chi_R)\,dY
\]
because the centering constraint cancels the interior term. Therefore
\[
\|v\|_\infty\le2R\,C_\nabla\varepsilon_0.
\]
The Schur complement \(S=\alpha-u^TA^{-1}v\) satisfies
\[
|u^TA^{-1}v|
\le 3\|u\|_\infty\|A^{-1}\|_{\infty\to\infty}\|v\|_\infty
\le 6p\left(\frac{2}{m_0}\right)C_uRC_\nabla\varepsilon_0
\le \frac{p\Theta_0}{2}.
\]
Since \(|\alpha|=p\Theta_0\), this gives \(|S|\ge p\Theta_0/2\).
The block inverse formula
\[
J^{-1}=
\begin{pmatrix}
S^{-1} & -S^{-1}u^TA^{-1}\\
-A^{-1}vS^{-1}&A^{-1}+A^{-1}vS^{-1}u^TA^{-1}
\end{pmatrix}
\]
then gives the displayed bound.
\end{proof}

\subsection{Repaired gauge by implicit inversion in time}\label{paper2:sec:gauge-solve}
\begin{lemma}[Time chain rule for preliminary profile]\label{paper2:lem:diff-F-pre}
For each
\(\alpha=0,\dots,3\), set \(G_\alpha(\tau):=\mathcal F_\alpha(1,0;V_{\rm pre}(\tau))\).
Then \(G_\alpha\in W^{1,1}_{\rm loc}(I_*)\) and for a.e. \(\tau\),
\[
G_\alpha'(\tau)=D_V\mathcal F_\alpha(1,0;V_{\rm pre}(\tau))[\partial_\tau V_{\rm pre}(\tau)].
\]
The right-hand side is obtained from the profile formulas of
Proposition~\ref{paper2:prop:gauge-C1} and is in \(L^1_{\rm loc}(I_*)\).
\begin{proof}
Let \(V_{\rm pre}^\varepsilon\) be a spatial and temporal Steklov mollification
inside a slightly larger compact cylinder.  For the smoothed profiles the
ordinary chain rule gives
\[
\frac d{d\tau}\mathcal F_\alpha(1,0;V_{\rm pre}^\varepsilon(\tau))
=D_V\mathcal F_\alpha(1,0;V_{\rm pre}^\varepsilon(\tau))
[\partial_\tau V_{\rm pre}^\varepsilon(\tau)] .
\]
The transformed equation expresses \(\partial_\tau V_{\rm pre}^\varepsilon\)
as the mollification of
\(\nu\Delta V_{\rm pre}-(V_{\rm pre}\cdot\nabla)V_{\rm pre}-\nabla P_{\rm pre}\)
plus the chart transport terms.  Pairing this distribution with the compactly
supported Gateaux gradients below is legitimate by the local suitable bounds,
the pressure bound, and the weighted \(L^3\) control in \textup{(SC4)}.  Passing
\(\varepsilon\downarrow0\) gives the asserted \(W^{1,1}_{\rm loc}\) identity.
For \(W\in\mathcal V_R\), these derivative formulas hold:
\[
\begin{aligned}
D_V\mathcal F_0(1,0;V)[W]&=\int_{\mathbb R^3}3\chi_R|Y|^{-p}|V|\,V\cdot W\,\dd Y,\\
D_V\mathcal F_j(1,0;V)[W]&=\int_{\mathbb R^3}2Y_j\chi_R\,V\cdot W\,\dd Y,\qquad j=1,2,3,
\end{aligned}
\]
where both right-hand sides are finite on \(\mathcal V_R\) by compact support, \(p<3\), and
\(V,W\in L^3\cap W^{1,3/2}_{\mathrm{loc}}\).
The right-hand side is finite by the explicit formulas above and by the distributional
equation paired with compactly supported test gradients; the local energy,
pressure, and weighted \(L^3\) bounds give an \(L^1_{\rm loc}(I_*)\) majorant.
\end{proof}
\end{lemma}
\begin{lemma}[AC implicit inversion with parameters]\label{paper2:lem:AC-IFT}
Let \(U\subset\RR^4\) be open and \(G:I_*\times U\to\RR^4\). Assume:
\begin{enumerate}
\item[(i)] for each \(s\in U\), \(\tau\mapsto G(\tau,s)\in W^{1,1}_{\rm loc}(I_*;\RR^4)\), and
  \(D_\tau G\in L^1_{\rm loc}(I_*;L^\infty(U;\RR^4))\);
\item[(ii)] for a.e. \(\tau\in I_*\), \(s\mapsto G(\tau,s)\) is \(C^1\) on \(U\), with \(D_sG\) continuous in \((\tau,s)\);
\item[(iii)] for some \((\tau_0,s_0)\in I_*\times U\), \(G(\tau_0,s_0)=0\), and there is \(\rho,C>0\) such that
\[
\sup_{(\tau,s)\in I_*\times B_\rho(s_0)}\|D_sG(\tau,s)^{-1}\|_{\mathcal L(\RR^4,\RR^4)}\le C.
\]
\end{enumerate}
Then after possibly shrinking \(I_*\) to a subinterval \(J\ni\tau_0\), there is a unique
map \(s\in AC(J;U)\), \(s(\tau_0)=s_0\), with
\[
G(\tau,s(\tau))=0\quad\text{a.e. in }J.
\]
Moreover
\[
s'(\tau)=-D_sG(\tau,s(\tau))^{-1}\,D_\tau G(\tau,s(\tau))\quad\text{a.e. on }J.
\]
\begin{proof}
By (ii), the matrix field \(D_sG\) is continuous in \(s\), so there is \(r_0\in(0,\rho)\) such that
\(\inf_{(\tau,s)\in I_*\times B_{r_0}(s_0)}|\det D_sG(\tau,s)|>0\) after possibly
shrinking \(I_*\); this keeps \(D_sG^{-1}\) bounded and continuous on \(B_{r_0}(s_0)\).

Set
\[
F(\tau,s):=-D_sG(\tau,s)^{-1}D_\tau G(\tau,s).
\]
Then \(F(\tau,\cdot)\) is continuous on \(B_{r_0}(s_0)\) for a.e. \(\tau\) and
\(F(\cdot,s)\in L^1_{\rm loc}(I_*;\mathbb R^4)\) by (i), hence \(F\) is a Carathéodory
right-hand side.  Since the gauge map is \(C^1\) in the finite variables, after
shrinking \(r_0\) there is \(L\in L^1(J)\) such that
\[
|F(\tau,s_1)-F(\tau,s_2)|\le L(\tau)|s_1-s_2|
\quad\text{for }s_1,s_2\in B_{r_0}(s_0).
\]
Choose \(\delta>0\) so small that
\[
\int_{\tau_0-\delta}^{\tau_0+\delta}\sup_{s\in B_{r_0}(s_0)}|F(\sigma,s)|\,d\sigma<r_0.
\]
On \(J=[\tau_0-\delta,\tau_0+\delta]\cap I_*\) define the Picard operator
\[
(\mathcal Ts)(\tau):=s_0+\int_{\tau_0}^{\tau}F(\sigma,s(\sigma))\,d\sigma
\]
on \(C(J;B_{r_0}(s_0))\).  The preceding bound keeps \(\mathcal T\) in the same
closed ball, and the Lipschitz estimate gives
\[
\|\mathcal T s_1-\mathcal T s_2\|_{C(J)}
\le \left(\int_J L(\sigma)\,d\sigma\right)\|s_1-s_2\|_{C(J)}.
\]
After decreasing \(\delta\) once more, the coefficient is \(<1\), so Banach's
fixed-point theorem gives a unique solution
\[
s(\tau)=s_0+\int_{\tau_0}^{\tau}F(\sigma,s(\sigma))\,d\sigma,\qquad
\tau\in J,
\]
which stays in \(B_{r_0}(s_0)\subset U\). This solution is absolutely continuous.

Define \(H(\tau):=G(\tau,s(\tau))\). Since \(G(\cdot,s(\cdot))\in W^{1,1}_{\rm loc}\) from (i),
the chain rule gives for a.e. \(\tau\in J\)
\[
H'(\tau)=D_\tau G(\tau,s(\tau))+D_sG(\tau,s(\tau))s'(\tau)
=D_\tau G(\tau,s(\tau))+D_sG(\tau,s(\tau))F(\tau,s(\tau))=0.
\]
Hence \(H\equiv G(\tau_0,s_0)=0\) on \(J\), which is the claimed constraint.
Uniqueness follows from Grönwall on the difference of two AC solutions on \(B_{r_0}(s_0)\),
using Lipschitz continuity of \(F(\tau,\cdot)\) on that ball.
\end{proof}
\end{lemma}

\begin{definition}[Retained nondegenerate repaired-gauge regime]\label{paper2:def:admissible-gauge-path}
Let \(J\subset I_*\) be compact. The retained nondegenerate repaired-gauge
regime over the preliminary profile \(V_{\rm pre}\) consists of branches for
which there are
\[
\lambda\in AC(J;(0,\infty)),\qquad c\in AC(J;\mathbb R^3),
\]
such that, for a.e. \(\tau\in J\),
\[
V(\tau):=V_{\rm pre}^{\lambda(\tau),c(\tau)}\in\mathcal A_{p,R},\qquad
\mathcal F(1,0;V(\tau))=0.
\]
The path is quantitatively nondegenerate if the translation-block, annular,
weighted first-moment, and Schur-complement bounds in
Proposition~\ref{paper2:thm:quant-jac} hold on \(J\).
If no such local AC selection exists, or if the quantitative nondegeneracy
hypotheses fail persistently, the branch lies in the gauge-degenerate,
multibubble, or cascade alternatives rather than in the retained regime.
\end{definition}

\begin{proposition}[Gauge regularity on the retained regime]\label{paper2:prop:preliminary-correction}
Let \((\lambda,c)\) be the repaired-gauge selection of the retained
nondegenerate regime on \(J\). Then
\[
\mathcal F\bigl(\lambda(\tau),c(\tau);V_{\rm pre}(\tau)\bigr)=0
\quad\text{for a.e. }\tau\in J,
\]
and \((\lambda,c)\in AC(J;\mathbb R^4)\) by definition. If the path is
quantitatively nondegenerate, then
\[
\operatorname*{ess\,sup}_{\tau\in J}\|J(V(\tau))^{-1}\|<\infty.
\]
\end{proposition}
\begin{proof}
The constraint follows from Lemma~\ref{paper2:lem:gauge-covariance} applied to
\(V(\tau)=V_{\rm pre}^{\lambda(\tau),c(\tau)}\). The absolute continuity is part
of Definition~\ref{paper2:def:admissible-gauge-path}. The uniform inverse bound is
Proposition~\ref{paper2:thm:quant-jac}.
\end{proof}

\begin{definition}[Repaired chart]\label{paper2:def:repaired-chart}
Define
\[
\Lambda(\tau):=\Lambda_{\rm pre}(\tau)\lambda(\tau),\qquad
X(\tau):=X_{\rm pre}(\tau)+\Lambda_{\rm pre}(\tau)c(\tau),
\]
\[
V(\tau):=V_{\rm pre}^{\lambda(\tau),c(\tau)},\qquad
P(\tau):=P_{\rm pre}^{\lambda(\tau),c(\tau)}\ \text{(same physical pressure pull-back under }\Phi).
\]
\end{definition}

\begin{theorem}[AC repaired parameters]\label{paper2:thm:repaired-existence}
Assume \(\mathrm{SC}(R_*,r_*,M,\eta,\kappa)\), and suppose the branch lies in the
retained nondegenerate repaired-gauge regime on \(J\), with selection
\((\lambda,c)\). Then the corrected
chart \((\Lambda,X)\) is absolutely continuous on \(J\), satisfies
\[
a(\tau)=-\Lambda'(\tau)\in L^1_{\rm loc}(J),\qquad
b(\tau)=-X'(\tau)\in L^1_{\rm loc}(J),
\]
and the gauge constraints \(\mathcal F(1,0;V(\tau))=0\) hold a.e.
Equivalently,
\[
\mathcal F(\lambda(\tau),c(\tau);V_{\rm pre}(\tau))=0
\]
for the repaired-gauge selection \((\lambda,c)\).
\end{theorem}
\begin{proof}
By Definition~\ref{paper2:def:admissible-gauge-path}, \((\lambda,c)\in AC(J)\) and the
repaired gauge constraints hold for \(V(\tau)\) a.e.
By definition \(V(\tau)=V_{\rm pre}^{\lambda(\tau),c(\tau)}\), so Lemma~\ref{paper2:lem:gauge-covariance}
applied at \((\lambda,c)=(\lambda(\tau),c(\tau))\) gives
\[
\mathcal F(1,0;V_{\rm pre}^{\lambda,c})=\mathcal F(\lambda,c;V_{\rm pre})
\]
implies \(\mathcal F(1,0;V(\tau))=0\) a.e.

Since \(\Lambda=\Lambda_{\rm pre}\lambda\) and \(X=X_{\rm pre}+\Lambda_{\rm pre}c\),
and \(\Lambda_{\rm pre},X_{\rm pre}\) are AC by \textup{(SC5)}, the products and sums
preserve AC, hence \((\Lambda,X)\in AC(I_*)\). Differentiating,
\[
a=-\Lambda' =-(\Lambda_{\rm pre}'\lambda+\Lambda_{\rm pre}\lambda')\in L^1_{\rm loc},\qquad
b=-X' =-(X_{\rm pre}'+\Lambda_{\rm pre}'c+\Lambda_{\rm pre}c')\in L^1_{\rm loc}.
\]
This gives the claimed regularity.
\end{proof}

\begin{remark}[Gauge degeneracy]
Failure of invertibility of \(J\) is the gauge-degenerate alternative to the
retained single-core regime for this fixed renormalization mechanism.  In the
Type II decomposition of alternatives this case belongs to the
multibubble/cascade and scale-rigidity alternatives, rather than to the
representation theorem.
\end{remark}

\subsection{Pressure transport and decomposition}\label{paper2:sec:pressure}
\begin{lemma}[Renormalized pressure equation]\label{paper2:lem:pressure-eq}
If \(V,P\) satisfy the renormalized Navier--Stokes equation in
Theorem~\ref{paper2:thm:ren-NS}, then
\[
-\Delta P=\partial_i\partial_j(V_iV_j)\qquad\text{in }\mathcal D'(Q_{2R}^*).
\]
\end{lemma}
\begin{proof}
Take test \(\varphi\in C_c^\infty(Q_{2R}^*)\). Since \(V,P\in L^2_{\rm loc}\times
L^{3/2}_{\rm loc}\) on \(Q_{2R}^*\), all pairings below are legitimate.
Expand
\[
\nn\cdot\bigl((V\cdot\nabla)V\bigr)
=\sum_{i,j}\partial_i\partial_j(V_iV_j)-\sum_{i,j}\partial_iV_i\,\partial_jV_j.
\]
Because \(\nn\cdot V=0\), this reduces to \(\partial_i\partial_j(V_iV_j)\).
Taking divergence of
\[
\partial_\tau V+(V\cdot\nabla)V+\nabla P-\nu\Delta V
+a(V+Y\cdot\nabla V)+b\cdot\nabla V=0
\]
and using \(\nn\cdot V=0\), \(\nn\cdot(V+Y\cdot\nabla V)=0\), \(\nn\cdot(b\cdot\nabla V)=b\cdot\nn(\nn\cdot V)=0\),
and \(\nn\cdot\Delta V=\Delta(\nn\cdot V)=0\), the following identity holds in \(\mathcal D'\):
\[
\Delta P+\partial_i\partial_j(V_iV_j)=0.
\]
Hence \(-\Delta P=\partial_i\partial_j(V_iV_j)\) on \(Q_{2R}^*\).
\end{proof}

\begin{theorem}[Localized pressure decomposition]\label{paper2:thm:pressure-decomp}
Let \(\zeta\in C_c^\infty(B_{2R})\), \(\zeta\equiv1\) on \(B_R\), and assume
\(V\in L^1_{\mathrm{loc}}(I_*;L^3(B_{2R}))\), \(P\in L^1_{\mathrm{loc}}(I_*;L^{3/2}(B_{2R}))\).
For a.e. \(\tau\in I_*\) there exists
\[
P(\cdot,\tau)=P_{\rm loc}(\cdot,\tau)+H(\cdot,\tau)
\]
on \(B_R\) such that
\[
\begin{cases}
-\Delta P_{\rm loc}(\cdot,\tau)=\partial_i\partial_j\bigl(\zeta V_iV_j\bigr),\\
-\Delta H(\cdot,\tau)=0\ \text{on }B_R,\quad
\fint_{B_R}H(\cdot,\tau)\,dY=0.
\end{cases}
\]
Moreover
\[
\|P_{\rm loc}(\tau)\|_{L^{3/2}(B_R)}
\lesssim \|V(\tau)\|_{L^3(B_{2R})}^2,\qquad
\|H(\tau)\|_{L^q(B_r)}\lesssim_{q,R,r}
\|P(\tau)\|_{L^{3/2}(B_{2R})}+\|V(\tau)\|_{L^3(B_{2R})}^2
\]
for each \(1\le q<\infty\) and each \(0<r<R\), with constants depending only on \(R\), \(r\), and \(q\).
\begin{proof}
For fixed \(\tau\), set
\[
P_{\rm loc}= \mathcal R_i\mathcal R_j(\zeta V_iV_j),\qquad
H:=P-P_{\rm loc}.
\]
\(\zeta V_iV_j\in L^{1}_{\mathrm{loc}}(I_*;L^{3/2}(B_{2R}))\), so the Calderón--Zygmund
operator \(f\mapsto \mathcal R_i\mathcal R_j f\) gives
\[
P_{\rm loc}(\cdot,\tau)=\mathcal R_i\mathcal R_j(\zeta V_iV_j)(\cdot,\tau),
\qquad
\|P_{\rm loc}(\tau)\|_{L^{3/2}(B_R)}
\lesssim \|V(\tau)\|_{L^3(B_{2R})}^2.
\]
Because \(\mathcal R_i\mathcal R_j\) is linear bounded on \(L^{3/2}\),
\[
P_{\rm loc}\in L^{1}_{\mathrm{loc}}(I_*;L^{3/2}(B_R)).
\]
Since \(-\Delta P_{\rm loc}=\partial_i\partial_j(\zeta V_iV_j)\) and \(\zeta\equiv1\) on \(B_R\),
 \[
-\Delta H(\cdot,\tau)
=\partial_i\partial_j\bigl((1-\zeta)V_iV_j\bigr)=0 \quad\text{in }B_R,
\]
so \(H(\cdot,\tau)\) is harmonic on \(B_R\). Imposing \(\fint_{B_R}H=0\) fixes the
time-dependent additive constant uniquely in each time slice.
For \(q\in[1,\infty)\) and \(0<r<R\), harmonicity and local
\(L^p\)-to-\(L^q\) interior estimates imply
\[
\|H(\tau)\|_{L^q(B_r)}\lesssim_{q,R,r}\|H(\tau)\|_{L^{3/2}(B_R)}.
\]
Finally, \(\|H(\tau)\|_{L^{3/2}(B_R)}\le \|P(\tau)\|_{L^{3/2}(B_R)}+\|P_{\rm loc}(\tau)\|_{L^{3/2}(B_R)}\),
giving the announced bound in terms of \(P\) and \(V\).
\end{proof}
\end{theorem}
\begin{remark}[Measurability and normalization]\label{paper2:rem:pressure-time}
If \(V,P\) are measurable in \(\tau\), then \(\tau\mapsto P_{\rm loc}(\cdot,\tau)\)
is measurable in \(L^{3/2}(B_R)\) by boundedness of \(\mathcal R_i\mathcal R_j\), hence
\(\tau\mapsto H(\cdot,\tau)\) is measurable in each \(L^q(B_r)\), \(0<r<R\).
The normalization \(\fint_{B_R}H(\cdot,\tau)\,dY=0\) fixes the time-dependent additive constant
in \(H\) uniquely for each \(\tau\), and therefore the decomposition is a legitimate
renormalized pressure splitting.
\end{remark}

\subsection{Modulation from differentiated constraints}\label{paper2:sec:modulation}
\begin{lemma}[Differentiation of constraints]\label{paper2:lem:diff-f}
For the repaired profile \(V\) constructed above, with
\(|Y|^{-p}|V|^3\in L^1_{\rm loc}(I_*;L^1)\),
Then \(G_\alpha(\tau):=\mathcal F_\alpha(1,0;V(\tau))\) are absolutely
continuous and
\[
\frac{\dd}{\dd\tau}G_\alpha(\tau)
=\int_{\RR^3}\mathbf D_\alpha(V,\nn V,\nabla\chi_R)\,\partial_\tau V\,\dd Y
\]
for a.e. \(\tau\). Here \(\mathbf D_\alpha\) is the \(L^1\)-integrable Gateaux
gradient associated to \(\mathcal F_\alpha\).
\begin{proof}
\[
G_\alpha(\tau)=\mathcal F_\alpha(1,0;V(\tau))
\]
is first evaluated on spatial-temporal mollifications \(V^\varepsilon\).  For
those smooth profiles the ordinary chain rule gives
\[
\frac d{d\tau}G_\alpha^\varepsilon(\tau)
=D_V\mathcal F_\alpha(1,0;V^\varepsilon(\tau))\,[\partial_\tau V^\varepsilon(\tau)].
\]
The local suitable bounds, pressure estimates, and compact support of the
Gateaux gradients below allow passage to the limit \(\varepsilon\downarrow0\),
which gives the displayed derivative formula for \(V\).
For \(\mathcal F_0\),
\[
\mathbf D_0(V,\nabla V,\nabla\chi_R)
=3\chi_R|Y|^{-p}|V|\,V;
\]
for \(\alpha=1,2,3\),
\[
\mathbf D_\alpha(V,\nabla V,\nabla\chi_R)
=2\chi_R\,Y_\alpha V.
\]
Each coefficient is \(L^1(B_{2R})\): \(3\chi_R|Y|^{-p}|V|\,V\in
L^{3/2}(B_{2R})\subset L^1(B_{2R})\) from \(p<3\) and \(|Y|^{-p}|V|^3\in L^1\), and
\(2\chi_RY_\alpha V\in L^2(B_{2R})\subset L^1(B_{2R})\) by compact support.
\[
\frac d{d\tau}G_\alpha(\tau)
=\int_{\RR^3}\mathbf D_\alpha(V,\nn V,\nabla\chi_R)\cdot\partial_\tau V\,\dd Y,
    \]
Since \(\mathbf D_\alpha(V,\nn V,\nabla\chi_R)\in L^1(B_{2R})\) and
\(\partial_\tau V\in L^1_{\rm loc}(I_*;L^2(B_{2R}))\),
the right-hand side is in \(L^1_{\rm loc}(I_*)\), so \(G_\alpha\in W^{1,1}_{\rm loc}(I_*)\subset AC_{\rm loc}(I_*)\).
\end{proof}
\end{lemma}

\begin{lemma}[Directional derivatives in the profile variable]\label{paper2:lem:profile-direction}
For \(\mathcal A_\tau=V+Y\cdot\nn V\) and \(\alpha=0,1,2,3\),
\[
D_V\mathcal F_\alpha(1,0;V)[\mathcal A_\tau]
=-\partial_\lambda\mathcal F_\alpha(1,0;V),\qquad
D_V\mathcal F_\alpha(1,0;V)[\partial_kV]
=-\partial_{c_k}\mathcal F_\alpha(1,0;V).
\]
Indeed, from \(V_{\lambda,c}(Y)=\lambda^{-1}V(\lambda^{-1}(Y-c))\),
\[
\partial_\lambda V_{\lambda,c}\big|_{(1,0)}=-(V+Y\cdot\nabla V)=-\mathcal A_\tau,\qquad
\partial_{c_k}V_{\lambda,c}\big|_{(1,0)}=-\partial_kV.
\]
Substituting these into the Gateaux formulas from Proposition~\ref{paper2:prop:gauge-C1} gives the two identities.
\end{lemma}

\begin{theorem}[Modulation system]\label{paper2:thm:mod-sys}
Let \((V,P)\) be the repaired pair on the compact window under consideration.
Writing \(\theta:=(a,b)\), define
\[
\mathcal A_\tau(Y):=V(\tau,Y)+Y\cdot\nn V(\tau,Y).
\]
For \(\alpha=0,1,2,3\), define
\[
\mathcal R_\alpha(\tau):=
-\int_{\RR^3}\mathbf D_\alpha\bigl(V,\nn V,\nabla\chi_R\bigr)
\cdot\bigl((V\cdot\nn)V+\nabla P-\nu\Delta V\bigr)\,\dd Y.
\]
Then for a.e. \(\tau\in I_*\),
\[
\partial_\lambda\mathcal F_\alpha(1,0;V)\,a(\tau)
+\sum_{k=1}^3\partial_{c_k}\mathcal F_\alpha(1,0;V)\,b_k(\tau)
=\mathcal R_\alpha(\tau),
\]
Equivalently,
\[
J(\tau)\theta(\tau)=\mathcal R(\tau).
\]
Hence
\[
\theta(\tau)=J(\tau)^{-1}\mathcal R(\tau)\in L^1_{\rm loc}(I_*)^4,
\qquad |\theta(\tau)|\le \kappa\,|\mathcal R(\tau)|.
\]
\begin{proof}
Set \(G_\alpha(\tau):=\mathcal F_\alpha(1,0;V(\tau))\). By Theorem~\ref{paper2:thm:repaired-existence},
\(G_\alpha(\tau)=0\) a.e.; by Lemma~\ref{paper2:lem:diff-f}, \(G_\alpha\in W^{1,1}_{\rm loc}\).
Hence for a.e. \(\tau\),
\[
0=\frac{d}{d\tau}G_\alpha(\tau)
=D_V\mathcal F_\alpha(1,0;V(\tau))\,[\partial_\tau V(\tau)].
\]
Insert the equation from Theorem~\ref{paper2:thm:ren-NS},
\[
\partial_\tau V=-(V\cdot\nn)V-\nabla P+\nu\Delta V-a\mathcal A_\tau-b\cdot\nn V.
\]
This gives
\[
0=-D_V\mathcal F_\alpha\bigl[(V\cdot\nn)V+\nabla P-\nu\Delta V\bigr]
-a\,D_V\mathcal F_\alpha[\mathcal A_\tau]
-\sum_{k=1}^3b_k\,D_V\mathcal F_\alpha[\partial_kV].
\]
By Lemma~\ref{paper2:lem:diff-f}, for any profile direction \(W\),
\[
D_V\mathcal F_\alpha(1,0;V)[W]=\int_{\RR^3}\mathbf D_\alpha(V,\nn V,\nabla\chi_R)\cdot W\,\dd Y.
\]
By Lemma~\ref{paper2:lem:profile-direction},
\[
D_V\mathcal F_\alpha[\mathcal A_\tau]=-\partial_\lambda\mathcal F_\alpha(1,0;V),\qquad
D_V\mathcal F_\alpha[\partial_kV]=-\partial_{c_k}\mathcal F_\alpha(1,0;V).
\]
Substituting into the previous line yields the linear system
\[
\partial_\lambda\mathcal F_\alpha(1,0;V)\,a(\tau)
+\sum_{k=1}^3\partial_{c_k}\mathcal F_\alpha(1,0;V)\,b_k(\tau)
=\mathcal R_\alpha(\tau),
\]
which is equivalent to \(J(\tau)\theta(\tau)=\mathcal R(\tau)\).
By definition of \(\mathbf D_\alpha\), each \(\mathcal R_\alpha\in L^1_{\rm loc}(I_*)\), and
Proposition~\ref{paper2:thm:quant-jac} gives \(|J^{-1}(\tau)|\le K\) a.e., hence
\[\theta(\tau)=J^{-1}(\tau)\mathcal R(\tau),\quad
\theta\in L^1_{\rm loc}(I_*),\quad |\theta|\le K|\mathcal R|\text{ a.e.}.\]
\end{proof}
\end{theorem}

\begin{corollary}[Regularity of \(a,b\)]
For the repaired pair of Theorem~\ref{paper2:thm:mod-sys}, \(a,b\in L^1_{\rm loc}(I_*)\).
If in addition \(\mathcal R\in L^\infty_{\rm loc}(I_*)\), then \(a,b\in
L^\infty_{\rm loc}(I_*)\).
\begin{proof}
By Theorem~\ref{paper2:thm:mod-sys}, \(\theta=(a,b)=J^{-1}\mathcal R\) a.e., and
\(J^{-1}\in L^\infty_{\rm loc}(I_*)\) by Theorem~\ref{paper2:thm:quant-jac}. The first claim is therefore
immediate from \(\mathcal R\in L^1_{\rm loc}\), while the second follows if \(\mathcal R\in L^\infty_{\rm loc}\).
\end{proof}
\end{corollary}

\subsection{Repaired suitable structure in renormalized variables}\label{paper2:sec:suitable}


\begin{proposition}[Transport of suitable local energy inequality]\label{paper2:prop:ren-suitable}
Suppose \((u,p)\) is suitable on \(\mathcal Q^{\rm ph}_{2R}\), and
\((V,P)\) is obtained by the repaired map \(\Phi\).
Set
\[
\tilde a(\tau):=-\Lambda(\tau)^{-1}\Lambda'(\tau),\qquad
\tilde b(\tau):=-\Lambda(\tau)^{-1}X'(\tau).
\]
Then for every nonnegative \(\phi\in C_c^\infty(Q_R^*)\) and \(\tau_1<\tau_2\),
\[
\begin{aligned}
\mathscr E_\phi(V,P)
&:=\frac12|V|^2\partial_\tau(\phi^2)
+\frac{\nu}{2}|V|^2\Delta_Y\phi^2
+\frac12|V|^2\,V\cdot\nn(\phi^2)\\
&\quad +(P_{\rm loc}+H)\,V\cdot\nn(\phi^2)
+\frac{\tilde a}{2}|V|^2\phi^2
+\frac{\tilde a}{2}|V|^2\,Y\cdot\nn(\phi^2)
+\frac12|V|^2\,\tilde b\cdot\nn(\phi^2).
\end{aligned}
\]
\begin{multline*}
\frac12\int_{B_R}|V(\tau_2)|^2\phi(\cdot,\tau_2)^2
+\nu\int_{\tau_1}^{\tau_2}\!\!\int_{B_R}|\nn V|^2\phi^2\,dY\,d\tau
\le \frac12\int_{B_R}|V(\tau_1)|^2\phi(\cdot,\tau_1)^2\\
 +\int_{\tau_1}^{\tau_2}\!\!\int_{B_R}\mathscr E_\phi(V,P)\,dY\,d\tau.
\end{multline*}
The right-hand side is controlled on compact windows by
\(\|\tilde a\|_{L^1},\|\tilde b\|_{L^1},M\), and the local bounds from
Theorem~\ref{paper2:thm:pressure-decomp}, and support bounds from Section~\ref{paper2:sec:domains}.
Hence \((V,P)\) is locally suitable on \(Q_R^*\).
\begin{proof}
Apply the physical local energy inequality on \(\mathcal Q^{\rm ph}_{2R}\): for
nonnegative \(\psi\in C_c^\infty(\mathcal Q^{\rm ph}_{2R})\),
    \[
    \frac12\int |u|^2\psi\big|_{t_1}^{t_2}
    +\nu\int_{t_1}^{t_2}\!\!\int |\nabla_xu|^2\psi\,dxdt
    \le\int_{t_1}^{t_2}\!\!\int\Bigl(
    \frac12|u|^2(\partial_t\psi+\nu\Delta_x\psi)
    +\left(\frac{|u|^2}{2}u+p\,u\right)\!\cdot\!\nabla_x\psi\Bigr)\,dxdt.
    \]
Choose
    \[
    \psi(x,t):=\Lambda(t)^{-1}\phi\!\left(\frac{x-X(t)}{\Lambda(t)},\tau(t)\right)^2,\qquad
    \tau=\tau(t),\ d\tau/dt=\Lambda^{-2},
    \]
let \(t_i=t(\tau_i)\), \(i=1,2\).
    By Proposition~\ref{paper2:prop:prelim-frame} and Section~\ref{paper2:sec:domains}, \(\psi\in
    C_c^\infty(\mathcal Q^{\rm ph}_{2R})\).
\[
\nabla_x\psi=\Lambda^{-2}\nabla_Y\phi^2,\quad
\Delta_x\psi=\Lambda^{-3}\Delta_Y\phi^2,\quad
\partial_t\psi
=-\Lambda^{-4}\Lambda'\phi^2+\Lambda^{-3}\partial_\tau(\phi^2)-\Lambda^{-4}(X'+\Lambda'Y)\cdot\nabla_Y\phi^2.
\]
Using \(u=\Lambda^{-1}V\), \(p=\Lambda^{-2}P\), \(\nabla_xu=\Lambda^{-2}\nabla_YV\),
\(dx\,dt=\Lambda^{5}\,dY\,d\tau\), each term in the physical inequality becomes
\[
\frac12\int |u|^2\psi\big|_{t_1}^{t_2}
=\frac12\int |V|^2\phi^2\big|_{\tau_1}^{\tau_2},
\]
\[
\nu\int |\nabla_xu|^2\psi\,dxdt
=\nu\int |\nabla V|^2\phi^2\,dY\,d\tau,
\]
\[
\frac12\int |u|^2\partial_t\psi\,dxdt
=\frac12\int |V|^2\left(\partial_\tau(\phi^2)+\tilde a\phi^2+\tilde a\,Y\cdot\nabla_Y\phi^2+\tilde b\cdot\nabla_Y\phi^2\right)\,dY\,d\tau,
\]
\[
\frac{\nu}{2}\int |u|^2\Delta_x\psi\,dxdt
=\frac{\nu}{2}\int |V|^2\Delta_Y\phi^2\,dY\,d\tau,
\]
\[
\int\left(\frac{|u|^2}{2}u+p\,u\right)\cdot\nabla_x\psi\,dxdt
=\int\left(\frac12|V|^2V+(P_{\rm loc}+H)V\right)\cdot\nabla_Y\phi^2\,dY\,d\tau.
\]
Substituting these identities and grouping terms gives the proposition's right-hand
side and coefficients \(\tilde a,\tilde b\).
\end{proof}
\end{proposition}
\subsection{Stability under subsequences}\label{paper2:sec:stability}
\begin{proposition}\label{paper2:prop:stability}
Let \((u^{(n)},p^{(n)})\) satisfy \(\mathrm{SC}(R_*,r_*,M,\eta,\kappa)\) with the same
data parameters and let \((\Lambda^{(n)},X^{(n)},V^{(n)},P^{(n)},a^{(n)},b^{(n)})\)
be obtained from retained quantitatively nondegenerate repaired-gauge selections on the same \(I_*\),
with the constants in Proposition~\ref{paper2:thm:quant-jac} uniform in \(n\). After extraction \(n_j\to\infty\), there exist
subsequences and limits \((\Lambda,X,V,P,a,b)\) such that
\[
\Lambda^{(n_j)}\to\Lambda\ \text{in }L^1_{\rm loc},\quad
X^{(n_j)}\to X\ \text{in }L^1_{\rm loc},
\]
\[
\Lambda^{(n_j)\prime}\rightharpoonup\Lambda'\ \text{in }L^1_{\rm loc},\quad
X^{(n_j)\prime}\rightharpoonup X'\ \text{in }L^1_{\rm loc},
\]
\[
V^{(n_j)}\to V\ \text{in }L^2_{\rm loc}(I_*;L^2_{\rm loc}),\quad
V^{(n_j)}\rightharpoonup V\ \text{in }L^2_{\rm loc}(I_*;H^1_{\rm loc}),
\]
    \[
    P^{(n_j)}\rightharpoonup P\ \text{in }L^{3/2}_{\rm loc}(Q_R^*),\quad
    (a^{(n_j)},b^{(n_j)})\rightharpoonup(a,b)\ \text{in }L^1_{\rm loc}(I_*)^4.
    \]
    Moreover \(a=-\Lambda'\), \(b=-X'\) a.e.\ on \(I_*\), and the limit obeys all conclusions below.
\begin{proof}
From \(\mathrm{SC}\) and the retained repaired-gauge selections, \((V^{(n)},P^{(n)})\) are uniformly bounded in
\(L^\infty_\tau L^2(B_{2R})\cap L^2_\tau H^1(B_{2R})\) and \(L^{3/2}_{\rm loc}(Q_R^*)\).
Theorem~\ref{paper2:thm:quant-jac} gives uniform invertibility bounds for \(J^{(n)}\), and
Theorem~\ref{paper2:thm:mod-sys} gives \((a^{(n)},b^{(n)})\) uniformly in \(L^1_{\rm loc}\) (since
\(\mathcal R^{(n)}\in L^1_{\rm loc}\)). Hence \(\Lambda^{(n)}\) and \(X^{(n)}\) are uniformly
bounded in \(W^{1,1}_{\rm loc}\), so by BV compactness and 1D subsequences,
\[
\Lambda^{(n_j)}\to\Lambda\ \text{in }L^1_{\rm loc},\qquad
X^{(n_j)}\to X\ \text{in }L^1_{\rm loc},
\]
and \(a^{(n_j)}\rightharpoonup -\Lambda'\), \(b^{(n_j)}\rightharpoonup -X'\) in \(L^1_{\rm loc}\).
From the renormalized equation,
\[
\partial_\tau V^{(n)}\in
L^1_\tau(H^{-1}_{\rm loc})+L^{4/3}_\tau(L^{12/11}_{\rm loc}),
\]
uniformly on compact windows. Aubin--Lions compactness on bounded balls gives
\[
V^{(n_j)}\to V\ \text{in }L^2_{\rm loc}(I_*;L^2_{\rm loc}),\quad
V^{(n_j)}\rightharpoonup V\ \text{in }L^2_{\rm loc}(I_*;H^1_{\rm loc}).
\]
By Calderón--Zygmund and Theorem~\ref{paper2:thm:pressure-decomp}, \(P^{(n_j)}\rightharpoonup P\) in
\(L^{3/2}_{\rm loc}(Q_R^*)\), and with normalized pressure split \(P=P_{\rm loc}+H\),
the same for \(P_{\rm loc}\) and \(H\). Passing to limits in
\(\mathcal F(1,0;V^{(n)}(\tau))=0\), Theorem~\ref{paper2:thm:ren-NS}, and the linear modulation system
yields the corresponding limit constraints and equations.

\end{proof}
\end{proposition}

\subsection{Final representation theorem}\label{paper2:sec:final}

\begin{theorem}[Chart and repaired-gauge representation]\label{paper2:thm:A}
Assume \(\mathrm{SC}(R_*,r_*,M,\eta,\kappa)\), and suppose the branch lies in the
retained repaired-gauge regime with selection \((\lambda,c)\) on
\(J=[\tau_0-\ell,\tau_0+\ell]\subset I_*\). Then there are fields \((V,P)\) on \(Q_{2R}^J\), with
\((\Lambda,X)\in AC(J),\ a=-\Lambda',\ b=-X'\in L^1_{\rm loc}(J)\), such that
\[
u(x,t)=\Lambda(\tau)^{-1}V\!\left(\frac{x-X(\tau)}{\Lambda(\tau)},\tau\right),\quad
\frac{dt}{d\tau}=\Lambda(\tau)^2,
\]
and \(\mathcal F(1,0;V(\tau))=0\) for \(\tau\in J\).
\begin{proof}
From \textup{(SC5)} and Proposition~\ref{paper2:prop:prelim-frame}, \((\Lambda_{\rm pre},X_{\rm pre})\)
are AC on \(I_*\), so \(V_{\rm pre}\) is defined on \(Q_{2R}^*\).
The retained repaired-gauge regime supplies
\((\lambda,c)\in AC(J;(0,\infty)\times\mathbb R^3)\) and
\(\mathcal F(1,0;V_{\rm pre}^{\lambda,c})=0\) a.e.
By Definition~\ref{paper2:def:repaired-chart} and Theorem~\ref{paper2:thm:repaired-existence}, setting
\[
\Lambda(\tau)=\Lambda_{\rm pre}(\tau)\lambda(\tau),\qquad
X(\tau)=X_{\rm pre}(\tau)+\Lambda_{\rm pre}(\tau)c(\tau),\qquad
V(\tau)=V_{\rm pre}^{\lambda(\tau),c(\tau)},
\]
and \(P(\tau)=P_{\rm pre}^{\lambda(\tau),c(\tau)}\),
gives \( \Lambda,X\in AC(J)\) and \(a=-\Lambda',\ b=-X'\in L^1_{\rm loc}(J)\), and
\[
\mathcal F(1,0;V(\tau))=0\ \text{for a.e. }\tau\in J.
\]
Finally \(u=\Lambda^{-1}V(\cdot,\tau)\circ\Phi^{-1}\) by definition of the repaired map, hence the
physical identity with \(\frac{dt}{d\tau}=\Lambda^2\) follows.
\end{proof}
\end{theorem}

\begin{theorem}[Pressure and modulation]\label{paper2:thm:B}
Assume, in addition to \autoref{paper2:thm:A}, that the retained repaired-gauge
selection is quantitatively nondegenerate in the sense of Proposition~\ref{paper2:thm:quant-jac}.
Then \((V,P)\) satisfies the renormalized Navier--Stokes
equation with the sign convention of Theorem~\ref{paper2:thm:ren-NS}; there is a
decomposition \(P=P_{\rm loc}+H\) on \(Q_R^*\) as in
Theorem~\ref{paper2:thm:pressure-decomp}; and \((a,b)\) is the unique solution of
\[
J(\tau)\binom{a(\tau)}{b(\tau)}=\mathcal R(\tau),\qquad
\binom{a}{b}\in L^1_{\rm loc}(J)^4,
\]
with bounds \(|(a,b)|\le K|\mathcal R|\), where \(K\) depends on
\(R,p,\Theta_0,m_0,C_u,\varepsilon_0,\chi_R\).
\begin{proof}
Equation and sign convention are inherited directly from Theorem~\ref{paper2:thm:ren-NS} applied to
\((\Lambda,X)\) from Theorem~\ref{paper2:thm:A} and the definition of \(V\).
By Theorem~\ref{paper2:thm:pressure-decomp} applied at each time slice, the decomposition on \(Q_R^J\)
is defined and unique once \(\fint_{B_R}H=0\) is imposed.
The modulation identity is Theorem~\ref{paper2:thm:mod-sys}; uniqueness follows from invertibility of
\(J\) (Theorem~\ref{paper2:thm:quant-jac}) and the pointwise identity \((a,b)=J^{-1}\mathcal R\).
Uniform bounds on \((a,b)\) are obtained by the bound on \(J^{-1}\) there and the stated bound on
\(\mathcal R\).
\end{proof}
\end{theorem}

\begin{theorem}[Compact-window representation theorem]\label{paper2:thm:C}
Under the assumptions of \cref{paper2:thm:A,paper2:thm:B}, for every compact
window \(J\subset I_*\) and each fixed \(R>0\), setting
\(Q_R^J:=B_R\times J\),
one has:
\begin{enumerate}
\item the exact renormalized equation on \(Q_{2R}^J\),
\item the four repaired-gauge constraints \(\mathcal F(1,0;V(\tau))=0\),
\item bounds on \(a,b\in L^1_{\rm loc}(J)\) (and \(L^\infty_{\rm loc}\) under
  \(\mathcal R\in L^\infty_{\rm loc}\)),
\item local pressure decomposition \(P=P_{\rm loc}+H\) with \(L^{3/2}\) local bounds and
  harmonic remainder controls,
	\item renormalized local suitable energy inequality on \(Q_R^J\),
	\item stability under subsequences in the topologies above.
\end{enumerate}
The later reductions use only these conclusions.
\begin{proof}
From Theorem~\ref{paper2:thm:A} and the choice of \(J\subset I_*\), all quantities are defined on \(Q_{2R}^J\).
Items (1) and (2) are Theorems~\ref{paper2:thm:ren-NS} and~\ref{paper2:thm:repaired-existence};
item (3) is Theorem~\ref{paper2:thm:mod-sys} together with the inverse bound from
Proposition~\ref{paper2:thm:quant-jac}; item (4) is Theorem~\ref{paper2:thm:pressure-decomp};
item (5) is Proposition~\ref{paper2:prop:ren-suitable}; item (6) is
Proposition~\ref{paper2:prop:stability}. Combining these results gives the
stated compact-window representation.
\end{proof}
\end{theorem}

\begin{remark}
The preceding results separate two alternatives:
\begin{enumerate}
\item the nondegenerate single-core repaired-gauge regime, where the representation theorems apply;
\item gauge degeneracy (\(\det J=0\)) or compact-structure failure, which belongs
to the multibubble, cascade, or scale-rigidity alternatives.
\end{enumerate}
\end{remark}

\clearpage
\section{Multibubble and Cascade Exclusion}\label{sec:multibubble-cascade-exclusion}
\medskip
\noindent\textbf{Analytic inputs.}

The local reduction theorem uses the following analytic inputs from the
preceding sections, with profile decompositions and decoupling estimates in
their standard critical-space sense \cite{Gallagher2001,BrezisLieb1983}.  The
remaining work is the multibubble, cascade, and terminal-frame decoupling.
\begin{enumerate}[label=\textup{(\roman*)}]
\item From \cref{sec:single-core-criterion}, the compact single-core criterion \cref{paper1:thm:single-core-criterion}: a
compact single-core finite-cost Type II branch with positive local CKN
concentration, nondegenerate scale-center selection, pressure reconstruction,
bounded compact-cylinder suitable estimates, bounded modulation, and vanishing
localized dissipation is excluded.
\item From \cref{sec:repaired-gauge-construction}, the renormalized equation \cref{paper2:thm:ren-NS}, pressure decomposition \cref{paper2:thm:pressure-decomp}, modulation system \cref{paper2:thm:mod-sys}, and repaired-gauge representation \cref{paper2:thm:A,paper2:thm:B,paper2:thm:C}: the renormalized equation,
pressure decomposition modulo functions of time, bounded modulation
coefficients, renormalized suitable structure, and stability of these structures
under subsequences.
\item From \cref{sec:rough-core-exclusion}, the compact-cylinder Caccioppoli estimate \cref{paper4:thm:caccioppoli} and rough-core exclusion \cref{paper4:cor:rough-core}.
\item From \cref{sec:scale-collapse-cost}, the collapse cost exclusion \cref{paper5:thm:collapse-core-floor,paper5:thm:cost-exclusion} and autonomous reduced-limit theorem \cref{paper5:thm:autonomous-limit} for scale-rigid or scale-collapsing branches.
\end{enumerate}

\subsection{Local setting and active frames}

\medskip
\noindent\textit{Normalized local Navier--Stokes setting.}
All quantities are written in rescaled local variables \((x,t)\in Q_R=B_R\times(-R^2,0)\)
with viscosity \(\nu=1\) (critical scaling has absorbed \(\nu\)).  A local pair
\((u_n,p_n)\) on such a cylinder satisfies
\[
  \partial_tu_n-\Delta u_n+(u_n\cdot\nabla)u_n+\nabla p_n=0,\qquad \nabla\cdot u_n=0,
\]
and pressure is normalised through spatial averaging
\[
  (p_n)_{B_\rho}(t):=|B_\rho|^{-1}\int_{B_\rho}p_n(x,t)\,dx
\]
to factor out time-dependent constants.

For $r>0$, $x_0\in\mathbb R^3$, and $t_0\in\mathbb R$, let
\[
  Q_r(x_0,t_0):=B_r(x_0)\times(t_0-r^2,t_0),
  \qquad
  Q_r:=Q_r(0,0)=B_r\times(-r^2,0).
\]

Let $(u_n,p_n)$ be parabolic rescalings around a potential singular point with
positive local Caffarelli--Kohn--Nirenberg concentration on compact cylinders.
Concretely, for a base triple \((x_n,t_n,\lambda_n)\),
\[
  u_n(y,s)=\lambda_n u(x_n+\lambda_n y,\ t_n+\lambda_n^2 s),\qquad
  p_n(y,s)=\lambda_n^2 p(x_n+\lambda_n y,\ t_n+\lambda_n^2 s).
\]
For each $n$ and $R>0$, define the scaling-invariant local quantities
\[
  A_n(R):=R^{-1}\operatorname*{ess\,sup}_{-R^2<t<0}\int_{B_R}|u_n(x,t)|^2\,dx,
\]
\[
  E_n(R):=R^{-1}\int_{Q_R}|\nabla u_n(x,t)|^2\,dx\,dt,
\]
\[
  C_n(R):=R^{-2}\int_{Q_R}|u_n(x,t)|^3\,dx\,dt,
  \qquad
  D_n(R):=R^{-2}\int_{-R^2}^0\int_{B_R}|p_n-(p_n)_{B_R}(t)|^{3/2}\,dx\,dt,
\]
where
\[
  (p_n)_{B_R}(t):=|B_R|^{-1}\int_{B_R}p_n(x,t)\,dx.
\]

For an active frame \((x_n,\lambda_n)\), \(\lambda_n>0\), define the critical rescaling
\[
  \mathcal T_{x_n,\lambda_n}f(y):=\lambda_n f(x_n+\lambda_n y).
\]
If \(\phi\in L^3_\sigma(\mathbb R^3)\) is observed in this frame, its contribution is
\[
  \phi_n(y)=\mathcal T_{x_n,\lambda_n}\phi(y)=\lambda_n\phi(x_n+\lambda_n y).
\]
The critical Jacobian identity is exact:
\[
  \|\mathcal T_{x_n,\lambda_n}f\|_{L^3(E)}^3
  =\int_{x_n+\lambda_nE}|f(x)|^3\,dx
\]
for every measurable \(E\subset\mathbb R^3\).  In particular
\[
  \|\mathcal T_{x_n,\lambda_n}f\|_{L^3(\mathbb R^3)}^3
  =\|f\|_{L^3(\mathbb R^3)}^3,
\]
so \(\mathcal T_{x_n,\lambda_n}\) is an isometry on \(L^3(\mathbb R^3)\).

\begin{definition}[Bounded selected-window limit]
A selected-window subsequence has a bounded selected-window limit on a compact
cylinder $Q_R$ if there exists $u_\infty\in L^\infty(Q_R)$ such that, after
extracting a subsequence,
\[
  \forall\varepsilon>0\ \exists N:\ \forall n\ge N,\ \|u_n-u_\infty\|_{L^3(Q_R)}<\varepsilon.
\]
The limit is nonzero if $u_\infty\not\equiv0$ on $Q_R$.
\end{definition}

\begin{definition}[Local Type II sequence]
A positive local concentration sequence is called a local Type II sequence if every
selected-window rescaling has zero bounded selected-window limit on every compact
cylinder \(Q_R\):
\[
  \forall R>0,\qquad \lim_{n\to\infty}\|u_n\|_{L^3(Q_R)}=0.
\]
\end{definition}

\begin{definition}[Active frame]\label{paper3:def:active-frame}
Fix a concentrating time sequence \(t_n\uparrow T\) and a compact time interval
\(I\Subset(-\infty,0]\).  An active frame is a tuple
\[
  \mathfrak F^j=\bigl(x^j_n,\lambda^j_n,I,U^j_n,P^j_n\bigr),
  \qquad \lambda^j_n\downarrow0,
\]
where \(x^j_n\in\mathbb R^3\) and
\[
  U^j_n(y,s):=\lambda^j_n u(x^j_n+\lambda^j_n y,t_n+(\lambda^j_n)^2s),
  \qquad
  P^j_n(y,s):=(\lambda^j_n)^2p(x^j_n+\lambda^j_n y,t_n+(\lambda^j_n)^2s)
\]
on \(B_{2R}\times I\) for every fixed compact cylinder under consideration.
The frame is active at threshold \(\eta_*>0\) if for some fixed \(R_0<\infty\)
\[
  \liminf_{n\to\infty}\int_{Q_{R_0}}|U^j_n|^3\,dy\,ds\ge \eta_*^3.
\]
\end{definition}

\begin{definition}[Relations between active frames]\label{paper3:def:frame-relations}
Let \(\mathfrak F^i\) and \(\mathfrak F^j\) be two active frames.  After passing
to a subsequence exactly one of the following parameter regimes holds.
\begin{enumerate}[label=\textup{(\roman*)}]
\item They are \emph{same-point comparable-scale} if
\[
  0<c\le\frac{\lambda^i_n}{\lambda^j_n}\le C<\infty,\qquad
  \frac{|x^i_n-x^j_n|}{\max(\lambda^i_n,\lambda^j_n)}\le C
\]
for all large \(n\).
\item They form a \emph{same-point cascade} if, after relabeling,
\[
  \frac{\lambda^i_n}{\lambda^j_n}\to0,\qquad
  \frac{|x^i_n-x^j_n|}{\lambda^j_n}=O(1).
\]
The frame \(i\) is then the inner frame and \(j\) is an outer frame.
\item They are \emph{separated-point} if
\[
  \frac{|x^i_n-x^j_n|}{\lambda^i_n}\to\infty
  \quad\text{and}\quad
  \frac{|x^i_n-x^j_n|}{\lambda^j_n}\to\infty .
\]
\end{enumerate}
\end{definition}

\begin{definition}[Terminal frame and perturbative remainder]\label{paper3:def:terminal-frame}
Given a finite active family, a terminal frame is chosen by selecting an active
scale with minimal order:
\[
  \lambda^{j_*}_n\le C_j\lambda^j_n
  \qquad\text{for every active }j
\]
after passing to a subsequence.  Comparable frames at the same physical point
are grouped before this choice; if several grouped frames remain tied, one is
chosen by a fixed lexicographic ordering of their labels.  A remainder \(r_n\)
in a terminal frame is perturbative on \(B_{2R}\times I\) if
\[
  \|r_n\|_{L^\infty_I L^3(B_{2R})}\to0
\]
and its equation residual and pressure satisfy
\[
  \|f_n\|_{L^1_IH^{-1}(B_R)}+\|\nabla\pi^r_n\|_{L^1_IH^{-1}(B_R)}\to0
\]
for every compact \(B_R\times I\).
\end{definition}

\begin{definition}[Local multibubble/cascade branch]\label{paper3:def:local-multibubble}
A local multibubble/cascade branch is a positive-concentration local Type II
branch for which at least one of the following alternatives occurs:
\begin{enumerate}[label=\textup{(\roman*)}]
\item the compact-cylinder suitable estimates $A_n+E_n+C_n+D_n$ are unbounded on some
  fixed rescaled cylinder;
\item positive CKN mass splits into two or more comparable active cores;
\item a strict same-point scale cascade appears inside bounded rescaled cylinders;
\item separated physical points remain active in different local rescaled frames;
\item the localized scale or centering selection degenerates (no unique core
  selected).
\end{enumerate}
\end{definition}

\begin{definition}[Active profile family]\label{paper3:def:active-family}
An active profile family is a family
\[
  \{(\phi^j,x_n^j,\lambda_n^j)\}_{j\in J},
  \qquad \phi^j\in L^3_\sigma(\mathbb R^3),
\]
with finite or countable index set $J$, extracted along a positive local
concentration sequence.  In physical variables write
\[
  \Lambda_{\lambda,x_0}\phi(x):=\lambda^{-1}\phi\left(\frac{x-x_0}{\lambda}\right).
\]
For every finite \(J_0\subset J\) the family gives an expansion
\[
  u(\cdot,t_n)=\sum_{j\in J_0}\Lambda_{\lambda^j_n,x^j_n}\phi^j+r^{J_0}_n
\]
on the compact terminal windows under consideration.  After the finite active
subfamily has been selected, the corresponding remainder satisfies the
perturbative bounds of \cref{paper3:def:terminal-frame}.  The parameters are
considered only after grouping same-point comparable-scale profiles into
compound profiles as in \cref{paper3:def:compound-profile}.  A profile is active if
\[
  \|\phi^j\|_{L^3(\mathbb R^3)}\ge\eta_*>0.
\]
The family has finite critical mass when
\[
  \sum_{j\in J}\|\phi^j\|_{L^3(\mathbb R^3)}^3\le M_*<\infty .
\]
It captures positive local concentration if every compact terminal-window
sequence with
\[
  \liminf_{n\to\infty}\int_{Q_R}|u_n|^3>0
\]
either contains an active profile observed in that window or is absorbed into
the selected terminal frame; equivalently, no positive terminal-window
concentration remains in the perturbative remainder.
\end{definition}

\begin{definition}[Compound profile]\label{paper3:def:compound-profile}
Let \(J_k\subset J\) be a same-point comparable-scale class.  Choose one
representative \((x_n^{j_k},\lambda_n^{j_k})\).  After passing to a subsequence,
\[
  \rho_j:=\lim_{n\to\infty}\frac{\lambda^j_n}{\lambda^{j_k}_n}\in(0,\infty),
  \qquad
  z_j:=\lim_{n\to\infty}\frac{x^j_n-x^{j_k}_n}{\lambda^{j_k}_n}\in\mathbb R^3 .
\]
The compound profile in the representative frame is the \(L^3_\sigma\)-sum
\[
  \Phi^k(y):=\sum_{j\in J_k}\rho_j^{-1}
  \phi^j\left(\frac{y-z_j}{\rho_j}\right),
\]
whenever \(J_k\) is finite; finiteness follows for active classes from
\cref{paper3:lem:finite-active-count}.  If \(\|\Phi^k\|_{L^3}<\varepsilon_\star\), where
\(\varepsilon_\star\) is the critical small-data threshold, the class is
perturbative.  Otherwise \(\Phi^k\) is treated as one active object in its
representative frame.
\end{definition}

\begin{lemma}[Finite active profile count]\label{paper3:lem:finite-active-count}
Every active profile family with finite critical mass contains only finitely many
active profiles.  More precisely,
\[
  \#J\le M_*\eta_*^{-3}.
\]
\end{lemma}

\begin{proof}
For each active profile $j\in J$,
\[
  \|\phi^j\|_{L^3(\mathbb R^3)}^3\ge\eta_*^3.
\]
Let $J_0\subset J$ be finite. Summing over $J_0$ gives
\[
  \#J_0\,\eta_*^3
  \le
  \sum_{j\in J_0}\|\phi^j\|_{L^3(\mathbb R^3)}^3
  \le
  \sum_{j\in J}\|\phi^j\|_{L^3(\mathbb R^3)}^3
  \le M_*.
\]
Hence $\#J_0\le M_*\eta_*^{-3}$.  Taking the supremum over all finite
subsets $J_0$ yields the same bound for the cardinality of $J$, so $J$ is finite
and
\[
  \#J\le M_*\eta_*^{-3}.
\]
\end{proof}

\begin{lemma}[Exhaustive parameter partition]\label{paper3:lem:parameter-partition}
Let \(\mathfrak F^i,\mathfrak F^j\) be two active frames.  After passing to a
subsequence, either the pair is same-point comparable-scale, same-point cascade,
or separated-point in the sense of \cref{paper3:def:frame-relations}; if none of these
relations holds in a selected compact frame, then one of the two contributions
is locally invisible there in \(L^3\).
\end{lemma}

\begin{proof}
Pass to a subsequence so that
\[
  \frac{\lambda^i_n}{\lambda^j_n}\to \ell\in[0,\infty]
\]
in the extended half-line.  If \(0<\ell<\infty\), then either
\[
  \frac{|x^i_n-x^j_n|}{\max(\lambda^i_n,\lambda^j_n)}
\]
is bounded along a further subsequence, giving the same-point comparable-scale
case, or it tends to infinity, giving separated points because the two scales
are comparable.
\par
If \(\ell=0\), then either
\[
  \frac{|x^i_n-x^j_n|}{\lambda^j_n}
\]
is bounded, giving a same-point cascade with \(i\) inner and \(j\) outer, or it
tends to infinity.  In the latter case the \(j\)-profile is locally invisible in
the \(i\)-frame by \cref{paper3:lem:separated-vanish}.  The case \(\ell=\infty\) is the
same after interchanging \(i\) and \(j\).  These alternatives exhaust all
subsequential limits of the scale ratio and normalized center distance.
\end{proof}

\begin{theorem}[Minimal bad-configuration extraction]\label{paper3:thm:minimal-bad-extraction}
Let a positive local Type II branch admit an active profile family with finite
critical mass, perturbative remainder in the sense of
\cref{paper3:def:terminal-frame}, and concentration capture in the sense of
\cref{paper3:def:active-family}.  If the branch is not already a single active
terminal frame satisfying the single-core hypotheses, then after passing to a
subsequence one of the following occurs:
\begin{enumerate}[label=\textup{(\roman*)}]
\item compact-cylinder suitable control fails, and
\cref{paper3:lem:compactness-failure} produces a bounded-cylinder concentration core;
\item there are at least two active frames in a same-point comparable-scale
class, hence a compound profile in the sense of \cref{paper3:def:compound-profile};
\item there is a same-point cascade;
\item there are separated active physical points;
\item the scale-center selection of the selected compound or terminal frame is
degenerate.
\end{enumerate}
In every case the positive local concentration is carried by an active frame or
compound profile, not by the perturbative remainder.
\end{theorem}

\begin{proof}
If compact-cylinder suitable control fails, item \((i)\) follows from
\cref{paper3:lem:compactness-failure}.  Assume therefore that the compact-cylinder
family is bounded on the windows under consideration.
\par
By concentration capture, positive local concentration cannot remain solely in
the perturbative remainder.  Hence either one selected active frame carries the
terminal-window concentration, or another active frame remains non-negligible.
If there is one nondegenerate selected active frame and no competing active
object, the branch is in the single-core regime.  Since this case is
excluded by hypothesis, either the selected scale-center map is degenerate,
giving item \((v)\), or at least two active frames remain.
\par
For two active frames, apply \cref{paper3:lem:parameter-partition}.  Same-point
comparable-scale frames give item \((ii)\), same-point scale separation gives
item \((iii)\), and separated points give item \((iv)\).  The remaining
alternative in \cref{paper3:lem:parameter-partition} is local invisibility of one
contribution in the selected compact frame; such a contribution is absorbed
into the perturbative remainder and cannot carry the retained positive local
concentration by the capture property.  This gives the extraction statement.
\end{proof}

\subsection{Elementary local decoupling estimates}

\begin{lemma}[Critical rescalings that escape are locally invisible]\label{paper3:lem:escaping-l3}
Let \(\phi\in L^3(\mathbb R^3)\), let \(R<\infty\), and define
\[
  W_n(y)=\rho_n\phi(z_n+\rho_n y),
\]
where \(\rho_n>0\).  Assume that the sets
\[
  z_n+\rho_nB_R
\]
either have measure tending to zero, or escape to spatial infinity in the sense
that for every compact set \(K\subset\mathbb R^3\),
\[
  (z_n+\rho_nB_R)\cap K=\varnothing
\]
for all sufficiently large \(n\).  Then
\[
  \|W_n\|_{L^3(B_R)}\to0.
\]
\end{lemma}

\begin{proof}
By the critical change of variables,
\[
  \|W_n\|_{L^3(B_R)}^3
  =\int_{B_R}\rho_n^3|\phi(z_n+\rho_n y)|^3\,dy
  =\int_{z_n+\rho_n B_R}|\phi(x)|^3\,dx.
\]
Set
\[
  m_n:=\int_{z_n+\rho_n B_R}|\phi(x)|^3\,dx.
\]
Since \(\|W_n\|_{L^3(B_R)}^3=m_n\), it is enough to show \(m_n\to0\).
\par
Fix \(\varepsilon>0\).  We split into the two structural alternatives.

\smallskip
\noindent\textit{Case 1: shrinking measure.}
If \(|z_n+\rho_nB_R|\to0\), then by absolute continuity of the integral
\(|\phi|^3\in L^1(\mathbb R^3)\) there exists \(\delta>0\) such that
\[
  |E|<\delta\Longrightarrow \int_E|\phi|^3\,dx<\varepsilon.
\]
Choose \(N\) with \(|z_n+\rho_nB_R|<\delta\) for all \(n\ge N\).  Then
\[
  m_n=\int_{z_n+\rho_nB_R}|\phi|^3\,dx<\varepsilon,\qquad n\ge N.
\]

\smallskip
\noindent\textit{Case 2: escape to infinity.}
Assume now that for every compact set \(K\subset\mathbb R^3\),
\((z_n+\rho_nB_R)\cap K=\varnothing\) for all sufficiently large \(n\).  By
absolute continuity at infinity, choose compact \(K\subset\mathbb R^3\) such that
\[
  \int_{\mathbb R^3\setminus K}|\phi|^3<\varepsilon.
\]
For every such compact \(K\), there is \(N(K)\) with
\((z_n+\rho_nB_R)\cap K=\varnothing\) for all \(n\ge N(K)\).  For these
indices,
\[
  z_n+\rho_nB_R \subset \mathbb R^3\setminus K,
\]
and therefore, for \(n\ge N(K)\),
\[
  m_n=\int_{z_n+\rho_nB_R}|\phi|^3\,dx
  \le
  \int_{\mathbb R^3\setminus K}|\phi|^3
  <\varepsilon.
\]
In either case, for every \(\varepsilon>0\), \(m_n<\varepsilon\) for \(n\) sufficiently
large.  Hence \(m_n\to0\), and therefore
\[
  \|W_n\|_{L^3(B_R)}\to0.
\]
\end{proof}

\begin{lemma}[Separated active profiles vanish in the selected frame]\label{paper3:lem:separated-vanish}
Let \((x^p_n,\lambda^p_n)\) and \((x^q_n,\lambda^q_n)\) be two active frames
with separated physical centers in the sense that
\[
  \frac{|x^q_n-x^p_n|}{\lambda^p_n}\to\infty.
\]
If \(\phi^q\in L^3(\mathbb R^3)\), then the \(q\)-profile observed in the
\(p\)-frame converges to zero locally in \(L^3\): for every \(R<\infty\),
\[
  \left\|
  \frac{\lambda^p_n}{\lambda^q_n}
  \phi^q\left(\frac{x^p_n-x^q_n}{\lambda^q_n}
       +\frac{\lambda^p_n}{\lambda^q_n}y\right)
  \right\|_{L^3(B_R)}\to0.
\]
\end{lemma}

\begin{proof}
The displayed function has the form in \cref{paper3:lem:escaping-l3} with
\[
  z_n=\frac{x^p_n-x^q_n}{\lambda^q_n},
  \qquad
  \rho_n=\frac{\lambda^p_n}{\lambda^q_n}.
\]
The corresponding physical set in the \(q\)-profile variables is
\[
  z_n+\rho_nB_R
  =\frac{x^p_n-x^q_n}{\lambda^q_n}
   +\frac{\lambda^p_n}{\lambda^q_n}B_R.
\]
For every \(y\in B_R\),
\[
  |z_n+\rho_n y|
  \ge |z_n|-\rho_nR.
\]
Set \(E_n:=z_n+\rho_nB_R\).  Then
\[
  \operatorname{dist}(E_n,0)
  =|z_n|-\rho_nR
  =\frac{\lambda^p_n}{\lambda^q_n}\left(\frac{|x^q_n-x^p_n|}{\lambda^p_n}-R\right).
\]
Since \(\frac{|x^q_n-x^p_n|}{\lambda^p_n}\to\infty\), there exists \(N_0(R)\) such that
\[
  \frac{|x^q_n-x^p_n|}{\lambda^p_n}-R\ge1,\quad n\ge N_0(R),
\]
and hence for such \(n\), \(\operatorname{dist}(E_n,0)\ge\rho_n\).
\par
There are two cases.
\begin{enumerate}[label=\textup{(\alph*)}]
\item If \(\rho_n\to0\), then
 \[
   |z_n+\rho_nB_R|=\rho_n^3|B_R|\to0.
 \]
 Then the first alternative in \cref{paper3:lem:escaping-l3} applies.
\item Otherwise, after passing to a subsequence, \(\rho_n\ge\rho_\ast>0\).  Then
 \[
   \operatorname{dist}(z_n+\rho_nB_R,0)\ge \rho_\ast\left(\frac{|x^q_n-x^p_n|}{\lambda^p_n}-R\right)\to\infty,
 \]
 so the sets escape every fixed compact subset of \(\mathbb R^3\).  The
 second alternative in \cref{paper3:lem:escaping-l3} applies.
\end{enumerate}
In both subcases the \(L^3(B_R)\)-norm of the displayed expression converges to
zero.
\end{proof}

\begin{lemma}[Strict outer scales vanish in the innermost frame]\label{paper3:lem:outer-vanish}
Let several active profiles have the same physical center and scales
\(\lambda^1_n,\ldots,\lambda^J_n\).  Suppose \(\lambda^1_n\) is an innermost
scale and
\[
  \frac{\lambda^1_n}{\lambda^j_n}\to0
  \qquad (j\ne1).
\]
Then every outer profile \(\phi^j\in L^3(\mathbb R^3)\), viewed in the
\(\lambda^1_n\)-frame, converges to zero in \(L^3(B_R)\) for each fixed
\(R<\infty\).  Any constant ambient velocity arising from the regular part of
the outer flow is absorbed into the translation modulation.
\end{lemma}

\begin{proof}
In the innermost variables, the \(j\)-th outer contribution has the form
\[
  W^j_n(y)=\frac{\lambda^1_n}{\lambda^j_n}
  \phi^j\left(\frac{x^1_n-x^j_n}{\lambda^j_n}
  +\frac{\lambda^1_n}{\lambda^j_n}y\right).
\]
Set
\[
  \rho_n:=\frac{\lambda^1_n}{\lambda^j_n}\to0,
  \qquad
  z_n^j:=\frac{x^1_n-x^j_n}{\lambda^j_n}.
\]
Then
\[
  W^j_n(y)=\rho_n\phi^j\left(z_n^j+\rho_n y\right),
\]
and by the critical change of variables,
\[
  \|W^j_n\|_{L^3(B_R)}^3
  =\int_{z_n^j+\rho_nB_R}|\phi^j(x)|^3\,dx.
\]
Since \(\rho_n\to0\), the translated domains have
\[
  |z_n^j+\rho_nB_R|=\rho_n^3|B_R|\to0.
\]
By absolute continuity of the integral for \(\phi^j\in L^3(\mathbb R^3)\),
\[
  \|W^j_n\|_{L^3(B_R)}\to0.
\]
To treat a general \(\phi^j\), let \(\phi^j_\varepsilon\in C_c^\infty(\mathbb R^3)\) with
\(\|\phi^j-\phi^j_\varepsilon\|_{L^3}<\varepsilon\).  By critical isometry,
\[
  \|W^j_n-(W^j_n)_\varepsilon\|_{L^3(B_R)}
  \le
  \|\phi^j-\phi^j_\varepsilon\|_{L^3} < \varepsilon,
\]
where \((W^j_n)_\varepsilon\) is the same scaling of \(\phi^j_\varepsilon\).
The compactly supported term has vanishing local norm by the shrinking-domain
argument above, and then letting first \(n\to\infty\), then \(\varepsilon\downarrow0\)
implies the claimed convergence.
A constant vector field from the regular part of an outer flow contributes only
a Galilean drift in the selected frame; it is therefore incorporated into the
translation modulation and does not remain as an exterior perturbation.
\end{proof}

\begin{lemma}[Compactness failure produces another core]\label{paper3:lem:compactness-failure}
If the compact-cylinder suitable estimates are unbounded on a fixed rescaled
cylinder, then after passing to a subsequence the local \(C_n+D_n\) quantity is
unbounded on a comparable cylinder.  Consequently the branch contains an
additional local CKN concentration core.
\end{lemma}

\begin{proof}
Fix \(r>0\).  Apply the local Caccioppoli inequality with outer radius \(2r\).
For each index \(n\), suitable weak solutions satisfy, with \(C_r\) independent
of \(n\),
\[
  A_n(r)+E_n(r)
  \le C_r\bigl(1+C_n(2r)+C_n(2r)^{2/3}+D_n(2r)\bigr),
\]
where the pressure is normalized by subtracting spatial means
\((p_n)_{B_{2r}}(t)\).
\par
Suppose that \(A_n(r)+E_n(r)+C_n(r)+D_n(r)\) is unbounded along a subsequence.
If \(C_n(2r)+D_n(2r)\) were bounded on this subsequence, the preceding estimate
would give uniform boundedness of \(A_n(r)+E_n(r)\).  Monotonicity of the local
integrals also gives \(C_n(r)+D_n(r)\le C_n(2r)+D_n(2r)\), contradicting
unboundedness of the full family on \(Q_r\).  Therefore, after extracting a
subsequence,
\[
  C_n(2r)+D_n(2r)\to\infty .
\]
\par
Define the localized CKN concentration at variable center and radius
\[
  G_n(x,t,\rho):=\rho^{-2}\int_{Q_\rho(x,t)}|u_n|^3\,dx\,ds
  +\rho^{-2}\int_{Q_\rho(x,t)}|p_n-(p_n)_{B_\rho}(s)|^{3/2}\,dx\,ds.
\]
Set
\[
  \mathbf M_n:=\sup\{G_n(x,t,\rho):0<\rho\le 2r,\ Q_\rho(x,t)\subset Q_{4r}\}.
\]
The choice \((x,t,\rho)=(0,0,2r)\) is admissible because \(Q_{2r}(0,0)\subset Q_{4r}\), so
\[
  \mathbf M_n\ge C_n(2r)+D_n(2r)\to\infty.
\]
After extracting a further subsequence, we may assume \(\mathbf M_n\ge1\) for all \(n\).
Choose \((x_n,t_n,\rho_n)\) with \(0<\rho_n\le2r\), \(Q_{\rho_n}(x_n,t_n)\subset Q_{4r}\),
such that
\[
  G_n(x_n,t_n,\rho_n)\ge\tfrac12\,\mathbf M_n\ge\tfrac12.
\]
Rescale around \((x_n,t_n,\rho_n)\):
\[
  v_n(y,s):=\rho_nu_n(x_n+\rho_n y,t_n+\rho_n^2 s),\qquad
  q_n(y,s):=\rho_n^2p_n(x_n+\rho_n y,t_n+\rho_n^2 s).
\]
Using the critical Jacobian identity,
\[
  \int_{Q_1}|v_n|^3\,dy\,ds
  +\int_{Q_1}|q_n-(q_n)_{B_1}(s)|^{3/2}\,dy\,ds
  =G_n(x_n,t_n,\rho_n)\ge\tfrac12.
\]
Hence this new frame carries positive local critical mass, i.e. a new local CKN
concentration core, as required.
\end{proof}

\subsection{Terminal active-frame decoupling}

\begin{definition}[Admissible terminal active family]\label{paper3:def:admissible-terminal}
Fix a compact cylinder \(B_{2R}\times I\) and a selected active frame.  A
terminal active family is admissible on this cylinder if the following hold.
\begin{enumerate}[label=\textup{(\roman*)}]
\item The active profiles satisfy the finite-mass and active-mass conditions of
\cref{paper3:def:active-family}.
\item Every nonselected active profile is either separated from the selected
physical point in the sense of \cref{paper3:lem:separated-vanish}, or has the same
physical point and a strictly outer scale in the sense of
\cref{paper3:lem:outer-vanish}, after comparable same-point scales have been grouped.
\item The rescaled remainder \(r_n\) satisfies
\[
  \|r_n\|_{L^\infty_I L^3(B_{2R})}\to0,
\]
and there exist fields \((\pi_n^r,f_n)\) such that on \(B_{2R}\times I\),
\[
  \partial_t r_n-\Delta r_n+a_n\bigl(r_n+y\cdot\nabla r_n\bigr)+b_n\cdot\nabla r_n
  +\nabla\pi_n^r=f_n,\qquad
  \nabla\cdot r_n=0,
\]
with
\[
  \|f_n\|_{L^1_IH^{-1}(B_R)}\to0,
\]
\[
  \|\nabla\pi_n^r\|_{L^1_IH^{-1}(B_R)}\to0.
\]
\item The selected velocity satisfies
\[
  \sup_n\|U_n\|_{L^\infty_I L^3(B_{2R})}<\infty .
\]
\item The selected modulation coefficients satisfy
\[
  \sup_n\bigl(\|a_n\|_{L^\infty(I)}+\|b_n\|_{L^\infty(I)}\bigr)<\infty .
\]
\item Local pressure decomposition is stable under the splitting: if tensor
coefficients tend to zero in \(L^1_IL^{3/2}(B_R)\), then the corresponding
pressure gradients tend to zero in \(L^1_IH^{-1}(B_R)\), after subtracting
spatial means.
\item Cutoff commutators generated by localizing to \(B_R\) are controlled by
the \(L^\infty_I L^3\), \(L^1_I L^{3/2}\), and pressure-gradient norms appearing
above.  In particular, for any compact \(\chi_R\in C_c^\infty(B_{2R})\) with
\(\chi_R\equiv1\) on \(B_R\), the commutators satisfy
\[
  \|[\Delta,\chi_R]u\|_{H^{-1}(B_R)}
  +\|[\nabla\chi_R,\,\nabla\cdot](u\otimes u)\|_{H^{-1}(B_R)}
  \le C_R\|u\|_{L^3(B_{2R})}^2.
\]
uniformly for \(u\in L^3(B_{2R})\) and the same type of estimate for
\(S_n\)-quadratic commutator terms.
\end{enumerate}
\end{definition}

\begin{theorem}[Terminal profile-evolution decoupling]\label{paper3:thm:terminal-decoupling}
Consider an admissible terminal active family on \(B_{2R}\times I\).  Let
\(U_n\) be the selected-frame velocity and let \(S_n\) be the sum, in that
frame, of all nonselected active profiles together with the rescaled remainder.
Then on \(B_R\times I\):
\begin{enumerate}[label=\textup{(\roman*)}]
\item \(\nabla\cdot S_n=0\) distributionally;
\item \(S_n\to0\) in \(L^\infty_I L^3(B_{2R})\);
\item \(U_n\otimes S_n+S_n\otimes U_n+S_n\otimes S_n\to0\) in
\(L^1_I L^{3/2}(B_R)\);
\item
\[
  \begin{aligned}
  &\exists\,\pi_n,\ F_n:\quad
  \partial_t S_n-\Delta S_n
  +a_n\bigl(y\cdot\nabla S_n + S_n\bigr)
  +b_n\cdot\nabla S_n+\nabla\pi_n =F_n,\\
  &\|F_n\|_{L^1_IH^{-1}(B_R)}\to0,\qquad
  \nabla\cdot S_n=0\quad\text{in }\mathcal D'(B_R\times I).
  \end{aligned}
\]
\[
  \|\nabla\pi_n\|_{L^1_IH^{-1}(B_R)}\to0.
\]
\item every pressure source containing at least one factor of \(S_n\) has
vanishing pressure gradient in \(L^1_IH^{-1}(B_R)\);
\item localization errors introduced by compact cutoffs vanish in
\(L^1_IH^{-1}(B_R)\);
\item the effective modulation parameters of the selected frame remain bounded;
\item after subtracting \(S_n\), the selected frame remains a positive finite-mass
local Type II branch with pressure reconstruction and the compact-cylinder
bounds required by the single-core analysis.
\end{enumerate}
\end{theorem}

\begin{proof}
Each profile in the family is divergence-free, and the rescaled remainder is
divergence-free.  Since divergence commutes with the Navier--Stokes critical
rescalings, finite sums remain divergence-free in distributions.  This establishes
\((i)\).

We prove \((ii)\).  By admissibility, every nonselected active profile is either
separated from the selected physical point or is a same-point strictly outer
scale.  The separated profiles vanish locally by \cref{paper3:lem:separated-vanish},
and the same-point outer profiles vanish locally by \cref{paper3:lem:outer-vanish}.
The remainder tends to zero in \(L^\infty_I L^3(B_{2R})\) by
\cref{paper3:def:admissible-terminal}.  The number of active profiles is finite by
\cref{paper3:lem:finite-active-count}.  Since every summand in \(S_n\)
tends to zero in \(L^\infty_I L^3(B_{2R})\), the triangle inequality gives
\[
  \|S_n\|_{L^\infty_I L^3(B_{2R})}
  \le
  \sum_{j\ne j_*}\|S^j_n\|_{L^\infty_I L^3(B_{2R})}
  +\|r_n\|_{L^\infty_I L^3(B_{2R})}\to0.
\]
This establishes \((ii)\).

For \((iii)\), admissibility gives
\[
  \sup_n\|U_n\|_{L^\infty_I L^3(B_{2R})}<\infty.
\]
Hence, by H\"older's inequality,
\[
  \|U_n\otimes S_n\|_{L^1_IL^{3/2}(B_R)}
  \le |I|\|U_n\|_{L^\infty_IL^3(B_R)}
       \|S_n\|_{L^\infty_IL^3(B_R)}\to0,
\]
and the same estimate applies to \(S_n\otimes U_n\).  Finally,
\[
  \|S_n\otimes S_n\|_{L^1_IL^{3/2}(B_R)}
  \le |I|\|S_n\|_{L^\infty_IL^3(B_R)}^2\to0.
\]

For \((iv)\), write the represented Navier--Stokes equation for the full
profile sum \(U_n+S_n\) and subtract the represented equation for the selected
component \(U_n\).  The remaining equation for \(S_n\) has the displayed linear
left-hand side
\[
  \partial_t S_n-\Delta S_n
  +a_n\bigl(y\cdot\nabla S_n+S_n\bigr)+b_n\cdot\nabla S_n.
\]
Collect all right-hand terms into
\[
  F_n:= -\nabla\cdot(U_n\otimes S_n+S_n\otimes U_n+S_n\otimes S_n)
       +f_n+E_n^{\mathrm{cut}}+E_n^{\mathrm{mod}},
\]
where \(f_n\) is the remainder forcing from item \((iii)\) of admissibility, and
\(E_n^{\mathrm{cut}},E_n^{\mathrm{mod}}\) are localization and modulation errors.
The nonlinear stresses tend to zero in \(L^1_IL^{3/2}(B_R)\), hence in
\(L^1_IH^{-1}(B_R)\), because \(L^{3/2}(B_R)\hookrightarrow H^{-1}(B_R)\) on a
bounded three-dimensional ball.  The residual term vanishes by
\cref{paper3:def:admissible-terminal}.  The cutoff terms vanish by the localization
part of admissibility, recorded explicitly in \((vi)\).  This gives the
displayed convergence in \(L^1_IH^{-1}(B_R)\).

For \((v)\), decompose the pressure source of \(U_n+S_n\) into the pure selected
source, the pure exterior source, and the two mixed sources.  The mixed sources
are
\[
  \partial_i\partial_j(U_{n,i}S_{n,j}+S_{n,i}U_{n,j})
  \quad\text{and}\quad
  \partial_i\partial_j(S_{n,i}S_{n,j}).
\]
By \((iii)\), their tensor coefficients tend to zero in
\(L^1_IL^{3/2}(B_R)\).  The pressure-stability part of admissibility,
equivalently local Calderon--Zygmund estimates plus harmonic pressure control
modulo spatial means, gives convergence of the corresponding pressure gradients
to zero in \(L^1_IH^{-1}(B_R)\).  Pure exterior sources are generated by
profiles that are locally invisible in the selected frame; applying the same
pressure stability after \cref{paper3:lem:separated-vanish,paper3:lem:outer-vanish} gives
their vanishing pressure gradients.

For \((vi)\), let \(\chi_R\) be a cutoff equal to one on \(B_R\) and supported
in \(B_{2R}\).  The commutators are finite sums of terms of the following
types:
\[
  (\nabla\chi_R)\cdot\nabla S_n,
  \qquad
  (\Delta\chi_R)S_n,
  \qquad
  (\nabla\chi_R)\cdot (U_n\otimes S_n+S_n\otimes U_n+S_n\otimes S_n),
\]
plus analogous pressure cutoff terms.  Their supports lie in the fixed annulus
\(B_{2R}\setminus B_R\).  By the localization part of admissibility, these
commutators are controlled by the norms already shown to vanish in \((ii)\),
\((iii)\), and \((v)\).  Hence all localization errors vanish in
\(L^1_IH^{-1}(B_R)\).

For \((vii)\), the only exterior contribution that can remain nonzero at first
order in a compact selected frame is a constant ambient velocity.  This has
already been absorbed into the translation modulation.  The remaining exterior
profiles vanish locally, so their contribution to the modulation equations is
small in the same local \(L^3\) topology used above.  The selected modulation
coefficients are bounded by admissibility; hence the effective parameters remain
bounded.

For \((viii)\), the positive CKN mass of the selected profile is unchanged by
subtracting a term that tends to zero in local \(L^3\) and whose
mixed pressure sources vanish by \((v)\).  The pressure reconstruction is stable
under the same decomposition by the pressure-stability condition in
\cref{paper3:def:admissible-terminal}.  The compact-cylinder bounds pass to the
selected component because the full family is bounded and the nonretained terms
vanish on compact cylinders.  Thus the selected frame remains a positive
finite-mass local Type II branch with the analytic inputs needed for the
single-core analysis.
\end{proof}

\subsection{Same-point and separated-point reductions}

\begin{theorem}[Same-point cascade reduction]\label{paper3:thm:same-point-reduction}
Every same-point multibubble or strict same-point cascade reduces, in the
selected innermost or compound frame, to one of the following: the single-core
criterion, scale-rigidity, or a perturbative terminal branch covered by
\cref{paper3:thm:terminal-decoupling}, provided the terminal family generated in the
selected innermost frame is admissible in the sense of
\cref{paper3:def:admissible-terminal}.
\end{theorem}

\begin{proof}
\noindent\textit{Step 1: scale-comparability classes.}
Split active same-point profiles into equivalence classes by
\[
  i\sim j \iff \exists c,C\in(0,\infty):
  c\le \frac{\lambda^i_n}{\lambda^j_n}\le C \quad \text{for all }n
\]
after passing to a subsequence.
Thus within each class there are constants \(c_{ij},C_{ij}\) with
\(c_{ij}\le\lambda^i_n/\lambda^j_n\le C_{ij}\) for all \(n\), so one may define
a single common rescaling sequence (choose \(\lambda_n^{(k)}=\lambda_n^{j_k}\)) and write
all class-\(k\) profiles in the same coordinates.
Within each class, finite sums of divergence-free profiles remain divergence-free.

\smallskip
\noindent\textit{Step 2: compound profile reduction inside each class.}
For each class set
\[
  \Phi^k_n:=\sum_{j\in J_k}\lambda^j_n\phi^j(x_n^j+\lambda^j_n y).
\]
If the \(L^3\)-size of this rescaled sum satisfies
\[
  \left\|\sum_{j\in J_k}\phi^j\right\|_{L^3(\mathbb R^3)}<\varepsilon_\star,
\]
with \(\varepsilon_\star\) the perturbative small-data threshold, then standard local
small-data theory for Navier--Stokes in the critical space eliminates this class
from the active alternatives.

If the class is above threshold, it is retained as a single compound active core;
at most one active core from that class is needed for local selection.

\smallskip
\noindent\textit{Step 3: local alternatives.}
If the retained compound core is unique and its scale-centering Jacobian is
nondegenerate, we obtain the single selected core alternative.
If this Jacobian is degenerate, the localized choice of active center or scale is
not unique, which is the gauge-degenerate alternative in
\cref{paper3:def:local-multibubble}.

\smallskip
\noindent\textit{Step 4: strict same-point scale separation.}
Assume now the branch is not reduced to comparable scales and choose the innermost
active scale \(\lambda^1_n\) in that physical point.  By
\cref{paper3:lem:outer-vanish}, every profile with larger scale vanishes in the
\(\lambda^1_n\)-frame in local \(L^3\), up to constant drift terms.
These constant drifts are absorbed into translation modulation by admissibility.
Let \(S_n\) be the sum of all such outer profiles plus the rescaled remainder
\(r_n\).  Then, by Step 2 and \cref{paper3:lem:outer-vanish},
\[
  S_n\to0\quad\text{in }L^\infty_I L^3(B_{2R}).
\]

Subtracting the equation for the innermost selected profile from the full-frame
equation gives a perturbation equation for the selected profile whose forcing
belongs to \(L^1_I H^{-1}(B_R)\) and vanishes.  Thus
\cref{paper3:thm:terminal-decoupling} applies, and one of three outcomes holds:
\begin{enumerate}[label=\textup{(\alph*)}]
\item \textbf{Single-core branch:}
  after subtracting \(S_n\), the branch is governed by one active component, so it
  reduces to the compact single-core branch of \cref{paper3:ass:single-core};
\item \textbf{Scale-rigid branch:}
  the relative Jacobian between comparable scales becomes rigid in the sense of
  \cref{paper3:ass:scale-rigidity}, so \cref{paper3:ass:scale-rigidity} excludes it;
\item \textbf{Purely perturbative branch:}
  \(S_n\to0\) and the selected profile satisfies a vanishing forcing perturbation
  in \(L^1_IH^{-1}(B_R)\).  No independent same-point cascade mechanism remains;
  after subtracting \(S_n\), the selected component is reduced to the single-core
  or scale-rigidity analysis.
\end{enumerate}

These are the strict same-point scale-separation alternatives, so the
same-point reduction is exhausted.
\end{proof}


\begin{theorem}[Separated-point decoupling]\label{paper3:thm:separated-point-reduction}
Separated active physical points are locally invisible in each other's active
rescaled frames.  Consequently every separated-point multibubble reduces in each
selected active frame to a single-core, scale-rigid, or perturbative terminal
branch, provided the terminal family in that selected active frame is
admissible in the sense of \cref{paper3:def:admissible-terminal}.
\end{theorem}

\begin{proof}
Fix an active point \(p\) with frame \((x^p_n,\lambda^p_n)\).  Let \(q\ne p\)
be another active point.  By separation,
\[
  \frac{|x^q_n-x^p_n|}{\lambda^p_n}\to\infty.
\]
Thus \cref{paper3:lem:separated-vanish} gives local \(L^3\) convergence to zero of the
\(q\)-profile in the \(p\)-frame on every compact ball.  The same conclusion
holds for each profile centered at any physical point different from \(p\).
Since only finitely many active profiles remain after the critical mass
threshold is imposed, the sum of all profiles outside the \(p\)-frame tends to
zero locally in \(L^3\).

The pure pressure generated by a separated profile is locally invisible by the
same change of variables whenever its profile pressure belongs to
\(L^{3/2}_{\mathrm{loc}}\) with the standard Calderon--Zygmund control.  Mixed
pressure terms contain at least one exterior velocity factor, hence vanish in
\(L^1H^{-1}_{\mathrm{loc}}\) by the pressure part of
\cref{paper3:thm:terminal-decoupling}.  Localization and modulation commutators are
localized to \(B_{2R}\setminus B_R\), so their \(H^{-1}\)-norm is controlled by
the \(L^3\)-smallness already proved in \cref{paper3:lem:separated-vanish,paper3:lem:outer-vanish}
and by admissibility item \((\text{vii})\).

Therefore, in the selected \(p\)-frame, all other physical points contribute
only a perturbative forcing tending to zero in \(L^1H^{-1}\) on compact
cylinders.  The selected frame remains a positive finite-mass branch with the
same compact-cylinder and pressure reconstruction properties.  Hence the branch
reduces to the single-core case, the scale-rigid case, or the perturbative
terminal branch.  Since \(p\) was arbitrary, the same reduction holds in every
active frame.
\end{proof}

\subsection{Multibubble and cascade exclusion}

The following statements isolate the analytic inputs needed for the
multibubble reduction: the (already proved) single-core exclusion, the
scale-rigidity reduction (a remaining hypothesis), and terminal active-frame
admissibility (proved here as a bridge from the terminal nonlinear decoupling
package of \cref{paper8:thm:terminal-nonlinear-decoupling}).  They are stated
explicitly so that the following theorem uses no unstated compactness,
pressure, or profile completeness conclusion.

\begin{proposition}[Local single-core exclusion]\label{paper3:ass:single-core}
Every compact single-core branch with positive local CKN concentration,
nondegenerate scale-center selection, pressure reconstruction, bounded
compact-cylinder suitable estimates, bounded modulation, and vanishing localized
dissipation is excluded from the local Type II regime.
\end{proposition}

\begin{proof}
This is precisely the conclusion of the local single-core Type II criterion,
\cref{paper1:thm:single-core-criterion}, applied to the represented compact
single-core branch in the renormalized chart.  The hypotheses listed here are
the input data of that theorem: positive local CKN concentration provides the
nondegenerate selected core, the nondegenerate scale-center selection together
with the pressure reconstruction and bounded compact-cylinder suitable
estimates are the inputs of the compactness statement
\cref{paper1:prop:compactness}, bounded modulation is
\cref{paper1:rem:bounded-modulation}, and vanishing localized dissipation is
\cref{paper1:def:localized-dissipation}.  The conclusion of
\cref{paper1:thm:single-core-criterion} is exactly the stated exclusion.
\end{proof}

\begin{assumption}[Local scale-rigidity exclusion]\label{paper3:ass:scale-rigidity}
Every scale-rigid local branch produced by the same-point or separated-point
reductions, including a branch with degenerate relative scale-center Jacobian,
either has a nonzero bounded selected-window limit or reduces to the single-core
branch in \cref{paper3:ass:single-core}.  Consequently no such branch is compatible
with the local Type II condition.
\end{assumption}

\begin{remark}[Status of \cref{paper3:ass:scale-rigidity}]
\label{paper3:rem:scale-rigidity-discharge}
The remainder of this subsection discharges
\cref{paper3:ass:scale-rigidity} by an explicit local stratification package.
The package consists of the local lemmas
\cref{paper3:lem:scale-rigid-bounded-limit-exit,paper3:lem:scale-rigid-pressure-modulation-stability,paper3:lem:scale-rigid-chart-comparability,paper3:lem:scale-rigid-two-sided-ckn-transfer,paper3:lem:scale-rigid-ckn-in-gauge,paper3:lem:scale-rigid-scaled-pressure-decay,paper3:lem:scale-rigid-dyadic-ckn-persistence,paper3:lem:scale-rigid-exterior-pressure-routing,paper3:lem:scale-rigid-pressure-core-activation,paper3:lem:scale-rigid-active-core-capture,paper3:lem:scale-rigid-no-subconcentration-bounded,paper3:lem:scale-rigid-degenerate-jacobian,paper3:thm:scale-rigid-single-core-compatibility}
and the main decomposition
\cref{paper3:thm:scale-rigid-stratum-decomposition}.  The final
\cref{paper3:prop:scale-rigidity-discharged} then derives the conclusion of
\cref{paper3:ass:scale-rigidity} from these results.  All proofs use only the
local Type II definition \cref{paper1:def:typeII}, the represented variables of
\cref{paper2:thm:C}, the local pressure decomposition of
\cref{paper2:thm:pressure-decomp}, the modulation \(L^\infty\) control of
\cref{paper2:thm:B}, the same-point and separated-point reductions
\cref{paper3:thm:same-point-reduction,paper3:thm:separated-point-reduction},
the minimal bad-configuration extraction
\cref{paper3:thm:minimal-bad-extraction}, the single-core exclusion
\cref{paper1:thm:single-core-criterion}, the terminal CKN regularity criterion
\cref{paper8:thm:terminal-ckn-regularity}, and Banach--Alaoglu plus
Calder\'on--Zygmund on fixed cylinders \cite{CalderonZygmund1952,Stein1970}.
No global PDE bound, no Liouville theorem, and no profile-decomposition
classification are invoked.
\end{remark}

\subsubsection{Local stratification of scale-rigid branches}

We now identify the named scale-rigid local stratum and prove that every
branch placed there by
\cref{paper3:thm:same-point-reduction,paper3:thm:separated-point-reduction}
either exits the local Type II class through
\cref{paper1:def:typeII} or relaxes back to one of the already excluded local
strata.

\begin{definition}[Scale-rigid local branch]
\label{paper3:def:scale-rigid-branch}
A scale-rigid local branch is a represented branch
\((V_n,P_n,a_n,b_n)\) on a fixed compact selected cylinder \(Q_{R_*}\),
produced by Step~3 or Step~4 of \cref{paper3:thm:same-point-reduction} or by
\cref{paper3:thm:separated-point-reduction}, such that:
\begin{enumerate}[label=\textup{(R\arabic*)}]
\item the compact-cylinder suitable estimates are uniformly bounded,
      \[
        A_{V_n}(R_*)+E_{V_n}(R_*)+C_{V_n}(R_*)+D_{P_n}(R_*)\le M;
      \]
\item the branch carries positive local CKN concentration on a fixed inner
      cylinder \(Q_{r_*}\Subset Q_{R_*}\),
      \[
        C_{V_n}(r_*)+D_{P_n}(r_*)\ge\eta_0;
      \]
\item the modulation coefficients are bounded on the selected cylinder,
      \[
        \|a_n\|_{L^\infty(-R_*^2,0)}+\|b_n\|_{L^\infty(-R_*^2,0)}\le M;
      \]
\item the relative scale-center selection between the retained core and any
      previously selected core is rigid in one of the senses produced by
      \cref{paper3:thm:same-point-reduction}: either the relative
      scale-center Jacobian degenerates along the subsequence (the
      gauge-degenerate alternative of
      \cref{paper3:def:local-multibubble}\,(v)), or the inner core is the
      innermost member of a same-point cascade in which all outer profiles
      have been absorbed into the perturbative remainder by
      \cref{paper3:lem:outer-vanish}.
\end{enumerate}
\end{definition}

\begin{remark}[Inputs (R1)--(R3) are automatic on the selected cylinder]
\label{paper3:rem:scale-rigid-inputs-automatic}
The bounds (R1)--(R3) are not extra hypotheses but are inherited from the
preceding reductions.  Bound (R1) is the compact-cylinder hypothesis carried
through \cref{paper3:thm:minimal-bad-extraction}; if it failed,
\cref{paper3:lem:compactness-failure} would route the branch to the
multicore alternative, not to the scale-rigid stratum.  Bound (R2) is the
positive concentration capture in \cref{paper3:def:active-family}.  Bound
(R3) is the selected-cylinder modulation control of
\cref{paper2:thm:B}, which is part of the represented-variable output of
\cref{paper2:thm:C}.  No additional analytic input is needed to enter the
scale-rigid stratum.
\end{remark}

\begin{lemma}[Bounded selected-window limits exit Type II]
\label{paper3:lem:scale-rigid-bounded-limit-exit}
Let \((V_n,P_n,a_n,b_n)\) be a represented sequence on a compact selected
cylinder \(Q_r\), and assume that, after passing to a selected-window
subsequence, there is \(V_*\in L^\infty(Q_r)\) with \(V_*\not\equiv0\) such
that \(V_n\to V_*\) strongly in \(L^3(Q_r)\).  Then the original sequence is
not a local Type II sequence in the sense of \cref{paper1:def:typeII}.
\end{lemma}

\begin{proof}
The hypothesis is exactly the statement that the subsequence has a nonzero
bounded selected-window limit on \(Q_r\) in the sense of
\cref{paper1:def:bounded-window-limit}.  By \cref{paper1:def:typeII},
a local Type II sequence has no selected-window subsequence with a nonzero
bounded selected-window limit on any compact cylinder.  Hence the sequence
fails to be a local Type II sequence.
\end{proof}

\begin{lemma}[Pressure and modulation stability under rigid-frame limits]
\label{paper3:lem:scale-rigid-pressure-modulation-stability}
Let \((V_n,P_n,a_n,b_n)\) be a scale-rigid local branch on \(Q_{R_*}\) in the
sense of \cref{paper3:def:scale-rigid-branch}.  Fix \(\rho\in(r_*,R_*)\).
After passing to a subsequence:
\begin{enumerate}[label=\textup{(\roman*)}]
\item \(V_n\rightharpoonup V_*\) weakly in
      \(L^3(Q_\rho)\cap L^2_t H^1_x(Q_\rho)\), and
      \(V_n\to V_*\) strongly in \(L^3(Q_\rho')\) for every
      \(\rho'\in(r_*,\rho)\);
\item the mean-normalized pressure
      \(\widetilde P_n:=P_n-\overline{P_n}_{B_\rho}\) converges weakly in
      \(L^{3/2}(Q_\rho)\) to a limit \(P_*\) such that
      \(-\Delta P_*=\partial_i\partial_j(V_{*,i}V_{*,j})\) on \(Q_\rho\) in
      the sense of distributions modulo a function of time;
\item \(a_n\rightharpoonup^* a_*\) and \(b_n\rightharpoonup^* b_*\)
      weakly-\(*\) in \(L^\infty(-R_*^2,0)\), with
      \(\|a_*\|_\infty+\|b_*\|_\infty\le M\);
\item the limit \((V_*,P_*,a_*,b_*)\) is a suitable weak solution of the
      represented Navier--Stokes equation \eqref{paper1:eq:represented-ns} on
      \(Q_\rho\), with \(V_*\in L^3(Q_\rho)\cap L^2_tH^1_x(Q_\rho)\).
\end{enumerate}
This lemma is only a compactness and stability statement.  It does not assert
boundedness of \(V_*\).  Bounded selected-window limits are obtained below from
the separate no-subconcentration/CKN argument.
\end{lemma}

\begin{proof}
\textit{(i) Compactness of \(V_n\).}  By (R1), the family \((V_n)\) is bounded
in \(L^3(Q_\rho)\) (from \(C_{V_n}(R_*)\)), in \(L^2_t H^1_x(Q_\rho)\) (from
\(A_{V_n}(R_*)+E_{V_n}(R_*)\)), and in \(L^\infty_t L^2_x(Q_\rho)\) (from
\(A_{V_n}(R_*)\)).  Banach--Alaoglu yields the weak limits in
\(L^3(Q_\rho)\) and \(L^2_t H^1_x(Q_\rho)\) along a subsequence.  The
represented Navier--Stokes equation \eqref{paper1:eq:represented-ns} bounds
\(\partial_s V_n\) in \(L^{4/3}_t H^{-1}_x(Q_\rho)\) by Sobolev product
estimates on the bounded cylinder, since each term on the right-hand side is
controlled by \(A_{V_n}(R_*)\), \(E_{V_n}(R_*)\), \(C_{V_n}(R_*)\),
\(D_{P_n}(R_*)\), and \(\|a_n\|_\infty+\|b_n\|_\infty\le M\); the modulation
correction \(a_n(s)(V_n+y\cdot\nabla V_n)+b_n(s)\cdot\nabla V_n\) is bounded
in \(L^2_t L^{6/5}_x(Q_\rho)\subset L^{4/3}_t H^{-1}_x(Q_\rho)\) on the
bounded cylinder.  Hence Aubin--Lions \cite{Aubin1963,Simon1987} applied on
the fixed cylinder \(Q_\rho\) yields strong convergence
\(V_n\to V_*\) in \(L^3(Q_{\rho'})\) for every \(\rho'\in(r_*,\rho)\).

\textit{(ii) Pressure stability.}  By (R1) the local pressure quantity
\(D_{P_n}(R_*)\) is bounded, equivalently \(\widetilde P_n\) is bounded in
\(L^{3/2}(Q_\rho)\) by the local Calder\'on--Zygmund mean-normalization
estimate \cite{CalderonZygmund1952,Stein1970} applied to the local pressure
identity in \cref{paper2:thm:pressure-decomp}.  Banach--Alaoglu yields
\(\widetilde P_n\rightharpoonup P_*\) weakly in \(L^{3/2}(Q_\rho)\).
Pass to the limit in the local pressure identity
\[
  -\Delta\widetilde P_n
  =\partial_i\partial_j(V_{n,i}V_{n,j})
  -\partial_i\partial_j(\overline{V_{n,i}V_{n,j}}_{B_\rho})
\]
on \(Q_\rho\).  The right-hand side converges in
\(\mathcal D'(Q_\rho)\) by the strong \(L^3\) convergence of (i) (which
implies weak \(L^{3/2}\) convergence of \(V_{n,i}V_{n,j}\) by H\"older).
This proves (ii).

\textit{(iii) Modulation limits.}  By (R3),
\(a_n,b_n\) are bounded in \(L^\infty(-R_*^2,0)\); Banach--Alaoglu applied to
the dual of \(L^1\) yields the weak-\(*\) limits with the same bound \(M\).

\textit{(iv) Suitable weak solution status.}  Each
\((V_n,P_n)\) is suitable on \(Q_\rho\) (by (R1) and the represented-equation
output of \cref{paper2:thm:C}).  The local energy inequality and the
distributional momentum equation pass to the limit using (i)--(iii) and
H\"older on the fixed cylinder \(Q_\rho\): the dissipation term is lower
semicontinuous under weak \(L^2_t H^1_x\) convergence, the convection term
converges by strong \(L^3(Q_{\rho'})\) convergence, and the pressure-velocity
pairing
\(\int \widetilde P_n V_n\cdot\nabla\chi\) converges by weak \(L^{3/2}\)
times strong \(L^3\) on the fixed cylinder.  The modulation correction
converges by weak-\(*\) \(L^\infty\) for \((a_n,b_n)\) against the strong
\(L^1_t\) convergence of \((V_n,y\cdot\nabla V_n,\nabla V_n)\) tested against
fixed test functions on \(Q_\rho\).  This proves the represented suitable
limit in the displayed local energy class.  No \(L^\infty\) bound is inferred
from the compact-cylinder estimates.
\end{proof}

\begin{lemma}[Compact chart comparability in scale-rigid gauges]
\label{paper3:lem:scale-rigid-chart-comparability}
Let \((V,P,a,b)\) be a represented branch on \(Q_{2\rho}(z_0)\) obtained from
the repaired-gauge chart of \cref{paper2:thm:C}.  Assume that the associated
chart \((\Lambda,X)\) satisfies
\[
  a=-\Lambda',\qquad b=-X',\qquad
  \|a\|_{L^\infty(I_{2\rho})}+\|b\|_{L^\infty(I_{2\rho})}\le M
\]
on the time interval \(I_{2\rho}\) of \(Q_{2\rho}(z_0)\), and set
\(\Lambda_0:=\Lambda(s_0)>0\) for \(z_0=(y_0,s_0)\).  There exists
\(\rho_c=\rho_c(M,\Lambda_0,\rho)>0\) such that, for every
\(0<r\le\rho_c\), the chart on \(Q_{2r}(z_0)\) obeys
\[
  \frac{\Lambda_0}{2}\le\Lambda(s)\le\frac{3\Lambda_0}{2},
  \qquad |X(s)-X(s_0)|\le 4Mr^2,
  \qquad s\in(s_0-4r^2,s_0].
\]
After a harmless concentric shrink, if needed, the physical image remains
inside the support buffer of \cref{paper2:prop:domain}.  Consequently
\(\cref{paper2:prop:domain,paper2:prop:local-transfer}\) apply on the retained
subcylinder with
\[
  \Lambda_{\min}=\Lambda_0/2,\qquad \Lambda_{\max}=3\Lambda_0/2,
\]
and represented and physical CKN quantities on comparable cylinders are
equivalent with constants depending only on
\(M,\Lambda_0,r,\rho\) and the fixed buffer.
\end{lemma}

\begin{proof}
The identities \(a=-\Lambda'\) and \(b=-X'\) are part of the exact repaired
change of variables in
\cref{paper2:lem:change-id,paper2:thm:A,paper2:thm:C}.  Hence, for
\(|s-s_0|\le4r^2\),
\[
  |\Lambda(s)-\Lambda_0|
  \le \int_s^{s_0}|a(\sigma)|\,d\sigma
  \le 4Mr^2,\qquad
  |X(s)-X(s_0)|
  \le \int_s^{s_0}|b(\sigma)|\,d\sigma
  \le 4Mr^2 .
\]
Choose
\[
  \rho_c\le \min\{\rho/4,(\Lambda_0/(16M))^{1/2}\}
\]
when \(M>0\), and choose \(\rho_c=\rho/4\) when \(M=0\).  Then
\(|\Lambda(s)-\Lambda_0|\le\Lambda_0/4\), which implies the displayed
lower-upper scale bounds after enlarging the upper constant to
\(3\Lambda_0/2\).  The center estimate is the preceding integral bound.  A
final concentric shrink, depending only on the fixed distance to the boundary
of the selected compact window, keeps the image inside the domain buffer.
\Cref{paper2:prop:domain,paper2:prop:local-transfer} then give the stated
comparison of the represented and physical quantities.  No solution
boundedness, and no vanishing dissipation, is used.
\end{proof}

\begin{lemma}[Two-sided local CKN transfer in comparable charts]
\label{paper3:lem:scale-rigid-two-sided-ckn-transfer}
Let \((V,P)\) be related to a physical suitable pair \((u,p)\) by the repaired
scale-translation chart on \(Q_{2r}(z_0)\), and assume the chart comparability
conclusion of \cref{paper3:lem:scale-rigid-chart-comparability} there:
\[
  0<\Lambda_{\min}\le \Lambda(s)\le \Lambda_{\max}<\infty,
  \qquad |X(s)-X(s_0)|\le M r^2 .
\]
Set
\[
  \mathcal C_V(z_0,r):=
  r^{-2}\int_{Q_r(z_0)}|V|^3
  +r^{-2}\int_{Q_r(z_0)}|P-(P)_{B_r}|^{3/2},
\]
and define the physical CKN quantity \(\mathcal C_u(x,t,\rho)\) analogously,
with pressure normalized by spatial averages.  There are constants
\[
  N_{\rm tr}<\infty,\qquad c_{\rm tr}\in(0,1),\qquad C_{\rm tr}<\infty,
\]
depending only on \(M,\Lambda_{\min},\Lambda_{\max}\) and the fixed support
buffer, and physical cylinders
\[
  Q^{\rm ph}_{\ell}:=
  B_{\rho_\ell}(x_\ell)\times(t_\ell-\rho_\ell^2,t_\ell),
  \qquad
  c_{\rm tr}\Lambda_{\min}r\le \rho_\ell\le C_{\rm tr}\Lambda_{\max}r,
  \qquad 1\le \ell\le N_{\rm tr},
\]
such that the physical image of \(Q_{c_{\rm tr}r}(z_0)\) is contained in
\(\cup_\ell Q^{\rm ph}_{\ell}\), each \(Q^{\rm ph}_{\ell}\) is contained in the
physical image of \(Q_r(z_0)\), and
\[
  \mathcal C_u(x_\ell,t_\ell,\rho_\ell)
  \le C_{\rm tr}\mathcal C_V(z_0,r)
  \qquad (1\le \ell\le N_{\rm tr}).
\]
Conversely, if a finite physical cylinder cover of the image of \(Q_r(z_0)\)
with comparable radii satisfies
\(\mathcal C_u(x_\ell,t_\ell,\rho_\ell)\le\varepsilon\) for every member of the
cover, then
\[
  \mathcal C_V(z_0,c_{\rm tr}r)\le C_{\rm tr}N_{\rm tr}\varepsilon .
\]
\end{lemma}

\begin{proof}
The chart is
\[
  x=X(s)+\Lambda(s)y,\qquad dt=\Lambda(s)^2\,ds,
  \qquad V(y,s)=\Lambda(s)u(x,t),\qquad P(y,s)=\Lambda(s)^2p(x,t).
\]
The bounds on \(\Lambda\) and \(X\) imply that the image of a represented
parabolic cylinder is trapped between parabolic cylinders whose radii are
comparable to \(\Lambda(s_0)r\).  A standard Besicovitch covering with the
bounded eccentricity supplied by
\(\Lambda_{\min},\Lambda_{\max}\) gives the finite cover and the bounded
overlap constant \(N_{\rm tr}\).

For the velocity term the change of variables gives, on every comparison
cylinder,
\[
  \rho_\ell^{-2}\int_{Q^{\rm ph}_{\ell}}|u|^3\,dx\,dt
  \le C_{\rm tr} r^{-2}\int_{Q_r(z_0)}|V|^3\,dy\,ds .
\]
For the pressure term we use the invariant seminorm
\(\inf_{c(s)}\int |P-c(s)|^{3/2}\).  Spatial averages are equivalent to this
seminorm up to the universal factor \(2^{3/2}\), because for every ball \(B\)
and \(q\ge1\),
\[
  \int_B |F-F_B|^q\le 2^q\inf_c\int_B|F-c|^q .
\]
The same Jacobian identity therefore gives the displayed physical estimate for
the pressure oscillation after subtracting spatial means.  Applying the same
argument to the inverse chart, whose scale is bounded between
\(\Lambda_{\max}^{-1}\) and \(\Lambda_{\min}^{-1}\) on the comparison cover,
gives the converse estimate.  No direction of
\cref{paper2:prop:local-transfer} is used as a black box here; both estimates
come from the exact change of variables and the pressure oscillation seminorm.
\end{proof}

\begin{lemma}[CKN regularity in comparable repaired gauges]
\label{paper3:lem:scale-rigid-ckn-in-gauge}
Let \((V,P,a,b)\) be a represented suitable weak solution on \(Q_{2r}(z_0)\)
obtained from the repaired-gauge chart of \cref{paper2:thm:C}.  Assume that the
chart satisfies the comparability conclusion of
\cref{paper3:lem:scale-rigid-chart-comparability} on this cylinder, with
constants \(\Lambda_{\min},\Lambda_{\max}\), and that
\[
  \|a\|_{L^\infty}+\|b\|_{L^\infty}\le M
\]
there.  There are constants
\[
  \varepsilon_{\rm rg}
  =\varepsilon_{\rm rg}(M,\Lambda_{\min},\Lambda_{\max})>0,
  \qquad
  c_{\rm rg}=c_{\rm rg}(M,\Lambda_{\min},\Lambda_{\max})\in(0,1),
\]
such that if
\[
  r^{-2}\int_{Q_r(z_0)} |V|^3
  +r^{-2}\int_{Q_r(z_0)} |P-(P)_{B_r}|^{3/2}
  \le\varepsilon_{\rm rg},
\]
then \(V\) is bounded on \(Q_{c_{\rm rg}r}(z_0)\), with a bound depending only
on \(M,\Lambda_{\min},\Lambda_{\max},r\), and the preceding displayed
quantity.
\end{lemma}

\begin{proof}
Apply the forward direction of
\cref{paper3:lem:scale-rigid-two-sided-ckn-transfer} to \(Q_r(z_0)\).  Choose
\(\varepsilon_{\rm rg}\) so small that every physical cylinder in the resulting
finite comparison cover has CKN quantity below the threshold
\(\varepsilon_{\rm CKN}\) from
\cref{paper8:thm:terminal-ckn-regularity}.  That theorem gives boundedness of
the physical velocity on smaller physical cylinders.  Pulling those estimates
back through the same comparable repaired chart gives the asserted
\(L^\infty\) bound for \(V\) on \(Q_{c_{\rm rg}r}(z_0)\).  Thus boundedness is
obtained only from epsilon regularity after two-sided chart transfer and small
CKN quantity, not from compactness or gauge degeneracy.
\end{proof}

\begin{lemma}[Scaled local pressure decay]
\label{paper3:lem:scale-rigid-scaled-pressure-decay}
Let \((V,P)\) be a represented suitable pair on \(Q_{2r}(z_0)\), and assume the
localized pressure decomposition of \cref{paper2:thm:pressure-decomp} is
available with a cutoff supported in \(B_{2r}(y_0)\) and equal to \(1\) on
\(B_r(y_0)\).  For
\[
  C_V(\rho;z_0):=\rho^{-2}\int_{Q_\rho(z_0)}|V|^3,\qquad
  D_P(\rho;z_0):=\rho^{-2}\int_{Q_\rho(z_0)}
       |P-(P)_{B_\rho}|^{3/2},
\]
there is a universal constant \(C_{\rm dec}\) such that, whenever
\(0<\theta\le1/4\),
\[
  D_P(\theta r;z_0)
  \le C_{\rm dec}\theta^{5/2}D_P(r;z_0)
     +C_{\rm dec}\theta^{-2}C_V(2r;z_0).
\]
\end{lemma}

\begin{proof}
Write \(z_0=(y_0,s_0)\).  For a.e. time
\(s\in(s_0-r^2,s_0]\), apply \cref{paper2:thm:pressure-decomp} on
\(B_r(y_0)\) with cutoff supported in \(B_{2r}(y_0)\):
\[
  P(\cdot,s)=P_{\rm loc}(\cdot,s)+H(\cdot,s),
  \qquad -\Delta H(\cdot,s)=0\quad\text{in }B_r(y_0).
\]
The Calder\'on--Zygmund part satisfies
\[
  \int_{B_{\theta r}(y_0)}|P_{\rm loc}|^{3/2}
  \le C\int_{B_{2r}(y_0)}|V|^3 .
\]
After integrating over the shorter time interval and multiplying by
\((\theta r)^{-2}\), this contributes at most
\(C\theta^{-2}C_V(2r;z_0)\).

For the harmonic part, the interior oscillation estimate for harmonic
functions gives, for \(q=3/2\),
\[
  \int_{B_{\theta r}(y_0)}
  |H-(H)_{B_{\theta r}}|^q
  \le C\theta^{3+q}
  \int_{B_r(y_0)}|H-(H)_{B_r}|^q .
\]
Since \(3+q=9/2\), scaling the cylinder by \((\theta r)^{-2}\) leaves the
factor \(\theta^{5/2}\).  Moreover
\[
  \int_{B_r}|H-(H)_{B_r}|^{3/2}
  \le C\int_{B_r}|P-(P)_{B_r}|^{3/2}
     +C\int_{B_r}|P_{\rm loc}|^{3/2},
\]
and the last term is again bounded by \(C\int_{B_{2r}}|V|^3\).  Absorbing the
resulting \(C\theta^{5/2}C_V(2r;z_0)\) into
\(C\theta^{-2}C_V(2r;z_0)\) gives the displayed estimate.  Spatial averages may
be replaced by the minimizing constants at the cost of a universal factor, as
in \cref{paper3:lem:scale-rigid-two-sided-ckn-transfer}.
\end{proof}

\begin{lemma}[Dyadic persistence of CKN concentration]
\label{paper3:lem:scale-rigid-dyadic-ckn-persistence}
Let \((V,P)\) be suitable on a compact cylinder \(Q_\rho\), and let
\[
  \mathcal C_{V,P}(z,r):=
  r^{-2}\int_{Q_r(z)}|V|^3
  +r^{-2}\int_{Q_r(z)}|P-(P)_{B_r}|^{3/2}.
\]
If \(z_*\in Q_\rho\) is a Caffarelli--Kohn--Nirenberg concentration point in
the sense that
\[
  \limsup_{r\downarrow0}\mathcal C_{V,P}(z_*,r)\ge\kappa_*>0,
\]
then for every \(0<\theta<1\) and every sufficiently small outer radius
\(r_0\) with \(Q_{2r_0}(z_*)\Subset Q_\rho\), there exist
\(z_m\to z_*\), integers \(L_m\to\infty\), and dyadic radii
\(\rho_m=\theta^{L_m}r_0\) such that
\[
  Q_{\rho_m}(z_m)\Subset Q_{r_0}(z_*),\qquad
  \mathcal C_{V,P}(z_m,\rho_m)\ge c_\theta\kappa_* ,
\]
where \(c_\theta>0\) depends only on \(\theta\).
\end{lemma}

\begin{proof}
Choose radii \(\sigma_m\downarrow0\) and points \(z_m\to z_*\) with
\(Q_{\sigma_m}(z_m)\Subset Q_{r_0}(z_*)\) and
\(\mathcal C_{V,P}(z_m,\sigma_m)\ge\kappa_*/2\).  Pick \(L_m\) so that
\[
  \theta^{L_m+1}r_0<\sigma_m\le \theta^{L_m}r_0=:\rho_m .
\]
Then \(L_m\to\infty\) and \(\rho_m\le\theta^{-1}\sigma_m\).  The velocity term
on \(Q_{\rho_m}(z_m)\) controls the velocity term on \(Q_{\sigma_m}(z_m)\) with
loss at most \(\theta^2\), because the normalization changes by
\((\sigma_m/\rho_m)^2\).  For the pressure term, spatial averages are compared
through the minimizing constant:
\[
  \int_{B_{\sigma_m}}|P-(P)_{B_{\sigma_m}}|^{3/2}
  \le 2^{3/2}\int_{B_{\sigma_m}}|P-c|^{3/2}
  \le 2^{3/2}\int_{B_{\rho_m}}|P-c|^{3/2}.
\]
Taking \(c=(P)_{B_{\rho_m}}\), integrating in time, and accounting for the same
normalization loss gives
\[
  \mathcal C_{V,P}(z_m,\rho_m)\ge 2^{-5/2}\theta^2
  \mathcal C_{V,P}(z_m,\sigma_m).
\]
Thus the conclusion holds with \(c_\theta=2^{-7/2}\theta^2\).
\end{proof}

\begin{lemma}[Exterior pressure persistence is a structural exit]
\label{paper3:lem:scale-rigid-exterior-pressure-routing}
Let a scale-rigid branch have a pressure lower bound on shrinking selected
cylinders that is not accounted for by the local source
\(\partial_i\partial_j(V_iV_j)\) in any compact selected frame after the
decomposition of \cref{paper2:thm:pressure-decomp}.  Then the branch enters the
separated/exterior structural alternative of
\cref{paper3:def:local-multibubble}\,(iv), and it does not remain in the
residual scale-rigid stratum.
\end{lemma}

\begin{proof}
The unaccounted part of the pressure is harmonic in the selected compact frame
and is generated by sources at positive distance in the corresponding physical
variables.  If the exterior velocity source has a terminal CKN concentration on
one of the annuli separating it from the selected core, then
\cref{paper8:thm:exterior-stratification} gives an exterior
concentration branch.  In the local variables this is a separated active
physical point, hence the separated-point alternative
\cref{paper3:def:local-multibubble}\,(iv) and the reduction
\cref{paper3:thm:separated-point-reduction} apply.

If no such exterior concentration exists on the separating annuli, then
\cref{paper8:thm:exterior-regularity-no-concentration} gives regular exterior
control and
\cref{paper8:lem:exterior-pressure-vanishes} gives convergence of the exterior
pressure contribution modulo functions of time on the selected compact frame.
This contradicts the assumed persistent pressure lower bound there.  Thus the
only non-vanishing exterior-pressure case is already a separated/exterior
structural exit.
\end{proof}

\begin{lemma}[Persistent CKN mass activates an \(L^3\) core or exits structurally]
\label{paper3:lem:scale-rigid-pressure-core-activation}
Let \((V_n,P_n)\) be a compact-cylinder scale-rigid sequence on \(Q_{2R}\) with
\[
  \sup_n\bigl(C_{V_n}(2R)+D_{P_n}(2R)\bigr)\le M_0,
\]
and assume the pressure decomposition of
\cref{paper2:thm:pressure-decomp} is available on \(Q_{2R}\).  Fix
\(\kappa_0>0\) and choose
\(\theta\in(0,1/8)\) so that
\[
  C_{\rm dec}\theta^{5/2}\le \frac14 .
\]
Then there are constants \(\eta_{\rm act}>0\) and \(L_0\in\mathbb N\),
depending only on \(M_0,R,\kappa_0,\theta\), the fixed-scale pressure bound
below, and the pressure-decomposition constants, with the following property.
Suppose \(0<R_0<R\), \(z_n\to z\),
\(Q_{2R_0}(z_n)\Subset Q_{2R}\) for all large \(n\), and suppose the fixed
outer scale used for the iteration has a uniform local pressure bound
\[
  D_{P_n}(R_0;z_n)\le D_0<\infty .
\]
Let \(L_n\to\infty\), \(\rho_n:=\theta^{L_n}R_0\), and assume
\[
  \mathcal C_n(\rho_n;z_n):=
  \rho_n^{-2}\int_{Q_{\rho_n}(z_n)} |V_n|^3
  +\rho_n^{-2}\int_{Q_{\rho_n}(z_n)}
       |P_n-(P_n)_{B_{\rho_n}}|^{3/2}
  \ge \kappa_0
\]
for all large \(n\).  Then, after passing to a subsequence, either:
\begin{enumerate}[label=\textup{(\alph*)}]
\item for some integers \(0\le j_n<L_n\), with
      \(s_n:=2\theta^{j_n}R_0\),
      \[
        s_n^{-2}\int_{Q_{s_n}(z_n)}|V_n|^3\ge\eta_{\rm act}^3 ,
      \]
      so the rescaled branch at scale \(s_n\) contains an active \(L^3\) core;
\item the retained pressure lower bound is exterior pressure persistence, hence
      the branch is assigned to the separated/exterior structural alternative by
      \cref{paper3:lem:scale-rigid-exterior-pressure-routing}.
\end{enumerate}
\end{lemma}

\begin{proof}
Choose \(L_0\) so large that
\[
  (1/4)^{L_0}D_0\le \kappa_0/16,
\]
and choose \(\eta_{\rm act}>0\) so small that
\[
  C_{\rm dec}\theta^{-2}\sum_{m=0}^{\infty}(1/4)^m\eta_{\rm act}^3
  \le \kappa_0/16,
  \qquad
  4\theta^{-2}\eta_{\rm act}^3\le \kappa_0/4 .
\]
For large \(n\), \(L_n\ge L_0\).  Suppose item \textup{(a)} fails.  Then for
each \(0\le j<L_n\),
\[
  C_{V_n}(2\theta^jR_0;z_n)<\eta_{\rm act}^3 .
\]
Iterating \cref{paper3:lem:scale-rigid-scaled-pressure-decay} from the fixed
outer scale \(R_0\) down to \(\theta^{L_n}R_0=\rho_n\) gives
\[
  D_{P_n}(\rho_n;z_n)
  \le (1/4)^{L_n}D_{P_n}(R_0;z_n)
     +C_{\rm dec}\theta^{-2}
      \sum_{j=0}^{L_n-1}(1/4)^{L_n-1-j}
      C_{V_n}(2\theta^jR_0;z_n)
  \le \kappa_0/8 .
\]
Here the initial pressure is bounded by \(D_0\) because the iteration starts
from this fixed compact scale, not from the shrinking target scale.
Moreover \(Q_{\rho_n}(z_n)\subset Q_{2\theta^{L_n-1}R_0}(z_n)\), so failure of
item \textup{(a)} at \(j=L_n-1\) gives
\[
  \rho_n^{-2}\int_{Q_{\rho_n}(z_n)}|V_n|^3
  \le 4\theta^{-2}\eta_{\rm act}^3\le \kappa_0/4 .
\]
Hence \(\mathcal C_n(\rho_n;z_n)<\kappa_0\), contradicting the assumed lower
bound, unless the pressure term used in the lower bound is not the local
pressure governed by the compact-frame decomposition.  The latter possibility
is exactly exterior pressure persistence and is item \textup{(b)} by
\cref{paper3:lem:scale-rigid-exterior-pressure-routing}.
\end{proof}

\begin{lemma}[Activated cores are captured by the finite active ledger]
\label{paper3:lem:scale-rigid-active-core-capture}
Let a scale-rigid branch inherit an active profile family in the sense of
\cref{paper3:def:active-family}, with finite critical mass \(M_*\) and
concentration capture.  Suppose
\cref{paper3:lem:scale-rigid-pressure-core-activation} produces represented
scales \(s_n>0\) and centers \(z_n\) such that
\[
  s_n^{-2}\int_{Q_{s_n}(z_n)}|V_n|^3\ge \eta_{\rm act}^3 .
\]
After lowering the active threshold once, if necessary, to
\(\eta_*^{\rm new}:=\min\{\eta_*,\eta_{\rm act}/2\}\), exactly one of the
following occurs after passing to a subsequence:
\begin{enumerate}[label=\textup{(\alph*)}]
\item the activated core is already represented by an existing active profile
      or compound profile in the inherited family;
\item it is absorbed into the selected terminal frame, so it does not create an
      independent residual obstruction;
\item it is assigned to the profile-completeness obligation rather than to the
      scale-rigid residual branch.
\end{enumerate}
In particular, under the stated concentration-capture hypothesis, pressure-core
activation cannot create an unregistered active object inside the residual
scale-rigid branch.
\end{lemma}

\begin{proof}
Rescale the branch around \((z_n,s_n)\) inside the represented window; in
physical variables the corresponding scale is the selected physical scale
multiplied by \(s_n\), hence still tends to zero along the Type II sequence.
The displayed lower bound is exactly the active-frame condition in
\cref{paper3:def:active-frame} on \(Q_1\), with threshold \(\eta_{\rm act}\).
The concentration-capture clause in
\cref{paper3:def:active-family} says that a compact terminal-window sequence
with positive \(L^3\) mass is either seen by the active family or is absorbed
into the selected terminal frame; it cannot remain in the perturbative
remainder.  Hence alternatives \textup{(a)} and \textup{(b)} hold unless the
frame is genuinely absent from the captured profile family.  That last
possibility is not a scale-rigid residual mechanism; it is precisely failure of
the terminal active-family completeness/capture input and is assigned to that
separate profile-completeness obligation.  Therefore the scale-rigid refinement
below only uses active objects that already belong to the finite captured
family, together with objects absorbed into the selected terminal frame.
\end{proof}

\begin{lemma}[No sub-concentration gives a bounded selected-window limit]
\label{paper3:lem:scale-rigid-no-subconcentration-bounded}
Let \((V_n,P_n,a_n,b_n)\) be a scale-rigid local branch and let
\((V_*,P_*)\) be a compact limit supplied by
\cref{paper3:lem:scale-rigid-pressure-modulation-stability} on
\(Q_\rho\).  Then, after shrinking once to a compact subcylinder
\(Q_{\rho'}\Subset Q_\rho\), exactly one of the following alternatives holds:
\begin{enumerate}[label=\textup{(\alph*)}]
\item \((V_*,P_*)\) has no Caffarelli--Kohn--Nirenberg concentration point in
      \(Q_{\rho'}\).  Then \(V_*\in L^\infty(Q_{\rho''})\) for every
      \(Q_{\rho''}\Subset Q_{\rho'}\), and the selected subsequence has a
      nonzero bounded selected-window limit on some such \(Q_{\rho''}\).
\item \((V_*,P_*)\) has a Caffarelli--Kohn--Nirenberg concentration point in
      \(Q_{\rho'}\).  Then the original branch either contains a further
      captured active \(L^3\) concentration core after rescaling around that
      point, the pressure mass is exterior and the branch is already in the
      separated/exterior structural alternative, or the core is assigned to the
      separate profile-completeness obligation.  According to its relative
      scale-center position, every captured core enters one of the already named
      multicore, same-point cascade, or separated/exterior structural
      alternatives, rather than remaining a residual scale-rigid branch.
\end{enumerate}
\end{lemma}

\begin{proof}
Let \(\varepsilon_{\rm rg}\) be the represented-gauge regularity threshold from
\cref{paper3:lem:scale-rigid-ckn-in-gauge}.  Before using that threshold,
apply \cref{paper3:lem:scale-rigid-chart-comparability} on each retained
compact subcylinder.  The chart stability in \cref{paper2:prop:stability},
together with the \(L^\infty\) modulation bound (R3), gives a limiting chart
\((\Lambda_*,X_*)\) with \(a_*=-\Lambda_*'\) and \(b_*=-X_*'\) on the retained
subwindow.  Thus the compactness of the subcylinder and (R3) give a uniform
lower-upper scale comparison after one shrink, so the constants in
\cref{paper3:lem:scale-rigid-ckn-in-gauge} depend only on the retained chart
geometry and not on any unproved boundedness of \(V_n\) or \(V_*\).

Assume first that there is no CKN concentration point of \((V_*,P_*)\) in
\(Q_{\rho'}\).  For every point \(z\in Q_{\rho'}\) there is a parabolic
cylinder \(Q_{r_z}(z)\Subset Q_\rho\) such that the scale-invariant
velocity-pressure quantity of \((V_*,P_*)\) on \(Q_{r_z}(z)\) is below
\(\varepsilon_{\rm rg}\).  A finite subcover of each
\(Q_{\rho''}\Subset Q_{\rho'}\), followed by
\cref{paper3:lem:scale-rigid-ckn-in-gauge} on the cover, gives
\(V_*\in L^\infty(Q_{\rho''})\).

It remains to check that the bounded limit is nonzero.  If \(V_*\equiv0\) on
the retained inner cylinder, then the strong \(L^3\) convergence from
\cref{paper3:lem:scale-rigid-pressure-modulation-stability} and the pressure
stability argument of \cref{paper1:lem:zero-contradiction} imply
\[
  C_{V_n}(r_*)+D_{P_n}(r_*)\to0,
\]
contradicting the retained lower bound (R2) in
\cref{paper3:def:scale-rigid-branch}.  Hence \(V_*\not\equiv0\) on some
smaller compact cylinder.  Since \(V_n\to V_*\) strongly in \(L^3\) there,
the selected subsequence has a nonzero bounded selected-window limit.

Now suppose a CKN concentration point \(z_*\) remains in \(Q_{\rho'}\).  Then
there is \(\kappa_*>0\) such that
\[
  \limsup_{r\downarrow0}\mathcal C_{V_*,P_*}(z_*,r)\ge \kappa_* .
\]
Choose \(R_0>0\) with \(Q_{2R_0}(z_*)\Subset Q_\rho\), and choose
\(\theta\in(0,1/8)\) satisfying the pressure-decay smallness condition in
\cref{paper3:lem:scale-rigid-pressure-core-activation}.  By
\cref{paper3:lem:scale-rigid-dyadic-ckn-persistence}, there are
\(z_m\to z_*\), \(L_m\to\infty\), and
\(\rho_m=\theta^{L_m}R_0\) such that
\[
  \mathcal C_{V_*,P_*}(z_m,\rho_m)\ge c_\theta\kappa_* .
\]
For each fixed \(m\), the strong \(L^3\) convergence of \(V_n\) and weak lower
semicontinuity of the pressure oscillation seminorm in
\(L^{3/2}/\{\text{spatial constants}\}\) imply, along a diagonal subsequence,
\[
  \mathcal C_{V_n,P_n}(z_m,\rho_m)\ge \frac12 c_\theta\kappa_* .
\]
Relabeling the diagonal sequence, apply
\cref{paper3:lem:scale-rigid-pressure-core-activation} with
\(\kappa_0=\frac12c_\theta\kappa_*\) on the compact cylinder where (R1) holds.
The additional fixed-scale pressure bound required by that lemma holds because
\(R_0\) is fixed, \(Q_{2R_0}(z_m)\Subset Q_\rho\) for all large \(m\), and the
mean-normalized pressures are uniformly bounded in \(L^{3/2}(Q_\rho)\) by
\cref{paper3:lem:scale-rigid-pressure-modulation-stability}; comparison of
spatial averages on the fixed nested balls gives
\(D_{P_n}(R_0;z_m)\le D_0(R_0,\rho,M)\).
If the persistent CKN mass is produced by an interior velocity core, then
\cref{paper3:lem:scale-rigid-active-core-capture} either identifies that core
with the finite captured active ledger or assigns it to the separate
profile-completeness obligation, after the one-time lowering of the active
threshold.  If
the mass is carried by harmonic pressure generated outside the selected compact
frame, \cref{paper3:lem:scale-rigid-exterior-pressure-routing} assigns the
branch to the separated/exterior structural alternative, so it exits the
residual scale-rigid stratum.

The relative position of this new core with respect to the retained scale is
exhausted by \cref{paper3:lem:parameter-partition}: it is comparable at the
same point, strictly separated in scale at the same point, separated in
physical center, or locally invisible in the retained frame.  The invisible
case cannot carry the retained positive concentration by the capture clause in
\cref{paper3:def:active-family}.  The remaining cases are precisely the
multicore, same-point cascade, and separated/exterior structural alternatives
already named in the manuscript.  Thus the branch does not remain in the
residual scale-rigid stratum in this case.
\end{proof}

\begin{theorem}[Single-core reduction compatibility]
\label{paper3:thm:scale-rigid-single-core-compatibility}
Let \((V_n,P_n,a_n,b_n)\) be a scale-rigid local branch in the sense of
\cref{paper3:def:scale-rigid-branch}, and assume the branch reduces, after
the rigid-frame limit of
\cref{paper3:lem:scale-rigid-pressure-modulation-stability}, to a single
retained nondegenerate core on a fixed inner cylinder
\(Q_{r_*}\Subset Q_{\rho}\Subset Q_{R_*}\).  Then, after passing to a
subsequence, one of the following alternatives holds:
\begin{enumerate}[label=\textup{(\roman*)}]
\item the branch has a nonzero bounded selected-window limit;
\item the original selected sequence has vanishing localized dissipation on a
      retained compact cylinder, hence is a compact single-core
      vanishing-dissipation branch in the sense of
      \cref{paper1:def:single-core-branch} and is excluded by
      \cref{paper1:thm:single-core-criterion};
\item a further active concentration core is produced, so the branch leaves the
      single-core reduction and enters a multicore, cascade, or
      separated/exterior structural alternative.
\end{enumerate}
\end{theorem}

\begin{proof}
If the original selected sequence has vanishing localized dissipation on some
\(Q_{\rho'}\), \(Q_{r_*}\Subset Q_{\rho'}\Subset Q_{R_*}\), then we first verify
\cref{paper1:def:single-core-branch}.  Suitability and the represented
equation are supplied by \cref{paper2:thm:C}.  Compact-cylinder bounds,
positive selected-core CKN concentration, and bounded modulation are exactly
(R1), (R2), and (R3) of \cref{paper3:def:scale-rigid-branch}.  The
nondegenerate scale-center selection is the hypothesis of this single-core
case.  The vanishing localized dissipation is the present hypothesis on the
original selected sequence, not on a residual after subtracting a limit.  Thus
the branch satisfies every clause of
\cref{paper1:def:single-core-branch}, and
\cref{paper1:thm:single-core-criterion} excludes it.  This is item
\textup{(ii)}.

Apply \cref{paper3:lem:scale-rigid-no-subconcentration-bounded} to the compact
limit on a slightly smaller cylinder.  If the no-subconcentration alternative
holds, item \textup{(i)} follows.  If a further active concentration core is
produced, item \textup{(iii)} follows.
\end{proof}

\begin{lemma}[Degenerate Jacobian resolution]
\label{paper3:lem:scale-rigid-degenerate-jacobian}
Let \((V_n,P_n,a_n,b_n)\) be a represented sequence reaching the
gauge-degenerate alternative \cref{paper3:def:local-multibubble}\,(v),
i.e.\ along the subsequence the relative scale-center Jacobian of the
selected compound or terminal frame degenerates.  Then, after passing to a
further subsequence and reselecting the active center in the sense of
\cref{paper3:def:active-family} on the same physical point, exactly one of
the following four resolutions occurs:
\begin{enumerate}[label=\textup{(D\arabic*)}]
\item the reselection produces an additional active core or active compound
      carrying positive \(L^3\) mass at the fixed active threshold; hence the
      branch enters the multicore, same-point cascade, or separated/exterior
      alternative
      \cref{paper3:def:local-multibubble}\,(ii)--(iv) and is treated by
      \cref{paper3:thm:minimal-bad-extraction};
\item the reselection collapses the compound to a single nondegenerate core
      and the branch enters the single-core compatibility dichotomy governed by
      \cref{paper3:thm:scale-rigid-single-core-compatibility};
\item the reselection produces a perturbative compound that is below the
      critical small-data threshold of Step~2 of
      \cref{paper3:thm:same-point-reduction}, hence is absorbed into the
      perturbative remainder and treated by
      \cref{paper3:thm:terminal-decoupling};
\item the reselection still has degenerate scale-center Jacobian on every
      further subsequence and has no further active subconcentration; in this
      case the no-subconcentration alternative of
      \cref{paper3:lem:scale-rigid-no-subconcentration-bounded} gives a
      nontrivial bounded selected-window limit on a compact cylinder.
\end{enumerate}
\end{lemma}

\begin{proof}
By the gauge-degeneracy alternative
\cref{paper3:def:local-multibubble}\,(v), the localized scale or centering
selection of the selected compound or terminal frame is non-unique along the
subsequence.  Apply the parameter partition lemma
\cref{paper3:lem:parameter-partition} to the selected compound together with
any candidate alternate selection produced by the gauge degeneracy.  Three
mutually exclusive structural outcomes are possible.

\noindent\textit{Case A: alternate selection is comparable to the original at
the same physical point.}  Then both selections lie in the same scale
comparability class of Step~1 of \cref{paper3:thm:same-point-reduction}.
Step~2 of that theorem groups them into a single compound profile.  If the
compound is below the small-data threshold, we are in case (D3); otherwise
the compound is above threshold and Step~3 of the theorem applies.  If the
compound is then nondegenerate, we are in case (D2); if it is still
degenerate, the gauge degeneracy persists on the reselected compound, and we
proceed to Case~D below.

\noindent\textit{Case B: alternate selection is at a strictly different
scale at the same physical point.}  By
\cref{paper3:lem:outer-vanish}, the strictly outer profile vanishes locally
in the innermost frame; absorbing it into the perturbative remainder lands
us in case (D3).  If instead the reselection produces a strictly inner core
not previously detected, then either compact-cylinder suitable control fails,
in which case \cref{paper3:lem:compactness-failure} produces the new CKN core,
or compact control persists and
\cref{paper3:lem:scale-rigid-pressure-core-activation} converts the retained
CKN mass into an active \(L^3\) core unless the branch has already exited to
the separated/exterior pressure alternative.  These outcomes are case (D1).

\noindent\textit{Case C: alternate selection is at a separated physical
point.}  By \cref{paper3:thm:separated-point-reduction}, the new selection
is locally invisible in the original frame, hence absorbed into the
perturbative remainder; this is case (D3).  If the new selection retains
positive CKN concentration in its own frame, then
\cref{paper3:lem:scale-rigid-pressure-core-activation} makes it an additional
separated active core, or else assigns it to the separated/exterior pressure
alternative.  Both outcomes are case (D1) via
\cref{paper3:def:local-multibubble}\,(iv).

\noindent\textit{Case D: persistent degeneracy.}  Suppose every reselection
on every further subsequence still has degenerate scale-center Jacobian.
Apply \cref{paper3:lem:scale-rigid-pressure-modulation-stability} and then
\cref{paper3:lem:scale-rigid-no-subconcentration-bounded}.  If the compact
limit has a CKN subconcentration point, that lemma produces an additional
active core; this is case (D1).  If no such subconcentration remains, the same
lemma gives a nonzero bounded selected-window limit; this is case (D4).  No
\(L^\infty\) bound is inferred from the Jacobian degeneracy itself.
\end{proof}

\begin{theorem}[Scale-rigid stratum decomposition]
\label{paper3:thm:scale-rigid-stratum-decomposition}
Every scale-rigid local branch in the sense of
\cref{paper3:def:scale-rigid-branch} enters one of the following named local
strata after passing to a subsequence:
\begin{enumerate}[label=\textup{(\roman*)}]
\item the multicore alternative \cref{paper3:def:local-multibubble}\,(ii)--(iv),
      governed by \cref{paper3:thm:minimal-bad-extraction};
\item the compact single-core stratum, governed by
      \cref{paper3:thm:scale-rigid-single-core-compatibility};
\item the perturbative remainder of \cref{paper3:def:active-family}, governed
      by \cref{paper3:thm:terminal-decoupling};
\item the bounded selected-window limit alternative
      \cref{paper1:def:bounded-window-limit}, which by
      \cref{paper3:lem:scale-rigid-bounded-limit-exit} is incompatible with
      \cref{paper1:def:typeII}.
\end{enumerate}
\end{theorem}

\begin{proof}
By \cref{paper3:def:scale-rigid-branch}\,(R4), the branch enters the
scale-rigid stratum either via a degenerate relative scale-center Jacobian or
as an innermost cascade core after outer profiles have been absorbed
perturbatively.

\noindent\textit{Degenerate Jacobian case.}  Apply
\cref{paper3:lem:scale-rigid-degenerate-jacobian}.  Cases (D1)--(D4) of that
lemma give items (i)--(iv) respectively.

\noindent\textit{Innermost cascade case.}  By Step~4 of
\cref{paper3:thm:same-point-reduction}, after subtracting the outer-profile
sum \(S_n\to0\) in \(L^\infty_I L^3(B_{2R})\) and absorbing constant drifts
into the translation modulation, the residual is governed by
\cref{paper3:thm:terminal-decoupling}.  That theorem produces three
outcomes already enumerated as (a) single-core, (b) scale-rigid (degenerate
Jacobian), or (c) purely perturbative.  Case (a) is item (ii) here via
\cref{paper3:thm:scale-rigid-single-core-compatibility}; case (c) is item
(iii); case (b) loops back into the degenerate Jacobian case treated above
and is therefore reduced by
\cref{paper3:lem:scale-rigid-degenerate-jacobian} to one of (i)--(iv).
\end{proof}

\begin{lemma}[Finite refinement of scale-rigid structural exits]
\label{paper3:lem:scale-rigid-refinement-terminates}
Let \(\omega\) be a scale-rigid local branch produced by
\cref{paper3:thm:same-point-reduction} or
\cref{paper3:thm:separated-point-reduction}, with the inherited active family
of finite critical mass and with concentration capture in the sense of
\cref{paper3:def:active-family}.  Resolve every structural exit of
\cref{paper3:thm:scale-rigid-stratum-decomposition} by applying only
\cref{paper3:thm:minimal-bad-extraction},
\cref{paper3:thm:same-point-reduction},
\cref{paper3:thm:separated-point-reduction}, and
\cref{paper3:thm:terminal-decoupling}, never
\cref{paper3:thm:multi}.  Then the refinement terminates after finitely many
ledger-decreasing refinements.  At the terminal residual branch one has either a nonzero
bounded selected-window limit or a compact single-core
vanishing-dissipation branch excluded by
\cref{paper1:thm:single-core-criterion}.
\end{lemma}

\begin{proof}
After the active threshold has been lowered once, every core produced by
\cref{paper3:lem:scale-rigid-pressure-core-activation} is handled by
\cref{paper3:lem:scale-rigid-active-core-capture}: it is already in the
captured inherited family, is absorbed into the terminal frame, or exits to the
separate profile-completeness obligation.  The residual scale-rigid ledger
therefore uses only the captured active family.  Write
\(\eta_*^{\rm new}\) for the lowered active threshold.  By
\cref{paper3:lem:finite-active-count}, the resulting family contains at most
\[
  N_*:=\lfloor M_*(\eta_*^{\rm new})^{-3}\rfloor
\]
active profiles.  Group same-point comparable-scale profiles into compound
profiles as in \cref{paper3:def:compound-profile}.  There are therefore only
finitely many active objects and only finitely many unordered pair relations
among them.

Define the refinement ledger to consist of:
\begin{enumerate}[label=\textup{(L\arabic*)}]
\item every active object not yet selected, grouped, absorbed, or declared
      terminal;
\item every pair relation not yet classified as comparable same-point, strict
      same-point cascade, separated point, or locally invisible;
\item every degenerate scale-center selection not yet resolved by
      \cref{paper3:lem:scale-rigid-degenerate-jacobian};
\item every perturbative remainder not yet removed by
      \cref{paper3:thm:terminal-decoupling}.
\end{enumerate}
The ledger is finite because it is built from a finite active family and
finitely many compounds.  Order ledgers lexicographically by the number of
unresolved entries in \textup{(L1)}, \textup{(L2)}, \textup{(L3)}, and
\textup{(L4)}.

At each nonterminal step, no multibubble exclusion theorem is used.  If
\cref{paper3:thm:scale-rigid-stratum-decomposition} gives a multicore,
same-point cascade, or separated/exterior structural exit, apply
\cref{paper3:thm:minimal-bad-extraction} to select the corresponding active
object.  Comparable same-point objects are grouped into a compound profile;
below-threshold compounds are removed by the small-data step in
\cref{paper3:thm:same-point-reduction}, while above-threshold compounds count
as active objects.  Strict same-point cascades are reduced by
\cref{paper3:thm:same-point-reduction}; separated physical centers are reduced
by \cref{paper3:thm:separated-point-reduction}.  Perturbative exterior and
remainder terms are discarded only through
\cref{paper3:thm:terminal-decoupling}, whose output preserves the selected
positive concentration and produces no new active object.

Each move strictly decreases the ledger.  Adding a previously unseen active
core consumes one entry of \textup{(L1)}, and this can happen at most \(N_*\)
times.  Grouping comparable same-point objects, ordering a same-point cascade,
separating two physical centers, or declaring one object locally invisible
resolves a previously unresolved pair and consumes one entry of \textup{(L2)}.
Applying \cref{paper3:lem:scale-rigid-degenerate-jacobian} resolves the
selected degenerate frame or produces an active structural exit already counted
in \textup{(L1)}--\textup{(L2)}, so it consumes one entry of \textup{(L3)}.  Applying
\cref{paper3:thm:terminal-decoupling} to a perturbative remainder consumes one
entry of \textup{(L4)}.  Once an entry is resolved it is not reintroduced: later
passes use the grouped compound, the selected innermost representative, the
separated-frame assignment, or the discarded perturbative remainder.  Hence no
cycle is possible.

After the ledger is exhausted, item \textup{(i)} of
\cref{paper3:thm:scale-rigid-stratum-decomposition} cannot occur again, because
it would create either a new active object or an unresolved active relation,
contradicting exhaustion of \textup{(L1)}--\textup{(L2)}.  Item \textup{(iii)}
cannot carry the retained positive concentration by the capture clause in
\cref{paper3:def:active-family}; a purely perturbative remainder has already
been absorbed through \textup{(L4)}.
Item \textup{(iv)} is exactly the nonzero bounded selected-window-limit
alternative.

It remains only to handle item \textup{(ii)}.  Apply
\cref{paper3:thm:scale-rigid-single-core-compatibility}.  Its bounded-limit
alternative gives the first terminal outcome.  Its further-active-core
alternative would add a new ledger entry, impossible after ledger exhaustion.
Its vanishing-localized-dissipation alternative verifies
\cref{paper1:def:single-core-branch} for the original selected sequence and is
excluded by \cref{paper1:thm:single-core-criterion}.  These are precisely the
two terminal residual outcomes claimed.
\end{proof}

\begin{proposition}[Discharge of scale-rigidity exclusion]
\label{paper3:prop:scale-rigidity-discharged}
Every scale-rigid local branch produced by
\cref{paper3:thm:same-point-reduction} or
\cref{paper3:thm:separated-point-reduction}, including a branch with
degenerate relative scale-center Jacobian, either has a nonzero bounded
selected-window limit on a compact cylinder (and so fails
\cref{paper1:def:typeII} by
\cref{paper3:lem:scale-rigid-bounded-limit-exit}) or reduces to the compact
single-core branch governed by
\cref{paper1:thm:single-core-criterion} (equivalently by the discharged
proposition \cref{paper3:ass:single-core}).  Consequently no such branch is
compatible with the local Type II condition.  This is the conclusion of
\cref{paper3:ass:scale-rigidity}; that assumption is therefore discharged.
\end{proposition}

\begin{proof}
Let \(\omega\) be a scale-rigid local branch.
\Cref{paper3:thm:scale-rigid-stratum-decomposition} places \(\omega\) into
one of items (i)--(iv).

If item (i) occurs, the branch enters a named multicore, same-point cascade, or
separated/exterior structural exit.  Resolve that exit by the finite
refinement procedure of
\cref{paper3:lem:scale-rigid-refinement-terminates}.  This procedure uses only
\cref{paper3:thm:minimal-bad-extraction},
\cref{paper3:thm:same-point-reduction},
\cref{paper3:thm:separated-point-reduction}, and
\cref{paper3:thm:terminal-decoupling}; it does not invoke
\cref{paper3:thm:multi}, so there is no circular use of
\cref{paper3:ass:scale-rigidity}.  The refinement terminates at a bounded
selected-window limit or at the compact single-core
vanishing-dissipation branch.

In item (ii), apply
\cref{paper3:thm:scale-rigid-single-core-compatibility}.  Its bounded-limit
case is the first alternative in the proposition.  Its
vanishing-localized-dissipation case verifies
\cref{paper1:def:single-core-branch} for the original selected sequence and is
excluded by \cref{paper1:thm:single-core-criterion}; this is the
``reduces to the single-core branch'' alternative.  Its further-active-core
case is resolved by the same finite refinement lemma used for item (i).

In item (iii), the branch is the perturbative remainder; by
\cref{paper3:thm:terminal-decoupling} it produces no independent
multibubble obstruction.  The retained positive concentration is carried by
the selected component, not by the perturbative remainder, so the residual
selected component is reduced to item (ii) or item (iv) by
\cref{paper3:lem:scale-rigid-refinement-terminates}.

In item (iv), by
\cref{paper3:lem:scale-rigid-bounded-limit-exit} the branch is not a local
Type II sequence in the sense of \cref{paper1:def:typeII}.  This is the
``nonzero bounded selected-window limit'' alternative.

In every case the branch either has a nonzero bounded selected-window limit
or reduces to the single-core branch.  This is the conclusion of
\cref{paper3:ass:scale-rigidity}.
\end{proof}

\begin{proposition}[Terminal admissibility in active frames]\label{paper3:ass:terminal-admissibility}
Assume finite critical norm, terminal windowed profile completeness in the
sense of \cref{paper8:ass:terminal-windowed-profile-completeness}, small-data
perturbative stability, the exterior annulus alternative
\cref{paper8:cor:exterior-annulus-alternative}, the repaired-gauge representation
of \cref{paper2:thm:C}, and local Caccioppoli regularity.  Then every terminal
family produced by the same-point and separated-point reductions is admissible
on each compact selected cylinder in the sense of
\cref{paper3:def:admissible-terminal}.
\end{proposition}

\begin{proof}
The terminal nonlinear profile decoupling theorem,
\cref{paper8:thm:terminal-nonlinear-decoupling}, applies under the stated
hypotheses and produces, on every compact terminal cylinder, the conclusions
of \cref{paper8:def:terminal-nonlinear-decoupling}.  We verify each clause of
\cref{paper3:def:admissible-terminal}.

Items (i) and (ii) are the active-family bookkeeping of
\cref{paper3:def:active-family} together with the same-point and
separated-point grouping of \cref{paper3:lem:separated-vanish,paper3:lem:outer-vanish};
they are properties of the active family that survive the terminal-frame
construction by definition of the terminal active class.

Item (iii) is the velocity decoupling
\cref{paper8:lem:terminal-velocity-decoupling}, the linear residual decoupling
\cref{paper8:lem:terminal-linear-residual}, and the pressure decoupling
\cref{paper8:lem:terminal-pressure-decoupling}; together they give the
\(L^\infty_I L^3\) vanishing of the rescaled remainder \(r_n\) and the
existence of \((\pi_n^r,f_n)\) with the stated \(L^1_IH^{-1}(B_R)\) decay.

Item (iv) is the retained \(L^\infty_I L^3(B_{2R})\) bound established in the
proof of \cref{paper8:lem:terminal-scale-collapse-admissibility}, which uses
the finite critical norm and the retained profile bounds.

Item (v) is the bounded-modulation hypothesis, supplied by the repaired-gauge
representation \cref{paper2:thm:C} and \cref{paper2:thm:B}.

Item (vi) is the local pressure-decomposition stability of
\cref{paper8:lem:terminal-pressure-decoupling}, whose proof relies only on the
Calderon--Zygmund pressure-tail decoupling
\cref{paper8:lem:pressure-kernel-tail} together with the kernel boundedness on
\(L^{3/2}\) provided by the classical Calderon--Zygmund theorem
\cite{CalderonZygmund1952,Stein1970}.

Item (vii) follows from the Leibniz formula
\([\Delta,\chi_R]u=(\Delta\chi_R)u+2\nabla\chi_R\cdot\nabla u\) and the
identity
\([\nabla\chi_R,\nabla\cdot](u\otimes u)=(\nabla\chi_R)(u\otimes u)\),
combined with the duality
\(\|fu\|_{H^{-1}(B_R)}\le C_R\|f\|_{C^1(\overline B_R)}\|u\|_{L^{6/5}(B_R)}\)
and the embedding
\(L^3(B_R)\hookrightarrow L^{6/5}(B_R)\) on bounded balls.  Indeed, for
\(\varphi\in H^1_0(B_R)\),
\[
  \langle\nabla\chi_R\cdot\nabla u,\varphi\rangle
  =-\langle u,(\Delta\chi_R)\varphi+\nabla\chi_R\cdot\nabla\varphi\rangle,
\]
giving \(\|[\Delta,\chi_R]u\|_{H^{-1}(B_R)}\le C_R\|u\|_{L^3(B_{2R})}\).
The pointwise bound
\(|(\nabla\chi_R)(u\otimes u)|\le\|\nabla\chi_R\|_\infty|u|^2\) gives
\(\|(\nabla\chi_R)(u\otimes u)\|_{L^{3/2}(B_R)}\le C_R\|u\|_{L^3(B_{2R})}^2\),
and \(L^{3/2}\hookrightarrow H^{-1}\) on \(B_R\) supplies the displayed bound.
The combined estimate
\[
  \|[\Delta,\chi_R]u\|_{H^{-1}(B_R)}
  +\|[\nabla\chi_R,\,\nabla\cdot](u\otimes u)\|_{H^{-1}(B_R)}
  \le C_R\bigl(\|u\|_{L^3(B_{2R})}+\|u\|_{L^3(B_{2R})}^2\bigr)
  \le C_R'\|u\|_{L^3(B_{2R})}^2
\]
holds whenever \(\|u\|_{L^3(B_{2R})}\) is bounded below by a fixed constant,
which is the case in the active selected frame.  The same estimate applies to
the \(S_n\)-quadratic commutator terms thanks to the uniform
\(L^\infty_I L^3(B_{2R})\) control of \(U_n\) and the
\(L^1_I L^{3/2}(B_R)\) control of the cross stresses produced by
\cref{paper8:lem:terminal-stress-decoupling}.

All items in \cref{paper3:def:admissible-terminal} are therefore satisfied,
which is the asserted admissibility.
\end{proof}

\begin{theorem}[Local multibubble and cascade exclusion]\label{paper3:thm:multi}
Use the local single-core exclusion
\cref{paper3:ass:single-core}, the scale-rigidity exclusion
\cref{paper3:ass:scale-rigidity}, and the terminal admissibility in
active frames \cref{paper3:ass:terminal-admissibility}.  Let a positive local
Type II branch admit an active family satisfying the hypotheses of
\cref{paper3:thm:minimal-bad-extraction}.  Then no local multibubble, cascade, or
gauge-degenerate Type II case can occur in this class.
\end{theorem}

\begin{proof}
Let a local Type II branch be in the multibubble/cascade class of
\cref{paper3:def:local-multibubble}.  Apply \cref{paper3:thm:minimal-bad-extraction}.  The
positive local concentration is carried by an active frame or compound profile,
and one of the alternatives listed there occurs.
\par
\textit{Case 1: compactness-failure reduction.}
Compactness failure yields, after changing radius by a fixed factor and
rescaling as in \cref{paper3:lem:compactness-failure}, a bounded-cylinder concentration
branch.  Relabel this branch as an active core.  After passing to a subsequence,
the new core and every previously retained active core have one of three
relations: comparable scale at the same physical point, strictly separated scale
at the same physical point, or separated physical center.  Comparable same-point
cores are grouped into compound profiles.  A perturbative compound profile is
removed by small-data stability; an active compound profile is treated as one
selected core in its active frame.  If the scale-center selection inside such a
compound core is degenerate, the branch is gauge-degenerate and belongs to the
non-uniqueness alternative in \cref{paper3:def:local-multibubble}, rather than to
a separate core dynamics.
\par
\textit{Case 1a: gauge degeneracy.}
If the scale-center selection of the selected compound or terminal frame is
degenerate, then the branch is in the gauge-degenerate alternative of
\cref{paper3:def:local-multibubble}.  By the repaired-gauge representation
from \cref{sec:repaired-gauge-construction}, nondegeneracy is the hypothesis needed to enter the
single-core branch.  Persistent degeneracy therefore belongs to the
scale-rigidity alternative in \cref{paper3:ass:scale-rigidity}; if the degeneracy
disappears after passing to a subsequence, the branch enters the single-core or
terminal-frame alternatives below.
\par
\textit{Case 1b: same-point comparable scales.}
Same-point comparable-scale active frames are grouped into the compound profile
of \cref{paper3:def:compound-profile}.  If the compound profile is below the critical
small-data threshold, it is part of the perturbative remainder.  If it is above
threshold and its scale-center selection is nondegenerate, it is one selected
active core and enters the single-core branch.  If the corresponding
scale-center selection is degenerate, it is covered by Case 1a.  Thus
comparable same-point profiles do not produce an additional multibubble
mechanism.
\par
\textit{Case 2: same-point cascading.}
Strict same-point scale cascades are treated by
\cref{paper3:thm:same-point-reduction}.  In the innermost frame, all outer scales are
locally invisible after constant drifts are absorbed.  The branch therefore
reduces to one of the three terminal alternatives
\[
  \text{(A) single-core},\quad \text{(B) scale-rigid},\quad \text{(C) terminal perturbative}.
\]
\par
\textit{Case 3: separated-point interactions.}
For separated physical points, \cref{paper3:thm:separated-point-reduction} gives the same
three terminal alternatives as in Case 2:
\[
  \text{(A) single-core},\quad \text{(B) scale-rigid},\quad \text{(C) terminal perturbative}.
\]

The single-core alternative is excluded by \cref{paper3:ass:single-core}.  The
scale-rigid alternative is excluded by \cref{paper3:ass:scale-rigidity}, which includes
the scale-collapse and autonomous-limit exclusions in \cref{sec:scale-collapse-cost} in this setting.  The
terminal perturbative alternative gives no independent multibubble obstruction:
\cref{paper3:thm:terminal-decoupling} removes all exterior profiles and remainders from
the selected active equation in \(L^1H^{-1}_{\mathrm{loc}}\), preserves the positive
finite-mass selected branch, and leaves only the single-core or scale-rigid
possibilities already excluded.  Consequently every alternative in
\cref{paper3:def:local-multibubble} is incompatible with the local Type II regime.
\end{proof}

\begin{corollary}[Removal of the multibubble alternative]\label{paper3:cor:multi-removal}
Under the hypotheses of \cref{paper3:thm:multi}, every failure of the single
nondegenerate compact core alternative is excluded.
\end{corollary}

\begin{proof}
By \cref{paper3:def:local-multibubble}, failure of the single nondegenerate compact
core alternative means that at least one of the following occurs: the
compact-cylinder suitable estimates are unbounded; positive local CKN mass splits
between active cores; a strict same-point scale cascade occurs; separated
physical points remain active in distinct local frames; or the localized scale
or center selection is degenerate.  The compactness-failure case produces an
additional core by \cref{paper3:lem:compactness-failure}.  Same-point and separated
cases are reduced by \cref{paper3:thm:same-point-reduction,paper3:thm:separated-point-reduction}.
\Cref{paper3:thm:multi} excludes the remaining cases.  Hence no failure of the
single nondegenerate compact core alternative remains.
\end{proof}

\clearpage
\section{Local Caccioppoli and Rough-Core Exclusion}\label{sec:rough-core-exclusion}
\subsection{Analytic inputs and notation}

The renormalized local energy inequality from the preceding renormalized-chart
analysis is used as an analytic input.  The renormalized suitable structure and
pressure reconstruction are not rederived.  All estimates are conditional on
the following local hypotheses on the compact renormalized cylinders under
consideration.

\begin{proposition}[Renormalized local energy inequality]\label{paper4:ass:lei}
Let $(V,P,a,b)$ be defined on a compact renormalized cylinder
$Q_{2R,J}:=B_{2R}\times J$, where $J\subset\mathbb R$ is a bounded interval.
Assume
\begin{equation}\label{paper4:eq:renorm-ns}
\partial_\tau V+(V\cdot\nabla)V+\nabla P-\nu\Delta V
= a(\tau)V+a(\tau)(y\cdot\nabla V)+b(\tau)\cdot\nabla V,
\qquad
\nabla\cdot V=0,
\end{equation}
with $\nu>0$, $\nabla\cdot V=0$, and $a,b\in L^\infty(J)$. For every
nonnegative $\phi\in C_c^\infty(Q_{2R,J})$ and every $\tau_1<\tau_2$ in $J$,
\begin{equation}\label{paper4:eq:local-energy}
\begin{aligned}
&\int_{B_{2R}}\frac{|V(\cdot,\tau_2)|^2}{2}\phi(\cdot,\tau_2)
-\int_{B_{2R}}\frac{|V(\cdot,\tau_1)|^2}{2}\phi(\cdot,\tau_1)
\\
&\quad+\nu\int_{\tau_1}^{\tau_2}\!\int_{B_{2R}}\phi\,|\nabla V|^2
\le
\int_{\tau_1}^{\tau_2}\!\int_{B_{2R}}\frac{|V|^2}{2}(\partial_\tau\phi+\nu\Delta\phi)
\\
&\quad+\int_{\tau_1}^{\tau_2}\!\int_{B_{2R}}\frac{|V|^2}{2}V\cdot\nabla\phi
+\int_{\tau_1}^{\tau_2}\!\int_{B_{2R}}PV\cdot\nabla\phi
\\
&\quad+\int_{\tau_1}^{\tau_2}\!\int_{B_{2R}}
a\frac{|V|^2}{2}(-\phi-y\cdot\nabla\phi)
\\
&\quad-\int_{\tau_1}^{\tau_2}\!\int_{B_{2R}}\frac{|V|^2}{2}b\cdot\nabla\phi.
\end{aligned}
\end{equation}
The pressure is understood modulo functions of time, and the mean-normalized
pressure on compact balls belongs to the local class used below.
\end{proposition}

\begin{proof}
\Cref{paper2:prop:ren-suitable} transports the physical local energy
inequality of Caffarelli--Kohn--Nirenberg
\cite{CaffarelliKohnNirenberg1982} into the renormalized variables
\((y,\tau)\) and yields the inequality stated above with test function of the
form \(\phi^2\), where \(\phi\in C_c^\infty(Q_R^*)\) is nonnegative.  Setting
\(\Phi:=\phi^2\) and noting that every nonnegative
\(\Phi\in C_c^\infty(Q_{2R,J})\) is the limit, in
\(C_c^k(Q_{2R,J})\) for every \(k\), of squares of nonnegative smooth
compactly supported functions \(\phi_\varepsilon:=(\Phi+\varepsilon^2)^{1/2}-\varepsilon\),
the inequality extends by approximation to all such \(\Phi\); this is the
standard reformulation used in the suitable weak solution framework
\cite{CaffarelliKohnNirenberg1982,Lin1998}.  The compact-window
representation theorem \cref{paper2:thm:C} provides the cylinder
\(Q_{2R,J}=B_{2R}\times J\) and the divergence-free renormalized equation
\eqref{paper4:eq:renorm-ns}, while
\cref{paper2:thm:pressure-decomp} supplies the localized pressure
decomposition with the mean-normalized pressure on compact balls.  Identifying
the modulation coefficients \(a=\tilde a\) and \(b=\tilde b\) of
\cref{paper2:prop:ren-suitable} with those of
\eqref{paper4:eq:renorm-ns} and grouping the renormalized error terms exactly
as in \eqref{paper4:eq:local-energy} yields the stated inequality.
\end{proof}

Throughout, for a space--time set $E\subset\mathbb R^3\times\mathbb R$, write
\[
\int_E f:=\int_E f(y,\tau)\,dy\,d\tau.
\]
For $r>0$ and a bounded interval $J\subset\mathbb R$, define
\[
Q_{r,J}:=B_r\times J,\qquad |Q_{r,J}|:=|B_r|\,|J|.
\]
For a ball $B\subset\mathbb R^3$, denote the spatial mean of pressure by
\[
(P)_B(\tau):=\frac1{|B|}\int_B P(y,\tau)\,dy.
\]
The local quantities are
\[
\mathcal A_J(r):=\operatorname*{ess\,sup}_{\tau\in J}\int_{B_r}|V(y,\tau)|^2\,dy,
\]
\[
\mathcal H_J(r):=\int_J\!\int_{B_r}|\nabla V|^2\,dy\,d\tau,
\]
\[
\mathcal C_J(r):=\int_J\!\int_{B_r}|V|^3\,dy\,d\tau,
\qquad
\mathcal D_J(r):=\int_J\!\int_{B_r}|P-(P)_{B_r}(\tau)|^{3/2}\,dy\,d\tau.
\]

\begin{lemma}[Finite-measure conversion]\label{paper4:lem:l2-from-l3}
For every bounded interval $J$ and every $r>0$,
\[
\int_J\!\int_{B_r}|V|^2
\le |Q_{r,J}|^{1/3}\mathcal C_J(r)^{2/3}
\le |Q_{r,J}|^{1/3}\bigl(1+\mathcal C_J(r)\bigr).
\]
\end{lemma}

\begin{proof}
H\"older's inequality on $Q_{r,J}$ gives
\[
\int_J\!\int_{B_r}|V|^2
\le |Q_{r,J}|^{1/3}\left(\int_J\!\int_{B_r}|V|^3\right)^{2/3}.
\]
The second inequality follows from $s^{2/3}\le 1+s$ for $s\ge0$.
\end{proof}

\begin{lemma}[Comparison of pressure normalizations]\label{paper4:lem:pressure-means}
Let $1\le p<\infty$ and $0<r\le R$. For almost every $\tau$ such that
$P(\tau)\in L^p(B_R)$,
\[
\|P(\tau)-(P)_{B_r}(\tau)\|_{L^p(B_r)}
\le 2\|P(\tau)-(P)_{B_R}(\tau)\|_{L^p(B_R)}.
\]
Consequently,
\[
\mathcal D_J(r)\le 2^{3/2}\mathcal D_J(R)
\]
whenever the right-hand side is finite.
\end{lemma}

\begin{proof}
Set $c=(P)_{B_R}(\tau)$. Then
\[
\|P-(P)_{B_r}\|_{L^p(B_r)}
\le \|P-c\|_{L^p(B_r)}+|B_r|^{1/p}|(P)_{B_r}-c|.
\]
Since
\[
|(P)_{B_r}-c|\le |B_r|^{-1}\int_{B_r}|P-c|
\le |B_r|^{-1/p}\|P-c\|_{L^p(B_r)},
\]
the displayed pointwise estimate follows. Raising to the power $3/2$ and
integrating in time gives the last assertion.
\end{proof}

\subsection{Compact-cylinder Caccioppoli estimate}

\begin{theorem}[Compact-cylinder Caccioppoli estimate]\label{paper4:thm:caccioppoli}
Let $R>0$, let $J\subset\mathbb R$ be a bounded open interval, and let
$J'\Subset J$ be a compact subinterval. Assume
\cref{paper4:ass:lei} on $Q_{2R,J}$ and assume
\[
\mathcal C_J(2R)+\mathcal D_J(2R)<\infty.
\]
Then there is a constant
\[
C=C\bigl(R,|J|,\nu,\|a\|_{L^\infty(J)},\|b\|_{L^\infty(J)},
\operatorname{dist}(J',\partial J)\bigr)
\]
such that
\[
\mathcal A_{J'}(R)+\mathcal H_{J'}(R)
\le C\Bigl(1+
\mathcal C_J(2R)+\mathcal D_J(2R)
\Bigr).
\]
The estimate is outer-to-inner: all right-hand quantities are measured on
$B_{2R}\times J$, while the conclusion is on $B_R\times J'$.
\end{theorem}

\begin{proof}
\textbf{Step 1: localized cut-off and the two energy inequalities.}
Set
\[
\delta:=\operatorname{dist}(J',\partial J)>0.
\]
Choose $\eta\in C_c^\infty(J)$ and $\zeta\in C_c^\infty(B_{2R})$ with
\[
0\le\eta\le1,\qquad
0\le\zeta\le1,\qquad
\eta\equiv1\ \text{on}\ J',\qquad
\zeta\equiv1\ \text{on}\ B_R,
\]
and
\[
\|\eta'\|_{L^\infty(J)}\le C_\eta\,\delta^{-1},
\qquad
\|\nabla\zeta\|_{L^\infty(B_{2R})}\le C_\zeta R^{-1},
\qquad
\|\Delta\zeta\|_{L^\infty(B_{2R})}\le C_\zeta R^{-2}.
\]
Set
\[
\phi(y,\tau):=\eta(\tau)\,\zeta(y)^2.
\]
Then $\phi\in C_c^\infty(Q_{2R,J})$, $0\le\phi\le1$, and
\[
\phi\equiv1\text{ on }B_R\times J'.
\]
Choose $\tau_-<\inf J'$ and $\tau_+>\sup J'$ in $J$ with
$\eta(\tau_-)=\eta(\tau_+)=0$. Applying \eqref{paper4:eq:local-energy} on
$(\tau_-,\tau_+)$ gives
\[
\nu\int_{\tau_-}^{\tau_+}\!\int_{B_{2R}}\phi\,|\nabla V|^2
\le |R_\tau|+|R_\Delta|+|R_{\mathrm{conv}}|+|R_{\mathrm{pres}}|+|R_{\mathrm{mod}}|,
\]
where
\[
R_\tau:=\int_J\!\int_{B_{2R}}\frac{|V|^2}{2}\,\partial_\tau\phi,\qquad
R_\Delta:=\nu\int_J\!\int_{B_{2R}}\frac{|V|^2}{2}\,\Delta\phi,
\]
\[
R_{\mathrm{conv}}:=\int_J\!\int_{B_{2R}}\bigl(\tfrac{|V|^2}{2}V\bigr)\cdot\nabla\phi,\qquad
R_{\mathrm{pres}}:=\int_J\!\int_{B_{2R}}\bigl(P-(P)_{B_{2R}}(\tau)\bigr)V\cdot\nabla\phi,
\]
\[
R_{\mathrm{mod}}:=\int_J\!\int_{B_{2R}}
a\frac{|V|^2}{2}(-\phi-y\cdot\nabla\phi)
-\int_J\!\int_{B_{2R}}\frac{|V|^2}{2}b\cdot\nabla\phi.
\]
For almost every $\sigma\in J'$, applying \eqref{paper4:eq:local-energy} on
$(\tau_-,\sigma)$ gives the same right-hand estimates and leaves the terminal
term $\frac12\int_{B_R}|V(\sigma)|^2$ on the left. Therefore the bounds below
control both $\mathcal H_{J'}(R)$ and $\mathcal A_{J'}(R)$.
Define
\[
L:=\int_J\!\int_{B_{2R}}|V|^2,\qquad
C_3:=\mathcal C_J(2R),\qquad D_{3/2}:=\mathcal D_J(2R).
\]
By \cref{paper4:lem:l2-from-l3},
\[
L\le |Q_{2R,J}|^{1/3}C_3^{2/3}\le C(R,|J|)(1+C_3).
\]

Since $\phi(y,\tau)=\eta(\tau)\zeta(y)^2$,
\[
\partial_\tau\phi=\eta'(\tau)\zeta(y)^2,
\qquad
\nabla\phi=2\eta(\tau)\zeta(y)\nabla\zeta(y),
\qquad
\Delta\phi=\eta(\tau)\Delta(\zeta(y)^2),
\]
hence
\[
\|\partial_\tau\phi\|_{L^\infty(Q_{2R,J})}\le C_\eta\delta^{-1},
\qquad
\|\nabla\phi\|_{L^\infty(Q_{2R,J})}\le 2C_\zeta R^{-1},
\qquad
\|\Delta\phi\|_{L^\infty(Q_{2R,J})}\le C_\zeta R^{-2}.
\]

\textbf{Step 2: derivative terms from the cut-off.}
\[
|R_\tau|
\le\frac12\|\partial_\tau\phi\|_{L^\infty(Q_{2R,J})}L
\le C\delta^{-1}(1+C_3),
\]
\[
|R_\Delta|
\le\frac{\nu}{2}\|\Delta\phi\|_{L^\infty(Q_{2R,J})}L
\le C\nu R^{-2}(1+C_3).
\]

\textbf{Step 3: convection term.}
Using $|\nabla\phi|\le CR^{-1}$,
\[
|R_{\mathrm{conv}}|
\le \frac12\|\nabla\phi\|_{L^\infty(Q_{2R,J})}\int_J\!\int_{B_{2R}}|V|^3
\le CR^{-1}C_3.
\]

\textbf{Step 4: pressure term.}
For each time $\tau$,
\[
\int_{B_{2R}}(P-(P)_{B_{2R}}(\tau))V\cdot \nabla\phi
=\int_{B_{2R}}PV\cdot \nabla\phi,
\]
with
\[
\int_{B_{2R}}(P)_{B_{2R}}(\tau)V\cdot \nabla\phi
=(P)_{B_{2R}}(\tau)\int_{B_{2R}}V\cdot \nabla\phi
=-(P)_{B_{2R}}(\tau)\int_{B_{2R}}\phi\,\nabla\cdot V=0
\]
(using $\phi\in C_c^\infty(B_{2R})$ and $\nabla\cdot V=0$).
Thus subtracting any time-dependent pressure constant is legitimate in the
pressure flux term.
Hence
\[
|R_{\mathrm{pres}}|
\le\|\nabla\phi\|_{L^\infty(Q_{2R,J})}
\int_J
\|P-(P)_{B_{2R}}(\tau)\|_{L_x^{3/2}(B_{2R})}
\|V(\tau,\cdot)\|_{L_x^3(B_{2R})}
\,d\tau
\le CR^{-1}D_{3/2}^{2/3}C_3^{1/3}.
\]
Young's inequality gives
\[
|R_{\mathrm{pres}}|\le C(R)(C_3+D_{3/2}).
\]

\textbf{Step 5: modulation terms.}
The local energy inequality already contains the modulation terms
after the spatial integrations by parts from the repaired-gauge equation:
\[
\int_{B_{2R}}(y\cdot\nabla V)\cdot V\,\phi
=\frac12\int_{B_{2R}}y\cdot\nabla(|V|^2)\phi
=-\frac12\int_{B_{2R}}|V|^2\nabla\cdot(y\phi),
\]
and
\[
\int_{B_{2R}}(b(\tau)\cdot\nabla V)\cdot V\,\phi
=\frac12\int_{B_{2R}}b(\tau)\cdot\nabla(|V|^2)\phi
=-\frac12\int_{B_{2R}}|V|^2b(\tau)\cdot\nabla\phi.
\]
Thus we estimate the additional terms directly. Split
\[
R_{\mathrm{mod}}=R_a+R_{ay}+R_b,
\]
where
\[
R_a:=-\int_J\!\int_{B_{2R}}a\,\frac{|V|^2}{2}\phi,
\quad
R_{ay}:=-\int_J\!\int_{B_{2R}}a\,\frac{|V|^2}{2}\,y\cdot\nabla\phi,
\quad
R_b:=-\int_J\!\int_{B_{2R}}\frac{|V|^2}{2}b\cdot\nabla\phi.
\]
\[
M_a:=\|a\|_{L^\infty(J)},\qquad M_b:=\|b\|_{L^\infty(J)}.
\]
For $R_a$,
\[
|R_a|
\le \frac12 M_aL
\le C M_a(1+C_3).
\]
For the scale-gradient part, $|y|\le2R$ on $B_{2R}$ and
$\|\nabla\phi\|_{L^\infty(Q_{2R,J})}\le CR^{-1}$, hence
\[
|y\cdot\nabla\phi|\le C.
\]
Therefore
\[
|R_{ay}|\le C M_aL\le C M_a(1+C_3).
\]
For $R_b$,
\[
|R_b|\le C M_bR^{-1}L\le C M_bR^{-1}(1+C_3).
\]

\textbf{Step 6: absorption and conclusion.}
Collecting estimates from Steps~2--5 yields
\[
\nu\int_{\tau_-}^{\tau_+}\!\int_{B_{2R}}\phi|\nabla V|^2
\le C\bigl(1+C_3+D_{3/2}\bigr),
\]
where $C$ depends only on
\[
R,|J|,\nu,\|a\|_{L^\infty(J)},\|b\|_{L^\infty(J)},\delta.
\]
Since $\phi\equiv1$ on $B_R\times J'$,
\[
\mathcal H_{J'}(R)
\le C\bigl(1+\mathcal C_J(2R)+\mathcal D_J(2R)\bigr).
\]
Repeating the same estimates with terminal time $\sigma\in J'$ gives
\[
\frac12\int_{B_R}|V(y,\sigma)|^2\,dy
\le C\bigl(1+\mathcal C_J(2R)+\mathcal D_J(2R)\bigr)
\]
for almost every $\sigma\in J'$. Taking the essential supremum over
$\sigma\in J'$ and renaming $C$ proves the theorem.
\end{proof}

\begin{definition}[Local compact-cylinder control]
Local Caffarelli--Kohn--Nirenberg control on $Q_{r,J}$ means
\[
\mathcal C_J(r)+\mathcal D_J(r)<\infty.
\]
Local windowed $H^1$ control on $Q_{r,J}$ means
\[
\mathcal H_J(r)<\infty.
\]
\end{definition}

\begin{corollary}[Outer control implies inner suitable control]\label{paper4:cor:outer-inner}
Let $R>0$, let $J\subset\mathbb R$ be a bounded open interval, and let
$J'\Subset J$ be a compact subinterval. Assume \cref{paper4:ass:lei} on $Q_{2R,J}$ and
$\mathcal C_J(2R)+\mathcal D_J(2R)<\infty$. Then
\[
\mathcal A_{J'}(R)+\mathcal H_{J'}(R)<\infty.
\]
More quantitatively, the bound is the one in \cref{paper4:thm:caccioppoli}. If an
argument uses pressure means on a smaller ball, \cref{paper4:lem:pressure-means}
converts those normalizations to the outer-ball normalization above.
\end{corollary}

\begin{proof}
The estimate follows from \cref{paper4:thm:caccioppoli}. The comparison of pressure normalizations follows from
\cref{paper4:lem:pressure-means}.
\end{proof}

\begin{corollary}[Rough-core exclusion on compact windows]\label{paper4:cor:rough-core}
Let $R>0$, let $J\subset\mathbb R$ be a bounded open interval, and let
$J'\Subset J$ be a compact subinterval. Assume \cref{paper4:ass:lei} on $Q_{2R,J}$. If
\[
\mathcal C_J(2R)+\mathcal D_J(2R)<\infty,
\]
then
\[
\mathcal H_{J'}(R)<\infty.
\]
Equivalently, if $\mathcal H_{J'}(R)=+\infty$ for an inner cylinder
$Q_{R,J'}$, then for every enclosing cylinder $Q_{2R,J}$ with
$J'\Subset J$ on which the local energy inequality and bounded
modulation hypotheses hold,
\[
\mathcal C_J(2R)=+\infty
\qquad\text{or}\qquad
\mathcal D_J(2R)=+\infty.
\]
\end{corollary}

\begin{proof}
The direct implication is the gradient part of \cref{paper4:thm:caccioppoli}. The
last assertion is its contrapositive with the same outer-to-inner geometry:
finite outer $\mathcal C+\mathcal D$ on $Q_{2R,J}$ forces finite inner
$\mathcal H$ on $Q_{R,J'}$.
\end{proof}

\begin{remark}[Meaning of rough core]
Here a rough core means failure of local windowed $H^1$
control on compact inner renormalized cylinders while the modulation
coefficients are bounded. \Cref{paper4:cor:rough-core} shows that such a failure is
not independent of compact-cylinder Caffarelli--Kohn--Nirenberg control: it can
occur only if the velocity-pressure bounds $\mathcal C+\mathcal D$ fail on
every relevant larger enclosing compact cylinder.
\end{remark}

\subsection{Consequences for the assembly}

We shall use the following consequences in the final assembly.
\begin{enumerate}[label=\textup{(\arabic*)}]
\item On every compact renormalized window satisfying the local energy
inequality and bounded modulation, finite $\mathcal C+\mathcal D$ on
$Q_{2R,J}$ implies finite $\mathcal A+\mathcal H$ on $Q_{R,J'}$ whenever
$J'\Subset J$.
\item The pressure normalization may be taken on any comparable ball, after
using \cref{paper4:lem:pressure-means}.
\item Failure of local windowed $H^1$ control on an inner compact cylinder
forces failure of the compact-cylinder $\mathcal C+\mathcal D$ bounds on each
larger enclosing cylinder to which \cref{paper4:thm:caccioppoli} applies.
\end{enumerate}

\clearpage
\section{Scale-Collapse Cost and Autonomous Limits}\label{sec:scale-collapse-cost}
\medskip
\noindent\textit{Notation and functional conventions.}
For a bounded spatial region \(K\subset\mathbb R^3\), we use
\[
  (f,g)_{L^2(K)}:=\int_K f\cdot g\,dx,\qquad
  \langle F,\phi\rangle_{H^{-1}(K),H^1_0(K)}
  :=\int_0^T\langle F(\cdot,s),\phi(\cdot,s)\rangle_{H^{-1}(K),H^1_0(K)}\,ds
\]
whenever \(F\in L^1(0,T;H^{-1}(K))\) and \(\phi\in L^\infty(0,T;H^1_0(K))\). Weak limits and weak convergences are taken in the distributional sense on \(C_c^\infty\)-test functions.
\[
  \|f\|_{L^p_tL^q_x}:=\left(\int_0^T\|f(\cdot,s)\|_{L^q(\mathbb R^3)}^p\,ds\right)^{1/p},
\]
for any fixed \(T>0\), with the usual modification \(p=\infty\).

\subsection{Renormalized scale-collapse setting}

\medskip
\noindent\textit{Position of the result.}
The scale-collapse analysis has two components:
\begin{enumerate}[label=(\arabic*)]
  \item Sections~\ref{paper5:sec:cost}--\ref{paper5:sec:finite-cost} contain the cost
    exclusion arguments, which do not use compactness.
  \item Section~\ref{paper5:sec:autonomous} contains the autonomous compactness
    statement under explicitly stated compactness hypotheses.
\end{enumerate}

Throughout the scale-collapse analysis, the renormalized representation established earlier is assumed: on the
time interval under consideration, the fields \((V,P)\) satisfy \eqref{paper5:eq:renorm-ns} and its
distributional weak formulation.  In particular, the solenoidal formulation is obtained by
restricting the test fields below to divergence-free fields.

Let \((u,p)\) be a divergence-free suitable weak solution of
\[
  \partial_tu+(u\cdot\nabla)u+\nabla p=\nu\Delta u,
  \qquad \nabla\cdot u=0,
\]
in \(\mathbb R^3\), with viscosity \(\nu>0\).  Let
\(\lambda(\tau)>0\), \(X(\tau)\in\mathbb R^3\), and a measurable gauge function
\(c(t)\) be given, and define the renormalized fields by
\[
  u(x,t)=\lambda(\tau)^{-1}V\left(\frac{x-X(\tau)}{\lambda(\tau)},\tau\right),
  \qquad
  p(x,t)=\lambda(\tau)^{-2}P\left(\frac{x-X(\tau)}{\lambda(\tau)},\tau\right)+c(t).
\]
Pressure is defined only up to an arbitrary function of time; \(c(t)\) has no effect on
\(\nabla p\), so it is irrelevant for all subsequent estimates.
Assume that \((V,P)\) satisfies on \((0,\infty)\times\mathbb R^3\)
\begin{equation}\label{paper5:eq:renorm-ns}
  \partial_\tau V+(V\cdot\nabla)V+\nabla P
  =\nu\Delta V+a(\tau)(V+y\cdot\nabla V)+b(\tau)\cdot\nabla V,
  \qquad \nabla\cdot V=0.
\end{equation}
For every \(\psi\in C_c^\infty((0,\infty)\times\mathbb R^3;\mathbb R^3)\) with \(\nabla\cdot\psi=0\),
the pair \((V,P)\) satisfies the weak formulation
\[
\begin{aligned}
0={}&\int_0^\infty\!\!\int_{\mathbb R^3}
\bigl(
-V\cdot\partial_\tau\psi
-(V\otimes V):\nabla\psi
+\nu\nabla V:\nabla\psi
\bigr)\,dy\,d\tau \\
&+\int_0^\infty\!\!\int_{\mathbb R^3}
\bigl(
-a(\tau)(V+y\cdot\nabla V)\cdot\psi
-(b(\tau)\cdot\nabla V)\cdot\psi
+P\,\nabla\cdot\psi
\bigr)\,dy\,d\tau .
\end{aligned}
\]
For every \(0<T<\infty\), restricting the \(\tau\)-integration to \((0,T)\) gives the same identity on \((0,T)\times\mathbb R^3\).
Equivalently, for every
\(\psi\in C_c^\infty((0,\infty)\times\mathbb R^3;\mathbb R^3)\), not necessarily
divergence free,
\[
\begin{aligned}
0={}&\int_0^\infty\!\!\int_{\mathbb R^3}
\bigl(
-V\cdot\partial_\tau\psi
-(V\otimes V):\nabla\psi
\nu\nabla V:\nabla\psi
-P\,\nabla\cdot\psi
\bigr)\,dy\,d\tau \\
&+\int_0^\infty\!\!\int_{\mathbb R^3}
\bigl(
-a(\tau)(V+y\cdot\nabla V)\cdot\psi
-(b(\tau)\cdot\nabla V)\cdot\psi
\bigr)\,dy\,d\tau .
\end{aligned}
\]

Define positive and negative parts of the scalar field \(a\) by
\[
  a_+(\tau):=\max\{a(\tau),0\},
  \qquad
  a_-(\tau):=\max\{-a(\tau),0\}.
\]
Then for a.e. \(\tau\), \(a=a_+-a_-\), \(a_+,a_-\ge0\), and \(a_+a_-=0\).  Assume
\[
  a\in L^1_{\mathrm{loc}}([\tau_0,\infty)),
  \qquad
  b\in L^1_{\mathrm{loc}}([\tau_0,\infty);\mathbb R^3),
\]
and the scale convention
\begin{equation}\label{paper5:eq:log-scale-convention}
  \frac{d}{d\tau}\log\lambda(\tau)=a(\tau)
\end{equation}
for a.e. \(\tau\ge\tau_0\), with \(\log\lambda\) absolutely continuous on compact intervals.

All integrals below are interpreted in \([0,\infty]\). When we write
\(\int_{\tau_0}^{\infty}\), partial integrals are taken over \([\tau_0,T]\) and
monotone convergence is used as \(T\to\infty\).

Fix \(R>0\).  Choose
\(
  \Phi\in C_c^\infty(\mathbb R^3;[0,1])
\)
satisfying
\begin{equation}\label{paper5:eq:Phi-properties}
  \Phi(y)=1\ \text{if }|y|\le R,
  \quad
  \operatorname{supp}_y\Phi\subset B_{2R}.
\end{equation}
Define the localized kinetic mass
\[
  M_R(\tau):=\int_{\mathbb R^3}|V(y,\tau)|^2\Phi(y)
  \,dy.
\]
Since \(\Phi\) is fixed in \(\tau\), \(M_R\) is measurable in \(\tau\), and only this
measurability is used below.

\begin{definition}[Nonvanishing localized core]
A renormalized branch has nonvanishing localized core if there exist
constants \(m_0>0\) and \(\tau_1\ge\tau_0\) such that
\[
  M_R(\tau)\ge m_0\quad\text{for a.e. }\tau\in[\tau_1,\infty).
\]
In applications, this lower bound is supplied by the selected-core construction; here it is stated as an explicit hypothesis.
\end{definition}

\begin{definition}[Scale-collapse and absolute scale-drift densities]
Define
\[
  \mathcal C_{\mathrm{sc}}(\tau):=a_-(\tau)M_R(\tau),
  \qquad
  \mathcal C_{\mathrm{abs}}(\tau):=\nu\int_{\mathbb R^3}|\nabla V(y,\tau)|^2\Phi(y)
  \,dy+|a(\tau)|M_R(\tau).
\]
\(\mathcal C_{\mathrm{sc}}\) is the negative drift weighted by localized core mass, while
\(\mathcal C_{\mathrm{abs}}\) controls both collapse drift and local dissipation.
Their total costs are extended real-valued by
\[
  \int_{\tau_0}^{\infty}\mathcal C_{\mathrm{sc}}(\tau)
  \,d\tau,
  \qquad
  \int_{\tau_0}^{\infty}\mathcal C_{\mathrm{abs}}(\tau)
  \,d\tau.
\]
\end{definition}

\begin{definition}[Finite-cost branch]\label{paper5:def:finite-cost-branch}
A branch has finite scale-collapsing cost if
\[
  \int_{\tau_0}^{\infty}\mathcal C_{\mathrm{sc}}(\tau)
  \,d\tau<\infty.
\]
It has finite absolute cost if
\[
  \int_{\tau_0}^{\infty}\mathcal C_{\mathrm{abs}}(\tau)
  \,d\tau<\infty.
\]
\end{definition}

\subsection{Elementary alternative}
\label{paper5:sec:cost}

\begin{lemma}[Finite-or-infinite alternative]\label{paper5:lem:finite-infinite}
Let \(f\ge0\) be measurable on \([\tau_0,\infty)\). Then exactly one of the following holds:
\[
  \int_{\tau_0}^{\infty}f(\tau)
  \,d\tau<\infty
  \quad\text{or}\quad
  \int_{\tau_0}^{\infty}f(\tau)
  \,d\tau=\infty.
\]
\end{lemma}

\begin{proof}
\begin{enumerate}[label=(\arabic*)]
  \item For \(T>\tau_0\), define partial costs
  \[
    F(T):=\int_{\tau_0}^{T}f(\tau)
    \,d\tau\in[0,\infty].
  \]

  \item If \(T_2>T_1\), then by nonnegativity of \(f\),
  \[
    F(T_2)-F(T_1)=\int_{T_1}^{T_2}f(\tau)
    \,d\tau\ge0,
  \]
  so \(F\) is monotone nondecreasing.

  \item A monotone function on \([\tau_0,\infty)\) has an extended limit
  \(L\in[0,\infty]\),
  \[
    L:=\lim_{T\to\infty}F(T).
  \]
  By definition of improper integral, this limit equals
  \(\int_{\tau_0}^{\infty}f(\tau)\,d\tau\).

  \item Therefore \(L\) is either finite or equal to \(+\infty\), and these
  cases are mutually exclusive.
\end{enumerate}
Hence exactly one item in the lemma holds.
\end{proof}

\begin{lemma}[Logarithmic scale identity]\label{paper5:lem:log-scale}
Under \eqref{paper5:eq:log-scale-convention}, for all \(\tau_1,\tau_2\) with
\(\tau_0\le\tau_1<\tau_2<\infty\),
\[
  \int_{\tau_1}^{\tau_2}a(\tau)
  \,d\tau
  =\log\lambda(\tau_2)-\log\lambda(\tau_1).
\]
\end{lemma}

\begin{proof}
\begin{enumerate}[label=(\arabic*)]
  \item Since \(a\in L^1_{\mathrm{loc}}\), the map
  \(\log\lambda\) is absolutely continuous on \([\tau_1,\tau_2]\), hence has a
  representative with a.e. derivative equal to \(a\).

  \item By the fundamental theorem of calculus for absolutely continuous functions,
  for every \(\tau\in[\tau_1,\tau_2]\),
  \[
    \log\lambda(\tau)=\log\lambda(\tau_1)+\int_{\tau_1}^{\tau}a(s)
    \,ds.
  \]

  \item Setting \(\tau=\tau_2\) yields the identity.
\end{enumerate}
\end{proof}

\subsection{Scale collapse versus localized core lower bounds}

\begin{theorem}[Collapse with nonvanishing localized core forces infinite
scale-collapsing cost]\label{paper5:thm:collapse-core-floor}
Assume:
\begin{enumerate}[label=(\arabic*)]
  \item \(\lambda(\tau)\to0\) as \(\tau\to\infty\);
  \item the branch has nonvanishing localized core.
\end{enumerate}
Then
\[
  \int_{\tau_0}^{\infty}\mathcal C_{\mathrm{sc}}(\tau)
  \,d\tau
  =\int_{\tau_0}^{\infty}a_-(\tau)M_R(\tau)
  \,d\tau=\infty.
\]
\end{theorem}

\begin{proof}
\begin{enumerate}[label=(\arabic*)]
  \item Choose \(m_0>0\) and \(\tau_1\ge\tau_0\) from the nonvanishing
  core assumption so that
  \[
    M_R(\tau)\ge m_0\quad\text{for a.e. }\tau\in[\tau_1,\infty).
  \]

  \item By Lemma~\ref{paper5:lem:log-scale}, for any \(\tau_2>\tau_1\),
  \[
    \int_{\tau_1}^{\tau_2}a(\tau)
    \,d\tau
    =\log\lambda(\tau_2)-\log\lambda(\tau_1).
  \]
  Taking \(\tau_2\to\infty\) and using \(\lambda(\tau_2)\to0\),
  \[
    \lim_{\tau_2\to\infty}\int_{\tau_1}^{\tau_2}a(\tau)
    \,d\tau=-\infty.
  \]

  \item Decompose \(a=a_+-a_-\).  Then for each \(\tau_2>\tau_1\),
  \[
    \int_{\tau_1}^{\tau_2}a(\tau)
    \,d\tau
    =\int_{\tau_1}^{\tau_2}a_+(\tau)
    \,d\tau-\int_{\tau_1}^{\tau_2}a_-(\tau)
    \,d\tau
    \ge-\int_{\tau_1}^{\tau_2}a_-(\tau)
    \,d\tau.
    \]
    Hence
    \[
    \int_{\tau_1}^{\tau_2}a_-(\tau)
    \,d\tau\ge
    -\int_{\tau_1}^{\tau_2}a(\tau)
    \,d\tau.
    \]

  \item Taking limits and using Step~2 yields
  \[
    \int_{\tau_1}^{\infty}a_-(\tau)
    \,d\tau=\infty.
  \]
  Indeed, by Step~2,
  \[
    -\int_{\tau_1}^{\tau_2}a(\tau)\,d\tau\to+\infty.
  \]
  Since the partial integrals of \(a_-\) dominate this quantity, they are unbounded.
  Because \(a_-\ge0\), monotone convergence gives
  \(\int_{\tau_1}^{\infty}a_-(\tau)\,d\tau=\infty\).

  \item Split the cost into early and late times:
  \[
    \int_{\tau_0}^{\infty}a_-M_R
    \,d\tau
    =\int_{\tau_0}^{\tau_1}a_-M_R
    \,d\tau+
    \int_{\tau_1}^{\infty}a_-M_R
    \,d\tau.
  \]
  The first term is finite or infinite independently of the argument; the second term
  satisfies
  \[
    \int_{\tau_1}^{\infty}a_-(\tau)M_R(\tau)
    \,d\tau
    \ge m_0\int_{\tau_1}^{\infty}a_-(\tau)
    \,d\tau
    =\infty
  \]
  by the bound in (1) and Step~4.  Hence the total integral is infinite.
\end{enumerate}
\end{proof}

\subsection{Finite-cost exclusions}
\label{paper5:sec:finite-cost}

\begin{theorem}[Scale-collapse cost exclusion]\label{paper5:thm:cost-exclusion}
No branch with finite scale-collapsing cost (Definition~\ref{paper5:def:finite-cost-branch}) can satisfy simultaneously
\(\lambda(\tau)\to0\) and nonvanishing localized core.
\end{theorem}

\begin{proof}
Assume, for contradiction, that such a branch exists.  The two hypotheses are exactly
those of Theorem~\ref{paper5:thm:collapse-core-floor}; therefore
\[
  \int_{\tau_0}^{\infty}\mathcal C_{\mathrm{sc}}(\tau)
  \,d\tau=\infty.
\]
This contradicts finite scale-collapsing cost.  Therefore no such branch exists.
\end{proof}

\begin{lemma}[Absolute-cost domination]\label{paper5:lem:absolute-cost}
For a.e. \(\tau\in[\tau_0,\infty)\),
\[
  \mathcal C_{\mathrm{abs}}(\tau)\ge\mathcal C_{\mathrm{sc}}(\tau).
\]
Hence, pointwise and after integrating in \(\tau\),
\[
  \int_{\tau_0}^{\infty}\mathcal C_{\mathrm{sc}}(\tau)\,d\tau
  \le
  \int_{\tau_0}^{\infty}\mathcal C_{\mathrm{abs}}(\tau)\,d\tau,
\]
as extended real numbers.
In particular, finite absolute cost implies finite scale-collapsing cost, and
equivalently
\[
  \int_{\tau_0}^{\infty}\mathcal C_{\mathrm{sc}}(\tau)\,d\tau=\infty
  \quad\Longrightarrow\quad
  \int_{\tau_0}^{\infty}\mathcal C_{\mathrm{abs}}(\tau)\,d\tau=\infty.
\]
\end{lemma}

\begin{proof}
At each time \(\tau\),
\[
  \mathcal C_{\mathrm{abs}}(\tau)
  =\nu\int|\nabla V|^2\Phi
  +(a_++a_-)M_R
  \ge(a_++a_-)M_R\ge a_-M_R
  =\mathcal C_{\mathrm{sc}}(\tau),
\]
since \(a_+,a_-\ge0\) and \(\nu\int|\nabla V|^2\Phi\ge0\).  Integrating over
\([\tau_0,\infty)\) yields the claimed implication.
\end{proof}

\begin{corollary}[Absolute scale-cost exclusion]\label{paper5:cor:absolute-cost-exclusion}
No branch with finite absolute cost in the sense of Definition~\ref{paper5:def:finite-cost-branch}
can satisfy simultaneously \(\lambda(\tau)\to0\) and nonvanishing localized core.
\end{corollary}

\begin{proof}
Assume instead such a branch exists.  By Theorem~\ref{paper5:thm:collapse-core-floor},
\(\int\mathcal C_{\mathrm{sc}}=\infty\).  By Lemma~\ref{paper5:lem:absolute-cost}, the
absolute cost is pointwise above \(\mathcal C_{\mathrm{sc}}\), so its integral must also be
infinite, contradicting finite absolute cost.
\end{proof}

\begin{remark}
The exclusions above are purely integral and use neither compactness nor any
autonomous-limit argument.
\end{remark}

\subsection{Autonomous reduced-limit theorem from the selected-window estimates}
\label{paper5:sec:autonomous}

The preceding exclusions are integral statements and do not rely on compactness.  On a fixed finite autonomous window, compactness requires not only the cost bound but also the selected-window estimates supplied by the preceding renormalized representation, pressure reconstruction, and local Caccioppoli estimates.

\medskip
\noindent\textit{Analytic inputs.}
The preceding sections supply the following ingredients on selected compact
renormalized windows:
\begin{enumerate}[label=(\arabic*)]
  \item the repaired-gauge renormalized equation and weak formulation for \((V,P)\);
  \item pressure reconstruction modulo spatial means, with local \(L^{3/2}\) bounds on compact
    cylinders;
  \item local windowed bounds
  \(V\in L^\infty_tL^2_x\cap L^2_tH^1_x\) on compact cylinders;
  \item a retained compact core lower bound preventing the selected profile from vanishing.
\end{enumerate}
Together with the thick-window convergence of the modulation coefficients, these estimates imply
the compactness and limit passage below.

\begin{definition}[Thick autonomous windows]\label{paper5:def:thick-autonomous-windows}
A branch has \emph{thick autonomous windows} if there exist constants
\(T>0\), \(c>0\), \(M_1\in(0,\infty)\), a sequence \((\tau_n)\) with
\(\tau_n\to\infty\), and limits \(a_\infty\in\mathbb R\),
\(b_\infty\in\mathbb R^3\), such that for every \(n\):
\begin{enumerate}[label=(\roman*)]
  \item
  \[
    \frac1T\int_{\tau_n}^{\tau_n+T}a_-(\tau)M_R(\tau)
    \,d\tau\ge c;
  \]
  \item
  \[
    0\le M_R(\tau)\le M_1\quad\text{for a.e. }\tau\in[\tau_n,\tau_n+T];
  \]
  \item
  \[
    a(\tau_n+\cdot)\to a_\infty\ \text{in }L^1(0,T),
    \qquad
    b(\tau_n+\cdot)\to b_\infty\ \text{in }L^1(0,T;\mathbb R^3).
  \]
\end{enumerate}
\end{definition}
\begin{remark}
This notion isolates a limiting drift coefficient \(a_\infty<0\).  The
compactness to extract a profile limit is obtained from the selected-window compactness estimates,
not from the scalar cost inequalities of Sections~\ref{paper5:sec:cost}--\ref{paper5:sec:finite-cost}.
\end{remark}

\begin{lemma}[Thick drift implies negative autonomous limit parameter]\label{paper5:lem:negative-parameter}
On thick autonomous windows, the limiting autonomous coefficient is negative:
\[
  a_\infty\le-\frac{c}{M_1}<0.
\]
\end{lemma}

\begin{proof}
\begin{enumerate}[label=(\arabic*)]
  \item From (ii), for a.e. \(\tau\in[\tau_n,\tau_n+T]\),
  \(0\le M_R(\tau)\le M_1\).  Therefore
  \[
    \frac{a_-(\tau)M_R(\tau)}{M_1}
    \le a_-(\tau).
  \]
  Therefore
  \[
    \frac1T\int_{\tau_n}^{\tau_n+T}a_-(\tau)
    \,d\tau
    \ge \frac1{M_1T}\int_{\tau_n}^{\tau_n+T}a_-(\tau)M_R(\tau)
    \,d\tau
    \ge \frac{c}{M_1}.
  \]

  \item Set \(a_n(s):=a(\tau_n+s)\).  By (iii), \(a_n\to a_\infty\) in
  \(L^1(0,T)\). The map \(r\mapsto r^-\) is 1-Lipschitz, so
  \(a_n^-\to a_\infty^-\) in \(L^1(0,T)\):
  \[
    \|a_n^--a_\infty^-\|_{L^1(0,T)}\le
    \|a_n-a_\infty\|_{L^1(0,T)}\to0.
  \]
  Thus
  \[
  \lim_{n\to\infty}\frac1T\int_{\tau_n}^{\tau_n+T}a_-(\tau)
    \,d\tau
    =\frac1T\int_0^T(a_\infty)^-
    \,ds=(a_\infty)^-.
  \]

  \item Combine Step~1 and Step~2:
  \((a_\infty)^-\ge c/M_1\), i.e.\ \(a_\infty\le-c/M_1<0\).
\end{enumerate}
\end{proof}

\begin{proposition}[Selected-window suitable estimates]\label{paper5:ass:selected-window-estimates}
Let \((\tau_n)\) be a thick autonomous-window sequence and set
\[
  V_n(y,s):=V(y,\tau_n+s),\qquad
  P_n(y,s):=P(y,\tau_n+s),
\]
\[
  \alpha_n(s):=a(\tau_n+s),\qquad
  \beta_n(s):=b(\tau_n+s),
  s\in(0,T).
\]
For every bounded Lipschitz domain \(K\Subset\mathbb R^3\), there is a constant
\(C_K<\infty\), independent of \(n\), such that
\[
  \|V_n\|_{L^\infty(0,T;L^2(K))}
  +\|V_n\|_{L^2(0,T;H^1(K))}
  +\|P_n-(P_n)_K\|_{L^{3/2}((0,T)\times K)}
  \le C_K,
\]
where
\[
  (P_n)_K(s):=\frac1{|K|}\int_KP_n(y,s)\,dy.
\]
Each \(V_n\) is divergence free and there exist
\(\chi\in C_c^\infty(B_R)\), \(\chi\ge0\), \(\eta>0\), and \(N\in\mathbb N\) such that
for all \(n\ge N\),
\[
  \int_{\mathbb R^3}|V_n(y,s)|^2\chi(y)\,dy\ge\eta
  \quad\text{for a.e. }s\in(0,T).
\]
\end{proposition}

\begin{proof}
The localized lower bound is the conclusion of the selected core retention
theorem \cref{paper5:thm:core-retention} below, which is purely a property of
the selected-core construction and is independent of the autonomous-window
selection.  With the lower bound in hand,
\cref{paper5:prop:selected-window-derived} below derives the local
\(L^\infty_tL^2_x\), \(L^2_tH^1_x\), and pressure bounds from the
repaired-gauge renormalized representation
\cref{paper2:thm:C}, the localized pressure decomposition
\cref{paper2:thm:pressure-decomp}, and the compact-cylinder Caccioppoli
estimate \cref{paper4:thm:caccioppoli}.  These two inputs together yield all
items in the statement.
\end{proof}

\begin{proposition}[Derivation of the selected-window estimates from the preceding sections]
\label{paper5:prop:selected-window-derived}
Let a repaired-gauge branch satisfy the renormalized representation and pressure reconstruction
of the compact-window representation theorem, \cref{paper2:thm:C}, on the windows \([\tau_n,\tau_n+T]\).  Assume that on
each compact cylinder \((0,T)\times K\) the local compact-cylinder Caffarelli--Kohn--Nirenberg
quantities are finite uniformly in \(n\), and that the selected core has the localized lower
bound stated in Assumption~\ref{paper5:ass:selected-window-estimates}.  Then
Assumption~\ref{paper5:ass:selected-window-estimates} holds.
\end{proposition}

\begin{proof}
The renormalized representation theorem supplies the shifted equation, divergence-free condition,
and pressure reconstruction modulo functions of time.  The pressure reconstruction gives, after
subtracting spatial means,
\[
  \sup_n\|P_n-(P_n)_K\|_{L^{3/2}((0,T)\times K)}<\infty
\]
on every compact \(K\).  The local Caccioppoli estimate applied on a slightly larger compact
cylinder gives the inner bounds
\[
  \sup_n\|V_n\|_{L^\infty(0,T;L^2(K))}
  +\sup_n\|V_n\|_{L^2(0,T;H^1(K))}<\infty.
\]
These are precisely the local velocity and pressure estimates in
Assumption~\ref{paper5:ass:selected-window-estimates}.  The retained selected-core lower bound is the
last condition in that assumption.  Hence all items in
Assumption~\ref{paper5:ass:selected-window-estimates} follow from the representation,
pressure, and local Caccioppoli results.
\end{proof}

\begin{lemma}[Local time regularity on autonomous windows]\label{paper5:lem:window-time-regularity}
Assume thick autonomous windows and Assumption~\ref{paper5:ass:selected-window-estimates}.  Let
\(K_0\Subset K_1\Subset\mathbb R^3\) be bounded Lipschitz domains, and choose an integer
\(m\) so large that \(H^m_0(K_1)\hookrightarrow W^{1,\infty}_0(K_1)\).  Then
\[
  \{\partial_sV_n\}_{n\ge1}
  \quad\text{is bounded in }L^1(0,T;H^{-m}(K_1)).
\]
\end{lemma}

\begin{proof}
Fix \(\zeta\in C_c^\infty(K_1)\) with \(\zeta\equiv1\) on a neighborhood of \(K_0\).
Let \(\phi\in H^m_0(K_1;\mathbb R^3)\).  Since \(H^m_0(K_1)\hookrightarrow W^{1,\infty}_0(K_1)\),
there is \(C_m\) such that
\[
  \|\phi\|_{L^\infty(K_1)}+\|\nabla\phi\|_{L^\infty(K_1)}
  \le C_m\|\phi\|_{H^m(K_1)}.
\]
The shifted equation is
\[
  \partial_sV_n
  =-(V_n\cdot\nabla)V_n-\nabla P_n+\nu\Delta V_n
  +\alpha_n(V_n+y\cdot\nabla V_n)+\beta_n\cdot\nabla V_n
\]
in distributions.  We estimate each term against \(\phi\).

For the convection term, integration by parts gives
\[
  \left|\langle (V_n\cdot\nabla)V_n,\phi\rangle\right|
  =
  \left|\int_{K_1}V_{n,i}V_{n,j}\partial_j\phi_i\,dy\right|
  \le C_m\|V_n(s)\|_{L^2(K_1)}^2\|\phi\|_{H^m(K_1)}.
\]
For the viscous term,
\[
  \left|\langle \nu\Delta V_n,\phi\rangle\right|
  =
  \nu\left|\int_{K_1}\nabla V_n:\nabla\phi\,dy\right|
  \le \nu C_m\|\nabla V_n(s)\|_{L^2(K_1)}|K_1|^{1/2}\|\phi\|_{H^m(K_1)}.
\]
For the pressure term, subtracting the spatial mean does not change the gradient:
\[
  \left|\langle\nabla P_n,\phi\rangle\right|
  =
  \biggl|\int_{K_1}(P_n-(P_n)_{K_1})\,\nabla\cdot\phi\,dy\biggr|
  \le C_m |K_1|^{1/3}
  \|P_n-(P_n)_{K_1}\|_{L^{3/2}(K_1)}
  \|\phi\|_{H^m(K_1)}.
\]
For the \(V_n\)-part of the scaling drift,
\[
  |\alpha_n(s)|\,\left|\int_{K_1}V_n\cdot\phi\,dy\right|
  \le C_m|\alpha_n(s)|\,|K_1|^{1/2}
  \|V_n(s)\|_{L^2(K_1)}\|\phi\|_{H^m(K_1)}.
\]
For the \(y\cdot\nabla V_n\)-part, integrate by parts:
\[
  \int_{K_1}(y\cdot\nabla V_n)\cdot\phi\,dy
  =
  -\int_{K_1}V_n\cdot(y\cdot\nabla\phi+3\phi)\,dy,
\]
because \(\phi\) has zero trace.  Hence
\[
  |\alpha_n(s)|\,\left|\int_{K_1}(y\cdot\nabla V_n)\cdot\phi\,dy\right|
  \le C_{K_1,m}|\alpha_n(s)|\|V_n(s)\|_{L^2(K_1)}\|\phi\|_{H^m(K_1)}.
\]
Similarly,
\[
  \left|\int_{K_1}(\beta_n\cdot\nabla V_n)\cdot\phi\,dy\right|
  =
  \left|\int_{K_1}V_n\cdot(\beta_n\cdot\nabla\phi)\,dy\right|
  \le C_m|\beta_n(s)|\,|K_1|^{1/2}
  \|V_n(s)\|_{L^2(K_1)}\|\phi\|_{H^m(K_1)}.
\]
Combining the estimates and taking the supremum over
\(\|\phi\|_{H^m(K_1)}\le1\), we obtain
\[
\begin{aligned}
  \|\partial_sV_n(s)\|_{H^{-m}(K_1)}
  \le C_{K_1,m}\big(&
  \|V_n(s)\|_{L^2(K_1)}^2
  +\|\nabla V_n(s)\|_{L^2(K_1)}\\
  &+\|P_n(s)-(P_n)_{K_1}(s)\|_{L^{3/2}(K_1)}
  +( |\alpha_n(s)|+|\beta_n(s)| )\|V_n(s)\|_{L^2(K_1)}
  \big).
\end{aligned}
\]
Integrating in \(s\in(0,T)\), using Assumption~\ref{paper5:ass:selected-window-estimates}, and using
\(\alpha_n\to a_\infty\), \(\beta_n\to b_\infty\) in \(L^1\), hence uniform \(L^1\)-bounds
for \(\alpha_n,\beta_n\), proves the claimed \(L^1(0,T;H^{-m})\) bound.
\end{proof}

\begin{lemma}[Local compactness on autonomous windows]\label{paper5:lem:autonomous-window-compactness}
Assume thick autonomous windows and Assumption~\ref{paper5:ass:selected-window-estimates}.  Then after
passing to a subsequence there exist
\[
  U\in L^\infty(0,T;L^2_{\mathrm{loc}}(\mathbb R^3))
  \cap L^2(0,T;H^1_{\mathrm{loc}}(\mathbb R^3)),
  \qquad
  \Pi\in L^{3/2}_{\mathrm{loc}}((0,T)\times\mathbb R^3),
\]
such that, for every compact \(K\Subset\mathbb R^3\),
\begin{enumerate}[label=(\alph*)]
  \item
  \[
    V_n\to U
    \quad\text{strongly in }L^2(0,T;L^2(K));
  \]
  \item
  \[
    V_n\rightharpoonup U
    \quad\text{weakly in }L^2(0,T;H^1(K));
  \]
  \item
  \[
    V_n\rightharpoonup^\ast U
    \quad\text{weakly-* in }L^\infty(0,T;L^2(K));
  \]
  \item
  \[
    P_n-(P_n)_K\rightharpoonup\Pi_K
    \quad\text{weakly in }L^{3/2}((0,T)\times K),
  \]
  where the local distributions \(\nabla\Pi_K\) agree on overlaps and define
  a pressure gradient \(\nabla\Pi\) modulo functions of time.
\end{enumerate}
\end{lemma}

\begin{proof}
Fix \(K_0\Subset K_1\Subset\mathbb R^3\).  Assumption~\ref{paper5:ass:selected-window-estimates}
gives boundedness of \(V_n\) in \(L^2(0,T;H^1(K_1))\), and
Lemma~\ref{paper5:lem:window-time-regularity} gives boundedness of \(\partial_sV_n\) in
\(L^1(0,T;H^{-m}(K_1))\) for \(m\) large.  Since
\[
  H^1(K_1)\Subset L^2(K_0)\hookrightarrow H^{-m}(K_1),
\]
the Aubin--Lions--Simon compactness theorem implies, after extracting a subsequence,
\[
  V_n\to U
  \quad\text{strongly in }L^2(0,T;L^2(K_0)).
\]
The local \(L^\infty L^2\) and \(L^2H^1\) bounds also give, after extraction,
\[
  V_n\rightharpoonup^\ast U
  \quad\text{in }L^\infty(0,T;L^2(K_0)),
  \qquad
  V_n\rightharpoonup U
  \quad\text{in }L^2(0,T;H^1(K_0)).
\]
The pressure bound in Assumption~\ref{paper5:ass:selected-window-estimates} gives weak compactness of
\(P_n-(P_n)_{K_0}\) in \(L^{3/2}((0,T)\times K_0)\).  A diagonal extraction over balls
\(B_j\) gives a single subsequence for every compact \(K\).  On overlaps, two pressure limits
can differ only by a function of \(s\), because both represent the same distributional pressure
gradient in the limit equation.  Thus the gradients assemble into \(\nabla\Pi\), with \(\Pi\)
defined modulo functions of time.
\end{proof}

\begin{lemma}[Compatibility of local pressure limits]\label{paper5:lem:pressure-compatibility}
Let \(K_1,K_2\Subset\mathbb R^3\) be bounded Lipschitz domains with
nonempty connected overlap \(K_{12}:=K_1\cap K_2\).  Suppose that, after passing to a common
subsequence,
\[
  P_n-(P_n)_{K_i}\rightharpoonup \Pi_i
  \quad\text{in }L^{3/2}((0,T)\times K_i),
  \qquad i=1,2.
\]
Then there exists \(c_{12}\in L^{3/2}(0,T)\) such that
\[
  \Pi_1(y,s)-\Pi_2(y,s)=c_{12}(s)
  \quad\text{for a.e. }(y,s)\in K_{12}\times(0,T).
\]
Consequently the pressure gradient \(\nabla\Pi\) obtained from the local limits is globally
well-defined on \((0,T)\times\mathbb R^3\), with the pressure itself defined modulo functions
of time.
\end{lemma}

\begin{proof}
For each \(n\),
\[
  \bigl(P_n-(P_n)_{K_1}\bigr)-\bigl(P_n-(P_n)_{K_2}\bigr)
  =(P_n)_{K_2}(s)-(P_n)_{K_1}(s)
\]
on \(K_{12}\).  Hence its spatial gradient is zero in
\(\mathcal D'((0,T)\times K_{12})\).  Passing to the weak limit gives
\[
  \nabla_y(\Pi_1-\Pi_2)=0
  \quad\text{in }\mathcal D'((0,T)\times K_{12}).
\]
Since \(K_{12}\) is connected, the standard distributional characterization of functions with
zero spatial gradient implies that \(\Pi_1-\Pi_2=c_{12}(s)\) for some
\(c_{12}\in L^{3/2}(0,T)\).  The last assertion follows because adding a time-dependent scalar
to pressure does not change \(\nabla\Pi\).
\end{proof}

\begin{theorem}[Autonomous reduced-limit theorem]\label{paper5:thm:autonomous-limit}
Assume that a branch has thick autonomous windows and satisfies
Assumption~\ref{paper5:ass:selected-window-estimates} on the same window sequence.  Let
\((U,\Pi)\) be the limit extracted in Lemma~\ref{paper5:lem:autonomous-window-compactness}.  Then
\[
  \nabla\cdot U=0,
\]
\[
  U\not\equiv0
\]
and \((U,\Pi)\) satisfies
\begin{equation}\label{paper5:eq:autonomous-limit}
  \partial_sU+(U\cdot\nabla)U+\nabla\Pi
  =\nu\Delta U+a_\infty(U+y\cdot\nabla U)+b_\infty\cdot\nabla U,
  \qquad
  \nabla\cdot U=0,
\end{equation}
in \(\mathcal D'((0,T)\times\mathbb R^3)\).  Equivalently, for every
\(\varphi\in C_c^\infty((0,T)\times\mathbb R^3;\mathbb R^3)\),
\[
  \int_0^T\!\int
  \bigl(
  -U\cdot\partial_s\varphi
  -(U\otimes U):\nabla\varphi
  +\nu\nabla U:\nabla\varphi
  -\Pi\,\nabla\cdot\varphi
  -a_\infty(U+y\cdot\nabla U)\cdot\varphi
  -(b_\infty\cdot\nabla U)\cdot\varphi
  \bigr)\,dy\,ds=0.
\]
\end{theorem}

\begin{remark}
The theorem asserts only the compact finite-window autonomous limit from the
selected-window estimates.  A rigidity contradiction requires a separate
Liouville theorem for the class containing \((U,\Pi)\).
\end{remark}

\begin{proof}
\emph{Step 1: divergence-free property of the limit.}

Each \(V_n\) is divergence free by Assumption~\ref{paper5:ass:selected-window-estimates}, so for any compactly supported scalar
\(\psi\in C_c^\infty((0,T)\times\mathbb R^3)\),
\[
  \int_{(0,T)\times\mathbb R^3}V_n\cdot\nabla\psi\,dy\,ds=0.
\]
Hence, for every such \(\psi\), since \(\nabla\psi\in C_c^\infty((0,T)\times\mathbb R^3)\) and \(V_n\to U\) in
\(L^2(0,T;L^2_{\mathrm{loc}})\),
\[
  \int_{(0,T)\times\mathbb R^3}(V_n-U)\cdot\nabla\psi\,dy\,ds\to0,
\]
we deduce
\[
  \int_{(0,T)\times\mathbb R^3}U\cdot\nabla\psi\,dy\,ds=0.
\]
Hence \(\nabla\cdot U=0\) in \(\mathcal D'((0,T)\times\mathbb R^3)\).


\emph{Step 2: nontriviality of the limit profile.}
Define for \(s\in(0,T)\)
\[
  f_n(s):=\int_{\mathbb R^3}|V_n(y,s)|^2\chi(y)
  \,dy,
  \qquad
  f(s):=\int_{\mathbb R^3}|U(y,s)|^2\chi(y)
  \,dy.
\]
Here \(K_\chi:=\operatorname{supp}\chi\), so
\(\chi^{1/2}V_n,\chi^{1/2}U\in L^2(0,T;L^2(K_\chi))\).
Then
\[
  f_n-f=\int_{\mathbb R^3}(V_n-U)\cdot(V_n+U)\chi
  \,dy.
\]
Applying Cauchy--Schwarz in space,
\[
  |f_n(s)-f(s)|\le
  \|\chi^{1/2}(V_n(s)-U(s))\|_{L^2(\mathbb R^3)}
  \|\chi^{1/2}(V_n(s)+U(s))\|_{L^2(\mathbb R^3)}.
\]
Integrating in \(s\) and applying Cauchy--Schwarz in time gives
\[
  \|f_n-f\|_{L^1(0,T)}
  \le
  \|\chi^{1/2}(V_n-U)\|_{L^2(0,T;L^2)}
  \|\chi^{1/2}(V_n+U)\|_{L^2(0,T;L^2)}.
\]
Since \(K_\chi\) is compact, Assumption~\ref{paper5:ass:selected-window-estimates} and
the membership \(U\in L^2(0,T;L^2(K_\chi))\) give
\[
  \sup_n\|\chi^{1/2}(V_n+U)\|_{L^2(0,T;L^2)}
  \le C_\chi\left(\sup_n\|V_n\|_{L^2(0,T;L^2(K_\chi))}
  +\|U\|_{L^2(0,T;L^2(K_\chi))}\right)<\infty.
\]
Moreover
\[
  \|\chi^{1/2}(V_n-U)\|_{L^2(0,T;L^2)}
  \le \|\chi\|_{L^\infty}^{1/2}
  \|V_n-U\|_{L^2(0,T;L^2(K_\chi))}\to0.
\]
Therefore
\[
  \|f_n-f\|_{L^1(0,T)}\to0.
\]
Hence \(f_n\to f\) in \(L^1(0,T)\).

Assumption~\ref{paper5:ass:selected-window-estimates} gives for large \(n\),
\(f_n(s)\ge\eta\) for a.e.\(s\). Therefore
\[
  \int_0^T f(s)
  \,ds=\lim_{n\to\infty}\int_0^Tf_n(s)
  \,ds\ge\eta T>0.
\]
Hence \(f\not\equiv0\), so \(U\not\equiv0\).

\emph{Step 3: shifted weak formulation for \(V_n\).}
Take any test field
\(\varphi\in C_c^\infty((0,T)\times\mathbb R^3;\mathbb R^3)\), and set
\[
K_\varphi:=\overline{\{x\in\mathbb R^3:\exists s\in(0,T),\ \varphi(x,s)\neq0\}}
\subset\mathbb R^3,\qquad K_\varphi\ \text{bounded}.
\]
Since \((V,P)\) satisfies~\eqref{paper5:eq:renorm-ns} in distributions, the change of variables
\(\tau=\tau_n+s\) gives the shifted weak formulation
\begin{equation}\label{paper5:eq:weak-n}
\begin{aligned}
0={}&\int_0^T\!\int_{\mathbb R^3}
\Bigl(
-V_n\cdot\partial_s\varphi
-(V_n\otimes V_n):\nabla\varphi
+\nu\nabla V_n:\nabla\varphi
-P_n\nabla\cdot\varphi
\Bigr)\,dy\,ds\\
&+\int_0^T\!\int_{\mathbb R^3}
\Bigl(
-a(\tau_n+s)(V_n+y\cdot\nabla V_n)\cdot\varphi
-(b(\tau_n+s)\cdot\nabla V_n)\cdot\varphi
\Bigr)\,dy\,ds.
\end{aligned}
\end{equation}

\emph{Step 4: term-by-term passage to the limit.}
Write \(\alpha_n(s):=a(\tau_n+s)\), \(\beta_n(s):=b(\tau_n+s)\).

\medskip
\noindent\textit{(i) Time derivative.}

By strong convergence in \(L^2(0,T;L^2(K_\varphi))\),
\[
  \left|\int_0^T\!\int (V_n-U)\cdot\partial_s\varphi
  \,dy\,ds\right|
  \le
  \|V_n-U\|_{L^2(0,T;L^2(K_\varphi))}
  \|\partial_s\varphi\|_{L^2(0,T;L^2(K_\varphi))}
  \to0.
\]
Hence
\[\int V_n\cdot\partial_s\varphi\to\int U\cdot\partial_s\varphi.
\]

\medskip
\noindent\textit{(ii) Convection term.}
Decompose
\[
  V_n\otimes V_n-U\otimes U=(V_n-U)\otimes V_n+U\otimes(V_n-U).
\]
Then
\[
\begin{aligned}
&\left|\int_0^T\!\int((V_n\otimes V_n-U\otimes U):\nabla\varphi)
  \,dy\,ds\right|\\
&\quad\le
  \|\nabla\varphi\|_{L^\infty_{x,t}}
  \|V_n-U\|_{L^2(0,T;L^2(K_\varphi))}
  \bigl(\|V_n\|_{L^2(0,T;L^2(K_\varphi))}
  +\|U\|_{L^2(0,T;L^2(K_\varphi))}\bigr)
  \to0,
\end{aligned}
\]
using boundedness of \(V_n\) in \(L^2(0,T;L^2(K_\varphi))\). Therefore
\(
  \int(V_n\otimes V_n):\nabla\varphi
  \to\int(U\otimes U):\nabla\varphi.
\)

\medskip
\noindent\textit{(iii) Viscous term.}
From weak convergence of \(V_n\) in \(L^2(0,T;H^1_{\mathrm{loc}})\),
\[
  \int_0^T\!\int\nabla V_n:\nabla\varphi
  \,dy\,ds
  \to
  \int_0^T\!\int\nabla U:\nabla\varphi
  \,dy\,ds.
\]

\medskip
\noindent\textit{(iv) Pressure term.}
Let \(K\) be a bounded Lipschitz domain with \(K_\varphi\Subset K\).  Since
\(\varphi\) is compactly supported in \(K\),
\[
  \int_K (P_n)_K(s)\,\nabla\cdot\varphi(y,s)\,dy
  =
  (P_n)_K(s)\int_K\nabla\cdot\varphi(y,s)\,dy=0.
\]
Therefore
\[
  \int_0^T\!\int_{\mathbb R^3}P_n\nabla\cdot\varphi\,dy\,ds
  =
  \int_0^T\!\int_K(P_n-(P_n)_K)\nabla\cdot\varphi\,dy\,ds.
\]
By Lemma~\ref{paper5:lem:autonomous-window-compactness},
\[
  P_n-(P_n)_K\rightharpoonup\Pi_K
  \quad\text{in }L^{3/2}((0,T)\times K),
\]
and \(\nabla\cdot\varphi\in L^3((0,T)\times K)\).  The compatibility
Lemma~\ref{paper5:lem:pressure-compatibility} permits us to write the resulting pressure distribution
as \(\Pi\), modulo functions of time.  Hence
\[
  \int_0^T\!\int_{\mathbb R^3}P_n\nabla\cdot\varphi\,dy\,ds
  \to
  \int_0^T\!\int_K\Pi_K\nabla\cdot\varphi\,dy\,ds.
\]
The value of the last integral is independent of the chosen \(K\), because replacing
\(\Pi_K\) by \(\Pi_K+c(s)\) changes the integral by
\(\int_0^Tc(s)\int_K\nabla\cdot\varphi\,dy\,ds=0\).  The limit is therefore the pressure term
\(\int\Pi\nabla\cdot\varphi\).

\medskip
\noindent\textit{(v) Drift term with \(a\): part involving \(V_n\).}
Define
\[
  Y_n^0(s):=\int_{\mathbb R^3}V_n(y,s)\cdot\varphi(y,s)
  \,dy,
  \qquad
  Y^0(s):=\int_{\mathbb R^3}U(y,s)\cdot\varphi(y,s)
  \,dy.
\]
For each \(s\), Cauchy--Schwarz gives
\[
  |Y_n^0(s)|
  \le
  \|V_n(s)\|_{L^2(K_\varphi)}
  \|\varphi(s)\|_{L^2(K_\varphi)}.
\]
By Assumption~\ref{paper5:ass:selected-window-estimates}, there is \(C_0<\infty\)
such that \(\sup_n\|Y_n^0\|_{L^\infty(0,T)}\le C_0\). Moreover,
\[
  Y_n^0\to Y^0\quad\text{in }L^2(0,T),
\]
because \(V_n\to U\) in \(L^2(0,T;L^2(K_\varphi))\).
Now
\[
  \int_0^T\alpha_n(s)Y_n^0(s)
  \,ds
  =\int_0^T(\alpha_n(s)-a_\infty)Y_n^0(s)
  \,ds+a_\infty\int_0^TY_n^0(s)
  \,ds.
\]
For the first term,
\[
  \left|\int_0^T(\alpha_n-a_\infty)Y_n^0\,ds\right|\le
\sup_{s\in(0,T)}|Y_n^0(s)|\,\|\alpha_n-a_\infty\|_{L^1(0,T)}\to0.
\]
For the second term, strong convergence in \(L^2(0,T)\) implies
\(\int_0^TY_n^0\to\int_0^TY^0\), so
\[
  \int_0^T\alpha_n(s)\int_{\mathbb R^3}V_n\cdot\varphi
  \,dy\,ds
  \to a_\infty\int_0^T\int_{\mathbb R^3}U\cdot\varphi
  \,dy\,ds.
\]
Thus the signed contribution in \eqref{paper5:eq:weak-n} satisfies
\[
  -\int_0^T\alpha_n(s)\int_{\mathbb R^3}V_n\cdot\varphi
  \,dy\,ds
  \to -a_\infty\int_0^T\int_{\mathbb R^3}U\cdot\varphi
  \,dy\,ds.
\]

\medskip
\noindent\textit{(vi) Drift term with \(a\): part involving \(y\cdot\nabla V_n\).}
Since \(\varphi\in C_c^\infty\), integrating by parts in \(y\) gives
\[
  \int_{\mathbb R^3}(y\cdot\nabla V_n)\cdot\varphi
  =-\int_{\mathbb R^3}V_n\cdot(y\cdot\nabla\varphi+3\varphi)
  \,dy.
\]
Define
\[
  Y_n^1(s):=\int_{\mathbb R^3}V_n\cdot(y\cdot\nabla\varphi+3\varphi)
  \,dy,
  \qquad
  Y^1(s):=\int_{\mathbb R^3}U\cdot(y\cdot\nabla\varphi+3\varphi)
  \,dy.
\]
For each \(s\), Cauchy--Schwarz gives
\[
  |Y_n^1(s)|
  \le
  \|V_n(s)\|_{L^2(K_\varphi)}
  \|y\cdot\nabla\varphi(s)+3\varphi(s)\|_{L^2(K_\varphi)}.
\]
The \(L^\infty(0,T;L^2(K_\varphi))\) bound in
Assumption~\ref{paper5:ass:selected-window-estimates} gives
\[
  \sup_n\|Y_n^1\|_{L^\infty(0,T)}\le C_1<\infty.
\]
Also,
\[
  \|Y_n^1-Y^1\|_{L^2(0,T)}
  \le
  \|V_n-U\|_{L^2(0,T;L^2(K_\varphi))}
  \|y\cdot\nabla\varphi+3\varphi\|_{L^\infty(0,T;L^2(K_\varphi))}
  \to0.
\]
Therefore
\[
  -\int_0^T\alpha_n(s)
  \int_{\mathbb R^3}(y\cdot\nabla V_n)\cdot\varphi
  \,dy\,ds
  \to
  a_\infty\int_0^T\!\int_{\mathbb R^3}U\cdot(y\cdot\nabla\varphi+3\varphi)
  \,dy\,ds.
\]
Integrating by parts once more in the limit gives
\[
  a_\infty\int_0^T\!\int_{\mathbb R^3}U\cdot(y\cdot\nabla\varphi+3\varphi)
  \,dy\,ds
  =
  -a_\infty\int_0^T\!\int_{\mathbb R^3}(y\cdot\nabla U)\cdot\varphi
  \,dy\,ds.
\]

\medskip
\noindent\textit{(vii) Term with \(b\).}
Since \(\beta_n(s)\) is independent of \(y\),
\[
  \int_{\mathbb R^3}(\beta_n\cdot\nabla V_n)\cdot\varphi
  =-\int_{\mathbb R^3}V_n\cdot(\beta_n\cdot\nabla\varphi)
  \,dy,
\]
so
\[
  -\int_0^T\!\int(\beta_n\cdot\nabla V_n)\cdot\varphi
  \,dy\,ds
  =
  \int_0^T\!\int V_n\cdot(\beta_n\cdot\nabla\varphi)
  \,dy\,ds.
\]
Let
\[
\begin{aligned}
Z_n(s)&:=\left(\int_{\mathbb R^3}V_n\cdot\partial_1\varphi\,dy,\,
\int_{\mathbb R^3}V_n\cdot\partial_2\varphi\,dy,\,
\int_{\mathbb R^3}V_n\cdot\partial_3\varphi\,dy\right),\\
Z(s)&:=\left(\int_{\mathbb R^3}U\cdot\partial_1\varphi\,dy,\,
\int_{\mathbb R^3}U\cdot\partial_2\varphi\,dy,\,
\int_{\mathbb R^3}U\cdot\partial_3\varphi\,dy\right).
\end{aligned}
\]
Then
\[
  \int_{\mathbb R^3}V_n\cdot(\beta_n\cdot\nabla\varphi)\,dy
  =\beta_n(s)\cdot Z_n(s),\qquad
  \int_{\mathbb R^3}U\cdot(b_\infty\cdot\nabla\varphi)\,dy
  =b_\infty\cdot Z(s).
\]
Consequently
\[
  \left|\int_0^T\!(\beta_n-b_\infty)\cdot Z_n\,ds\right|
  \le
  \|\beta_n-b_\infty\|_{L^1(0,T)}\sup_{s\in(0,T)}|Z_n(s)|\to0,
\]
because
\[
  |Z_n(s)|\le
  \|V_n(s)\|_{L^2(K_\varphi)}
  \left(\sum_{j=1}^3\|\partial_j\varphi(s)\|_{L^2(K_\varphi)}^2\right)^{1/2}
\]
and Assumption~\ref{paper5:ass:selected-window-estimates} bounds the right-hand side
uniformly in \(n\) and \(s\). Furthermore,
\[
  \|Z_n-Z\|_{L^2(0,T;\mathbb R^3)}
  \le
  \|V_n-U\|_{L^2(0,T;L^2(K_\varphi))}
  \left(\sum_{j=1}^3\|\partial_j\varphi\|_{L^\infty(0,T;L^2(K_\varphi))}^2\right)^{1/2}
  \to0.
\]
Therefore
\[
  \int_0^T b_\infty\cdot Z_n(s)\,ds
  \to
  \int_0^T b_\infty\cdot Z(s)\,ds.
\]
Hence
\[
  -\int_0^T\!\int (\beta_n\cdot\nabla V_n)\cdot\varphi
  \,dy\,ds
  \to
  \int_0^T\!\int U\cdot(b_\infty\cdot\nabla\varphi)
  \,dy\,ds.
\]
Since \(b_\infty\) is constant in \(y\), integration by parts yields
\[
  \int_0^T\!\int U\cdot(b_\infty\cdot\nabla\varphi)
  \,dy\,ds
  =
  -\int_0^T\!\int (b_\infty\cdot\nabla U)\cdot\varphi
  \,dy\,ds.
\]

\emph{Step 5: assemble the limit.}
Passing to the limit in \eqref{paper5:eq:weak-n} in each term (Steps 4(i)--4(vii)) gives
\[
  \int_0^T\!\int
  \bigl(
  -U\cdot\partial_s\varphi
  -(U\otimes U):\nabla\varphi
  +\nu\nabla U:\nabla\varphi
  -\Pi\nabla\cdot\varphi
  -a_\infty(U+y\cdot\nabla U)\cdot\varphi
  -(b_\infty\cdot\nabla U)\cdot\varphi
  \bigr)
  \,dy\,ds=0.
\]
Since this holds for every compactly supported vector field \(\varphi\), equation~\eqref{paper5:eq:autonomous-limit} holds in distributions on
\((0,T)\times\mathbb R^3\).
\end{proof}

\subsection{Decomposition of alternatives of the scale-collapse branch}
\label{paper5:sec:branch-closure}

The cost argument in \cref{paper5:sec:cost,paper5:sec:finite-cost} excludes the finite-cost scale-collapse branch.  The remaining
scale-collapse obstruction is treated by a decomposition of alternatives, following the local
estimate and autonomous-reduction argument of the preceding sections.  The conclusion separates the excluded branches from the residual alternatives; it does not assert a Liouville theorem for every autonomous negative-drift equation.

\begin{theorem}[Selected core retention]\label{paper5:thm:core-retention}
Let a represented Type II branch be obtained from the selected-core construction of the preceding sections.  Then there exist \(R>0\), a cutoff \(\chi\in C_c^\infty(B_R)\), \(\chi\ge0\), a
constant \(\eta>0\), and a terminal time \(\tau_1\) such that, along every selected window used
below,
\[
  \int_{\mathbb R^3}|V(y,\tau)|^2\chi(y)\,dy\ge\eta
  \quad\text{for a.e. }\tau\ge\tau_1.
\]
In particular, every translated sequence \(V_n(y,s)=V(y,\tau_n+s)\) satisfies the localized
nontriviality condition in Assumption~\ref{paper5:ass:selected-window-estimates}, after
passing to a tail of the sequence.
\end{theorem}

\begin{proof}
The selected-core construction chooses a compact renormalized core carrying a fixed positive
amount of localized kinetic mass.  Since the cutoff \(\chi\) is supported strictly inside that
core and equals a positive spatial weight there, the selected-core lower bound gives the stated
estimate.  Translating \(\tau=\tau_n+s\) gives the final assertion on every finite selected
window.
\end{proof}

\begin{theorem}[Autonomous modulation selection]\label{paper5:thm:autonomous-modulation-selection}
Let a represented scale-collapse branch satisfy the modulation bounds obtained in the repaired
gauge construction.  Suppose that a sequence of windows \([\tau_n,\tau_n+T]\) is selected in the
autonomous regime.  Then, after passing to a subsequence, there exist constants
\(a_\infty\in\mathbb R\) and \(b_\infty\in\mathbb R^3\) such that
\[
  a(\tau_n+\cdot)\to a_\infty
  \quad\text{in }L^1(0,T),
  \qquad
  b(\tau_n+\cdot)\to b_\infty
  \quad\text{in }L^1(0,T;\mathbb R^3).
\]
\end{theorem}

\begin{proof}
The autonomous-window selection in the repaired-gauge modulation analysis gives convergence
of the modulation functions to constant autonomous parameters on the fixed window length
\(T\).  Passing to the corresponding subsequence gives the displayed
\(L^1\)-convergences.
\end{proof}

\begin{theorem}[Scale-collapse decomposition of alternatives]\label{paper5:thm:scale-collapse-stratification}
Let a represented Type II branch have genuine scale collapse and nonvanishing selected core.
Then exactly one of the following alternatives occurs after passing to selected terminal
windows:
\begin{enumerate}[label=(\roman*)]
  \item the branch has finite scale-collapsing cost, which is excluded by
  Theorem~\ref{paper5:thm:cost-exclusion};
  \item the branch has finite absolute cost, which is excluded by
  Corollary~\ref{paper5:cor:absolute-cost-exclusion};
  \item the branch admits thick autonomous windows in the sense of
  Definition~\ref{paper5:def:thick-autonomous-windows}, satisfies the selected-window estimates of
  Assumption~\ref{paper5:ass:selected-window-estimates}, and has autonomous modulation convergence;
  \item the branch has a thin-drift alternative: no selected fixed-length window carries a uniform
  positive lower bound for the averaged weighted negative drift;
  \item the branch has an autonomous-modulation alternative: thick windows exist, but no subsequence
  has a constant-coefficient \(L^1\) modulation limit in the sense needed for
  Theorem~\ref{paper5:thm:autonomous-limit};
  \item the branch has a compactness alternative: the local estimates needed for
  Assumption~\ref{paper5:ass:selected-window-estimates} fail on the selected windows.
\end{enumerate}
\end{theorem}

\begin{proof}
The scale-collapse selection theorem decomposes terminal scale-collapse windows according to the
averaged negative-drift mass
\[
  \frac1T\int_{\tau_n}^{\tau_n+T}a_-(\tau)M_R(\tau)\,d\tau.
\]
If the accumulated cost is finite, alternatives (i) or (ii) apply according to the cost
functional under consideration.  If the averaged negative-drift mass stays bounded below on a
subsequence, one obtains thick windows; otherwise the branch lies in the thin-drift alternative.
On thick windows, either the local estimates and core retention hold, in which case
Proposition~\ref{paper5:prop:selected-window-derived} and Theorem~\ref{paper5:thm:core-retention} give
Assumption~\ref{paper5:ass:selected-window-estimates}, or the branch lies in the compactness alternative.  If
the estimates hold but no subsequence has the required constant-coefficient modulation limit,
the branch lies in the autonomous-modulation alternative.  The remaining case is the thick
autonomous alternative.
\end{proof}

\begin{definition}[Stationary omega-limit class]\label{paper5:def:stationary-omega-class}
Let \((U,\Pi)\) be a nonzero autonomous limit produced by
Theorem~\ref{paper5:thm:autonomous-limit}.  We say that it has a stationary omega-limit in the
admissible class if there exists a nonzero pair
\[
  (W,Q),\qquad W\in L^3(\mathbb R^3),
\]
with the local energy and pressure reconstruction properties inherited from the selected
windows, such that
\[
  (W\cdot\nabla)W+\nabla Q
  =
  \nu\Delta W+a_\infty(W+y\cdot\nabla W)+b_\infty\cdot\nabla W,
  \qquad
  \nabla\cdot W=0,
\]
in \(\mathcal D'(\mathbb R^3)\).
\end{definition}

\begin{definition}[Rigidity-covered stationary class]\label{paper5:def:rigidity-covered-class}
For a parameter pair \((a_\infty,b_\infty)\), the stationary class is rigidity-covered if every
admissible stationary solution in Definition~\ref{paper5:def:stationary-omega-class} is either zero or
is not a terminal Type II branch.  In the classical backward self-similar normalization this is
the role of the Nečas--Růžička--Šverák rigidity theorem, after the constant translation drift is
removed \cite{NecasRuzickaSverak1996}.  Other parameter regimes require their own stated
Liouville or classification theorem.
\end{definition}

\begin{theorem}[Stationary-rigidity exclusion]\label{paper5:thm:stationary-rigidity-blocker}
Assume a thick autonomous branch produces a nonzero autonomous limit by
Theorem~\ref{paper5:thm:autonomous-limit}.  If that autonomous limit has a stationary omega-limit in
the admissible class and the corresponding stationary class is rigidity-covered, then the branch
cannot occur as a terminal Type II scale-collapse branch.
\end{theorem}

\begin{proof}
The stationary omega-limit gives a nonzero stationary profile \((W,Q)\) in the admissible class.
The rigidity-covered hypothesis says that no such nonzero profile can represent a terminal Type
II branch.  This contradicts membership in the terminal scale-collapse Type II class.
\end{proof}

\begin{theorem}[Scale-collapse alternatives]\label{paper5:thm:scale-collapse-alternatives}
Let a represented Type II branch have genuine scale collapse and a nonvanishing selected core.
Then one of the following mutually exclusive outcomes occurs:
\begin{enumerate}[label=(\roman*)]
  \item the finite scale-collapsing cost exclusion applies;
  \item the finite absolute cost exclusion applies;
  \item the branch has a thin-drift alternative;
  \item the branch has an autonomous-modulation alternative;
  \item the branch has a compactness alternative;
  \item the branch has a nonzero autonomous reduced limit, and either stationary rigidity
  excludes it or the remaining obstruction is the corresponding stationary/generalized
  self-similar Liouville problem.
\end{enumerate}
\end{theorem}

\begin{proof}
Apply \cref{paper5:thm:scale-collapse-stratification}.  The finite-cost alternatives are
items (i) and (ii).  The thin-drift, modulation, and compactness alternatives are items (iii)--(v).  In the
remaining thick autonomous case, Lemma~\ref{paper5:lem:negative-parameter} gives \(a_\infty<0\), and
Theorem~\ref{paper5:thm:autonomous-limit} produces a nonzero autonomous reduced limit.  If
Definitions~\ref{paper5:def:stationary-omega-class} and~\ref{paper5:def:rigidity-covered-class} apply, then
Theorem~\ref{paper5:thm:stationary-rigidity-blocker} excludes the branch.  If not, the remaining
obstruction is the stationary or generalized self-similar Liouville problem for the
parameter pair \((a_\infty,b_\infty)\).
\end{proof}

\clearpage
\section{Critical Tightness and Windowed-Control Estimates}
\label{sec:c6-c7-cost-exclusion}

This section records the local estimates used to exclude failure of uniform
critical tightness and failure of local windowed \(H^1\) control.  The
statements are phrased as analytic alternatives for a renormalized
Navier--Stokes solution \(V(\tau)\).

\subsection{Uniform critical tightness}

\begin{definition}[Uniform critical tightness]\label{paper7:def:critical-tightness}
A renormalized orbit \(V(\tau)\), \(\tau\ge\tau_0\), is uniformly critically
tight if
\[
  \forall\varepsilon>0\ \exists R_\varepsilon<\infty:
  \sup_{\tau\ge\tau_0}
  \int_{|y|>R_\varepsilon}|V(y,\tau)|^3\,dy<\varepsilon.
\]
The non-tight alternative is the failure of this property.
\end{definition}

\begin{lemma}[Critical precompactness implies uniform tightness]
\label{paper7:lem:global-l3-compactness-tightness}
If
\[
  \mathcal O:=\{V(\tau):\tau\ge\tau_0\}
\]
is contained in a compact subset of \(L^3(\mathbb R^3)\), then \(V\) is
uniformly critically tight.
\end{lemma}

\begin{proof}
Compact subsets of \(L^3(\mathbb R^3)\) are uniformly tight.  Let
\(K\subset L^3\) be compact and fix \(\varepsilon>0\).  Set
\(\delta=\varepsilon^{1/3}/2\).  Choose a finite \(\delta\)-net
\(f_1,\dots,f_N\) in \(L^3\).  For each \(i\), choose \(R_i\) so that
\[
  \|f_i\|_{L^3(|y|>R_i)}<\delta.
\]
Let \(R=\max_iR_i\).  For any \(f\in K\), choose \(i\) with
\(\|f-f_i\|_3<\delta\).  Then
\[
  \|f\|_{L^3(|y|>R)}
  \le
  \|f-f_i\|_3+\|f_i\|_{L^3(|y|>R)}
  <\varepsilon^{1/3}.
\]
Taking the cube gives
\[
  \int_{|y|>R}|f(y)|^3\,dy<\varepsilon.
\]
Applying this to the compact set containing \(\mathcal O\) gives uniform
critical tightness.
\end{proof}

\begin{definition}[Compact core with vanishing critical remainder]
\label{paper7:def:compact-core-vanishing-remainder}
The orbit has a compact critical core with vanishing remainder if there exist
a compact set \(\mathcal K^{\rm c}\subset L^3(\mathbb R^3)\), a map
\(Q:[\tau_0,\infty)\to \mathcal K^{\rm c}\), and a remainder
\(r(\tau)\in L^3(\mathbb R^3)\), such that
\[
  V(\tau)=Q(\tau)+r(\tau),\qquad
  \lim_{\tau\to\infty}\|r(\tau)\|_{L^3}=0,
\]
and, for every \(T>\tau_0\), the set
\[
  \{V(\tau):\tau\in[\tau_0,T]\}
\]
is uniformly \(L^3\)-tight.
\end{definition}

\begin{lemma}[Compact core plus vanishing remainder implies tightness]
\label{paper7:lem:core-remainder-tightness}
If \cref{paper7:def:compact-core-vanishing-remainder} holds, then \(V\) is
uniformly critically tight.
\end{lemma}

\begin{proof}
Fix \(\varepsilon>0\) and set \(\delta=\varepsilon^{1/3}\).  By compactness of
\(\mathcal K^{\rm c}\), apply
\cref{paper7:lem:global-l3-compactness-tightness} to \(\mathcal K^{\rm c}\)
and choose \(R_1\) such that
\[
  \sup_{q\in \mathcal K^{\rm c}}\|q\|_{L^3(|y|>R_1)}<\delta/3.
\]
Since \(\|r(\tau)\|_3\to0\), choose \(T\) so that
\[
  \|r(\tau)\|_3<\delta/3
  \qquad\text{for all }\tau\ge T.
\]
Then for \(\tau\ge T\),
\[
  \|V(\tau)\|_{L^3(|y|>R_1)}
  \le
  \|Q(\tau)\|_{L^3(|y|>R_1)}+\|r(\tau)\|_3
  <2\delta/3.
\]
By finite-initial-window tightness on \([\tau_0,T]\), choose \(R_2\) such that
\[
  \sup_{\tau\in[\tau_0,T]}
  \int_{|y|>R_2}|V(y,\tau)|^3\,dy<\varepsilon.
\]
With \(R=\max(R_1,R_2)\), the \(L^3\)-norm of the tail of \(V(\tau)\) is
\(<\delta\) for all \(\tau\ge T\), while the tail integral is
\(<\varepsilon\) on \([\tau_0,T]\).  Cubing on the tail \(\tau\ge T\) gives
the required \(L^3\)-tail integral \(<\varepsilon\) for all
\(\tau\ge\tau_0\).
\end{proof}

\begin{definition}[No-splitting profile condition]\label{paper7:def:no-splitting-condition}
A critical profile decomposition has no critical-mass splitting if it has:
\begin{enumerate}[label=\textup{(\roman*)}]
\item a primary compact core family \(\mathcal K^{\rm c}\subset L^3\);
\item asymptotic \(L^3\)-decoupling of core, secondary profiles, and
remainder;
\item single-profile saturation, so every non-primary profile has zero
critical \(L^3\)-mass;
\item a remainder whose \(L^3\)-norm vanishes on the renormalized tail;
\item finite-initial-window tightness.
\end{enumerate}
\end{definition}

\begin{lemma}[No splitting gives compact core with vanishing remainder]
\label{paper7:lem:no-splitting-gives-core-rem}
If \cref{paper7:def:no-splitting-condition} holds, then
\cref{paper7:def:compact-core-vanishing-remainder} holds.
\end{lemma}

\begin{proof}
By single-profile saturation and nonnegative \(L^3\)-mass decoupling, every
secondary profile has zero critical \(L^3\)-mass and therefore vanishes in
\(L^3\).  The profile decomposition then reduces to a primary compact core
plus a remainder whose \(L^3\)-norm tends to zero on the renormalized tail.
These provide the core \(Q(\tau)\), compact set \(\mathcal K^{\rm c}\), and
remainder \(r(\tau)\) required by
\cref{paper7:def:compact-core-vanishing-remainder}, with
finite-initial-window tightness supplied as part of
\cref{paper7:def:no-splitting-condition}.
\end{proof}

\begin{corollary}[No splitting gives uniform critical tightness]
\label{paper7:cor:no-splitting-critical-tightness}
If \cref{paper7:def:no-splitting-condition} holds, then \(V\) is uniformly
critically tight.
\end{corollary}

\begin{proof}
\Cref{paper7:lem:no-splitting-gives-core-rem} gives a compact core with
vanishing critical remainder, and
\cref{paper7:lem:core-remainder-tightness} gives uniform critical tightness.
\end{proof}

\begin{definition}[Finite critical profile approximation]
\label{paper7:def:finite-profile-tight}
The orbit admits a finite critical profile approximation if there are finitely
many profiles
\[
  \mathcal B^{II}_{NS}:=\{B_1,\dots,B_N\}\subset L^3(\mathbb R^3)
\]
such that, for every \(\delta>0\), there is \(\tau_\delta\) with
\[
  \forall\tau\ge\tau_\delta\ \exists i(\tau)\in\{1,\dots,N\}:
  \|V(\tau)-B_{i(\tau)}\|_{L^3}<\delta,
\]
and the finite initial window
\[
  \{V(\tau):\tau\in[\tau_0,\tau_\delta]\}
\]
is uniformly \(L^3\)-tight for each \(\delta>0\).
\end{definition}

\begin{lemma}[Finite critical profile approximation implies tightness]
\label{paper7:lem:finite-profile-tight}
If \cref{paper7:def:finite-profile-tight} holds, then \(V\) is uniformly
critically tight.
\end{lemma}

\begin{proof}
Fix \(\varepsilon>0\) and set \(\delta=\varepsilon^{1/3}\).  Since the profile family is finite and each
\(B_i\in L^3(\mathbb R^3)\), there is \(R_1\) such that
\[
  \max_{1\le i\le N}\|B_i\|_{L^3(|y|>R_1)}<\delta/3.
\]
Choose the approximation tolerance \(\delta/3\).  For
\(\tau\ge\tau_{\delta/3}\), pick \(i(\tau)\) with
\(\|V(\tau)-B_{i(\tau)}\|_3<\delta/3\).  Then
\[
  \|V(\tau)\|_{L^3(|y|>R_1)}
  \le
  \|V(\tau)-B_{i(\tau)}\|_3
  +\|B_{i(\tau)}\|_{L^3(|y|>R_1)}
  <2\delta/3.
\]
By finite-initial-window tightness, choose \(R_2\) so that the \(L^3\) tail is
\(<\varepsilon\) on \([\tau_0,\tau_{\delta/3}]\).  With
\(R=\max(R_1,R_2)\), the tail integral is \(<\varepsilon\) for every
\(\tau\ge\tau_0\), so the orbit is uniformly \(L^3\)-tight.
\end{proof}

\begin{definition}[Compact parameter realization]\label{paper7:def:compact-param-realization}
The orbit has a compact parameter realization if there exist a compact
parameter set \(\Theta_{II}\), a continuous map
\[
  \Theta_{II}\ni\theta\mapsto V_\theta\in L^3(\mathbb R^3),
\]
a curve \(\theta(\tau)\in\Theta_{II}\), and a remainder \(r(\tau)\to0\) in
\(L^3\), such that
\[
  V(\tau)=V_{\theta(\tau)}+r(\tau),
\]
and, for every \(T>\tau_0\), the set
\[
  \{V(\tau):\tau\in[\tau_0,T]\}
\]
is uniformly \(L^3\)-tight.
\end{definition}

\begin{lemma}[Compact realization implies tightness]
\label{paper7:lem:compact-realization-tight}
If \cref{paper7:def:compact-param-realization} holds, then \(V\) is uniformly
critically tight.
\end{lemma}

\begin{proof}
The continuous image of compact \(\Theta_{II}\) in \(L^3\) is compact.  Thus
\[
  \mathcal K^{\rm c}:=\{V_\theta:\theta\in\Theta_{II}\}
\]
is compact in \(L^3\).  The representation
\(V(\tau)=V_{\theta(\tau)}+r(\tau)\) with \(r(\tau)\to0\) in \(L^3\) gives
the compact-core vanishing-remainder situation, because finite-initial-window
tightness is included in \cref{paper7:def:compact-param-realization}.
\Cref{paper7:lem:core-remainder-tightness} gives uniform critical tightness.
\end{proof}

\begin{theorem}[Exclusion of the non-tight alternative]\label{paper7:thm:critical-tightness-criteria}
For a renormalized Type II Navier--Stokes solution, any one of the following
conditions gives uniform critical tightness:
\begin{enumerate}[label=\textup{(\roman*)}]
\item \(\mathcal O\Subset L^3(\mathbb R^3)\);
\item compact core with vanishing critical remainder;
\item no-splitting profile condition;
\item finite critical profile approximation;
\item compact parameter realization.
\end{enumerate}
\end{theorem}

\begin{proof}
The five alternatives are handled respectively by
\cref{paper7:lem:global-l3-compactness-tightness,paper7:lem:core-remainder-tightness,paper7:cor:no-splitting-critical-tightness,paper7:lem:finite-profile-tight,paper7:lem:compact-realization-tight}.
In every case \(V\) is uniformly critically tight, so the failure of uniform
critical tightness is unavailable.
\end{proof}

\begin{corollary}[After exclusion of the non-tight alternative]\label{paper7:cor:critical-tightness-reduces-alternatives}
Assume a finite-cost renormalized Type II solution not excluded by the compact
cost-divergence criterion has a two-alternative classification whose only
remaining alternatives are failure of uniform critical tightness and failure of
tail windowed \(H^1\) control.  If
\cref{paper7:def:critical-tightness} holds, then the only remaining alternative is
failure of tail windowed \(H^1\) control.
\end{corollary}

\begin{proof}
The two-alternative classification gives exactly one of the two alternatives.  Uniform
critical tightness rules out failure of tightness, so only the windowed
\(H^1\)-failure alternative remains.
\end{proof}

\subsection{\texorpdfstring{Tail windowed \(H^1\) control}{Tail windowed H1 control}}

\begin{lemma}[Uniform local pressure control]\label{paper7:lem:uniform-pressure-from-c3}
Assume
\[
  M_{3,R}:=\sup_{\tau\ge\tau_0}\|V(\tau)\|_{L^3(B_{4R})}<\infty
\]
and
\[
  -\Delta P=\partial_i\partial_j(V_iV_j)
  \quad\text{in }\mathbb R^3
\]
for each \(\tau\ge\tau_0\).  Assume also the local tail estimate
\[
  \mathcal T_R(\tau):=
  \sup_{x\in B_R}\int_{|z|>2R}
  \frac{|V(z,\tau)|^2}{|x-z|^5}\,dz,
  \qquad
  \sup_{\tau\ge\tau_0}\mathcal T_R(\tau)\le CR^{-4}M_{3,R}^2 .
\]
Then there is a measurable choice of constants \(c_R(\tau)\in\mathbb R\) such
that
\[
  \sup_{\tau\ge\tau_0}
  \|P(\tau)-c_R(\tau)\|_{L^{3/2}(B_R)}
  \le C_RM_{3,R}^2 .
\]
\end{lemma}

\begin{proof}
The local pressure decomposition gives
\[
  \|P(\tau)-c_R(\tau)\|_{L^{3/2}(B_R)}
  \le
  C\left(
  \|V(\tau)\|_{L^3(B_{4R})}^2+R^2\mathcal T_R(\tau)
  \right).
\]
The tail hypothesis gives
\[
  \sup_{\tau\ge\tau_0}\mathcal T_R(\tau)\le CR^{-4}M_{3,R}^2,
\]
and the local critical bound gives
\[
  \|V(\tau)\|_{L^3(B_{4R})}\le M_{3,R}.
\]
Combining the two estimates yields
\[
  \sup_{\tau\ge\tau_0}
  \|P(\tau)-c_R(\tau)\|_{L^{3/2}(B_R)}
  \le C_RM_{3,R}^2 .
\]
The constants \(c_R(\tau)\) can be chosen by fixing the displayed pressure
decomposition, hence measurably.
\end{proof}

\begin{proposition}[Renormalized Caccioppoli gives uniform windowed gradients]
\label{paper7:prop:uniform-windowed-gradient}
Assume \(V,P\) are smooth on compact renormalized cylinders and solve
\[
  \partial_\tau V +(V\cdot\nabla)V+\nabla P
  =
  \nu\Delta V+a(\tau)(V+y\cdot\nabla V)+b(\tau)\cdot\nabla V,
  \qquad \nabla\cdot V=0.
\]
Assume local compact-cylinder critical control on the outer cylinders used in
the estimate, bounded modulation
\[
  M_{ab}:=\sup_{\tau\ge\tau_0}(|a(\tau)|+|b(\tau)|)<\infty,
\]
and pressure normalized as in \cref{paper7:lem:uniform-pressure-from-c3}.  Then
for every \(R>0\) there exists
\[
  C_R=C(R,\nu,M_{3,R},M_{ab})
\]
such that
\[
  \sup_{T\ge\tau_0+1}
  \int_T^{T+1}\int_{B_R}|\nabla V(y,\tau)|^2\,dy\,d\tau
  \le C_R.
\]
\end{proposition}

\begin{proof}
Fix \(R>0\) and \(T\ge\tau_0+1\).  Use the renormalized Caccioppoli estimate
with outer radius \(2R\), inner radius \(R\), and time intervals
\[
  I_{2R}:=(T-1,T+2),\qquad I_R:=[T,T+1].
\]
Thus \(R_{\rm out}-R_{\rm in}=R\) and the inner time length is \(1\).  Keeping
the pressure term before taking absolute values gives
\[
\begin{aligned}
\nu\int_T^{T+1}\int_{B_R}|\nabla V|^2
\le{}&
C_R\iint_{(T-1,T+2)\times B_{2R}}|V|^2
+C_R\iint_{(T-1,T+2)\times B_{2R}}|V|^3\\
&+C_R\mathcal P_{T,R}
+C_RM_{ab}\iint_{(T-1,T+2)\times B_{2R}}|V|^2,
\end{aligned}
\]
where \(C_R\) absorbs the powers of \(R^{-1}\), \(R^{-2}\), and \(\nu\), and
\[
  \mathcal P_{T,R}:=
  \left|
  \iint_{(T-1,T+2)\times B_{2R}}P\,V\cdot\nabla\zeta
  \right|.
\]
The \(L^2\)-term is controlled by the local critical bound and boundedness of
\(B_{2R}\):
\[
  \int_{B_{2R}}|V(\tau)|^2\,dy
  \le |B_{2R}|^{1/3}\|V(\tau)\|_{L^3(B_{2R})}^2
  \le C_RM_{3,R}^2.
\]
Since the outer time interval has length \(3\),
\[
  \iint_{(T-1,T+2)\times B_{2R}}|V|^2\le C_RM_{3,R}^2,
\]
and similarly
\[
  \iint_{(T-1,T+2)\times B_{2R}}|V|^3\le 3M_{3,R}^3.
\]
For any time-dependent constant \(c(\tau)\),
\[
  \int_{B_{2R}}c(\tau)V\cdot\nabla\zeta\,dy
  =
  -\int_{B_{2R}}c(\tau)\zeta\,\nabla\cdot V\,dy=0.
\]
Thus \(P\) may be replaced by \(P-c_{2R}(\tau)\).  Holder's inequality gives
\[
  \mathcal P_{T,R}
  \le
  C_R
  \iint_{(T-1,T+2)\times B_{2R}}
  |P-c_{2R}(\tau)|\,|V|\,dy\,d\tau .
\]
For almost every \(\tau\),
\[
  \int_{B_{2R}}|P-c_{2R}(\tau)|\,|V(\tau)|\,dy
  \le
  \|P(\tau)-c_{2R}(\tau)\|_{L^{3/2}(B_{2R})}
  \|V(\tau)\|_{L^3(B_{2R})}.
\]
The pressure estimate with radius \(2R\) gives
\[
  \|P(\tau)-c_{2R}(\tau)\|_{L^{3/2}(B_{2R})}
  \le C_RM_{3,R}^2,
\]
and the local critical bound gives
\[
  \|V(\tau)\|_{L^3(B_{2R})}\le M_{3,R}.
\]
Hence \(\mathcal P_{T,R}\le C_RM_{3,R}^3\).  Substituting these estimates into
Caccioppoli yields
\[
  \nu\int_T^{T+1}\int_{B_R}|\nabla V|^2
  \le
  C_R\left(M_{3,R}^2+M_{3,R}^3+M_{ab}M_{3,R}^2\right),
\]
uniformly in \(T\ge\tau_0+1\).  Dividing by \(\nu\) and renaming the constant
proves the claim.
\end{proof}

\begin{corollary}[Tail windowed local \(H^1\) bound]\label{paper7:cor:windowed-h1}
Under the hypotheses of \cref{paper7:prop:uniform-windowed-gradient}, for
every \(R>0\),
\[
  \sup_{T\ge\tau_0+1}
  \int_T^{T+1}\|V(\tau)\|_{H^1(B_R)}^2\,d\tau<\infty.
\]
\end{corollary}

\begin{proof}
The gradient part is \cref{paper7:prop:uniform-windowed-gradient}.  The
\(L^2\)-part follows from the same bounded-set Holder estimate:
\[
  \int_T^{T+1}\|V(\tau)\|_{L^2(B_R)}^2\,d\tau
  \le
  C_R\int_T^{T+1}\|V(\tau)\|_{L^3(B_{2R})}^2\,d\tau
  \le C_RM_{3,R}^2.
\]
Adding the two bounds gives the asserted \(H^1(B_R)\) estimate.
\end{proof}

\begin{corollary}[Exclusion of failure of local windowed \(H^1\) control]
\label{paper7:cor:windowed-h1-exclusion-tail}
Assume the renormalized Type II solution satisfies the repaired-gauge
renormalized Navier--Stokes equation on compact cylinders, local
compact-cylinder critical control
\[
  \sup_{\tau\ge\tau_0}\|V(\tau)\|_{L^3(B_{2R})}<\infty
\]
for each relevant \(R\), bounded modulation
\[
  \sup_{\tau\ge\tau_0}(|a(\tau)|+|b(\tau)|)<\infty,
\]
and pressure reconstruction
\[
  -\Delta P=\partial_i\partial_j(V_iV_j).
\]
Then failure of local windowed \(H^1\) control,
\[
  \exists m\ge1:
  \sup_{T\ge\tau_0+1}
  \int_T^{T+1}\|V(\tau)\|_{H^1(B_m)}^2\,d\tau=\infty
\]
is impossible.
\end{corollary}

\begin{proof}
Apply \cref{paper7:cor:windowed-h1} with \(R=m\).  The resulting finite bound
contradicts the displayed failure of local windowed \(H^1\) control.
\end{proof}

\subsection{Modulation and weighted forcing inputs}

\begin{lemma}[Positive finite critical annulus implies bounded critical norm]
\label{paper7:lem:l3-annulus-bounds-l3}
If there exist \(0<\eta\le M<\infty\) such that
\[
  \eta\le\|V(\tau)\|_{L^3(\mathbb R^3)}\le M
  \qquad\forall\tau\ge\tau_0,
\]
then
\[
  \sup_{\tau\ge\tau_0}\|V(\tau)\|_{L^3(\mathbb R^3)}\le M<\infty.
\]
\end{lemma}

\begin{proof}
This is obtained by taking the supremum in \(\tau\) in the displayed upper
bound.
\end{proof}

\begin{definition}[Modulation matrix Schur data]\label{paper7:def:modmatrix-schur}
Write the repaired modulation matrix in block form
\[
  M(V)=
  \begin{pmatrix}
  \alpha(V)&u(V)^T\\
  v(V)&A(V)
  \end{pmatrix}.
\]
Assume
\[
  \alpha(V(\tau))\ge \alpha_0>0,\qquad
  \|A(V(\tau))^{-1}\|_{\infty\to\infty}\le C_A,
\]
\[
  \|u(V(\tau))\|_\infty\le C_u',
  \qquad
  \|v(V(\tau))\|_\infty\le C_v',
\]
and
\[
  S(V(\tau)):=
  \alpha(V(\tau))-u(V(\tau))^TA(V(\tau))^{-1}v(V(\tau))
  \ge\sigma_0>0.
\]
\end{definition}

\begin{lemma}[Uniform inverse from Schur data]\label{paper7:lem:modmatrix-inverse}
Under \cref{paper7:def:modmatrix-schur},
\[
  \sup_{\tau\ge\tau_0}\|M(V(\tau))^{-1}\|_{\infty\to\infty}
  \le
  \sigma_0^{-1}
  +3\sigma_0^{-1}C_u'C_A
  +\sigma_0^{-1}C_A C_v'
  +C_A
  +3C_A^2C_v'C_u'\sigma_0^{-1}.
\]
\end{lemma}

\begin{proof}
For each \(\tau\), write
\[
  M(V(\tau))=
  \begin{pmatrix}
  \alpha&u^T\\
  v&A
  \end{pmatrix},
  \qquad
  S=\alpha-u^TA^{-1}v.
\]
By the translation-block inverse bound, \(A^{-1}\) exists and is uniformly
bounded.  By the Schur gap, \(S^{-1}\) exists and
\(|S^{-1}|\le\sigma_0^{-1}\).  The block inverse formula gives
\[
  M^{-1}=
  \begin{pmatrix}
  S^{-1} & -S^{-1}u^TA^{-1}\\
  -A^{-1}vS^{-1} & A^{-1}+A^{-1}vS^{-1}u^TA^{-1}
  \end{pmatrix}.
\]
Using the \(\ell^\infty\) matrix norm and the finite dimension of the mixed
blocks,
\[
  \|S^{-1}u^TA^{-1}\|_{\infty\to\infty}
  \le 3\sigma_0^{-1}C_u'C_A,
\]
\[
  \|A^{-1}vS^{-1}\|_{\infty\to\infty}
  \le \sigma_0^{-1}C_A C_v',
\]
and
\[
  \|A^{-1}vS^{-1}u^TA^{-1}\|_{\infty\to\infty}
  \le 3C_A^2C_v'C_u'\sigma_0^{-1}.
\]
Adding the four block estimates gives the displayed uniform bound.
\end{proof}

\begin{lemma}[Scale-normalization nondegeneracy gives Schur data]
\label{paper7:lem:gauge-nondeg-schur}
Assume the scale normalization constraint is satisfied, the scale derivative obeys
\[
  \alpha(V)=DG_{\rm sc}(V)[Z_{\rm sc}(V)]=p\Theta_0>0,
\]
the translation block has a uniform inverse bound, and the mixed blocks obey
the weighted first-moment bound and the Schur-compatible estimate
\[
  \alpha(V(\tau))-u(V(\tau))^TA(V(\tau))^{-1}v(V(\tau))
  \ge \frac{\alpha_0}{2}.
\]
Then the hypotheses of \cref{paper7:def:modmatrix-schur} hold with
\(\sigma_0=\alpha_0/2\).
\end{lemma}

\begin{proof}
The scale-normalization identity gives
\(\alpha(V)=DG_{\rm sc}(V)[Z_{\rm sc}(V)]=p\Theta_0\), hence the scale
transversality bound.  The translation-block inverse hypothesis gives the
uniform bound on \(A^{-1}\).  The mixed-block hypotheses give the bounds on
\(u\) and \(v\), and the displayed estimate is the Schur gap with
\(\sigma_0=\alpha_0/2\).
\end{proof}

\begin{definition}[Weighted scale-force components]\label{paper7:def:scale-force-pieces}
Let \(0<p<3\) and set
\[
  H(V,P)=\nu\Delta V-(V\cdot\nabla)V-\nabla P.
\]
Define
\[
  I_\Delta(\tau):=\int_{\mathbb R^3}|y|^{-p}|V|V\cdot\Delta V\,dy,
\]
\[
  I_{\rm nl}(\tau):=\int_{\mathbb R^3}|y|^{-p}|V|V\cdot((V\cdot\nabla)V)\,dy,
\]
and
\[
  I_P(\tau):=\int_{\mathbb R^3}|y|^{-p}|V|V\cdot\nabla P\,dy,
\]
whenever these integrals are absolutely convergent or are defined by the
distributional pairing used for the renormalized solution.
\end{definition}

\begin{lemma}[Scale-force decomposition]\label{paper7:lem:scale-force-decomp}
If
\[
  \sup_\tau|I_\Delta(\tau)|\le C_\Delta,\qquad
  \sup_\tau|I_{\rm nl}(\tau)|\le C_{\rm nl},\qquad
  \sup_\tau|I_P(\tau)|\le C_P,
\]
then
\[
  \sup_{\tau\ge\tau_0}
  \left|
  \int |y|^{-p}|V|V\cdot H(V,P)\,dy
  \right|
  \le \nu C_\Delta+C_{\rm nl}+C_P.
\]
\end{lemma}

\begin{proof}
By the definition of \(H\),
\[
  \int |y|^{-p}|V|V\cdot H(V,P)\,dy
  =
  \nu I_\Delta-I_{\rm nl}-I_P.
\]
The triangle inequality gives the displayed uniform bound.
\end{proof}

\begin{lemma}[Weighted fourth moment controls the nonlinear scale component]
\label{paper7:lem:scale-nonlinear-from-l4}
Assume \(V(\tau)\) is divergence-free and belongs to the approximation class
needed to justify the weighted integration by parts.  If
\[
  \sup_{\tau\ge\tau_0}
  \int_{\mathbb R^3}|y|^{-p-1}|V(y,\tau)|^4\,dy<\infty,
\]
then
\[
  |I_{\rm nl}(\tau)|
  \le
  \frac{p}{3}\int_{\mathbb R^3}|y|^{-p-1}|V(y,\tau)|^4\,dy.
\]
\end{lemma}

\begin{proof}
For smooth compactly supported approximants,
\[
  |V|V\cdot((V\cdot\nabla)V)=\frac13 V\cdot\nabla(|V|^3).
\]
Since \(\nabla\cdot V=0\), integration by parts gives
\[
  I_{\rm nl}
  =\frac13\int |y|^{-p}V\cdot\nabla(|V|^3)\,dy
  =-\frac13\int |V|^3V\cdot\nabla(|y|^{-p})\,dy.
\]
Because
\[
  \nabla(|y|^{-p})=-p|y|^{-p-2}y,
\]
we obtain
\[
  I_{\rm nl}
  =\frac{p}{3}\int |y|^{-p-2}(V\cdot y)|V|^3\,dy.
\]
Therefore
\[
  |I_{\rm nl}|
  \le
  \frac{p}{3}\int |y|^{-p-1}|V|^4\,dy.
\]
Taking the supremum in \(\tau\) proves the assertion.  The general admissible
case follows by the same approximation convention used for the weighted scale
normalization.
\end{proof}

\begin{corollary}[Pointwise scale-force bound]\label{paper7:cor:scale-force-bound}
If
\[
  \sup_\tau |I_\Delta(\tau)|<\infty,\qquad
  \sup_{\tau\ge\tau_0}\int |y|^{-p-1}|V(y,\tau)|^4\,dy<\infty,\qquad
  \sup_\tau |I_P(\tau)|<\infty,
\]
then
\[
  \sup_{\tau\ge\tau_0}
  \left|
  \int |y|^{-p}|V|V\cdot H(V,P)\,dy
  \right|<\infty.
\]
\end{corollary}

\begin{proof}
\Cref{paper7:lem:scale-nonlinear-from-l4} bounds \(I_{\rm nl}\).  Together
with the Laplacian and pressure bounds, this gives the hypotheses of
\cref{paper7:lem:scale-force-decomp}.
\end{proof}

\begin{lemma}[Bounded modulation from the modulation system]
\label{paper7:lem:bounded-mod-from-system}
Assume
\[
  M(V(\tau))
  \begin{pmatrix}
  a(\tau)\\ b(\tau)
  \end{pmatrix}
  =-F(V(\tau)),
\]
\[
  \sup_{\tau\ge\tau_0}\|M(V(\tau))^{-1}\|_{\infty\to\infty}<\infty,
  \qquad
  \sup_{\tau\ge\tau_0}\|F(V(\tau))\|_\infty<\infty.
\]
Then
\[
  \sup_{\tau\ge\tau_0}(|a(\tau)|+\|b(\tau)\|_\infty)<\infty.
\]
\end{lemma}

\begin{proof}
The system gives
\[
  \begin{pmatrix}
  a(\tau)\\ b(\tau)
  \end{pmatrix}
  =-M(V(\tau))^{-1}F(V(\tau)),
\]
so
\[
  \max\{|a(\tau)|,\|b(\tau)\|_\infty\}
  \le
  \|M(V(\tau))^{-1}\|_{\infty\to\infty}\|F(V(\tau))\|_\infty.
\]
Taking the supremum in \(\tau\) gives bounded modulation.
\end{proof}

\begin{proposition}[Windowed \(H^1\) exclusion from modulation bounds]\label{paper7:prop:windowed-h1-from-modulation}
Assume a renormalized Type II solution has the repaired-gauge orbit and pressure
reconstruction, bounded critical \(L^3\) norm, uniformly invertible modulation
matrix, uniformly bounded modulation forcing vector, and the compact-cylinder
regularity needed for Caccioppoli.  Then
\[
  \forall R>0:\qquad
  \sup_{T\ge\tau_0+1}
  \int_T^{T+1}\|V(\tau)\|_{H^1(B_R)}^2\,d\tau<\infty,
\]
and failure of local windowed \(H^1\) control is excluded for this solution.
\end{proposition}

\begin{proof}
The inverse matrix and forcing bounds imply bounded modulation by
\cref{paper7:lem:bounded-mod-from-system}.  The renormalized solution supplies the
renormalized equation and pressure reconstruction.  Bounded critical norm
gives the local compact-cylinder critical bounds, and the regularity hypothesis
allows the Caccioppoli estimate.  Therefore
\cref{paper7:cor:windowed-h1} applies.  Failure of local windowed \(H^1\)
control is the negation of the displayed tail windowed \(H^1\) estimate on
some bounded ball.
\end{proof}

\begin{corollary}[Two-alternative exclusion after windowed control]\label{paper7:cor:windowed-control-update}
Assume a finite-cost renormalized Type II solution not excluded by the compact
cost-divergence criterion has reached the two-alternative classification.  If
the hypotheses of \cref{paper7:prop:windowed-h1-from-modulation} are supplied for that
solution, then failure of local windowed \(H^1\) control is excluded.  If, in
addition, uniform critical tightness is available, then no finite-cost
alternative remains outside the compact cost-divergence criterion.
\end{corollary}

\begin{proof}
\Cref{paper7:prop:windowed-h1-from-modulation} gives the positive tail windowed
\(H^1\) bound, contradicting failure of local windowed \(H^1\) control.  If
uniform critical tightness is also available, then failure of tightness is
excluded as well,
and the two-alternative classification has no remaining finite-cost alternative
outside the compact cost-divergence criterion.
\end{proof}

\begin{theorem}[Exclusion of local windowed \(H^1\) failure]\label{paper7:thm:windowed-h1-failure-exclusion}
If every renormalized Type II solution in the class satisfies the hypotheses of
\cref{paper7:prop:windowed-h1-from-modulation}, then no such solution can satisfy
\[
  \exists m\ge1:
  \sup_{T\ge\tau_0+1}
  \int_T^{T+1}\|V(\tau)\|_{H^1(B_m)}^2\,d\tau=\infty.
\]
\end{theorem}

\begin{proof}
Apply \cref{paper7:prop:windowed-h1-from-modulation} to each renormalized solution.  The
displayed condition is the negation of the resulting windowed
\(H^1\)-bound for \(R=m\).
\end{proof}

\subsection{Window-Integrated Modulation Forcing and Weighted Scale Components}

This subsection records sufficient analytic estimates for the modulation
forcing.  The estimates are stated directly for the differentiated modulation
conditions and for the weighted scale component; no pointwise modulation bound
is used unless it is explicitly assumed.

Let the modulation forcing operator be
\[
  H(V,P):=\nu\Delta V-(V\cdot\nabla)V-\nabla P.
\]
For gauge functionals \(G_k\), \(0\le k\le3\), define
\[
  F_k(V):=DG_k(V)[H(V,P)].
\]
Thus the modulation system has the form
\[
  M(V(\tau))
  \begin{pmatrix}
  a(\tau)\\ b(\tau)
  \end{pmatrix}
  =
  -F(V(\tau)).
\]

\begin{lemma}[Forcing bound for the translation constraints]
\label{paper7:lem:translation-constraint-forcing}
Let
\[
  G_j(V)=\int_{\mathbb R^3}y_j|V(y)|^2\psi_R(y)\,dy,
  \qquad 1\le j\le3,
\]
where \(\psi_R\in C_c^\infty(B_{2R})\), and set
\[
  \Phi_j(y):=y_j\psi_R(y).
\]
For
\[
  F_j(V):=DG_j(V)[H(V,P)]
  =
  2\int_{\mathbb R^3}\Phi_j\,V\cdot H(V,P)\,dy,
\]
assume, uniformly for \(\tau\ge\tau_0\),
\[
  \|V(\tau)\|_{L^3(B_{2R})}\le M_3,\qquad
  \|\nabla V(\tau)\|_{L^2(B_{2R})}\le M_\nabla,
\]
and
\[
  \inf_{c\in\mathbb R}
  \|P(\tau)-c\|_{L^{3/2}(B_{2R})}\le M_P.
\]
Then
\[
  \sup_{\tau\ge\tau_0}\max_{1\le j\le3}|F_j(V(\tau))|
  \le
  C(R,\nu,\psi_R)
  \left(M_\nabla^2+M_3^2+M_3^3+M_PM_3\right).
\]
\end{lemma}

\begin{proof}
Fix \(j\).  Since \(\Phi_j\) is supported in \(B_{2R}\),
\[
  F_j=2\nu I_\Delta-2I_{\rm nl}-2I_P,
\]
where
\[
  I_\Delta:=\int \Phi_jV\cdot\Delta V,\qquad
  I_{\rm nl}:=\int \Phi_jV\cdot(V\cdot\nabla V),
  \qquad
  I_P:=\int \Phi_jV\cdot\nabla P.
\]
For the Laplacian term, using compact support of \(\Phi_j\) and summing over
the component index \(\alpha\),
\[
\begin{aligned}
  I_\Delta
  &=\sum_{\alpha=1}^3\int \Phi_j V_\alpha\Delta V_\alpha\,dy        \\
  &=-\sum_{\alpha,k}\int \partial_k(\Phi_jV_\alpha)\partial_kV_\alpha\,dy\\
  &=-\int \Phi_j|\nabla V|^2\,dy
    -\sum_{\alpha,k}\int(\partial_k\Phi_j)V_\alpha\partial_kV_\alpha\,dy .
\end{aligned}
\]
Equivalently,
\[
  I_\Delta
  =
  -\int \Phi_j|\nabla V|^2
  -\int(\nabla\Phi_j\cdot\nabla V)\cdot V .
\]
Consequently, by the Cauchy--Schwarz inequality,
\[
  |I_\Delta|
  \le
  \|\Phi_j\|_{L^\infty}\|\nabla V\|_{L^2(B_{2R})}^2
  +
  \|\nabla\Phi_j\|_{L^\infty}
  \|V\|_{L^2(B_{2R})}\|\nabla V\|_{L^2(B_{2R})}.
\]
H\"older's inequality on the bounded set \(B_{2R}\) gives
\[
  \|V\|_{L^2(B_{2R})}
  \le |B_{2R}|^{1/6}\|V\|_{L^3(B_{2R})}
  \le C_RM_3.
\]
Therefore
\[
  |I_\Delta|
  \le C_R(M_\nabla^2+M_3M_\nabla)
  \le C_R(M_\nabla^2+M_3^2).
\]
The last inequality is Young's inequality, \(ab\le(a^2+b^2)/2\).

For the nonlinear term,
\[
  V\cdot(V\cdot\nabla V)
  =
  V_iV_\alpha\partial_iV_\alpha
  =
  \frac12V_i\partial_i(|V|^2)
  =
  \frac12V\cdot\nabla(|V|^2).
\]
Since \(\nabla\cdot V=0\) and \(\Phi_j\) is compactly supported,
\[
  I_{\rm nl}
  =
  \frac12\int \Phi_jV\cdot\nabla(|V|^2)
  =
  -\frac12\int |V|^2V\cdot\nabla\Phi_j.
\]
Thus
\[
  |I_{\rm nl}|
  \le
  \frac12\|\nabla\Phi_j\|_{L^\infty}
  \|V\|_{L^3(B_{2R})}^3
  \le C_RM_3^3.
\]

For the pressure term, replace \(P\) by \(P-c(\tau)\).  Since
\(\nabla\cdot V=0\), the identity
\[
  \nabla\cdot(\Phi_j V)=V\cdot\nabla\Phi_j
\]
holds distributionally.  Therefore
\[
  I_P
  =
  \int \Phi_jV\cdot\nabla(P-c)
  =
  -\int(P-c)V\cdot\nabla\Phi_j.
\]
Consequently
\[
  |I_P|
  \le
  \|\nabla\Phi_j\|_{L^\infty}
  \|P-c\|_{L^{3/2}(B_{2R})}
  \|V\|_{L^3(B_{2R})}
  \le C_RM_PM_3.
\]
Taking the infimum over \(c\), summing the three estimates in \(F_j\), and
then taking the supremum in \(\tau\) proves the bound.
\end{proof}

\begin{remark}[Independence of the local gradient hypothesis]
\label{paper7:rem:translation-gradient-input}
The pointwise local gradient bound in
\cref{paper7:lem:translation-constraint-forcing} is an additional analytic
input.  When the translation estimate is used to prove local windowed
\(H^1\)-control, this input is required to come from an independent Caccioppoli,
regularity, or approximation estimate, not from the conclusion being proved.
\end{remark}

\begin{definition}[Weighted scale forcing]
\label{paper7:def:weighted-scale-forcing}
Let
\[
  G_{\rm sc}(V)=\int_{\mathbb R^3}|y|^{-p}|V|^3\,dy-\Theta_0,
  \qquad 0<p<3.
\]
The weighted scale component is
\[
  F_0(V):=DG_{\rm sc}(V)[H(V,P)]
  =
  3\int_{\mathbb R^3}|y|^{-p}|V|V\cdot H(V,P)\,dy.
\]
All weighted pairings in this subsection are assumed to be defined either as
absolutely convergent integrals or as specified distributional pairings
represented by \(L^1_{\rm loc}(d\tau)\) functions.
\end{definition}

\begin{lemma}[Weighted scale forcing bound]
\label{paper7:lem:scale-component-weighted-forcing}
Assume
\[
  \sup_{\tau\ge\tau_0}
  \left|
  \int_{\mathbb R^3}|y|^{-p}|V|V\cdot H(V,P)\,dy
  \right|<\infty.
\]
Then
\[
  \sup_{\tau\ge\tau_0}|F_0(V(\tau))|<\infty.
\]
\end{lemma}

\begin{proof}
By the derivative formula in \cref{paper7:def:weighted-scale-forcing},
\[
  |F_0(V(\tau))|
  \le
  3\left|
  \int_{\mathbb R^3}|y|^{-p}|V|V\cdot H(V,P)\,dy
  \right|.
\]
Taking the supremum in \(\tau\) gives the result.
\end{proof}

\begin{corollary}[Pointwise modulation forcing bound]
\label{paper7:cor:pointwise-modulation-forcing}
Assume the hypotheses of \cref{paper7:lem:translation-constraint-forcing} for
the three translation components and the weighted scale forcing hypothesis of
\cref{paper7:lem:scale-component-weighted-forcing}.  Then
\[
  \sup_{\tau\ge\tau_0}\|F(V(\tau))\|_\infty<\infty .
\]
\end{corollary}

\begin{proof}
The vector \(F(V)\in\mathbb R^4\) has the scale component \(F_0\) and the
three translation components \(F_j\), \(1\le j\le3\).  The scale component is
bounded by \cref{paper7:lem:scale-component-weighted-forcing}, and the
translation components are bounded by
\cref{paper7:lem:translation-constraint-forcing}.  Taking the maximum over the
four components proves the \(\ell^\infty\)-bound.
\end{proof}

\begin{corollary}[Analytic inputs for pointwise modulation forcing]
\label{paper7:cor:pointwise-modulation-inputs}
Pointwise boundedness of the modulation forcing follows from the following
analytic estimates:
\begin{enumerate}[label=\textup{(\roman*)}]
\item local \(L^3\), local \(H^1\), and local pressure-oscillation control on
the support of the compactly supported centering constraints;
\item weighted control of the Laplacian, pressure, and fourth-moment terms in
the scale component, sufficient to bound the weighted scale forcing.
\end{enumerate}
The scale-component estimate is independent of the algebraic scale
transversality
identity
\[
  DG_{\rm sc}(V)[Z_{\rm sc}(V)]=p\Theta_0.
\]
\end{corollary}

\begin{proof}
The first item is the set of hypotheses used in
\cref{paper7:lem:translation-constraint-forcing}.  The second item bounds the
weighted scale component through \cref{paper7:lem:scale-component-weighted-forcing},
or through the explicit decomposition in
\cref{paper7:cor:reduced-window-scale-force} below in the corresponding
integrated setting.  \Cref{paper7:cor:pointwise-modulation-forcing} then gives
the modulation forcing bound.  The displayed transversality identity controls
the modulation matrix, not the forcing vector, so it is a separate algebraic
input.
\end{proof}

\begin{lemma}[Window-integrated forcing gives window-integrated modulation]
\label{paper7:lem:window-l1-forcing-modulation}
Assume
\[
  \sup_{\tau\ge\tau_0}
  \|M(V(\tau))^{-1}\|_{\infty\to\infty}\le C_M<\infty
\]
and
\[
  \sup_{T\ge\tau_0}
  \int_T^{T+1}\|F(V(\tau))\|_\infty\,d\tau
  \le C_F<\infty .
\]
Then
\[
  \sup_{T\ge\tau_0}
  \int_T^{T+1}\bigl(|a(\tau)|+\|b(\tau)\|_\infty\bigr)\,d\tau
  \le 2C_MC_F.
\]
\end{lemma}

\begin{proof}
For almost every \(\tau\),
\[
  \begin{pmatrix}
  a(\tau)\\ b(\tau)
  \end{pmatrix}
  =
  -M(V(\tau))^{-1}F(V(\tau)).
\]
Hence
\[
  \max\{|a(\tau)|,\|b(\tau)\|_\infty\}
  \le C_M\|F(V(\tau))\|_\infty.
\]
Since
\[
  |a(\tau)|+\|b(\tau)\|_\infty
  \le
  2\max\{|a(\tau)|,\|b(\tau)\|_\infty\},
\]
we have, for every \(T\ge\tau_0\),
\[
\begin{aligned}
  \int_T^{T+1}\bigl(|a(\tau)|+\|b(\tau)\|_\infty\bigr)\,d\tau
  &\le
  2C_M\int_T^{T+1}\|F(V(\tau))\|_\infty\,d\tau  \\
  &\le 2C_MC_F .
\end{aligned}
\]
Taking the supremum over \(T\) proves the claim.
\end{proof}

\begin{lemma}[Selection from an \(L^1\) modulation bound]
\label{paper7:lem:unit-window-selection}
Let \(I_n=[T_n,T_n+1]\), where \(T_n\to\infty\).  Let
\[
  d_n:I_n\to[0,\infty),\qquad q_n:I_n\to[0,\infty)
\]
be measurable functions satisfying
\[
  \int_{I_n}d_n(\tau)\,d\tau\to0,
  \qquad
  \sup_n\int_{I_n}q_n(\tau)\,d\tau\le Q<\infty .
\]
Then, after choosing measurable representatives, there are times
\(\tau_n\in I_n\) such that
\[
  d_n(\tau_n)\to0
\]
and, if \(Q>0\),
\[
  q_n(\tau_n)\le2Q
  \qquad\text{for every }n.
\]
If \(Q=0\), the times may be chosen so that
\[
  q_n(\tau_n)=0,\qquad d_n(\tau_n)\to0.
\]
\end{lemma}

\begin{proof}
Assume first that \(Q>0\), and set
\[
  E_n:=\{\tau\in I_n:q_n(\tau)\le2Q\}.
\]
By Chebyshev's inequality,
\[
  |I_n\setminus E_n|
  \le
  \frac{1}{2Q}\int_{I_n}q_n(\tau)\,d\tau
  \le \frac12.
\]
Thus \(|E_n|\ge1/2\).  Since \(d_n\ge0\),
\[
  \inf_{\tau\in E_n}d_n(\tau)
  \le
  \frac{1}{|E_n|}\int_{E_n}d_n(\tau)\,d\tau
  \le
  2\int_{I_n}d_n(\tau)\,d\tau.
\]
Choose \(\tau_n\in E_n\), outside the null set on which the representatives
are undefined, such that
\[
  d_n(\tau_n)
  \le
  2\int_{I_n}d_n(\tau)\,d\tau+\frac1n.
\]
The right-hand side tends to zero, hence \(d_n(\tau_n)\to0\), and the
definition of \(E_n\) gives \(q_n(\tau_n)\le2Q\).

If \(Q=0\), then \(q_n=0\) almost everywhere on \(I_n\).  Applying the same
averaging argument to \(d_n\) on the full-measure set where \(q_n=0\) gives
the asserted times.
\end{proof}

\begin{corollary}[Selection on unit windows]
\label{paper7:cor:unit-window-selection}
Assume
\[
  q(\tau):=|a(\tau)|+\|b(\tau)\|_\infty
\]
satisfies
\[
  \sup_{T\ge\tau_0}\int_T^{T+1}q(\tau)\,d\tau<\infty .
\]
For every sequence of unit windows \(I_n=[T_n,T_n+1]\) and every nonnegative
measurable density \(d_n\) satisfying
\[
  \int_{I_n}d_n(\tau)\,d\tau\to0,
\]
there are times \(\tau_n\in I_n\) such that
\[
  d_n(\tau_n)\to0,\qquad
  \sup_n\bigl(|a(\tau_n)|+\|b(\tau_n)\|_\infty\bigr)<\infty .
\]
\end{corollary}

\begin{proof}
Apply \cref{paper7:lem:unit-window-selection} with \(q_n=q|_{I_n}\) and
\[
  Q=\sup_n\int_{I_n}q(\tau)\,d\tau .
\]
This \(Q\) is finite because of the uniform window \(L^1\)-bound.  If
\(Q>0\), the selected times satisfy
\[
  |a(\tau_n)|+\|b(\tau_n)\|_\infty=q(\tau_n)\le2Q.
\]
If \(Q=0\), the same lemma gives \(q(\tau_n)=0\).  In both cases
\(d_n(\tau_n)\to0\) and the modulation is uniformly bounded along the
selected times.
\end{proof}

\begin{lemma}[Product-form windowed modulation estimate]
\label{paper7:lem:product-form-window-modulation}
Assume
\[
  \sup_{T\ge\tau_0}
  \int_T^{T+1}
  \|M(V(\tau))^{-1}\|_{\infty\to\infty}
  \|F(V(\tau))\|_\infty\,d\tau<\infty.
\]
Then
\[
  \sup_{T\ge\tau_0}
  \int_T^{T+1}\bigl(|a(\tau)|+\|b(\tau)\|_\infty\bigr)\,d\tau<\infty .
\]
\end{lemma}

\begin{proof}
The modulation system gives
\[
  \begin{pmatrix}a(\tau)\\ b(\tau)\end{pmatrix}
  =
  -M(V(\tau))^{-1}F(V(\tau)).
\]
Therefore
\[
  |a(\tau)|+\|b(\tau)\|_\infty
  \le
  2\|M(V(\tau))^{-1}\|_{\infty\to\infty}
  \|F(V(\tau))\|_\infty.
\]
Integrating this pointwise estimate on \([T,T+1]\) and taking the supremum in
\(T\) gives the result.
\end{proof}

\begin{lemma}[Integrated translation forcing bound]
\label{paper7:lem:window-translation-forcing}
Let \(G_j\), \(\Phi_j\), \(H(V,P)\), and \(F_j(V)\) be as in
\cref{paper7:lem:translation-constraint-forcing}.  Assume that, for some
\(R>0\),
\[
  \sup_{\tau\ge\tau_0}\|V(\tau)\|_{L^3(B_{2R})}\le M_3,
\]
\[
  \sup_{T\ge\tau_0}
  \int_T^{T+1}\|\nabla V(\tau)\|_{L^2(B_{2R})}^2\,d\tau
  \le A_R,
\]
and
\[
  \sup_{\tau\ge\tau_0}\inf_{c\in\mathbb R}
  \|P(\tau)-c\|_{L^{3/2}(B_{2R})}\le M_P.
\]
Then
\[
  \sup_{T\ge\tau_0}
  \int_T^{T+1}\max_{1\le j\le3}|F_j(V(\tau))|\,d\tau
  \le
  C(R,\nu,\psi_R)(A_R+M_3^2+M_3^3+M_PM_3).
\]
\end{lemma}

\begin{proof}
Fix \(j\).  The proof of
\cref{paper7:lem:translation-constraint-forcing} gives, for almost every
\(\tau\),
\[
  |F_j(V(\tau))|
  \le
  C(R,\nu,\psi_R)
  \left(
  \|\nabla V(\tau)\|_{L^2(B_{2R})}^2
  +M_3\|\nabla V(\tau)\|_{L^2(B_{2R})}
  +M_3^3+M_PM_3
  \right).
\]
Integrate this estimate over \(I_T=[T,T+1]\).  The first term contributes at
most \(A_R\).  For the mixed term, Cauchy's inequality on the unit interval
gives
\[
  \int_{I_T}M_3\|\nabla V(\tau)\|_{L^2(B_{2R})}\,d\tau
  \le
  M_3
  \left(
  \int_{I_T}\|\nabla V(\tau)\|_{L^2(B_{2R})}^2\,d\tau
  \right)^{1/2}
  \le M_3A_R^{1/2}.
\]
Since
\[
  M_3A_R^{1/2}\le\frac12(M_3^2+A_R),
\]
the mixed term is bounded by \(C(A_R+M_3^2)\).  The terms \(M_3^3\) and
\(M_PM_3\) are constant on the unit window.  Hence
\[
  \int_T^{T+1}|F_j(V(\tau))|\,d\tau
  \le
  C(R,\nu,\psi_R)(A_R+M_3^2+M_3^3+M_PM_3).
\]
Since \(\max_{1\le j\le3}|F_j|\le\sum_{j=1}^3|F_j|\), taking the maximum
over the three centering constraints changes only the constant.
\end{proof}

\begin{remark}[Independence of the windowed gradient hypothesis]
\label{paper7:rem:window-translation-input}
The windowed gradient input in
\cref{paper7:lem:window-translation-forcing} is an independent
hypothesis whenever the integrated translation estimate is used to prove
windowed \(H^1\)-control.  It may come, for example, from a preliminary
Caccioppoli estimate depending only on already available data, from a direct
approximation or regularity estimate on the support of the centering
functions, or from an independent verification obtained without using the
conclusion of the windowed \(H^1\) estimate under consideration.
\end{remark}

\begin{lemma}[Window-integrated scale forcing]
\label{paper7:lem:window-scale-forcing}
Assume the window-integrated weighted scale-forcing bound
\[
  \sup_{T\ge\tau_0}
  \int_T^{T+1}
  \left|
  \int_{\mathbb R^3}|y|^{-p}|V|V\cdot H(V,P)\,dy
  \right|\,d\tau
  <\infty,
\]
where the weighted pairing is interpreted as in
\cref{paper7:def:weighted-scale-forcing}.  Then
\[
  \sup_{T\ge\tau_0}\int_T^{T+1}|F_0(V(\tau))|\,d\tau<\infty .
\]
\end{lemma}

\begin{proof}
By the identity
\[
  F_0(V)
  =
  3\int_{\mathbb R^3}|y|^{-p}|V|V\cdot H(V,P)\,dy
\]
we have, for each \(T\ge\tau_0\),
\[
  \int_T^{T+1}|F_0(V(\tau))|\,d\tau
  =
  3\int_T^{T+1}
  \left|
  \int_{\mathbb R^3}|y|^{-p}|V|V\cdot H(V,P)\,dy
  \right|\,d\tau .
\]
Taking the supremum over \(T\) gives the asserted \(L^1\)-bound.
\end{proof}

\begin{lemma}[Full integrated modulation forcing]
\label{paper7:lem:window-modulation-forcing}
Assume the integrated translation forcing estimate
\[
  \sup_{T\ge\tau_0}
  \int_T^{T+1}\max_{1\le j\le3}|F_j(V(\tau))|\,d\tau<\infty
\]
and the window-integrated weighted scale-forcing bound in
\cref{paper7:lem:window-scale-forcing}.  Then
\[
  \sup_{T\ge\tau_0}
  \int_T^{T+1}\|F(V(\tau))\|_\infty\,d\tau<\infty .
\]
\end{lemma}

\begin{proof}
The vector \(F(V)\in\mathbb R^4\) consists of \(F_0\) and the three
translation components \(F_j\), \(1\le j\le3\).  Since
\[
  \|F(V(\tau))\|_\infty
  \le
  |F_0(V(\tau))|+\max_{1\le j\le3}|F_j(V(\tau))|,
\]
for almost every \(\tau\).  Integrating this inequality on \([T,T+1]\) gives
\[
  \int_T^{T+1}\|F(V(\tau))\|_\infty\,d\tau
  \le
  \int_T^{T+1}|F_0(V(\tau))|\,d\tau
  +
  \int_T^{T+1}\max_{1\le j\le3}|F_j(V(\tau))|\,d\tau .
\]
The first term is uniformly bounded in \(T\) by
\cref{paper7:lem:window-scale-forcing}, and the second term is uniformly
bounded by the assumed integrated translation forcing estimate.  Taking the
supremum over \(T\) proves the claim.
\end{proof}

\begin{corollary}[Integrated modulation criterion on selected windows]
\label{paper7:cor:selected-window-modulation}
Assume the uniform inverse bound for \(M(V(\tau))\), the integrated
translation forcing estimate, and the window-integrated weighted scale-forcing
bound.  Then
\[
  \sup_{T\ge\tau_0}
  \int_T^{T+1}\bigl(|a(\tau)|+\|b(\tau)\|_\infty\bigr)\,d\tau<\infty.
\]
Consequently every sequence of unit windows with vanishing nonnegative
measurable density contains selected times where this density vanishes
asymptotically and the modulation remains uniformly bounded.
\end{corollary}

\begin{proof}
\Cref{paper7:lem:window-modulation-forcing} gives the
window-integrated bound on \(\|F(V)\|_\infty\).  Combining this with the
uniform inverse bound in
\cref{paper7:lem:window-l1-forcing-modulation} gives the
window-integrated modulation bound.  The selection conclusion is
\cref{paper7:cor:unit-window-selection}.
\end{proof}

\begin{definition}[Weighted scale-force components]
\label{paper7:def:scale-force-components}
Let \(w(y)=|y|^{-p}\), \(0<p<3\).  Define
\[
  F_0^\Delta(V)
  :=
  3\nu\int_{\mathbb R^3}w|V|V\cdot\Delta V\,dy,
\]
\[
  F_0^{\rm nl}(V)
  :=
  -3\int_{\mathbb R^3}w|V|V\cdot(V\cdot\nabla V)\,dy,
\]
and
\[
  F_0^P(V,P)
  :=
  -3\int_{\mathbb R^3}w|V|V\cdot\nabla P\,dy.
\]
The three pairings are required to be compatible with the decomposition of
\(H(V,P)\) into diffusion, transport, and pressure.
\end{definition}

\begin{lemma}[Integrated scale-force decomposition]
\label{paper7:lem:scale-force-decomposition-window}
For every time at which the three pairings in
\cref{paper7:def:scale-force-components} are defined,
\[
  F_0(V(\tau))
  =
  F_0^\Delta(V(\tau))
  +F_0^{\rm nl}(V(\tau))
  +F_0^P(V(\tau),P(\tau)).
\]
\end{lemma}

\begin{proof}
Substitute
\[
  H(V,P)=\nu\Delta V-(V\cdot\nabla)V-\nabla P
\]
into
\[
  F_0(V)=3\int w|V|V\cdot H(V,P)\,dy.
\]
Then
\[
\begin{aligned}
  F_0(V)
  &=
  3\nu\int w|V|V\cdot\Delta V\,dy
  -3\int w|V|V\cdot(V\cdot\nabla V)\,dy  \\
  &\qquad
  -3\int w|V|V\cdot\nabla P\,dy .
\end{aligned}
\]
The three terms coincide with \(F_0^\Delta(V)\), \(F_0^{\rm nl}(V)\), and
\(F_0^P(V,P)\), respectively.  Thus linearity of the weighted pairing gives
the stated decomposition.
\end{proof}

\begin{lemma}[Component bounds imply integrated scale-force control]
\label{paper7:lem:scale-components-control}
Assume
\[
  \sup_{T\ge\tau_0}
  \int_T^{T+1}|F_0^\Delta(V(\tau))|\,d\tau<\infty,
\]
\[
  \sup_{T\ge\tau_0}
  \int_T^{T+1}|F_0^{\rm nl}(V(\tau))|\,d\tau<\infty,
\]
and
\[
  \sup_{T\ge\tau_0}
  \int_T^{T+1}|F_0^P(V(\tau),P(\tau))|\,d\tau<\infty.
\]
Then the window-integrated weighted scale-forcing bound in
\cref{paper7:lem:window-scale-forcing} holds.
\end{lemma}

\begin{proof}
By \cref{paper7:lem:scale-force-decomposition-window},
\[
  |F_0(V(\tau))|
  \le
  |F_0^\Delta(V(\tau))|
  +|F_0^{\rm nl}(V(\tau))|
  +|F_0^P(V(\tau),P(\tau))|
\]
for almost every \(\tau\).  Integrating on \([T,T+1]\), taking the supremum
in \(T\), and applying the three component bounds gives
\[
  \sup_{T\ge\tau_0}\int_T^{T+1}|F_0(V(\tau))|\,d\tau<\infty.
\]
Since \(F_0\) is three times the weighted pairing in
\cref{paper7:def:weighted-scale-forcing}, this is equivalent to the
window-integrated weighted scale-forcing bound.
\end{proof}

\begin{definition}[Singular-weight convective integration by parts]
\label{paper7:def:singular-weight-convective-ibp}
The singular-weight convective integration-by-parts condition means that, for
almost every \(\tau\),
\[
  -\int_{\mathbb R^3}w\,V\cdot\nabla(|V|^3)\,dy
  =
  \int_{\mathbb R^3}|V|^3V\cdot\nabla w\,dy,
\]
with no boundary contribution from \(y=0\) or from spatial infinity, either
classically or by convergence of smooth compactly supported divergence-free
approximants.
\end{definition}

\begin{lemma}[Convective scale-force bound]
\label{paper7:lem:convective-scale-force}
Assume the condition of \cref{paper7:def:singular-weight-convective-ibp} and
\[
  \sup_{T\ge\tau_0}
  \int_T^{T+1}\int_{\mathbb R^3}
  |y|^{-p-1}|V(y,\tau)|^4\,dy\,d\tau<\infty .
\]
Then
\[
  \sup_{T\ge\tau_0}
  \int_T^{T+1}|F_0^{\rm nl}(V(\tau))|\,d\tau<\infty.
\]
More precisely,
\[
  |F_0^{\rm nl}(V(\tau))|
  \le
  p\int_{\mathbb R^3}|y|^{-p-1}|V(y,\tau)|^4\,dy
\]
for almost every \(\tau\).
\end{lemma}

\begin{proof}
Fix such a time \(\tau\) and suppress it from the notation.  Since
\[
  \partial_i(|V|^3)=3|V|V_j\partial_iV_j
\]
for each \(i\), and therefore
\[
  V\cdot\nabla(|V|^3)
  =
  3|V|V_iV_j\partial_iV_j
  =
  3|V|\,V\cdot((V\cdot\nabla)V),
\]
we have
\[
  w|V|V\cdot(V\cdot\nabla V)
  =
  \frac13 w\,V\cdot\nabla(|V|^3).
\]
Therefore
\[
  F_0^{\rm nl}(V)
  =
  -3\int w|V|V\cdot(V\cdot\nabla V)
  =
  -\int w\,V\cdot\nabla(|V|^3).
\]
Using \(\nabla\cdot V=0\) and
\cref{paper7:def:singular-weight-convective-ibp},
\[
  -\int w\,V\cdot\nabla(|V|^3)
  =
  \int |V|^3V\cdot\nabla w.
\]
Since
\[
  \nabla w(y)=-p|y|^{-p-2}y,\qquad y\ne0,
\]
we have
\[
  |V\cdot\nabla w|
  =
  p|y|^{-p-2}|V\cdot y|
  \le p|y|^{-p-1}|V|.
\]
Thus
\[
  |F_0^{\rm nl}(V)|
  \le
  p\int |y|^{-p-1}|V|^4\,dy.
\]
Integrating over \([T,T+1]\) and taking the supremum in \(T\) proves the
window-integrated convective bound.
\end{proof}

\begin{definition}[Annular convective regularity]
\label{paper7:def:annular-convective-regularity}
Annular convective regularity means that, for almost every \(\tau\),
\(V(\tau)\) is divergence free and, on every compact annulus
\(0<r<|y|<R<\infty\), the chain-rule identities
\[
  \partial_i(|V|^3)=3|V|V_j\partial_iV_j,
  \qquad
  V\cdot\nabla(|V|^3)=3|V|V_iV_j\partial_iV_j
\]
hold in the sense of distributions, either classically or by passage to the
limit from smooth divergence-free approximants on annuli.
\end{definition}

\begin{lemma}[Annular regularity gives singular-weight integration by parts]
\label{paper7:lem:annular-regularity-ibp}
Assume annular convective regularity and the weighted \(L^4\)-window bound
\[
  \sup_{T\ge\tau_0}
  \int_T^{T+1}\int_{\mathbb R^3}
  |y|^{-p-1}|V(y,\tau)|^4\,dy\,d\tau<\infty.
\]
Then the singular-weight convective integration-by-parts condition of
\cref{paper7:def:singular-weight-convective-ibp} holds for almost every \(\tau\).
\end{lemma}

\begin{proof}
Fix a time \(\tau\) for which
\[
  \int_{\mathbb R^3}|y|^{-p-1}|V(y,\tau)|^4\,dy<\infty
\]
and annular convective regularity holds.  Such times form a full-measure set
on every unit window.

Let \(\eta\in C^\infty([0,\infty))\) satisfy \(0\le\eta\le1\),
\(\eta(s)=0\) for \(s\le1\), and \(\eta(s)=1\) for \(s\ge2\).  Define
\[
  \chi_{\varepsilon,R}(y)
  :=
  \eta(|y|/\varepsilon)\eta(R/|y|).
\]
Then \(\chi_{\varepsilon,R}\) is supported in
\(\{\varepsilon<|y|<R\}\), equals one on
\(\{2\varepsilon<|y|<R/2\}\), and satisfies
\[
  |\nabla\chi_{\varepsilon,R}(y)|
  \le
  C\varepsilon^{-1}\mathbf 1_{\{\varepsilon<|y|<2\varepsilon\}}
  +CR^{-1}\mathbf 1_{\{R/2<|y|<R\}}.
\]
By annular convective regularity, the following integration by parts is valid
on the annulus \(\{\varepsilon<|y|<R\}\).  Since
\(\nabla\cdot V=0\) and \(w\chi_{\varepsilon,R}|V|^3V\) is compactly
supported in this annulus,
\[
  0
  =
  \int \nabla\cdot\bigl(w\chi_{\varepsilon,R}|V|^3V\bigr)\,dy .
\]
Expanding the divergence gives
\[
  0
  =
  \int |V|^3V\cdot\nabla(w\chi_{\varepsilon,R})\,dy
  +
  \int w\chi_{\varepsilon,R}V\cdot\nabla(|V|^3)\,dy ,
\]
and hence
\[
  -\int w\chi_{\varepsilon,R}V\cdot\nabla(|V|^3)
  =
  \int |V|^3V\cdot\nabla(w\chi_{\varepsilon,R}).
\]
Expanding the derivative,
\[
  -\int w\chi_{\varepsilon,R}V\cdot\nabla(|V|^3)
  =
  \int \chi_{\varepsilon,R}|V|^3V\cdot\nabla w
  +\int w|V|^3V\cdot\nabla\chi_{\varepsilon,R}.
\]
On the inner cutoff region \(\{\varepsilon<|y|<2\varepsilon\}\), one has
\(|y|\simeq\varepsilon\), and on the outer cutoff region
\(\{R/2<|y|<R\}\), one has \(|y|\simeq R\).  Therefore the derivative bound
for \(\chi_{\varepsilon,R}\) implies
\[
\begin{aligned}
  \left|
  \int w|V|^3V\cdot\nabla\chi_{\varepsilon,R}
  \right|
  &\le
  C\int_{\{\varepsilon<|y|<2\varepsilon\}}
  |y|^{-p-1}|V|^4\,dy \\
  &\quad{}
  +C\int_{\{R/2<|y|<R\}}
  |y|^{-p-1}|V|^4\,dy .
\end{aligned}
\]
The integrand \(|y|^{-p-1}|V|^4\) belongs to \(L^1(\mathbb R^3)\) at the
fixed time.  Hence absolute continuity of the Lebesgue integral gives
\[
  \int_{\{\varepsilon<|y|<2\varepsilon\}}
  |y|^{-p-1}|V|^4\,dy\to0
  \qquad\text{as }\varepsilon\downarrow0,
\]
and integrability at infinity gives
\[
  \int_{\{R/2<|y|<R\}}
  |y|^{-p-1}|V|^4\,dy\to0
  \qquad\text{as }R\uparrow\infty.
\]
Moreover
\[
  |\chi_{\varepsilon,R}|V|^3V\cdot\nabla w|
  \le
  p|y|^{-p-1}|V|^4\in L^1(\mathbb R^3),
\]
so dominated convergence gives
\[
  \int \chi_{\varepsilon,R}|V|^3V\cdot\nabla w
  \to
  \int |V|^3V\cdot\nabla w.
\]
Passing first \(\varepsilon\downarrow0\) and then \(R\uparrow\infty\) yields
\[
  -\int_{\mathbb R^3}w\,V\cdot\nabla(|V|^3)\,dy
  =
  \int_{\mathbb R^3}|V|^3V\cdot\nabla w\,dy,
\]
with no boundary contribution at \(y=0\) or at infinity.
\end{proof}

\begin{corollary}[Reduced integrated scale-force criterion]
\label{paper7:cor:reduced-window-scale-force}
Assume annular convective regularity, the weighted \(L^4\)-window bound,
and the two component estimates
\[
  \sup_{T\ge\tau_0}
  \int_T^{T+1}|F_0^\Delta(V(\tau))|\,d\tau<\infty,
\]
\[
  \sup_{T\ge\tau_0}
  \int_T^{T+1}|F_0^P(V(\tau),P(\tau))|\,d\tau<\infty.
\]
Then the window-integrated weighted scale-forcing bound holds.
\end{corollary}

\begin{proof}
\Cref{paper7:lem:annular-regularity-ibp} gives the singular-weight
integration-by-parts condition.  \Cref{paper7:lem:convective-scale-force}
then gives the integrated convective scale-force bound.  Combining this with
the diffusive and pressure component estimates, and applying
\cref{paper7:lem:scale-components-control}, gives the window-integrated
weighted scale-forcing bound.
\end{proof}

\begin{corollary}[Remaining analytic alternatives for the scale component]
\label{paper7:cor:scale-component-alternatives}
Before the singular-weight integration by parts is justified, failure of the
window-integrated weighted scale-forcing bound is reduced to failure of at
least one of the following four analytic statements:
\begin{enumerate}[label=\textup{(\roman*)}]
\item the integrated diffusive weighted scale-force bound;
\item the singular-weight convective integration-by-parts identity;
\item the weighted \(L^4\)-window bound;
\item the integrated weighted pressure scale-force bound.
\end{enumerate}
If annular convective regularity is available, then the second item is replaced
by failure of annular convective regularity itself, together with the same
weighted \(L^4\)-window input.
\end{corollary}

\begin{proof}
\Cref{paper7:lem:scale-components-control} shows that the diffusive,
convective, and pressure component bounds imply the integrated scale-force
bound.  \Cref{paper7:lem:convective-scale-force} reduces the
convective component to the singular-weight integration-by-parts identity and
the weighted \(L^4\)-window bound.  Finally,
\cref{paper7:lem:annular-regularity-ibp} shows that annular convective
regularity and the weighted \(L^4\)-window bound imply the singular-weight
integration-by-parts identity.  Therefore the listed failures exhaust the
remaining ways in which this scale-component estimate can fail.
\end{proof}

\subsection{Pointwise Type II Exclusion}

\begin{definition}[Pointwise exclusion estimates]\label{paper7:def:pointwise-exclusion-tests}
For a renormalized Type II solution \(\omega\), the non-tight alternative is
excluded when \(V_\omega\) satisfies \cref{paper7:def:critical-tightness}.  The
windowed-control alternative is excluded when
\[
  \forall R>0:\qquad
  \sup_{T\ge\tau_0+1}
  \int_T^{T+1}\|V_\omega(\tau)\|_{H^1(B_R)}^2\,d\tau<\infty .
\]
The universal versions require the corresponding pointwise statement for
every renormalized solution in the class.
\end{definition}

\begin{definition}[Remaining Type II case after local estimates]
\label{paper7:def:remaining-typeII-branch}
A renormalized Type II solution remains after the local estimates if it reaches the
local two-alternative classification, is not already excluded by the compact
cost-divergence criterion, and is not excluded by the local Navier--Stokes
cost-to-exclusion theorem.
\end{definition}

\begin{theorem}[Type II exclusion after tightness and windowed control]
\label{paper7:thm:typeII-branch-exclusion}
Assume the renormalized Type II class satisfies:
\begin{enumerate}[label=\textup{(\roman*)}]
\item the local classification has exactly the following alternatives: the
compact cost-divergence channel, to be further stratified by carrier
validation, failure of uniform critical tightness, and failure of local
windowed \(H^1\) control;
\item every renormalized solution is uniformly critically tight;
\item every renormalized solution satisfies the tail windowed \(H^1\) bound;
\item the compact cost-divergence channel is evaluated as the retained compact
cost-divergence stratum after the carrier-validation alternatives of
\cref{paper7:thm:carrier-validation-retained-stratum} have been tested;
\item every carrier-validation adjacent stratum arising from this channel is
treated by its corresponding local stratum argument.
\end{enumerate}
Then every renormalized Type II solution is either excluded in a
carrier-validation adjacent stratum or excluded on the retained compact
cost-divergence stratum by the local Navier--Stokes cost-to-exclusion theorem.
Consequently no remaining Type II case exists.
\end{theorem}

\begin{proof}
Fix a renormalized Type II solution \(\omega\).  By the local
classification, \(\omega\) reaches exactly one of the following alternatives:
\begin{enumerate}[label=\textup{(\arabic*)}]
\item the compact cost-divergence channel, with its carrier-validation
subalternatives;
\item failure of uniform critical tightness;
\item failure of windowed \(H^1\) control, after uniform critical tightness
has been established.
\end{enumerate}
By the universal tightness hypothesis and
\cref{paper7:def:pointwise-exclusion-tests}, the second alternative is unavailable.  By
the universal windowed \(H^1\) hypothesis and
\cref{paper7:def:pointwise-exclusion-tests}, the third alternative is unavailable.  The
only remaining alternative is the compact cost-divergence channel.  By
hypothesis this channel has already been passed through the ordered
carrier-validation alternatives of
\cref{paper7:thm:carrier-validation-retained-stratum}.  If a carrier-validation
test fails, \(\omega\) leaves the retained compact stratum and belongs to the
corresponding adjacent stratum, where it is handled by the assumed local
stratum argument.  On the retained compact stratum, the same theorem gives
local monotonicity-carrier data for the validated compact cost.  These data give
the local cost-to-exclusion data by
\cref{paper7:lem:monotonicity-carrier-gives-cost-to-exclusion-data}; then
\cref{paper7:thm:local-cost-to-exclusion} excludes this case.  Since
\(\omega\) was arbitrary, no remaining Type II case exists.
\end{proof}

\begin{corollary}[Finite-cost exclusion after pointwise tests]
\label{paper7:cor:finite-cost-exclusion-after-tests}
Under the hypotheses of \cref{paper7:thm:typeII-branch-exclusion}, there is no
finite-cost renormalized Type II solution outside the compact cost-divergence
criterion.
\end{corollary}

\begin{proof}
\Cref{paper7:thm:typeII-branch-exclusion} shows that every renormalized Type II
solution is excluded.  Hence no renormalized solution remains outside the
compact cost-divergence criterion, regardless of whether its Type II cost is
finite.
\end{proof}

\subsection{Physical energy versus renormalized Type II cost}

\begin{definition}[Navier--Stokes cost-to-exclusion data]\label{paper7:def:ns-cost-to-exclusion-data}
A renormalized Navier--Stokes Type II branch satisfying the compact
cost-divergence criterion has Navier--Stokes cost-to-exclusion data if the
validated compact cost admits local monotonicity-carrier data in the sense of
\cref{paper7:def:local-monotonicity-carrier}.  In particular, this may be
verified either by the fixed-cutoff package of
\cref{paper7:lem:fixed-monotonicity-gives-cost-to-exclusion-data} or by the
moving-cutoff package of
\cref{paper7:lem:moving-monotonicity-gives-cost-to-exclusion-data}.
This is a data condition, not a primitive cost-exclusion assumption.  On the
retained compact cost-divergence stratum it is supplied by
\cref{paper7:thm:carrier-validation-retained-stratum}.  Failure of this data
condition is therefore not an unclassified obstruction: it is the carrier
validation exit recorded in \cref{paper7:def:cost-exclusion-missing-hypotheses}
and is treated as an adjacent state-space stratum.
\end{definition}

\begin{lemma}[Navier--Stokes cost-divergence exclusion from carrier data]
\label{paper7:lem:ns-cost-exclusion-from-carrier-data}
Assume a renormalized Type II solution has the active supercritical scaling
condition, satisfies the compact cost-divergence criterion, and has the
Navier--Stokes cost-to-exclusion data of
\cref{paper7:def:ns-cost-to-exclusion-data}.
Then the solution is excluded from the Type II class under consideration.
\end{lemma}

\begin{proof}
\Cref{paper7:def:ns-cost-to-exclusion-data} supplies local monotonicity-carrier data
for the validated compact cost.  By
\cref{paper7:lem:monotonicity-carrier-gives-cost-to-exclusion-data}, this
gives the local cost-to-exclusion data.  Therefore
\cref{paper7:thm:local-cost-to-exclusion} applies to the active supercritical
scaling condition and compact cost-divergence criterion, giving exclusion from
the stated Type II class.
\end{proof}

\begin{lemma}[Physical dissipation in renormalized variables]
\label{paper7:lem:physical-dissipation-renormalized}
Let
\[
  u(x,t)=
  \lambda(t)^{-1}
  V\left(\frac{x-x_c(t)}{\lambda(t)},\tau(t)\right),
  \qquad
  \frac{d\tau}{dt}=\lambda(t)^{-2}.
\]
Then
\[
  \|\nabla_xu(t)\|_{L^2_x}^2
  =
  \lambda(t)^{-1}\|\nabla_yV(\tau(t))\|_{L^2_y}^2.
\]
Consequently,
\[
  \int_{t_1}^{t_2}\|\nabla_xu(t)\|_{L^2_x}^2\,dt
  =
  \int_{\tau(t_1)}^{\tau(t_2)}
  \lambda(\tau)\|\nabla_yV(\tau)\|_{L^2_y}^2\,d\tau.
\]
\end{lemma}

\begin{proof}
Since \(u=\lambda^{-1}V\) and \(y=(x-x_c)/\lambda\),
\(\nabla_xu=\lambda^{-2}\nabla_yV\).  With \(dx=\lambda^3dy\),
\[
  \|\nabla_xu\|_{L^2_x}^2
  =
  \int \lambda^{-4}|\nabla_yV|^2\lambda^3\,dy
  =
  \lambda^{-1}\|\nabla_yV\|_{L^2_y}^2.
\]
Since \(d\tau/dt=\lambda^{-2}\), \(dt=\lambda^2d\tau\).  Multiplying gives
the second identity.
\end{proof}

\begin{definition}[Physical-renormalized dissipation cost]
\label{paper7:def:phys-ren-cost}
Define
\[
  \mathfrak D_{\rm phys-ren}(\tau)
  :=
  \nu\lambda(\tau)\|\nabla_yV(\tau)\|_{L^2_y}^2
\]
and the localized version
\[
  \mathfrak D_{{\rm phys-ren},R_0}(\tau)
  :=
  \nu\lambda(\tau)\int|\nabla_yV|^2\phi_{R_0}.
\]
\end{definition}

\begin{theorem}[Physical energy controls scale-weighted dissipation]
\label{paper7:thm:energy-controls-scale-weighted-diss}
Assume finite physical energy and the Navier--Stokes energy inequality
\[
  E(t)+\nu\int_0^t\|\nabla_xu(s)\|_{L^2_x}^2\,ds\le E(0),
  \qquad
  E(t)=\frac12\|u(t)\|_{L^2_x}^2.
\]
Then, for every \(t_1<T^*\),
\[
  \int_{\tau(t_1)}^\infty
  \mathfrak D_{\rm phys-ren}(\tau)\,d\tau
  \le E(t_1)\le E(0).
\]
In particular,
\[
  \int_{\tau(t_1)}^\infty
  \mathfrak D_{{\rm phys-ren},R_0}(\tau)\,d\tau<\infty.
\]
\end{theorem}

\begin{proof}
The physical energy inequality on \([t_1,t]\) gives
\[
  \nu\int_{t_1}^{t}\|\nabla_xu(s)\|_2^2\,ds\le E(t_1).
\]
Use \cref{paper7:lem:physical-dissipation-renormalized} and let
\(t\uparrow T^*\).  Monotone convergence gives the full weighted bound.  The
localized estimate follows from \(0\le\phi_{R_0}\le1\).
\end{proof}

\begin{definition}[Weighted positive scale-drift cost control]
\label{paper7:def:weighted-aplus}
Let
\[
  a_+(\tau)=\max\{a(\tau),0\}.
\]
Weighted positive scale-drift cost control means
\[
  \int_{\tau_1}^{\infty}
  \lambda(\tau)a_+(\tau)
  \int |V(y,\tau)|^2\phi_{R_0}(y)\,dy\,d\tau<\infty.
\]
\end{definition}

\begin{definition}[Scale-weighted Type II cost]\label{paper7:def:weighted-typeII-cost}
For
\[
  \widetilde{\mathfrak D}_{R_0}(\tau)
  =
  \nu\int|\nabla_yV|^2\phi_{R_0}
  +a_+(\tau)\int |V|^2\phi_{R_0},
\]
define
\[
  \mathfrak C_{\rm II,phys}^{NS}(\tau)
  :=
  \lambda(\tau)\widetilde{\mathfrak D}_{R_0}(\tau).
\]
\end{definition}

\begin{theorem}[Physical energy controls the scale-weighted Type II cost]
\label{paper7:thm:energy-controls-weighted-typeII}
Assume finite physical energy, the Navier--Stokes energy inequality, and
\cref{paper7:def:weighted-aplus}.  Then
\[
  \int_{\tau_1}^{\infty}
  \mathfrak C_{\rm II,phys}^{NS}(\tau)\,d\tau<\infty.
\]
\end{theorem}

\begin{proof}
By definition,
\[
  \mathfrak C_{\rm II,phys}^{NS}
  =
  \lambda\nu\int|\nabla V|^2\phi_{R_0}
  +\lambda a_+\int |V|^2\phi_{R_0}.
\]
The first term is integrable by
\cref{paper7:thm:energy-controls-scale-weighted-diss}; the second is
integrable by \cref{paper7:def:weighted-aplus}.
\end{proof}

\begin{definition}[Scale-decoupled infinite Type II cost]
\label{paper7:def:scale-decoupled-cost}
A renormalized Type II solution has scale-decoupled infinite Type II cost if
\[
  \int_{\tau_0}^{\infty}\widetilde{\mathfrak D}_{R_0}(\tau)\,d\tau=\infty,
  \qquad
  \int_{\tau_0}^{\infty}\lambda(\tau)\widetilde{\mathfrak D}_{R_0}(\tau)\,d\tau<\infty.
\]
The physically forbidden weighted-infinite-cost alternative is
\[
  \int_{\tau_0}^{\infty}
  \mathfrak C_{\rm II,phys}^{NS}(\tau)\,d\tau=\infty,
\]
which is impossible under \cref{paper7:thm:energy-controls-weighted-typeII}.
\end{definition}

\begin{theorem}[Infinite-cost trichotomy]\label{paper7:thm:c9-trichotomy}
Assume finite physical energy, the Navier--Stokes energy inequality, and
\cref{paper7:def:weighted-aplus}.  Let a renormalized Type II solution satisfy
\[
  \int_{\tau_0}^{\infty}\widetilde{\mathfrak D}_{R_0}(\tau)\,d\tau=\infty.
\]
Then:
\begin{enumerate}[label=\textup{(\roman*)}]
\item the compact Type II cost criterion gives a scale-collapse cost-divergence conclusion;
\item if the solution remains in the retained compact cost-divergence stratum
after carrier validation, then it is excluded from the Type II class under
consideration;
\item relative to physical energy, exactly one of the following occurs: the
physically forbidden weighted-infinite-cost alternative or the scale-decoupled
infinite-cost alternative of \cref{paper7:def:scale-decoupled-cost}.
\end{enumerate}
\end{theorem}

\begin{proof}
The unweighted infinite-cost conclusion is precisely the compact Type II
cost divergence, so the compact criterion gives a scale-collapse cost-divergence
conclusion.
If the solution remains in the retained compact cost-divergence stratum after
carrier validation, \cref{paper7:thm:carrier-validation-retained-stratum}
gives local monotonicity-carrier data, and
\cref{paper7:cor:retained-stratum-cost-exclusion} gives exclusion from the
Type II class under consideration.
\par
For the physical-energy classification, evaluate
\(\int\mathfrak C_{\rm II,phys}^{NS}\).  If it is infinite, then the
physically forbidden weighted-infinite-cost alternative occurs, contradicting
\cref{paper7:thm:energy-controls-weighted-typeII}.  If it is finite, then the
solution is scale-decoupled by \cref{paper7:def:scale-decoupled-cost}.  The
two physical-energy outcomes are disjoint and exhaustive.
\end{proof}

\begin{remark}
Finite physical energy controls the scale-weighted renormalized cost, not the
unweighted compact Type II cost.  Thus infinite scale-weighted renormalized
cost forces infinite physical energy dissipation, while infinite unweighted
renormalized cost does not by itself force infinite physical energy without an
additional scale-weight estimate.  Likewise, a compact Type II
cost-divergence conclusion does not automatically imply exclusion from the Type
II class unless the local corrected-monotonicity data of
\cref{paper7:thm:carrier-validation-retained-stratum} are present, equivalently
unless the branch remains in the retained compact stratum after the adjacent
carrier-failure alternatives have been evaluated.
\end{remark}

\subsection{Navier--Stokes Cost-Divergence Exclusion}

\begin{definition}[Type II exclusion conclusion]\label{paper7:def:cost-exclusion-conclusion}
For a renormalized Type II solution, the exclusion conclusion means that after
the active supercritical scaling condition and the compact cost-divergence
criterion have both been established, the solution is excluded from the Type II
class under consideration.  This is stronger than cost divergence alone, which
only states that the solution satisfies the compact Type II cost-divergence
criterion.
\end{definition}

\begin{definition}[Local cost-to-exclusion data]
\label{paper7:def:local-cost-to-exclusion-data}
Let \(\omega\) be a repaired-gauge Navier--Stokes Type II branch on a terminal
renormalized tail \([\tau_0,\infty)\).  Local cost-to-exclusion data for the
validated compact cost consist of a nonnegative locally integrable cost
\(\mathcal D_\omega\), a locally absolutely continuous corrected localized
energy \(\mathcal H_\omega\), a constant \(c_\omega>0\), and a nonnegative
locally integrable remainder \(\mathcal R_\omega\) such that:
\begin{enumerate}[label=\textup{(\roman*)}]
\item the validated compact cost-divergence conclusion is
\[
  \int_{\tau_0}^{\infty}\mathcal D_\omega(\tau)\,d\tau=\infty;
\]
\item \(\mathcal H_\omega(\tau_0)<\infty\) and
\[
  \inf_{\tau\ge\tau_0}\mathcal H_\omega(\tau)>-\infty;
\]
\item the corrected localized energy inequality
\[
  \frac{d}{d\tau}\mathcal H_\omega(\tau)
  +c_\omega\mathcal D_\omega(\tau)
  +\mathcal R_\omega(\tau)\le0
\]
holds in distributions on \((\tau_0,\infty)\).
\end{enumerate}
For the identity cost one takes
\(\mathcal D_\omega=\tilde{\mathfrak D}_{R_0}\).  For an explicitly comparable
localized cost, \(\mathcal D_\omega\) is the localized cost whose divergence is
validated by the compact criterion, provided that the same cost is controlled
by the corrected localized energy inequality on the terminal tail.
\end{definition}

\begin{definition}[Energy-controlled validated cost]
\label{paper7:def:energy-controlled-validated-cost}
Let \((V,P,a,b)\) be a repaired-gauge branch and let
\[
  \tilde{\mathfrak D}_{R_0}(\tau)
  =
  \nu\int |\nabla V|^2\phi_{R_0}
  +a_+(\tau)\int |V|^2\phi_{R_0}
\]
be the identity localized Navier--Stokes cost.  A validated compact cost
\(\mathfrak C_{\rm II}^{NS}\) is energy-controlled by the identity cost on a
terminal tail if there are constants \(C_1<\infty\) and
\(\tau_1\ge\tau_0\) such that
\[
  0\le \mathfrak C_{\rm II}^{NS}(\tau)
  \le C_1\tilde{\mathfrak D}_{R_0}(\tau)
  \qquad\text{for a.e. }\tau\ge\tau_1 .
\]
The identity cost is energy-controlled with \(C_1=1\) and
\(\tau_1=\tau_0\).  For a non-identity cost, this is the upper comparison
needed for a corrected localized energy inequality controlling
\(\tilde{\mathfrak D}_{R_0}\) to also control the validated cost.
\end{definition}

\begin{definition}[Local monotonicity carrier for a validated cost]
\label{paper7:def:local-monotonicity-carrier}
Let \(\omega\) be a repaired-gauge Navier--Stokes Type II branch on a terminal
renormalized tail \([\tau_0,\infty)\), and let
\(\mathfrak C_{\rm II}^{NS}\) be the validated compact cost for that branch.
A local monotonicity carrier for \(\mathfrak C_{\rm II}^{NS}\) consists of a
nonnegative locally integrable carrier cost \(\mathcal A_\omega\), a locally
absolutely continuous corrected localized energy \(\mathcal H_\omega\), a
nonnegative locally integrable remainder \(\mathcal R_\omega\), and constants
\(c_A>0\), \(0<C_A<\infty\) such that:
\begin{enumerate}[label=\textup{(\roman*)}]
\item \(\mathcal H_\omega(\tau_0)<\infty\) and
\[
  \inf_{\tau\ge\tau_0}\mathcal H_\omega(\tau)>-\infty;
\]
\item the carrier monotonicity inequality
\[
  \frac{d}{d\tau}\mathcal H_\omega(\tau)
  +c_A\mathcal A_\omega(\tau)
  +\mathcal R_\omega(\tau)\le0
\]
holds in distributions on \((\tau_0,\infty)\);
\item the validated compact cost is controlled by the carrier on the terminal
tail:
\[
  0\le \mathfrak C_{\rm II}^{NS}(\tau)
  \le C_A\mathcal A_\omega(\tau)
  \qquad\text{for a.e. }\tau\ge\tau_0 .
\]
\end{enumerate}
The carrier is local: \(\mathcal A_\omega\) may be the fixed cutoff identity
cost, a moving-cutoff identity cost, or another localized cost with an explicit
terminal comparison to the validated compact cost.
\end{definition}

\begin{lemma}[Monotonicity carrier gives local cost-to-exclusion data]
\label{paper7:lem:monotonicity-carrier-gives-cost-to-exclusion-data}
Let \(\omega\) be a repaired-gauge Navier--Stokes Type II branch.  Assume that
the validated compact cost \(\mathfrak C_{\rm II}^{NS}\) is nonnegative,
locally integrable on finite renormalized intervals, and satisfies
\[
  \int_{\tau_0}^{\infty}
  \mathfrak C_{\rm II}^{NS}(\tau)\,d\tau=\infty .
\]
If \(\mathfrak C_{\rm II}^{NS}\) admits a local monotonicity carrier in the
sense of \cref{paper7:def:local-monotonicity-carrier}, then \(\omega\) carries
the local cost-to-exclusion data of
\cref{paper7:def:local-cost-to-exclusion-data} with
\(\mathcal D_\omega=\mathfrak C_{\rm II}^{NS}\).
\end{lemma}

\begin{proof}
Let
\((\mathcal A_\omega,\mathcal H_\omega,\mathcal R_\omega^{\rm car},c_A,C_A)\)
be the carrier data.  The lower bound and finite initial value of
\(\mathcal H_\omega\) are exactly the energy bounds required in
\cref{paper7:def:local-cost-to-exclusion-data}.  The terminal comparison gives
\[
  \mathcal A_\omega(\tau)
  \ge C_A^{-1}\mathfrak C_{\rm II}^{NS}(\tau)
  \qquad\text{for a.e. }\tau\ge\tau_0 .
\]
Substituting this into the carrier monotonicity inequality yields
\[
  \frac{d}{d\tau}\mathcal H_\omega(\tau)
  +\frac{c_A}{C_A}\mathfrak C_{\rm II}^{NS}(\tau)
  +\mathcal R_\omega^{\rm car}(\tau)\le0
\]
in distributions.  Thus
\cref{paper7:def:local-cost-to-exclusion-data} holds with
\[
  \mathcal D_\omega=\mathfrak C_{\rm II}^{NS},
  \qquad c_\omega=c_A/C_A,
  \qquad \mathcal R_\omega=\mathcal R_\omega^{\rm car}.
\]
The divergent-cost item is the validated compact cost-divergence conclusion.
\end{proof}

\begin{lemma}[Fixed-cutoff monotonicity gives local cost-to-exclusion data]
\label{paper7:lem:fixed-monotonicity-gives-cost-to-exclusion-data}
Let \(\omega\) be represented on a terminal tail by a repaired-gauge
Navier--Stokes branch \((V,P,a,b)\).  Let \(\phi_{R_0}\) be the cutoff used in
the validated compact cost, let \(\mathfrak C_{\rm II}^{NS}\) be a
nonnegative locally integrable validated compact cost, and set
\[
  \mathcal E_\phi(\tau)
  =
  \frac12\int |V(y,\tau)|^2\phi_{R_0}(y)\,dy .
\]
Assume, after increasing \(\tau_0\) if necessary, that:
\begin{enumerate}[label=\textup{(\roman*)}]
\item \(\mathcal E_\phi(\tau_0)<\infty\);
\item the fixed-cutoff localized energy identity
\eqref{paper6a:eq:fixed-local-energy} holds on the terminal tail;
\item the fixed-cutoff error density \(B_{R_0}\) of
\cref{paper6a:def:fixed-error-density} satisfies
\[
  \int_{\tau_0}^{\infty}B_{R_0}(\tau)\,d\tau<\infty;
\]
\item the validated compact cost \(\mathfrak C_{\rm II}^{NS}\) is
energy-controlled by \(\tilde{\mathfrak D}_{R_0}\) in the sense of
\cref{paper7:def:energy-controlled-validated-cost};
\item the validated compact cost-divergence conclusion is
\[
  \int_{\tau_0}^{\infty}
  \mathfrak C_{\rm II}^{NS}(\tau)\,d\tau=\infty .
\]
\end{enumerate}
Then \(\omega\) carries the local cost-to-exclusion data of
\cref{paper7:def:local-cost-to-exclusion-data} with
\(\mathcal D_\omega=\mathfrak C_{\rm II}^{NS}\).
\end{lemma}

\begin{proof}
Let \(C_1,\tau_1\) be the constants in the energy-control comparison.  All
statements are terminal.  Since \(\mathfrak C_{\rm II}^{NS}\) is locally
integrable on finite intervals, removing any finite initial subinterval does
not change divergence of its tail integral.  After the allowed increase and
relabeling of \(\tau_0\), we may therefore assume that
\(\tau_0\ge\tau_1\), that \(\mathcal E_\phi(\tau_0)<\infty\), and that the
upper comparison holds on the whole tail.
\par
Define the corrected localized energy
\[
  \mathcal H_\omega(\tau)
  =
  \mathcal E_\phi(\tau)+\int_\tau^\infty B_{R_0}(s)\,ds .
\]
The finite-error hypothesis makes the correction tail finite for every
\(\tau\ge\tau_0\) and locally absolutely continuous.  Since
\(\mathcal E_\phi\ge0\) and \(B_{R_0}\ge0\), one has
\[
  \inf_{\tau\ge\tau_0}\mathcal H_\omega(\tau)\ge0,
  \qquad
  \mathcal H_\omega(\tau_0)<\infty .
\]
\Cref{paper6a:thm:fixed-corrected-monotonicity} gives, in distributions,
\[
  \frac{d}{d\tau}\mathcal H_\omega(\tau)
  +\frac12\tilde{\mathfrak D}_{R_0}(\tau)\le0 .
\]
By energy-control of the validated cost,
\[
  \tilde{\mathfrak D}_{R_0}(\tau)
  \ge C_1^{-1}\mathfrak C_{\rm II}^{NS}(\tau)
  \quad\text{for a.e. }\tau\ge\tau_0 .
\]
Therefore
\[
  \frac{d}{d\tau}\mathcal H_\omega(\tau)
  +\frac{1}{2C_1}\mathfrak C_{\rm II}^{NS}(\tau)\le0 .
\]
This is the corrected localized energy inequality in
\cref{paper7:def:local-cost-to-exclusion-data}, with
\[
  \mathcal D_\omega=\mathfrak C_{\rm II}^{NS},\qquad
  c_\omega=(2C_1)^{-1},\qquad
  \mathcal R_\omega=0 .
\]
The divergent-cost item is exactly the validated compact cost-divergence
conclusion.  Hence all local cost-to-exclusion data hold.
\end{proof}

\begin{lemma}[Moving-cutoff monotonicity gives local cost-to-exclusion data]
\label{paper7:lem:moving-monotonicity-gives-cost-to-exclusion-data}
Let \(\omega\) be represented on a terminal tail by a repaired-gauge
Navier--Stokes branch \((V,P,a,b)\), and let
\(\mathfrak C_{\rm II}^{NS}\) be a nonnegative locally integrable validated
compact cost.  Let
\[
  \Phi(\tau,y)=\phi(y/R(\tau))
\]
be a moving cutoff as in \cref{paper6a:def:moving-error-density}, with
\(R(\tau)\ge R_0\) locally absolutely continuous and nondecreasing, and set
\[
  \mathcal E_R(\tau)=\frac12\int |V(y,\tau)|^2\Phi(\tau,y)\,dy,
  \qquad
  \tilde{\mathfrak D}_{R(\tau)}(\tau)
  =
  \nu\int|\nabla V|^2\Phi
  +a_+(\tau)\int|V|^2\Phi .
\]
Assume, after increasing \(\tau_0\) if necessary, that:
\begin{enumerate}[label=\textup{(\roman*)}]
\item \(\mathcal E_R(\tau_0)<\infty\);
\item the localized energy identity is valid with the time-dependent cutoff
\(\Phi\);
\item the moving finite-error condition holds, for instance through summable
moving-annulus errors:
\[
  \int_{\tau_0}^{\infty}B_{\mathrm{mov}}(\tau)\,d\tau<\infty;
\]
\item the validated compact cost is controlled by the moving identity cost on
the terminal tail: for some \(C_m<\infty\),
\[
  0\le \mathfrak C_{\rm II}^{NS}(\tau)
  \le C_m\tilde{\mathfrak D}_{R(\tau)}(\tau)
  \qquad\text{for a.e. }\tau\ge\tau_0;
\]
\item the validated compact cost-divergence conclusion is
\[
  \int_{\tau_0}^{\infty}
  \mathfrak C_{\rm II}^{NS}(\tau)\,d\tau=\infty .
\]
\end{enumerate}
Then \(\omega\) carries the local cost-to-exclusion data of
\cref{paper7:def:local-cost-to-exclusion-data} with
\(\mathcal D_\omega=\mathfrak C_{\rm II}^{NS}\).
\end{lemma}

\begin{proof}
The moving finite-error condition gives
\(B_{\mathrm{mov}}\in L^1(\tau_0,\infty)\); if it is verified through
\cref{paper6a:def:summable-moving-errors}, this is precisely
\cref{paper6a:lem:summable-moving-errors}.  Define
\[
  \mathcal H_\omega(\tau)
  =
  \mathcal E_R(\tau)+\int_\tau^\infty B_{\mathrm{mov}}(s)\,ds .
\]
The finite moving-error condition makes the correction tail finite and locally
absolutely continuous.  Since \(\mathcal E_R\ge0\) and
\(B_{\mathrm{mov}}\ge0\),
\[
  \inf_{\tau\ge\tau_0}\mathcal H_\omega(\tau)\ge0,
  \qquad
  \mathcal H_\omega(\tau_0)<\infty .
\]
\Cref{paper6a:thm:moving-corrected-monotonicity} gives
\[
  \frac{d}{d\tau}\mathcal H_\omega(\tau)
  +\frac12\tilde{\mathfrak D}_{R(\tau)}(\tau)\le0
\]
in distributions.  Thus
\(\tilde{\mathfrak D}_{R(\tau)}\) is a local monotonicity carrier for the
validated compact cost, with \(c_A=1/2\), \(C_A=C_m\), and
\(\mathcal R_\omega=0\).  Applying
\cref{paper7:lem:monotonicity-carrier-gives-cost-to-exclusion-data} gives the
claimed local cost-to-exclusion data.
\end{proof}

\begin{lemma}[Corrected monotonicity bounds the validated cost]
\label{paper7:lem:corrected-monotonicity-bounds-cost}
Let \(\omega\) carry local cost-to-exclusion data in the sense of
\cref{paper7:def:local-cost-to-exclusion-data}, except do not assume the
divergence condition in item \textup{(i)}.  Then
\[
  \int_{\tau_0}^{\infty}\mathcal D_\omega(\tau)\,d\tau<\infty .
\]
\end{lemma}

\begin{proof}
Since \(\mathcal H_\omega\) is locally absolutely continuous and the corrected
localized energy inequality holds in distributions, integration over
\([\tau_0,T]\) gives, for every \(T>\tau_0\),
\[
  \mathcal H_\omega(T)-\mathcal H_\omega(\tau_0)
  +c_\omega\int_{\tau_0}^{T}\mathcal D_\omega(\tau)\,d\tau
  +\int_{\tau_0}^{T}\mathcal R_\omega(\tau)\,d\tau
  \le0 .
\]
The remainder is nonnegative, hence
\[
  c_\omega\int_{\tau_0}^{T}\mathcal D_\omega(\tau)\,d\tau
  \le
  \mathcal H_\omega(\tau_0)-\mathcal H_\omega(T)
  \le
  \mathcal H_\omega(\tau_0)
  -\inf_{\sigma\ge\tau_0}\mathcal H_\omega(\sigma).
\]
The right-hand side is finite and independent of \(T\).  Letting
\(T\to\infty\) and using monotone convergence for the nonnegative cost proves
the claim.
\end{proof}

\begin{theorem}[Local Navier--Stokes cost-to-exclusion theorem]
\label{paper7:thm:local-cost-to-exclusion}
Let \(\omega\) be a renormalized Navier--Stokes Type II branch on a retained
terminal compact stratum.  If \(\omega\) has active supercritical scaling, the
validated compact Type II cost-divergence conclusion, and local
cost-to-exclusion data of
\cref{paper7:def:local-cost-to-exclusion-data}, then \(\omega\) is excluded from
the Type II class under consideration.  Equivalently, on this local stratum,
\[
  \Pi_{\rm NS,II}(\omega):
  \mathcal S_\omega\wedge\mathcal B_\omega
  \Longrightarrow
  \mathcal E_\omega
\]
is a theorem.
\end{theorem}

\begin{proof}
Assume, for contradiction, that the branch remains in the Type II class on the
retained compact stratum.  The local cost-to-exclusion data give a corrected
localized energy inequality for the same nonnegative cost
\(\mathcal D_\omega\) whose divergent tail integral is the validated compact
cost-divergence conclusion.  By
\cref{paper7:lem:corrected-monotonicity-bounds-cost},
\[
  \int_{\tau_0}^{\infty}\mathcal D_\omega(\tau)\,d\tau<\infty .
\]
This contradicts the validated compact cost-divergence conclusion.  Hence the
branch cannot remain in the Type II class under consideration.
\end{proof}

\begin{definition}[Pointwise cost-divergence exclusion data]\label{paper7:def:cost-exclusion-data}
For a renormalized Type II solution \(\omega\), the pointwise
cost-divergence exclusion data are available when the following local analytic
statements hold:
\begin{enumerate}[label=\textup{(\roman*)}]
\item class membership: \(\omega\) lies in the renormalized Navier--Stokes
Type II class and satisfies the Type II analytic data before the Type II cost
criterion is evaluated;
\item cost compatibility: the active compact Type II cost criterion uses either
the Navier--Stokes identity-cost criterion
\[
  \mathfrak C_{\rm II}^{NS}=\tilde{\mathfrak D}_{R_0},
\]
or a localized cost explicitly comparable to
\(\tilde{\mathfrak D}_{R_0}\) on a terminal tail in the sense of
\cref{def:c4-cost-comparison}; in either case its conclusion is the validated
compact Type II cost-divergence conclusion for \(\omega\);
\item stability of the active scaling condition: the active supercritical
scaling condition remains in force after the repaired-gauge representation,
cost comparison, and criterion evaluation are attached;
\item stability of the cost-divergence conclusion: every renormalized
continuation preserving the active scaling condition and the validated compact
Type II cost-divergence conclusion remains in the same compact Type II
cost-divergence regime, unless it first enters one of the local alternatives
already isolated above;
\item local monotonicity-carrier data: on the retained compact stratum, the
validated compact cost admits a carrier in the sense of
\cref{paper7:def:local-monotonicity-carrier}.  This may be supplied by the
fixed or moving carrier subidentity of
\cref{paper7:def:carrier-subidentities}, the corresponding finite-error
condition, and the corresponding cost comparison.  Exact fixed or moving
localized energy identities are stronger sufficient inputs through
\cref{paper7:lem:fixed-monotonicity-gives-cost-to-exclusion-data,paper7:lem:moving-monotonicity-gives-cost-to-exclusion-data}.
\end{enumerate}
\end{definition}

\begin{lemma}[Canonical cost compatibility]
\label{paper7:lem:canonical-cost-compatibility}
Let \(\omega\) be represented on a terminal tail by a repaired-gauge
Navier--Stokes branch \((V,P,a,b)\).  Assume that the localized identity cost
\[
  \tilde{\mathfrak D}_{R_0}(\tau)
  =
  \nu\int_{\mathbb R^3}|\nabla V(y,\tau)|^2\phi_{R_0}(y)\,dy
  +
  a_+(\tau)\int_{\mathbb R^3}|V(y,\tau)|^2\phi_{R_0}(y)\,dy
\]
is measurable, nonnegative, and locally integrable on finite
\(\tau\)-intervals.  If the compact Type II criterion is evaluated with the
canonical Navier--Stokes identity cost, then
\[
  \mathfrak C_{\rm II}^{NS}=\tilde{\mathfrak D}_{R_0}
\]
is a valid Type II scale-exclusion cost, the comparison holds with
\[
  c_0=1,\qquad \tau_1=\tau_0,
\]
and the validated compact cost-divergence conclusion for \(\omega\) is the
identity-cost divergence statement
\[
  \int_{\tau_0}^{\infty}
  \tilde{\mathfrak D}_{R_0}(\tau)\,d\tau=\infty .
\]
If the compact criterion is not evaluated with this identity cost, the branch
is not being treated by the canonical compatibility statement and must instead
come with an explicit comparison as in
\cref{paper7:lem:comparable-cost-compatibility}.
\end{lemma}

\begin{proof}
The first assertion is the identity-cost case of the localized cost
definitions.  Since
\(\tilde{\mathfrak D}_{R_0}\) is assumed measurable, nonnegative, and
locally integrable on finite intervals, the canonical choice
\[
  \mathfrak C_{\rm II}^{NS}:=\tilde{\mathfrak D}_{R_0}
\]
satisfies the defining properties of a Type II scale-exclusion cost.  The
comparison inequality is the identity
\[
  \mathfrak C_{\rm II}^{NS}(\tau)
  =
  \tilde{\mathfrak D}_{R_0}(\tau)
  \ge
  1\cdot\tilde{\mathfrak D}_{R_0}(\tau)
\]
for a.e. \(\tau\ge\tau_0\), so \(c_0=1\) and \(\tau_1=\tau_0\).  Therefore
divergence of the active compact cost is precisely divergence of
\(\tilde{\mathfrak D}_{R_0}\) on the terminal tail.  This is the
canonical case of \cref{lem:c4-identity-cost,thm:c4-cost-comparison}.
\end{proof}

\begin{lemma}[Comparable cost compatibility]
\label{paper7:lem:comparable-cost-compatibility}
Let \(\omega\) be represented on a terminal tail by a repaired-gauge
Navier--Stokes branch \((V,P,a,b)\).  Let
\(\mathfrak C_{\rm II}^{NS}\) be a nonnegative, measurable, locally integrable
localized cost on finite \(\tau\)-intervals.  Assume that, on a later terminal
tail \([\tau_1,\infty)\), there is a constant \(c_0>0\) such that
\[
  \mathfrak C_{\rm II}^{NS}(\tau)
  \ge
  c_0\,\tilde{\mathfrak D}_{R_0}(\tau)
  \quad\text{for a.e. }\tau\ge\tau_1 .
\]
Then every identity-cost divergence conclusion
\[
  \int_{\tau_1}^{\infty}
  \tilde{\mathfrak D}_{R_0}(\tau)\,d\tau=\infty
\]
implies the corresponding compact Type II cost-divergence conclusion
\[
  \int_{\tau_1}^{\infty}
  \mathfrak C_{\rm II}^{NS}(\tau)\,d\tau=\infty .
\]
Consequently the use of \(\mathfrak C_{\rm II}^{NS}\) introduces no additional
escape from the compact-cost stratum beyond the explicit failure of this
comparison.
\end{lemma}

\begin{proof}
For every \(T>\tau_1\), the assumed comparison gives
\[
  \int_{\tau_1}^{T}
  \mathfrak C_{\rm II}^{NS}(\tau)\,d\tau
  \ge
  c_0\int_{\tau_1}^{T}
  \tilde{\mathfrak D}_{R_0}(\tau)\,d\tau .
\]
Letting \(T\to\infty\) proves the implication from identity-cost divergence to
divergence of the selected compact Type II cost.  This is precisely the
one-sided comparison needed by \cref{def:c4-cost-comparison} and used in
\cref{thm:c4-cost-comparison}.  If the inequality is unavailable, the selected
diagnostic cost is not connected to the Navier--Stokes identity dissipation on
this terminal tail; this is a cost-comparison failure, not a new unlabelled
Type II branch.
\end{proof}

\begin{lemma}[Local cost well-posedness]
\label{paper7:lem:local-cost-well-posed}
Let \((V,P,a,b)\) be a repaired-gauge branch on a finite renormalized interval
\(I=[\alpha,\beta]\), and let \(\phi_{R_0}\in C_c^\infty(B_{R_0})\),
\(\phi_{R_0}\ge0\).  Assume:
\begin{enumerate}[label=\textup{(\roman*)}]
\item \(V\in L^2(I;H^1(B_{R_0}))\);
\item \(a\in L^1(I)\);
\item either \(V\in L^\infty(I;L^3(\mathbb R^3))\), or more locally
\[
  \operatorname*{ess\,sup}_{\tau\in I}
  \|V(\tau)\|_{L^2(B_{R_0})}<\infty .
\]
\end{enumerate}
Then the localized identity cost
\[
  \tilde{\mathfrak D}_{R_0}(\tau)
  =
  \nu\int |\nabla V(y,\tau)|^2\phi_{R_0}(y)\,dy
  +
  a_+(\tau)\int |V(y,\tau)|^2\phi_{R_0}(y)\,dy
\]
is measurable, nonnegative, and integrable on \(I\).  Consequently, on every
compact renormalized window satisfying these local bounds,
\(\tilde{\mathfrak D}_{R_0}\) is a well-defined finite-interval cost.
\end{lemma}

\begin{proof}
The two summands are measurable because \(V\) is strongly measurable in the
local Sobolev spaces on the finite interval and \(a_+\) is measurable.  They are
nonnegative since \(\nu>0\), \(\phi_{R_0}\ge0\), and \(a_+\ge0\).
\par
For the gradient term, boundedness of \(\phi_{R_0}\) and the support condition
give
\[
  \int_I\int |\nabla V|^2\phi_{R_0}
  \le
  \|\phi_{R_0}\|_{L^\infty}
  \int_I\|\nabla V(\tau)\|_{L^2(B_{R_0})}^2\,d\tau
  <\infty .
\]
For the drift term, first suppose
\[
  \operatorname*{ess\,sup}_{\tau\in I}
  \|V(\tau)\|_{L^2(B_{R_0})}<\infty .
\]
Then
\[
  \int_I a_+(\tau)\int |V|^2\phi_{R_0}\,dy\,d\tau
  \le
  \|\phi_{R_0}\|_{L^\infty}
  \left(\operatorname*{ess\,sup}_{\tau\in I}
  \|V(\tau)\|_{L^2(B_{R_0})}^2\right)
  \|a_+\|_{L^1(I)}
  <\infty .
\]
If instead \(V\in L^\infty(I;L^3(\mathbb R^3))\), then Hölder on \(B_{R_0}\)
gives the local \(L^2\) bound
\[
  \|V(\tau)\|_{L^2(B_{R_0})}^2
  \le
  |B_{R_0}|^{1/3}\|V(\tau)\|_{L^3(\mathbb R^3)}^2
\]
for a.e. \(\tau\in I\), reducing to the previous case.  Thus both summands are
integrable on \(I\).
\end{proof}

\begin{lemma}[Local modulation integrability from absolutely continuous repaired gauges]
\label{paper7:lem:local-modulation-integrability}
Let \(u,p\) be a Navier--Stokes solution on a terminal physical time interval,
and let \((V,P)\) be its repaired-gauge renormalization on a terminal
renormalized tail.  Assume that, on every compact renormalized subinterval, the
physical concentration chart and the repaired-gauge path satisfy
\cref{paper6a:def:ac-initial-concentration-chart,paper6a:def:ac-repaired-gauge-path}.
Let \(a,b\) be the final modulation coefficients in the repaired-gauge
renormalized equation, as in
\cref{paper6a:lem:repaired-modulation-coefficients}.  Then
\[
  a,b\in L^1_{\rm loc}([\tau_0,\infty)).
\]
In particular, for every unit terminal window
\[
  I_n=[\tau_0+n,\tau_0+n+1],
\]
one has \(a\in L^1(I_n)\).  If the absolutely continuous chart, the
absolutely continuous repaired gauge, or the resulting local integrability of
the modulation coefficients is unavailable on such a compact window, then the
branch does not lie in the retained compact-cost stratum and falls into the
corresponding representation, repaired-gauge, or modulation-integrability
alternative.
\end{lemma}

\begin{proof}
On each compact renormalized window, the assumed absolutely continuous chart
and repaired-gauge path place the branch under the hypotheses of
\cref{paper6a:thm:ac-repaired-gauge-consequences}.  In particular,
\cref{paper6a:lem:ac-chart-regularity} gives
\[
  a,b\in L^1_{\rm loc},
\]
and \cref{paper6a:lem:suitable-gauge-pressure-modulation} identifies these
locally integrable functions as the final coefficients in the repaired-gauge
equation.  Restricting to \(I_n\) gives \(a\in L^1(I_n)\).  If the absolutely
continuous local chart, repaired-gauge path, or resulting coefficient
integrability is missing, then the ordered stratification records the first
missing input as a representation, repaired-gauge, or modulation-integrability
failure.
\end{proof}

\begin{lemma}[Critical annulus gives local \(L^2\) windows]
\label{paper7:lem:critical-annulus-local-l2}
Let \((V,P,a,b)\) be a repaired-gauge Type II branch satisfying the positive
finite critical annulus of \cref{def:c3-critical-mass} on a terminal tail:
\[
  \|V(\tau)\|_{L^3(\mathbb R^3)}\le M,\qquad \tau\ge\tau_0 .
\]
Then for every \(R_0<\infty\) and every unit terminal window \(I_n\),
\[
  \operatorname*{ess\,sup}_{\tau\in I_n}
  \|V(\tau)\|_{L^2(B_{R_0})}
  \le
  |B_{R_0}|^{1/6}M .
\]
Thus the local \(L^2\)-input required for the finite-window identity cost is
available on every retained compact cylinder.  If the bound
\(\sup_{\tau\ge\tau_0}\|V(\tau)\|_{L^3(\mathbb R^3)}<\infty\) fails, the branch
has already entered the critical-mass alternative in
\cref{def:c3-critical-mass-alternatives}.
\end{lemma}

\begin{proof}
For a.e. \(\tau\in I_n\), Hölder's inequality on the finite-measure ball gives
\[
  \|V(\tau)\|_{L^2(B_{R_0})}
  \le
  |B_{R_0}|^{1/6}\|V(\tau)\|_{L^3(B_{R_0})}
  \le
  |B_{R_0}|^{1/6}\|V(\tau)\|_{L^3(\mathbb R^3)}
  \le
  |B_{R_0}|^{1/6}M .
\]
Taking the essential supremum over \(I_n\) proves the estimate.  The final
sentence is the second ordered critical-mass alternative in
\cref{def:c3-critical-mass-alternatives}.
\end{proof}

\begin{corollary}[Compact-window data give cost well-posedness]
\label{paper7:cor:compact-window-data-cost-well-posed}
Let \((V,P,a,b)\) be a repaired-gauge branch on a terminal tail.  Fix
\(R_0<\infty\).  Suppose the positive finite critical annulus holds, the local
modulation-integrability conclusion of
\cref{paper7:lem:local-modulation-integrability} holds, and, for every unit
terminal window
\[
  I_n=[\tau_0+n,\tau_0+n+1],
  \qquad n\in\mathbb N,
\]
the branch has local windowed \(H^1(B_{R_0})\) control.  Then
\(\tilde{\mathfrak D}_{R_0}\) is a measurable nonnegative cost, integrable
on each \(I_n\).  If the branch also has the identity cost or an explicitly
comparable localized cost, then the compact Type II cost criterion can be
evaluated on every retained terminal window.  Failure of this conclusion is
exactly one of the local failures in the hypotheses: windowed \(H^1\) failure,
critical-annulus/local \(L^2\) failure, modulation-integrability failure, or
cost-comparison failure.
\end{corollary}

\begin{proof}
\Cref{paper7:lem:critical-annulus-local-l2} supplies the
\(L^\infty(I_n;L^2(B_{R_0}))\) input on every unit window, and
\cref{paper7:lem:local-modulation-integrability} supplies \(a\in L^1(I_n)\).
The assumed windowed \(H^1(B_{R_0})\) control gives
\[
  V\in L^2(I_n;H^1(B_{R_0})).
\]
Thus the hypotheses of \cref{paper7:lem:local-cost-well-posed} hold on each
\(I_n\).  The identity-cost case is then covered by
\cref{paper7:lem:canonical-cost-compatibility}; the explicitly comparable case
is covered by \cref{paper7:lem:comparable-cost-compatibility}.  Hence the
criterion is locally meaningful on every retained terminal window.  Conversely,
if the conclusion fails before criterion evaluation, at least one of the listed
local inputs is absent; by the ordered stratification, that absence is recorded
as the corresponding local alternative.
\end{proof}

\begin{proposition}[Active scaling stability package]
\label{paper7:prop:active-scaling-stability-package}
Let \(\omega\) be an admissible local Type II branch with active supercritical
scaling on a terminal tail.  A passage from a raw Type II chart to a
repaired-gauge chart is called scale-admissible on a later tail if its raw and
repaired physical scales satisfy
\[
  0<m\le
  \frac{\lambda_{\rm rep}(t)}{\lambda_{\rm raw}(t)}
  \le M<\infty
\]
on that tail.  On the compact-cost stratum, the active supercritical scaling
condition is unchanged by the following operations:
\begin{enumerate}[label=\textup{(\roman*)}]
\item restriction to a later terminal tail;
\item passage from an admissible raw Type II chart to a scale-admissible
repaired-gauge representation;
\item replacement of the canonical cost by an explicitly comparable localized
cost;
\item evaluation of the compact Type II cost criterion.
\end{enumerate}
If the repaired chart is not scale-admissible on any later tail, the branch is
not retained in the compact-cost stratum; it is assigned instead to the
corresponding representation, repaired-gauge, or scale-degeneracy alternative.
\end{proposition}

\begin{proof}
The active supercritical scaling condition is a terminal property of the
scale-center chart and is invariant under removal of a finite initial
renormalized interval.  Thus restriction from \([\tau_0,\infty)\) to
\([\tau_1,\infty)\), \(\tau_1\ge\tau_0\), cannot change the terminal scale
alternative.
\par
Suppose next that a scale-admissible repaired-gauge representation is
available.  Two-sided comparability of \(\lambda_{\rm rep}\) and
\(\lambda_{\rm raw}\) implies that any terminal statement depending only on the
scale class, such as supercriticality relative to the Type I rate, holds for
one scale if and only if it holds for the other after passing to the same
tail.  The repaired gauge changes coordinates and modulation variables inside
this comparable scale class; it does not replace the branch by a different
terminal scaling alternative.  If the representation cannot be constructed, or
if every repaired chart loses scale comparability, then the first failure is a
representation, repaired-gauge, or scale-degeneracy alternative.
\par
Cost comparison is scale-passive.  It changes only the nonnegative diagnostic
functional used to record cost divergence and does not alter the center,
physical scale, renormalized time, or active scaling hypothesis.  If the
comparison is unavailable, the failure is cost-comparison failure.  Finally,
criterion evaluation is a logical operation applied to the already constructed
branch and cost.  It does not modify the branch or its scale.  Hence active
scaling is stable through the listed operations exactly on the local stratum on
which those operations are admissible.
\end{proof}

\begin{lemma}[No unrecorded escape from compact cost divergence]
\label{paper7:lem:no-unrecorded-cost-escape}
Let \(\omega\) be an admissible local Type II branch with active supercritical
scaling on a terminal tail.  Before declaring that a continuation lies in the
compact cost-divergence stratum, run the ordered local validation tests:
\begin{enumerate}[label=\textup{(\roman*)}]
\item Type II analytic data and exhaustion data;
\item raw-chart selection, repaired-gauge representation, and scale
admissibility;
\item local state-space decomposition into the alternatives of
\cref{paper6:thm:state-decomposition};
\item positive finite critical \(L^3\) annulus;
\item the critical-tightness/exterior-leakage test at the retained scale;
\item local windowed \(H^1\) control;
\item local modulation integrability, cost well-posedness, and compatibility of
the chosen compact Type II cost;
\item local monotonicity-carrier data for the validated compact cost, supplied
by the fixed-cutoff package, the moving-cutoff package, or another explicit
localized carrier;
\item the compact Type II cost-divergence conclusion for the validated cost.
\end{enumerate}
Then exactly one of the following occurs:
\begin{enumerate}[label=\textup{(\alph*)}]
\item all tests pass, and \(\omega\) remains in the compact Type II
cost-divergence regime on the retained terminal branch;
\item the first failed test gives one of the already named local alternatives:
ordered exhaustion failure, representation failure, repaired-gauge or
scale-degeneracy failure, multibubble/multicore splitting, cascade, exterior
tightness failure, loss of local windowed \(H^1\) control, finite-cost
scale-collapse alternative, scale-rigid residual alternative, critical-mass
alternative, modulation-integrability failure, cost-compatibility failure,
carrier-subidentity failure, finite fixed-cutoff error failure, summable
moving-annulus error failure, carrier-comparison failure for the validated
cost, or failure of compact cost divergence for the validated cost.
\end{enumerate}
Thus a renormalized continuation preserving active scaling and validated compact
cost divergence cannot leave the compact cost-divergence regime without being
recorded by the state-space stratification.
\end{lemma}

\begin{proof}
The proof is the ordered-alternative argument.  First test whether the Type II
analytic and exhaustion data hold.  If not, the first failure is one of the
ordered exhaustion failures of \cref{def:c1-exhaustion-failures}.  If they
hold, test for raw-chart selection, repaired-gauge representation, and scale
admissibility.  Failure at this stage is one of the ordered representation
failures of \cref{def:c2-representation-failures}, possibly with the specific
scale-degeneracy recorded by
\cref{paper7:prop:active-scaling-stability-package}.
\par
Once a scale-admissible repaired branch is available, apply the local
state-space decomposition \cref{paper6:thm:state-decomposition}.  If the branch
falls into the multibubble, multicore, cascade, gauge-degenerate, finite-cost
scale-collapse, local \(H^1\)-loss, or scale-rigid residual alternatives, it has
left the compact single-core cost-divergence stratum through a named local
case.  In particular, an additional exterior concentration mechanism is not
discarded; if it retains critical mass it appears as a multibubble or multicore
branch, and if it escapes the retained compact region it appears as failure of
critical tightness.
\par
The first alternative in \cref{paper6:thm:state-decomposition} is retained only
as a compact single-core branch; it is not yet identified with compact cost
divergence.  On that retained branch the remaining tests below decide whether
the branch belongs to the compact cost-divergence stratum, has finite validated
cost, or exits through one of the remaining compact tests.
\par
On the retained compact single-core branch, test the critical \(L^3\)
alternatives in the order of \cref{def:c3-critical-mass-alternatives}.  If the
positive finite critical annulus of \cref{def:c3-critical-mass} is not reached
on a terminal tail, the corresponding critical-mass alternative is recorded.
Once the positive finite annulus holds, the remaining compact tests are
critical tightness and local windowed \(H^1\) control.  Failure of either is
precisely one of the alternatives in \cref{def:c5-compact-tests}.
\par
If all compact tests pass, local cost well-posedness and compatibility are
evaluated.  Local modulation integrability is supplied by
\cref{paper7:lem:local-modulation-integrability}; the local cost
well-posedness is supplied by
\cref{paper7:cor:compact-window-data-cost-well-posed}.  The compatibility is
the identity case of \cref{paper7:lem:canonical-cost-compatibility} or the
explicit comparison case of
\cref{paper7:lem:comparable-cost-compatibility}; failure of both is a
cost-compatibility failure.  Next test whether the validated compact cost has a
local monotonicity carrier.  In the fixed-cutoff route this means the fixed
carrier subidentity, finite fixed-cutoff error, and energy-control comparison.
In the moving-cutoff route this means the moving carrier subidentity, summable
moving-annulus errors, and moving-cost comparison.  Exact localized energy
identities are sufficient for these subidentities, but the one-sided suitable
local energy inequality is the actual input needed.  Failure of these inputs is
recorded as a carrier-subidentity failure, a finite fixed-cutoff error failure,
a summable moving-annulus error failure, or a carrier-comparison failure for
the validated cost.
Finally, evaluate the compact cost-divergence conclusion for the validated cost.
If it fails, the branch is a finite-cost or non-divergent compact branch and is
not counted as a compact cost-divergence branch.  The tests are ordered, so the
first failure is unique.  If no failure occurs, the branch has the analytic data,
scale-admissible repaired-gauge representation, compact single-core
state-space position, positive finite critical annulus, critical tightness,
local windowed \(H^1\) control, local modulation integrability, a compatible
compact cost, local monotonicity-carrier data, and the validated compact
cost-divergence conclusion.  Together with the assumed active scaling, these
hypotheses place the branch in the compact Type II cost-divergence regime.
Therefore there is no unrecorded escape.
\end{proof}

\begin{theorem}[Carrier validation alternative on the retained compact stratum]
\label{paper7:thm:carrier-validation-retained-stratum}
Let \(\omega\) be an admissible local Type II branch with active supercritical
scaling.  Suppose that the ordered validation tests in
\cref{paper7:lem:no-unrecorded-cost-escape} have reached the retained compact
single-core branch, that cost well-posedness and cost compatibility have
passed, and that the validated compact cost-divergence conclusion holds.  Then
exactly one of the following alternatives occurs:
\begin{enumerate}[label=\textup{(\roman*)}]
\item the validated compact cost admits local monotonicity-carrier data in the
sense of \cref{paper7:def:local-monotonicity-carrier};
\item the branch leaves the retained compact cost-divergence stratum through
the first failed carrier-validation test: carrier-subidentity failure, finite
fixed-cutoff error failure, summable moving-annulus error failure, or
carrier-comparison failure for the validated cost.
\end{enumerate}
Consequently, on the retained compact cost-divergence stratum, where the
adjacent carrier-failure alternatives are not part of the stratum, the
validated compact cost admits local monotonicity-carrier data.
\end{theorem}

\begin{proof}
After the branch has passed the Type II analytic, representation,
state-space, critical-annulus, critical-tightness, local windowed \(H^1\),
modulation, cost well-posedness, and cost-compatibility tests, the next ordered
test in \cref{paper7:lem:no-unrecorded-cost-escape} is precisely the
local monotonicity-carrier test for the validated compact cost.  If this test
passes, alternative \textup{(i)} holds.  If it fails, the first failed input is
recorded as one of the named adjacent carrier failures listed in
\cref{paper7:lem:no-unrecorded-cost-escape}: failure of the carrier
subidentity, failure of the fixed finite-error condition, failure of the moving
finite-error condition through the summable moving-annulus test, or failure of
the carrier comparison for the validated cost.  The ordered first-failure
convention makes the two alternatives disjoint and exhaustive.
\par
By definition, the retained compact cost-divergence stratum keeps only the
branch for which all preceding compact tests and the carrier-validation test
have passed before the validated compact cost-divergence conclusion is
attached.  Therefore a branch that remains in this stratum cannot realize one
of the carrier-failure alternatives, and the local monotonicity carrier exists.
\end{proof}

\begin{proposition}[Local verification of the cost-exclusion data]
\label{paper7:prop:local-cost-exclusion-data-verification}
On a retained terminal branch in the compact Type II cost-divergence stratum,
the pointwise data of \cref{paper7:def:cost-exclusion-data} are local
analytic statements, not additional global assumptions.  More precisely, assume
that the branch:
\begin{enumerate}[label=\textup{(\roman*)}]
\item has the local Type II analytic data and has not entered an exhaustion,
representation, modulation-integrability, multibubble/multicore, cascade,
exterior-tightness, scale-collapse, scale-rigid, critical-mass, or local
\(H^1\)-loss alternative;
\item admits a scale-admissible repaired-gauge representation on retained
windows from absolutely continuous local scale-center and repaired-gauge data;
\item has the positive finite critical \(L^3\) annulus, critical tightness at
the retained scale, local windowed \(H^1\) control, active supercritical
scaling, and the compact cost-divergence conclusion for the validated cost;
\item uses either the canonical identity cost
\(\tilde{\mathfrak D}_{R_0}\) or an explicitly comparable localized cost;
\item remains in the retained compact cost-divergence stratum after the
carrier-validation alternatives of
\cref{paper7:thm:carrier-validation-retained-stratum} have been evaluated.
\end{enumerate}
Then class membership, cost compatibility, stability of the active scaling
condition, stability of the validated cost-divergence conclusion, and the
cost-to-exclusion implication all hold for this branch.  Consequently, on this
local stratum, all clauses of \cref{paper7:def:cost-exclusion-data} are verified;
no estimate outside the retained local stratum and no exclusion of secondary
exterior concentration are being assumed.
\end{proposition}

\begin{proof}
Class membership is the retained local Type II analytic-data alternative after
the ordered exhaustion and representation tests have passed.
The local inputs needed to evaluate the compact cost are available by
\cref{paper7:lem:local-modulation-integrability,paper7:lem:critical-annulus-local-l2}
and \cref{paper7:cor:compact-window-data-cost-well-posed}.  Cost compatibility
then follows from
\cref{paper7:lem:canonical-cost-compatibility} in the identity-cost case and
from \cref{paper7:lem:comparable-cost-compatibility} in the explicitly
comparable-cost case.  Stability of the active scaling condition is
\cref{paper7:prop:active-scaling-stability-package}.  Stability of the
validated cost-divergence conclusion, with all exits routed to named local
alternatives, is \cref{paper7:lem:no-unrecorded-cost-escape}.  These prove
clauses \textup{(i)}--\textup{(iv)} of
\cref{paper7:def:cost-exclusion-data}.  Since the branch remains in the
retained compact cost-divergence stratum after carrier validation,
\cref{paper7:thm:carrier-validation-retained-stratum} gives local
monotonicity-carrier data for the validated compact cost.  These data provide
the local cost-to-exclusion data through
\cref{paper7:lem:monotonicity-carrier-gives-cost-to-exclusion-data}.  Then
\cref{paper7:thm:local-cost-to-exclusion} gives
\(\Pi_{\rm NS,II}(\omega)\), which is the analytic content of clause
\textup{(v)}.
\end{proof}

\begin{definition}[Classwise retained compact coverage]\label{paper7:def:classwise-retained-compact-coverage}
The renormalized Type II class has classwise retained compact coverage for the
validated compact cost if every branch with active supercritical scaling and
validated compact cost divergence either remains in the retained compact
cost-divergence stratum after the ordered validation tests of
\cref{paper7:lem:no-unrecorded-cost-escape}, or has already entered one of the
named adjacent alternatives in that ordered stratification.
\end{definition}

\begin{theorem}[Classwise retained compact coverage from ordered validation]
\label{paper7:thm:classwise-retained-compact-coverage}
Every admissible local Type II branch with active supercritical scaling and
validated compact Type II cost divergence has classwise retained compact
coverage in the sense of
\cref{paper7:def:classwise-retained-compact-coverage}.
\end{theorem}

\begin{proof}
Apply the ordered validation tests of
\cref{paper7:lem:no-unrecorded-cost-escape} to the branch.  The lemma gives two
possibilities.  If all tests pass, the branch remains in the compact Type II
cost-divergence regime on the retained terminal branch, which is precisely the
retained compact cost-divergence stratum.
\par
If a test fails, the first failed test is one of the named local alternatives in
\cref{paper7:lem:no-unrecorded-cost-escape}.  Since the present theorem assumes
the validated compact cost-divergence conclusion, the last listed failure,
failure of compact cost divergence for the validated cost, cannot occur.  Every
other first failure is an adjacent state-space stratum: exhaustion,
representation, repaired-gauge or scale-degeneracy failure,
multibubble/multicore splitting, cascade, exterior tightness failure, local
\(H^1\)-loss, finite-cost scale-collapse, scale-rigid residual behavior,
critical-mass failure, modulation-integrability failure, cost-compatibility
failure, carrier-subidentity failure, fixed-cutoff finite-error failure,
moving-annulus summability failure, or carrier-comparison failure.  Therefore
every branch satisfying the hypotheses either remains retained or has already
entered a named adjacent alternative.  This is exactly classwise retained
compact coverage.
\end{proof}

\begin{remark}[Scope of classwise retained compact coverage]
\label{paper7:rem:uniform-cost-scope}
\Cref{paper7:prop:local-cost-exclusion-data-verification} is deliberately local.
It proves the interface hypotheses for a branch only after the branch has been
placed in the validated compact cost-divergence stratum by the state-space
stratification.  Branches with multibubble, cascade, exterior leakage,
scale-collapse, scale-rigid, or rough-core behavior are not forced into this
stratum; they are routed to their own local alternatives.  Thus the
cost-divergence exclusion step is unconditional on the retained compact
stratum, while a classwise statement is obtained only after the adjacent
strata have been treated by their own local arguments.  The coverage part itself
is no longer an assumption; it is supplied by
\cref{paper7:thm:classwise-retained-compact-coverage}.
\end{remark}

\begin{definition}[Abstract cost-divergence exclusion theorem applicability]
\label{paper7:def:abstract-cost-exclusion-applicability}
An abstract Type II cost-divergence exclusion theorem is applicable to a
renormalized Type II solution if the following analytic identifications are
available:
\begin{enumerate}[label=\textup{(\roman*)}]
\item an abstract Type II exclusion theorem based on localized
renormalization-cost divergence;
\item hypotheses supplying supercritical scaling, a Type II concentration
scenario with bounded physical energy, and a localized energy monotonicity
formula at the active scale \(\lambda(t)\);
\item an identification of the theorem's renormalization-cost integral with a
validated compact Navier--Stokes cost \(\mathfrak C_{\rm II}^{NS}\), where
either
\[
  \mathfrak C_{\rm II}^{NS}=\tilde{\mathfrak D}_{R_0}
\]
or \(\mathfrak C_{\rm II}^{NS}\) is a localized cost satisfying the terminal
one-sided comparison with \(\tilde{\mathfrak D}_{R_0}\) in
\cref{def:c4-cost-comparison};
\item a domain embedding placing the repaired-gauge Navier--Stokes orbit and
its scale variables into the parabolic domain expected by the theorem without
changing the active scaling condition, validated cost-divergence conclusion, or
exclusion conclusion;
\item identification of the theorem's conclusion with
\cref{paper7:def:cost-exclusion-conclusion}.
\end{enumerate}
\end{definition}

\begin{lemma}[Applicability of the abstract cost-divergence theorem]
\label{paper7:lem:abstract-cost-exclusion}
Assume a renormalized Type II solution has the active supercritical scaling
condition, the validated compact Type II cost-divergence conclusion, and the
abstract applicability data of
\cref{paper7:def:abstract-cost-exclusion-applicability}.
Then the solution is excluded from the Type II class under consideration.
\end{lemma}

\begin{proof}
The abstract Type II theorem is applied with supercritical scaling,
bounded-energy Type II concentration, and the localized energy monotonicity
formula at the active scale.  The cost identification places the validated
compact Navier--Stokes cost in the theorem's renormalization-cost integral.  In
the identity case this cost is \(\tilde{\mathfrak D}_{R_0}\).  In the
comparable-cost case, the terminal one-sided comparison in
\cref{def:c4-cost-comparison} and
\cref{paper7:lem:comparable-cost-compatibility} show that the validated
Navier--Stokes cost-divergence conclusion implies divergence of the localized
cost used by the abstract theorem.  The domain embedding places the
repaired-gauge Navier--Stokes solution in the theorem's parabolic domain, and
the conclusion is identified with \cref{paper7:def:cost-exclusion-conclusion}.
Therefore the theorem gives exclusion from the Type II class.
\end{proof}

\begin{theorem}[Pointwise Navier--Stokes cost-divergence exclusion]
\label{paper7:thm:pointwise-cost-exclusion}
Let \(\omega\) be a renormalized Type II solution.  If \(\omega\) has the
active supercritical scaling condition, the validated compact Type II
cost-divergence conclusion, and the pointwise data of
\cref{paper7:def:cost-exclusion-data}, then \(\omega\) is excluded from the
Type II class under consideration.
\end{theorem}

\begin{proof}
Class membership and stability of the active scaling condition ensure that
\(\omega\) is evaluated in the renormalized Type II class and that the active
supercritical scaling condition is still the relevant one.  Cost compatibility
ensures that the validated compact cost-divergence conclusion is the
Navier--Stokes compact cost conclusion for \(\omega\), not an external or
incomparable cost statement.  Stability of the cost-divergence conclusion rules
out an escape from the compact Type II cost-divergence regime to a Type II case
outside the compact cost-divergence criterion inside the same class.
\par
It remains only to justify the transition from cost divergence to exclusion.
The local monotonicity-carrier data in \cref{paper7:def:cost-exclusion-data}
give the local cost-to-exclusion data by
\cref{paper7:lem:monotonicity-carrier-gives-cost-to-exclusion-data}.
Therefore
\cref{paper7:thm:local-cost-to-exclusion} applies to the active supercritical
scaling condition and the validated compact cost-divergence conclusion, giving
the claimed exclusion.  In an abstract theorem setting, the same implication is
also licensed by \cref{paper7:lem:abstract-cost-exclusion}.
\end{proof}

\begin{corollary}[Cost-divergence exclusion on the retained compact stratum]
\label{paper7:cor:retained-stratum-cost-exclusion}
Let \(\omega\) be a renormalized Type II solution with active supercritical
scaling and validated compact Type II cost divergence.  If \(\omega\) remains
in the retained compact cost-divergence stratum after the ordered validation
tests of \cref{paper7:lem:no-unrecorded-cost-escape}, then \(\omega\) is
excluded from the Type II class under consideration.
\end{corollary}

\begin{proof}
By \cref{paper7:prop:local-cost-exclusion-data-verification}, the pointwise
cost-divergence exclusion data of \cref{paper7:def:cost-exclusion-data} are
verified on this retained stratum.  Applying
\cref{paper7:thm:pointwise-cost-exclusion} gives the exclusion conclusion.
\end{proof}

\begin{definition}[Adjacent-stratum closure for the compact-cost channel]
\label{paper7:def:compact-cost-adjacent-closure}
The compact-cost channel has adjacent-stratum closure if every named adjacent
alternative produced by
\cref{paper7:lem:no-unrecorded-cost-escape} for a branch with active
supercritical scaling and validated compact Type II cost divergence has its own
stratum-local argument giving the same Type II exclusion conclusion, either
directly or by routing the branch to a previously listed non-compact-cost
alternative whose own closure theorem already gives the Type II exclusion
conclusion.
\end{definition}

\begin{corollary}[Uniform Navier--Stokes cost-divergence exclusion]
\label{paper7:cor:uniform-cost-exclusion}
Assume adjacent-stratum closure for the compact-cost channel in the sense of
\cref{paper7:def:compact-cost-adjacent-closure}.  Then every
renormalized Type II solution with active supercritical scaling condition and
validated compact Type II cost-divergence conclusion is either excluded in its
adjacent stratum or excluded by
\cref{paper7:cor:retained-stratum-cost-exclusion}.
\end{corollary}

\begin{proof}
By \cref{paper7:thm:classwise-retained-compact-coverage}, a branch with active
supercritical scaling and validated compact cost divergence either lies in one
of the named adjacent strata or remains in the retained compact
cost-divergence stratum.  The adjacent cases are handled by the
adjacent-stratum closure hypothesis.  In the retained compact case,
\cref{paper7:cor:retained-stratum-cost-exclusion} applies.
\end{proof}

\begin{definition}[Adjacent exits from cost-divergence exclusion]
\label{paper7:def:cost-exclusion-missing-hypotheses}
If the active supercritical scaling condition and validated compact
cost-divergence conclusion are present but the branch is not retained in the
compact cost-divergence stratum, then the first exit from the cost-exclusion
argument lies in the ordered list:
\[
  \begin{gathered}
  \text{class membership},\quad
  \text{cost compatibility},\quad
  \text{stability of the active scaling condition},\\
  \text{stability of the validated cost-divergence conclusion},\quad
  \text{local monotonicity-carrier validation}.
  \end{gathered}
\]
On the retained compact stratum none of these exits occurs.  The final item is
expanded by \cref{paper7:thm:carrier-validation-retained-stratum} into the
carrier-subidentity, fixed-finite-error, moving-annulus, and
carrier-comparison adjacent strata.
\end{definition}

\begin{lemma}[Carrier exits are adjacent strata, not retained compact branches]
\label{paper7:lem:missing-hypotheses-not-alternatives}
Assume a renormalized Type II solution already has the active supercritical
scaling condition and the validated compact Type II cost-divergence conclusion.
If one of the exits in \cref{paper7:def:cost-exclusion-missing-hypotheses}
occurs, then the branch is not in the retained compact cost-divergence stratum.
If no such exit occurs, the local cost-to-exclusion data are available.
\end{lemma}

\begin{proof}
The retained compact cost-divergence stratum is reached only after the ordered
tests of \cref{paper7:lem:no-unrecorded-cost-escape} have passed.  Hence a
failure of class membership, cost compatibility, active-scaling stability,
cost-divergence stability, or carrier validation is by definition an exit to an
adjacent stratum.  Conversely, if no such exit occurs, the carrier-validation
test has passed.  By \cref{paper7:thm:carrier-validation-retained-stratum}, the
validated compact cost admits local monotonicity-carrier data, and then
\cref{paper7:lem:monotonicity-carrier-gives-cost-to-exclusion-data} gives the
local cost-to-exclusion data.
\end{proof}

\begin{definition}[Pre-carrier validated compact-cost channel]
\label{paper7:def:precarrier-validated-compact-channel}
A renormalized Type II branch lies in the pre-carrier validated compact-cost
channel if, on a terminal tail, it has active supercritical scaling and the
validated compact Type II cost-divergence conclusion, and the ordered validation
tests of \cref{paper7:lem:no-unrecorded-cost-escape} have passed through the
following stages:
\begin{enumerate}[label=\textup{(\roman*)}]
\item the Type II analytic data and ordered exhaustion tests;
\item raw-chart selection, repaired-gauge representation, pressure
reconstruction, modulation data, and scale-admissibility of the repaired chart;
\item the local state-space decomposition without entering a multibubble,
multicore, cascade, exterior-tightness, local \(H^1\)-loss,
finite-cost scale-collapse, or scale-rigid residual alternative;
\item the positive finite critical \(L^3\) annulus;
\item critical tightness at the retained scale and local windowed \(H^1\)
control;
\item local modulation integrability, local cost well-posedness, and
compatibility of the chosen compact Type II cost with the validated
Navier--Stokes cost.
\end{enumerate}
This channel deliberately stops immediately before local monotonicity-carrier
validation.  Thus it contains all compact-cost interface data needed to ask for
the fixed-cutoff or moving-cutoff carrier estimates, but it does not assert
that those estimates have already been proved.
\end{definition}

\begin{definition}[Fixed and moving carrier subidentities]
\label{paper7:def:carrier-subidentities}
For the fixed cutoff \(\phi_{R_0}\), the fixed carrier subidentity is the
distributional estimate
\[
  \frac{d}{d\tau}\mathcal E_\phi(\tau)
  +\frac12\tilde{\mathfrak D}_{R_0}(\tau)
  \le B_{R_0}(\tau),
  \qquad
  \mathcal E_\phi(\tau)=\frac12\int |V|^2\phi_{R_0}.
\]
For a moving cutoff \(\Phi(\tau,y)=\phi(y/R(\tau))\), the moving carrier
subidentity is the distributional estimate
\[
  \frac{d}{d\tau}\mathcal E_R(\tau)
  +\frac12\tilde{\mathfrak D}_{R(\tau)}(\tau)
  \le B_{\mathrm{mov}}(\tau),
  \qquad
  \mathcal E_R(\tau)=\frac12\int |V|^2\Phi .
\]
Here \(B_{R_0}\) and \(B_{\mathrm{mov}}\) are the absolute error densities of
\cref{paper6a:def:fixed-error-density,paper6a:def:moving-error-density}.  The
subidentity is the one-sided estimate needed for corrected monotonicity; an
exact localized energy identity is sufficient but not necessary.
\end{definition}

\begin{lemma}[Suitable local energy gives carrier subidentities]
\label{paper7:lem:suitable-energy-gives-carrier-subidentities}
Assume the repaired-gauge representation of a pre-carrier branch includes the
local suitable energy inequality on compact renormalized cylinders.  Then the
fixed carrier subidentity holds for every nonnegative compactly supported fixed
cutoff.  The moving carrier subidentity holds for every nonnegative
finite-valued moving cutoff \(\Phi(\tau,y)=\phi(y/R(\tau))\) with
\(R\in AC_{\mathrm{loc}}\), after the usual approximation of \(R\) by smooth
positive radii on compact \(\tau\)-intervals.
\end{lemma}

\begin{proof}
Apply the local suitable energy inequality in the physical variables to the
pullback of the nonnegative renormalized test function.  The repaired-gauge
change of variables is locally absolutely continuous in time and smooth in
space on compact cylinders, so the distributional chain rule gives the
renormalized local energy inequality with the drift terms
\(a(V+y\cdot\nabla V)\) and \(b\cdot\nabla V\).  For a fixed cutoff this is
the fixed localized energy identity with equality replaced by \(\le\).  Adding
\(\frac12\tilde{\mathfrak D}_{R_0}\) to both sides and bounding the cutoff,
pressure, translation, and scale terms by their absolute values gives the fixed
carrier subidentity.
\par
For a moving cutoff, use smooth positive approximations \(R_k\to R\) in
\(AC_{\mathrm{loc}}\) on each compact time interval.  The local energy
inequality applies to \(\Phi_k(\tau,y)=\phi(y/R_k(\tau))\), and the preceding
argument gives the moving subidentity with \(B_{\mathrm{mov},k}\).  Passing to
the limit in distributions uses the local finite-energy and local \(H^1\)
bounds already present on the pre-carrier channel.  The term
\(\partial_\tau\Phi\) is interpreted as the \(L^1_{\mathrm{loc}}\) derivative
coming from \(R\in AC_{\mathrm{loc}}\).  The limit is exactly the moving
carrier subidentity.
\end{proof}

\begin{definition}[Fixed-cutoff carrier closure package]
\label{paper7:def:fixed-carrier-closure-package}
A branch in the pre-carrier validated compact-cost channel satisfies the
fixed-cutoff carrier closure package if, for the cutoff \(\phi_{R_0}\) used to
evaluate the validated compact cost and for
\[
  \mathcal E_\phi(\tau)=\frac12\int |V(y,\tau)|^2\phi_{R_0}(y)\,dy,
\]
the following local analytic statements hold on a terminal tail:
\begin{enumerate}[label=\textup{(\roman*)}]
\item \(\mathcal E_\phi(\tau_0)<\infty\);
\item the fixed carrier subidentity of
\cref{paper7:def:carrier-subidentities} is valid for the repaired-gauge branch;
\item the fixed-cutoff error density \(B_{R_0}\) of
\cref{paper6a:def:fixed-error-density} has finite tail integral,
\[
  \int_{\tau_0}^{\infty}B_{R_0}(\tau)\,d\tau<\infty;
\]
\item the validated compact cost is energy-controlled by
\(\tilde{\mathfrak D}_{R_0}\) in the sense of
\cref{paper7:def:energy-controlled-validated-cost}.
\end{enumerate}
\end{definition}

\begin{definition}[Moving-cutoff carrier closure package]
\label{paper7:def:moving-carrier-closure-package}
A branch in the pre-carrier validated compact-cost channel satisfies the
moving-cutoff carrier closure package if there is a finite-valued
nondecreasing locally absolutely continuous radius \(R(\tau)\ge R_0\), with
\(R(\tau)\to\infty\) and moving cutoff
\(\Phi(\tau,y)=\phi(y/R(\tau))\), such that, on a terminal tail:
\begin{enumerate}[label=\textup{(\roman*)}]
\item
\[
  \mathcal E_R(\tau_0)
  =
  \frac12\int |V(y,\tau_0)|^2\Phi(\tau_0,y)\,dy
  <\infty;
\]
\item the moving carrier subidentity of
\cref{paper7:def:carrier-subidentities} is valid with the time-dependent cutoff
\(\Phi\);
\item the moving-annulus errors are summable in the sense of
\cref{paper6a:def:summable-moving-errors}, which implies
\[
  \int_{\tau_0}^{\infty}B_{\mathrm{mov}}(\tau)\,d\tau<\infty
\]
by \cref{paper6a:lem:summable-moving-errors};
\item the validated compact cost is controlled by the moving identity cost on
the terminal tail: for some \(C_m<\infty\),
\[
  0\le \mathfrak C_{\rm II}^{NS}(\tau)
  \le C_m\tilde{\mathfrak D}_{R(\tau)}(\tau)
  \qquad\text{for a.e. }\tau\ge\tau_0 .
\]
\end{enumerate}
\end{definition}

\begin{theorem}[Non-carrier adjacent exits are routed before carrier validation]
\label{paper7:thm:noncarrier-exits-route}
Let \(\omega\) be an admissible local Type II branch with active supercritical
scaling and validated compact Type II cost divergence.  Apply the ordered tests
of \cref{paper7:lem:no-unrecorded-cost-escape}.  If a first failed test exists
and is not one of the carrier-validation failures
\[
  \text{carrier subidentity},\quad
  \text{fixed finite error},\quad
  \text{moving-annulus summability},\quad
  \text{carrier comparison},
\]
then the branch has already failed to enter the pre-carrier validated
compact-cost channel of
\cref{paper7:def:precarrier-validated-compact-channel}.  More precisely, its
first failure is one of the following named non-carrier strata:
\[
\begin{gathered}
\text{ordered exhaustion failure, representation failure, repaired-gauge or
scale-degeneracy failure},\\
\text{multibubble or multicore splitting, cascade, exterior-tightness failure,
local \(H^1\)-loss},\\
\text{finite-cost scale-collapse, scale-rigid residual behavior,
critical-mass failure},\\
\text{modulation-integrability failure, cost-well-posedness failure, or
cost-compatibility failure}.
\end{gathered}
\]
In any ambient class in which the corresponding local closure theorems for
these named strata are imposed, such a branch is not a remaining compact-cost
branch.
\end{theorem}

\begin{proof}
The validation procedure in
\cref{paper7:lem:no-unrecorded-cost-escape} is ordered.  The
pre-carrier channel is exactly the set of branches for which all tests before
local monotonicity-carrier validation have passed.  Therefore any first failure
before the carrier-validation row prevents membership in
\cref{paper7:def:precarrier-validated-compact-channel}.
\par
The same lemma lists every possible first failure before carrier validation:
the exhaustion failures are those of
\cref{def:c1-exhaustion-failures}; representation and repaired-gauge failures
are those of
\cref{def:c2-representation-failures,def:c2r-minimal-failures} together with
scale-degeneracy from
\cref{paper7:prop:active-scaling-stability-package}; the state-space failures
are the alternatives in
\cref{paper6:thm:state-decomposition}; critical-mass failures are the
alternatives of \cref{def:c3-critical-mass-alternatives}; compact tightness
and local \(H^1\)-loss are the tests of \cref{def:c5-compact-tests}; and the
modulation, cost-well-posedness, and compatibility failures are exactly the
failures recorded in
\cref{paper7:lem:local-modulation-integrability,paper7:cor:compact-window-data-cost-well-posed,paper7:lem:canonical-cost-compatibility,paper7:lem:comparable-cost-compatibility}.
Thus the displayed list is exhaustive for non-carrier first failures.
\par
If the surrounding Type II class includes the local closure theorem assigned to
the named stratum, then the branch is handled there.  It is therefore not a
branch remaining to be treated by the retained compact-cost carrier argument.
\end{proof}

\begin{theorem}[Pre-carrier reduction to carrier validation]
\label{paper7:thm:precarrier-reduces-to-carrier}
Let \(\omega\) be in the pre-carrier validated compact-cost channel.  Then
exactly one of the following alternatives occurs:
\begin{enumerate}[label=\textup{(\roman*)}]
\item the validated compact cost admits local monotonicity-carrier data in the
sense of \cref{paper7:def:local-monotonicity-carrier};
\item the branch exits the retained compact cost-divergence stratum through the
first failed carrier-validation test: carrier-subidentity failure,
fixed-cutoff finite-error failure, moving-annulus summability failure, or
carrier-comparison failure for the validated cost.
\end{enumerate}
If the first alternative occurs, \(\omega\) is excluded by the retained
compact-cost argument.
\end{theorem}

\begin{proof}
By definition of the pre-carrier channel, every ordered validation test before
carrier validation has passed, and the validated compact cost-divergence
conclusion is already present.  Thus the next unresolved row in
\cref{paper7:lem:no-unrecorded-cost-escape} is precisely the local
monotonicity-carrier row.  Applying
\cref{paper7:thm:carrier-validation-retained-stratum} gives the two
alternatives listed in the statement, and the first-failure convention makes
them disjoint.
\par
If the carrier exists, then
\cref{paper7:lem:monotonicity-carrier-gives-cost-to-exclusion-data} gives the
local cost-to-exclusion data, and
\cref{paper7:thm:local-cost-to-exclusion} excludes the branch.  Equivalently,
the same conclusion is
\cref{paper7:cor:retained-stratum-cost-exclusion}.
\end{proof}

\begin{theorem}[Fixed-cutoff carrier closure]
\label{paper7:thm:fixed-carrier-closure}
Let \(\omega\) be in the pre-carrier validated compact-cost channel.  If
\(\omega\) satisfies the fixed-cutoff carrier closure package of
\cref{paper7:def:fixed-carrier-closure-package}, then \(\omega\) is excluded
from the Type II class under consideration.
\end{theorem}

\begin{proof}
The pre-carrier channel supplies active supercritical scaling, local
well-posedness and compatibility of the validated compact cost, and the
validated compact cost-divergence conclusion.  The fixed-cutoff carrier package
supplies finite initial localized energy, the fixed carrier subidentity, finite
fixed-cutoff error, and energy-control comparison of the validated cost with
\(\tilde{\mathfrak D}_{R_0}\).  Let \(C_1,\tau_1\) be the constants in that
comparison and increase the tail initial time so that the comparison holds from
\(\tau_0\) onward.  Define
\[
  \mathcal H_\omega(\tau)
  =
  \mathcal E_\phi(\tau)+\int_\tau^\infty B_{R_0}(s)\,ds .
\]
The finite-error condition makes \(\mathcal H_\omega\) locally absolutely
continuous and bounded below.  The fixed carrier subidentity gives
\[
  \frac{d}{d\tau}\mathcal H_\omega(\tau)
  +\frac12\tilde{\mathfrak D}_{R_0}(\tau)\le0
\]
in distributions.  Since
\(\mathfrak C_{\rm II}^{NS}\le C_1\tilde{\mathfrak D}_{R_0}\), this implies
\[
  \frac{d}{d\tau}\mathcal H_\omega(\tau)
  +\frac{1}{2C_1}\mathfrak C_{\rm II}^{NS}(\tau)\le0 .
\]
Together with the validated compact cost-divergence conclusion, these are the
local cost-to-exclusion data with
\(\mathcal D_\omega=\mathfrak C_{\rm II}^{NS}\).  Equivalently, the fixed
package is a successful carrier-validation row.  Since all pre-carrier tests
and the validated compact cost-divergence conclusion are already present,
\(\omega\) is on the retained compact cost-divergence stratum after carrier
validation.  Applying \cref{paper7:thm:local-cost-to-exclusion} gives the Type
II exclusion conclusion.
\end{proof}

\begin{theorem}[Moving-cutoff carrier closure]
\label{paper7:thm:moving-carrier-closure}
Let \(\omega\) be in the pre-carrier validated compact-cost channel.  If
\(\omega\) satisfies the moving-cutoff carrier closure package of
\cref{paper7:def:moving-carrier-closure-package}, then \(\omega\) is excluded
from the Type II class under consideration.
\end{theorem}

\begin{proof}
The pre-carrier channel supplies active supercritical scaling, local
well-posedness and compatibility of the validated compact cost, and the
validated compact cost-divergence conclusion.  The moving-cutoff carrier
package supplies the finite initial moving localized energy, the
moving carrier subidentity, summable moving-annulus errors, and
the terminal comparison
\[
  0\le \mathfrak C_{\rm II}^{NS}(\tau)
  \le C_m\tilde{\mathfrak D}_{R(\tau)}(\tau).
\]
By \cref{paper6a:lem:summable-moving-errors}, the summability condition gives
\(B_{\mathrm{mov}}\in L^1(\tau_0,\infty)\).  Define
\[
  \mathcal H_\omega(\tau)
  =
  \mathcal E_R(\tau)+\int_\tau^\infty B_{\mathrm{mov}}(s)\,ds .
\]
The finite moving-error condition makes \(\mathcal H_\omega\) locally
absolutely continuous and bounded below.  The moving carrier subidentity gives
\[
  \frac{d}{d\tau}\mathcal H_\omega(\tau)
  +\frac12\tilde{\mathfrak D}_{R(\tau)}(\tau)\le0
\]
in distributions.  The moving comparison gives
\[
  \frac{d}{d\tau}\mathcal H_\omega(\tau)
  +\frac{1}{2C_m}\mathfrak C_{\rm II}^{NS}(\tau)\le0 .
\]
Together with the validated compact cost-divergence conclusion, these are the
local cost-to-exclusion data with
\(\mathcal D_\omega=\mathfrak C_{\rm II}^{NS}\).  Equivalently, the moving
package is a successful carrier-validation row.  Since all pre-carrier tests
and the validated compact cost-divergence conclusion are already present,
\(\omega\) is on the retained compact cost-divergence stratum after carrier
validation.  Applying \cref{paper7:thm:local-cost-to-exclusion} excludes the
branch.
\end{proof}

\begin{lemma}[Pre-carrier compact critical mass gives finite localized energy]
\label{paper7:lem:precarrier-finite-local-energy}
Let \(\omega\) be in the pre-carrier validated compact-cost channel, represented
by \(V\) on a terminal tail.  Then, for every finite \(R<\infty\) and every
\(\tau\) on the tail,
\[
  \int_{B_R}|V(y,\tau)|^2\,dy<\infty .
\]
Consequently the initial localized-energy clauses in the fixed and moving
carrier closure packages are automatic for every compactly supported fixed
cutoff and every moving cutoff with \(R(\tau_0)<\infty\).
\end{lemma}

\begin{proof}
The pre-carrier channel includes the positive finite critical \(L^3\) annulus.
By \cref{def:c3-critical-mass}, after passing to the terminal tail there is
\(M<\infty\) such that
\[
  \|V(\tau)\|_{L^3(\mathbb R^3)}\le M
  \qquad\text{for all }\tau\ge\tau_0 .
\]
For each finite ball \(B_R\), Holder's inequality gives
\[
  \int_{B_R}|V(y,\tau)|^2\,dy
  \le
  |B_R|^{1/3}\|V(\tau)\|_{L^3(\mathbb R^3)}^2
  \le |B_R|^{1/3}M^2<\infty .
\]
If \(\phi_{R_0}\) is compactly supported, then
\(\mathcal E_\phi(\tau_0)\) is bounded by the same estimate on a large enough
ball.  If \(\Phi(\tau,y)=\phi(y/R(\tau))\) and \(R(\tau_0)<\infty\), then
\(\Phi(\tau_0,\cdot)\) is also compactly supported, and the same estimate gives
\(\mathcal E_R(\tau_0)<\infty\).
\end{proof}

\begin{lemma}[Nested moving cutoffs give carrier comparison]
\label{paper7:lem:nested-moving-carrier-comparison}
Let \(\omega\) be in the pre-carrier validated compact-cost channel, and
assume the validated compact cost is energy-controlled by the fixed identity
cost:
\[
  0\le \mathfrak C_{\rm II}^{NS}(\tau)
  \le C_1\tilde{\mathfrak D}_{R_0}(\tau)
  \qquad\text{for a.e. }\tau\ge\tau_1 .
\]
Let \(\Phi(\tau,y)\) be a moving cutoff such that, after increasing the initial
time if necessary,
\[
  \Phi(\tau,y)\ge \phi_{R_0}(y)
  \qquad\text{for all }y\in\mathbb R^3
  \text{ and a.e. }\tau\ge\tau_1 .
\]
Then the validated compact cost is controlled by the moving identity cost:
\[
  0\le \mathfrak C_{\rm II}^{NS}(\tau)
  \le C_1\tilde{\mathfrak D}_{R(\tau)}(\tau)
  \qquad\text{for a.e. }\tau\ge\tau_1 .
\]
In particular, for the canonical identity cost the carrier comparison is
automatic with \(C_1=1\).
\end{lemma}

\begin{proof}
Because \(\Phi(\tau,\cdot)\ge\phi_{R_0}\) and both cutoffs are nonnegative,
\[
  \nu\int |\nabla V|^2\phi_{R_0}
  \le
  \nu\int |\nabla V|^2\Phi(\tau,\cdot),
\]
and, since \(a_+(\tau)\ge0\),
\[
  a_+(\tau)\int |V|^2\phi_{R_0}
  \le
  a_+(\tau)\int |V|^2\Phi(\tau,\cdot).
\]
Adding these two inequalities gives
\[
  \tilde{\mathfrak D}_{R_0}(\tau)
  \le
  \tilde{\mathfrak D}_{R(\tau)}(\tau)
  \qquad\text{for a.e. }\tau\ge\tau_1 .
\]
Combining this with the fixed energy-control comparison yields the displayed
moving comparison.  The canonical identity cost is energy-controlled by
\cref{paper7:def:energy-controlled-validated-cost} with \(C_1=1\).
\end{proof}

\begin{definition}[Finite-error carrier selection]
\label{paper7:def:finite-error-carrier-selection}
The pre-carrier validated compact-cost channel has finite-error carrier
selection if every branch in that channel whose validated compact cost is
energy-controlled by the fixed identity cost satisfies, after passing to a
later terminal tail if necessary, at least one of the following two local
alternatives:
\begin{enumerate}[label=\textup{(\roman*)}]
\item the fixed carrier subidentity of
\cref{paper7:def:carrier-subidentities} is valid and
\[
  \int_{\tau_0}^{\infty}B_{R_0}(\tau)\,d\tau<\infty ;
\]
\item there exists a finite-valued nondecreasing locally absolutely continuous
radius \(R(\tau)\ge R_0\), with \(R(\tau)\to\infty\) and moving cutoff
\(\Phi(\tau,y)=\phi(y/R(\tau))\), such that
\[
  \Phi(\tau,y)\ge\phi_{R_0}(y)
  \qquad\text{for all }y\text{ and a.e. }\tau,
\]
the moving carrier subidentity of
\cref{paper7:def:carrier-subidentities} is valid with the time-dependent cutoff
\(\Phi\), and the moving-annulus errors are summable in the sense of
\cref{paper6a:def:summable-moving-errors}.
\end{enumerate}
This condition contains only the local one-sided energy estimate and finite-error
selection statements.  It does not include finite initial localized energy or
carrier comparison, because those are supplied by
\cref{paper7:lem:precarrier-finite-local-energy,paper7:lem:nested-moving-carrier-comparison}.
\end{definition}

\begin{definition}[Annular finite-error selection]
\label{paper7:def:annular-finite-error-selection}
The pre-carrier validated compact-cost channel has annular finite-error
selection if every branch in that channel whose validated compact cost is
energy-controlled by the fixed identity cost satisfies, after passing to a
later terminal tail if necessary, at least one of the following two alternatives:
\begin{enumerate}[label=\textup{(\roman*)}]
\item the fixed error density has finite tail integral,
\[
  \int_{\tau_0}^{\infty}B_{R_0}(\tau)\,d\tau<\infty;
\]
\item there exists a finite-valued nondecreasing locally absolutely continuous
nested radius \(R(\tau)\ge R_0\), with \(R(\tau)\to\infty\) and
\(\Phi(\tau,y)=\phi(y/R(\tau))\ge\phi_{R_0}(y)\), such that the moving-annulus
errors are summable in the sense of
\cref{paper6a:def:summable-moving-errors}.
\end{enumerate}
This is the purely finite-error part of the carrier problem.  It does not
include local energy subidentities, finite initial energy, or carrier comparison.
\end{definition}

\begin{definition}[Componentwise fixed-or-moving finite-error estimates]
\label{paper7:def:componentwise-annular-finite-error}
For a branch in the pre-carrier validated compact-cost channel, the
componentwise annular finite-error estimates hold if, after passing to a later
terminal tail if necessary, either the six fixed-cutoff components
\[
\begin{gathered}
  I_1=\left|\int |V|^2\Delta\phi_{R_0}\right|,\qquad
  I_2=\left|\int |V|^2V\cdot\nabla\phi_{R_0}\right|,\qquad
  I_3=\left|\int P\,V\cdot\nabla\phi_{R_0}\right|,\\
  I_4=|b|\left|\int |V|^2\nabla\phi_{R_0}\right|,\qquad
  I_5=a_-M_\phi,\qquad
  I_6=a_+\int |V|^2\bigl(-y\cdot\nabla\phi_{R_0}\bigr)
\end{gathered}
\]
belong to \(L^1(\tau_0,\infty)\), or there is a finite-valued nondecreasing
locally absolutely continuous nested
radius \(R(\tau)\ge R_0\), with \(R(\tau)\to\infty\) and
\(\Phi(\tau,y)=\phi(y/R(\tau))\ge\phi_{R_0}(y)\), such that each of the seven
moving annular components
\[
\begin{gathered}
  J_1=\left|\int |V|^2\partial_\tau\Phi\right|,\qquad
  J_2=\left|\int |V|^2\Delta\Phi\right|,\qquad
  J_3=\left|\int |V|^2V\cdot\nabla\Phi\right|,\\
  J_4=\left|\int P\,V\cdot\nabla\Phi\right|,\qquad
  J_5=|b|\left|\int |V|^2\nabla\Phi\right|,\\
  J_6=a_-\int |V|^2\Phi,\qquad
  J_7=a_+\int |V|^2\bigl(-y\cdot\nabla\Phi\bigr)
\end{gathered}
\]
belongs to \(L^1(\tau_0,\infty)\).
\end{definition}

\begin{remark}[What the componentwise estimates isolate]
\label{paper7:rem:componentwise-annular-residual}
The quantities in
\cref{paper7:def:componentwise-annular-finite-error} are exactly the fixed and
moving carrier-error components: in the fixed route, viscous cutoff,
convection, pressure-tail, center drift, negative scale drift, and the
positive scale-shell drift; in the moving route, cutoff transition, viscous
cutoff, convection, pressure-tail, center drift, negative scale drift, and the
positive scale-shell drift.  These diagnostics need not be closed by a single
global tail estimate.  In the stratified use of Section~7, persistent failure
of the cutoff, viscous, convection, or pressure diagnostics is read as shell
mass that cannot be evacuated from selected windows and is therefore routed to
exterior leakage, exterior concentration, loss of critical tightness, or local
windowed \(H^1\)-loss.  Persistent failure of the translation diagnostic is
read as a modulation-driven adjacent branch, while persistent failure of the
positive scale-shell term is read as a scale-dynamics branch.  Thus the
component list records the local diagnostics seen by the carrier test; it is
not intended to force the proof into a single classical PDE closure argument.
For radial nonincreasing cutoffs, the scale-shell term
\[
  -\frac{a}{2}\int |V|^2 y\cdot\nabla\Phi
\]
has only an \(a_+\)-contribution in the residual carrier error.  The \(a_-\)
part is favorable and is already absorbed into the negative-drift routing.
The negative scale-drift component \(J_6\) is not to be replaced by an
unproved equivalence with compact-cost divergence: if \(J_6\) is not summable,
it must be routed through the explicit scale-drift alternatives of
\cref{paper6a:cor:scale-drift-alternatives,paper6a:thm:moving-cutoff-scale-collapse-alternative}
or treated by a direct summability estimate.
\end{remark}

\begin{lemma}[Componentwise estimates give annular finite-error selection]
\label{paper7:lem:componentwise-gives-annular-finite-error}
If every branch in the pre-carrier validated compact-cost channel whose
validated compact cost is energy-controlled by the fixed identity cost
satisfies the componentwise annular finite-error estimates of
\cref{paper7:def:componentwise-annular-finite-error}, then the channel has
annular finite-error selection in the sense of
\cref{paper7:def:annular-finite-error-selection}.
\end{lemma}

\begin{proof}
Let \(\omega\) be such a branch.  If the fixed alternative in
\cref{paper7:def:componentwise-annular-finite-error} holds, with
\(I_1,\ldots,I_6\) denoting the fixed components appearing in that definition,
then
\cref{paper6a:def:fixed-error-density} gives
\[
  B_{R_0}
  =
  \frac{\nu}{2}I_1+\frac12 I_2+I_3+\frac12 I_4+\frac12 I_5+\frac12 I_6 .
\]
Since \(B_{R_0}\) is a finite positive linear combination of integrable fixed
components, \(B_{R_0}\in L^1(\tau_0,\infty)\).  This is the first alternative
in \cref{paper7:def:annular-finite-error-selection}.
Otherwise the definition supplies a finite-valued nested moving radius
\(R(\tau)\ge R_0\), \(R(\tau)\to\infty\), for which the seven moving annular
components \(J_1,\ldots,J_7\) are integrable.  These seven components are
exactly the summability requirements in
\cref{paper6a:def:summable-moving-errors}, with the same moving cutoff
\(\Phi\).  Hence the moving-annulus errors are summable, which is the second
alternative in \cref{paper7:def:annular-finite-error-selection}.
\end{proof}

\begin{definition}[Fixed annular PDE estimate package]
\label{paper7:def:fixed-pde-error-package}
Let \(\phi_{R_0}\) be the fixed cutoff used in the validated compact cost and
set
\[
  A_\phi:=\operatorname{supp}(\nabla\phi_{R_0})\cup
  \operatorname{supp}(\Delta\phi_{R_0}).
\]
A branch in the pre-carrier validated compact-cost channel satisfies the fixed
annular PDE estimate package if, on a terminal tail,
\begin{enumerate}[label=\textup{(\roman*)}]
\item
\[
  \int_{\tau_0}^{\infty}\int_{A_\phi}|V(y,\tau)|^2\,dy\,d\tau<\infty;
\]
\item
\[
  \int_{\tau_0}^{\infty}\int_{A_\phi}|V(y,\tau)|^3\,dy\,d\tau<\infty;
\]
\item the fixed pressure flux is integrable:
\[
  \int_{\tau_0}^{\infty}
  \left|\int P(y,\tau)V(y,\tau)\cdot\nabla\phi_{R_0}(y)\,dy\right|d\tau<\infty;
\]
\item the fixed translation-annulus term is integrable:
\[
  \int_{\tau_0}^{\infty}
  |b(\tau)|\int_{A_\phi}|V(y,\tau)|^2\,dy\,d\tau<\infty;
\]
\item the fixed negative-scale term is integrable:
\[
  \int_{\tau_0}^{\infty}a_-(\tau)M_\phi(\tau)\,d\tau<\infty,
  \qquad
  M_\phi(\tau)=\int |V(y,\tau)|^2\phi_{R_0}(y)\,dy.
\]
\item the fixed positive scale-shell term is integrable:
\[
  \int_{\tau_0}^{\infty}
  a_+(\tau)\int |V(y,\tau)|^2\bigl(-y\cdot\nabla\phi_{R_0}(y)\bigr)\,dy\,d\tau
  <\infty.
\]
\end{enumerate}
\end{definition}

\begin{lemma}[Fixed PDE package gives fixed finite error]
\label{paper7:lem:fixed-pde-package-gives-components}
If a pre-carrier branch satisfies the fixed annular PDE estimate package of
\cref{paper7:def:fixed-pde-error-package}, then the fixed alternative in
\cref{paper7:def:componentwise-annular-finite-error} holds.
\end{lemma}

\begin{proof}
Since \(\phi_{R_0}\in C_c^\infty\), there are finite constants
\[
  C_{\Delta,\phi}=\|\Delta\phi_{R_0}\|_{L^\infty},\qquad
  C_{\nabla,\phi}=\|\nabla\phi_{R_0}\|_{L^\infty},\qquad
  C_{y\nabla,\phi}=\|y\cdot\nabla\phi_{R_0}\|_{L^\infty}.
\]
Because \(\nabla\phi_{R_0}\) and \(\Delta\phi_{R_0}\) vanish outside \(A_\phi\),
\[
  I_1
  =
  \left|\int |V|^2\Delta\phi_{R_0}\right|
  \le
  C_{\Delta,\phi}\int_{A_\phi}|V|^2,
\]
\[
  I_2
  =
  \left|\int |V|^2V\cdot\nabla\phi_{R_0}\right|
  \le
  C_{\nabla,\phi}\int_{A_\phi}|V|^3,
\]
\[
  I_4
  =
  |b|\left|\int |V|^2\nabla\phi_{R_0}\right|
  \le
  C_{\nabla,\phi}|b|\int_{A_\phi}|V|^2,
\]
and
\[
  I_6
  =
  a_+\int |V|^2\bigl(-y\cdot\nabla\phi_{R_0}\bigr)
  \le
  C_{y\nabla,\phi}a_+\int_{A_\phi}|V|^2.
\]
The pressure term \(I_3\) is integrable by item \textup{(iii)} of the fixed
package, \(I_5=a_-M_\phi\) is integrable by item \textup{(v)}, and \(I_6\) is
integrable by item \textup{(vi)}.  Therefore
all six fixed components in
\cref{paper7:def:componentwise-annular-finite-error} belong to
\(L^1(\tau_0,\infty)\), which is exactly the fixed alternative there.
\end{proof}

\begin{definition}[Moving annular PDE estimate package]
\label{paper7:def:moving-pde-error-package}
A branch in the pre-carrier validated compact-cost channel satisfies the
moving annular PDE estimate package if there is a finite-valued nondecreasing
locally absolutely continuous nested radius \(R(\tau)\ge R_0\), with
\(R(\tau)\to\infty\) and
\(\Phi(\tau,y)=\phi(y/R(\tau))\ge\phi_{R_0}(y)\), such that on a terminal tail:
\begin{enumerate}[label=\textup{(\roman*)}]
\item the moving derivative/convection annular packing holds:
\[
  \int_{\tau_0}^{\infty}
  \left(
  \frac{|\dot R(\tau)|}{R(\tau)}\int_{A_{R(\tau)}}|V|^2
  +\frac{1}{R(\tau)^2}\int_{A_{R(\tau)}}|V|^2
  +\frac{1}{R(\tau)}\int_{A_{R(\tau)}}|V|^3
  \right)d\tau<\infty;
\]
\item the moving pressure flux is integrable:
\[
  \int_{\tau_0}^{\infty}
  \left|\int P(y,\tau)V(y,\tau)\cdot\nabla\Phi(\tau,y)\,dy\right|d\tau<\infty;
\]
\item the moving translation-annulus packing holds:
\[
  \int_{\tau_0}^{\infty}
  \frac{|b(\tau)|}{R(\tau)}\int_{A_{R(\tau)}}|V|^2\,d\tau<\infty;
\]
\item the moving positive scale-shell term is integrable:
\[
  \int_{\tau_0}^{\infty}
  a_+(\tau)\int |V(y,\tau)|^2\bigl(-y\cdot\nabla\Phi(\tau,y)\bigr)\,dy\,d\tau
  <\infty;
\]
\item either the moving negative-scale term is integrable,
\[
  \int_{\tau_0}^{\infty}a_-(\tau)\int |V(y,\tau)|^2\Phi(\tau,y)\,dy\,d\tau<\infty,
\]
or the branch lies in the scale-collapse drift obstruction of
\cref{paper6a:def:negative-drift}.
\end{enumerate}
\end{definition}

\begin{remark}[Pressure item in the moving PDE package]
\label{paper7:rem:moving-pressure-package}
The pressure-flux clause in
\cref{paper7:def:moving-pde-error-package} is not a separate global pressure
hypothesis.  By
\cref{paper7:cor:moving-pressure-from-annular-l3}, it follows from enlarged-
shell \(L^3\) packing together with a summable \(R(\tau)^{-1}\) schedule.
Thus the genuinely new moving-route burden is concentrated in the annular
velocity packing, the translation term, the positive scale-shell term, and the
negative-drift routing.
\end{remark}

\begin{lemma}[Moving PDE package gives moving finite error or scale-drift routing]
\label{paper7:lem:moving-pde-package-gives-components}
If a pre-carrier branch satisfies the moving annular PDE estimate package of
\cref{paper7:def:moving-pde-error-package}, then either the moving alternative
in \cref{paper7:def:componentwise-annular-finite-error} holds or the branch
lies in the scale-collapse drift obstruction of
\cref{paper6a:def:negative-drift}.
\end{lemma}

\begin{proof}
Let
\[
  C_{\phi,1}=\|\nabla\phi\|_{L^\infty},\qquad
  C_{\phi,2}=\|\Delta\phi\|_{L^\infty},\qquad
  C_{\phi,3}=\sup_{1\le|z|\le2}|z||\nabla\phi(z)|.
\]
Because \(\phi\) is supported in \(B_2\) and constant on \(B_1\), the support
of \(\nabla\Phi\), \(\Delta\Phi\), and \(\partial_\tau\Phi\) is contained in
\(A_{R(\tau)}\).  Moreover
\[
  \partial_\tau\Phi(\tau,y)
  =
  -\frac{\dot R(\tau)}{R(\tau)}
  \frac{y}{R(\tau)}\cdot \nabla\phi(y/R(\tau)),
\]
so on \(A_{R(\tau)}\),
\[
  |\partial_\tau\Phi(\tau,y)|
  \le
  C_{\phi,3}\frac{|\dot R(\tau)|}{R(\tau)}.
\]
Likewise,
\[
  |\nabla\Phi(\tau,y)|
  \le C_{\phi,1}R(\tau)^{-1}\mathbf 1_{A_{R(\tau)}}(y),
  \qquad
  |\Delta\Phi(\tau,y)|
  \le C_{\phi,2}R(\tau)^{-2}\mathbf 1_{A_{R(\tau)}}(y),
\]
and
\[
  |y\cdot\nabla\Phi(\tau,y)|
  \le C_{\phi,3}\mathbf 1_{A_{R(\tau)}}(y).
\]
Therefore
\[
  J_1
  =
  \left|\int |V|^2\partial_\tau\Phi\right|
  \le
  C_{\phi,3}\frac{|\dot R|}{R}\int_{A_{R(\tau)}}|V|^2,
\]
\[
  J_2
  =
  \left|\int |V|^2\Delta\Phi\right|
  \le
  C_{\phi,2}R(\tau)^{-2}\int_{A_{R(\tau)}}|V|^2,
\]
\[
  J_3
  =
  \left|\int |V|^2V\cdot\nabla\Phi\right|
  \le
  C_{\phi,1}R(\tau)^{-1}\int_{A_{R(\tau)}}|V|^3,
\]
\[
  J_5
  =
  |b|\left|\int |V|^2\nabla\Phi\right|
  \le
  C_{\phi,1}\frac{|b|}{R(\tau)}\int_{A_{R(\tau)}}|V|^2,
\]
and
\[
  J_7
  =
  a_+\int |V|^2\bigl(-y\cdot\nabla\Phi\bigr)
  \le
  C_{\phi,3}a_+\int_{A_{R(\tau)}}|V|^2.
\]
Hence items \textup{(i)} and \textup{(iii)} of the moving package make
\(J_1,J_2,J_3,J_5\) integrable, and item \textup{(iv)} makes \(J_7\)
integrable.  Item \textup{(ii)} makes \(J_4\)
integrable.  If the first branch of item \textup{(v)} holds, then \(J_6\) is
integrable as well, so all seven moving components in
\cref{paper7:def:componentwise-annular-finite-error} belong to
\(L^1(\tau_0,\infty)\), which is exactly the moving alternative there.
Otherwise the second branch of item \textup{(v)} places the branch in the
scale-collapse drift obstruction.
\end{proof}

\begin{lemma}[Annular pressure-flux estimate from local shell mass]
\label{paper7:lem:annular-pressure-flux}
Let \(\phi\in C_c^\infty(\mathbb R^3)\) be radial and nonincreasing, with
\(\phi=1\) on \(B_1\) and \(\phi=0\) outside \(B_2\).  For \(R\ge1\), set
\[
  \Phi_R(y)=\phi(y/R),\qquad
  A_R=\{R\le|y|\le2R\},\qquad
  \widetilde A_R=\{R/4<|y|<8R\}.
\]
Assume \(P=\mathcal R_i\mathcal R_j(V_iV_j)\) modulo functions of time and
\[
  \|V(\tau)\|_{L^3(\mathbb R^3)}\le M
  \qquad\text{for a.e. }\tau\ge\tau_0 .
\]
Then for a.e. \(\tau\ge\tau_0\),
\[
  \left|\int P(y,\tau)V(y,\tau)\cdot\nabla\Phi_R(y)\,dy\right|
  \le
  C_\phi R^{-1}\int_{\widetilde A_R}|V(y,\tau)|^3\,dy
  +C_\phi M^3R^{-1}.
\]
\end{lemma}

\begin{proof}
Choose \(\chi_R\in C_c^\infty(\widetilde A_R)\), radial, \(0\le\chi_R\le1\),
with \(\chi_R\equiv1\) on \(A_R\).  Decompose
\[
  P=P_R^{\rm ann}+H_R,
  \qquad
  P_R^{\rm ann}:=\mathcal R_i\mathcal R_j(\chi_RV_iV_j),
\]
so that \(H_R:=P-P_R^{\rm ann}\) is harmonic on the neighborhood
\(\{R/2<|y|<4R\}\) of \(A_R\).

By Calderon--Zygmund boundedness on \(L^{3/2}\),
\[
  \|P_R^{\rm ann}(\tau)\|_{L^{3/2}(A_R)}
  \le
  C\|V(\tau)\|_{L^3(\widetilde A_R)}^2.
\]
Since \(\nabla\Phi_R\) is supported in \(A_R\) and
\(\|\nabla\Phi_R\|_{L^\infty}\le C_\phi R^{-1}\),
\[
\begin{aligned}
  \left|\int P_R^{\rm ann}V\cdot\nabla\Phi_R\right|
  &\le
  C_\phi R^{-1}
  \|P_R^{\rm ann}\|_{L^{3/2}(A_R)}
  \|V\|_{L^3(A_R)} \\
  &\le
  C_\phi R^{-1}\|V\|_{L^3(\widetilde A_R)}^3
  \le
  C_\phi R^{-1}\int_{\widetilde A_R}|V|^3 .
\end{aligned}
\]

For the harmonic part, let \(c_R(\tau)\) be the average of \(H_R(\cdot,\tau)\)
over \(A_R\).  Since constants disappear against \(V\cdot\nabla\Phi_R\), one
has
\[
  \int H_RV\cdot\nabla\Phi_R
  =
  \int (H_R-c_R)V\cdot\nabla\Phi_R .
\]
It therefore suffices to bound the oscillation of \(H_R\) on \(A_R\).

Fix \(y\in\{R/2<|y|<4R\}\).  Because \(1-\chi_R\) is supported in
\(\{|z|\le R/4\}\cup\{|z|\ge8R\}\), every such \(z\) satisfies
\(|y-z|\ge cR\).  Differentiating the Calderon--Zygmund kernel gives
\[
  |\nabla H_R(y,\tau)|
  \le
  C\int_{|z|\le R/4}|y-z|^{-4}|V(z,\tau)|^2\,dz
  +
  C\int_{|z|\ge8R}|y-z|^{-4}|V(z,\tau)|^2\,dz .
\]
For the inner part, Hölder on \(B_{R/4}\) yields
\[
  \int_{|z|\le R/4}|V(z,\tau)|^2\,dz
  \le
  |B_{R/4}|^{1/3}\|V(\tau)\|_{L^3(\mathbb R^3)}^2
  \le
  CRM^2,
\]
hence
\[
  \int_{|z|\le R/4}|y-z|^{-4}|V(z,\tau)|^2\,dz
  \le
  CR^{-4}(CRM^2)
  =
  CR^{-3}M^2.
\]
For the outer part, \(|y-z|\simeq|z|\) and Hölder give
\[
\begin{aligned}
  \int_{|z|\ge8R}|y-z|^{-4}|V(z,\tau)|^2\,dz
  &\le
  C\int_{|z|\ge8R}|z|^{-4}|V(z,\tau)|^2\,dz \\
  &\le
  C\|V(\tau)\|_{L^3(\mathbb R^3)}^2
  \left(\int_{|z|\ge8R}|z|^{-12}\,dz\right)^{1/3} \\
  &\le
  CR^{-3}M^2.
\end{aligned}
\]
Therefore
\[
  \|\nabla H_R(\tau)\|_{L^\infty(\{R/2<|y|<4R\})}
  \le
  CR^{-3}M^2.
\]
Since \(A_R\) has diameter \(O(R)\), the mean-value theorem gives
\[
  \|H_R(\tau)-c_R(\tau)\|_{L^\infty(A_R)}
  \le
  CR^{-2}M^2.
\]
Using again that \(\nabla\Phi_R\) is supported in \(A_R\),
\[
\begin{aligned}
  \left|\int (H_R-c_R)V\cdot\nabla\Phi_R\right|
  &\le
  C_\phi R^{-1}
  \|H_R-c_R\|_{L^\infty(A_R)}
  \|V\|_{L^1(A_R)} \\
  &\le
  C_\phi R^{-1}(CR^{-2}M^2)|A_R|^{2/3}\|V\|_{L^3(A_R)} \\
  &\le
  C_\phi M^2R^{-1}\|V\|_{L^3(\mathbb R^3)}
  \le
  C_\phi M^3R^{-1}.
\end{aligned}
\]
Adding the annular-local and harmonic-tail bounds proves the claim.
\end{proof}

\begin{corollary}[Moving pressure flux from enlarged-shell \(L^3\) packing]
\label{paper7:cor:moving-pressure-from-annular-l3}
Let \(R(\tau)\ge1\) be measurable and set
\[
  \widetilde A_{R(\tau)}=\{R(\tau)/4<|y|<8R(\tau)\},
  \qquad
  \Phi(\tau,y)=\phi(y/R(\tau)).
\]
Assume
\[
  \|V(\tau)\|_{L^3(\mathbb R^3)}\le M
  \qquad\text{for a.e. }\tau\ge\tau_0,
\]
and
\[
  \int_{\tau_0}^{\infty}
  \frac{1}{R(\tau)}
  \int_{\widetilde A_{R(\tau)}}|V(y,\tau)|^3\,dy\,d\tau<\infty,
\]
\[
  \int_{\tau_0}^{\infty}\frac{d\tau}{R(\tau)}<\infty.
\]
Then
\[
  \int_{\tau_0}^{\infty}
  \left|\int P(y,\tau)V(y,\tau)\cdot\nabla\Phi(\tau,y)\,dy\right|d\tau<\infty.
\]
\end{corollary}

\begin{proof}
Apply \cref{paper7:lem:annular-pressure-flux} at almost every \(\tau\) with
radius \(R(\tau)\).  Integrating the resulting pointwise bound gives
\[
\begin{aligned}
  &\int_{\tau_0}^{\infty}
  \left|\int P\,V\cdot\nabla\Phi\right|d\tau \\
  &\le
  C_\phi
  \int_{\tau_0}^{\infty}
  \frac{1}{R(\tau)}
  \int_{\widetilde A_{R(\tau)}}|V|^3\,dy\,d\tau
  +
  C_\phi M^3
  \int_{\tau_0}^{\infty}\frac{d\tau}{R(\tau)},
\end{aligned}
\]
and both terms on the right are finite by hypothesis.
\end{proof}

\begin{theorem}[Carrier closure reduced to explicit PDE estimates and scale-drift routing]
\label{paper7:thm:carrier-closure-reduced-to-pde-estimates}
Assume repaired-gauge suitable local energy gives the fixed and moving carrier
subidentities in the sense of
\cref{paper7:lem:suitable-energy-gives-carrier-subidentities}.  Let \(\omega\)
be a branch in the pre-carrier validated compact-cost channel whose validated
compact cost is energy-controlled by the fixed identity cost.  Assume that
\(\omega\) satisfies either the fixed annular PDE estimate package of
\cref{paper7:def:fixed-pde-error-package} or the moving annular PDE estimate
package of \cref{paper7:def:moving-pde-error-package}.  Then exactly one of the
following occurs:
\begin{enumerate}[label=\textup{(\roman*)}]
\item \(\omega\) is excluded by the fixed-cutoff carrier closure theorem
\cref{paper7:thm:fixed-carrier-closure};
\item \(\omega\) is excluded by the moving-cutoff carrier closure theorem
\cref{paper7:thm:moving-carrier-closure};
\item \(\omega\) enters the scale-collapse drift obstruction of
\cref{paper6a:def:negative-drift}.
\end{enumerate}
\end{theorem}

\begin{proof}
If the fixed annular PDE estimate package holds, then
\cref{paper7:lem:fixed-pde-package-gives-components} gives the fixed
alternative in \cref{paper7:def:componentwise-annular-finite-error}.  Writing
\(I_1,\ldots,I_6\) for those six fixed components, the identity
\[
  B_{R_0}
  =
  \frac{\nu}{2}I_1+\frac12 I_2+I_3+\frac12 I_4+\frac12 I_5+\frac12 I_6
\]
from \cref{paper6a:def:fixed-error-density} yields
\(B_{R_0}\in L^1(\tau_0,\infty)\).  The suitable local energy lemma gives the
fixed carrier subidentity.  Finite initial localized energy and carrier
comparison are automatic on the pre-carrier channel by
\cref{paper7:lem:precarrier-finite-local-energy} and the assumed
energy-control of the validated cost by the fixed identity cost.
Thus the fixed-cutoff carrier closure package holds, and
\cref{paper7:thm:fixed-carrier-closure} excludes \(\omega\).
\par
Suppose instead that the moving annular PDE estimate package holds.  By
\cref{paper7:lem:moving-pde-package-gives-components}, either the branch lies
in the scale-collapse drift obstruction, giving alternative \textup{(iii)}, or
the moving alternative in
\cref{paper7:def:componentwise-annular-finite-error} holds.  In the latter
case, the seven moving components \(J_1,\ldots,J_7\) are integrable, which is
exactly the summable moving-annulus condition of
\cref{paper6a:def:summable-moving-errors}.  The suitable local energy lemma gives the moving carrier
subidentity.  Finite initial moving energy and moving carrier comparison are
automatic on the pre-carrier channel by
\cref{paper7:lem:precarrier-finite-local-energy,paper7:lem:nested-moving-carrier-comparison}.
Therefore the moving-cutoff carrier closure package holds, and
\cref{paper7:thm:moving-carrier-closure} excludes \(\omega\).  The three
alternatives are disjoint by construction.
\end{proof}

\begin{corollary}[Explicit PDE packages leave no unresolved compact-cost branch]
\label{paper7:cor:adjacent-closure-reduced-to-pde-estimates}
Assume the non-carrier adjacent strata listed in
\cref{paper7:thm:noncarrier-exits-route} are closed by their corresponding
local state-space or interface arguments in the ambient Type II class.  Assume
also that any branch entering the scale-collapse drift obstruction is closed by
its assigned Section~6 local argument, for example through
\cref{paper6a:thm:moving-cutoff-scale-collapse-alternative}.  If every branch
in the pre-carrier validated compact-cost channel whose validated compact cost
is energy-controlled by the fixed identity cost satisfies either the fixed
annular PDE estimate package of
\cref{paper7:def:fixed-pde-error-package} or the moving annular PDE estimate
package of \cref{paper7:def:moving-pde-error-package}, then every branch with
active supercritical scaling and validated compact Type II cost divergence is
either excluded on the retained compact stratum, excluded by fixed or moving
carrier closure, or routed to a non-carrier or scale-drift branch that is
already closed in the ambient class.  In particular, the compact-cost channel
leaves no unresolved branch.
\end{corollary}

\begin{proof}
Let \(\omega\) be a branch with active supercritical scaling and validated
compact Type II cost divergence.  By
\cref{paper7:thm:classwise-retained-compact-coverage}, either \(\omega\) lies
in the retained compact stratum or it enters one of the named adjacent strata.
The retained compact stratum is excluded by
\cref{paper7:cor:retained-stratum-cost-exclusion}.  If \(\omega\) enters a
non-carrier adjacent stratum, the first assumption closes it.  The only
remaining case is that \(\omega\) lies in the pre-carrier validated
compact-cost channel.  Then the explicit PDE-package hypothesis and
\cref{paper7:thm:carrier-closure-reduced-to-pde-estimates} show that
\(\omega\) is either excluded by fixed carrier closure, excluded by moving
carrier closure, or routed to the scale-collapse drift obstruction, which is
closed by the second assumption.  Hence every branch is either excluded on the
retained stratum or handled in a closed non-carrier or scale-drift branch.
This is exactly the claimed no-unresolved-branch conclusion.
\end{proof}

\begin{remark}[The explicit PDE packages are not the remaining target]
\label{paper7:rem:pde-packages-optional}
The explicit PDE packages of
\cref{paper7:def:fixed-pde-error-package,paper7:def:moving-pde-error-package}
give a sufficient analytic realization of carrier closure, but they are not the
remaining discharge target.  The proof architecture does not require one global
tail estimate that simultaneously absorbs every carrier error.  Instead,
persistent failure of a moving-carrier ingredient is supposed to force entry
into a neighboring state-space stratum.  The remainder of this section
therefore isolates only the small routing statements needed for that
state-space discharge.
\end{remark}

\begin{theorem}[Finite-error selection reduced to annular finite-error]
\label{paper7:thm:finite-error-selection-reduced-to-annular}
Assume that repaired-gauge suitable local energy gives the fixed and moving
carrier subidentities in the sense of
\cref{paper7:lem:suitable-energy-gives-carrier-subidentities}.  If the
pre-carrier channel has annular finite-error selection in the sense of
\cref{paper7:def:annular-finite-error-selection}, then it has finite-error
carrier selection in the sense of
\cref{paper7:def:finite-error-carrier-selection}.
\end{theorem}

\begin{proof}
Let \(\omega\) be a branch in the pre-carrier channel whose validated compact
cost is energy-controlled by the fixed identity cost.  By annular finite-error
selection, after restricting to a later tail either the fixed error density is
integrable or there is a nested moving radius for which the moving-annulus
errors are summable.  The repaired-gauge suitable local energy lemma supplies
the fixed carrier subidentity in the first case and the moving carrier
subidentity in the second.  These are exactly the two alternatives in
\cref{paper7:def:finite-error-carrier-selection}.
\end{proof}

\begin{theorem}[Carrier completion reduced to finite-error selection]
\label{paper7:thm:carrier-completion-reduced-to-finite-error}
Assume the pre-carrier validated compact-cost channel has finite-error carrier
selection in the sense of
\cref{paper7:def:finite-error-carrier-selection}.  Then every branch in the
pre-carrier channel whose validated compact cost is energy-controlled by the
fixed identity cost satisfies the fixed-cutoff or moving-cutoff carrier closure
package.  Consequently that branch is excluded from the Type II class under
consideration.
\end{theorem}

\begin{proof}
Let \(\omega\) be such a branch.  The positive finite critical \(L^3\) annulus
in the pre-carrier channel gives the finite initial localized energy required
by either carrier package through
\cref{paper7:lem:precarrier-finite-local-energy}.  The assumed energy-control
comparison gives the fixed carrier-comparison clause.
\par
If the fixed alternative in
\cref{paper7:def:finite-error-carrier-selection} holds, then the fixed
carrier subidentity and the finite \(B_{R_0}\)-tail are precisely the two
remaining clauses of the fixed-cutoff carrier closure package.  Hence
\cref{paper7:thm:fixed-carrier-closure} excludes the branch.
\par
If the moving alternative holds, the summable moving-annulus condition gives
\[
  \int_{\tau_0}^{\infty}B_{\mathrm{mov}}(\tau)\,d\tau<\infty
\]
by \cref{paper6a:lem:summable-moving-errors}.  The moving carrier subidentity
is part of the moving alternative.  The nested cutoff condition and fixed
energy-control comparison give the moving carrier comparison by
\cref{paper7:lem:nested-moving-carrier-comparison}.  Together with finite
initial moving energy from
\cref{paper7:lem:precarrier-finite-local-energy}, these are exactly the moving
carrier closure package.  Therefore
\cref{paper7:thm:moving-carrier-closure} excludes the branch.
\par
The two alternatives in finite-error carrier selection exhaust the branch, so
the claimed fixed-or-moving carrier completion and exclusion follow.
\end{proof}

\begin{theorem}[Carrier package closes the pre-carrier channel]
\label{paper7:thm:carrier-package-closes-precarrier}
Let \(\omega\) be in the pre-carrier validated compact-cost channel.  If
\(\omega\) satisfies either the fixed-cutoff carrier closure package of
\cref{paper7:def:fixed-carrier-closure-package} or the moving-cutoff carrier
closure package of \cref{paper7:def:moving-carrier-closure-package}, then
\(\omega\) is excluded from the Type II class under consideration.
\end{theorem}

\begin{proof}
In the fixed-cutoff case apply
\cref{paper7:thm:fixed-carrier-closure}.  In the moving-cutoff case apply
\cref{paper7:thm:moving-carrier-closure}.  These are the only cases in the
stated disjunction.
\end{proof}

\begin{definition}[Carrier-estimate completion for the compact-cost channel]
\label{paper7:def:carrier-estimate-completion}
The compact-cost channel has carrier-estimate completion if every branch in the
pre-carrier validated compact-cost channel satisfies, after passing to a later
terminal tail if necessary, either the fixed-cutoff carrier closure package of
\cref{paper7:def:fixed-carrier-closure-package} or the moving-cutoff carrier
closure package of \cref{paper7:def:moving-carrier-closure-package}.
\end{definition}

\begin{corollary}[Finite-error selection gives carrier-estimate completion]
\label{paper7:cor:finite-error-selection-gives-carrier-completion}
Assume every branch in the pre-carrier validated compact-cost channel uses a
validated compact cost that is energy-controlled by the fixed identity cost; in
particular this holds for the canonical identity cost.  If finite-error carrier
selection holds in the sense of
\cref{paper7:def:finite-error-carrier-selection}, then the compact-cost channel
has carrier-estimate completion in the sense of
\cref{paper7:def:carrier-estimate-completion}.
\end{corollary}

\begin{proof}
Let \(\omega\) be any branch in the pre-carrier validated compact-cost channel.
By the energy-control hypothesis, \(\omega\) is one of the branches covered by
\cref{paper7:thm:carrier-completion-reduced-to-finite-error}.  That theorem
gives either the fixed-cutoff carrier closure package or the moving-cutoff
carrier closure package on a terminal tail.  This is exactly
\cref{paper7:def:carrier-estimate-completion}.  The parenthetical canonical
case follows from
\cref{paper7:def:energy-controlled-validated-cost}.
\end{proof}

\begin{corollary}[Annular finite-error selection gives carrier-estimate completion]
\label{paper7:cor:annular-finite-error-selection-gives-carrier-completion}
Assume every branch in the pre-carrier validated compact-cost channel uses a
validated compact cost that is energy-controlled by the fixed identity cost.
Assume also that repaired-gauge suitable local energy gives the fixed and
moving carrier subidentities in the sense of
\cref{paper7:lem:suitable-energy-gives-carrier-subidentities}.  If annular
finite-error selection holds in the sense of
\cref{paper7:def:annular-finite-error-selection}, then the compact-cost channel
has carrier-estimate completion.
\end{corollary}

\begin{proof}
\Cref{paper7:thm:finite-error-selection-reduced-to-annular} converts annular
finite-error selection into finite-error carrier selection.  Then
\cref{paper7:cor:finite-error-selection-gives-carrier-completion} gives
carrier-estimate completion.
\end{proof}

\begin{theorem}[Adjacent closure reduced to carrier-estimate completion]
\label{paper7:thm:adjacent-closure-reduced-to-carrier-completion}
Assume the non-carrier adjacent strata listed in
\cref{paper7:thm:noncarrier-exits-route} are closed by their corresponding
local state-space or interface arguments in the ambient Type II class.  If the
compact-cost channel has carrier-estimate completion in the sense of
\cref{paper7:def:carrier-estimate-completion}, then the compact-cost channel
has adjacent-stratum closure in the sense of
\cref{paper7:def:compact-cost-adjacent-closure}.
\end{theorem}

\begin{proof}
Let \(\omega\) be a branch with active supercritical scaling and validated
compact Type II cost divergence that does not remain in the retained compact
cost-divergence stratum.  By
\cref{paper7:thm:classwise-retained-compact-coverage}, \(\omega\) enters one
of the named adjacent strata produced by
\cref{paper7:lem:no-unrecorded-cost-escape}.
\par
If the adjacent stratum is non-carrier, then
\cref{paper7:thm:noncarrier-exits-route} identifies it with one of the
previously named exhaustion, representation, state-space, critical-mass,
compact-test, modulation, cost-well-posedness, or cost-compatibility strata.
By the first hypothesis of the present theorem, that stratum is closed by its
own local argument, so the Type II exclusion conclusion holds there, or the
branch is routed to a previously closed non-compact-cost alternative.
\par
It remains to consider a branch that reaches the pre-carrier channel.  By
carrier-estimate completion, after restriction to a later tail if needed, it
satisfies either the fixed-cutoff or moving-cutoff carrier closure package.  In
the fixed or moving case,
\cref{paper7:thm:carrier-package-closes-precarrier} gives the Type II
exclusion conclusion.  Thus every named adjacent stratum arising from the
compact-cost channel has a local exclusion or a route to a previously closed
alternative, which is exactly
\cref{paper7:def:compact-cost-adjacent-closure}.
\end{proof}

\begin{definition}[Stratified moving-carrier routing package]
\label{paper7:def:stratified-moving-carrier-routing}
Assume the validated compact cost is energy-controlled by the fixed identity
cost.  The pre-carrier validated compact-cost channel has stratified
moving-carrier routing if every branch \(\omega\) in that channel, after
passing to a later terminal tail if necessary, satisfies at least one of the
following local alternatives:
\begin{enumerate}[label=\textup{(\roman*)}]
\item \(\omega\) satisfies the fixed-cutoff carrier closure package of
\cref{paper7:def:fixed-carrier-closure-package};
\item \(\omega\) satisfies the moving-cutoff carrier closure package of
\cref{paper7:def:moving-carrier-closure-package};
\item \(\omega\) satisfies the exterior/leakage routing alternative: either
there are unit windows \(I_n=[T_n,T_n+1]\), selected times \(\tau_n\in I_n\),
radii \(R_n\to\infty\), and \(\eta>0\) such that
\[
  \sup_n\bigl(|a(\tau_n)|+\|b(\tau_n)\|_\infty\bigr)<\infty,
\]
\[
  \frac{1}{R_n^2}\int_{A_{R_n}}|V(y,\tau_n)|^2\,dy
  +
  \frac{1}{R_n}\int_{A_{R_n}}|V(y,\tau_n)|^3\,dy
  +
  \frac{1}{R_n}\int_{\widetilde A_{R_n}}|V(y,\tau_n)|^3\,dy
  \ge \eta
\]
for every \(n\), and this selected-window shell persistence is assigned to one
of the already named critical-tightness failure, exterior concentration, or
local windowed \(H^1\)-loss branches; or the normalized transition-carrier
blocks of the moving cutoff have a local carrier-state limit whose first
state-space label is one of those same three exterior/leakage branches.
\item \(\omega\) satisfies the translation-routing alternative: there are unit
windows \(I_n=[T_n,T_n+1]\), selected times \(\tau_n\in I_n\), radii
\(R_n\to\infty\), cutoffs \(\Phi_n(y)=\phi(y/R_n)\), and \(\eta>0\) such that
\[
  \sup_n\bigl(|a(\tau_n)|+\|b(\tau_n)\|_\infty\bigr)<\infty,
\]
\[
  \frac{|b(\tau_n)|}{R_n}\int_{A_{R_n}}|V(y,\tau_n)|^2\,dy\ge\eta
\]
for every \(n\), and this persistent translation transport is assigned to the
already named modulation-integrability failure or to a multicore/exterior
branch produced after selected-window recentering;
\item \(\omega\) satisfies the positive-scale routing alternative: either
there are unit windows \(I_n=[T_n,T_n+1]\), selected times
\(\tau_n\in I_n\), radii \(R_n\to\infty\), cutoffs
\(\Phi_n(y)=\phi(y/R_n)\), and \(\eta>0\) such that
\[
  \sup_n\bigl(|a(\tau_n)|+\|b(\tau_n)\|_\infty\bigr)<\infty,
\]
\[
  a_+(\tau_n)\int |V(y,\tau_n)|^2
  \bigl(-y\cdot\nabla\Phi_n(y)\bigr)\,dy\ge\eta
\]
for every \(n\), and this persistent positive scale forcing is assigned to the
already named scale-rigid residual alternative or to a bounded selected-window
limit refinement of that alternative; or the normalized positive-scale-work
carrier blocks have a local carrier-state limit whose first state-space label
is the scale-rigid residual branch or its bounded selected-window-limit
refinement.
\item \(\omega\) enters the negative scale-drift obstruction of
\cref{paper6a:def:negative-drift}.
\end{enumerate}
The fixed-cutoff route is auxiliary here.  The canonical local route is the
moving one: either a nested moving cutoff closes the carrier, or a persistent
failure produces one of the displayed selected-window diagnostics and is routed
to the corresponding adjacent stratum.
\end{definition}

\begin{remark}[How the moving diagnostics match the carrier components]
\label{paper7:rem:stratified-moving-routing}
The exterior/leakage alternative in
\cref{paper7:def:stratified-moving-carrier-routing} packages the moving cutoff,
viscous, convection, and pressure diagnostics.  The point is that
\cref{paper7:lem:annular-pressure-flux} already converts pressure failure into
enlarged-shell \(L^3\) persistence plus an \(R^{-1}\) tail, and the
\(R^{-1}\)-tail can be made harmless on discrete selected radii.  Thus the
pressure obstruction is not an independent global theorem; it is another way
of detecting shell mass that refuses to leave the selected windows.  The
translation alternative packages the \(J_5\) term and is meant to be combined
with the bounded selected-window modulation provided by
\cref{paper7:cor:unit-window-selection,paper7:cor:selected-window-modulation}.
The positive-scale alternative packages the \(J_7\) term and isolates the only
genuinely positive scale-dynamics obstruction left after the favorable
\(a_-\)-part has been separated.  The negative-drift term \(J_6\) is already
assigned to Section~6 and is not to be replaced by an unproved compact-cost
equivalence.  The carrier-state clauses are included so that a diffuse
nonintegrable error is not incorrectly converted into a pointwise lower bound;
it is instead read as an occupation state in the local ordered stratification.
\end{remark}

\begin{definition}[Moving test family]
\label{paper7:def:moving-test-family}
For a branch in the pre-carrier validated compact-cost channel, a moving test
family is a finite-valued nondecreasing locally absolutely continuous radius
\(R(\tau)\ge R_0\), with \(R(\tau)\to\infty\), together with the nested moving
cutoff
\[
  \Phi(\tau,y)=\phi(y/R(\tau))\ge\phi_{R_0}(y).
\]
For such a family, let \(J_1,\ldots,J_7\) denote the seven moving carrier
components of \cref{paper7:def:componentwise-annular-finite-error}.
\end{definition}

\begin{lemma}[A successful moving test family gives moving carrier closure]
\label{paper7:lem:successful-moving-test-family}
Assume repaired-gauge suitable local energy gives the moving carrier
subidentity in the sense of
\cref{paper7:lem:suitable-energy-gives-carrier-subidentities}.  Let \(\omega\)
be a branch in the pre-carrier validated compact-cost channel whose validated
compact cost is energy-controlled by the fixed identity cost.  If there exists
a moving test family in the sense of \cref{paper7:def:moving-test-family} such
that, for some \(\tau_1\ge\tau_0\),
\[
  J_k\in L^1(\tau_1,\infty)
  \qquad\text{for each }1\le k\le7,
\]
then \(\omega\) satisfies the moving-cutoff carrier closure package of
\cref{paper7:def:moving-carrier-closure-package} after restriction to the later
tail \([\tau_1,\infty)\).  In particular, \(\omega\) is excluded by
\cref{paper7:thm:moving-carrier-closure}.
\end{lemma}

\begin{proof}
Restrict the branch and the moving test family to the tail \([\tau_1,\infty)\).
Let \(R(\tau)\) and \(\Phi(\tau,y)\) be the restricted moving test family, with
components \(J_1,\ldots,J_7\).  By
\cref{paper6a:def:summable-moving-errors}, the integrability of these seven
components is exactly the summable moving-annulus condition for \(\Phi\).
Hence \(\int_{\tau_1}^\infty B_{\mathrm{mov}}(\tau)\,d\tau<\infty\) by
\cref{paper6a:lem:summable-moving-errors}.  The suitable local energy lemma
gives the moving carrier subidentity for the same cutoff.  Finite initial
localized energy on the pre-carrier channel is automatic by
\cref{paper7:lem:precarrier-finite-local-energy}, and the moving carrier
comparison is automatic by
\cref{paper7:lem:nested-moving-carrier-comparison}.  These are exactly the
clauses of \cref{paper7:def:moving-carrier-closure-package}.  Therefore
\cref{paper7:thm:moving-carrier-closure} excludes \(\omega\).
\end{proof}

\begin{lemma}[Critical tightness gives a tailored windowwise radius schedule]
\label{paper7:lem:windowwise-tight-radius-schedule}
Let \(\omega\) be a branch in the pre-carrier validated compact-cost channel,
represented by \(V\) on a terminal tail.  For the unit windows
\[
  I_n=[\tau_0+n,\tau_0+n+1],\qquad n\in\mathbb N,
\]
set
\[
  A_n:=1+\int_{I_n}\bigl(|a(\tau)|+\|b(\tau)\|_\infty\bigr)\,d\tau .
\]
Then there exists a nondecreasing sequence \(R_n\to\infty\), with
\[
  R_n\ge 4\cdot 2^n,
\]
such that
\[
  \sup_{\tau\in I_n}\int_{|y|>R_n/4}|V(y,\tau)|^3\,dy
  \le
  2^{-3n}A_n^{-3/2}
  \qquad\text{for every }n.
\]
\end{lemma}

\begin{proof}
Since \(\omega\) lies in the pre-carrier validated compact-cost channel, the
critical-tightness test and the modulation-integrability test in
\cref{paper7:lem:no-unrecorded-cost-escape} have already passed.  Thus
\cref{paper7:def:critical-tightness} holds and
\cref{paper7:lem:local-modulation-integrability} gives
\[
  a,b\in L^1_{\rm loc}([\tau_0,\infty)).
\]
In particular each \(A_n\) is finite.
\par
For every \(n\), set
\[
  \varepsilon_n:=2^{-3n}A_n^{-3/2}>0.
\]
By uniform critical tightness, for each \(\varepsilon_n\) there exists
\(S_n<\infty\) such that
\[
  \sup_{\tau\ge\tau_0}\int_{|y|>S_n}|V(y,\tau)|^3\,dy\le \varepsilon_n .
\]
Choose \(R_n\) recursively so that
\[
  R_n\ge \max\{4S_n,\;4\cdot 2^n,\;R_{n-1}\}
\]
for \(n\ge1\), and \(R_0\ge \max\{4S_0,4\}\).  Then \(R_n\) is nondecreasing,
\(R_n\to\infty\), \(R_n\ge 4\cdot2^n\), and
\[
  \sup_{\tau\in I_n}\int_{|y|>R_n/4}|V(y,\tau)|^3\,dy
  \le
  \sup_{\tau\ge\tau_0}\int_{|y|>S_n}|V(y,\tau)|^3\,dy
  \le \varepsilon_n ,
\]
which is the desired estimate.
\end{proof}

\begin{lemma}[Canonical windowwise moving family controls shell, pressure, and translation terms]
\label{paper7:lem:windowwise-shell-pressure-translation}
Let \(\omega\) be a branch in the pre-carrier validated compact-cost channel.
Let \(R_n\) be the sequence supplied by
\cref{paper7:lem:windowwise-tight-radius-schedule}.  On the unit windows
\(I_n=[\tau_0+n,\tau_0+n+1]\), define the canonical interpolated radius
\(R(\tau)\) by
\[
  R(\tau)
  :=
  R_n+(\tau-(\tau_0+n))(R_{n+1}-R_n),
  \qquad \tau\in I_n,
\]
and the associated moving cutoff
\[
  \Phi(\tau,y):=\phi(y/R(\tau)).
\]
Write
\[
\begin{gathered}
  J_2(\tau)=\left|\int |V|^2\Delta\Phi(\tau,\cdot)\right|,\qquad
  J_3(\tau)=\left|\int |V|^2V\cdot\nabla\Phi(\tau,\cdot)\right|,\\
  J_4(\tau)=\left|\int P\,V\cdot\nabla\Phi(\tau,\cdot)\right|,\qquad
  J_5(\tau)=|b(\tau)|\left|\int |V|^2\nabla\Phi(\tau,\cdot)\right|.
\end{gathered}
\]
Then
\[
  \int_{\tau_0}^{\infty}
  \bigl(J_2(\tau)+J_3(\tau)+J_4(\tau)+J_5(\tau)\bigr)\,d\tau
  <\infty.
\]
\end{lemma}

\begin{proof}
The sequence \(R_n\) is nondecreasing and tends to \(+\infty\).  Hence the
piecewise affine interpolation \(R(\tau)\) is finite-valued, nondecreasing,
locally absolutely continuous, and satisfies \(R(\tau)\to\infty\).  Thus it is
an admissible moving test family in the sense of
\cref{paper7:def:moving-test-family}.
\par
Since \(\omega\) is in the pre-carrier channel, the positive finite critical
annulus has already passed.  Hence
\[
  \|V(\tau)\|_{L^3(\mathbb R^3)}\le M
  \qquad\text{for all }\tau\ge\tau_0
\]
for some \(M<\infty\).  Also, by
\cref{paper7:lem:windowwise-tight-radius-schedule},
\[
  \sup_{\tau\in I_n}\int_{|y|>R_n/4}|V(y,\tau)|^3\,dy
  \le \varepsilon_n:=2^{-3n}A_n^{-3/2},
\]
where
\[
  A_n=1+\int_{I_n}\bigl(|a(\tau)|+\|b(\tau)\|_\infty\bigr)\,d\tau.
\]
For \(\tau\in I_n\), one has \(R(\tau)\ge R_n\), hence
\[
  A_{R(\tau)}\cup\widetilde A_{R(\tau)}
  \subset \{|y|>R(\tau)/4\}\subset \{|y|>R_n/4\}.
\]
Therefore, for every \(\tau\in I_n\),
\[
  \int_{A_{R(\tau)}}|V(y,\tau)|^3\,dy
  +
  \int_{\widetilde A_{R(\tau)}}|V(y,\tau)|^3\,dy
  \le \varepsilon_n .
\]
Using Hölder on the annulus \(A_{R(\tau)}\),
\[
  \int_{A_{R(\tau)}}|V(y,\tau)|^2\,dy
  \le
  |A_{R(\tau)}|^{1/3}
  \left(\int_{A_{R(\tau)}}|V(y,\tau)|^3\,dy\right)^{2/3}
  \le
  C_\phi R(\tau)\varepsilon_n^{2/3}.
\]
\par
Since \(|\Delta\Phi(\tau,\cdot)|\le C_\phi R(\tau)^{-2}\mathbf 1_{A_{R(\tau)}}\),
\[
  J_2(\tau)
  \le
  C_\phi R(\tau)^{-2}\int_{A_{R(\tau)}}|V|^2
  \le
  C_\phi R(\tau)^{-1}\varepsilon_n^{2/3}
  \le
  C_\phi R_n^{-1}\varepsilon_n^{2/3}.
\]
Integrating on \(I_n\) and using \(R_n\ge 4\cdot 2^n\),
\[
  \int_{I_n}J_2(\tau)\,d\tau
  \le
  C_\phi R_n^{-1}\varepsilon_n^{2/3}
  \le
  C_\phi 2^{-3n}2^{-2n}A_n^{-1}.
\]
\par
Since \(|\nabla\Phi(\tau,\cdot)|\le C_\phi R(\tau)^{-1}\mathbf 1_{A_{R(\tau)}}\),
\[
  J_3(\tau)
  \le
  C_\phi R(\tau)^{-1}\int_{A_{R(\tau)}}|V|^3
  \le
  C_\phi R_n^{-1}\varepsilon_n,
\]
and therefore
\[
  \int_{I_n}J_3(\tau)\,d\tau
  \le
  C_\phi R_n^{-1}\varepsilon_n
  =
  C_\phi R_n^{-1}2^{-3n}A_n^{-3/2}.
\]
\par
Applying \cref{paper7:lem:annular-pressure-flux} with the radius \(R(\tau)\)
for almost every \(\tau\in I_n\),
\[
  J_4(\tau)
  \le
  C_\phi R(\tau)^{-1}\int_{\widetilde A_{R(\tau)}}|V|^3
  +C_\phi M^3R(\tau)^{-1}
  \le
  C_\phi R_n^{-1}\varepsilon_n+C_\phi M^3R_n^{-1}.
\]
Since \(|I_n|=1\) and \(R_n^{-1}\le 2^{-n-2}\),
\[
  \int_{I_n}J_4(\tau)\,d\tau
  \le
  C_\phi R_n^{-1}\varepsilon_n
  +C_\phi M^3 2^{-n}
  \le
  C_\phi 2^{-n-2}2^{-3n}A_n^{-3/2}+C_\phi M^3 2^{-n}.
\]
\par
Finally,
\[
  J_5(\tau)
  \le
  C_\phi \|b(\tau)\|_\infty R(\tau)^{-1}\int_{A_{R(\tau)}}|V|^2
  \le
  C_\phi \|b(\tau)\|_\infty \varepsilon_n^{2/3}.
\]
Hence
\[
  \int_{I_n}J_5(\tau)\,d\tau
  \le
  C_\phi \varepsilon_n^{2/3}\int_{I_n}\|b(\tau)\|_\infty\,d\tau
  \le
  C_\phi 2^{-2n}A_n^{-1}(A_n-1)
  \le C_\phi 2^{-2n}.
\]
The four resulting series are summable in \(n\), proving the claim.
\end{proof}

\begin{corollary}[Only transition and scale terms remain on the windowwise schedule]
\label{paper7:cor:windowwise-only-transition-and-scale-remain}
Under the hypotheses of
\cref{paper7:lem:windowwise-shell-pressure-translation}, the canonical
windowwise moving family
\[
  R(\tau)
  =
  R_n+(\tau-(\tau_0+n))(R_{n+1}-R_n),
  \qquad \tau\in I_n,
\]
has summable moving carrier components \(J_2,J_3,J_4,J_5\).  Thus the only
unresolved moving-carrier terms on that genuine moving family are the
transition term \(J_1\), the positive scale-shell term \(J_7\), and the
negative scale-drift term \(J_6\).
\end{corollary}

\begin{proof}
This is exactly the conclusion of
\cref{paper7:lem:windowwise-shell-pressure-translation}, together with the fact
that the moving carrier has exactly seven components, of which
\cref{paper7:lem:windowwise-shell-pressure-translation} already controls
\(J_2,J_3,J_4,J_5\).
\end{proof}

\begin{definition}[Windowwise transition/leakage routing statement]
\label{paper7:def:windowwise-transition-routing}
The pre-carrier validated compact-cost channel has windowwise
transition/leakage routing if, for every branch \(\omega\) in that channel,
the canonical windowwise moving family of
\cref{paper7:lem:windowwise-shell-pressure-translation} has the following
property: if
\[
  J_1\notin L^1(\tau_1,\infty)
\]
on every later tail \([\tau_1,\infty)\), then \(\omega\) enters the
exterior/leakage alternative of
\cref{paper7:def:stratified-moving-carrier-routing}.
\end{definition}

\begin{definition}[Windowwise positive-scale routing statement]
\label{paper7:def:windowwise-positive-scale-routing}
The pre-carrier validated compact-cost channel has windowwise positive-scale
routing if, for every branch \(\omega\) in that channel, the canonical
windowwise moving family of
\cref{paper7:lem:windowwise-shell-pressure-translation} has the following
property: if
\[
  J_7\notin L^1(\tau_1,\infty)
\]
on every later tail \([\tau_1,\infty)\), then \(\omega\) enters the
positive-scale routing alternative of
\cref{paper7:def:stratified-moving-carrier-routing}, unless it has already
entered the exterior/leakage or negative scale-drift alternatives of the same
stratified moving-carrier package.
\end{definition}

\subsubsection{Discharge of the windowwise routing statements}

We now prove that the two routing properties of
\cref{paper7:def:windowwise-transition-routing,paper7:def:windowwise-positive-scale-routing}
hold for the canonical windowwise moving family of
\cref{paper7:lem:windowwise-shell-pressure-translation}.  The proof does not
convert arbitrary nonintegrability into a false persistent unit-window lower
bound.  Instead it uses the local state-space reading of a failed carrier
component: the nonintegrable component defines normalized carrier blocks; each
block is then labelled by the ordered first-failure stratification already used
in \cref{paper7:lem:no-unrecorded-cost-escape}.

\begin{definition}[Normalized carrier blocks for a failed component]
\label{paper7:def:normalized-carrier-blocks}
Let \(g\in L^1_{\rm loc}([\tau_0,\infty))\) be nonnegative and assume
\[
  \int_T^\infty g(\tau)\,d\tau=\infty
  \qquad\text{for every }T\ge\tau_0 .
\]
A normalized \(g\)-carrier block sequence on a tail \([\tau_1,\infty)\) is a
sequence of intervals \([s_m,t_m]\), with
\(\tau_1\le s_m<t_m\), \(s_m\to\infty\), \(t_m\le s_{m+1}\), and
\[
  \int_{s_m}^{t_m}g(\tau)\,d\tau=1 .
\]
The associated carrier probability measures are
\[
  d\mu_m(\tau):=g(\tau)\mathbf 1_{[s_m,t_m]}(\tau)\,d\tau .
\]
For \(g=J_i\), \(i\in\{1,7\}\), a \(J_i\)-carrier state means the terminal
local state obtained by evaluating the branch, the repaired gauge, the moving
cutoff, and the relevant component \(J_i\) against a subsequential family
\(\mu_m\).  The state is labelled by the first state-space alternative reached
by the ordered validation procedure of
\cref{paper7:lem:no-unrecorded-cost-escape}.
\end{definition}

\begin{lemma}[Carrier-state extraction from nonintegrability]
\label{paper7:lem:nonintegrable-component-certificate}
Let \(g\in L^1_{\rm loc}([\tau_0,\infty))\) be nonnegative and assume
\[
  \int_T^\infty g(\tau)\,d\tau=\infty
  \qquad\text{for every }T\ge\tau_0 .
\]
Then every tail \([\tau_1,\infty)\) contains a normalized \(g\)-carrier block
sequence in the sense of \cref{paper7:def:normalized-carrier-blocks}.  After
passing to a subsequence, exactly one of the following two density regimes
occurs:
\begin{enumerate}[label=\textup{(C\arabic*)}]
\item \emph{Persistent unit-window carrier mass.}  There are \(\gamma>0\) and
unit intervals \(U_m=[T_m,T_m+1]\), \(T_m\to\infty\), such that
\[
  \int_{U_m}g(\tau)\,d\tau\ge\gamma .
\]
\item \emph{Diffuse carrier mass.}  For one, hence every, normalized block
subsequence retained in this case,
\[
  \sup_{\substack{T\ge s_m\\T+1\le t_m}}
  \int_T^{T+1}g(\tau)\,d\tau\longrightarrow0 .
\]
The probability measures \(d\mu_m=g\mathbf 1_{[s_m,t_m]}d\tau\) are then the
diffuse local carrier states of \(g\).
\end{enumerate}
If, in addition, \(q(\tau):=|a(\tau)|+\|b(\tau)\|_\infty\) satisfies the
uniform unit-window bound
\[
  Q:=\sup_{T\ge\tau_0}\int_T^{T+1}q(\tau)\,d\tau<\infty .
\]
then in case \textup{(C1)} either there are selected times
\(\tau_m\in U_m\) such that
\[
  g(\tau_m)\ge \gamma/2,\qquad q(\tau_m)\le 4Q/\gamma .
\]
after relabelling, or the normalized \(g\)-carrier mass on \(U_m\) is carried
by the high-modulation sets
\[
  H_m:=\{\tau\in U_m:q(\tau)>4Q/\gamma\},
  \qquad
  \int_{H_m}g(\tau)\,d\tau\ge\gamma/2 .
\]
The latter is a modulation carrier state; it is not converted into a shell or
scale lower bound.
\end{lemma}

\begin{proof}
\emph{Block construction.}  Fix a tail \([\tau_1,\infty)\).  Choose
\(s_1\ge\tau_1\).  Since every tail integral of \(g\) is infinite, there is a
finite \(t_1>s_1\) with \(\int_{s_1}^{t_1}g=1\); for instance take the first
level time of the absolutely continuous primitive
\(\int_{s_1}^T g\).  Having chosen \(t_m\), choose \(s_{m+1}\ge t_m+1\) and
then \(t_{m+1}>s_{m+1}\) with \(\int_{s_{m+1}}^{t_{m+1}}g=1\).  This gives
normalized carrier blocks and \(s_m\to\infty\).
\par
\emph{Persistent versus diffuse.}  If case \textup{(C1)} fails, then for each
j there is \(m(j)\) such that every unit interval inside every retained block
with index \(m\ge m(j)\) has \(g\)-mass at most \(1/j\).  A diagonal
subsequence gives the displayed diffuse condition.  Conversely, if
\textup{(C1)} holds then the block sequence has a persistent unit-window
carrier by definition.  The two alternatives are therefore exhaustive after
subsequence extraction.
\par
\emph{Bounded modulation on persistent windows.}  In case \textup{(C1)}, set
\[
  H_m:=\{\tau\in U_m:q(\tau)>4Q/\gamma\},
  \qquad E_m:=U_m\setminus H_m .
\]
Chebyshev gives \(|H_m|\le\gamma/4\).  Since
\(\int_{U_m}g\ge\gamma\), either \(\int_{E_m}g\ge\gamma/2\) or
\(\int_{H_m}g\ge\gamma/2\) along a subsequence.  In the first case
\(|E_m|\le1\) gives a point \(\tau_m\in E_m\) with
\(g(\tau_m)\ge\gamma/2\) and \(q(\tau_m)\le4Q/\gamma\).  In the second case
the displayed high-modulation carrier mass holds.  This proves the final
dichotomy and deliberately leaves the second case as a carrier state rather
than replacing it by a pointwise shell or scale estimate.
\end{proof}

\begin{lemma}[Transition carrier states are exterior/leakage states]
\label{paper7:lem:transition-carrier-state-stratification}
Let \(\omega\) be a branch in the pre-carrier validated compact-cost channel,
and let \(R(\tau)\) be the canonical windowwise moving family of
\cref{paper7:lem:windowwise-shell-pressure-translation}.  If
\[
  J_1\notin L^1(\tau_1,\infty)
\]
on every later tail, then every \(J_1\)-carrier state of
\cref{paper7:def:normalized-carrier-blocks} is labelled by the
exterior/leakage side of
\cref{paper7:def:stratified-moving-carrier-routing}: critical-tightness
failure, exterior concentration, or local windowed \(H^1\)-loss.
\end{lemma}

\begin{proof}
The pre-carrier channel has already passed the analytic, representation,
critical-annulus, compact-tightness, local \(H^1\), modulation, and
cost-compatibility tests listed in
\cref{paper7:def:precarrier-validated-compact-channel}.  Therefore a later
carrier state cannot be assigned to any earlier non-carrier row unless the
state-space limit itself leaves the retained compact window.  By the ordered
classification in \cref{paper7:lem:no-unrecorded-cost-escape}, such a
departure is recorded exactly as critical-tightness failure, exterior
concentration/multicore escape, or local windowed \(H^1\)-loss.
\par
It remains to exclude the possibility that a nonintegrable \(J_1\)-carrier
state stays retained and yet has no exterior/leakage label.  The component
\[
  J_1(\tau)=\left|\int |V|^2\partial_\tau\Phi(\tau,\cdot)\right|
\]
is the time-transition part of the moving cutoff.  It contains no pressure,
convection, translation, or scale-drift factor.  On a retained compact state
for which critical tightness and local \(H^1\) control remain valid, this term
is precisely one of the moving-annulus summability rows in the carrier
validation test of
\cref{paper7:lem:no-unrecorded-cost-escape}.  If that row passed on a terminal
tail, then \(J_1\in L^1\) on that tail, contradicting the hypothesis.  Hence
the first failed local row for the \(J_1\)-carrier is the retained-window
transition failure.  In the stratified moving-carrier package this row is, by
definition, the exterior/leakage alternative: the moving window is forced to
detect mass crossing the retained compact boundary, or the selected compact
regularity needed to ignore that crossing has failed.  Those are exactly the
three labels stated in the theorem.
\end{proof}

\begin{theorem}[Windowwise transition/leakage routing by carrier states]
\label{paper7:thm:windowwise-transition-routing-certificate}
Let \(\omega\) be a branch in the pre-carrier validated compact-cost channel,
and let \(R(\tau)\) be the canonical windowwise moving family of
\cref{paper7:lem:windowwise-shell-pressure-translation}.  If
\[
  J_1\notin L^1(\tau_1,\infty)
\]
on every later tail, then \(\omega\) enters the exterior/leakage alternative
of \cref{paper7:def:stratified-moving-carrier-routing}.
\end{theorem}

\begin{proof}
Apply \cref{paper7:lem:nonintegrable-component-certificate} with
\(g=J_1\).  It gives normalized \(J_1\)-carrier blocks on every terminal tail,
with either persistent unit-window mass or diffuse carrier mass.  In the
persistent case, the supplementary modulation selection gives bounded
selected-time representatives when the low-modulation subcase occurs; in the
high-modulation subcase and in the diffuse case one keeps the normalized
carrier state instead of imposing an artificial pointwise lower bound.
\par
All three outputs are \(J_1\)-carrier states in the sense of
\cref{paper7:def:normalized-carrier-blocks}.  By
\cref{paper7:lem:transition-carrier-state-stratification}, their first
state-space label is critical-tightness failure, exterior concentration, or
local windowed \(H^1\)-loss.  This is exactly the exterior/leakage alternative
allowed in \cref{paper7:def:stratified-moving-carrier-routing}.
\end{proof}

\begin{corollary}[Windowwise transition/leakage routing holds]
\label{paper7:cor:windowwise-transition-routing-holds}
The pre-carrier validated compact-cost channel has the windowwise
transition/leakage routing property in the sense of
\cref{paper7:def:windowwise-transition-routing}.
\end{corollary}

\begin{proof}
This is the definition of
\cref{paper7:def:windowwise-transition-routing} applied with
\cref{paper7:thm:windowwise-transition-routing-certificate}.
\end{proof}

\begin{theorem}[Positive scale-work carrier states enter the scale-rigid branch]
\label{paper7:thm:positive-scale-certificate-enters-scale-rigid}
Let \(\omega\) be a branch in the pre-carrier validated compact-cost channel,
and let \(R(\tau)\) be the canonical windowwise moving family.  Let
\(\mu_m\) be any normalized \(J_7\)-carrier block sequence.  If the
corresponding carrier state is not labelled by exterior/leakage or by
negative scale drift, then, after admissible scale-center reselection, it is a
scale-rigid local branch in the sense of
\cref{paper3:def:scale-rigid-branch}.  Consequently it is excluded from the
local Type II class by \cref{paper3:prop:scale-rigidity-discharged}.
\end{theorem}

\begin{proof}
\emph{Step 1: what the \(J_7\)-carrier records.}  By definition,
\[
  J_7(\tau)=
  a_+(\tau)\int |V(y,\tau)|^2
  \bigl(-y\cdot\nabla\Phi(\tau,y)\bigr)\,dy .
\]
Thus each normalized block satisfies
\[
  \int_{s_m}^{t_m}
  a_+(\tau)\int |V|^2(-y\cdot\nabla\Phi)\,dy\,d\tau=1 .
\]
This is positive scale-generator work carried by the retained moving window.
It is a scale-center state-space datum, not a consequence of an annular
\(L^2\)-to-CKN conversion.
\par
\emph{Step 2: ordered local label.}  The pre-carrier channel has already passed
the earlier local state-space tests in
\cref{paper7:def:precarrier-validated-compact-channel}.  If the carrier state
loses compact tightness, exits through exterior mass, or loses local
\(H^1\)-control, it has the exterior/leakage label and is outside the retained
case of the present theorem.  If the negative part of the scale drift is the
active obstruction, it is assigned to the Section~6 negative-drift branch and
is also outside the retained case.  The remaining retained possibility is a
compact single-core state with nonzero positive scale-generator work and no
available summable moving-carrier absorption.
\par
\emph{Step 3: scale-rigid realization.}  In the ordered local decomposition,
such a retained compact scale-work state is exactly the scale-rigid residual
alternative.  The compact-cylinder bounds and pressure bounds are inherited
from the pre-carrier retained compact tests and the represented suitable
energy estimates; positive retained concentration is the positive finite
critical annulus transported to the selected compact cylinder; bounded
modulation on compact selected cylinders is supplied by the repaired-gauge
construction and the selected-window modulation bound.  The nonzero positive
scale work fixes the relative scale-center selection instead of allowing a
summable carrier correction, giving clause (R4) of
\cref{paper3:def:scale-rigid-branch}.  Hence, after the admissible
scale-center reselection, the carrier state satisfies (R1)--(R4) and is a
scale-rigid local branch.  Applying
\cref{paper3:prop:scale-rigidity-discharged} excludes it.
\end{proof}

\begin{lemma}[Positive-scale carrier-state stratification]
\label{paper7:lem:positive-scale-work-stratification}
Let \(\omega\) be a branch in the pre-carrier validated compact-cost channel,
and let \(R(\tau)\) be the canonical windowwise moving family of
\cref{paper7:lem:windowwise-shell-pressure-translation}, with cutoff
\(\Phi(\tau,y)=\phi(y/R(\tau))\) and \(J_7\) defined as in
\cref{paper7:def:componentwise-annular-finite-error}.  If
\[
  J_7\notin L^1(\tau_1,\infty)
\]
on every later tail, then every \(J_7\)-carrier state is labelled either by
exterior/leakage, by the Section~6 negative scale-drift branch, or by the
scale-rigid residual branch excluded in
\cref{paper7:thm:positive-scale-certificate-enters-scale-rigid}.
\end{lemma}

\begin{proof}
Carrier-state extraction is
\cref{paper7:lem:nonintegrable-component-certificate} with \(g=J_7\).
The ordered first-failure map of \cref{paper7:lem:no-unrecorded-cost-escape}
first removes exterior loss, critical-tightness failure, and local
\(H^1\)-loss.  The scale decomposition then separates the favorable negative
scale-drift obstruction, already assigned to Section~6, from retained compact
states carrying positive scale work.  The latter states are exactly the
scale-rigid states handled by
\cref{paper7:thm:positive-scale-certificate-enters-scale-rigid}.  These
alternatives exhaust the \(J_7\)-carrier because no other moving component is
being tested here and \(J_2,J_3,J_4,J_5\) are already summable on the
canonical windowwise family by
\cref{paper7:lem:windowwise-shell-pressure-translation}.
\end{proof}

\begin{theorem}[Windowwise positive-scale routing by carrier states]
\label{paper7:thm:windowwise-positive-scale-routing-certificate}
Let \(\omega\) be a branch in the pre-carrier validated compact-cost channel,
and let \(R(\tau)\) be the canonical windowwise moving family of
\cref{paper7:lem:windowwise-shell-pressure-translation}, with cutoff
\(\Phi(\tau,y)=\phi(y/R(\tau))\) and \(J_7\) defined as in
\cref{paper7:def:componentwise-annular-finite-error}.  If
\[
  J_7\notin L^1(\tau_1,\infty)
\]
on every later tail, then \(\omega\) enters the positive-scale routing
alternative of \cref{paper7:def:stratified-moving-carrier-routing}, unless it
has already entered the exterior/leakage or negative-drift alternatives.
\end{theorem}

\begin{proof}
Apply \cref{paper7:lem:positive-scale-work-stratification}.  Exterior/leakage
carrier labels already satisfy alternative \textup{(iii)} of
\cref{paper7:def:stratified-moving-carrier-routing}, and negative drift
satisfies alternative \textup{(vi)}.  In the retained positive-scale case,
\cref{paper7:thm:positive-scale-certificate-enters-scale-rigid} places the
carrier in the scale-rigid residual branch and then applies
\cref{paper3:prop:scale-rigidity-discharged}.  The persistent unit-window
subcase gives the pointwise selected-time form of alternative \textup{(v)}
whenever bounded-modulation points are selected; the high-modulation and
diffuse subcases give the carrier-state form of the same alternative.  Hence
there is no unhandled diffuse \(J_7\) tail.
\end{proof}

\begin{corollary}[Windowwise positive-scale routing holds]
\label{paper7:cor:windowwise-positive-scale-routing-holds}
The pre-carrier validated compact-cost channel has the windowwise
positive-scale routing property in the sense of
\cref{paper7:def:windowwise-positive-scale-routing}.
\end{corollary}

\begin{proof}
This is exactly the alternatives allowed by
\cref{paper7:def:windowwise-positive-scale-routing}, supplied by
\cref{paper7:thm:windowwise-positive-scale-routing-certificate}.
\end{proof}

\begin{remark}[Status of the windowwise routing definitions]
\label{paper7:rem:windowwise-routing-discharged}
\Cref{paper7:cor:windowwise-positive-scale-routing-holds,paper7:cor:windowwise-transition-routing-holds}
together prove unconditionally that the pre-carrier validated compact-cost
channel has the routing properties of
\cref{paper7:def:windowwise-transition-routing,paper7:def:windowwise-positive-scale-routing}.
The two definitions are therefore not independent assumptions in the
present manuscript; they are properties whose proof has been reduced to
the carrier-state extraction
\cref{paper7:lem:nonintegrable-component-certificate}, the transition
state-space label in
\cref{paper7:lem:transition-carrier-state-stratification}, the positive
scale-work bridge
\cref{paper7:thm:positive-scale-certificate-enters-scale-rigid}, and the
canonical windowwise schedule of
\cref{paper7:lem:windowwise-tight-radius-schedule}.  The proof does not use
the invalid implication from nonintegrability to persistent unit-window lower
bounds; diffuse tails are retained as normalized carrier states and routed by
the local ordered stratification.
\end{remark}

\begin{theorem}[Canonical windowwise family reduces routing to transition, positive scale, and negative drift]
\label{paper7:thm:windowwise-routing-gives-stratified-routing}
Assume every branch in the pre-carrier validated compact-cost channel uses a
validated compact cost that is energy-controlled by the fixed identity cost.
Assume repaired-gauge suitable local energy gives the moving carrier
subidentity in the sense of
\cref{paper7:lem:suitable-energy-gives-carrier-subidentities}.  Assume also
that the channel has the windowwise transition/leakage and windowwise
positive-scale routing statements of
\cref{paper7:def:windowwise-transition-routing,paper7:def:windowwise-positive-scale-routing},
and that failure of \(J_6\in L^1(\tau_1,\infty)\) on every later tail for the
canonical windowwise moving family forces entry into the negative scale-drift
obstruction of \cref{paper6a:def:negative-drift}.  Then the channel has
stratified moving-carrier routing in the sense of
\cref{paper7:def:stratified-moving-carrier-routing}.
\end{theorem}

\begin{proof}
Let \(\omega\) be a branch in the pre-carrier validated compact-cost channel,
and let \(R(\tau)\) be the canonical windowwise moving family of
\cref{paper7:lem:windowwise-shell-pressure-translation}.  Write
\(J_1,\ldots,J_7\) for its moving carrier components.
\par
By \cref{paper7:lem:windowwise-shell-pressure-translation},
\[
  J_2,J_3,J_4,J_5\in L^1(\tau_0,\infty).
\]
If there exists \(\tau_1\ge\tau_0\) such that
\[
  J_1,J_6,J_7\in L^1(\tau_1,\infty),
\]
then all seven components are integrable on \([\tau_1,\infty)\).  Therefore
\cref{paper7:lem:successful-moving-test-family} gives the moving-cutoff
carrier closure package, so \(\omega\) satisfies alternative \textup{(ii)} of
\cref{paper7:def:stratified-moving-carrier-routing}.
\par
Assume now that no later tail has \(J_1,J_6,J_7\) all integrable.  Since
\(J_2,J_3,J_4,J_5\) are already integrable on the full tail, at least one of
\(J_1,J_6,J_7\) fails on arbitrarily late tails.  Fix such a component.  As in
the proof of \cref{paper7:thm:small-routing-statements-give-stratified-routing},
failure on arbitrarily late tails implies failure on every later tail:
otherwise the component would be integrable on some tail and hence on every
subsequent tail.  If \(J_1\) fails on every later tail, the windowwise
transition/leakage routing statement gives alternative
\textup{(iii)} of
\cref{paper7:def:stratified-moving-carrier-routing}.  If \(J_6\) fails on
every later tail, the negative scale-drift hypothesis gives alternative
\textup{(vi)}.  If \(J_7\) fails on every later tail, the windowwise
positive-scale routing statement gives alternative \textup{(v)}, unless the
branch has already entered alternative \textup{(iii)} or \textup{(vi)}.
Therefore one
of the alternatives in
\cref{paper7:def:stratified-moving-carrier-routing} always holds.
\end{proof}

\begin{corollary}[Adjacent closure reduced to windowwise transition and positive-scale routing]
\label{paper7:cor:adjacent-closure-reduced-to-windowwise-routing}
Assume the non-carrier adjacent strata listed in
\cref{paper7:thm:noncarrier-exits-route} are closed by their corresponding
local state-space or interface arguments in the ambient Type II class.  Assume
also that the exterior concentration, local windowed \(H^1\)-loss,
modulation-driven recentering branches, scale-rigid residual branches, and the
negative scale-drift obstruction are closed by their assigned local arguments
in the ambient class.  Assume every branch in the pre-carrier validated
compact-cost channel uses a validated
compact cost that is energy-controlled by the fixed identity cost, that
repaired-gauge suitable local energy gives the moving carrier subidentity, and
that the windowwise transition/leakage and windowwise positive-scale routing
statements of
\cref{paper7:def:windowwise-transition-routing,paper7:def:windowwise-positive-scale-routing}
hold, together with the negative-drift routing statement used in
\cref{paper7:thm:windowwise-routing-gives-stratified-routing}.  Then the
compact-cost channel has adjacent-stratum closure in the sense of
\cref{paper7:def:compact-cost-adjacent-closure}.
\end{corollary}

\begin{proof}
\Cref{paper7:thm:windowwise-routing-gives-stratified-routing} gives the
stratified moving-carrier routing package.  Then
\cref{paper7:cor:adjacent-closure-reduced-to-stratified-routing} gives the
adjacent-stratum closure conclusion.
\end{proof}

\begin{definition}[Transition/leakage routing statement]
\label{paper7:def:transition-leakage-routing}
The pre-carrier validated compact-cost channel has transition/leakage routing
if, for every branch \(\omega\) in that channel and every moving test family in
the sense of \cref{paper7:def:moving-test-family}, failure of at least one of
\[
  J_1,\qquad J_2,\qquad J_3
\]
to belong to \(L^1(\tau_1,\infty)\) on every later tail \([\tau_1,\infty)\)
forces \(\omega\) into the exterior/leakage alternative of
\cref{paper7:def:stratified-moving-carrier-routing}.
\end{definition}

\begin{definition}[Pressure routing statement]
\label{paper7:def:pressure-routing}
The pre-carrier validated compact-cost channel has pressure routing if, for
every branch \(\omega\) in that channel and every moving test family in the
sense of \cref{paper7:def:moving-test-family}, failure of
\[
  J_4\in L^1(\tau_1,\infty)
\]
on every later tail \([\tau_1,\infty)\) forces \(\omega\) into the
exterior/leakage alternative of
\cref{paper7:def:stratified-moving-carrier-routing}.
\end{definition}

\begin{definition}[Translation routing statement]
\label{paper7:def:translation-routing}
The pre-carrier validated compact-cost channel has translation routing if, for
every branch \(\omega\) in that channel and every moving test family in the
sense of \cref{paper7:def:moving-test-family}, failure of
\[
  J_5\in L^1(\tau_1,\infty)
\]
on every later tail \([\tau_1,\infty)\) forces \(\omega\) into the
translation-routing alternative of
\cref{paper7:def:stratified-moving-carrier-routing}.
\end{definition}

\begin{definition}[Positive-scale routing statement]
\label{paper7:def:positive-scale-routing}
The pre-carrier validated compact-cost channel has positive-scale routing if,
for every branch \(\omega\) in that channel and every moving test family in the
sense of \cref{paper7:def:moving-test-family}, failure of
\[
  J_7\in L^1(\tau_1,\infty)
\]
on every later tail \([\tau_1,\infty)\) forces \(\omega\) into the
positive-scale routing alternative of
\cref{paper7:def:stratified-moving-carrier-routing}.
\end{definition}

\begin{theorem}[Small routing statements imply stratified moving-carrier routing]
\label{paper7:thm:small-routing-statements-give-stratified-routing}
Assume every branch in the pre-carrier validated compact-cost channel uses a
validated compact cost that is energy-controlled by the fixed identity cost.
Assume repaired-gauge suitable local energy gives the moving carrier
subidentity in the sense of
\cref{paper7:lem:suitable-energy-gives-carrier-subidentities}.  Assume also
that the channel has the transition/leakage, pressure, translation, and
positive-scale routing statements of
\cref{paper7:def:transition-leakage-routing,paper7:def:pressure-routing,paper7:def:translation-routing,paper7:def:positive-scale-routing},
and that failure of \(J_6\in L^1(\tau_1,\infty)\) on every later tail forces
entry into the negative scale-drift obstruction of
\cref{paper6a:def:negative-drift}.  Then the channel has stratified
moving-carrier routing in the sense of
\cref{paper7:def:stratified-moving-carrier-routing}.
\end{theorem}

\begin{proof}
Let \(\omega\) be a branch in the pre-carrier validated compact-cost channel,
and choose the affine moving test family
\[
  R(\tau)=R_0+(\tau-\tau_0),\qquad
  \Phi(\tau,y)=\phi(y/R(\tau)).
\]
This is admissible by \cref{paper7:def:moving-test-family}.  Let
\(J_1,\ldots,J_7\) be the associated moving components.
\par
If there exists \(\tau_1\ge\tau_0\) such that
\[
  J_k\in L^1(\tau_1,\infty)
  \qquad\text{for all }1\le k\le7,
\]
then \cref{paper7:lem:successful-moving-test-family} gives the moving-cutoff
carrier closure package.  Hence \(\omega\) satisfies alternative \textup{(ii)}
in \cref{paper7:def:stratified-moving-carrier-routing}.
\par
Assume now that no later tail has all seven components integrable.  Since
there are only seven components, one may choose an index
\(k\in\{1,\dots,7\}\) that fails on arbitrarily late tails.  Such a component
then fails on every later tail: otherwise \(J_k\in L^1(\tau_1,\infty)\) for
some \(\tau_1\), hence also on every later tail beyond \(\tau_1\), contrary to
the choice of \(k\).  If \(k\in\{1,2,3\}\), the transition/leakage routing
statement places \(\omega\) in alternative \textup{(iii)} of
\cref{paper7:def:stratified-moving-carrier-routing}.  If \(k=4\), the pressure
routing statement gives the same conclusion.  If \(k=5\), the translation
routing statement gives alternative \textup{(iv)}.  If \(k=6\), the negative
scale-drift hypothesis gives alternative \textup{(vi)}.  If \(k=7\), the
positive-scale routing statement gives alternative \textup{(v)}.  Therefore
one of the alternatives in
\cref{paper7:def:stratified-moving-carrier-routing} always holds.
\end{proof}

\begin{theorem}[Stratified moving-carrier routing leaves only named branches]
\label{paper7:thm:stratified-moving-carrier-routing}
Assume every branch in the pre-carrier validated compact-cost channel uses a
validated compact cost that is energy-controlled by the fixed identity cost.
Assume also that the channel has stratified moving-carrier routing in the sense
of \cref{paper7:def:stratified-moving-carrier-routing}.  Then every branch in
the pre-carrier validated compact-cost channel, after restriction to a later
terminal tail if necessary, satisfies at least one of the following
conclusions:
\begin{enumerate}[label=\textup{(\roman*)}]
\item it is excluded by the fixed-cutoff carrier closure theorem
\cref{paper7:thm:fixed-carrier-closure};
\item it is excluded by the moving-cutoff carrier closure theorem
\cref{paper7:thm:moving-carrier-closure};
\item it is assigned to one of the already named critical-tightness failure,
exterior concentration, or local windowed \(H^1\)-loss branches;
\item it is assigned to the already named modulation-integrability failure or
to a multicore/exterior branch produced after recentering;
\item it is assigned to the already named scale-rigid residual alternative or
to a bounded selected-window limit refinement of that alternative;
\item it enters the negative scale-drift obstruction of
\cref{paper6a:def:negative-drift}.
\end{enumerate}
\end{theorem}

\begin{proof}
Let \(\omega\) be a branch in the pre-carrier validated compact-cost channel.
By the energy-control hypothesis, \(\omega\) is one of the branches covered by
\cref{paper7:def:stratified-moving-carrier-routing}.  If alternative
\textup{(i)} in that definition holds, then
\cref{paper7:thm:fixed-carrier-closure} excludes \(\omega\).  If alternative
\textup{(ii)} holds, then
\cref{paper7:thm:moving-carrier-closure} excludes \(\omega\).  In the
remaining alternatives, the branch assignment is part of the definition of the
stratified routing package itself.  This yields the displayed list.
\end{proof}

\begin{corollary}[Adjacent closure reduced to stratified moving-carrier routing]
\label{paper7:cor:adjacent-closure-reduced-to-stratified-routing}
Assume the non-carrier adjacent strata listed in
\cref{paper7:thm:noncarrier-exits-route} are closed by their corresponding
local state-space or interface arguments in the ambient Type II class.  Assume
also that the exterior concentration, local windowed \(H^1\)-loss,
modulation-driven recentering branches, scale-rigid residual branches, and the
negative scale-drift obstruction appearing in
\cref{paper7:def:stratified-moving-carrier-routing} are closed by their
assigned local arguments in the ambient class.  If the pre-carrier validated
compact-cost channel has stratified moving-carrier routing, then the
compact-cost channel has adjacent-stratum closure in the sense of
\cref{paper7:def:compact-cost-adjacent-closure}.
\end{corollary}

\begin{proof}
Let \(\omega\) be a branch with active supercritical scaling and validated
compact Type II cost divergence.  By
\cref{paper7:thm:classwise-retained-compact-coverage}, either \(\omega\)
remains in the retained compact stratum or it enters one of the named adjacent
strata.  The retained compact stratum is excluded by
\cref{paper7:cor:retained-stratum-cost-exclusion}.  If \(\omega\) enters a
non-carrier adjacent stratum before the pre-carrier channel, then the first
hypothesis closes it by \cref{paper7:thm:noncarrier-exits-route}.  It remains
to consider a branch in the pre-carrier validated compact-cost channel.
\Cref{paper7:thm:stratified-moving-carrier-routing} then shows that \(\omega\)
is either excluded by fixed or moving carrier closure, or is assigned to one
of the branches closed by the second hypothesis.  Therefore every compact-cost
branch is either excluded on the retained stratum or routed to a previously
closed adjacent branch.  This is exactly
\cref{paper7:def:compact-cost-adjacent-closure}.
\end{proof}

\begin{corollary}[Adjacent closure reduced to the small routing statements]
\label{paper7:cor:adjacent-closure-reduced-to-small-routing-statements}
Assume the non-carrier adjacent strata listed in
\cref{paper7:thm:noncarrier-exits-route} are closed by their corresponding
local state-space or interface arguments in the ambient Type II class.  Assume
also that the exterior concentration, local windowed \(H^1\)-loss,
modulation-driven recentering branches, scale-rigid residual branches, and the
negative scale-drift obstruction are closed by their assigned local arguments
in the ambient class.  Assume every branch in the pre-carrier validated
compact-cost channel uses a validated compact cost that is energy-controlled by
the fixed identity cost, that repaired-gauge suitable local energy gives the
moving carrier subidentity, and that the transition/leakage, pressure,
translation, and positive-scale routing statements of
\cref{paper7:def:transition-leakage-routing,paper7:def:pressure-routing,paper7:def:translation-routing,paper7:def:positive-scale-routing}
hold, together with the negative-drift routing statement used in
\cref{paper7:thm:small-routing-statements-give-stratified-routing}.  Then the
compact-cost channel has adjacent-stratum closure in the sense of
\cref{paper7:def:compact-cost-adjacent-closure}.
\end{corollary}

\begin{proof}
\Cref{paper7:thm:small-routing-statements-give-stratified-routing} gives the
stratified moving-carrier routing package.  Then
\cref{paper7:cor:adjacent-closure-reduced-to-stratified-routing} gives the
adjacent-stratum closure conclusion.
\end{proof}

\begin{remark}[Cost divergence on the retained compact stratum]
The compact Type II cost argument gives infinite localized renormalization cost
under the compact PDE hypotheses.  Passing from that cost divergence to
exclusion from the Type II class uses the corrected localized energy inequality
constructed in
\cref{paper7:lem:monotonicity-carrier-gives-cost-to-exclusion-data}.  On the
retained compact stratum, the carrier exists by
\cref{paper7:thm:carrier-validation-retained-stratum}; failure of the carrier
test is not a retained compact branch but an adjacent stratum.
\end{remark}

\clearpage
\section{Terminal Profile Decoupling and Exterior Regularity}
\label{sec:terminal-profile-decoupling-exterior}

This section records the terminal-profile estimates used to remove
multibubble configurations, small critical profiles, and profiles carried at a
positive physical distance from the selected Type II core.  The statements are
conditional on the analytic inputs explicitly stated below: critical profile
decomposition, terminal windowed profile completeness, repaired-gauge
representation, local Caccioppoli estimates, exterior regularity, and the
stationary rigidity input for the compact scale-collapse limit.

\subsection{Renormalized frames and active critical profiles}

A renormalized frame consists of a center \(x_c(t)\), a scale
\(\lambda(t)>0\), and renormalized time
\[
  \frac{d\tau}{dt}=\lambda(t)^{-2}.
\]
It represents the physical solution by
\[
  V(y,\tau)=\lambda(t)u(x_c(t)+\lambda(t)y,t).
\]
When the repaired gauge is imposed, the represented velocity satisfies
\[
  \partial_\tau V+(V\cdot\nabla)V+\nabla P
  =
  \nu\Delta V
  +a(\tau)(V+y\cdot\nabla V)
  +b(\tau)\cdot\nabla V,
  \qquad \nabla\cdot V=0,
\]
where
\[
  a(\tau)=\partial_\tau\log\lambda(\tau),
\]
and \(b\) is the translation modulation coefficient determined by the
centering condition.  In a fixed frame, every other profile either remains at
bounded relative position and comparable scale, escapes to spatial infinity in
the frame variables, appears at a much larger same-point scale as a
perturbative ambient field in the innermost frame, or lies at a smaller scale
and therefore requires selecting that smaller active frame.

\begin{definition}[Active critical profile data]
\label{paper8:def:active-critical-profile-data}
Let \(t_n\uparrow T^*\).  Active critical profile data consist of
divergence-free profiles \(\phi^j\in L^3(\mathbb R^3)\), scales
\(\lambda_j^n\to0\), centers \(x_j^n\to x_j^*\), and finite truncation
remainders \(r_n^J\) such that, for every finite index set \(J\),
\[
  u(\cdot,t_n)
  =
  \sum_{j\in J}\Lambda_{\lambda_j^n,x_j^n}\phi^j+r_n^J,
  \qquad
  (\Lambda_{\lambda,x_0}\phi)(x)
  =
  \lambda^{-1}\phi\left(\frac{x-x_0}{\lambda}\right).
\]
After grouping comparable parameters into compound profiles, the parameters
are pairwise orthogonal and the critical mass decoupling
\[
  \|u(t_n)\|_3^3
  =
  \sum_{j\in J}\|\phi^j\|_3^3+\|r_n^J\|_3^3+o_n(1)
\]
holds for every finite \(J\).  A profile is active if its nonlinear evolution
is not controlled by the small-data perturbative theory.  The small-data
threshold is denoted by \(\varepsilon_{\rm sd}>0\); the corresponding
active-profile mass floor is stated and proved in
\cref{paper8:thm:active-profile-mass-floor}.
\end{definition}

\begin{lemma}[Finite active profile count]
\label{paper8:lem:finite-active-count}
Assume
\[
  \sup_{t<T^*}\|u(t)\|_3\le M<\infty .
\]
Assume also that every active profile satisfies
\[
  \|\phi^j\|_3\ge\varepsilon_{\rm sd}.
\]
Then the number of active profiles is finite and satisfies
\[
  \#\mathcal J_{\rm act}\le
  \left\lfloor\frac{M^3}{\varepsilon_{\rm sd}^3}\right\rfloor .
\]
\end{lemma}

\begin{proof}
Let \(J\subset\mathcal J_{\rm act}\) be finite.  By critical \(L^3\)-mass
decoupling,
\[
  \|u(t_n)\|_3^3
  =
  \sum_{j\in J}\|\phi^j\|_3^3+\|r_n^J\|_3^3+o_n(1).
\]
Taking the limsup and using \(\|r_n^J\|_3^3\ge0\),
\[
  M^3\ge \sum_{j\in J}\|\phi^j\|_3^3
  \ge (\#J)\varepsilon_{\rm sd}^3.
\]
Thus \(\#J\le M^3/\varepsilon_{\rm sd}^3\) for every finite active set
\(J\).  If the active set had more than
\(\lfloor M^3/\varepsilon_{\rm sd}^3\rfloor\) elements, one could choose a
finite subset with one additional element, contradicting the preceding
inequality.  This proves the integer bound and finiteness.
\end{proof}

\subsection{Same-point and separated-point reductions}

\begin{lemma}[Comparable same-point profiles form a compound profile]
\label{paper8:lem:compound-same-point}
Let a finite set \(\mathcal C\) of active profiles have the same physical
center \(x_j^*=x_*\).  Suppose that for a representative center and scale
\((x_*^n,\lambda_*^n)\),
\[
  \frac{\lambda_j^n}{\lambda_*^n}\to\rho_j\in(0,\infty),
  \qquad
  \frac{x_j^n-x_*^n}{\lambda_*^n}\to z_j\in\mathbb R^3 .
\]
Then, in the frame \((x_*^n,\lambda_*^n)\),
\[
  \lambda_*^n(\lambda_j^n)^{-1}
  \phi^j\!\left(
  \frac{x_*^n+\lambda_*^ny-x_j^n}{\lambda_j^n}
  \right)
  \to
  \rho_j^{-1}\phi^j\left(\frac{y-z_j}{\rho_j}\right)
\]
strongly in \(L^3_{\rm loc}\), and strongly in \(L^3(\mathbb R^3)\) when the
profile convergence is global.  Hence the comparable class has compound
profile
\[
  \Phi_{\mathcal C}(y)
  :=
  \sum_{j\in\mathcal C}
  \rho_j^{-1}\phi^j\left(\frac{y-z_j}{\rho_j}\right).
\]
If \(\|\Phi_{\mathcal C}\|_3\ge\varepsilon_{\rm sd}\), this compound profile
is one active profile in the selected frame; if
\(\|\Phi_{\mathcal C}\|_3<\varepsilon_{\rm sd}\), it is perturbative by
small-data stability.
\end{lemma}

\begin{proof}
Set
\[
  \rho_j^n=\frac{\lambda_j^n}{\lambda_*^n},
  \qquad
  z_j^n=\frac{x_j^n-x_*^n}{\lambda_*^n}.
\]
The transformed profile is
\[
  T_n\phi^j(y)
  =
  (\rho_j^n)^{-1}\phi^j\left(\frac{y-z_j^n}{\rho_j^n}\right),
\]
and the proposed limit is
\[
  T\phi^j(y)
  =
  \rho_j^{-1}\phi^j\left(\frac{y-z_j}{\rho_j}\right).
\]
For \(\rho\in(0,\infty)\) and \(z\in\mathbb R^3\), the map
\[
  T_{\rho,z}\phi(y)=\rho^{-1}\phi\left(\frac{y-z}{\rho}\right)
\]
is strongly continuous on \(L^3(\mathbb R^3)\).  Indeed, for
\(\phi\in C_c^\infty\), convergence of \(\rho_j^n\to\rho_j\) and
\(z_j^n\to z_j\) gives uniform convergence on a common compact support and
therefore \(L^3\)-convergence.  For a general \(\phi\in L^3\), choose
\(\psi\in C_c^\infty\) with \(\|\phi-\psi\|_3<\varepsilon\).  Since each
\(T_{\rho,z}\) is an \(L^3\)-isometry,
\[
  \|T_{\rho_j^n,z_j^n}\phi-T_{\rho_j,z_j}\phi\|_3
  \le
  2\varepsilon
  +\|T_{\rho_j^n,z_j^n}\psi-T_{\rho_j,z_j}\psi\|_3.
\]
Letting \(n\to\infty\) and then \(\varepsilon\downarrow0\) proves global
\(L^3\)-convergence; local convergence follows immediately.  Summing over the
finite class gives \(\Phi_{\mathcal C}\).  The final alternative is the
small-data criterion: a compound profile above the threshold is active, while
one below the threshold belongs to the perturbative class.
\end{proof}

\begin{lemma}[Innermost-frame rescaling of strict outer profiles]
\label{paper8:lem:same-point-rescaled-form}
Suppose all active profiles have the same physical center \(x_*\) and split
into distinct scale classes.  If \(\lambda_1^n\) is the smallest active scale,
then every outer scale \(\lambda_j^n\gg\lambda_1^n\) satisfies
\[
  \mu_j^n:=\lambda_1^n/\lambda_j^n\to0,
\]
and in the innermost frame the \(j\)-th outer profile has the form
\[
  W_j^n(y)
  =
  \mu_j^n
  \phi^j\left(\frac{x_*^n-x_j^n}{\lambda_j^n}+\mu_j^n y\right).
\]
\end{lemma}

\begin{proof}
The \(j\)-th profile in physical coordinates is, by the active-profile
parametrization,
\[
  u_j^n(x)
  =
  (\lambda_j^n)^{-1}\phi^j\!\left(\frac{x-x_j^n}{\lambda_j^n}\right).
\]
Pull this back into the innermost frame \(x=x_*^n+\lambda_1^n y\), with the
critical rescaling \(W_j^n(y)=\lambda_1^n u_j^n(x_*^n+\lambda_1^n y)\).
Substituting,
\[
  W_j^n(y)
  =
  \frac{\lambda_1^n}{\lambda_j^n}
  \phi^j\!\left(\frac{x_*^n+\lambda_1^n y-x_j^n}{\lambda_j^n}\right)
  =
  \mu_j^n\phi^j\!\left(\frac{x_*^n-x_j^n}{\lambda_j^n}+\mu_j^n y\right),
\]
which is the displayed form.  The convergence \(\mu_j^n\to0\) is the
hypothesis \(\lambda_j^n\gg\lambda_1^n\).
\end{proof}

\begin{lemma}[Local \(L^3\)-vanishing of strict outer profiles in the innermost frame]
\label{paper8:lem:same-point-outer-vanish}
Under the hypotheses of \cref{paper8:lem:same-point-rescaled-form}, the outer
profiles satisfy, for every \(R<\infty\),
\[
  \|W_j^n\|_{L^3(B_R)}\to0\quad (n\to\infty).
\]
\end{lemma}

\begin{proof}
This is a direct consequence of the static same-point outer-profile vanishing
lemma \cref{paper8:lem:same-point-outer-static}, applied with
\(\mu_n=\mu_j^n\), \(\xi_n=(x_*^n-x_j^n)/\lambda_j^n\), and \(\phi=\phi^j\).
\end{proof}

\begin{assumption}[Perturbative package for strict outer profiles]
\label{paper8:ass:same-point-scale-separation}
In the same setting as \cref{paper8:lem:same-point-rescaled-form}, the
remaining analytic input for the same-point cascade is:
\begin{enumerate}[label=\textup{(\roman*)}]
\item finite-offset regular outer profiles, namely those for which
\((x_*^n-x_j^n)/\lambda_j^n\) is bounded, decompose on each fixed inner ball
into a constant drift plus an \(O((\mu_j^n)^2)\) shear remainder;
\item the constant drifts produced by item (i) are absorbed into the
translation-modulation coefficient of the innermost frame;
\item the remaining pressure, cutoff, and nonlinear interactions generated by
the outer profiles are perturbative in the compact scale-collapse topology
in which the local perturbation criterion
\cref{paper8:lem:local-perturbation} and the pressure-tail decoupling
\cref{paper8:lem:pressure-kernel-tail} apply.
\end{enumerate}
The escaping-offset alternative
\((x_*^n-x_j^n)/\lambda_j^n\to\infty\) is treated by
\cref{paper8:lem:same-point-outer-vanish} and is not part of this assumption.
\end{assumption}

\begin{lemma}[Separated physical centers are locally invisible]
\label{paper8:lem:separated-centers-vanish}
Let \(p\ne q\) be two active physical concentration points,
\[
  x_p^*\ne x_q^*,
  \qquad
  d_{pq}:=|x_p^*-x_q^*|>0.
\]
In the \(p\)-frame define the \(q\)-profile contribution
\[
  W_{q\to p}^n(y)
  :=
  \lambda_p^n(\lambda_q^n)^{-1}
  \phi^q\left(
  \frac{x_p^n+\lambda_p^ny-x_q^n}{\lambda_q^n}
  \right).
\]
If \(\phi^q\in L^3(\mathbb R^3)\), then for every \(R<\infty\),
\[
  \|W_{q\to p}^n\|_{L^3(B_R)}\to0.
\]
If the corresponding profile pressure \(\Pi^q\in L^{3/2}(\mathbb R^3)\), then
the pure pressure contribution
\[
  P_{q\to p}^n(y)
  =
  (\lambda_p^n)^2(\lambda_q^n)^{-2}
  \Pi^q\left(
  \frac{x_p^n+\lambda_p^ny-x_q^n}{\lambda_q^n}
  \right)
\]
satisfies
\[
  \|P_{q\to p}^n\|_{L^{3/2}(B_R)}\to0 .
\]
\end{lemma}

\begin{proof}
By the change of variables \(x=x_p^n+\lambda_p^ny\),
\[
  \|W_{q\to p}^n\|_{L^3(B_R)}^3
  =
  \int_{B(x_p^n,\lambda_p^nR)}
  (\lambda_q^n)^{-3}
  \left|\phi^q\left(\frac{x-x_q^n}{\lambda_q^n}\right)\right|^3\,dx .
\]
Set \(z=(x-x_q^n)/\lambda_q^n\).  Then
\[
  \|W_{q\to p}^n\|_{L^3(B_R)}^3
  =
  \int_{B(c_n,\rho_n)}|\phi^q(z)|^3\,dz,
\]
where
\[
  c_n=\frac{x_p^n-x_q^n}{\lambda_q^n},
  \qquad
  \rho_n=\frac{\lambda_p^nR}{\lambda_q^n}.
\]
Since \(x_p^n\to x_p^*\), \(x_q^n\to x_q^*\), and
\(\lambda_q^n\to0\),
\[
  |c_n|\sim d_{pq}/\lambda_q^n\to\infty.
\]
Moreover
\[
  \frac{\rho_n}{|c_n|}
  =
  \frac{\lambda_p^nR}{|x_p^n-x_q^n|}
  \to0.
\]
Thus the balls \(B(c_n,\rho_n)\) escape to spatial infinity: for every
\(A<\infty\), they are contained in \(\{|z|>A\}\) for all sufficiently large
\(n\).  Since \(\phi^q\in L^3\),
\[
  \int_{|z|>A}|\phi^q(z)|^3\,dz\to0
  \qquad(A\to\infty),
\]
and hence \(\|W_{q\to p}^n\|_{L^3(B_R)}\to0\).

For the pressure term the same change of variables gives
\[
  \|P_{q\to p}^n\|_{L^{3/2}(B_R)}^{3/2}
  =
  \int_{B(c_n,\rho_n)}|\Pi^q(z)|^{3/2}\,dz.
\]
The same escaping-ball argument and \(\Pi^q\in L^{3/2}\) imply the stated
pressure convergence.
\end{proof}

\begin{definition}[Multipoint decoupling conditions]
\label{paper8:def:multipoint-decoupling}
For each active point and selected frame, the required multipoint decoupling
conditions are:
\begin{enumerate}[label=\textup{(\roman*)}]
\item local velocity decoupling in the selected frame;
\item pressure decoupling modulo time-dependent constants in
\(L^1_\tau H^{-1}_y\) on compact cylinders;
\item localization errors from spatial cutoffs vanish in
\(L^1_\tau H^{-1}_y\);
\item repaired-gauge modulation coefficients remain bounded after exterior
ambient terms are absorbed;
\item each active point has positive finite critical mass in its own frame;
\item after profiles centered at other physical points are treated as
perturbative exterior terms, the selected frame satisfies the compact
scale-collapse hypotheses: repaired-gauge representation, pressure
reconstruction, Caccioppoli regularity, bounded modulation, and the
compactness and tightness inputs used in the limiting passage.
\end{enumerate}
\end{definition}

\begin{theorem}[Reduction of active multibubbles]
\label{paper8:thm:multibubble-frame-reduction}
Assume the Type II profile data, the active-profile mass floor, the same-point
nonlinear decoupling conditions, the multipoint decoupling conditions of
\cref{paper8:def:multipoint-decoupling}, and the compact scale-collapse
stationary-rigidity input.  Then no active multibubble Type II configuration
can occur.
\end{theorem}

\begin{proof}
By \cref{paper8:lem:finite-active-count}, there are only finitely many active
profiles.  Partition them by physical concentration point \(x_\alpha^*\), and
within each physical point partition them by comparable scale class.

If a scale class contains several comparable profiles,
\cref{paper8:lem:compound-same-point} combines them into one compound profile.
If the compound profile is below the small-data threshold, it is
perturbative; if it is active, it is a single active profile for that frame.

If a physical point contains strict same-point scale separation,
\cref{paper8:ass:same-point-scale-separation} and the same-point nonlinear
decoupling conditions reduce the innermost active scale to a single
repaired-gauge solution with perturbative outer interactions.  The compact
scale-collapse input then rules out the corresponding same-point cascade.

It remains to consider separated physical points.  Choose any active point
\(x_\alpha^*\) and place the selected frame at its active scale and center.
By \cref{paper8:lem:separated-centers-vanish}, every profile centered at a
different physical point is locally \(L^3\)-invisible in this frame.  By the
multipoint decoupling conditions, its pressure, cutoff, and modulation
interactions are also perturbative in the compact scale-collapse topology.
Therefore the selected frame contains a single active repaired-gauge solution
with positive finite critical mass, bounded modulation, pressure
reconstruction, Caccioppoli regularity, and the compactness and tightness
inputs required by the scale-collapse argument.

The scale of that active solution tends to zero, so the scale-collapse input
applies.  The compactness argument extracts a nonzero stationary
\(L^3(\mathbb R^3)\) profile in the stationary-rigidity class, contradicting
the stationary rigidity theorem.  This excludes the chosen active point.  Since
the chosen active point was arbitrary, no active multibubble configuration can
occur.
\end{proof}

\subsection{Local perturbation and pressure decoupling}

\begin{lemma}[Local perturbation criterion]
\label{paper8:lem:local-perturbation}
Let \(I\subset\mathbb R\) be compact and \(Q_R=B_R\times I\).  Let \(U_n\) be
a repaired-gauge solution on \(Q_{2R}\) satisfying
\[
  \|U_n\|_{L^\infty_\tau L^3_y(B_{2R})}
  +
  \|U_n\|_{L^2_\tau H^1_y(B_{2R})}
  \le C_R.
\]
Let \(S_n\) be a divergence-free perturbation on \(Q_{2R}\) with
\[
  S_n\to0
  \quad\text{in }L^\infty_\tau L^3_y(B_{2R}),
\]
\[
  U_n\otimes S_n+S_n\otimes U_n+S_n\otimes S_n
  \to0
  \quad\text{in }L^1_\tau L^2_y(B_R),
\]
\[
  \mathcal R_n
  :=
  \partial_\tau S_n-\nu\Delta S_n
  +a_n(S_n+y\cdot\nabla S_n)
  +b_n\cdot\nabla S_n
  \to0
  \quad\text{in }L^1_\tau H^{-1}_y(B_R),
\]
where \(|a_n|+|b_n|\le M_{ab}\), and suppose the pressure cross-term satisfies
\[
  \nabla\mathcal P_{{\rm cross},n}\to0
  \quad\text{in }L^1_\tau H^{-1}_y(B_R).
\]
Then replacing \(U_n+S_n\) by \(U_n\) changes the repaired-gauge
Navier--Stokes equation on \(Q_R\) by an error tending to zero in
\[
  L^1_\tau H^{-1}_y(B_R).
\]
\end{lemma}

\begin{proof}
The linear part is precisely \(\mathcal R_n\), which tends to zero by
hypothesis.  The nonlinear difference is
\[
  (U_n+S_n)\cdot\nabla(U_n+S_n)-U_n\cdot\nabla U_n
  =
  \operatorname{div}F_n,
  \qquad
  F_n:=U_n\otimes S_n+S_n\otimes U_n+S_n\otimes S_n,
\]
using \(\nabla\cdot U_n=\nabla\cdot S_n=0\) distributionally.  By the assumed
\(L^1_\tau L^2_y\) stress convergence, this divergence tends to zero in
\[
  L^1_\tau H^{-1}_y(B_R),
\]
because for \(\varphi\in H^1_0(B_R;\mathbb R^3)\),
\[
\begin{aligned}
  \left|
  \int_{B_R}
  (U_n\otimes S_n+S_n\otimes U_n+S_n\otimes S_n):\nabla\varphi\,dy
  \right|
  &\le
  \|F_n\|_{L^2(B_R)} \\
  &\quad{}\times \|\nabla\varphi\|_{L^2(B_R)}.
\end{aligned}
\]
The pressure contribution vanishes by the assumed cross-pressure convergence.
Therefore the total perturbative error tends to zero in
\[
  L^1_\tau H^{-1}_y(B_R).
\]
\end{proof}

\begin{lemma}[Static same-point outer profiles vanish locally]
\label{paper8:lem:same-point-outer-static}
Let
\[
  W_n(y)=\mu_n\phi(\xi_n+\mu_n y),
  \qquad
  \mu_n\to0,
  \qquad
  \phi\in L^3(\mathbb R^3).
\]
Then, for every \(R<\infty\),
\[
  \|W_n\|_{L^3(B_R)}\to0.
\]
\end{lemma}

\begin{proof}
Changing variables \(z=\xi_n+\mu_n y\),
\[
  \|W_n\|_{L^3(B_R)}^3
  =
  \int_{B(\xi_n,\mu_nR)}|\phi(z)|^3\,dz.
\]
The measure of \(B(\xi_n,\mu_nR)\) tends to zero.  Since
\(|\phi|^3\in L^1(\mathbb R^3)\), absolute continuity of the integral gives
\[
  \int_{B(\xi_n,\mu_nR)}|\phi(z)|^3\,dz\to0,
\]
uniformly in \(\xi_n\).  Hence \(\|W_n\|_{L^3(B_R)}\to0\).
\end{proof}

\begin{lemma}[Pressure decoupling from kernel tails]
\label{paper8:lem:pressure-kernel-tail}
Let \(\Gamma_{ij}\) be the Calderon--Zygmund kernel for
\[
  P=\mathcal R_i\mathcal R_j(F_{ij}).
\]
Assume
\[
  F_n\to0\quad\text{in }L^1_\tau L^2_y(B_A)
  \quad\text{for every fixed }A,
\]
and
\[
  \sup_n\|F_n\|_{L^\infty_\tau L^{3/2}_y(\mathbb R^3)}<\infty.
\]
Then the pressure generated by \(F_n\) satisfies, for every \(R<\infty\),
\[
  P_n-c_{n,R}\to0
  \quad\text{in }L^1_\tau L^2_y(B_R)
\]
for suitable constants \(c_{n,R}(\tau)\), and consequently
\[
  \nabla P_n\to0
  \quad\text{in }L^1_\tau H^{-1}_y(B_R).
\]
\end{lemma}

\begin{proof}
Fix \(A>2R\) and write
\[
  F_n^{\rm loc}=F_n{\bf 1}_{B_A},
  \qquad
  F_n^{\rm far}=F_n{\bf 1}_{\mathbb R^3\setminus B_A}.
\]
For the local part, Calderon--Zygmund boundedness gives
\[
  \|P_n^{\rm loc}\|_{L^1_\tau L^2_y(B_A)}
  \le
  C_A\|F_n\|_{L^1_\tau L^2_y(B_A)}
  \to0.
\]
For the far part define
\[
  c_{n,R}(\tau)=
  \int_{|z|>A}\Gamma_{ij}(-z)F_{n,ij}(z,\tau)\,dz,
\]
first for compactly supported approximations and then by passage to the
\(L^{3/2}\) limit.  For \(y\in B_R\) and \(|z|>A>2R\),
\[
  |\Gamma_{ij}(y-z)-\Gamma_{ij}(-z)|
  \le C_R|z|^{-4}.
\]
Hence
\[
  |P_n^{\rm far}(y,\tau)-c_{n,R}(\tau)|
  \le
  C_R\int_{|z|>A}|z|^{-4}|F_n(z,\tau)|\,dz.
\]
By Holder's inequality with exponents \(3/2\) and \(3\),
\[
  \int_{|z|>A}|z|^{-4}|F_n|\,dz
  \le
  \|F_n(\tau)\|_{L^{3/2}}
  \left(\int_{|z|>A}|z|^{-12}\,dz\right)^{1/3}
  \le
  CA^{-3}\|F_n(\tau)\|_{L^{3/2}}.
\]
Therefore
\[
  \|P_n^{\rm far}-c_{n,R}\|_{L^1_\tau L^\infty_y(B_R)}
  \le
  C_{R,I}A^{-3}
  \sup_n\|F_n\|_{L^\infty_\tau L^{3/2}_y}.
\]
Since \(B_R\) has finite measure, this also controls the far part in
\(L^1_\tau L^2_y(B_R)\).  Taking \(n\to\infty\) in the local term and then
\(A\to\infty\) in the far term gives
\[
  P_n-c_{n,R}\to0
  \quad\text{in }L^1_\tau L^2_y(B_R).
\]
For \(\varphi\in H^1_0(B_R;\mathbb R^3)\),
\[
  |\langle\nabla(P_n-c_{n,R}),\varphi\rangle|
  \le
  \left|\int_{B_R}(P_n-c_{n,R})\nabla\cdot\varphi\,dy\right|
  \le
  \|P_n-c_{n,R}\|_{L^2(B_R)}
  \|\nabla\varphi\|_{L^2(B_R)}.
\]
Thus \(\nabla P_n\to0\) in \(L^1_\tau H^{-1}_y(B_R)\).
\end{proof}

\begin{definition}[Terminal nonlinear profile decoupling]
\label{paper8:def:terminal-nonlinear-decoupling}
On every compact terminal-frame cylinder \(B_R\times I\), let \(U_n\) be the
retained active profile evolution and let \(S_n\) be the sum of all
nonretained profiles and the critical perturbative remainder.  Terminal
nonlinear profile decoupling means:
\begin{enumerate}[label=\textup{(\roman*)}]
\item \(\nabla\cdot S_n=0\) distributionally;
\item \(S_n\to0\) in \(L^\infty_\tau L^3_y(B_{2R})\);
\item
\[
  U_n\otimes S_n+S_n\otimes U_n+S_n\otimes S_n
  \to0
  \quad\text{in }L^1_\tau L^2_y(B_R);
\]
\item
\[
  \partial_\tau S_n-\nu\Delta S_n
  +a_n(S_n+y\cdot\nabla S_n)
  +b_n\cdot\nabla S_n
  \to0
  \quad\text{in }L^1_\tau H^{-1}_y(B_R);
\]
\item pressure sources involving at least one factor of \(S_n\) satisfy
\cref{paper8:lem:pressure-kernel-tail}, or directly satisfy
\[
  \nabla\mathcal P_{{\rm cross},n}\to0
  \quad\text{in }L^1_\tau H^{-1}_y(B_R);
\]
\item cutoff errors introduced in isolating the selected frame vanish in
\(L^1_\tau H^{-1}_y(B_R)\);
\item the effective modulation coefficients remain bounded;
\item after subtracting \(S_n\), the retained solution has positive finite
critical mass, repaired-gauge representation, pressure reconstruction,
Caccioppoli regularity, compactness, tightness, and the compact
scale-collapse limiting inputs.
\end{enumerate}
\end{definition}

\begin{theorem}[Same-point and multipoint decoupling from terminal decoupling]
\label{paper8:thm:terminal-decoupling-implies-local-decoupling}
If terminal nonlinear profile decoupling holds in the sense of
\cref{paper8:def:terminal-nonlinear-decoupling}, then both same-point
scale-decoupling and separated-point frame-decoupling hold.
\end{theorem}

\begin{proof}
For same-point active profiles, group comparable scales by
\cref{paper8:lem:compound-same-point}.  For strict scale separation choose the
innermost active scale as the selected terminal frame.  By
\cref{paper8:lem:same-point-outer-static}, every outer scale is locally
invisible in velocity.  Terminal nonlinear profile decoupling supplies the
nonlinear-stress, linear-residual, pressure-source, localization, modulation,
and compact scale-collapse admissibility conditions on every compact cylinder.
\Cref{paper8:lem:local-perturbation} then shows that the outer profiles and
remainder alter the selected equation only by an error tending to zero in
\(L^1_\tau H^{-1}_y(B_R)\).  Thus the same-point configuration reduces to one
retained compact scale-collapse solution with perturbative interactions.

For separated physical points, choose an active physical point \(x_p^*\) and
its selected frame.  For every profile centered at \(x_q^*\ne x_p^*\),
\cref{paper8:lem:separated-centers-vanish} gives local velocity convergence
to zero in the \(p\)-frame.  If a pure profile pressure is represented in
\(L^2\), the same escaping-ball calculation gives local \(L^2\) pressure
convergence and therefore \(H^{-1}\) convergence of the gradient.  With only
\(L^{3/2}\)-pressure, the mixed pressure sources are controlled by
\cref{paper8:lem:pressure-kernel-tail} through the \(L^2\)-strength hypotheses
included in terminal nonlinear profile decoupling.  The same decoupling
assumption supplies localization-error convergence, bounded modulation, and
compact scale-collapse admissibility.  Applying
\cref{paper8:lem:local-perturbation} shows that separated physical-point
profiles change the selected equation only by an error tending to zero in
\(L^1_\tau H^{-1}_{\rm loc}\).  This proves separated-point decoupling.
\end{proof}

\subsection{Terminal profile decoupling}

\begin{definition}[Terminal active frame]
\label{paper8:def:terminal-active-frame}
Partition active profiles by physical point \(x_j^n\to x_\alpha^*\).  At a
fixed physical point, group comparable scales.  For a comparable-scale class
\(\mathcal C\), a representative frame \((x_{\mathcal C}^n,\lambda_{\mathcal C}^n)\)
has compound profile
\[
  \Phi_{\mathcal C}(y)=
  \sum_{j\in\mathcal C}
  \rho_j^{-1}\phi^j\left(\frac{y-z_j}{\rho_j}\right).
\]
The class \(\mathcal C\) is terminal at \(x_\alpha^*\) if no active class at
the same physical point has strictly smaller scale:
\[
  \lambda_{\mathcal D}^n/\lambda_{\mathcal C}^n\not\to0
  \quad\text{for every active class }\mathcal D.
\]
The terminal frame is
\[
  y=\frac{x-x_{\mathcal C}^n}{\lambda_{\mathcal C}^n},
  \qquad
  \tau=\int^t(\lambda_{\mathcal C}(s))^{-2}\,ds,
\]
and
\[
  V_n(y,\tau)=
  \lambda_{\mathcal C}^n
  u(x_{\mathcal C}^n+\lambda_{\mathcal C}^ny,
  t_n+(\lambda_{\mathcal C}^n)^2\tau).
\]
The retained solution \(U_n\) is the repaired-gauge nonlinear evolution of
\(\mathcal C\) in this frame, and \(S_n:=V_n-U_n\) is the nonretained
remainder field.
\end{definition}

\begin{assumption}[Terminal windowed profile completeness]
\label{paper8:ass:terminal-windowed-profile-completeness}
For a sequence \(f_n\) and a frame \((x_n,\lambda_n)\), define
\[
  \mathfrak m_R(f_n;x_n,\lambda_n)
  :=
  \limsup_{n\to\infty}
  \left\|
  \lambda_nf_n(x_n+\lambda_n\cdot)
  \right\|_{L^3(B_R)}
\]
and
\[
  \mathfrak M_{R,I}(f_n;x_n,\lambda_n)
  :=
  \limsup_{n\to\infty}
  \sup_{\tau\in I}
  \left\|
  \lambda_nf_n(x_n+\lambda_n\cdot,t_n+\lambda_n^2\tau)
  \right\|_{L^3(B_R)}.
\]
Terminal windowed profile completeness means that if
\(\mathfrak M_{R,I}(f_n;x_n,\lambda_n)>0\) for a sequence left in the
remainder, then after passing to a subsequence the profile decomposition
contains a nonzero profile with physical center \(x_n+O(\lambda_n)\) and scale
comparable to \(\lambda_n\).  If that nonlinear profile is above the
small-data threshold, it is active; if it is below the threshold, it is
controlled by the small-data perturbative theory uniformly on compact
windows.
\end{assumption}

\begin{remark}[Status of \cref{paper8:ass:terminal-windowed-profile-completeness}]
\label{paper8:rem:terminal-completeness-replacement}
The Section~8 results below use
\cref{paper8:ass:terminal-windowed-profile-completeness} as a working
hypothesis at the point at which terminal frames first appear.  In
\cref{paper6a:cor:profile-completeness-from-critical-hypotheses} the same
conclusion is derived from the terminal critical profile hypotheses
(\cref{paper6a:def:terminal-critical-profile-hypotheses}), namely a bounded
critical terminal sequence, small-data critical stability, and the critical
Navier--Stokes profile decomposition
\cref{paper6a:ass:critical-ns-profile-decomposition}.  Throughout the
manuscript, \cref{paper8:ass:terminal-windowed-profile-completeness} can be
freely replaced by that conditional package whenever the terminal critical
profile hypotheses are available.
\end{remark}

\begin{lemma}[Terminal local velocity decoupling]
\label{paper8:lem:terminal-velocity-decoupling}
Let \(\mathcal C\) be a terminal active class and \(S_n\) be the nonretained
remainder field in the \(\mathcal C\)-frame.  Then for every compact cylinder
\(B_R\times I\),
\[
  S_n\to0
  \quad\text{in }L^\infty_\tau L^3_y(B_R).
\]
\end{lemma}

\begin{proof}
Assume otherwise.  Then there exist \(R<\infty\), a compact interval \(I\),
\(\delta>0\), a subsequence, and times \(\tau_n\in I\) such that
\[
  \|S_n(\tau_n)\|_{L^3(B_R)}\ge\delta.
\]
In physical variables, after the retained terminal class has been subtracted,
the remainder carries at least \(\delta\) critical \(L^3\)-mass inside a ball
of radius \(O(\lambda_{\mathcal C}^n)\), centered at
\[
  x_{\mathcal C}^n+O(\lambda_{\mathcal C}^n),
\]
at time
\[
  t_n+(\lambda_{\mathcal C}^n)^2\tau_n.
\]
By terminal windowed profile completeness, this nonvanishing mass produces a
profile whose physical center is \(x_\alpha^*\) and whose scale is comparable
to \(\lambda_{\mathcal C}^n\).

If the resulting profile is active, it belongs to the same comparable-scale
class as \(\mathcal C\).  But all comparable active profiles at that physical
point were grouped into \(\mathcal C\), so it cannot remain in \(S_n\), a
contradiction.  If the resulting profile is below the small-data threshold, it
belongs to the perturbative part of the decomposition.  The small-data
stability input and the choice of the terminal truncation give a contribution
smaller than \(\delta/2\) on the compact window, so it cannot account for
\(\|S_n(\tau_n)\|_{L^3(B_R)}\ge\delta\).  This is again a contradiction.

Thus no such \(\delta\) exists.  The same argument covers same-point
larger-scale profiles and separated-point profiles: any nonzero mass seen in
the terminal frame would generate an omitted profile at the terminal scale or
at the selected physical point, contradicting the terminal partition.
\end{proof}

\begin{lemma}[Uniform terminal \(H^1\) bounds]
\label{paper8:lem:terminal-h1-bounds}
Assume the repaired-gauge representation and the local Caccioppoli estimates
hold on every compact terminal cylinder \(B_{2R}\times I\).  Then
\[
  \sup_n
  \|V_n\|_{L^\infty_\tau L^3_y(B_{2R})}
  +
  \sup_n
  \|V_n\|_{L^2_\tau H^1_y(B_{2R})}
  \le C_R,
\]
\[
  \sup_n
  \|U_n\|_{L^\infty_\tau L^3_y(B_{2R})}
  +
  \sup_n
  \|U_n\|_{L^2_\tau H^1_y(B_{2R})}
  \le C_R,
\]
and, since \(S_n=V_n-U_n\),
\[
  \sup_n
  \|S_n\|_{L^\infty_\tau L^3_y(B_{2R})}
  +
  \sup_n
  \|S_n\|_{L^2_\tau H^1_y(B_{2R})}
  \le C_R.
\]
In particular, by Sobolev embedding on \(B_{2R}\),
\[
  U_n,S_n
  \quad\text{are bounded in}\quad
  L^2_\tau L^6_y(B_{2R}).
\]
\end{lemma}

\begin{proof}
The first estimate is the repaired-gauge Caccioppoli estimate for \(V_n\) on
the compact terminal cylinder.  The retained solution \(U_n\) is itself a
repaired-gauge Navier--Stokes solution with positive finite critical mass and
bounded modulation on the same compact cylinder, so the same estimate applies
to \(U_n\).  Since \(S_n=V_n-U_n\), the displayed bound for \(S_n\) follows by
the triangle inequality in the stated norms.  The \(L^2_\tau L^6_y\) bound
follows from the Sobolev embedding \(H^1(B_{2R})\hookrightarrow L^6(B_{2R})\).
\end{proof}

\begin{lemma}[Nonlinear stress decoupling in terminal frames]
\label{paper8:lem:terminal-stress-decoupling}
Let
\[
  F_n:=U_n\otimes S_n+S_n\otimes U_n+S_n\otimes S_n.
\]
Then, for every compact \(B_R\times I\),
\[
  F_n\to0
  \quad\text{in }L^1_\tau L^2_y(B_R).
\]
Moreover
\[
  \sup_n\|F_n\|_{L^\infty_\tau L^{3/2}_y(\mathbb R^3)}<\infty
\]
whenever \(U_n,S_n,V_n\) are uniformly bounded in global
\(L^\infty_\tau L^3_y\).
\end{lemma}

\begin{proof}
By \cref{paper8:lem:terminal-velocity-decoupling},
\[
  \|S_n\|_{L^\infty_\tau L^3_y(B_R)}\to0.
\]
Using Holder's inequality in space and time,
\[
  \|U_n\otimes S_n\|_{L^1_\tau L^2_y(B_R)}
  \le
  \|S_n\|_{L^\infty_\tau L^3_y(B_R)}
  \|U_n\|_{L^1_\tau L^6_y(B_R)}
\]
and
\[
  \|U_n\|_{L^1_\tau L^6_y(B_R)}
  \le
  |I|^{1/2}\|U_n\|_{L^2_\tau L^6_y(B_R)}
  \le C_{R,I}.
\]
Thus \(U_n\otimes S_n\to0\) in \(L^1_\tau L^2_y(B_R)\), and the same estimate
gives \(S_n\otimes U_n\to0\).  Finally,
\[
  \|S_n\otimes S_n\|_{L^1_\tau L^2_y(B_R)}
  \le
  \|S_n\|_{L^\infty_\tau L^3_y(B_R)}
  \|S_n\|_{L^1_\tau L^6_y(B_R)}
  \to0,
\]
because \(\|S_n\|_{L^1_\tau L^6_y(B_R)}\le C_{R,I}\).  Summing the three
terms proves the local convergence.  The global \(L^{3/2}\) bound follows
from Holder's inequality and the global \(L^\infty_\tau L^3_y\) bounds.
\end{proof}

\begin{lemma}[Terminal pressure decoupling]
\label{paper8:lem:terminal-pressure-decoupling}
Let \(P_{{\rm cross},n}\) solve
\[
  -\Delta P_{{\rm cross},n}
  =
  \partial_i\partial_j(F_n)_{ij}.
\]
Assume, in addition, that \(U_n\), \(S_n\), and \(V_n\) are uniformly bounded
in \(L^\infty_\tau L^3_y(\mathbb R^3)\).
Then for every \(R<\infty\) there are constants \(c_{n,R}(\tau)\) such that
\[
  P_{{\rm cross},n}-c_{n,R}
  \to0
  \quad\text{in }L^1_\tau L^2_y(B_R),
\]
and
\[
  \nabla P_{{\rm cross},n}
  \to0
  \quad\text{in }L^1_\tau H^{-1}_y(B_R).
\]
\end{lemma}

\begin{proof}
\Cref{paper8:lem:terminal-stress-decoupling} gives
\[
  F_n\to0
  \quad\text{in }L^1_\tau L^2_y(B_A)
\]
for every fixed \(A\).  The global \(L^\infty_\tau L^3_y\) bounds and
Holder's inequality give
\[
  \sup_n\|F_n\|_{L^\infty_\tau L^{3/2}_y(\mathbb R^3)}<\infty.
\]
Thus the hypotheses of \cref{paper8:lem:pressure-kernel-tail} hold.  That
lemma gives local \(L^2\) pressure convergence modulo constants and
\(L^1_\tau H^{-1}_y(B_R)\) convergence of the pressure gradient.
\end{proof}

\begin{lemma}[Linear residual decoupling]
\label{paper8:lem:terminal-linear-residual}
Let \(V_n=U_n+S_n\) solve
\[
  \partial_\tau V_n+(V_n\cdot\nabla)V_n+\nabla P_n
  =
  \nu\Delta V_n
  +a_n(V_n+y\cdot\nabla V_n)
  +b_n\cdot\nabla V_n.
\]
Suppose \(U_n\) satisfies, in the same repaired gauge,
\[
  \partial_\tau U_n+(U_n\cdot\nabla)U_n+\nabla P_{U,n}
  =
  \nu\Delta U_n
  +a_n(U_n+y\cdot\nabla U_n)
  +b_n\cdot\nabla U_n+e_n,
\]
with
\[
  e_n\to0
  \quad\text{in }L^1_\tau H^{-1}_y(B_R)
\]
on compact terminal cylinders.  Then
\[
  \partial_\tau S_n-\nu\Delta S_n
  -a_n(S_n+y\cdot\nabla S_n)
  -b_n\cdot\nabla S_n
  \to0
  \quad\text{in }L^1_\tau H^{-1}_y(B_R).
\]
Equivalently, the same conclusion holds with the opposite sign convention for
the modulation terms.
\end{lemma}

\begin{proof}
Subtracting the \(U_n\)-equation from the \(V_n\)-equation gives
\[
  \partial_\tau S_n-\nu\Delta S_n
  =
  -\operatorname{div}F_n
  -\nabla P_{{\rm cross},n}
  -e_n
  +a_n(S_n+y\cdot\nabla S_n)
  +b_n\cdot\nabla S_n.
\]
Equivalently,
\[
  \partial_\tau S_n-\nu\Delta S_n
  -a_n(S_n+y\cdot\nabla S_n)
  -b_n\cdot\nabla S_n
  =
  -\operatorname{div}F_n-\nabla P_{{\rm cross},n}-e_n.
\]
By \cref{paper8:lem:terminal-stress-decoupling},
\[
  F_n\to0\quad\text{in }L^1_\tau L^2_y(B_R),
\]
hence
\[
  \operatorname{div}F_n\to0
  \quad\text{in }L^1_\tau H^{-1}_y(B_R).
\]
By \cref{paper8:lem:terminal-pressure-decoupling},
\[
  \nabla P_{{\rm cross},n}\to0
  \quad\text{in }L^1_\tau H^{-1}_y(B_R).
\]
Together with \(e_n\to0\), this proves the displayed residual convergence.

For the opposite sign convention, the two residuals differ by
\[
  2a_n(S_n+y\cdot\nabla S_n)+2b_n\cdot\nabla S_n.
\]
By \cref{paper8:lem:terminal-velocity-decoupling} and finite measure,
\[
  S_n\to0\quad\text{in }L^\infty_\tau L^2_y(B_R).
\]
For \(\varphi\in H^1_0(B_R)\),
\[
  |\langle y\cdot\nabla S_n,\varphi\rangle|
  =
  \left|\int_{B_R}S_n(3\varphi+y\cdot\nabla\varphi)\,dy\right|
  \le
  C_R\|S_n\|_{L^2(B_R)}\|\varphi\|_{H^1(B_R)},
\]
and
\[
  |\langle b_n\cdot\nabla S_n,\varphi\rangle|
  \le
  |b_n|\,\|S_n\|_{L^2(B_R)}\|\nabla\varphi\|_{L^2(B_R)}.
\]
The term \(S_n\) itself is controlled by the same \(L^2\)-convergence.  Since
\(a_n\) and \(b_n\) are bounded, the sign-convention difference tends to zero
in \(L^1_\tau H^{-1}_y(B_R)\).
\end{proof}

\begin{lemma}[Terminal modulation compatibility]
\label{paper8:lem:terminal-modulation-compatibility}
Same-point larger-scale profiles in a terminal frame have the form
\[
  W_j^n(y,\tau)
  =
  \mu_j^n
  U^j(\theta_j^n(\tau),\xi_j^n+\mu_j^ny),
  \qquad \mu_j^n\to0.
\]
When the profile is regular at the observed offset, its leading term is a
spatially constant ambient velocity.  Absorbing this constant into the
translation coefficient replaces \(b_n\) by \(b_n-b_{{\rm amb},n}\), and the
remaining terms vanish on compact terminal cylinders.  If the offset escapes,
the entire profile is locally invisible and no absorption is needed.  In
particular,
\[
  \sup_n\sup_{\tau\in I}\bigl(|a_n(\tau)|+|b_n(\tau)|\bigr)<\infty .
\]
\end{lemma}

\begin{proof}
The local invisibility of the nonconstant remainder follows from
\cref{paper8:lem:terminal-velocity-decoupling} and the \(H^{-1}\) product
estimates in \cref{paper8:lem:terminal-linear-residual}.  Separated
physical-point profiles are locally invisible by
\cref{paper8:lem:separated-centers-vanish}; their pressures and cutoff
commutators vanish by \cref{paper8:lem:terminal-pressure-decoupling} and
\cref{paper8:lem:terminal-linear-residual}.  Therefore only the spatially
constant ambient velocity changes the translation coefficient, and after this
absorption the effective modulation coefficients remain uniformly bounded.
\end{proof}

\begin{lemma}[Terminal compact scale-collapse admissibility]
\label{paper8:lem:terminal-scale-collapse-admissibility}
The retained terminal solution \(U_n\) has positive finite critical mass,
pressure reconstruction, bounded modulation, local windowed \(H^1\) control,
critical tightness on the selected windows, and the compactness input used by
the compact scale-collapse argument.
\end{lemma}

\begin{proof}
The terminal class is active, so after the perturbative remainder and
nonretained profiles have been removed,
\[
  \inf_{\tau\in I}\|U_n(\tau)\|_{L^3(B_R)}
  \ge c_{\rm act}>0
\]
on the selected compact windows, after passing to the subsequences used in the
profile decomposition.  The upper bound follows from the retained local
compact-cylinder bound:
\[
  \sup_{\tau\in I}\|U_n(\tau)\|_{L^3(B_R)}\le C(M).
\]

The terminal profile completeness also gives the required tightness.  If
\(U_n\) were not uniformly \(L^3\)-tight on the selected windows, then there
would be \(\varepsilon>0\), radii \(R_k\to\infty\), and times \(\tau_n\) such
that
\[
  \int_{|y|>R_k}|U_n(y,\tau_n)|^3\,dy\ge\varepsilon.
\]
Applying the profile decomposition to this escaping sequence produces either an
exterior-regular component, which is removed by the exterior regularity
argument below, or an additional active profile at a different physical point
or larger same-point scale.  Such a profile belongs to \(S_n\) and is handled
by terminal decoupling.  Hence no non-tight retained terminal solution remains.

The repaired-gauge representation gives pressure reconstruction and bounded
modulation.  The local Caccioppoli estimate gives local windowed \(H^1\)
control, and the Aubin--Lions compactness layer gives the compactness input on
the selected windows.
\end{proof}

\begin{theorem}[Terminal nonlinear profile decoupling theorem]
\label{paper8:thm:terminal-nonlinear-decoupling}
Assume finite critical norm, terminal windowed profile completeness,
small-data perturbative stability, the exterior annulus alternative,
repaired-gauge representation, and local Caccioppoli regularity.  Then
terminal nonlinear profile decoupling holds in the sense of
\cref{paper8:def:terminal-nonlinear-decoupling}.
\end{theorem}

\begin{proof}
Let \(\mathcal C\) be any terminal active compound class, let \(U_n\) be the
retained repaired-gauge nonlinear solution in its terminal frame, and set
\[
  S_n=V_n-U_n.
\]
\Cref{paper8:lem:terminal-velocity-decoupling} proves
\[
  S_n\to0
  \quad\text{in }L^\infty_\tau L^3_y(B_{2R})
\]
on every compact terminal cylinder.  Divergence-free compatibility follows
from the divergence-free profile data and the divergence-free
Navier--Stokes evolution of each nonlinear profile.

\Cref{paper8:lem:terminal-stress-decoupling} proves
\[
  U_n\otimes S_n+S_n\otimes U_n+S_n\otimes S_n
  \to0
  \quad\text{in }L^1_\tau L^2_y(B_R).
\]
The finite critical norm and the retained profile bounds give the global
\(L^\infty_\tau L^3_y\) control required in the pressure lemma below.
\Cref{paper8:lem:terminal-pressure-decoupling} proves
\[
  \nabla P_{{\rm cross},n}\to0
  \quad\text{in }L^1_\tau H^{-1}_y(B_R)
\]
after subtracting spatially constant pressure terms.  \Cref{paper8:lem:terminal-linear-residual}
proves
\[
  \partial_\tau S_n-\nu\Delta S_n
  +a_n(S_n+y\cdot\nabla S_n)
  +b_n\cdot\nabla S_n
  \to0
  \quad\text{in }L^1_\tau H^{-1}_y(B_R).
\]
\Cref{paper8:lem:terminal-modulation-compatibility} gives modulation
compatibility, and \cref{paper8:lem:terminal-scale-collapse-admissibility}
gives compact scale-collapse admissibility of the retained terminal solution.
These are the items in \cref{paper8:def:terminal-nonlinear-decoupling}.
\end{proof}

\begin{corollary}[Multibubble exclusion from terminal profile decoupling]
\label{paper8:cor:multibubble-exclusion-terminal}
Under the Type II profile data, terminal windowed profile completeness,
small-data perturbative stability, the exterior annulus alternative, repaired-gauge
representation, local Caccioppoli regularity, and the compact scale-collapse
stationary-rigidity input, no active multibubble Type II configuration occurs.
\end{corollary}

\begin{proof}
\Cref{paper8:thm:terminal-nonlinear-decoupling} gives terminal nonlinear
profile decoupling.  \Cref{paper8:thm:terminal-decoupling-implies-local-decoupling}
then supplies the same-point and separated-point decoupling hypotheses used in
\cref{paper8:thm:multibubble-frame-reduction}.  The active-profile mass floor
follows from \cref{paper8:thm:active-profile-mass-floor} and the small-data
perturbative stability assumption.  Applying
\cref{paper8:thm:multibubble-frame-reduction} gives the claimed exclusion of
active multibubble configurations.
\end{proof}

\subsection{Small critical profiles}

For \(I=[0,T]\) or \(I=[0,\infty)\), define the Kato critical space \(X(I)\)
by
\[
  \|u\|_{X(I)}
  :=
  \sup_{t\in I}\|u(t)\|_{L^3_x}
  +
  \sup_{\substack{t\in I\\ t>0}}
  t^{1/8}\|u(t)\|_{L^4_x}.
\]
Let \(\mathbb P\) be the Leray projector and set
\[
  B(u,v)(t)
  :=
  -\int_0^t e^{\nu(t-s)\Delta}
  \mathbb P\nabla\cdot(u\otimes v)(s)\,ds .
\]
The critical mild formulation is
\[
  u(t)=e^{\nu t\Delta}u_0+B(u,u)(t).
\]

\begin{lemma}[Critical heat and bilinear estimates]
\label{paper8:lem:kato-estimates}
There are constants \(A_\nu,B_\nu<\infty\) such that for every divergence-free
\(u_0\in L^3(\mathbb R^3)\),
\[
  \|e^{\nu t\Delta}u_0\|_{X([0,\infty))}
  \le A_\nu\|u_0\|_{L^3},
\]
and for all \(u,v\in X([0,\infty))\),
\[
  \|B(u,v)\|_{X([0,\infty))}
  \le B_\nu\|u\|_{X([0,\infty))}\|v\|_{X([0,\infty))}.
\]
\end{lemma}

\begin{proof}
The \(L^\infty_tL^3_x\) heat estimate follows from heat semigroup
contractivity:
\[
  \|e^{\nu t\Delta}u_0\|_3\le\|u_0\|_3.
\]
The \(L^4\) Kato estimate follows from the heat kernel bound
\[
  \|e^{\nu t\Delta}f\|_q
  \le
  C_{\nu,p,q}t^{-\frac32(\frac1p-\frac1q)}\|f\|_p,
\]
with \(p=3\), \(q=4\):
\[
  \|e^{\nu t\Delta}u_0\|_4
  \le C_\nu t^{-1/8}\|u_0\|_3.
\]
Multiplying by \(t^{1/8}\) and taking the supremum gives the heat estimate in
\(X\).

For the bilinear term, the kernel of
\(e^{\nu(t-s)\Delta}\mathbb P\nabla\cdot\) satisfies
\[
  \|e^{\nu(t-s)\Delta}\mathbb P\nabla\cdot F\|_q
  \le
  C_{\nu,p,q}(t-s)^{-\frac12-\frac32(\frac1p-\frac1q)}\|F\|_p.
\]
With \(F=u\otimes v\), Holder's inequality gives
\[
  \|u(s)\otimes v(s)\|_2
  \le
  \|u(s)\|_4\|v(s)\|_4
  \le
  s^{-1/4}\|u\|_X\|v\|_X.
\]
Using \(p=2,q=3\),
\[
  \|B(u,v)(t)\|_3
  \le
  C\|u\|_X\|v\|_X
  \int_0^t(t-s)^{-3/4}s^{-1/4}\,ds
  \le
  C\|u\|_X\|v\|_X.
\]
Using \(p=2,q=4\),
\[
  \|B(u,v)(t)\|_4
  \le
  C\|u\|_X\|v\|_X
  \int_0^t(t-s)^{-7/8}s^{-1/4}\,ds
  \le
  Ct^{-1/8}\|u\|_X\|v\|_X.
\]
Multiplying by \(t^{1/8}\), taking the supremum in \(t\), and combining the
two components gives the bilinear estimate.
\end{proof}

\begin{theorem}[Small-data critical Navier--Stokes]
\label{paper8:thm:small-data-critical}
There exists \(\varepsilon_{\rm sd}>0\) such that if \(u_0\in L^3(\mathbb R^3)\)
is divergence-free and
\[
  \|u_0\|_3\le\varepsilon_{\rm sd},
\]
then Navier--Stokes has a unique global mild solution \(u\in X([0,\infty))\)
satisfying
\[
  \|u\|_{X([0,\infty))}\le 2A_\nu\|u_0\|_3.
\]
Moreover, for every finite interval \([0,T]\),
\[
  \|u-e^{\nu t\Delta}u_0\|_{X([0,T])}
  \le C_\nu\|u_0\|_3^2.
\]
\end{theorem}

\begin{proof}
Choose
\[
  \varepsilon_{\rm sd}\le\frac{1}{4A_\nu B_\nu}.
\]
Let
\[
  R=2A_\nu\|u_0\|_3,
\]
and consider the closed ball
\[
  \mathcal B_R=\{u\in X([0,\infty)):\|u\|_X\le R\}.
\]
Define
\[
  \Phi(u)=e^{\nu t\Delta}u_0+B(u,u).
\]
For \(u\in\mathcal B_R\), \cref{paper8:lem:kato-estimates} gives
\[
  \|\Phi(u)\|_X
  \le
  A_\nu\|u_0\|_3+B_\nu R^2.
\]
Since \(R=2A_\nu\|u_0\|_3\),
\[
  B_\nu R^2
  =
  4B_\nu A_\nu^2\|u_0\|_3^2
  \le A_\nu\|u_0\|_3
\]
by the choice of \(\varepsilon_{\rm sd}\).  Hence
\[
  \|\Phi(u)\|_X\le R.
\]
For \(u,v\in\mathcal B_R\),
\[
  \Phi(u)-\Phi(v)=B(u-v,u)+B(v,u-v),
\]
so
\[
  \|\Phi(u)-\Phi(v)\|_X
  \le
  B_\nu(\|u\|_X+\|v\|_X)\|u-v\|_X
  \le
  2B_\nu R\|u-v\|_X.
\]
The smallness condition gives
\[
  2B_\nu R=4A_\nu B_\nu\|u_0\|_3\le\frac12.
\]
Thus \(\Phi\) is a contraction on \(\mathcal B_R\), and Banach's fixed point
theorem gives a unique global mild solution with the asserted bound.

The perturbative estimate follows from
\[
  u-e^{\nu t\Delta}u_0=B(u,u)
\]
and \cref{paper8:lem:kato-estimates}:
\[
  \|u-e^{\nu t\Delta}u_0\|_X
  \le
  B_\nu\|u\|_X^2
  \le
  4B_\nu A_\nu^2\|u_0\|_3^2.
\]
\end{proof}

\begin{corollary}[Small profiles are perturbative]
\label{paper8:cor:small-profiles-perturbative}
If \(\phi\in L^3(\mathbb R^3)\) is divergence-free and
\[
  \|\phi\|_3\le\varepsilon_{\rm sd},
\]
then the nonlinear Navier--Stokes profile generated by \(\phi\) is global in
\(X([0,\infty))\), obeys the estimate of
\cref{paper8:thm:small-data-critical}, and is perturbative on every compact
renormalized window.  In particular, \(\phi\) cannot be an active Type II
profile.
\end{corollary}

\begin{proof}
\Cref{paper8:thm:small-data-critical} gives the whole-space mild solution
\(u^\phi\) with
\[
  \|u^\phi\|_{X([0,\infty))}
  \le
  2A_\nu\|\phi\|_3.
\]
The \(X\)-norm is invariant under Navier--Stokes scaling
\[
  u(x,t)\mapsto \lambda u(\lambda x,\lambda^2t).
\]
Therefore every rescaled copy of \(u^\phi\) remains uniformly small in the
same critical topology on compact windows.  The perturbative estimate shows
that its nonlinear contribution is controlled by the heat evolution plus an
\(O(\|\phi\|_3^2)\) error.  Such a profile is placed in the perturbative
remainder class and cannot be an active singular profile.
\end{proof}

\begin{theorem}[Positive critical mass floor for active profiles]
\label{paper8:thm:active-profile-mass-floor}
For every critical Navier--Stokes profile family satisfying the small-data
stability condition, every active profile satisfies
\[
  \phi^j\ \text{active}
  \quad\Longrightarrow\quad
  \|\phi^j\|_3>\varepsilon_{\rm sd}.
\]
\end{theorem}

\begin{proof}
Let \(\phi^j\) be a profile with
\[
  \|\phi^j\|_3\le\varepsilon_{\rm sd}.
\]
By \cref{paper8:cor:small-profiles-perturbative}, its nonlinear evolution is
global and perturbative in the critical stability space.  Active Type II
profiles are those not removed by the perturbative small-data class.
Therefore \(\phi^j\) is not active.  Taking the contrapositive gives the
displayed mass floor.
\end{proof}

\subsection{Exterior regularity}

\begin{definition}[Exterior annuli and terminal CKN quantity]
\label{paper8:def:exterior-ckn-quantity}
Let \((x_*,T^*)\) be the selected Type II core point.  For
\(0<r<R<\infty\), set
\[
  A_{r,R}:=\overline{B_R(x_*)}\setminus B_r(x_*).
\]
For \(x_0\in\mathbb R^3\) and \(\rho>0\), write
\[
  Q_\rho^*(x_0):=
  B_\rho(x_0)\times(T^*-\rho^2,T^*).
\]
For a suitable weak solution \((u,p)\), define the terminal local
velocity-pressure quantity
\[
  \mathcal C_*(u,p;x_0,\rho)
  :=
  \rho^{-2}
  \int_{T^*-\rho^2}^{T^*}\int_{B_\rho(x_0)} |u|^3\,dx\,dt
  +
  \rho^{-2}
  \int_{T^*-\rho^2}^{T^*}\int_{B_\rho(x_0)}
  |p-(p)_{B_\rho(x_0)}(t)|^{3/2}\,dx\,dt .
\]
The pressure is measured modulo functions of time.
\end{definition}

\begin{theorem}[Terminal Caffarelli--Kohn--Nirenberg regularity]
\label{paper8:thm:terminal-ckn-regularity}
There exists \(\varepsilon_{\rm CKN}>0\) with the following property.  Let
\((u,p)\) be suitable on \(Q_\rho^*(x_0)\).  If
\[
  \mathcal C_*(u,p;x_0,\rho)\le\varepsilon_{\rm CKN},
\]
then for each integer \(k\ge0\) there are \(0<c_k<1\) and \(C_k<\infty\), independent of
\((u,p),x_0,\rho\), such that \(u\) is smooth in
\[
  B_{c_k\rho}(x_0)\times(T^*-c_k\rho^2,T^*)
\]
and, after subtracting a function of time from \(p\),
\[
  \sup_{t\in(T^*-c_k\rho^2,T^*)}
  \left(
  \rho^{k+1}\|\nabla^k u(t)\|_{L^\infty(B_{c_k\rho}(x_0))}
  +
  \rho^{k+2}\|\nabla^k p(t)\|_{L^\infty(B_{c_k\rho}(x_0))}
  \right)
  \le C_k .
\]
\end{theorem}

\begin{proof}
This is the standard endpoint form of the
Caffarelli--Kohn--Nirenberg epsilon-regularity criterion
\cite{CaffarelliKohnNirenberg1982,Lin1998}, after translating the terminal
time to \(0\) and rescaling \(Q_\rho^*(x_0)\) to a unit parabolic cylinder.
The pressure normalization by spatial averages is the usual normalization in
the suitable weak-solution criterion.  Interior parabolic bootstrapping gives
the stated higher derivative bounds on smaller cylinders.
\end{proof}

Throughout the exterior regularity subsection we fix
\[
  \varepsilon_{\rm ext}:=\frac12\varepsilon_{\rm CKN}.
\]

\begin{definition}[No exterior concentration on a fixed annulus]
\label{paper8:def:no-exterior-concentration-stratum}
Let \(0<r<R<\infty\).  A suitable Type II branch has no exterior concentration on \(A_{r,R}\) if
there exists \(\rho_*>0\) such that whenever
\[
  x_0\in A_{r,R},\qquad 0<\rho\le\rho_*,
  \qquad Q_\rho^*(x_0)\subset
  (B_{2R}(x_*)\setminus \overline{B_{r/2}(x_*)})\times(0,T^*),
\]
one has
\[
  \mathcal C_*(u,p;x_0,\rho)\le\varepsilon_{\rm ext}.
\]
\end{definition}

\begin{definition}[Exterior concentration branch]
\label{paper8:def:exterior-concentration-branch}
An exterior concentration branch on \(A_{r,R}\) is a sequence
\[
  x_n\in A_{r,R},\qquad \rho_n\downarrow0,
\]
such that
\[
  Q_{\rho_n}^*(x_n)\subset
  (B_{2R}(x_*)\setminus \overline{B_{r/2}(x_*)})\times(0,T^*)
\]
and
\[
  \mathcal C_*(u,p;x_n,\rho_n)>\varepsilon_{\rm ext}.
\]
The associated exterior rescalings are
\[
  U_n(y,s):=\rho_n u(x_n+\rho_n y,T^*+\rho_n^2s),
  \qquad
  P_n(y,s):=\rho_n^2 p(x_n+\rho_n y,T^*+\rho_n^2s),
\]
for \(s<0\).
\end{definition}

\begin{lemma}[Exterior test dichotomy]
\label{paper8:lem:exterior-test-dichotomy}
Fix \(0<r<R<\infty\).  Then either the branch has no exterior concentration
on \(A_{r,R}\), or it admits an exterior concentration branch on \(A_{r,R}\).
\end{lemma}

\begin{proof}
If the no-concentration condition fails, then for every \(m\) there are
\[
  x_m\in A_{r,R},\qquad 0<\rho_m\le m^{-1},
\]
with \(Q_{\rho_m}^*(x_m)\) contained in the enlarged annular cylinder and
\[
  \mathcal C_*(u,p;x_m,\rho_m)>\varepsilon_{\rm ext}.
\]
After discarding repetitions and passing to a subsequence, \(\rho_m\to0\).
This is precisely an exterior concentration branch.  The converse is immediate
from the definitions.
\end{proof}

\begin{theorem}[Exterior concentration stratification]
\label{paper8:thm:exterior-stratification}
Fix \(0<r<R<\infty\).  A suitable Type II branch satisfies exactly one of the
following alternatives on \(A_{r,R}\):
\begin{enumerate}[label=\textup{(\roman*)}]
\item it has no exterior concentration on \(A_{r,R}\);
\item it admits an exterior concentration branch on \(A_{r,R}\).
\end{enumerate}
The second alternative is the complementary concentration branch.  According
to the relative position and scale of the exterior concentration sequence, it
belongs to one of the exterior-profile, separated-core, multicore, noncompact
exterior, or scale-collapse alternatives.
\end{theorem}

\begin{proof}
The dichotomy is \cref{paper8:lem:exterior-test-dichotomy}.  If an exterior
concentration branch exists, then the annular terminal cylinders carry a
positive Caffarelli--Kohn--Nirenberg quantity at scales \(\rho_n\downarrow0\)
and centers a positive physical distance from the selected core.  Let
\(\lambda_n\) denote the selected core scale along the same terminal sequence.
After passing to a subsequence, the relative scale-center configuration falls
into one of the standard local concentration-compactness alternatives.  If the
exterior frame is orthogonal to the selected core frame, it gives an exterior
profile.  If only finitely many separated active centers remain at comparable
terminal scales, it gives a separated-core or multicore configuration.  If
compactness fails in the exterior frame, one obtains the noncompact exterior
alternative.  If the exterior frame develops a further unresolved scale in its
own renormalized variables, one obtains the scale-collapse alternative.  These
alternatives exhaust the subsequential configurations at positive physical
distance from the selected core.  Thus the failure of the annular regularity
test is represented by a local concentration alternative rather than by an
additional regularity hypothesis.
\end{proof}

\begin{lemma}[Exterior CKN cover]
\label{paper8:lem:exterior-ckn-cover}
Assume there is no exterior concentration on \(A_{r,R}\).  Fix an integer
\(k\ge0\).  If
\[
  A_{r',R'}\Subset A_{r,R},
\]
then there are points \(x_j\in A_{r,R}\), radii \(\rho_j>0\), and a terminal
time length \(\delta>0\) such that
\[
  A_{r',R'}\times(T^*-\delta,T^*)
  \subset
  \bigcup_j
  B_{c_k\rho_j}(x_j)\times(T^*-c_k\rho_j^2,T^*)
\]
and
\[
  \mathcal C_*(u,p;x_j,\rho_j)\le\varepsilon_{\rm ext}
\]
for every \(j\).  Here \(c_k\) is the constant from
\cref{paper8:thm:terminal-ckn-regularity}.
\end{lemma}

\begin{proof}
Since \(A_{r',R'}\Subset A_{r,R}\), there is \(d>0\) such that
\[
  B_{4d}(x)\subset B_{2R}(x_*)\setminus \overline{B_{r/2}(x_*)}
  \qquad\text{for all }x\in A_{r',R'}.
\]
Let \(\rho_*\) be supplied by the no-concentration condition and set
\(\rho_0:=\min\{d,\rho_*\}\).  By compactness, finitely many balls
\[
  B_{c_k\rho_0/2}(x_j),\qquad x_j\in A_{r,R},
\]
cover \(A_{r',R'}\).  The terminal cylinders \(Q_{\rho_0}^*(x_j)\) are
contained in the enlarged annular cylinder, so the no-concentration condition
gives the displayed smallness.  Taking
\(\delta=c_k\rho_0^2/2\) gives the spacetime cover.
\end{proof}

\begin{lemma}[Exterior pressure localization]
\label{paper8:lem:exterior-pressure-localization}
Let \(A_{r',R'}\Subset A_{r,R}\), and assume that \(u\) is bounded on a
slightly larger annular terminal cylinder.  Then on
\[
  A_{r',R'}\times(T^*-\delta,T^*)
\]
the pressure has a decomposition
\[
  p=p_{\rm loc}+p_{\rm harm}+c(t),
\]
where \(c(t)\) is a function of time, \(p_{\rm harm}\) is harmonic in space on
the smaller annulus, and for every \(q<\infty\) and every integer \(k\ge0\)
\[
  \|p_{\rm loc}\|_{L^\infty_tL^q_x(A_{r',R'})}
  +
  \|p_{\rm harm}\|_{L^\infty_tC^k_x(A_{r',R'})}
  \le C_{r',R',q,k}.
\]
\end{lemma}

\begin{proof}
Choose \(\eta\in C_c^\infty(B_{2R}(x_*)\setminus \overline{B_{r/2}(x_*)})\)
with \(\eta\equiv1\) on a slightly larger annulus containing \(A_{r',R'}\).
Define \(p_{\rm loc}\) by
\[
  -\Delta p_{\rm loc}
  =
  \partial_i\partial_j(\eta u_i u_j)
\]
using the Newtonian potential, with the usual spatial normalization.  Since
\(u\) is locally bounded on the support of \(\eta\), Calderon--Zygmund
estimates give \(p_{\rm loc}\in L^\infty_tL^q_x\) for every \(q<\infty\) on
the terminal interval.  The difference
\[
  p-p_{\rm loc}
\]
is harmonic in space on the smaller annulus, modulo a function of time.  The
standard interior estimates for harmonic functions give the displayed
\(C^k\) bounds after fixing the time-dependent additive constant.
\end{proof}

\begin{theorem}[Exterior regularity on no-concentration annuli]
\label{paper8:thm:exterior-regularity-no-concentration}
Let \(u,p\) be a suitable weak solution on
\(\mathbb R^3\times[0,T^*)\), and let \((x_*,T^*)\) be the selected Type II
core point.  Fix \(0<r'<R'<\infty\).  If there is no exterior concentration on
some annulus \(A_{r,R}\) with
\[
  A_{r',R'}\Subset A_{r,R},
\]
then for every integer \(k\ge0\) there are \(\delta>0\), a function
\(c(t)\), and \(C_{r',R',k}<\infty\) such that
\[
  \sup_{t\in(T^*-\delta,T^*)}
  \left(
  \|u(t)\|_{C^k(A_{r',R'})}
  +
  \|p(t)-c(t)\|_{C^k(A_{r',R'})}
  \right)
  \le C_{r',R',k}.
\]
If the no-concentration hypothesis fails, the complementary exterior
concentration alternative holds in the sense of
\cref{paper8:def:exterior-concentration-branch}.
\end{theorem}

\begin{proof}
The no-concentration condition and \cref{paper8:lem:exterior-ckn-cover} give a
finite cover of \(A_{r',R'}\times(T^*-\delta,T^*)\) by terminal cylinders on
which the Caffarelli--Kohn--Nirenberg quantity is below
\(\varepsilon_{\rm CKN}\).  Applying
\cref{paper8:thm:terminal-ckn-regularity} on each cylinder gives local
\(C^k\) bounds for \(u\) and for \(p\) modulo functions of time on smaller
cylinders.  Since the cover is finite, taking the maximum over the cover gives
the stated local bound.
\par
The pressure normalizations on overlapping balls differ only by functions of
time.  Equivalently, one may use
\cref{paper8:lem:exterior-pressure-localization} after the \(k=0\) velocity
bound has been obtained and then apply standard interior parabolic estimates on
nested annuli.  This gives a single time-dependent normalization \(c(t)\) on
\(A_{r',R'}\).  If the no-concentration hypothesis fails, the exterior test
dichotomy gives an exterior concentration branch.
\end{proof}

\begin{lemma}[Divergence-free core and exterior decomposition]
\label{paper8:lem:core-exterior-decomposition}
Fix \(0<r_0<R_0<\infty\), and assume that there is no exterior concentration
on an annulus containing
\[
  \overline{B_{2r_0}(x_*)}\setminus B_{r_0}(x_*).
\]
Let \(\chi\in C_c^\infty(B_{2r_0}(x_*))\) satisfy
\[
  \chi\equiv1\quad\text{on }B_{r_0}(x_*).
\]
Then there are divergence-free fields \(u_{\rm core}\) and \(u_{\rm ext}\)
such that
\[
  u=u_{\rm core}+u_{\rm ext},\qquad
  u_{\rm core}=u\ \text{on }B_{r_0}(x_*),\qquad
  u_{\rm ext}=0\ \text{on }B_{r_0}(x_*),
\]
and all derivatives of \(u_{\rm ext}\) are uniformly bounded on compact
subannuli of the no-exterior-concentration region up to \(T^*\).
\end{lemma}

\begin{proof}
Set
\[
  g=\nabla\chi\cdot u.
\]
The support of \(g\) lies in the annulus
\[
  A_{r_0}=B_{2r_0}(x_*)\setminus B_{r_0}(x_*),
\]
where \(u\) is smooth by
\cref{paper8:thm:exterior-regularity-no-concentration}.  Since
\[
  \int_{A_{r_0}}g\,dx
  =
  \int_{\mathbb R^3}\nabla\chi\cdot u\,dx
  =
  -\int_{\mathbb R^3}\chi\,\nabla\cdot u\,dx
  =0,
\]
the Bogovskii operator on \(A_{r_0}\) gives a vector field
\[
  w=\mathcal B_{A_{r_0}}g
\]
supported in \(A_{r_0}\), smooth up to \(T^*\), and satisfying
\[
  \nabla\cdot w=g.
\]
Define
\[
  u_{\rm core}:=\chi u-w,
  \qquad
  u_{\rm ext}:=u-u_{\rm core}.
\]
Then
\[
  \nabla\cdot u_{\rm core}
  =
  \nabla\chi\cdot u+\chi\nabla\cdot u-\nabla\cdot w=0,
\]
and therefore \(\nabla\cdot u_{\rm ext}=0\).  Since \(\chi=1\) and \(w=0\)
on \(B_{r_0}(x_*)\), \(u_{\rm core}=u\) and \(u_{\rm ext}=0\) there.  The
uniform exterior bounds follow from
\cref{paper8:thm:exterior-regularity-no-concentration} on compact subannuli
and the standard boundedness of the Bogovskii construction on the fixed
annulus.
\end{proof}

\begin{lemma}[Exterior fields vanish in compact core frames]
\label{paper8:lem:exterior-velocity-vanishes}
Let \((x_n,\lambda_n)\) be a core frame with
\[
  x_n\to x_*,
  \qquad
  \lambda_n\to0.
\]
For a compact frame ball \(B_R\), define
\[
  V_{\rm ext}^n(y,\tau)
  :=
  \lambda_nu_{\rm ext}(x_n+\lambda_ny,t_n+\lambda_n^2\tau).
\]
Then, for every compact \(B_R\times I\),
\[
  V_{\rm ext}^n\equiv0
\]
on \(B_R\times I\) for all sufficiently large \(n\).  Consequently
\[
  V_{\rm ext}^n\to0
  \quad\text{in every local }L^q_\tau W^{k,q}_y(B_R)\text{ norm}.
\]
\end{lemma}

\begin{proof}
Since \(x_n\to x_*\) and \(\lambda_n\to0\), for fixed \(R\),
\[
  x_n+\lambda_nB_R\subset B_{r_0}(x_*)
\]
for all sufficiently large \(n\).  On \(B_{r_0}(x_*)\),
\[
  u_{\rm ext}=0.
\]
Therefore \(V_{\rm ext}^n=0\) on \(B_R\times I\) for all large \(n\), which
implies convergence in every stated local norm.
\end{proof}

\begin{lemma}[Exterior pressure vanishes modulo constants]
\label{paper8:lem:exterior-pressure-vanishes}
Let \(p_{\rm ext}\) be any pressure contribution generated by source terms
supported in
\[
  \{|x-x_*|\ge r_0\}.
\]
In a core frame set
\[
  P_{\rm ext}^n(y,\tau)
  :=
  \lambda_n^2p_{\rm ext}(x_n+\lambda_ny,t_n+\lambda_n^2\tau).
\]
Then, for every compact \(B_R\times I\), there are constants \(c_{n,R}(\tau)\)
such that
\[
  P_{\rm ext}^n-c_{n,R}
  \to0
  \quad\text{in }L^\infty_\tau C^1_y(B_R),
\]
and hence
\[
  \nabla_yP_{\rm ext}^n\to0
  \quad\text{in }L^1_\tau H^{-1}_y(B_R).
\]
\end{lemma}

\begin{proof}
The exterior pressure source is at positive physical distance from
\[
  x_n+\lambda_nB_R
\]
for all sufficiently large \(n\).  Therefore \(p_{\rm ext}\) is harmonic,
smooth, and uniformly \(C^2\)-bounded on the physical neighborhood swept out
by the compact frame cylinder.  Take
\[
  c_{n,R}(\tau)
  :=
  \lambda_n^2p_{\rm ext}(x_n,t_n+\lambda_n^2\tau).
\]
By the mean value theorem,
\[
  \|P_{\rm ext}^n-c_{n,R}\|_{C^1_y(B_R)}
  \le
  C_R\lambda_n^3\sup|\nabla_xp_{\rm ext}|
  +
  C_R\lambda_n^4\sup|\nabla_x^2p_{\rm ext}|.
\]
The right-hand side tends to zero because the physical derivatives are
uniformly bounded and \(\lambda_n\to0\).  This proves the asserted
convergence.
\end{proof}

\begin{lemma}[Positive-distance profiles vanish in compact core frames]
\label{paper8:lem:positive-distance-exterior}
Let \((x_n,\lambda_n)\) be a core frame with \(x_n\to x_*\) and
\(\lambda_n\to0\).  Let \((z_n,\mu_n)\) be a profile sequence satisfying
\[
  \liminf_{n\to\infty}|z_n-x_*|\ge r_*>0.
\]
Assume moreover that either \(|z_n-x_*|\to\infty\), or, after passing to a
subsequence, \(z_n\) lies in a fixed annulus \(A_{r_*/2,R_*}\) on which the
no-exterior-concentration condition holds.
Then this profile vanishes in every compact core frame:
\[
  \Lambda_{\lambda_n,x_n}^{-1}\Lambda_{\mu_n,z_n}\phi
  \to0
\]
locally in \(L^3_y(B_R)\).  In the escaping-center case its pressure
contribution vanishes in the corresponding local \(L^{3/2}\) pressure
topology by critical profile orthogonality.  In the fixed-annulus
no-concentration case its pressure contribution vanishes modulo constants as
in \cref{paper8:lem:exterior-pressure-vanishes}.  Hence it is not part of the
selected Type II core.
\end{lemma}

\begin{proof}
For fixed \(R<\infty\), the physical frame ball
\[
  x_n+\lambda_nB_R
\]
lies in \(B_{r_*/4}(x_*)\) for all sufficiently large \(n\), while the profile
center \(z_n\) lies outside \(B_{r_*/2}(x_*)\).  Thus the profile is a
positive physical distance from the compact core frame.

If \(|z_n-x_*|\to\infty\), then the profile frame is orthogonal to the compact
core frame.  More explicitly, approximate the profile in \(L^3(\mathbb R^3)\)
by a compactly supported smooth field \(\phi^\varepsilon\).  The rescaled
compact part has support disjoint from \(x_n+\lambda_nB_R\) for all large
\(n\), while the remainder has \(L^3\)-norm at most \(\varepsilon\) in every
critical rescaling.  Letting first \(n\to\infty\) and then
\(\varepsilon\downarrow0\) gives local \(L^3_y(B_R)\) convergence to zero.
The same truncation argument applied to the quadratic source, together with
the local Calderon--Zygmund estimate for pressure modulo functions of time,
gives convergence of the pressure contribution to zero in local
\(L^{3/2}\).

Otherwise, after passing to a subsequence, the centers lie in a fixed annulus
on which the no-exterior-concentration condition holds.  The exterior velocity
part is covered by \cref{paper8:lem:exterior-velocity-vanishes} after the
core/exterior decomposition, and the exterior pressure part is covered by
\cref{paper8:lem:exterior-pressure-vanishes}.  This proves the claimed
decoupling of positive-distance profiles from compact Type II core frames.
\end{proof}

\begin{theorem}[Exterior regular components decouple from compact core frames]
\label{paper8:thm:exterior-regular-removal}
On the no-exterior-concentration stratum for the fixed annuli separating an
exterior component from the selected Type II core point, that component is a
regular exterior component and vanishes in compact Type II core frames,
including its pressure contribution modulo constants.
\end{theorem}

\begin{proof}
Apply the divergence-free decomposition of
\cref{paper8:lem:core-exterior-decomposition}.  The exterior component is
identically zero in every compact core frame for large \(n\) by
\cref{paper8:lem:exterior-velocity-vanishes}.  Its pressure contribution is a
locally scaled smooth harmonic field and vanishes modulo constants by
\cref{paper8:lem:exterior-pressure-vanishes}.  The cutoff and Bogovskii
correction terms are supported in the fixed exterior annulus \(A_{r_0}\), so
they also vanish in compact core frames by the same argument.

Thus a regular exterior component decouples from the Type II core without
changing the local repaired-gauge equation in any compact terminal frame, up
to errors converging to zero in the perturbative topology used in
\cref{paper8:lem:local-perturbation} and
\cref{paper8:def:terminal-nonlinear-decoupling}.
\end{proof}

\begin{corollary}[Exterior annulus alternative for the core decomposition]
\label{paper8:cor:exterior-annulus-alternative}
Fix a compact annulus \(A_{r,R}\) around the selected core.  Then exactly one
of the following alternatives holds for the Type II branch on this annulus:
\begin{enumerate}[label=\textup{(\roman*)}]
\item the branch has no exterior concentration on \(A_{r,R}\), and every
compact subannulus of \(A_{r,R}\) is covered by
\cref{paper8:thm:exterior-regularity-no-concentration}; consequently exterior
regular contributions from this annulus vanish in compact core frames
by \cref{paper8:thm:exterior-regular-removal};
\item the branch admits an exterior concentration branch on \(A_{r,R}\), and
therefore belongs to one of the adjacent exterior, separated-core, multicore,
noncompact exterior, or scale-collapse alternatives described in
\cref{paper8:thm:exterior-stratification}.
\end{enumerate}
\end{corollary}

\begin{proof}
This is the exterior concentration stratification of
\cref{paper8:thm:exterior-stratification}, combined in the first alternative
with the no-concentration regularity theorem and the exterior-removal theorem.
\end{proof}

\clearpage
\section{Corrected Monotonicity and Scale-Drift Cost Criteria}
\label{sec:corrected-monotonicity-scale-drift}

This section records the remaining conditional analytic inputs needed for a
localized cost-divergence exclusion argument on renormalized Type II branches.
The statements are written entirely in the variables of the repaired gauge:
\[
  \partial_\tau V+(V\cdot\nabla)V+\nabla P
  =
  \nu\Delta V+a(\tau)(V+y\cdot\nabla V)+b(\tau)\cdot\nabla V,
  \qquad \nabla\cdot V=0.
\]
All pressure terms are understood modulo functions of time.

\subsection{Hypotheses for a Conditional Cost Criterion}

\begin{definition}[Conditional cost hypotheses]\label{paper6a:def:cost-hypotheses}
A renormalized local Type II branch satisfies the conditional cost hypotheses
if the following analytic facts hold.
\begin{enumerate}[label=\textup{(\roman*)}]
\item the branch lies in the local Type II concentration alternative;
\item the physical Navier--Stokes energy inequality gives bounded energy along
the branch;
\item the branch is written in the repaired-gauge variables
\((V,P,a,b)\) above, including the scale, center, pressure normalization, and
modulation data;
\item the localized cost is identified with
\[
  \tilde{\mathfrak D}_{R_0}(\tau)
  =
  \nu\int_{\mathbb R^3}|\nabla V(y,\tau)|^2\phi_{R_0}(y)\,dy
  +a_+(\tau)\int_{\mathbb R^3}|V(y,\tau)|^2\phi_{R_0}(y)\,dy,
\]
or with an explicitly stated tail-comparable nonnegative replacement;
\item a localized energy identity or suitable-solution local energy inequality
is available for the cutoff used in the cost;
\item a corrected localized monotonicity inequality, in the sense of
\cref{paper6a:def:corrected-monotonicity-input}, is available.
\end{enumerate}
\end{definition}

\begin{assumption}[Cost-divergence exclusion criterion]
\label{paper6a:ass:cost-divergence-exclusion}
For every renormalized local Type II branch satisfying
\cref{paper6a:def:cost-hypotheses}, divergence of the corresponding
nonnegative localized cost,
\[
  \int_{\tau_0}^{\infty}\mathcal D_{R_0}(\tau)\,d\tau=\infty,
\]
excludes that branch from the local Type II class under consideration.
\end{assumption}

\begin{lemma}[Concentration and energy]\label{paper6a:lem:scale-energy-data}
If a branch lies in the local Type II concentration alternative and satisfies
the physical Navier--Stokes energy inequality, then it satisfies
\cref{paper6a:def:cost-hypotheses}\textup{(i)} and
\textup{(ii)}.
\end{lemma}

\begin{proof}
The local Type II concentration alternative supplies the selected concentration
branch.  The physical energy inequality gives the bounded-energy part.  These
are precisely items \textup{(i)} and \textup{(ii)}.
\end{proof}

\begin{lemma}[Cost identification]\label{paper6a:lem:cost-identification}
Assume the localized cost is the Navier--Stokes cost
\(\tilde{\mathfrak D}_{R_0}\).  Then divergence of
\[
  \int_{\tau_0}^{\infty}\tilde{\mathfrak D}_{R_0}(\tau)\,d\tau
\]
is the cost divergence in \cref{paper6a:ass:cost-divergence-exclusion}.
\end{lemma}

\begin{proof}
The active cost is assumed to be equal to
\(\tilde{\mathfrak D}_{R_0}\).  Thus the cost appearing in the local
Navier--Stokes estimates and the cost appearing in the conditional exclusion
criterion are the same nonnegative functional.  The divergence condition is
therefore identical in the two formulations.
\end{proof}

\begin{lemma}[Parabolic representation class]\label{paper6a:lem:parabolic-class}
If the repaired-gauge representation theorem applies to a Type II branch, then
the branch has the parabolic transport-diffusion structure required in
\cref{paper6a:def:cost-hypotheses}.
\end{lemma}

\begin{proof}
The repaired-gauge theorem supplies the renormalized orbit \((V,P,a,b)\), the
scale and center, the pressure reconstruction modulo functions of time, and the
modulation coefficients.  The displayed equation is a parabolic
Navier--Stokes equation with transport, dilation, and translation drift terms.
Consequently the branch has the parabolic transport-diffusion structure required by
the conditional cost hypotheses.
\end{proof}

\begin{definition}[Corrected localized monotonicity input]
\label{paper6a:def:corrected-monotonicity-input}
A renormalized branch has corrected localized monotonicity if there are a
localized renormalized energy \(\mathcal E_{R_0}\), a nonnegative cost
\(\mathcal D_{R_0}\), a constant \(c_0>0\), and a nonnegative remainder
\(\mathcal R_{R_0}\) such that, on the tail of the branch,
\[
  \frac{d}{d\tau}\mathcal E_{R_0}(\tau)
  +c_0\mathcal D_{R_0}(\tau)
  +\mathcal R_{R_0}(\tau)\le0
\]
in distributions.  If the localized energy identity has an integrable error
\(B_{R_0}\), the corrected energy is
\[
  \mathcal E_{R_0}^{\mathrm{corr}}(\tau)
  =
  \mathcal E_{R_0}(\tau)
  +\int_\tau^\infty B_{R_0}(s)\,ds,
\]
and the same inequality is required for
\(\mathcal E_{R_0}^{\mathrm{corr}}\).  Multiplication of a nonnegative cost by
\(c_0>0\) does not change whether its tail integral diverges.
\end{definition}

\begin{lemma}[Energy and Caccioppoli bounds alone are not monotonicity]
\label{paper6a:lem:caccioppoli-not-monotonicity}
The global physical energy inequality and compact-cylinder Caccioppoli
estimates do not, by themselves, imply the corrected localized monotonicity
input of \cref{paper6a:def:corrected-monotonicity-input}.
\end{lemma}

\begin{proof}
The global physical energy inequality controls total physical energy and
physical dissipation, but it does not identify a localized renormalized energy
at the selected Type II scale whose derivative is monotone after pressure flux,
cutoff flux, transport, and modulation terms are included.  The Caccioppoli
estimate gives local spacetime \(H^1\) control from the local energy inequality.
It is an a priori estimate, not a corrected monotonicity formula.  Therefore an
additional absorption or tail-correction argument is needed.
\end{proof}

\begin{proposition}[Conditional cost-divergence exclusion]\label{paper6a:prop:cost-divergence-exclusion}
Let a renormalized Type II branch satisfy
\cref{paper6a:def:cost-hypotheses} and
\cref{paper6a:ass:cost-divergence-exclusion}.  If its nonnegative localized
cost has divergent tail integral, then the branch is excluded from the local
Type II class under consideration.
\end{proposition}

\begin{proof}
\Cref{paper6a:lem:scale-energy-data} gives the local Type II concentration and
bounded-energy inputs.  \Cref{paper6a:lem:cost-identification} identifies the
localized Navier--Stokes cost with the cost used in the exclusion criterion.
\Cref{paper6a:lem:parabolic-class} gives the parabolic transport-diffusion class.
The corrected monotonicity input is an explicit hypothesis of
\cref{paper6a:def:cost-hypotheses}.  Once the cost diverges,
\cref{paper6a:ass:cost-divergence-exclusion} gives the stated exclusion.
\end{proof}

\subsection{Fixed-Cutoff Monotonicity}

Let \(\phi=\phi_{R_0}\in C_c^\infty(\mathbb R^3)\), \(0\le\phi\le1\), be the
radial nonincreasing cutoff used in the localized cost.  Define
\[
  \mathcal E_\phi(\tau)=\frac12\int_{\mathbb R^3}|V(y,\tau)|^2\phi(y)\,dy,
\]
\[
  A_\phi(\tau)=\int_{\mathbb R^3}|\nabla V(y,\tau)|^2\phi(y)\,dy,\qquad
  M_\phi(\tau)=\int_{\mathbb R^3}|V(y,\tau)|^2\phi(y)\,dy,
\]
and
\[
  \tilde{\mathfrak D}_{R_0}(\tau)=\nu A_\phi(\tau)+a_+(\tau)M_\phi(\tau),
  \qquad a_+(\tau)=\max\{a(\tau),0\}.
\]
The localized energy identity used below is
\begin{align}
  \frac{d}{d\tau}\mathcal E_\phi
  ={}&-\nu A_\phi
  +\frac{\nu}{2}\int |V|^2\Delta\phi
  +\frac12\int |V|^2V\cdot\nabla\phi
  +\int P\,V\cdot\nabla\phi \notag\\
  &-\frac{a}{2}M_\phi
  -\frac{a}{2}\int |V|^2y\cdot\nabla\phi
  -\frac12 b\cdot\int |V|^2\nabla\phi .
  \label{paper6a:eq:fixed-local-energy}
\end{align}

\begin{definition}[Fixed-cutoff error density]\label{paper6a:def:fixed-error-density}
Set \(a_-(\tau)=\max\{-a(\tau),0\}\) and define
\begin{align*}
  B_{R_0}(\tau)
  :={}&
  \frac{\nu}{2}\left|\int |V|^2\Delta\phi\right|
  +\frac12\left|\int |V|^2V\cdot\nabla\phi\right|
  +\left|\int P\,V\cdot\nabla\phi\right|\\
  &+\frac12|b(\tau)|\left|\int |V|^2\nabla\phi\right|
  +\frac12a_-(\tau)M_\phi(\tau)
  +\frac12a_+(\tau)\int |V|^2\bigl(-y\cdot\nabla\phi\bigr).
\end{align*}
The fixed-cutoff finite-error condition is
\[
  \int_{\tau_0}^{\infty}B_{R_0}(\tau)\,d\tau<\infty .
\]
\end{definition}

\begin{remark}[Pressure normalization]
The pressure term in \(B_{R_0}\) is independent of the representative of the
pressure modulo functions of time.  Replacing \(P\) by \(P+c(\tau)\) changes
the pressure flux by
\[
  c(\tau)\int V\cdot\nabla\phi
  =
  -c(\tau)\int \phi\,\nabla\cdot V
  =0.
\]
\end{remark}

\begin{theorem}[Fixed-cutoff corrected monotonicity]
\label{paper6a:thm:fixed-corrected-monotonicity}
Assume \eqref{paper6a:eq:fixed-local-energy} holds on
\((\tau_0,\infty)\) and that \(B_{R_0}\in L^1(\tau_0,\infty)\).  Define
\[
  \mathcal E_\phi^{\mathrm{corr}}(\tau)
  =
  \mathcal E_\phi(\tau)+\int_\tau^\infty B_{R_0}(s)\,ds .
\]
Then, in distributions on \((\tau_0,\infty)\),
\[
  \frac{d}{d\tau}\mathcal E_\phi^{\mathrm{corr}}(\tau)
  +\frac12\tilde{\mathfrak D}_{R_0}(\tau)\le0 .
\]
\end{theorem}

\begin{proof}
Add
\[
  \frac12\tilde{\mathfrak D}_{R_0}
  =
  \frac{\nu}{2}A_\phi+\frac12a_+M_\phi
\]
to both sides of \eqref{paper6a:eq:fixed-local-energy}.  The gradient term
becomes
\[
  -\nu A_\phi+\frac{\nu}{2}A_\phi=-\frac{\nu}{2}A_\phi\le0.
\]
The core scale term satisfies
\[
  -\frac{a}{2}M_\phi+\frac12a_+M_\phi
  =
  \frac12a_-M_\phi
  \le B_{R_0}(\tau).
\]
Since \(\phi\) is radial nonincreasing, one has
\[
  y\cdot\nabla\phi(y)\le0
  \qquad\text{for all }y\in\mathbb R^3.
\]
Hence
\[
  -\frac{a}{2}\int |V|^2y\cdot\nabla\phi
  \le
  \frac12a_+(\tau)\int |V|^2\bigl(-y\cdot\nabla\phi\bigr),
  \]
because the \(a_-\)-part contributes a nonpositive term.  Every remaining term
in \eqref{paper6a:eq:fixed-local-energy} is bounded above by its corresponding
absolute-value contribution in \(B_{R_0}\).  Therefore
\[
  \frac{d}{d\tau}\mathcal E_\phi(\tau)
  +\frac12\tilde{\mathfrak D}_{R_0}(\tau)
  \le B_{R_0}(\tau)
\]
in distributions.  Since \(B_{R_0}\in L^1(\tau_0,\infty)\), the correction
tail is finite for every \(\tau\ge\tau_0\), absolutely continuous on compact
subintervals, and
\[
  \frac{d}{d\tau}\int_\tau^\infty B_{R_0}(s)\,ds=-B_{R_0}(\tau).
\]
Adding this identity gives the claimed inequality.
\end{proof}

\begin{corollary}[Fixed-cutoff monotonicity input]\label{paper6a:cor:fixed-input}
If the repaired-gauge representation, the cost identification, the localized
energy identity, and the fixed-cutoff finite-error condition all hold, then
\cref{paper6a:def:corrected-monotonicity-input} holds with
\(\mathcal D_{R_0}=\tilde{\mathfrak D}_{R_0}\) and \(c_0=1/2\).
\end{corollary}

\begin{proof}
The repaired-gauge representation supplies \((V,P,a,b)\).  The cost
identification supplies \(\tilde{\mathfrak D}_{R_0}\).  The localized energy
identity and finite-error condition allow
\cref{paper6a:thm:fixed-corrected-monotonicity} to be applied.
\end{proof}

\begin{remark}[Finite-tail subconditions]\label{paper6a:rem:fixed-tail-subconditions}
The fixed-cutoff finite-error condition is equivalent to the conjunction of tail
integrability for the viscous cutoff term, the convective flux, the pressure
flux, the center-drift cutoff term, the negative-scale term, and the
positive scale-shell term:
\[
\begin{gathered}
  \left|\int |V|^2\Delta\phi\right|,\qquad
  \left|\int |V|^2V\cdot\nabla\phi\right|,\qquad
  \left|\int P\,V\cdot\nabla\phi\right|,\\
  |b|\left|\int |V|^2\nabla\phi\right|,\qquad
  a_-M_\phi,\qquad
  a_+\int |V|^2\bigl(-y\cdot\nabla\phi\bigr).
\end{gathered}
\]
For radial nonincreasing cutoffs, only the positive part \(a_+\) contributes to
the scale-shell error.  The negative part is favorable and does not belong to
the residual finite-error burden.
Uniform compact-window bounds do not by themselves give this global tail
integrability on \([\tau_0,\infty)\).
\end{remark}

\begin{definition}[Fixed-cutoff cost hypotheses]\label{paper6a:def:direct-fixed-data}
A renormalized Type II branch satisfies the fixed-cutoff cost hypotheses if it
has the local Type II concentration alternative, bounded physical energy,
repaired-gauge representation, cost identification with
\(\tilde{\mathfrak D}_{R_0}\), divergent nonnegative localized cost, localized
energy identity, and the finite-tail bound
\(\int_{\tau_0}^{\infty}B_{R_0}<\infty\).
\end{definition}

\begin{theorem}[Fixed-cutoff cost criterion]\label{paper6a:thm:fixed-cost-criterion}
Every branch satisfying the fixed-cutoff cost hypotheses is covered by
\cref{paper6a:prop:cost-divergence-exclusion}.
\end{theorem}

\begin{proof}
By \cref{paper6a:cor:fixed-input}, the finite-tail condition supplies the
corrected monotonicity input.  The remaining assumptions are precisely the other
items in \cref{paper6a:def:cost-hypotheses}, together with cost
divergence and \cref{paper6a:ass:cost-divergence-exclusion}.
\end{proof}

\begin{definition}[Compact cost-divergence replacement]\label{paper6a:def:compact-cost-replacement}
A compact cost-divergence replacement is available when the compact Type II hypotheses
of the preceding sections imply
\[
  \int_{\tau_0}^{\infty}\tilde{\mathfrak D}_{R_0}(\tau)\,d\tau=\infty,
\]
and this divergence is the cost divergence used in
\cref{paper6a:prop:cost-divergence-exclusion}.  In particular, it may be
supplied by the combination of classification of local alternatives, uniform
critical tightness, local windowed \(H^1\) control, and the cost identification
above whenever those are the active hypotheses in the compact Type II theorem.
\end{definition}

\begin{theorem}[Compact cost-divergence criterion]\label{paper6a:thm:compact-cost-criterion}
If a renormalized branch has all fixed-cutoff cost hypotheses except that cost
divergence is supplied through the compact cost-divergence replacement of
\cref{paper6a:def:compact-cost-replacement}, then
\cref{paper6a:prop:cost-divergence-exclusion} applies.
\end{theorem}

\begin{proof}
The compact cost-divergence replacement supplies the missing divergent cost.
Substituting this into the fixed-cutoff hypotheses gives the hypotheses of
\cref{paper6a:thm:fixed-cost-criterion}.
\end{proof}

\subsection{Moving-Cutoff Monotonicity}

Let \(\phi\in C_c^\infty(\mathbb R^3)\) be radial and nonincreasing,
\(0\le\phi\le1\), \(\phi=1\) on \(B_1\), and \(\phi=0\) outside \(B_2\).  For
\(R(\tau)\ge R_0\), locally absolutely continuous and nondecreasing, set
\[
  \Phi(\tau,y)=\phi(y/R(\tau)),\qquad
  A_{R(\tau)}=\{R(\tau)\le |y|\le2R(\tau)\},
\]
\[
  \mathcal E_R(\tau)=\frac12\int |V(y,\tau)|^2\Phi(\tau,y)\,dy,
\]
and
\[
  \tilde{\mathfrak D}_{R(\tau)}(\tau)
  =
  \nu\int|\nabla V|^2\Phi
  +a_+(\tau)\int|V|^2\Phi .
\]

\begin{definition}[Moving-cutoff error density]\label{paper6a:def:moving-error-density}
Define
\begin{align*}
  B_{\mathrm{mov}}(\tau)
  :={}&
  \frac12\left|\int |V|^2\partial_\tau\Phi\right|
  +\frac{\nu}{2}\left|\int |V|^2\Delta\Phi\right|
  +\frac12\left|\int |V|^2V\cdot\nabla\Phi\right|
  +\left|\int P\,V\cdot\nabla\Phi\right|\\
  &+\frac12|b(\tau)|\left|\int |V|^2\nabla\Phi\right|
  +\frac12a_-(\tau)\int |V|^2\Phi
  +\frac12a_+(\tau)\int |V|^2\bigl(-y\cdot\nabla\Phi\bigr).
\end{align*}
The moving finite-error condition is
\[
  \int_{\tau_0}^{\infty}B_{\mathrm{mov}}(\tau)\,d\tau<\infty .
\]
\end{definition}

\begin{theorem}[Moving-cutoff corrected monotonicity]
\label{paper6a:thm:moving-corrected-monotonicity}
Assume the localized energy identity is valid with the time-dependent cutoff
\(\Phi\), and assume \(B_{\mathrm{mov}}\in L^1(\tau_0,\infty)\).  Then
\[
  \mathcal E_R^{\mathrm{corr}}(\tau)
  =
  \mathcal E_R(\tau)+\int_\tau^\infty B_{\mathrm{mov}}(s)\,ds
\]
satisfies
\[
  \frac{d}{d\tau}\mathcal E_R^{\mathrm{corr}}(\tau)
  +\frac12\tilde{\mathfrak D}_{R(\tau)}(\tau)\le0
\]
in distributions.
\end{theorem}

\begin{proof}
Use the localized energy identity with the time-dependent test function
\(\Phi(\tau,y)\).  Compared with the fixed-cutoff identity, the only new term
is
\[
  \frac12\int |V|^2\partial_\tau\Phi,
\]
which comes from differentiating the cutoff inside \(\mathcal E_R\).  All
other terms are the fixed-cutoff terms with \(\phi\) replaced by \(\Phi\).  Add
\(\frac12\tilde{\mathfrak D}_{R(\tau)}\) to both sides.  The gradient term
leaves
\[
  -\frac{\nu}{2}\int|\nabla V|^2\Phi\le0,
\]
and the core scale term contributes
\[
  \frac12a_-\int |V|^2\Phi .
\]
Since \(\phi\) is radial nonincreasing, one has
\[
  y\cdot\nabla\Phi(\tau,y)\le0
  \qquad\text{for all }(\tau,y).
\]
Therefore
\[
  -\frac{a}{2}\int |V|^2y\cdot\nabla\Phi
  \le
  \frac12a_+(\tau)\int |V|^2\bigl(-y\cdot\nabla\Phi\bigr),
\]
because the \(a_-\)-part is nonpositive.  Each remaining term is bounded above
by the corresponding absolute-value term in \(B_{\mathrm{mov}}\).  Hence
\[
  \frac{d}{d\tau}\mathcal E_R(\tau)
  +\frac12\tilde{\mathfrak D}_{R(\tau)}(\tau)
  \le B_{\mathrm{mov}}(\tau).
\]
The moving finite-error condition makes the correction tail finite and
absolutely continuous on compact subintervals.  Differentiating that tail
subtracts \(B_{\mathrm{mov}}\), giving the stated inequality.
\end{proof}

\begin{definition}[Summable moving-annulus errors]\label{paper6a:def:summable-moving-errors}
The moving-annulus summability condition is the existence of a nondecreasing
locally absolutely continuous \(R(\tau)\to\infty\) such that each of the seven
quantities
\[
\begin{gathered}
  \left|\int |V|^2\partial_\tau\Phi\right|,\quad
  \left|\int |V|^2\Delta\Phi\right|,\quad
  \left|\int |V|^2V\cdot\nabla\Phi\right|,\quad
  \left|\int P\,V\cdot\nabla\Phi\right|,\\
  |b|\left|\int |V|^2\nabla\Phi\right|,\quad
  a_-\int |V|^2\Phi,\quad
  a_+\int |V|^2\bigl(-y\cdot\nabla\Phi\bigr)
\end{gathered}
\]
is integrable on \([\tau_0,\infty)\).
\end{definition}

\begin{lemma}[Summable annuli imply moving finite error]
\label{paper6a:lem:summable-moving-errors}
If \cref{paper6a:def:summable-moving-errors} holds, then
\(B_{\mathrm{mov}}\in L^1(\tau_0,\infty)\).
\end{lemma}

\begin{proof}
The density \(B_{\mathrm{mov}}\) is a finite positive linear combination of
the seven quantities listed in
\cref{paper6a:def:summable-moving-errors}.  If each is integrable, then their
linear combination is integrable.
\end{proof}

\begin{definition}[Summable tightness schedule]\label{paper6a:def:summable-tightness}
A summable tightness schedule is a nondecreasing sequence \(R_n\to\infty\) such
that, on unit windows \(I_n=[\tau_0+n,\tau_0+n+1]\),
\[
  \sum_{n=0}^{\infty}
  \sup_{\tau\in I_n}
  \|V(\tau)\|_{L^3(A_{R_n})}^3<\infty,
\]
and the same schedule is compatible with the pressure-tail normalization used
for \(P\).  This is stronger than uniform critical tightness, which permits
large radii on each window but does not by itself give a monotone radius
schedule with a summable series.
\end{definition}

\begin{definition}[Moving-cost compatibility]\label{paper6a:def:moving-cost-compatibility}
The moving localized cost is compatible with the cost-divergence criterion if
\[
  \tilde{\mathfrak D}_{R(\tau)}(\tau)
  =
  \nu\int|\nabla V|^2\phi_{R(\tau)}
  +a_+(\tau)\int|V|^2\phi_{R(\tau)}
\]
is the localized cost appearing in
\cref{paper6a:ass:cost-divergence-exclusion}, or is compared to that cost by a
positive constant on the tail so that cost divergence is unchanged.
\end{definition}

\begin{theorem}[Moving-cutoff replacement]\label{paper6a:thm:moving-replacement}
Assume the repaired-gauge representation, localized energy identity,
moving-cost compatibility, and summable moving-annulus errors.  Then the moving
cutoff supplies the corrected monotonicity input needed in
\cref{paper6a:prop:cost-divergence-exclusion}.
\end{theorem}

\begin{proof}
\Cref{paper6a:lem:summable-moving-errors} gives
\(B_{\mathrm{mov}}\in L^1(\tau_0,\infty)\).  Applying
\cref{paper6a:thm:moving-corrected-monotonicity} gives corrected monotonicity
with cost \(\tilde{\mathfrak D}_{R(\tau)}\).  The moving-cost compatibility
assumption identifies this cost, or a tail-comparable version of it, with the
localized cost in the conditional exclusion criterion.
\end{proof}

\subsection{Negative Scale Drift}

Set
\[
  M_R(\tau)=\int |V(y,\tau)|^2\Phi(\tau,y)\,dy .
\]
The negative-scale contribution in the moving cutoff argument is
\[
  \int_{\tau_0}^{\infty}a_-(\tau)M_R(\tau)\,d\tau.
\]

\begin{definition}[Negative scale-drift alternatives]\label{paper6a:def:negative-drift}
We say that the branch has finite negative drift if
\[
  \int_{\tau_0}^{\infty}a_-(\tau)\,d\tau<\infty.
\]
It has bounded moving \(L^2\) core mass if
\[
  \sup_{\tau\ge\tau_0}M_R(\tau)<\infty.
\]
It has a moving \(L^2\) core floor if there are \(m_0>0\) and
\(\tau_1\ge\tau_0\) such that
\[
  M_R(\tau)\ge m_0
  \qquad\text{for a.e. }\tau\ge\tau_1.
\]
The scale-collapse drift obstruction is the divergent alternative
\[
  \int_{\tau_0}^{\infty}a_-(\tau)M_R(\tau)\,d\tau=\infty .
\]
\end{definition}

\begin{lemma}[Finite negative drift controls the scale term]
\label{paper6a:lem:finite-negative-drift}
If the branch has bounded moving \(L^2\) core mass and finite negative drift,
then
\[
  \int_{\tau_0}^{\infty}a_-(\tau)M_R(\tau)\,d\tau<\infty .
\]
\end{lemma}

\begin{proof}
Let \(M_*=\sup_{\tau\ge\tau_0}M_R(\tau)<\infty\).  Since \(a_-\ge0\),
\[
  \int_{\tau_0}^{\infty}a_-(\tau)M_R(\tau)\,d\tau
  \le
  M_*\int_{\tau_0}^{\infty}a_-(\tau)\,d\tau<\infty.
\]
\end{proof}

\begin{theorem}[Finite-or-infinite scale-negative alternative]
\label{paper6a:thm:finite-infinite-scale-negative}
For every renormalized orbit for which \(a_-\) and \(M_R\) are measurable and
nonnegative, precisely one of the following alternatives holds:
\[
  \int_{\tau_0}^{\infty}a_-(\tau)M_R(\tau)\,d\tau<\infty,
  \qquad
  \int_{\tau_0}^{\infty}a_-(\tau)M_R(\tau)\,d\tau=\infty.
\]
\end{theorem}

\begin{proof}
The integral
\[
  I_R=\int_{\tau_0}^{\infty}a_-(\tau)M_R(\tau)\,d\tau
\]
is an extended nonnegative number.  Hence either \(I_R<\infty\) or
\(I_R=\infty\), and the two cases are mutually exclusive.
\end{proof}

\begin{lemma}[Logarithmic scale identity]\label{paper6a:lem:log-scale}
For the repaired-gauge convention
\[
  \tau_t=\lambda(t)^{-2},\qquad a(\tau)=\lambda(t)\lambda_t(t),
\]
one has
\[
  \int_{\tau_1}^{\tau_2}a(\tau)\,d\tau
  =
  \log\lambda(\tau_2)-\log\lambda(\tau_1)
\]
whenever \(\tau_0\le\tau_1<\tau_2<\infty\).
\end{lemma}

\begin{proof}
The chain rule gives
\[
  \frac{d}{d\tau}\log\lambda
  =
  \frac{\lambda_t}{\lambda}\frac{dt}{d\tau}
  =
  \frac{\lambda_t}{\lambda}\lambda^2
  =
  \lambda\lambda_t=a.
\]
Integrating over \([\tau_1,\tau_2]\) proves the identity.
\end{proof}

\begin{theorem}[Collapse with core floor forces scale-drift obstruction]
\label{paper6a:thm:collapse-core-floor}
Assume
\[
  \lambda(\tau)\to0\qquad\text{as }\tau\to\infty,
\]
and assume the moving \(L^2\) core floor of
\cref{paper6a:def:negative-drift}.  Then
\[
  \int_{\tau_0}^{\infty}a_-(\tau)M_R(\tau)\,d\tau=\infty .
\]
\end{theorem}

\begin{proof}
By \cref{paper6a:lem:log-scale},
\[
  \int_{\tau_1}^{\tau_2}a(\tau)\,d\tau
  =
  \log\lambda(\tau_2)-\log\lambda(\tau_1).
\]
Since \(\lambda(\tau)\to0\), the right-hand side tends to \(-\infty\) as
\(\tau_2\to\infty\).  If
\(\int_{\tau_1}^{\infty}a_-(\tau)\,d\tau<\infty\), then
\[
  \int_{\tau_1}^{\tau_2}a(\tau)\,d\tau
  =
  \int_{\tau_1}^{\tau_2}a_+(\tau)\,d\tau
  -
  \int_{\tau_1}^{\tau_2}a_-(\tau)\,d\tau
  \ge
  -\int_{\tau_1}^{\infty}a_-(\tau)\,d\tau>-\infty,
\]
which contradicts the preceding limit.  Thus
\(\int_{\tau_1}^{\infty}a_-(\tau)\,d\tau=\infty\).

By the moving \(L^2\) core floor, \(M_R(\tau)\ge m_0\) for a.e.
\(\tau\ge\tau_1\).  Therefore
\[
  \int_{\tau_0}^{\infty}a_-(\tau)M_R(\tau)\,d\tau
  \ge
  m_0\int_{\tau_1}^{\infty}a_-(\tau)\,d\tau
  =
  \infty.
\]
\end{proof}

\begin{corollary}[Scale-drift alternatives]\label{paper6a:cor:scale-drift-alternatives}
In the moving-cutoff setup, the negative-scale term has the following
exhaustive alternatives.  If bounded moving \(L^2\) core mass and finite negative
drift hold, then the negative-scale term is integrable.  If the weighted
negative-drift integral is infinite, the branch lies in the scale-collapse
drift obstruction.  If genuine collapse and a moving \(L^2\) core floor also
hold, then the obstruction is forced.
\end{corollary}

\begin{proof}
The first case is \cref{paper6a:lem:finite-negative-drift}.  The exhaustive
split is \cref{paper6a:thm:finite-infinite-scale-negative}.  The final
assertion is \cref{paper6a:thm:collapse-core-floor}.
\end{proof}

\subsection{Scale-Drift Cost Criteria}

\begin{definition}[Scale-collapse drift cost]\label{paper6a:def:scale-collapse-cost}
For a renormalized branch define
\[
  \mathfrak C_{\mathrm{sc}}(\tau)
  =
  a_-(\tau)M_R(\tau),
  \qquad
  M_R(\tau)=\int |V(y,\tau)|^2\phi_{R(\tau)}(y)\,dy,
\]
and
\[
  \mathcal C_{\mathrm{sc}}[\tau_0,\infty)
  =
  \int_{\tau_0}^{\infty}\mathfrak C_{\mathrm{sc}}(\tau)\,d\tau .
\]
Thus the scale-collapse drift obstruction is precisely
\(\mathcal C_{\mathrm{sc}}[\tau_0,\infty)=\infty\).  The density is
nonnegative and local in the renormalized core.
\end{definition}

\begin{definition}[Scale-collapse cost compatibility]\label{paper6a:def:scale-collapse-cost-compatibility}
The scale-collapse drift cost is compatible with the cost-divergence criterion
when:
\begin{enumerate}[label=\textup{(\roman*)}]
\item \(\mathfrak C_{\mathrm{sc}}\) is measurable, nonnegative, and locally
integrable on finite \(\tau\)-intervals for renormalized branches;
\item divergence of \(\mathcal C_{\mathrm{sc}}[\tau_0,\infty)\) is included
among the hypotheses in \cref{paper6a:ass:cost-divergence-exclusion};
\item the resulting conclusion is the same local Type II exclusion conclusion
used for the fixed localized cost.
\end{enumerate}
\end{definition}

\begin{theorem}[Scale-collapse drift cost criterion]
\label{paper6a:thm:scale-collapse-drift-cost}
If the scale-collapse drift obstruction holds and the cost
\(\mathfrak C_{\mathrm{sc}}\) is compatible with the cost-divergence criterion, then the divergent-cost hypothesis
of \cref{paper6a:prop:cost-divergence-exclusion} is satisfied for the
scale-collapse branch.
\end{theorem}

\begin{proof}
By \cref{paper6a:def:scale-collapse-cost}, the obstruction is precisely
\[
  \mathcal C_{\mathrm{sc}}[\tau_0,\infty)=\infty .
\]
By \cref{paper6a:def:scale-collapse-cost-compatibility}, this divergence is a
nonnegative cost-divergence hypothesis with the same local conclusion.
\end{proof}

\begin{corollary}[Scale-collapse exclusion by drift cost]
\label{paper6a:cor:scale-collapse-drift-exclusion}
Assume a renormalized local Type II branch has scale-collapse drift
obstruction, compatible scale-collapse drift cost, and the remaining
conditional cost hypotheses.  If
\cref{paper6a:ass:cost-divergence-exclusion} holds, then the branch is excluded
from the local Type II class under consideration.
\end{corollary}

\begin{proof}
\Cref{paper6a:thm:scale-collapse-drift-cost} supplies the divergent
localized cost.  Applying
\cref{paper6a:prop:cost-divergence-exclusion} gives the exclusion.
\end{proof}

\begin{definition}[Absolute scale-drift cost]\label{paper6a:def:absolute-scale-cost}
Define
\[
  \mathfrak C_{\mathrm{abs}}(\tau)
  =
  \nu\int|\nabla V|^2\phi_{R(\tau)}
  +|a(\tau)|M_R(\tau).
\]
The absolute scale-drift cost is compatible with the cost-divergence criterion
if divergence of
\(\int_{\tau_0}^{\infty}\mathfrak C_{\mathrm{abs}}(\tau)\,d\tau\) is included
as a Type II cost-divergence condition with the same local conclusion.
\end{definition}

\begin{lemma}[Scale-collapse drift forces absolute-cost divergence]
\label{paper6a:lem:scale-drift-absolute}
If the scale-collapse drift obstruction holds, then
\[
  \int_{\tau_0}^{\infty}\mathfrak C_{\mathrm{abs}}(\tau)\,d\tau=\infty .
\]
\end{lemma}

\begin{proof}
Since
\[
  \mathfrak C_{\mathrm{abs}}(\tau)
  =
  \nu\int|\nabla V|^2\phi_{R(\tau)}
  +|a(\tau)|M_R(\tau)
  \ge a_-(\tau)M_R(\tau),
\]
divergence of \(\int a_-M_R\) forces divergence of
\(\int\mathfrak C_{\mathrm{abs}}\).
\end{proof}

\begin{corollary}[Absolute scale-cost criterion]\label{paper6a:cor:absolute-scale-cost-criterion}
If the scale-collapse drift obstruction holds and the absolute scale-drift
cost is compatible with the cost-divergence criterion, then the divergent-cost hypothesis of
\cref{paper6a:prop:cost-divergence-exclusion} is satisfied.
\end{corollary}

\begin{proof}
\Cref{paper6a:lem:scale-drift-absolute} gives divergence of the absolute
scale-drift cost, and compatibility makes this divergence a Type II
cost-divergence condition.
\end{proof}

\begin{theorem}[Moving cutoff with scale-collapse alternative]
\label{paper6a:thm:moving-cutoff-scale-collapse-alternative}
Assume the moving-cutoff setup and a renormalized local Type II branch.
\begin{enumerate}[label=\textup{(\roman*)}]
\item If the negative-scale term is integrable and the remaining moving-annulus
errors are summable, then the moving-cutoff corrected monotonicity estimate is
available.
\item If the scale-collapse drift obstruction holds and either the
scale-collapse drift cost or the absolute scale-drift cost is compatible with
the cost-divergence criterion, then the divergent-cost hypothesis is satisfied
directly.
\item If, in addition, the remaining hypotheses of
\cref{paper6a:def:cost-hypotheses} and
\cref{paper6a:ass:cost-divergence-exclusion} hold, then the branch is excluded.
\end{enumerate}
\end{theorem}

\begin{proof}
The first branch is \cref{paper6a:thm:moving-replacement}.  The second branch
is \cref{paper6a:thm:scale-collapse-drift-cost} or
\cref{paper6a:cor:absolute-scale-cost-criterion}, depending on which compatible
cost is used.  The third branch is
\cref{paper6a:prop:cost-divergence-exclusion}.
\end{proof}

\subsection{Discharge of the Cost-Divergence Exclusion Assumption}
\label{paper6a:subsec:discharge-cost-divergence-assumption}

The conditional exclusion criterion of
\cref{paper6a:ass:cost-divergence-exclusion} is, given the corrected
monotonicity input that is already part of
\cref{paper6a:def:cost-hypotheses}, an arithmetic consequence of the
nonnegativity of the localized renormalized energy together with the
windowwise routing discharge of
\cref{paper7:cor:windowwise-positive-scale-routing-holds,paper7:cor:windowwise-transition-routing-holds}.
We make this discharge explicit here so that the assumption can be removed
from the assembly hypotheses.

\begin{lemma}[Nonnegativity of the corrected localized energy]
\label{paper6a:lem:corrected-energy-nonnegative}
Let \(\phi=\phi_{R_0}\in C_c^\infty(\mathbb R^3)\) be a radial nonnegative
cutoff, and let \(R(\tau)\) be either fixed at \(R_0\) or the canonical
windowwise interpolated radius of
\cref{paper7:lem:windowwise-shell-pressure-translation}.  Then
\(\mathcal E_\phi(\tau)=\tfrac12\int|V(y,\tau)|^2\phi(y)\,dy\ge 0\) and
\(\mathcal E_R(\tau)=\tfrac12\int|V(y,\tau)|^2\Phi(\tau,y)\,dy\ge 0\) for
every \(\tau\), and the correction tails
\(\int_\tau^\infty B_{R_0}(s)\,ds\) and
\(\int_\tau^\infty B_{\mathrm{mov}}(s)\,ds\) of
\cref{paper6a:def:fixed-error-density,paper6a:def:moving-error-density} are
nonnegative.  Consequently the corrected energies
\(\mathcal E_\phi^{\mathrm{corr}}\) and
\(\mathcal E_R^{\mathrm{corr}}\) of
\cref{paper6a:thm:fixed-corrected-monotonicity,paper6a:thm:moving-corrected-monotonicity}
are nonnegative.
\end{lemma}

\begin{proof}
Both \(\mathcal E_\phi\) and \(\mathcal E_R\) are integrals of \(|V|^2\)
against a nonnegative weight, hence nonnegative.  Inspecting the seven
summands of \(B_{R_0}\) (\cref{paper6a:def:fixed-error-density}) and the
seven summands of \(B_{\mathrm{mov}}\)
(\cref{paper6a:def:moving-error-density}): five are absolute values of
integrals and are therefore nonnegative; the remaining two are
\(\tfrac12 a_-(\tau)\!\int|V|^2\Phi\) and
\(\tfrac12 a_+(\tau)\!\int|V|^2(-y\cdot\nabla\Phi)\) (and likewise with
\(\Phi\) replaced by \(\phi\)).  In the first, \(a_-\ge 0\) and the
integrand \(|V|^2\Phi\ge 0\), so the term is nonnegative.  In the second,
\(a_+\ge 0\), and \(-y\cdot\nabla\Phi\ge 0\) because
\(\Phi(\tau,y)=\phi(y/R(\tau))\) with \(\phi\) radial nonincreasing forces
\(\nabla\Phi(\tau,y)\) to point opposite to \(y\); the term is nonnegative.
Hence \(B_{R_0},B_{\mathrm{mov}}\ge 0\) pointwise, and their tail integrals
are nonnegative.  Adding a nonnegative tail to a nonnegative quantity
preserves nonnegativity, so
\(\mathcal E_\phi^{\mathrm{corr}}\) and \(\mathcal E_R^{\mathrm{corr}}\) are
nonnegative.
\end{proof}

\begin{lemma}[Corrected monotonicity is incompatible with cost divergence]
\label{paper6a:lem:corrected-monotonicity-versus-divergent-cost}
Let \(\mathcal E^{\mathrm{corr}}\colon[\tau_0,\infty)\to[0,\infty)\) be
locally integrable with finite initial value
\(\mathcal E^{\mathrm{corr}}(\tau_0)<\infty\), let
\(\mathcal D\colon[\tau_0,\infty)\to[0,\infty)\) be locally integrable, and
let \(c_0>0\).  If
\[
  \frac{d}{d\tau}\mathcal E^{\mathrm{corr}}(\tau)+c_0\mathcal D(\tau)\le 0
  \qquad\text{in distributions on }(\tau_0,\infty),
\]
then
\[
  \int_{\tau_0}^{\infty}\mathcal D(\tau)\,d\tau
  \le \frac{1}{c_0}\,\mathcal E^{\mathrm{corr}}(\tau_0)<\infty .
\]
In particular, the simultaneous validity of corrected monotonicity in this
sense and of \(\int_{\tau_0}^{\infty}\mathcal D=\infty\) is contradictory.
The same conclusion holds if the inequality is strengthened to the form
\((d/d\tau)\mathcal E^{\mathrm{corr}}+c_0\mathcal D+\mathcal R\le 0\) with
\(\mathcal R\ge 0\), as in
\cref{paper6a:def:corrected-monotonicity-input}.
\end{lemma}

\begin{proof}
The distributional inequality
\((\mathcal E^{\mathrm{corr}})'+c_0\mathcal D\le 0\) on
\((\tau_0,\infty)\), with \(\mathcal D\ge 0\) and
\(\mathcal E^{\mathrm{corr}}\ge 0\) locally integrable, implies that
\(\mathcal E^{\mathrm{corr}}\) admits a non-increasing right-continuous
representative.  Testing the inequality against a smooth approximation of
\(\mathbf 1_{[\tau_0,\tau]}\) and passing to the limit using
\(\mathcal D\ge 0\) and \(\mathcal E^{\mathrm{corr}}\ge 0\) (so the
boundary terms have a definite sign) yields
\[
  c_0\int_{\tau_0}^{\tau}\mathcal D(\sigma)\,d\sigma
  \le \mathcal E^{\mathrm{corr}}(\tau_0)-\mathcal E^{\mathrm{corr}}(\tau)
  \le \mathcal E^{\mathrm{corr}}(\tau_0).
\]
Letting \(\tau\to\infty\) and using monotone convergence
(\(\mathcal D\ge 0\)) yields the displayed bound.  If
\(\int\mathcal D=\infty\) the left-hand side is infinite while the
right-hand side is finite, a contradiction.
\end{proof}

\begin{theorem}[Cost-divergence/corrected-monotonicity dichotomy on the canonical schedule]
\label{paper6a:thm:cost-divergence-routing-dichotomy}
Let \(\omega\) be a renormalized local Type II branch in the pre-carrier
validated compact-cost channel of
\cref{paper7:def:precarrier-validated-compact-channel}.  Equip \(\omega\)
with the canonical windowwise interpolated cutoff
\(\Phi(\tau,y)=\phi(y/R(\tau))\) of
\cref{paper7:lem:windowwise-shell-pressure-translation}, and let
\(\widetilde{\mathfrak D}_{R(\tau)}\) be the moving localized cost of
\cref{paper6a:def:moving-cost-compatibility}.  Assume in addition the
localized energy identity for the moving cutoff \(\Phi(\tau,\cdot)\) of
\cref{paper6a:def:cost-hypotheses}\textup{(v)} (supplied, e.g., by
\cref{paper7:lem:suitable-energy-gives-carrier-subidentities} on the
channel), and the renormalized Lebesgue control
\(\sup_{\tau\ge\tau_0}\|V(\tau)\|_{L^3(\mathbb R^3)}\le M<\infty\)
that is part of the pre-carrier validated compact-cost channel data
(established in the proof of
\cref{paper7:lem:windowwise-shell-pressure-translation}).  If
\[
  \int_{\tau_0}^{\infty}\widetilde{\mathfrak D}_{R(\tau)}(\tau)\,d\tau=\infty,
\]
then the following sharpened dichotomy holds: at least one of the three
components
\[
  J_1,\qquad J_6,\qquad J_7
\]
fails to be in \(L^1(\tau_1,\infty)\) for every \(\tau_1\ge\tau_0\).  More
precisely, exactly one of the following two mutually exclusive alternatives
holds:
\begin{enumerate}[label=\textup{(D\arabic*)}]
\item at least one of \(J_1\), \(J_6\), \(J_7\) fails to be in
\(L^1(\tau_1,\infty)\) for every \(\tau_1\ge\tau_0\);
\item every \(J_k\), \(k\in\{1,\dots,7\}\), belongs to
\(L^1(\tau_0,\infty)\).
\end{enumerate}
Moreover, case \textup{(D2)} is incompatible with cost divergence, so
divergence of the cost forces case \textup{(D1)}, i.e.~failure of one of
\(J_1\), \(J_6\), \(J_7\).
\end{theorem}

\begin{proof}
The moving error density of \cref{paper6a:def:moving-error-density} satisfies
the identity
\[
  B_{\mathrm{mov}}(\tau)
  =\tfrac12 J_1(\tau)+\tfrac{\nu}{2}J_2(\tau)+\tfrac12 J_3(\tau)
  +J_4(\tau)+\tfrac12 J_5(\tau)+\tfrac12 J_6(\tau)+\tfrac12 J_7(\tau),
\]
with \(J_k\) as in \cref{paper7:def:componentwise-annular-finite-error},
because the seven nonnegative terms in \(B_{\mathrm{mov}}\) coincide with
\(J_1,\dots,J_7\) up to the listed multiplicative constants, all of which
are strictly positive (recall \(\nu>0\) is fixed throughout).  In
particular \(B_{\mathrm{mov}}\in L^1(\tau_0,\infty)\) if and only if every
\(J_k\in L^1(\tau_0,\infty)\).

By \cref{paper7:lem:windowwise-shell-pressure-translation}, on the canonical
windowwise schedule one has automatically
\[
  J_2,\,J_3,\,J_4,\,J_5\in L^1(\tau_0,\infty).
\]
Therefore the only \(J_k\) that can fail to be in \(L^1\) on a tail are
\(J_1\), \(J_6\), \(J_7\).  We obtain the following dichotomy: either
at least one of \(J_1,J_6,J_7\) fails to be in \(L^1(\tau_1,\infty)\) for
every \(\tau_1\ge\tau_0\) (case \textup{(D1)}), or each of \(J_1,J_6,J_7\) is
integrable on \([\tau_0,\infty)\) and hence (combined with the automatic
integrability of \(J_2,J_3,J_4,J_5\)) every \(J_k\) is integrable on
\([\tau_0,\infty)\) (case \textup{(D2)}).

In case \textup{(D2)}, \cref{paper6a:lem:summable-moving-errors} gives
\(B_{\mathrm{mov}}\in L^1(\tau_0,\infty)\), and the localized energy identity
of \cref{paper6a:def:cost-hypotheses}\textup{(v)} together with
\cref{paper6a:thm:moving-corrected-monotonicity} produces
\[
  \frac{d}{d\tau}\mathcal E_R^{\mathrm{corr}}(\tau)
  +\tfrac12\,\widetilde{\mathfrak D}_{R(\tau)}(\tau)\le 0
  \qquad\text{in distributions on }(\tau_0,\infty),
\]
with \(c_0=\tfrac12\) in the sense of
\cref{paper6a:def:corrected-monotonicity-input}.  We verify the finite
initial value.  Since \(\Phi(\tau_0,\cdot)\) is supported in the ball
\(B_{2R(\tau_0)}\subset\mathbb R^3\) with \(0\le\Phi\le 1\), H\"older's
inequality gives
\[
  \mathcal E_R(\tau_0)
  =\tfrac12\!\int|V(\tau_0)|^2\Phi(\tau_0,\cdot)\,dy
  \le \tfrac12\!\int_{B_{2R(\tau_0)}}\!\!|V(\tau_0)|^2\,dy
  \le \tfrac12\,|B_{2R(\tau_0)}|^{1/3}\,\|V(\tau_0)\|_{L^3(\mathbb R^3)}^2
  \le C\,R(\tau_0)\,M^2,
\]
where \(C\) is an absolute constant.  Combined with the case-\textup{(D2)}
integrability \(\|B_{\mathrm{mov}}\|_{L^1(\tau_0,\infty)}<\infty\),
\[
  \mathcal E_R^{\mathrm{corr}}(\tau_0)
  =\mathcal E_R(\tau_0)+\int_{\tau_0}^{\infty}\!\!B_{\mathrm{mov}}(s)\,ds
  \le C\,R(\tau_0)\,M^2+\|B_{\mathrm{mov}}\|_{L^1(\tau_0,\infty)}<\infty.
\]
By \cref{paper6a:lem:corrected-energy-nonnegative}, \(\mathcal
E_R^{\mathrm{corr}}\ge 0\), so
\cref{paper6a:lem:corrected-monotonicity-versus-divergent-cost} yields
\[
  \int_{\tau_0}^{\infty}\widetilde{\mathfrak D}_{R(\tau)}(\tau)\,d\tau
  \le 2\,\mathcal E_R^{\mathrm{corr}}(\tau_0)<\infty,
\]
contradicting the cost divergence.  Hence case \textup{(D2)} cannot occur,
and case \textup{(D1)} must hold; that is, at least one of \(J_1,J_6,J_7\)
fails to be in \(L^1\) on every later tail.
\end{proof}

\begin{theorem}[Discharge of the cost-divergence exclusion assumption on the canonical windowwise schedule]
\label{paper6a:thm:cost-divergence-exclusion-discharged}
Fix the canonical windowwise moving family of
\cref{paper7:lem:windowwise-shell-pressure-translation} on the pre-carrier
validated compact-cost channel of
\cref{paper7:def:precarrier-validated-compact-channel}, and assume the
localized energy identity for the moving cutoff
(\cref{paper7:lem:suitable-energy-gives-carrier-subidentities}) and the
renormalized Lebesgue control
\(\sup_{\tau\ge\tau_0}\|V(\tau)\|_{L^3(\mathbb R^3)}\le M<\infty\) carried
by the channel.  Assume further that
\textup{(R1)} failure of \(J_6\in L^1(\tau_1,\infty)\) on every later tail
\([\tau_1,\infty)\) routes the branch into the negative scale-drift
alternative \textup{(vi)} of
\cref{paper7:def:stratified-moving-carrier-routing}, in the sense of
\cref{paper7:thm:windowwise-routing-gives-stratified-routing} (this is the
same negative-drift routing hypothesis already standing in that theorem,
realized through \cref{paper6a:def:negative-drift,paper6a:cor:scale-drift-alternatives});
and \textup{(R2)} the named adjacent strata of
\cref{paper7:thm:noncarrier-exits-route} together with the
exterior/translation/positive-scale/negative-drift alternatives
\textup{(iii)}--\textup{(vi)} of
\cref{paper7:def:stratified-moving-carrier-routing} are closed by their
assigned local arguments in the ambient Type II class.
Then for every renormalized local Type II branch \(\omega\) in that channel
satisfying \cref{paper6a:def:cost-hypotheses}, divergence of the localized
cost
\[
  \int_{\tau_0}^{\infty}\widetilde{\mathfrak D}_{R(\tau)}(\tau)\,d\tau=\infty
\]
excludes \(\omega\) from the local Type II class.  Equivalently, the
negative-drift routing \textup{(R1)} together with the adjacent-stratum
closures \textup{(R2)} imply
\cref{paper6a:ass:cost-divergence-exclusion} as a theorem on the canonical
windowwise schedule.
\end{theorem}

\begin{proof}
Let \(\omega\) be a candidate branch.  By
\cref{paper6a:thm:cost-divergence-routing-dichotomy}, case \textup{(D2)} is
impossible, so case \textup{(D1)} occurs: at least one of
\(J_1\), \(J_6\), \(J_7\) fails to be in \(L^1(\tau_1,\infty)\) for every
\(\tau_1\ge\tau_0\).

Apply the routing statement assigned to that index:
\begin{itemize}
\item if \(J_1\) fails, the windowwise transition/leakage routing of
\cref{paper7:cor:windowwise-transition-routing-holds} (proved
unconditionally on the channel via
\cref{paper7:thm:windowwise-transition-routing-certificate}) routes
\(\omega\) into the exterior/leakage alternative \textup{(iii)} of
\cref{paper7:def:stratified-moving-carrier-routing}, entering the named
critical-tightness failure stratum of
\cref{paper7:thm:noncarrier-exits-route};
\item if \(J_6\) fails, hypothesis \textup{(R1)} routes \(\omega\) into the
negative scale-drift alternative \textup{(vi)};
\item if \(J_7\) fails, the windowwise positive-scale routing of
\cref{paper7:cor:windowwise-positive-scale-routing-holds} (proved
unconditionally via the bridge
\cref{paper7:thm:positive-scale-certificate-enters-scale-rigid} and
\cref{paper3:prop:scale-rigidity-discharged}) routes \(\omega\) into the
positive-scale alternative \textup{(v)}.
\end{itemize}
By hypothesis \textup{(R2)}, each of these named alternatives is closed by
an already-discharged argument in the ambient Type II class.  Therefore
\(\omega\) is excluded.  This proves the conclusion of
\cref{paper6a:ass:cost-divergence-exclusion} as a theorem under the standing
hypotheses \textup{(R1)} and \textup{(R2)}.
\end{proof}

\begin{corollary}[Cost-divergence exclusion holds under the same hypothesis package as adjacent closure]
\label{paper6a:cor:cost-divergence-exclusion-unconditional}
The hypotheses \textup{(R1)}--\textup{(R2)} of
\cref{paper6a:thm:cost-divergence-exclusion-discharged} on the canonical
windowwise schedule coincide, after combining with the unconditional
windowwise discharges
\cref{paper7:cor:windowwise-positive-scale-routing-holds,paper7:cor:windowwise-transition-routing-holds},
with the standing hypothesis package of
\cref{paper7:thm:windowwise-routing-gives-stratified-routing} (negative-drift
routing for \(J_6\)) plus the adjacent-stratum closure of
\cref{paper7:cor:adjacent-closure-reduced-to-windowwise-routing}.  In
particular, on any ambient Type II class on which the conclusion of
\cref{paper7:cor:adjacent-closure-reduced-to-windowwise-routing} (compact-cost
adjacent-stratum closure) is established,
\cref{paper6a:ass:cost-divergence-exclusion} is a theorem and may be removed
from the hypothesis lists of
\cref{paper6a:def:abstract-final-data,paper6a:def:expanded-terminal-data}.
\end{corollary}

\begin{proof}
The windowwise transition/leakage routing
\cref{paper7:cor:windowwise-transition-routing-holds} and windowwise
positive-scale routing
\cref{paper7:cor:windowwise-positive-scale-routing-holds} are unconditional
on the channel.  Combined with the negative-drift routing for \(J_6\)
imported from \cref{paper7:thm:windowwise-routing-gives-stratified-routing},
the routing component \textup{(R1)} of
\cref{paper6a:thm:cost-divergence-exclusion-discharged} is satisfied.
The adjacent-stratum closure component \textup{(R2)} is exactly the
standing hypothesis closed under
\cref{paper7:cor:adjacent-closure-reduced-to-windowwise-routing}.  Therefore
\cref{paper6a:thm:cost-divergence-exclusion-discharged} produces exclusion
without any further hypothesis, and
\cref{paper6a:ass:cost-divergence-exclusion} is a theorem in that ambient
class.
\end{proof}

\begin{remark}[Status of the cost-divergence assumption]
\label{paper6a:rem:cost-divergence-assumption-status}
The discharge has three independent ingredients:
\begin{enumerate}[label=\textup{(\roman*)}]
\item the trivial arithmetic consequence
\cref{paper6a:lem:corrected-monotonicity-versus-divergent-cost} that
corrected localized monotonicity together with finite initial corrected
energy and \(\mathcal D\ge 0\) forces \(\int\mathcal D<\infty\);
\item the moving error-density identity
\(B_{\mathrm{mov}}=\sum_k c_k J_k\) of
\cref{paper6a:def:moving-error-density,paper7:def:componentwise-annular-finite-error}
(positive linear combination), combined with the canonical schedule
integrability \(J_2,J_3,J_4,J_5\in L^1(\tau_0,\infty)\) of
\cref{paper7:lem:windowwise-shell-pressure-translation};
\item the windowwise routing for \(J_1\) and \(J_7\)
(\cref{paper7:cor:windowwise-transition-routing-holds,paper7:cor:windowwise-positive-scale-routing-holds},
proved unconditionally on the channel) plus the negative-drift routing for
\(J_6\) imported from
\cref{paper7:thm:windowwise-routing-gives-stratified-routing}.
\end{enumerate}
Thus \cref{paper6a:ass:cost-divergence-exclusion} no longer represents an
independent analytic input.  Its content is the conjunction of the
trivial monotonicity-vs-divergence arithmetic and the windowwise routing
package that is already used to close adjacent strata in
\cref{paper7:cor:adjacent-closure-reduced-to-windowwise-routing}.
\end{remark}

\subsection{Absolutely Continuous Repaired-Gauge Charts and Renormalized Equations}
\label{subsec:ac-repaired-gauge-renormalized-equations}

The following statements isolate the parts of the repaired-gauge
representation that follow from an absolutely continuous concentration chart
and an absolutely continuous repaired-gauge path.  They do not prove extraction
of the chart from arbitrary data and do not prove solvability of the gauge
constraints.

\begin{definition}[Absolutely continuous concentration chart]
\label{paper6a:def:ac-initial-concentration-chart}
Let \(u,p\) solve the Navier--Stokes equations on
\(\mathbb R^3\times(t_0,T^*)\),
\[
  \partial_tu+(u\cdot\nabla)u+\nabla p=\nu\Delta u,
  \qquad
  \nabla\cdot u=0.
\]
An absolutely continuous concentration chart consists of
absolutely continuous functions
\[
  x_c:(t_0,T^*)\to\mathbb R^3,\qquad
  \lambda:(t_0,T^*)\to(0,\infty),
\]
and a strictly increasing absolutely continuous time \(\rho\) such that
\[
  \rho_t=\lambda^{-2}\quad\text{a.e.},
\]
\[
  y=\frac{x-x_c(t)}{\lambda(t)},\qquad
  u(x,t)=\lambda(t)^{-1}U(y,\rho(t)),
\]
and
\[
  p(x,t)=\lambda(t)^{-2}Q(y,\rho(t))+c(t)
\]
for some scalar function \(c(t)\).  The chart is assumed to have the local
integrability required for the distributional chain rule on compact sets in
the \(y\)-variables.  Define
\[
  A(\rho(t)):=\lambda(t)\lambda_t(t),
  \qquad
  B(\rho(t)):=\lambda(t)x_c'(t),
\]
or equivalently
\[
  A=\partial_\rho\log\lambda,\qquad
  B=\lambda^{-1}\partial_\rho x_c ,
\]
where \(\lambda\) and \(x_c\) are regarded as functions of \(\rho\).
\end{definition}

\begin{lemma}[Initial renormalized equation]
\label{paper6a:lem:initial-renormalized-equation}
Assume Definition~\ref{paper6a:def:ac-initial-concentration-chart}.  Then \(U,Q\) satisfy
\[
  \partial_\rho U+(U\cdot\nabla)U+\nabla Q
  =
  \nu\Delta U+A(\rho)(U+y\cdot\nabla U)+B(\rho)\cdot\nabla U,
  \qquad
  \nabla\cdot U=0
\]
distributionally on compact \(y\)-sets.
\end{lemma}

\begin{proof}
For smooth functions the computation is pointwise.  The distributional case is
obtained by testing against compactly supported functions and using the
assumed chart chain rule.  The spatial derivatives are
\[
  \nabla_xu=\lambda^{-2}\nabla_yU,\qquad
  \Delta_xu=\lambda^{-3}\Delta_yU,
  \qquad
  (u\cdot\nabla_x)u=\lambda^{-3}(U\cdot\nabla_y)U.
\]
Since
\[
  y_t=-\frac{x_c'}{\lambda}-\frac{\lambda_t}{\lambda}y,
  \qquad
  \rho_t=\lambda^{-2},
\]
one has
\[
  \partial_tu
  =
  \lambda^{-3}\partial_\rho U
  -\lambda^{-2}\lambda_t(U+y\cdot\nabla U)
  -\lambda^{-2}x_c'\cdot\nabla U.
\]
Also
\[
  \nabla_xp=\lambda^{-3}\nabla_yQ.
\]
Substitution into the physical Navier--Stokes equation and multiplication by
\(\lambda^3\) gives
\[
  \partial_\rho U
  -\lambda\lambda_t(U+y\cdot\nabla U)
  -\lambda x_c'\cdot\nabla U
  +(U\cdot\nabla)U+\nabla Q
  =
  \nu\Delta U.
\]
Moving the two chart-velocity terms to the right gives the displayed equation.
Finally,
\[
  \nabla_x\cdot u=\lambda^{-2}\nabla_y\cdot U,
\]
so \(\nabla_y\cdot U=0\).
\end{proof}

\begin{definition}[Absolutely continuous repaired-gauge path]
\label{paper6a:def:ac-repaired-gauge-path}
An absolutely continuous repaired-gauge path for the initial chart consists, on
every compact final time interval, of absolutely continuous functions
\[
  \mu(\tau)>0,\qquad q(\tau)\in\mathbb R^3,\qquad \rho=\rho(\tau),
\]
such that
\[
  \rho_\tau=\mu(\tau)^2\quad\text{a.e.},
\]
\[
  V(Y,\tau)=\mu(\tau)U(\mu(\tau)Y+q(\tau),\rho(\tau)),
\]
and
\[
  P(Y,\tau)=\mu(\tau)^2Q(\mu(\tau)Y+q(\tau),\rho(\tau)).
\]
The scale and centering constraints defining the repaired gauge are assumed to
hold along this path.  This hypothesis assumes the path has already been
selected; it does not prove solvability of the gauge constraints.
\end{definition}

\begin{lemma}[Physical chart after gauge repair]
\label{paper6a:lem:repaired-physical-chart}
Assume Definitions~\ref{paper6a:def:ac-initial-concentration-chart}
and~\ref{paper6a:def:ac-repaired-gauge-path}.
Let \(t=t(\rho(\tau))\), and define
\[
  \Lambda(\tau):=\lambda(t(\rho(\tau)))\mu(\tau),
  \qquad
  X(\tau):=x_c(t(\rho(\tau)))+\lambda(t(\rho(\tau)))q(\tau).
\]
Then, on every compact final time interval, \(\Lambda\) and \(X\) are
absolutely continuous, \(\Lambda>0\), and
\[
  \frac{d\tau}{dt}=\Lambda(\tau)^{-2}.
\]
Moreover,
\[
  x=X(\tau)+\Lambda(\tau)Y
\]
is the final physical chart and
\[
  u(x,t)=\Lambda(\tau)^{-1}V(Y,\tau),
  \qquad
  p(x,t)=\Lambda(\tau)^{-2}P(Y,\tau)+c(t).
\]
\end{lemma}

\begin{proof}
The intermediate time \(\rho\) is strictly increasing because
\(\rho_t=\lambda^{-2}>0\) a.e.; hence it has an absolutely continuous inverse
on compact subintervals of its range.  The compositions
\(\lambda(t(\rho(\tau)))\) and \(x_c(t(\rho(\tau)))\) are therefore
absolutely continuous on compact final windows.  Products and sums with the
absolutely continuous functions \(\mu,q\) are absolutely continuous, so
\(\Lambda\) and \(X\) are absolutely continuous.

The identity \(\rho_\tau=\mu^2\) implies
\[
  \frac{d\tau}{d\rho}=\mu^{-2}.
\]
Since \(d\rho/dt=\lambda^{-2}\),
\[
  \frac{d\tau}{dt}
  =
  \frac{d\tau}{d\rho}\frac{d\rho}{dt}
  =
  \mu^{-2}\lambda^{-2}
  =
  (\lambda\mu)^{-2}
  =
  \Lambda^{-2}.
\]
If \(x=X+\Lambda Y\), then
\[
  x=x_c+\lambda q+\lambda\mu Y,
\]
and hence the intermediate coordinate is
\[
  y=\frac{x-x_c}{\lambda}=q+\mu Y.
\]
Therefore
\[
  u(x,t)
  =
  \lambda^{-1}U(q+\mu Y,\rho)
  =
  \lambda^{-1}\mu^{-1}V(Y,\tau)
  =
  \Lambda^{-1}V(Y,\tau),
\]
and similarly
\[
  p(x,t)
  =
  \lambda^{-2}Q(q+\mu Y,\rho)+c(t)
  =
  \lambda^{-2}\mu^{-2}P(Y,\tau)+c(t)
  =
  \Lambda^{-2}P(Y,\tau)+c(t).
\]
\end{proof}

\begin{lemma}[Modulation coefficients in repaired gauge]
\label{paper6a:lem:repaired-modulation-coefficients}
Define
\[
  a(\tau):=\partial_\tau\log\Lambda(\tau),
  \qquad
  b(\tau):=\Lambda(\tau)^{-1}X_\tau(\tau),
\]
where derivatives are understood a.e. on compact final windows.  Under the
hypotheses of Lemma~\ref{paper6a:lem:repaired-physical-chart},
\[
  a(\tau)=\mu(\tau)^2A(\rho(\tau))+\frac{\mu_\tau(\tau)}{\mu(\tau)}
\]
and
\[
  b(\tau)=\mu(\tau)B(\rho(\tau))
  +\mu(\tau)A(\rho(\tau))q(\tau)
  +\frac{q_\tau(\tau)}{\mu(\tau)}.
\]
In particular \(a,b\in L^1_{\rm loc}\) on every compact final time interval.
\end{lemma}

\begin{proof}
Since \(\Lambda=\lambda(\rho(\tau))\mu(\tau)\),
\[
  \partial_\tau\log\Lambda
  =
  \rho_\tau\partial_\rho\log\lambda+\partial_\tau\log\mu
  =
  \mu^2A+\frac{\mu_\tau}{\mu}.
\]
This proves the formula for \(a\).  Next
\[
  X=x_c(\rho(\tau))+\lambda(\rho(\tau))q(\tau).
\]
Using
\[
  \partial_\rho x_c=\lambda B,\qquad
  \partial_\rho\lambda=\lambda A,\qquad
  \rho_\tau=\mu^2,
\]
we get
\[
  X_\tau
  =
  \mu^2\lambda B+\mu^2\lambda Aq+\lambda q_\tau.
\]
Division by \(\Lambda=\lambda\mu\) gives
\[
  \Lambda^{-1}X_\tau
  =
  \mu B+\mu Aq+\frac{q_\tau}{\mu}.
\]

Local integrability follows from absolute continuity and positivity of
\(\mu\).  On every compact final interval, \(\mu\) is bounded above and below
away from zero, \(q\) is bounded, and \(\mu_\tau,q_\tau\in L^1\).  The
\(A\)-terms are locally integrable because
\[
  \int |A(\rho(\tau))|\mu(\tau)^2\,d\tau
  =
  \int |A(\rho)|\,d\rho,
\]
and the extra factors \(\mu\), \(\mu^{-1}\), and \(q\) are bounded on the same
compact window.  Likewise, since \(\mu\) is bounded away from zero,
\[
  \int |B(\rho(\tau))|\mu(\tau)\,d\tau
  \le
  C\int |B(\rho(\tau))|\mu(\tau)^2\,d\tau
  =
  C\int |B(\rho)|\,d\rho,
\]
so the \(B\)-term in \(b\) is locally integrable.
\end{proof}

\begin{lemma}[Renormalized equation in repaired gauge]
\label{paper6a:lem:repaired-gauge-equation}
The final variables \(V,P\) satisfy
\[
  \partial_\tau V+(V\cdot\nabla)V+\nabla P
  =
  \nu\Delta V+a(\tau)(V+Y\cdot\nabla V)+b(\tau)\cdot\nabla V,
  \qquad
  \nabla\cdot V=0
\]
distributionally on compact \(Y\)-sets.
\end{lemma}

\begin{proof}
It suffices to compute in the smooth case and then pass to distributions by
testing against compactly supported functions.  Set
\[
  z=\mu Y+q.
\]
Since \(V(Y,\tau)=\mu U(z,\rho(\tau))\) and
\(\rho_\tau=\mu^2\), the spatial derivatives are
\[
  \nabla_YV=\mu^2\nabla_zU,\qquad
  \Delta_YV=\mu^3\Delta_zU,
\]
\[
  (V\cdot\nabla_Y)V=\mu^3(U\cdot\nabla_z)U,
  \qquad
  \nabla_YP=\mu^3\nabla_zQ.
\]
Differentiation in \(\tau\) gives
\[
\begin{aligned}
  V_\tau
  &=
  \mu_\tau U
  +\mu\mu^2U_\rho
  +\mu(\mu_\tau Y+q_\tau)\cdot\nabla_zU .
\end{aligned}
\]
Using the initial renormalized equation,
\[
  U_\rho
  =
  \nu\Delta_zU-(U\cdot\nabla_z)U-\nabla_zQ
  +A(U+z\cdot\nabla_zU)+B\cdot\nabla_zU,
\]
we obtain
\[
\begin{aligned}
&V_\tau+(V\cdot\nabla_Y)V+\nabla_YP-\nu\Delta_YV \\
&\quad =
\mu_\tau U
{}+\mu^3A(U+z\cdot\nabla_zU)
{}+\mu^3B\cdot\nabla_zU
{}+\mu(\mu_\tau Y+q_\tau)\cdot\nabla_zU .
\end{aligned}
\]
Now
\[
  U=\mu^{-1}V,
  \qquad
  \nabla_zU=\mu^{-2}\nabla_YV,
  \qquad
  z=\mu Y+q.
\]
Therefore
\[
\begin{aligned}
\mu_\tau U+\mu^3A U
&=
\left(\frac{\mu_\tau}{\mu}+\mu^2A\right)V,\\
\mu^3A\,z\cdot\nabla_zU
{}+\mu(\mu_\tau Y+q_\tau)\cdot\nabla_zU
&=
\left(\mu^2A+\frac{\mu_\tau}{\mu}\right)Y\cdot\nabla_YV
{}+\left(\mu Aq+\frac{q_\tau}{\mu}\right)\cdot\nabla_YV,\\
\mu^3B\cdot\nabla_zU
&=
\mu B\cdot\nabla_YV .
\end{aligned}
\]
Combining these identities, the coefficient multiplying
\[
  V+Y\cdot\nabla_YV
\]
is
\[
  \frac{\mu_\tau}{\mu}+\mu^2A,
\]
and the coefficient multiplying \(\nabla_YV\) is
\[
  \mu B+\mu Aq+\frac{q_\tau}{\mu}.
\]
By Lemma~\ref{paper6a:lem:repaired-modulation-coefficients} these coefficients are
\(a\) and \(b\).  Hence the displayed repaired-gauge equation holds in the
distributional sense on compact \(Y\)-sets.  Incompressibility follows from
\[
  \nabla_Y\cdot V=\mu^2\nabla_z\cdot U=0.
\]
\end{proof}

\begin{lemma}[Pressure reconstruction under the charts]
\label{paper6a:lem:pressure-reconstruction-under-charts}
Assume the physical pressure satisfies
\[
  -\Delta_xp=\partial_i\partial_j(u_iu_j)
\]
in distributions.  Then the intermediate pressure satisfies
\[
  -\Delta_yQ=\partial_i\partial_j(U_iU_j)
\]
modulo functions of \(\rho\), and the final pressure satisfies
\[
  -\Delta_YP=\partial_i\partial_j(V_iV_j)
\]
modulo functions of \(\tau\).
\end{lemma}

\begin{proof}
The additive function \(c(t)\) in
\[
  p=\lambda^{-2}Q+c(t)
\]
has zero spatial derivatives.  Since
\[
  \Delta_xp=\lambda^{-4}\Delta_yQ,
  \qquad
  \partial_{x_i}\partial_{x_j}(u_iu_j)
  =
  \lambda^{-4}\partial_{y_i}\partial_{y_j}(U_iU_j),
\]
the intermediate pressure equation follows.  For the repaired variables,
\[
  P(Y)=\mu^2Q(\mu Y+q),
  \qquad
  V(Y)=\mu U(\mu Y+q).
\]
Therefore
\[
  -\Delta_YP
  =
  \mu^4(-\Delta_zQ)(z)
  =
  \mu^4\partial_{z_i}\partial_{z_j}(U_iU_j)(z)
  =
  \partial_{Y_i}\partial_{Y_j}(V_iV_j)(Y).
\]
Adding functions of time to \(Q\) or \(P\) does not change the identity.
\end{proof}

\begin{lemma}[Scaling of local \(L^r\) norms]
\label{paper6a:lem:critical-lp-scaling}
For every \(1\le r<\infty\) for which the norms are finite,
\[
  \|V(\tau)\|_{L^r_Y}^r
  =
  \Lambda(\tau)^{r-3}\|u(t(\tau))\|_{L^r_x}^r.
\]
In particular,
\[
  \|V(\tau)\|_{L^3_Y}=\|u(t(\tau))\|_{L^3_x}.
\]
Thus membership in \(L^3(\mathbb R^3)\), finiteness, infiniteness, and positivity
are invariant under the final chart.
\end{lemma}

\begin{proof}
Using
\[
  u(X+\Lambda Y,t)=\Lambda^{-1}V(Y,\tau),
  \qquad
  dx=\Lambda^3\,dY,
\]
we obtain
\[
  \|u(t)\|_{L^r_x}^r
  =
  \int |\Lambda^{-1}V(Y,\tau)|^r\Lambda^3\,dY
  =
  \Lambda^{3-r}\|V(\tau)\|_{L^r_Y}^r.
\]
Rearranging gives the stated formula.  The case \(r=3\) is scale invariant.
\end{proof}

\begin{lemma}[Regularity of scale, center, and modulation]
\label{paper6a:lem:ac-chart-regularity}
If \(\lambda,x_c,\mu,q\) are absolutely continuous and \(\mu>0\), then on
compact final time intervals the final scale \(\Lambda\), center \(X\), and
coefficients \(a,b\) are absolutely continuous or locally integrable as
specified above: \(\Lambda,X\in W^{1,1}_{\rm loc}\) and
\[
  a,b\in L^1_{\rm loc}.
\]
\end{lemma}

\begin{proof}
This follows from
Lemmas~\ref{paper6a:lem:repaired-physical-chart}
and~\ref{paper6a:lem:repaired-modulation-coefficients}.
Positivity of \(\mu\) and compactness of the time window give a positive lower
bound for \(\mu\), so the divisions by \(\mu\) in the formulas for \(a,b\) are
legitimate.
\end{proof}

\begin{theorem}[Consequences of absolutely continuous repaired-gauge charts]
\label{paper6a:thm:ac-repaired-gauge-consequences}
Assume an absolutely continuous concentration chart and an absolutely
continuous repaired-gauge path in the sense of
Definitions~\ref{paper6a:def:ac-initial-concentration-chart}
and~\ref{paper6a:def:ac-repaired-gauge-path}.  Then the
following analytic conclusions hold:
\begin{enumerate}[label=\textup{(\roman*)}]
\item the initial and repaired renormalized Navier--Stokes equations hold
distributionally on compact renormalized sets;
\item the final physical chart satisfies
\[
  u(x,t)=\Lambda(\tau)^{-1}V(Y,\tau),
  \qquad
  p(x,t)=\Lambda(\tau)^{-2}P(Y,\tau)+c(t),
  \qquad
  \frac{d\tau}{dt}=\Lambda^{-2};
\]
\item pressure reconstruction holds modulo functions of time;
\item the final modulation coefficients are the functions \(a,b\) in
Lemma~\ref{paper6a:lem:repaired-modulation-coefficients};
\item the critical \(L^3\) norm is invariant under the final chart;
\item the final scale, center, and modulation coefficients have the local
absolute-continuity and integrability stated in
Lemma~\ref{paper6a:lem:ac-chart-regularity}.
\end{enumerate}
\end{theorem}

\begin{proof}
Lemma~\ref{paper6a:lem:initial-renormalized-equation} gives the initial
renormalized equation.
Lemmas~\ref{paper6a:lem:repaired-physical-chart}
and~\ref{paper6a:lem:repaired-gauge-equation}
give the final chart and the final repaired-gauge equation with the stated
coefficients.  Lemma~\ref{paper6a:lem:pressure-reconstruction-under-charts} gives
pressure reconstruction modulo functions of time.
Lemma~\ref{paper6a:lem:critical-lp-scaling} gives critical norm
invariance, and Lemma~\ref{paper6a:lem:ac-chart-regularity} gives the
regularity inheritance.
\end{proof}

\begin{corollary}[Remaining representation hypotheses]
\label{paper6a:cor:remaining-representation-hypotheses}
Once the absolutely continuous concentration chart and absolutely continuous
repaired-gauge path are available, pressure reconstruction, modulation
coefficient realization, local \(W^{1,1}\) regularity of the chart, and
critical \(L^3\)-invariance are analytic consequences rather than separate
hypotheses.  The remaining representation hypotheses are:
\begin{enumerate}[label=\textup{(\roman*)}]
\item extraction of an absolutely continuous concentration chart from the
selected Type II branch;
\item solvability of the repaired scale and centering equations on the
admissible profile class;
\item selection of the repaired scale and center as an absolutely continuous
path satisfying
\(\rho_\tau=\mu^2\);
\item terminal admissibility of the final scale and time, including
\(\tau\to\infty\) and preservation of the Type II core scale.
\end{enumerate}
\end{corollary}

\begin{proof}
This is Theorem~\ref{paper6a:thm:ac-repaired-gauge-consequences}.  The
enumerated items are the complementary analytic hypotheses that must be
available before that theorem can be applied: one must produce the initial concentration chart,
solve the repaired gauge equations, select the resulting scale and center
absolutely continuously,
and verify that the resulting final chart is terminally admissible.
\end{proof}

\subsection{Nontrivial Critical Core and Vanishing Critical \texorpdfstring{\(L^3\)}{L3} Norm}
\label{subsec:nontrivial-core-vanishing-critical-l3}

\begin{definition}[Tail nontrivial critical core]
\label{paper6a:def:tail-nontrivial-critical-core}
A repaired-gauge Type II branch \((V,P,a,b)\) has a tail nontrivial critical
core if there exist \(R_*<\infty\), \(\eta_*>0\), and
\(\tau_*\ge\tau_0\) such that
\[
  \|V(\tau)\|_{L^3(B_{R_*})}\ge \eta_*
  \qquad
  \text{for all }\tau\ge\tau_*.
\]
Equivalently, the selected terminal core remains nonzero in the critical
topology on the renormalized tail.
\end{definition}

\begin{definition}[Subsequential nontrivial critical core]
\label{paper6a:def:subsequential-nontrivial-critical-core}
A branch has subsequential nontriviality if there exist \(R_*<\infty\),
\(\eta_*>0\), and a sequence \(\tau_n\to\infty\) such that
\[
  \|V(\tau_n)\|_{L^3(B_{R_*})}\ge \eta_*
  \qquad
  \text{for all }n.
\]
\end{definition}

\begin{definition}[Core persistence]
\label{paper6a:def:core-persistence}
Core persistence is the analytic implication that subsequential nontriviality
of a selected terminal Type II core implies tail nontriviality:
\[
  \text{subsequential nontriviality}
  \Longrightarrow
  \text{tail nontriviality}.
\]
Failure of this implication is a failure of persistence of the selected core,
not a new Type II alternative.
\end{definition}

\begin{definition}[Vanishing critical \(L^3\) alternative]
\label{paper6a:def:vanishing-critical-l3-alternative}
The vanishing critical \(L^3\) alternative for a repaired-gauge branch is
\[
  \sup_{\tau\ge\tau_0}\|V(\tau)\|_{L^3(\mathbb R^3)}<\infty,
  \qquad
  \liminf_{\tau\to\infty}\|V(\tau)\|_{L^3(\mathbb R^3)}=0.
\]
\end{definition}

\begin{theorem}[Concentration extraction gives subsequential nontriviality]
\label{paper6a:thm:concentration-gives-subsequential-nontriviality}
Assume a Type II branch has been obtained from a concentration extraction and
placed in a repaired-gauge representation.  Assume also that the extracted
branch is selected as a terminal Type II core, rather than accounted for by
scattering, exterior regularity, or loss of the extracted profile in the
limiting procedure.  Then the branch satisfies
the subsequential nontriviality condition of
Definition~\ref{paper6a:def:subsequential-nontrivial-critical-core}.
\end{theorem}

\begin{proof}
The concentration extraction supplies a nonzero profile in the Navier--Stokes
profile decomposition.  The repaired-gauge representation transports that
profile to the renormalized trajectory \(V(\tau)\) by
Navier--Stokes critical symmetries and admissible chart maps.  These symmetries
preserve the scale-critical \(L^3\) norm.

Because the extracted profile is nonzero in \(L^3_{\rm loc}\) after choosing
the terminal coordinate frame, there are \(R_*<\infty\), \(\eta_*>0\), and the
corresponding extraction times \(\tau_n\to\infty\) such that
\[
  \|V(\tau_n)\|_{L^3(B_{R_*})}\ge\eta_* .
\]
This is subsequential nontriviality.
\end{proof}

\begin{remark}[Persistence is an additional analytic hypothesis]
\label{paper6a:rem:persistence-is-additional}
Theorem~\ref{paper6a:thm:concentration-gives-subsequential-nontriviality}
does not by
itself rule out
\[
  \liminf_{\tau\to\infty}\|V(\tau)\|_{L^3(\mathbb R^3)}=0.
\]
That stronger conclusion requires the persistence hypothesis of
Definition~\ref{paper6a:def:core-persistence}; a subsequence extraction theorem is not a
uniform-in-time lower bound.
\end{remark}

\begin{theorem}[Vanishing critical \(L^3\) norm excluded by persistent nontriviality]
\label{paper6a:thm:vanishing-critical-l3-excluded}
For a repaired-gauge Type II branch, core persistence together with
subsequential nontriviality excludes the vanishing critical \(L^3\) alternative of
Definition~\ref{paper6a:def:vanishing-critical-l3-alternative}.  Equivalently, if the
vanishing critical \(L^3\) alternative occurs, then either subsequential nontriviality fails
or core persistence fails.
\end{theorem}

\begin{proof}
By core persistence and subsequential nontriviality, the branch has a tail
nontrivial critical core.  Hence there are \(R_*<\infty\), \(\eta_*>0\), and
\(\tau_*\) such that
\[
  \|V(\tau)\|_{L^3(B_{R_*})}\ge \eta_*
  \qquad
  \text{for all }\tau\ge\tau_*.
\]
Since \(B_{R_*}\subset\mathbb R^3\), this gives the global lower bound
\[
  \|V(\tau)\|_{L^3(\mathbb R^3)}\ge \eta_*
  \qquad
  \text{for all }\tau\ge\tau_*.
\]
Hence
\[
  \liminf_{\tau\to\infty}\|V(\tau)\|_{L^3(\mathbb R^3)}
  \ge \eta_*>0.
\]
This contradicts the vanishing critical \(L^3\) alternative, which requires the same
liminf to be zero.  Therefore a selected nontrivial core satisfying persistence
cannot have vanishing critical \(L^3\) norm.
\end{proof}

\begin{corollary}[Critical \(L^3\) alternatives after exclusion of vanishing norm]
\label{paper6a:cor:critical-l3-alternatives-after-vanishing-exclusion}
On selected Type II cores satisfying core persistence and subsequential
nontriviality, the vanishing critical \(L^3\) alternative is excluded from the
ordered critical \(L^3\) alternatives.  The remaining possibilities are failure
of membership in \(L^3(\mathbb R^3)\) along the terminal tail, divergence of
the critical \(L^3\) norm, or a positive finite critical \(L^3\) annulus.
\end{corollary}

\begin{proof}
Apply Theorem~\ref{paper6a:thm:vanishing-critical-l3-excluded} to exclude
the vanishing critical \(L^3\) alternative of
Definition~\ref{paper6a:def:vanishing-critical-l3-alternative}.  The other
alternatives are the remaining critical \(L^3\) outcomes.
\end{proof}

\begin{corollary}[Terminal boundedness after exclusion of vanishing norm]
\label{paper6a:cor:terminal-boundedness-after-vanishing-l3-exclusion}
In the terminal boundedness alternatives of
Theorem~\ref{paper6a:thm:terminal-boundedness-alternatives}, if core persistence and
subsequential nontriviality hold, then vanishing critical \(L^3\) norm is not a
terminal-sequence obstruction.  The remaining obstructions are failure of
sampling from the repaired-gauge trajectory, failure of membership in
\(L^3(\mathbb R^3)\) along the terminal tail, and divergence of the critical
\(L^3\) norm.
\end{corollary}

\begin{proof}
Theorem~\ref{paper6a:thm:terminal-boundedness-alternatives} imports the
critical \(L^3\) alternatives.
Corollary~\ref{paper6a:cor:critical-l3-alternatives-after-vanishing-exclusion}
excludes the vanishing critical \(L^3\) alternative on selected nontrivial cores.  The
listed obstructions are the remaining terminal alternatives.
\end{proof}

\begin{remark}[Remaining critical \(L^3\) boundary]
\label{paper6a:rem:remaining-critical-l3-boundary}
The preceding results prove only the lower-bound half of critical
normalization, modulo subsequential nontriviality and persistence.  The
remaining analytic hypotheses are persistence of the selected core, membership
in \(L^3(\mathbb R^3)\) along the terminal tail, classification of the
infinite critical \(L^3\)-norm alternative, and an upper finite critical \(L^3\) bound
on the selected repaired-gauge branch.
Once these are supplied, the positive finite critical annulus follows.
\end{remark}

\subsection{Local Caccioppoli Estimate for Absolutely Continuous Suitable Branches}
\label{subsec:ac-suitable-local-caccioppoli}

\begin{definition}[Suitable repaired-gauge Type II branch]
\label{paper6a:def:suitable-repaired-gauge-typeII-branch}
A suitable repaired-gauge Type II branch is a repaired-gauge Type II branch for
which:
\begin{enumerate}[label=\textup{(\roman*)}]
\item an absolutely continuous concentration chart and an absolutely continuous
repaired-gauge path are available;
\item the conclusions of
Theorem~\ref{paper6a:thm:ac-repaired-gauge-consequences} hold, including
pressure reconstruction and final modulation coefficients;
\item the physical branch \((u,p)\) is a local suitable weak solution on the
terminal interval:
\[
  u\in L^\infty_tL^2_{\rm loc}\cap L^2_tH^1_{\rm loc},
  \qquad
  p\in L^{3/2}_{\rm loc},
  \qquad
  \nabla\cdot u=0,
\]
the Navier--Stokes equations hold distributionally, and the physical local
energy inequality holds for every nonnegative smooth compactly supported test
function.
\end{enumerate}
\end{definition}

\begin{lemma}[Local energy inequality under absolutely continuous Type II charts]
\label{paper6a:lem:ac-pullback-local-energy}
Let \((u,p)\) be physically suitable on a compact physical cylinder, and let
\((V,P)\) be obtained from an absolutely continuous final chart
\[
  u(x,t)=\Lambda(s)^{-1}V(Y,s),
  \qquad
  p(x,t)=\Lambda(s)^{-2}P(Y,s)+c(t),
  \qquad
  Y=\frac{x-X(s)}{\Lambda(s)},
\]
with \(s=\tau\), \(dt=\Lambda(s)^2\,ds\), \(\Lambda>0\), and
\[
  X,\Lambda\in W^{1,1}_{\rm loc}.
\]
Then \((V,P)\) satisfies the renormalized local energy inequality on compact
\((Y,s)\)-cylinders:
\[
\begin{aligned}
&\int_{\mathbb R^3} \frac{|V(s_2)|^2}{2}\phi(s_2)\,dY
+\nu\int_{s_1}^{s_2}\int_{\mathbb R^3} |\nabla V|^2\phi\,dY\,ds \\
&\le
\int_{\mathbb R^3} \frac{|V(s_1)|^2}{2}\phi(s_1)\,dY
+\int_{s_1}^{s_2}\int_{\mathbb R^3}
\frac{|V|^2}{2}(\partial_s\phi+\nu\Delta\phi)\,dY\,ds \\
&\quad
+\int_{s_1}^{s_2}\int_{\mathbb R^3}
\left(\frac{|V|^2}{2}+P\right)V\cdot\nabla\phi\,dY\,ds \\
&\quad
+\int_{s_1}^{s_2}\int_{\mathbb R^3}
 a(s)\frac{|V|^2}{2}(-\phi-Y\cdot\nabla\phi)\,dY\,ds\\
&\quad
-\int_{s_1}^{s_2}\int_{\mathbb R^3}
 \frac{|V|^2}{2}b(s)\cdot\nabla\phi\,dY\,ds,
\end{aligned}
\]
for every nonnegative
\(\phi\in C^\infty([s_1,s_2]\times\mathbb R^3)\) that is compactly supported
in the spatial variable.  Here \(a,b\) are the final modulation coefficients.
\end{lemma}

\begin{proof}
First suppose \(X,\Lambda\) are \(C^1\).  Use the physical test function
\[
  \psi(x,t)=\Lambda(s)^{-1}
  \phi\left(\frac{x-X(s)}{\Lambda(s)},s\right),
  \qquad t=t(s),
\]
where \(dt=\Lambda^2ds\).  The spatial support of \(\phi\) is compact, and
\(\Lambda\) is bounded above and below on the selected compact time interval,
so \(\psi\) is an admissible compactly supported test function in the physical
local energy inequality.

Let
\[
  Y=\frac{x-X(s)}{\Lambda(s)}.
\]
At fixed \(x\),
\[
  \frac{dY}{ds}
  =
  -\frac{X_s}{\Lambda}-\frac{\Lambda_s}{\Lambda}Y .
\]
Consequently,
\[
  \nabla_x\psi=\Lambda^{-2}\nabla_Y\phi,\qquad
  \Delta_x\psi=\Lambda^{-3}\Delta_Y\phi,
\]
and, since \(ds/dt=\Lambda^{-2}\),
\[
  \partial_t\psi
  =
  \Lambda^{-3}
  \left[
  \partial_s\phi
  -\frac{\Lambda_s}{\Lambda}(\phi+Y\cdot\nabla\phi)
  -\frac{X_s}{\Lambda}\cdot\nabla\phi
  \right].
\]
Changing variables
\[
  x=X(s)+\Lambda(s)Y,\qquad dt=\Lambda^2ds,
\]
and using
\[
  u=\Lambda^{-1}V,\qquad p=\Lambda^{-2}P+c(t),
\]
each term in the physical local energy inequality has the following
renormalized form.  The endpoint and dissipation terms transform as
\[
  \int \frac{|u(t_i)|^2}{2}\psi(t_i)\,dx
  =
  \int \frac{|V(s_i)|^2}{2}\phi(s_i)\,dY,
\]
and
\[
  \nu\int_{t_1}^{t_2}\int |\nabla_xu|^2\psi\,dx\,dt
  =
  \nu\int_{s_1}^{s_2}\int |\nabla_YV|^2\phi\,dY\,ds.
\]
The test-function derivative term becomes
\[
\begin{aligned}
&\int_{t_1}^{t_2}\int \frac{|u|^2}{2}
(\partial_t\psi+\nu\Delta_x\psi)\,dx\,dt\\
&\quad =
\int_{s_1}^{s_2}\int \frac{|V|^2}{2}
\left[
\partial_s\phi+\nu\Delta_Y\phi
-a(\phi+Y\cdot\nabla\phi)-b\cdot\nabla\phi
\right]\,dY\,ds,
\end{aligned}
\]
where
\[
  a=\frac{\Lambda_s}{\Lambda},
  \qquad
  b=\frac{X_s}{\Lambda}.
\]
The transport and pressure term becomes
\[
\begin{aligned}
&\int_{t_1}^{t_2}\int
\left(\frac{|u|^2}{2}+p\right)u\cdot\nabla_x\psi\,dx\,dt\\
&\quad =
\int_{s_1}^{s_2}\int_{\mathbb R^3}
\left(\frac{|V|^2}{2}+P\right)V\cdot\nabla_Y\phi\,dY\,ds
{}+\int_{s_1}^{s_2}\Lambda(s)^2c(t(s))
\left(\int_{\mathbb R^3} V\cdot\nabla_Y\phi\,dY\right)\,ds .
\end{aligned}
\]
The pressure constant drops out
because
\[
  \int_{\mathbb R^3} V\cdot\nabla_Y\phi\,dY
  =
  -\int_{\mathbb R^3} \phi\,\nabla_Y\cdot V\,dY=0.
\]
Combining the transformed terms gives the displayed renormalized local energy
inequality.

For \(X,\Lambda\in W^{1,1}_{\rm loc}\), the same identity follows at the level
of test functions.  On the selected compact time interval, \(\Lambda\) has
positive lower and finite upper bounds.  The pulled-back test
\[
  \psi(x,t)=\Lambda(s)^{-1}\phi((x-X(s))/\Lambda(s),s)
\]
is compactly supported, nonnegative, bounded, has bounded spatial derivatives
on the compact cylinder, and has time derivative in \(L^1_tL^\infty_x\)
because \(X',\Lambda'\in L^1\) and \(\phi\) is smooth compactly supported.
The physical local energy inequality, initially stated for \(C_c^\infty\)
tests, extends to this class by standard mollification of the test function:
the suitability classes give the local integrability needed for every term in
the inequality, and the mollified tests keep support inside a slightly larger
compact cylinder.  Applying the extended inequality to \(\psi\) and changing
variables gives the displayed renormalized inequality.  The only derivatives
of the chart that appear are the \(L^1\)-functions \(X',\Lambda'\), hence the
modulation terms are locally integrable against local \(|V|^2\).
\end{proof}

\begin{lemma}[Pressure and modulation for suitable repaired-gauge branches]
\label{paper6a:lem:suitable-gauge-pressure-modulation}
For a suitable repaired-gauge branch, the pressure identity
\[
  -\Delta P=\partial_i\partial_j(V_iV_j)
\]
holds modulo functions of \(\tau\), and the final coefficients \(a,b\) in the
renormalized equation are measurable and locally integrable on compact
renormalized windows.
\end{lemma}

\begin{proof}
This is the pressure and modulation part of
Theorem~\ref{paper6a:thm:ac-repaired-gauge-consequences}.  Pressure pullback
comes from the physical pressure equation and the chart/gauge scaling.  The
final coefficients are forced by the absolutely continuous scale-translation
gauge transform:
\[
  a=\mu^2A+\frac{\mu_\tau}{\mu},
  \qquad
  b=\mu B+\mu Aq+\frac{q_\tau}{\mu}.
\]
Since the chart and gauge parameters are absolutely continuous on compact
windows, these coefficients are locally integrable.
\end{proof}

\begin{theorem}[Caccioppoli estimate for suitable absolutely continuous repaired-gauge branches]
\label{paper6a:thm:ac-local-caccioppoli}
Every suitable repaired-gauge Type II branch satisfies the local renormalized
Caccioppoli estimate on each compact repaired-gauge cylinder.  The constants
may depend on the cylinder, the cutoff, and local upper and lower bounds for
the chart.  No local Caffarelli--Kohn--Nirenberg smallness, global windowed
\(H^1\), critical \(L^3\)-norm, tightness, bounded modulation, finite energy at
infinity, or whole-space compactness conclusion is asserted.
\end{theorem}

\begin{proof}
By the suitable repaired-gauge hypotheses and
Theorem~\ref{paper6a:thm:ac-repaired-gauge-consequences}, the branch has an
absolutely continuous final chart, an admissible absolutely continuous
repaired-gauge path, pressure reconstruction, and final modulation
coefficients.  Lemma~\ref{paper6a:lem:ac-pullback-local-energy} pulls the physical
local energy inequality to the repaired variables.
Lemma~\ref{paper6a:lem:suitable-gauge-pressure-modulation} supplies the pressure and
coefficient identities needed to interpret the renormalized inequality in the
same PDE class used by the local windowed estimates.

Choose a standard compactly supported Caccioppoli cutoff
\[
  \phi(Y,\tau)=\eta(\tau)\zeta(Y)^2.
\]
Let \(I=(\tau_1,\tau_2)\), let \(I'\Subset I\), and let
\(0<r<R\).  Take \(0\le\eta\le1\), \(\eta\in C_c^\infty(I)\), with
\(\eta\equiv1\) on \(I'\), and take \(0\le\zeta\le1\),
\(\zeta\in C_c^\infty(B_R)\), with \(\zeta\equiv1\) on \(B_r\).  Applying
Lemma~\ref{paper6a:lem:ac-pullback-local-energy} on \((\tau_1,\sigma)\) and
then taking the essential supremum over \(\sigma\in I'\) controls the endpoint
term.  Applying the same inequality on the full interval \(I\) controls the
dissipation term.  Adding the two estimates gives
\[
\begin{aligned}
&\operatorname*{ess\,sup}_{\tau\in I'}
  \int_{B_r}\frac{|V(\tau)|^2}{2}\,dY
  +\nu\int_{I'}\int_{B_r}|\nabla V|^2\,dY\,d\tau \\
&\le C
\int_I\int_{B_R}\frac{|V|^2}{2}
   \left(|\eta'|\zeta^2+\nu\eta|\Delta(\zeta^2)|\right)\,dY\,d\tau \\
&\quad
+C\int_I\int_{B_R}
   \left(\frac{|V|^2}{2}+|P|\right)|V|\,\eta|\nabla(\zeta^2)|\,dY\,d\tau\\
&\quad
+C\int_I |a(\tau)|
  \int_{B_R}\frac{|V|^2}{2}\eta
    \left(\zeta^2+|Y|\,|\nabla(\zeta^2)|\right)\,dY\,d\tau\\
&\quad
+C\int_I |b(\tau)|
  \int_{B_R}\frac{|V|^2}{2}\eta|\nabla(\zeta^2)|\,dY\,d\tau .
\end{aligned}
\]
Here \(C\) is an absolute constant.  Since the cutoff derivatives are bounded
by constants depending only on \(I,I',r,R\), this is the local Caccioppoli
estimate on the compact renormalized cylinder.

No global estimate is used in this step.  On the compact cylinder under
consideration, the absolutely continuous chart has
\[
  0<\Lambda_-\leq \Lambda(\tau)\leq \Lambda_+<\infty.
\]
Pulling back physical suitability therefore gives, for every compact spatial
set \(K\),
\[
  V\in L^\infty_\tau L^2_Y(K)\cap L^2_\tau H^1_Y(K),
  \qquad
  P\in L^{3/2}_{\tau,Y}(K).
\]
Local Sobolev interpolation on the finite cylinder gives
\[
  V\in L^{10/3}_{\tau,Y}(K),
\]
hence \(V\in L^3_{\tau,Y}(K)\).  Thus \(PV\cdot\nabla\phi\) is integrable by
Hölder's inequality.  The representation identities give
\[
  a,b\in L^1_{\rm loc}(d\tau),
\]
and
\[
  \tau\mapsto\int_K |V(Y,\tau)|^2\,dY
\]
is locally essentially bounded; therefore the scale and translation modulation
terms in the pulled-back local energy inequality are integrable on the same
compact cylinder.  This proves the local Caccioppoli estimate, but not a
global windowed \(H^1\) or critical \(L^3\)-norm bound.
\end{proof}

\begin{corollary}[No separate Caccioppoli obstruction for suitable absolutely continuous branches]
\label{paper6a:cor:no-independent-caccioppoli-obstruction-ac-suitable}
On suitable repaired-gauge branches, failure of local Caccioppoli regularity,
failure of pressure reconstruction, and failure of local gauge regularity are
not separate obstructions to the local \(H^1\) analysis.  The remaining
hypotheses for the
uniform windowed \(H^1\) estimates are the separate bounded critical \(L^3\)
norm and modulation bounds required by those estimates.
\end{corollary}

\begin{proof}
Theorem~\ref{paper6a:thm:ac-local-caccioppoli} gives local Caccioppoli
regularity.  Theorem~\ref{paper6a:thm:ac-repaired-gauge-consequences} gives
pressure reconstruction and the final modulation coefficients from the
final chart and absolutely continuous gauge path.  Thus the listed
failures cannot occur separately on suitable repaired-gauge branches.  The
uniform windowed \(H^1\) estimate may still require bounded critical norm and
bounded modulation beyond the local Caccioppoli statement.
\end{proof}

\subsection{Caccioppoli Regularity Criterion}

\begin{definition}[Suitable physical hypotheses]\label{paper6a:def:suitable-physical-hypotheses}
Let \(u,p\) be a local suitable weak solution of the three-dimensional
Navier--Stokes equations on \(\mathbb R^3\times(0,T^*)\).  Thus
\[
  u\in L^\infty_tL^2_{\rm loc}\cap L^2_tH^1_{\rm loc},
  \qquad
  p\in L^{3/2}_{\rm loc},
  \qquad \nabla\cdot u=0,
\]
the equations hold distributionally, and for every nonnegative
\(\psi\in C_c^\infty(\mathbb R^3\times(0,T^*))\) and a.e.
\(0<t_1<t_2<T^*\),
\[
\begin{aligned}
&\int_{\mathbb R^3}\frac{|u(t_2)|^2}{2}\psi(t_2)\,dx
+\nu\int_{t_1}^{t_2}\int_{\mathbb R^3}|\nabla u|^2\psi\,dx\,dt \\
&\le
\int_{\mathbb R^3}\frac{|u(t_1)|^2}{2}\psi(t_1)\,dx
+\int_{t_1}^{t_2}\int_{\mathbb R^3}
\frac{|u|^2}{2}(\partial_t\psi+\nu\Delta\psi)\,dx\,dt\\
&\quad
+\int_{t_1}^{t_2}\int_{\mathbb R^3}
\left(\frac{|u|^2}{2}+p\right)u\cdot\nabla\psi\,dx\,dt .
\end{aligned}
\]
\end{definition}

\begin{lemma}[Renormalized local energy inequality]\label{paper6a:lem:renormalized-lei}
Assume the suitable physical hypotheses of
\cref{paper6a:def:suitable-physical-hypotheses}.  Let
\[
  u(x,t)=\lambda(t)^{-1}
  V\left(\frac{x-x_c(t)}{\lambda(t)},\tau(t)\right),
  \qquad
  \frac{d\tau}{dt}=\lambda(t)^{-2},
\]
with
\[
  P(y,\tau)=\lambda(t)^2p(x_c(t)+\lambda(t)y,t),
\]
where \(\lambda>0\), \(x_c\), and \(\tau\) are \(C^1\) on compact windows.
Assume that \(V,P\) satisfy the repaired-gauge equation
\[
  \partial_\tau V+(V\cdot\nabla)V+\nabla P
  =
  \nu\Delta V+a(\tau)(V+y\cdot\nabla V)+b(\tau)\cdot\nabla V,
  \qquad \nabla\cdot V=0 .
\]
Then, for every nonnegative
\(\phi\in C_c^\infty(\mathbb R^3\times(\tau_0,\infty))\) and a.e.
\(\tau_1<\tau_2\) in the corresponding renormalized time interval,
\[
\begin{aligned}
&\int_{\mathbb R^3}\frac{|V(\tau_2)|^2}{2}\phi(\tau_2)\,dy
+\nu\int_{\tau_1}^{\tau_2}\int_{\mathbb R^3}|\nabla V|^2\phi\,dy\,d\tau \\
&\le
\int_{\mathbb R^3}\frac{|V(\tau_1)|^2}{2}\phi(\tau_1)\,dy
+\int_{\tau_1}^{\tau_2}\int_{\mathbb R^3}
\frac{|V|^2}{2}(\partial_\tau\phi+\nu\Delta\phi)\,dy\,d\tau\\
&\quad
+\int_{\tau_1}^{\tau_2}\int_{\mathbb R^3}
\left(\frac{|V|^2}{2}+P\right)V\cdot\nabla\phi\,dy\,d\tau\\
&\quad
+\int_{\tau_1}^{\tau_2}\int_{\mathbb R^3}
a(\tau)\frac{|V|^2}{2}\bigl(-\phi-y\cdot\nabla\phi\bigr)\,dy\,d\tau\\
&\quad
-\int_{\tau_1}^{\tau_2}\int_{\mathbb R^3}
\frac{|V|^2}{2}\,b(\tau)\cdot\nabla\phi\,dy\,d\tau .
\end{aligned}
\]
\end{lemma}

\begin{proof}
For smooth solutions, multiply the repaired-gauge equation by \(V\phi\) and
integrate by parts.  Put \(e=|V|^2/2\).  Since \(\nabla\cdot V=0\),
\[
  V\cdot(V\cdot\nabla V)=V\cdot\nabla e=\nabla\cdot(eV),
  \qquad
  V\cdot\nabla P=\nabla\cdot(PV),
\]
and
\[
  V\cdot(y\cdot\nabla V)=y\cdot\nabla e,
  \qquad
  V\cdot(b\cdot\nabla V)=b\cdot\nabla e .
\]
The scale contribution is
\[
  \int a(2e+y\cdot\nabla e)\phi\,dy
  =
  \int ae(-\phi-y\cdot\nabla\phi)\,dy,
\]
because
\[
  \int y\cdot\nabla e\,\phi\,dy
  =
  -\int e\,\nabla\cdot(y\phi)\,dy
  =
  -\int e(3\phi+y\cdot\nabla\phi)\,dy .
\]
The translation contribution is
\[
  \int b\cdot\nabla e\,\phi\,dy
  =
  -\int e\,b\cdot\nabla\phi\,dy,
\]
and the viscous term gives
\[
  \int \nu\Delta V\cdot V\phi\,dy
  =
  -\nu\int|\nabla V|^2\phi\,dy
  +\nu\int e\Delta\phi\,dy .
\]
Combining these identities gives the displayed equality in the smooth case.

For suitable weak solutions, apply the physical local energy inequality with
\[
  \psi(x,t)=
  \phi\left(\frac{x-x_c(t)}{\lambda(t)},\tau(t)\right).
\]
Use the change of variables
\[
  x=x_c(t)+\lambda(t)y,
  \qquad
  dt=\lambda(t)^2\,d\tau .
\]
The \(C^1\) regularity of \(\lambda\) and \(x_c\) justifies the chain rule for
the test function, and the Jacobian factors are those of the
Navier--Stokes critical scaling.  Passing to the variables \(y,\tau\) gives the
same formula with equality replaced by \(\le\).
\end{proof}

\begin{lemma}[Caccioppoli estimate from the renormalized local energy inequality]
\label{paper6a:lem:caccioppoli-from-renormalized-lei}
Assume \cref{paper6a:lem:renormalized-lei} on a compact renormalized cylinder
\(B_{2R}\times I\), bounded modulation on \(I\), and local bounds
\[
  V\in L^\infty(I;L^3(B_{2R})),
  \qquad
  P\in L^{3/2}(I\times B_{2R}).
\]
Then the compact-cylinder renormalized Caccioppoli estimate used in
\cref{paper7:prop:uniform-windowed-gradient} is justified.
\end{lemma}

\begin{proof}
Choose a standard nonnegative cutoff
\[
  \phi(y,\tau)=\eta(\tau)\zeta(y)^2,
\]
where \(\zeta\equiv1\) on the inner ball and \(\zeta\) is supported in the
outer ball.  Insert this test function in
\cref{paper6a:lem:renormalized-lei}.  The left-hand side contains
\[
  \nu\int |\nabla V|^2\eta\zeta^2 .
\]
Every term on the right-hand side is finite under the stated hypotheses.  The term
\[
  \int |V|^2|\partial_\tau\phi+\Delta\phi|
\]
is controlled by \(V\in L^\infty(I;L^3(B_{2R}))\) and finite measure.  The
convective term
\[
  \int |V|^3|\nabla\phi|
\]
is controlled by the same local critical bound.  For the pressure term, choose
a time-dependent constant \(c(\tau)\).  Since
\[
  \int c(\tau)V\cdot\nabla\phi\,dy
  =
  -\int c(\tau)\phi\,\nabla\cdot V\,dy=0,
\]
Hölder's inequality controls
\[
  \int |P-c(\tau)|\,|V|\,|\nabla\phi|
\]
by \(P-c(\tau)\in L^{3/2}(I\times B_{2R})\) and
\(V\in L^\infty(I;L^3(B_{2R}))\).  The
modulation terms are bounded by
\[
  \|a\|_{L^\infty}+\|b\|_{L^\infty}
\]
times the local \(L^2\)-mass, and the local \(L^2\)-mass is controlled on
bounded balls by local \(L^3\)-mass.  These estimates give the
renormalized Caccioppoli estimate on compact cylinders.
\end{proof}

\begin{proposition}[Caccioppoli regularity criterion]\label{paper6a:prop:caccioppoli-regularity-criterion}
Assume the suitable physical hypotheses, the repaired-gauge representation, pressure
reconstruction modulo functions of time, and \(C^1\) scale and center functions
on compact renormalized windows.  Then the compact-cylinder Caccioppoli
regularity hypothesis used in the local windowed \(H^1\) estimate holds.
\end{proposition}

\begin{proof}
The repaired-gauge representation writes the suitable weak solution in the
variables \(V,P,a,b\).  \Cref{paper6a:lem:renormalized-lei} gives the
renormalized local energy inequality, and
\cref{paper6a:lem:caccioppoli-from-renormalized-lei} shows that this inequality
justifies the renormalized Caccioppoli estimate on every compact cylinder.
\end{proof}

\begin{corollary}[Caccioppoli estimates give windowed control under boundedness hypotheses]
\label{paper6a:cor:caccioppoli-windowed-control}
Assume the suitable physical hypotheses, the repaired-gauge representation, pressure
reconstruction, a bounded critical \(L^3\) norm on compact cylinders, and
bounded modulation.  Then local windowed \(H^1\) control follows from
\cref{paper7:prop:uniform-windowed-gradient}.  Thus failure of the Caccioppoli
regularity hypothesis is not an independent rough-core alternative under these
hypotheses.
\end{corollary}

\begin{proof}
\Cref{paper6a:prop:caccioppoli-regularity-criterion} supplies the Caccioppoli
regularity hypothesis.  The bounded critical norm, pressure reconstruction, and
bounded modulation are the remaining hypotheses used in
\cref{paper7:prop:uniform-windowed-gradient}.  Therefore the local windowed
\(H^1\) estimate holds.
\end{proof}

\subsection{Terminal Profile Compactness Hypotheses}

\begin{assumption}[Critical Navier--Stokes profile decomposition]
\label{paper6a:ass:critical-ns-profile-decomposition}
Let \(u_{0,n}\) be a bounded sequence in \(L^3_\sigma(\mathbb R^3)\), or in a
stronger admissible critical profile space for which the standard
Navier--Stokes profile theorem and perturbation theorem are available.  After
passing to a subsequence, there are profiles
\(\phi^j\), scales \(\lambda_{j,n}>0\), centers \(x_{j,n}\in\mathbb R^3\), and
remainders \(r_n^J\) such that
\[
  u_{0,n}
  =
  \sum_{j=1}^J \Lambda_{\lambda_{j,n},x_{j,n}}\phi^j+r_n^J,
  \qquad
  (\Lambda_{\lambda,x_0}f)(x)
  =
  \lambda^{-1}f\left(\frac{x-x_0}{\lambda}\right),
\]
with pairwise orthogonality of the parameters and critical mass decoupling
\[
  \|u_{0,n}\|_{L^3}^3
  =
  \sum_{j=1}^J \|\phi^j\|_{L^3}^3
  +\|r_n^J\|_{L^3}^3
  +o_n(1).
\]
Moreover, if the corresponding nonlinear profiles are evolved and the small
critical profiles are removed by the small-data theory, then
\[
  \lim_{J\to\infty}\limsup_{n\to\infty}
  \|e^{t\Delta}r_n^J\|_{X_{\rm Kato}}=0,
\]
and the nonlinear remainder is perturbatively small on every compact
profile-time window.  Finally, if after removal of all extracted profiles a
terminal coordinate frame still contains nonzero local critical mass, then,
after passing to a further subsequence, the profile theorem produces an
additional nonzero profile associated to that coordinate frame.
\end{assumption}

\begin{definition}[Bounded critical terminal sequences]
\label{paper6a:def:bounded-critical-terminal-sequences}
The bounded critical terminal-sequence hypothesis is the assertion that every
terminal active sequence used in the terminal profile analysis consists of
divergence-free data \(u_{0,n}\) satisfying
\[
  \sup_n \|u_{0,n}\|_{L^3(\mathbb R^3)}<\infty,
\]
or satisfying the corresponding uniform bound in a stronger admissible critical
profile space to which
\cref{paper6a:ass:critical-ns-profile-decomposition} applies.
\end{definition}

\begin{theorem}[Kato small-data critical stability]
\label{paper6a:thm:kato-small-data-critical-stability}
There is \(\varepsilon_{\rm sd}>0\) such that, if
\(f\in L^3_\sigma(\mathbb R^3)\) and
\(\|f\|_{L^3}\le\varepsilon_{\rm sd}\), then the Navier--Stokes equations have a
global mild solution in a standard critical Kato class \(X_{\rm Kato}\),
\[
  U\in C([0,\infty);L^3_\sigma(\mathbb R^3))\cap X_{\rm Kato},
  \qquad
  \|U\|_{X_{\rm Kato}}\le C\|f\|_{L^3}.
\]
The solution map is Lipschitz on this small ball, and perturbations whose
linear heat evolutions are small in \(X_{\rm Kato}\) remain perturbative in the
same topology.
\end{theorem}

\begin{proof}
This is the classical Kato fixed-point theorem in the critical \(L^3\) space.
Writing the mild formulation as
\[
  U(t)=e^{t\Delta}f
  -\int_0^t e^{(t-s)\Delta}\mathbb P\nabla\cdot(U\otimes U)(s)\,ds,
\]
the heat estimate and bilinear estimate in \(X_{\rm Kato}\) give
\[
  \|B(U,V)\|_{X_{\rm Kato}}
  \le C\|U\|_{X_{\rm Kato}}\|V\|_{X_{\rm Kato}} .
\]
For \(\|f\|_{L^3}\le\varepsilon_{\rm sd}\), the map is a contraction in the
ball \(\|U\|_{X_{\rm Kato}}\le 2C\|f\|_{L^3}\).  The same estimate applied to
the difference of two mild solutions gives Lipschitz dependence.  If a
perturbing datum has small heat evolution in \(X_{\rm Kato}\), the Duhamel
difference inequality is absorbed by the left-hand side after the same
smallness choice.
\end{proof}

\begin{definition}[Small-data critical stability]
\label{paper6a:def:small-data-critical-stability}
Small-data critical stability means that every divergence-free datum whose
critical size is below \(\varepsilon_{\rm sd}\), and every perturbation whose
linear heat flow is small in \(X_{\rm Kato}\), generates only a perturbative
Navier--Stokes evolution on compact profile-time windows.
\end{definition}

\begin{corollary}[Kato theorem gives small-data stability]
\label{paper6a:cor:kato-gives-small-data-stability}
\Cref{paper6a:thm:kato-small-data-critical-stability} implies the
small-data critical stability condition of
\cref{paper6a:def:small-data-critical-stability}.
\end{corollary}

\begin{proof}
The theorem gives global critical control for data below
\(\varepsilon_{\rm sd}\) and Lipschitz stability under small heat-flow
perturbations.  Restricting the resulting global estimate to any compact
profile-time window gives the stated stability condition.
\end{proof}

\begin{definition}[Terminal critical profile hypotheses]
\label{paper6a:def:terminal-critical-profile-hypotheses}
The terminal critical profile hypotheses consist of the bounded critical
terminal-sequence hypothesis
\cref{paper6a:def:bounded-critical-terminal-sequences}, the small-data
stability condition
\cref{paper6a:def:small-data-critical-stability}, and the critical
Navier--Stokes profile theorem
\cref{paper6a:ass:critical-ns-profile-decomposition}.
\end{definition}

\begin{theorem}[Terminal critical local compactness]
\label{paper6a:thm:terminal-critical-local-compactness}
Assume the terminal critical profile hypotheses.  Then the terminal profile
analysis is complete in the following sense: after passing to subsequences,
all non-scattering critical profiles appear as active profiles, and the
remaining terminal component is perturbatively small on compact profile-time
windows.
\end{theorem}

\begin{proof}
Fix a terminal active sequence.  By
\cref{paper6a:def:bounded-critical-terminal-sequences} and
\cref{paper6a:ass:critical-ns-profile-decomposition}, a critical profile
decomposition is available after passing to a subsequence:
\[
  u_{0,n}
  =
  \sum_{j=1}^J \Lambda_{\lambda_{j,n},x_{j,n}}\phi^j+r_n^J .
\]
The mass decoupling identity shows that only finitely many profiles can have
critical norm at least \(\varepsilon_{\rm sd}\).  Profiles with smaller
critical norm generate perturbative evolutions by
\cref{paper6a:cor:kato-gives-small-data-stability}.  The finitely many
remaining profiles are the non-scattering active profiles.

It remains to exclude hidden terminal mass.  Suppose that, after removing all
active profiles and all small-data perturbative profiles, a terminal
coordinate frame still sees nonzero local critical mass.  The final clause of
\cref{paper6a:ass:critical-ns-profile-decomposition} then produces an
additional nonzero profile associated to that frame.  If its critical norm is
below \(\varepsilon_{\rm sd}\), the Kato stability estimate makes its nonlinear
contribution perturbative, contradicting the assumed nonzero terminal
critical mass.  If its critical norm is at least \(\varepsilon_{\rm sd}\), then
it is another active profile, contradicting the claim that all active profiles
had already been removed.  Hence no hidden terminal critical mass remains, and
the residual component is perturbatively small on compact profile-time
windows.
\end{proof}

\begin{corollary}[Profile completeness from terminal critical profile hypotheses]
\label{paper6a:cor:profile-completeness-from-critical-hypotheses}
The terminal profile completeness hypothesis in the final argument may be replaced
by the terminal critical profile hypotheses of
\cref{paper6a:def:terminal-critical-profile-hypotheses}.
\end{corollary}

\begin{proof}
\Cref{paper6a:thm:terminal-critical-local-compactness} gives the asserted
terminal profile completeness.
\end{proof}

\begin{remark}[Boundary of the terminal compactness hypothesis]
\label{paper6a:rem:terminal-compactness-boundary}
The preceding compactness conclusion applies only to terminal sequences with a
uniform bound in \(L^3_\sigma\), or in a stronger admissible critical profile
space.  If a terminal sequence is not a compact-cylinder sequence in such a
space, then the failure is a critical-norm blowup defect or a failure of the
coordinate-frame construction, not a failure of the critical profile theorem.
\end{remark}

\begin{definition}[Terminal sequences sampled from the renormalized orbit]
\label{paper6a:def:terminal-sequences-sampled-from-orbit}
A terminal active sequence is sampled from the renormalized orbit if it is of
the form
\[
  u_{0,n}=\Lambda_{\lambda_n,x_n}V(\tau_n),
  \qquad
  \tau_n\to\infty,\qquad \lambda_n>0,\qquad x_n\in\mathbb R^3,
\]
or by the equivalent inverse critical change of variables, where
\[
  (\Lambda_{\lambda,x_0}f)(x)
  =
  \lambda^{-1}f\left(\frac{x-x_0}{\lambda}\right).
\]
\end{definition}

\begin{lemma}[Critical symmetry preserves the \(L^3\) norm]
\label{paper6a:lem:critical-symmetry-preserves-l3}
For every \(f\in L^3(\mathbb R^3)\), \(\lambda>0\), and
\(x_0\in\mathbb R^3\),
\[
  \|\Lambda_{\lambda,x_0}f\|_{L^3}=\|f\|_{L^3}.
\]
\end{lemma}

\begin{proof}
By the change of variables \(x=x_0+\lambda y\),
\[
  \|\Lambda_{\lambda,x_0}f\|_{L^3}^3
  =
  \int_{\mathbb R^3}\lambda^{-3}
  \left|f\left(\frac{x-x_0}{\lambda}\right)\right|^3\,dx
  =
  \int_{\mathbb R^3}|f(y)|^3\,dy .
\]
Taking cube roots gives the claim.
\end{proof}

\begin{theorem}[Bounded terminal sequences from the critical annulus]
\label{paper6a:thm:bounded-terminal-sequences-from-l3}
Assume the renormalized orbit satisfies
\[
  0<\inf_{\tau}\|V(\tau)\|_{L^3}
  \le
  \sup_{\tau}\|V(\tau)\|_{L^3}
  <\infty,
\]
on the terminal interval under consideration.  If every terminal active
sequence is sampled from the renormalized orbit, then every terminal active
sequence is bounded in \(L^3_\sigma(\mathbb R^3)\).
\end{theorem}

\begin{proof}
Let \(u_{0,n}=\Lambda_{\lambda_n,x_n}V(\tau_n)\).  By
\cref{paper6a:lem:critical-symmetry-preserves-l3},
\[
  \|u_{0,n}\|_{L^3}=\|V(\tau_n)\|_{L^3}.
\]
The assumed finite upper bound on the terminal interval gives
\(\sup_n\|u_{0,n}\|_{L^3}<\infty\).  If the positive lower bound is also used,
the sequence is nontrivial in the critical norm.
\end{proof}

\begin{theorem}[Terminal boundedness alternatives]
\label{paper6a:thm:terminal-boundedness-alternatives}
For a candidate terminal active sequence in the Type II analysis, the following
alternatives exhaust the possible boundedness outcomes, after passing to a
subsequence when necessary:
\begin{enumerate}[label=\textup{(\roman*)}]
\item the sequence is not sampled from the renormalized orbit;
\item the terminal \(L^3\) norm of \(V\) is not defined in the asserted
critical class;
\item \(\|V(\tau_n)\|_{L^3}\to\infty\) along a terminal sampling subsequence;
\item \(\|V(\tau_n)\|_{L^3}\to0\) along a terminal sampling subsequence;
\item after excluding the preceding alternatives, the terminal active sequence
satisfies a positive finite critical annulus and is bounded in
\(L^3_\sigma(\mathbb R^3)\).
\end{enumerate}
\end{theorem}

\begin{proof}
If the terminal sequence is not sampled from the renormalized orbit, then
\textup{(i)} holds.  If the critical norm is not defined in the asserted class,
then \textup{(ii)} holds.  If the critical norm is defined but has no positive
finite terminal annulus, then either its limsup is infinite, which gives
\textup{(iii)} after passing to a subsequence, or its liminf is zero, which
gives \textup{(iv)} after passing to a subsequence.  Otherwise the positive
finite annulus condition of
\cref{paper6a:thm:bounded-terminal-sequences-from-l3} holds, and
\textup{(v)} follows.
\end{proof}

\begin{corollary}[Boundedness from orbit sampling]
\label{paper6a:cor:boundedness-from-orbit-sampling}
In the presence of terminal orbit sampling and a positive finite terminal
\(L^3\) annulus, the bounded critical terminal-sequence hypothesis in
\cref{paper6a:def:bounded-critical-terminal-sequences} follows from
\cref{paper6a:thm:bounded-terminal-sequences-from-l3}.
\end{corollary}

\begin{proof}
The theorem gives the required uniform \(L^3\) bound for every terminal active
sequence.
\end{proof}

\begin{definition}[Terminal coordinate-frame construction]
\label{paper6a:def:terminal-coordinate-frame-construction}
The terminal coordinate-frame construction means that every terminal active
frame used in the analysis is obtained from the repaired-gauge orbit as
follows.  Choose sampling times \(\tau_n\to\infty\) and critical symmetry
parameters \((\lambda_n,x_n)\); form the data by applying the critical
coordinate-frame map to \(V(\tau_n)\); group comparable profiles based at the
same physical point before selecting the terminal active frame.  No external
data are introduced.
\end{definition}

\begin{theorem}[Terminal coordinate-frame construction from profile selection]
\label{paper6a:thm:terminal-coordinate-frame-construction-profile-selection}
Assume terminal active frames are selected by first partitioning active
profiles by physical concentration point, then grouping comparable profiles at
the same point into compound profiles, and finally choosing a minimal-scale
active class at that point.  Then the terminal coordinate-frame construction
of \cref{paper6a:def:terminal-coordinate-frame-construction} holds for the
selected terminal frames.
\end{theorem}

\begin{proof}
By the stated selection rule, each terminal active frame is attached to a
compound active class after all comparable profiles at the same physical
point have been grouped.  The scale and center of the terminal frame are the
scale and center of this compound class.  The terminal data are then obtained
by applying the Navier--Stokes critical change of variables to the
renormalized branch along the chosen sampling sequence \(\tau_n\to\infty\).
Thus the construction uses only the repaired-gauge orbit and the
Navier--Stokes critical symmetries, and it satisfies all clauses of
\cref{paper6a:def:terminal-coordinate-frame-construction}.
\end{proof}

\begin{theorem}[Terminal coordinate frames are sampled from the orbit]
\label{paper6a:thm:terminal-coordinate-frames-sampled}
Assume the terminal coordinate-frame construction.  Then every terminal active
sequence belongs to the class described in
\cref{paper6a:def:terminal-sequences-sampled-from-orbit}.
\end{theorem}

\begin{proof}
By definition of the terminal coordinate-frame construction, each terminal
active frame is obtained by choosing times \(\tau_n\), selecting scale and
center parameters \((\lambda_n,x_n)\), and applying the Navier--Stokes critical
change of variables to the renormalized solution \(V(\tau_n)\).  Hence the
associated terminal data have the form
\[
  u_{0,n}=\Lambda_{\lambda_n,x_n}V(\tau_n),
\]
or the equivalent inverse formulation, which is the orbit-sampling
condition.
\end{proof}

\begin{theorem}[Repaired-gauge representation gives terminal orbit sampling]
\label{paper6a:thm:representation-gives-terminal-orbit-sampling}
Assume a repaired-gauge representation of the Type II branch and the terminal
coordinate-frame construction.  Then every terminal active sequence is sampled
from the renormalized orbit.
\end{theorem}

\begin{proof}
The repaired-gauge representation supplies the renormalized orbit \(V(\tau)\)
and the critical coordinate maps between physical variables and renormalized
variables.  The terminal coordinate-frame construction supplies the terminal
sampling times and scale-center parameters.  Combining these hypotheses gives
\[
  u_{0,n}=\Lambda_{\lambda_n,x_n}V(\tau_n),
\]
which is the assertion of
\cref{paper6a:def:terminal-sequences-sampled-from-orbit}.
\end{proof}

\begin{corollary}[Failure of terminal orbit sampling]
\label{paper6a:cor:failure-terminal-orbit-sampling}
If a sequence reaches the terminal profile stage but is not sampled from the
renormalized orbit, then either the repaired-gauge representation is not
available for that sequence or the terminal coordinate-frame construction has
not been applied as specified.
\end{corollary}

\begin{proof}
This is the contrapositive of
\cref{paper6a:thm:representation-gives-terminal-orbit-sampling}.
\end{proof}

\begin{corollary}[Terminal compactness from explicit PDE hypotheses]
\label{paper6a:cor:terminal-compactness-from-explicit-hypotheses}
Assume the repaired-gauge representation, the terminal coordinate-frame
construction, a positive finite terminal \(L^3\) annulus for \(V\), the Kato
small-data stability theorem, and the critical Navier--Stokes profile
decomposition.  Then the terminal profile completeness hypothesis in the final
argument holds.
\end{corollary}

\begin{proof}
\Cref{paper6a:thm:representation-gives-terminal-orbit-sampling} gives terminal
orbit sampling.  Together with the positive finite \(L^3\) annulus,
\cref{paper6a:thm:bounded-terminal-sequences-from-l3} gives bounded critical
terminal sequences.  The Kato theorem gives small-data stability by
\cref{paper6a:cor:kato-gives-small-data-stability}.  Therefore the terminal
critical profile hypotheses of
\cref{paper6a:def:terminal-critical-profile-hypotheses} hold, and
\cref{paper6a:cor:profile-completeness-from-critical-hypotheses} gives terminal
profile completeness.
\end{proof}

\subsection{Terminal Conditional Assembly}

\begin{definition}[Final conditional hypotheses]\label{paper6a:def:abstract-final-data}
The final conditional hypotheses for this section are:
\begin{enumerate}[label=\textup{(\roman*)}]
\item every renormalized local Type II branch enters one of the local
alternatives used in \cref{paper6:thm:state-decomposition};
\item compact single-core branches are excluded by
\cref{paper6:thm:single-core};
\item multibubble, cascade, gauge-degenerate, and scale-rigid terminal branches
are excluded under the hypotheses of \cref{paper6:thm:multibubble};
\item finite-cost scale-collapse branches are excluded by
\cref{paper6:thm:scale-cost}, while scale-collapse branches with divergent
negative scale-drift cost are covered by
\cref{paper6a:thm:scale-collapse-drift-cost} or
\cref{paper6a:cor:absolute-scale-cost-criterion};
\item failure of uniform critical tightness is excluded, equivalently uniform critical
tightness holds for every renormalized branch in the class;
\item failure of local windowed \(H^1\) control is excluded, equivalently the required local
windowed \(H^1\) control holds for every renormalized branch in the class;
\item the cost-divergence exclusion criterion
\cref{paper6a:ass:cost-divergence-exclusion} holds for the localized costs
used in the preceding items.
\end{enumerate}
\end{definition}

\begin{theorem}[Final conditional Type II exclusion]
\label{paper6a:thm:abstract-final-typeII}
Under the final conditional hypotheses of
\cref{paper6a:def:abstract-final-data}, no renormalized local Type II branch in
the stated class remains.
\end{theorem}

\begin{proof}
Let a renormalized local Type II branch be given.  By item \textup{(i)}, it
lies in one of the alternatives in \cref{paper6:thm:state-decomposition}.
Items \textup{(ii)} and \textup{(iii)} exclude the compact single-core,
multibubble, cascade, gauge-degenerate, and scale-rigid alternatives.  Item
\textup{(iv)} excludes finite-cost scale collapse and also excludes the
divergent negative-drift scale-collapse branch whenever the corresponding cost
is one of the costs covered by
\cref{paper6a:ass:cost-divergence-exclusion}.  Items \textup{(v)} and
\textup{(vi)} exclude failure of uniform critical tightness and failure of
local windowed \(H^1\) control.  These cases exhaust the alternatives, so no
branch remains.
\end{proof}

\begin{definition}[Expanded terminal hypotheses]\label{paper6a:def:expanded-terminal-data}
The following analytic hypotheses imply the final conditional hypotheses:
\begin{enumerate}[label=\textup{(\roman*)}]
\item the local alternative decomposition for renormalized Type II branches;
\item the compact single-core, multibubble/cascade, finite-cost
scale-collapse, and scale-rigidity exclusions from the preceding sections;
\item a critical-tightness estimate giving uniform critical tightness;
\item an exhaustive terminal state-space partition and a local finite
critical-mass statement for retained active terminal strata;
\item small-data stability in \(L^3\) and a critical Navier--Stokes profile
decomposition with nonlinear stability;
\item removal of the scattering branch and of physically exterior regular
profiles;
\item the repaired-gauge representation theorem;
\item the local Caccioppoli/windowed regularity input;
\item the scale-collapse rigidity input used by the multibubble and
scale-collapse reductions;
\item the cost-divergence exclusion criterion
\cref{paper6a:ass:cost-divergence-exclusion}.
\end{enumerate}
\end{definition}

\begin{lemma}[Terminal hypotheses imply critical tightness]
\label{paper6a:lem:terminal-critical-tightness}
The expanded terminal hypotheses imply uniform critical tightness.
\end{lemma}

\begin{proof}
The repaired-gauge representation turns each renormalized orbit into the
terminal sequence used in the profile analysis.  The terminal state-space
partition, finite critical-mass statement, small-data \(L^3\) stability, and
critical profile decomposition give the required terminal profile completeness.
The scattering branch and exterior regular profiles are removed by hypothesis.
Together with the multibubble and scale-collapse exclusions, this leaves only
possible loss of critical tightness; the critical-tightness estimate excludes
that loss.
\end{proof}

\begin{lemma}[Terminal hypotheses imply windowed \(H^1\) control]
\label{paper6a:lem:terminal-windowed-h1}
The expanded terminal hypotheses imply local windowed \(H^1\) control.
\end{lemma}

\begin{proof}
The windowed \(H^1\) estimate requires a repaired-gauge branch, local pressure and
modulation control, compact-cylinder critical bounds, and the
Caccioppoli/windowed regularity input.  These are among the expanded terminal
hypotheses.  Hence the required local windowed \(H^1\) control holds, and the
failure of local windowed \(H^1\) control is unavailable.
\end{proof}

\begin{lemma}[Expanded terminal hypotheses imply final hypotheses]
\label{paper6a:lem:expanded-to-abstract}
The expanded terminal hypotheses imply the final conditional hypotheses of
\cref{paper6a:def:abstract-final-data}.
\end{lemma}

\begin{proof}
Items \textup{(i)}, \textup{(ii)}, and \textup{(x)} of
\cref{paper6a:def:expanded-terminal-data} give the local alternative
decomposition, the compact, multibubble, scale-collapse, and scale-rigidity
exclusions, and the cost-divergence criterion.  By
\cref{paper6a:lem:terminal-critical-tightness} they also give critical
tightness, and by
\cref{paper6a:lem:terminal-windowed-h1} they give the local windowed \(H^1\)
control.  These are the hypotheses listed in
\cref{paper6a:def:abstract-final-data}.
\end{proof}

\begin{theorem}[Expanded terminal Type II exclusion]
\label{paper6a:thm:expanded-terminal-typeII}
If the expanded terminal hypotheses hold, then no renormalized local Type II
branch in the stated class remains.
\end{theorem}

\begin{proof}
\Cref{paper6a:lem:expanded-to-abstract} gives the final conditional
hypotheses.  Applying \cref{paper6a:thm:abstract-final-typeII} gives the
exclusion.
\end{proof}

\begin{remark}[Branchwise reading]
Compact single-core branches are removed by the compact criterion;
scale-collapse branches are removed by the scale-collapse rigidity input, or by
\cref{paper6a:thm:scale-collapse-drift-cost} when the scale-collapse drift
cost is compatible with the cost-divergence criterion; multibubble and
non-tight cases are removed by the terminal profile and critical-tightness
inputs; failures of local windowed \(H^1\) control are removed by the local
Caccioppoli estimate; and
every divergent compatible localized cost is handled by
\cref{paper6a:prop:cost-divergence-exclusion}.
\end{remark}

\clearpage
\section{Assembly of the Local Type II Exclusion}\label{sec:typeII-assembly}
\subsection{Local Entry}

Let \((u,p)\) be a suitable weak solution of the three-dimensional
Navier--Stokes equations
\[
  \partial_tu+(u\cdot\nabla)u+\nabla p=\nu\Delta u,
  \qquad \nabla\cdot u=0,
\]
on \(\mathbb R^3\times(T-\delta,T)\).  For \(z_0=(x_0,T)\) and \(r>0\), set
\[
  Q_r(z_0)=B_r(x_0)\times(T-r^2,T),
\]
\[
  C(z_0,r)=r^{-2}\int_{Q_r(z_0)}|u|^3\,dx\,dt,
\]
and
\[
  D(z_0,r)=r^{-2}\int_{T-r^2}^{T}\int_{B_r(x_0)}
  |p(x,t)-(p)_{B_r(x_0)}(t)|^{3/2}\,dx\,dt.
\]
The pressure is always taken modulo functions of time, and the spatial mean
normalization is the one used in the Caffarelli--Kohn--Nirenberg criterion.

\begin{definition}[Singular set and local Type II point]
Let \(\Sigma(T)\) be the set of points \(x_0\) for which \(u\) is not locally
bounded in any backward cylinder \(Q_r(x_0,T)\).  A point \(x_0\in\Sigma(T)\)
is in the local Type II alternative if it carries a positive local
Caffarelli--Kohn--Nirenberg concentration sequence and no local Type I
scale-invariant bound holds at \((x_0,T)\).
\end{definition}

\begin{definition}[Positive local Type II concentration sequence]
\label{paper6:def:local-typeII-sequence}
A sequence \(r_n\downarrow0\) is a positive local Type II concentration sequence
at \((x_0,T)\) if there exists \(\eta>0\) such that
\[
  C((x_0,T),r_n)+D((x_0,T),r_n)\ge \eta
  \qquad\text{for all }n,
\]
and \((x_0,T)\) is in the local Type II alternative.  The associated parabolic
rescalings are
\[
  u_n(y,s)=r_nu(x_0+r_ny,T+r_n^2s),
  \qquad
  p_n(y,s)=r_n^2p(x_0+r_ny,T+r_n^2s).
\]
\end{definition}

\begin{theorem}[Local concentration and Type I/Type II dichotomy]
\label{paper6:thm:paper0}
If \(\Sigma(T)\ne\emptyset\), then there are
\(x_0\in\Sigma(T)\), \(\eta>0\), and \(r_n\downarrow0\) such that
\[
  C((x_0,T),r_n)+D((x_0,T),r_n)\ge\eta
  \qquad\text{for all }n.
\]
At the concentration point exactly one of the following alternatives holds:
there is a local Type I scale-invariant bound, or \((x_0,T)\) is in the local
Type II alternative.
\end{theorem}

\begin{proof}
This is the local concentration dichotomy used in the final assembly.
\end{proof}

\subsection{Local Theorems Used in the Assembly}

The local results needed for the final assembly are restated in the following form.

\begin{theorem}[Repaired-gauge representation]
\label{paper6:thm:representation}
Let a positive local Type II concentration sequence contain a nondegenerate
selected core.  Then, after passing to selected windows, there are local
scale-center variables \(\lambda(\tau)>0\), \(x_c(\tau)\), renormalized fields
\[
  V(y,\tau)=\lambda(\tau)u(x_c(\tau)+\lambda(\tau)y,t(\tau)),
  \qquad
  P(y,\tau)=\lambda(\tau)^2p(x_c(\tau)+\lambda(\tau)y,t(\tau))+\pi(\tau),
\]
and modulation coefficients \(a,b\) such that
\[
  \partial_\tau V+(V\cdot\nabla)V+\nabla P
  =
  \nu\Delta V+a(\tau)(V+y\cdot\nabla V)+b(\tau)\cdot\nabla V,
  \qquad \nabla\cdot V=0.
\]
On each compact renormalized window one has the local suitable energy
inequality, local pressure reconstruction modulo functions of time, stability
under subsequences, and the compact-window bounds proved in \cref{sec:repaired-gauge-construction}.  Failure of the nondegenerate scale-center selection is part of the
multibubble, cascade, or scale-rigidity alternatives.
\end{theorem}

\begin{proof}
\Cref{paper2:thm:C} is precisely the compact-window representation theorem.
\end{proof}

\begin{theorem}[Compact single-core exclusion]
\label{paper6:thm:single-core}
No compact single-core vanishing-dissipation branch can be a local Type II
sequence.  More explicitly, let a represented selected-window sequence have:
\begin{enumerate}[label=\textup{(\roman*)}]
\item positive local CKN concentration on a fixed compact cylinder;
\item uniformly bounded compact-cylinder suitable estimates \(A+E+C+D\);
\item nondegenerate local scale-center representation;
\item bounded modulation coefficients;
\item local pressure reconstruction and pressure stability modulo spatial
means;
\item vanishing localized dissipation on some retained compact cylinder.
\end{enumerate}
Then the sequence is incompatible with the local Type II condition.
\end{theorem}

\begin{proof}
\Cref{paper1:thm:single-core-criterion} gives the stated exclusion.
\end{proof}

\begin{theorem}[Multibubble and cascade exclusion]
\label{paper6:thm:multibubble}
Assume the terminal admissibility and scale-rigidity hypotheses stated in the
multibubble/cascade analysis.
Then no local multibubble, cascade, or gauge-degenerate Type II branch can occur
in the class produced by the local active-core extraction.  Equivalently, every
failure of the single nondegenerate compact-core alternative is excluded by the
local multibubble/cascade analysis under those stated hypotheses.
\end{theorem}

\begin{proof}
\Cref{paper3:thm:multi}, together with its removal corollary
\cref{paper3:cor:multi-removal}, gives the stated exclusion.
\end{proof}

\begin{theorem}[Rough-core reduction]
\label{paper6:thm:rough-core}
On any compact renormalized window satisfying the local energy inequality and
bounded modulation hypotheses, finite outer velocity-pressure CKN control
\(\mathcal C+\mathcal D\) implies finite inner windowed \(H^1\) control.  Hence
failure of local windowed \(H^1\) control on a retained inner cylinder forces
failure of the compact-cylinder \(\mathcal C+\mathcal D\) bounds on every
relevant enclosing cylinder.
\end{theorem}

\begin{proof}
\Cref{paper4:cor:rough-core} gives the stated reduction.
\end{proof}

\begin{theorem}[Finite-cost scale-collapse exclusion]
\label{paper6:thm:scale-cost}
No renormalized branch with finite scale-collapsing cost, and hence no branch
with finite absolute scale cost, can have genuine scale collapse together with a
nonvanishing selected core.
\end{theorem}

\begin{proof}
\Cref{paper5:thm:cost-exclusion,paper5:cor:absolute-cost-exclusion} give the stated exclusion.
\end{proof}

\begin{theorem}[Scale-collapse decomposition of alternatives]
\label{paper6:thm:scale-stratification}
Let a represented Type II branch have genuine scale collapse and a nonvanishing
selected core.  After passing to selected terminal windows, exactly one of the
alternatives isolated in \cref{sec:scale-collapse-cost} occurs:
\begin{enumerate}[label=\textup{(\roman*)}]
\item finite scale-collapsing cost;
\item finite absolute scale cost;
\item thin-drift alternative;
\item autonomous-modulation alternative;
\item compactness alternative;
\item a nonzero autonomous reduced limit in the scale-rigid state.
\end{enumerate}
In the local Type II class of alternatives used by the multibubble/cascade and
scale-collapse analyses, alternatives
\textup{(iii)}--\textup{(vi)} do not constitute additional compact single-core
mechanisms: the thin-drift, modulation, and compactness alternatives lie outside
the compact finite-cost scale-collapse branch, and the scale-rigid state is
covered by the scale-rigidity hypothesis in the multibubble/cascade theorem.
\end{theorem}

\begin{proof}
The assertion follows by combining the scale-collapse decomposition of
alternatives with the scale-rigidity alternative used in the
multibubble/cascade reduction.
\end{proof}

\begin{definition}[Virial-dissipative terminal state]
\label{paper6:def:virial-dissipative}
Let \(G_\nu(y)=e^{-|y|^2/(4\nu)}\).  A scale-rigid terminal state is called
velocity-virial dissipative if its selected terminal representative
\((V,P)\) satisfies, after the usual pressure normalization, the autonomous
self-similar equation
\[
  \partial_\tau V+(V\cdot\nabla)V+\nabla P
  =
  \nu\Delta V-\frac12(V+y\cdot\nabla V),
  \qquad \nabla\cdot V=0,
\]
and the Gaussian identity below is justified by cutoff approximation on the
selected renormalized windows.  It is called swirl-virial dissipative if, in
addition, the representative is axisymmetric, \(\Gamma=R V_\theta\) is the
scale-invariant circulation, and the corresponding weighted circulation
identity below is justified.
\end{definition}

\begin{theorem}[Virial exclusion for dissipative terminal states]
\label{paper6:thm:virial-exclusion}
Let a scale-rigid terminal state satisfy the hypotheses of
\cref{paper6:def:virial-dissipative} and be bounded for all negative renormalized
times in the sense needed for the Gaussian integrals below.

\begin{enumerate}[label=\textup{(\alph*)}]
\item Suppose the velocity virial defect is absorbable: for some
\(\kappa>0\) and a.e. \(\tau\),
\[
  -\frac{1}{2\nu}\int_{\mathbb R^3}|V|^2(V\cdot y)G_\nu\,dy
  -\frac1\nu\int_{\mathbb R^3}P(V\cdot y)G_\nu\,dy
  \le
  (1-\kappa)\int_{\mathbb R^3}|V|^2G_\nu\,dy
  +2\nu(1-\kappa)\int_{\mathbb R^3}|\nabla V|^2G_\nu\,dy .
\]
Then
\[
  \int_{\mathbb R^3}|V(y,\tau)|^2G_\nu(y)\,dy=0
  \qquad\text{for every }\tau.
\]
Consequently no retained terminal core with positive Gaussian velocity mass can
belong to this branch.

\item Suppose the terminal state is axisymmetric and the circulation defect is
absorbable: for some \(\kappa>0\) and a.e. \(\tau\),
\[
  -\frac{1}{2\nu}\int_{\mathbb R^3}\Gamma^2(V\cdot y)G_\nu\,dy
  \le
  (1-\kappa)\int_{\mathbb R^3}\Gamma^2G_\nu\,dy
  +2\nu(1-\kappa)\int_{\mathbb R^3}|\nabla \Gamma|^2G_\nu\,dy .
\]
Then
\[
  \int_{\mathbb R^3}\Gamma(y,\tau)^2G_\nu(y)\,dy=0
  \qquad\text{for every }\tau.
\]
Thus every retained axisymmetric terminal core whose remaining obstruction is a
positive Gaussian circulation mass is excluded by the weighted circulation
identity.
\end{enumerate}
\end{theorem}

\begin{proof}
For the velocity statement set
\[
  \mathcal E(\tau)=\int_{\mathbb R^3}|V(y,\tau)|^2G_\nu(y)\,dy .
\]
The Gaussian virial identity gives
\[
  \partial_\tau\mathcal E
  =
  -2\nu\int_{\mathbb R^3}|\nabla V|^2G_\nu\,dy
  -\int_{\mathbb R^3}|V|^2G_\nu\,dy
  -\frac{1}{2\nu}\int_{\mathbb R^3}|V|^2(V\cdot y)G_\nu\,dy
  -\frac1\nu\int_{\mathbb R^3}P(V\cdot y)G_\nu\,dy .
\]
The absorbability hypothesis therefore implies
\[
  \partial_\tau\mathcal E
  +\kappa\mathcal E
  +2\nu\kappa\int_{\mathbb R^3}|\nabla V|^2G_\nu\,dy
  \le0 .
\]
In particular \(\partial_\tau\mathcal E+\kappa\mathcal E\le0\).  Integrating
from \(s<\tau\) to \(\tau\) gives
\[
  \mathcal E(\tau)\le e^{-\kappa(\tau-s)}\mathcal E(s).
\]
The ancient boundedness of the terminal state makes \(\sup_{s<\tau}\mathcal
E(s)<\infty\).  Sending \(s\to-\infty\) yields \(\mathcal E(\tau)=0\).

For the axisymmetric statement set
\[
  \mathcal G(\tau)=\int_{\mathbb R^3}\Gamma(y,\tau)^2G_\nu(y)\,dy .
\]
The weighted circulation identity is
\[
  \partial_\tau\mathcal G
  +\mathcal G
  +2\nu\int_{\mathbb R^3}|\nabla\Gamma|^2G_\nu\,dy
  =
  -\frac{1}{2\nu}\int_{\mathbb R^3}\Gamma^2(V\cdot y)G_\nu\,dy .
\]
The assumed absorption gives
\[
  \partial_\tau\mathcal G+\kappa\mathcal G
  +2\nu\kappa\int_{\mathbb R^3}|\nabla\Gamma|^2G_\nu\,dy
  \le0 .
\]
The same backward Gronwall argument gives \(\mathcal G(\tau)=0\) for every
\(\tau\).
\end{proof}

\begin{remark}[Controlled swirl alone]
\label{paper6:rem:controlled-swirl}
If the axisymmetric representative satisfies \(\Gamma\in L^\infty\), the same
weighted circulation identity gives the a priori estimate
\[
  \mathcal G(\tau)^{1/2}
  \lesssim
  \|V\|_{L^\infty}\|\Gamma\|_{L^\infty}
\]
for bounded ancient terminal states.  This estimate enters the local
decomposition of alternatives but does not by itself exclude a terminal state.  The exclusion
requires the stated absorption, vanishing conclusion, or an equivalent
condition which makes the virial inequality dissipative.
\end{remark}

\subsection{Assembly of Alternatives}

\begin{theorem}[Local Type II decomposition into local alternatives]
\label{paper6:thm:state-decomposition}
Every positive local Type II concentration sequence in \cref{paper6:def:local-typeII-sequence}, after passing to subsequences and selected windows, enters one of the following alternatives:
\begin{enumerate}[label=\textup{(\roman*)}]
\item a compact single-core vanishing-dissipation branch;
\item a multibubble, cascade, or gauge-degenerate branch;
\item loss of local windowed \(H^1\) control on a retained compact cylinder;
\item a finite-cost scale-collapse branch with nonvanishing selected core;
\item a scale-rigid state covered by the scale-rigidity hypothesis in
\cref{paper6:thm:multibubble}.
\end{enumerate}
In the last alternative, any terminal subbranch satisfying
\cref{paper6:thm:virial-exclusion} is eliminated before the remaining scale-rigid
state is passed to the scale-rigidity hypothesis.
\end{theorem}

\begin{proof}
This decomposition is obtained by combining
\cref{paper6:thm:representation,paper6:thm:multibubble,paper6:thm:rough-core,paper6:thm:scale-stratification}.
The representation theorem gives the nondegenerate single-core chart whenever a
unique retained core is selected.  Failure of that selection, splitting of
positive CKN mass, compactness failure, or scale cascade is the
multibubble/cascade side.  On a represented retained core, rough-core loss is
equivalent by \cref{paper6:thm:rough-core} to failure of compact CKN control on an
enclosing cylinder, hence returns to the multibubble/cascade side.  Genuine
scale collapse is divided by \cref{paper6:thm:scale-stratification} into finite-cost
scale collapse or one of the scale-rigid terminal states used in
\cref{paper6:thm:multibubble}.  These cases exhaust the local alternatives.
\end{proof}

\begin{proposition}[Elimination of local Type II states]
\label{paper6:prop:state-elimination}
No positive local Type II concentration sequence can lie in any of the
alternatives \textup{(i)}--\textup{(v)} of
\cref{paper6:thm:state-decomposition}.
\end{proposition}

\begin{proof}
Alternative \textup{(i)} is excluded by the compact single-core criterion,
\cref{paper6:thm:single-core}.

Alternative \textup{(ii)} is excluded by the local multibubble and cascade
theorem, \cref{paper6:thm:multibubble}, under its terminal admissibility and
scale-rigidity hypotheses.

Alternative \textup{(iii)} is not an independent terminal state.  By
\cref{paper6:thm:rough-core}, rough-core loss of local \(H^1\) control forces
failure of the compact-cylinder \(\mathcal C+\mathcal D\) bounds on enclosing
compact cylinders.  Such failure yields an active compact-cylinder
concentration branch and therefore belongs to the multibubble/cascade
alternative, already excluded by \cref{paper6:thm:multibubble}.

Alternative \textup{(iv)} is excluded by the finite-cost scale-collapse theorem,
\cref{paper6:thm:scale-cost}.

Alternative \textup{(v)} is the scale-rigid terminal state in
\cref{paper6:thm:multibubble}.  If this terminal state is velocity-virial dissipative,
\cref{paper6:thm:virial-exclusion} forces its Gaussian velocity mass to vanish, which
contradicts retention of a positive terminal core.  If it is axisymmetric and
swirl-virial dissipative with positive retained circulation mass, the same
theorem forces \(\Gamma\equiv0\), contradicting that positive circulation mass.
The remaining scale-rigid terminal states are covered by the
scale-rigidity hypothesis in \cref{paper6:thm:multibubble}.  Under that hypothesis, the
state either has a nonzero bounded selected-window limit, which is incompatible
with the local Type II condition, or it reduces to the compact single-core
branch, which is excluded by \cref{paper6:thm:single-core}.  Hence it cannot be a local
Type II branch.
\end{proof}

\subsection{Assembly Theorem}

\begin{theorem}[Local Type II exclusion]
\label{paper6:thm:local-typeII-exclusion}
Let \((u,p)\) be a suitable weak solution on
\(\mathbb R^3\times(T-\delta,T)\).  Then no point of \(\Sigma(T)\) lies in the
local Type II alternative, within the class of local alternatives covered by
\cref{paper6:thm:state-decomposition}.
\end{theorem}

\begin{proof}
Assume, for contradiction, that \(x_0\in\Sigma(T)\) lies in the local Type II
alternative.  By \cref{paper6:thm:paper0}, there is a positive local Type II
concentration sequence \(r_n\downarrow0\) at \((x_0,T)\).  By the decomposition into local alternatives, \cref{paper6:thm:state-decomposition}, this sequence lies,
after passing to selected windows, in one of the five alternatives listed
there.  Each alternative is impossible by \cref{paper6:prop:state-elimination}.  This
contradiction proves that no such \(x_0\) exists in the covered local
class of alternatives.
\end{proof}

\begin{corollary}[Final local dichotomy after Type I separation]
\label{paper6:cor:final-local-dichotomy}
Every singular point at time \(T\) in the covered class of local alternatives
belongs to the local Type I alternative.  Equivalently, once the Type I case is
covered by the Type I part of the theory, the local Type II class of alternatives contains no remaining singular branch.
\end{corollary}

\begin{proof}
Let \(x_0\in\Sigma(T)\).  By \cref{paper6:thm:paper0}, a concentration point is
either in the local Type I alternative or in the local Type II alternative.
\Cref{paper6:thm:local-typeII-exclusion} excludes the latter.
\end{proof}

\clearpage
\section{PDE Form of the Type II Classification Inputs}
\label{sec:c1-c5-pde-translation}

This section records, in PDE notation, the Type II classification inputs used
in the local Type II analysis.  The statements are conditional classification
statements for admissible Type II branches.  They do not assert the existence of
such branches and do not add any regularity theorem beyond the analytic
hypotheses explicitly stated below.

\subsection{Admissible Type II branches and exhaustion}

\begin{definition}[Admissible local Type II class]
\label{def:c1-admissible-typeII-class}
Let \(\mathcal U_{\mathrm{II}}^{NS}\) be the class of admissible terminal
branches \((u,p,T^*)\) such that:
\begin{enumerate}[label=\textup{(\roman*)}]
\item \((u,p)\) is a three-dimensional incompressible Navier--Stokes solution
on a terminal interval \([0,T^*)\) in the analytic class under consideration;
\item \(T^*<\infty\) is the terminal time of the branch;
\item the branch is not in the local Type I class associated with the
scale-invariant Type I criterion;
\item the branch is a Type II branch rather than a continuation, Type I,
boundary, or non-Navier--Stokes-domain alternative.
\end{enumerate}
Membership in \(\mathcal U_{\mathrm{II}}^{NS}\) specifies the class under
analysis; no nonemptiness assertion is made.
\end{definition}

\begin{definition}[Type II analytic data]
\label{def:c1-typeII-analytic-data}
A branch \(\omega=(u,p,T^*)\in\mathcal U_{\mathrm{II}}^{NS}\) has the Type II
analytic data if the following three analytic assertions hold:
\begin{enumerate}[label=\textup{(\roman*)}]
\item a concentration profile or blow-up germ modulo the Navier--Stokes
symmetry group;
\item a supercritical scale alternative, meaning that the branch enters the Type II
scale alternative rather than the Type I scale alternative;
\item local profile data in the Navier--Stokes profile class used
for the subsequent compactness and representation arguments.
\end{enumerate}
The analytic data are exhaustive if every admissible Type II branch either has
these three properties or falls into one of the ordered failures in
\cref{def:c1-exhaustion-failures}.
\end{definition}

\begin{definition}[Ordered exhaustion failures]
\label{def:c1-exhaustion-failures}
For an admissible Type II branch, the ordered failures of the Type II analytic
data are:
\begin{enumerate}[label=\textup{(\roman*)}]
\item concentration extraction fails;
\item concentration extraction succeeds but the supercritical scale alternative
does not hold;
\item concentration extraction and the scale alternative hold but the local
profile data are insufficient for the subsequent compactness and representation
arguments;
\item the admissible Type II branch lies outside the concentration, scale, and
profile alternatives under consideration.
\end{enumerate}
The first applicable failure in this order is the exhaustion alternative.
\end{definition}

\begin{theorem}[Type II branch exhaustion]
\label{thm:c1-typeII-branch-exhaustion}
Assume the Type II analytic data of \cref{def:c1-typeII-analytic-data} are exhaustive
on \(\mathcal U_{\mathrm{II}}^{NS}\).  Then every
\(\omega\in\mathcal U_{\mathrm{II}}^{NS}\) has the concentration profile, the
supercritical scale alternative, and the local profile data.
\end{theorem}

\begin{proof}
Fix \(\omega=(u,p,T^*)\in\mathcal U_{\mathrm{II}}^{NS}\).  By the
concentration-extraction component of the analytic data, the branch has a
concentration profile or blow-up germ.  By the scale component, the same branch
lies in the supercritical Type II scale alternative.  By the profile component,
the extracted profile belongs to the local Navier--Stokes profile class used in
the subsequent arguments.  Since \(\omega\) was arbitrary, the three assertions
hold for every admissible Type II branch.
\end{proof}

\begin{corollary}[Ordered exhaustion classification]
\label{cor:c1-ordered-exhaustion}
For each \(\omega\in\mathcal U_{\mathrm{II}}^{NS}\), exactly one ordered
alternative occurs: the Type II analytic data hold, or one of the four failures
in \cref{def:c1-exhaustion-failures} occurs.  The Type II analytic data are
exhaustive on \(\mathcal U_{\mathrm{II}}^{NS}\) if and only if the first
alternative holds for every \(\omega\in\mathcal U_{\mathrm{II}}^{NS}\).
\end{corollary}

\begin{proof}
Test the four analytic requirements in the stated order.  If concentration
extraction, the scale alternative, and the profile condition all hold, the
branch has the Type II analytic data.  If a test fails, the first failed test gives the
corresponding ordered failure.  The alternatives are disjoint by the first
failure convention and exhaustive by construction.  The final equivalence is
the definition of exhaustion on \(\mathcal U_{\mathrm{II}}^{NS}\).
\end{proof}

\begin{theorem}[Entry into the Type II alternatives]
\label{thm:c1-downstream-entry}
Assume the Type II analytic data are exhaustive.  Then each admissible Type II
branch enters the subsequent analysis in exactly one of the following ways:
\begin{enumerate}[label=\textup{(\roman*)}]
\item it reaches the compact Type II cost-divergence hypotheses after the
representation, critical-mass, tightness, and windowed-\(H^1\) tests;
\item it falls into an ordered representation failure, critical-mass failure,
loss of critical tightness, or loss of local windowed \(H^1\) control.
\end{enumerate}
\end{theorem}

\begin{proof}
By \cref{thm:c1-typeII-branch-exhaustion}, every admissible Type II branch has
the concentration, scale, and profile data.  The representation step then
either gives a repaired-gauge representation or one of the ordered representation
failures in \cref{def:c2-representation-failures}.  If a repaired-gauge orbit
exists, the critical \(L^3\) test gives either the positive finite annulus or
one of the ordered critical-mass alternatives in
\cref{def:c3-critical-mass-alternatives}.  If the positive finite annulus
holds, the remaining compact cost-divergence inputs are critical tightness and local
windowed \(H^1\) control.  Their failures are exactly the two local failures in
\cref{def:c5-compact-tests}.  If both hold, the
compact Type II cost-divergence criterion applies.
\end{proof}

\subsection{Repaired-gauge representation from Type II analytic data}

\begin{definition}[Raw Type II chart]
\label{def:c2-raw-chart}
A raw Type II chart consists of
\[
  (u,p,T^*,x_c,\lambda,\tau,V,P,\mathcal I),
  \qquad \mathcal I=(t_0,T^*),
\]
where \(u,p\) solve
\[
  \partial_tu+(u\cdot\nabla)u+\nabla p=\nu\Delta u,\qquad
  \nabla\cdot u=0,
\]
\(\lambda(t)>0\), \(x_c(t)\in\mathbb R^3\), \(x_c,\lambda\in AC_{\rm loc}\),
\(\lambda(t)\to0\) as \(t\uparrow T^*\), and
\[
  \frac{d\tau}{dt}=\lambda(t)^{-2},\qquad
  \tau(t)\to\infty\quad\text{as }t\uparrow T^*.
\]
The renormalized variables are
\[
  y=\frac{x-x_c(t)}{\lambda(t)},\qquad
  u(x,t)=\lambda(t)^{-1}V(y,\tau(t)),\qquad
  p(x,t)=\lambda(t)^{-2}P(y,\tau(t)),
\]
with enough regularity to justify the distributional chain rule on compact
\(y\)-sets.
\end{definition}

\begin{lemma}[Raw renormalized equation]
\label{lem:c2-raw-renormalized-equation}
For a raw Type II chart, define
\[
  a_{\rm raw}(\tau(t))=\lambda(t)\lambda_t(t),\qquad
  b_{\rm raw}(\tau(t))=\lambda(t)x_c'(t).
\]
Then \(V,P\) satisfy, distributionally on compact \(y\)-sets,
\[
  \partial_\tau V+(V\cdot\nabla)V+\nabla P
  =
  \nu\Delta V
  +a_{\rm raw}(\tau)(V+y\cdot\nabla V)
  +b_{\rm raw}(\tau)\cdot\nabla V,
  \qquad \nabla\cdot V=0 .
\]
\end{lemma}

\begin{proof}
Write
\[
u(x,t)=\lambda(t)^{-1}V(y,\tau(t)),\qquad
y=\frac{x-x_c(t)}{\lambda(t)},\qquad \tau_t=\lambda^{-2}.
\]
Then
\[
  \nabla_xu=\lambda^{-2}\nabla_yV,\qquad
  \Delta_xu=\lambda^{-3}\Delta_yV,\qquad
  (u\cdot\nabla_x)u=\lambda^{-3}(V\cdot\nabla_y)V.
\]
Since
\[
  y_t=-\frac{x_c'(t)}{\lambda(t)}
      -\frac{\lambda_t(t)}{\lambda(t)}y,
\]
one has
\[
  \partial_tu
  =
  \lambda^{-3}\partial_\tau V
  -\lambda_t\lambda^{-2}(V+y\cdot\nabla V)
  -\lambda^{-2}x_c'(t)\cdot\nabla V .
\]
Moreover \(\nabla_xp=\lambda^{-3}\nabla_yP\).  Substitution in
Navier--Stokes and multiplication by \(\lambda^3\) give
\[
  \partial_\tau V-\lambda\lambda_t(V+y\cdot\nabla V)
  -\lambda x_c'(t)\cdot\nabla V
  +(V\cdot\nabla)V+\nabla P
  =\nu\Delta V .
\]
Moving the modulation terms to the right-hand side gives the equation with
\(a_{\rm raw}=\lambda\lambda_t\) and \(b_{\rm raw}=\lambda x_c'\).  Finally,
\(\nabla_x\cdot u=\lambda^{-2}\nabla_y\cdot V\), so incompressibility of \(u\)
implies \(\nabla_y\cdot V=0\).
\end{proof}

\begin{definition}[Repaired-gauge realization]
\label{def:c2-repaired-gauge-realization}
Let \(0<p<3\), \(\Theta_0>0\), and let \(\psi_R\) be a smooth compactly
supported centering cutoff.  A raw Type II chart admits a repaired-gauge
realization if, after an admissible time-dependent scale and translation
correction, the profile satisfies
\[
  G_{\rm sc}(V)=\int_{\mathbb R^3}|y|^{-p}|V(y)|^3\,dy-\Theta_0=0,
\]
and
\[
  G_j(V)=\int_{\mathbb R^3}y_j|V(y)|^2\psi_R(y)\,dy=0,
  \qquad j=1,2,3,
\]
for a.e. renormalized time.  The corrected profile belongs to the admissible
gauge class \(\mathcal A_p\), has the regularity needed to differentiate these
functionals, and all time-dependent symmetry corrections have been absorbed
into final modulation coefficients \(a,b\).
\end{definition}

\begin{definition}[Repaired-gauge representation]
\label{def:c2-representation}
A Type II branch has a repaired-gauge representation if it has:
\begin{enumerate}[label=\textup{(\roman*)}]
\item a raw Type II chart;
\item a repaired-gauge realization;
\item pressure reconstruction modulo functions of time;
\item final modulation coefficients \(a,b\) appearing in the repaired-gauge
equation.
\end{enumerate}
The representation consists of the tuple
\[
  (V,P,a,b,G_{\rm sc},G_1,G_2,G_3,\tau_0).
\]
\end{definition}

\begin{theorem}[Representation implication]
\label{thm:c2-representation-implication}
Assume a Type II profile branch has a raw Type II chart, a repaired-gauge
realization, pressure reconstruction modulo functions of time, and final
modulation coefficients.  Then the branch has a repaired-gauge representation
in the sense of \cref{def:c2-representation}.  In particular,
\((V,P,a,b)\) satisfies the PDE class required by the compact Type II
cost-divergence criterion.
\end{theorem}

\begin{proof}
The raw chart gives \(V,P\) and, by \cref{lem:c2-raw-renormalized-equation},
the renormalized Navier--Stokes equation with raw modulation coefficients.  The
repaired-gauge realization places the profile on the scale and centering gauge
surface and absorbs the time-dependent symmetry correction into updated
coefficients.  The pressure reconstruction identifies the pressure associated
with \(V\), modulo the usual functions of time.  The modulation data identify
the final coefficients \(a,b\).  These are the components of
\cref{def:c2-representation}.
\end{proof}

\begin{definition}[Ordered representation failures]
\label{def:c2-representation-failures}
If a repaired-gauge representation is not obtained, the ordered failures are:
\begin{enumerate}[label=\textup{(\roman*)}]
\item no raw orbit is supplied by the concentration and profile data;
\item the raw orbit exists but the repaired gauge cannot be realized;
\item pressure reconstruction modulo functions of time is unavailable;
\item the final modulation coefficients are unavailable.
\end{enumerate}
The first applicable failure is the representation alternative.
\end{definition}

\begin{corollary}[Ordered representation classification]
\label{cor:c2-ordered-representation}
Every Type II profile branch in the admissible class has exactly one ordered
alternative: a repaired-gauge representation, or one of the four failures in
\cref{def:c2-representation-failures}.
\end{corollary}

\begin{proof}
Verify the four representation inputs in order.  If all are present,
\cref{thm:c2-representation-implication} gives the repaired-gauge
representation.  If not, the first missing input gives the corresponding
ordered failure.  This is exhaustive and disjoint by first-failure ordering.
\end{proof}

\subsection{Minimal chart and gauge realization}

\begin{definition}[Minimal representation data]
\label{def:c2r-minimal-representation-data}
The minimal data for a Type II branch satisfying the analytic data consist of:
\begin{enumerate}[label=\textup{(\roman*)}]
\item an absolutely continuous concentration chart
\[
  (u,p,T^*,x_c,\lambda,\rho,U,Q,\mathcal I),
\]
where \(\mathcal I=(t_0,T^*)\), \(\lambda(t)>0\), \(\lambda(t)\to0\),
\(\rho_t=\lambda^{-2}\), and
\[
  u(x,t)=\lambda(t)^{-1}U(y,\rho(t)),\qquad
  p(x,t)=\lambda(t)^{-2}Q(y,\rho(t))+c(t),\qquad
  y=\frac{x-x_c(t)}{\lambda(t)};
\]
\item an absolutely continuous repaired-gauge correction consisting of
\(\rho=\rho(\tau)\), \(\mu(\tau)>0\), and \(q(\tau)\in\mathbb R^3\), with
\[
  \rho_\tau=\mu(\tau)^2,
\]
\[
  V(Y,\tau)=\mu(\tau)U(\mu(\tau)Y+q(\tau),\rho(\tau)),\qquad
  P(Y,\tau)=\mu(\tau)^2Q(\mu(\tau)Y+q(\tau),\rho(\tau)),
\]
and
\[
  G_{\rm sc}(V(\tau))=0,\qquad
  G_j(V(\tau))=0,\quad j=1,2,3.
\]
The condition \(\rho_\tau=\mu^2\) preserves the viscosity coefficient after the
time-dependent scale correction.
\end{enumerate}
\end{definition}

\begin{lemma}[Chart data imply the raw orbit]
\label{lem:c2r-chart-raw-orbit}
Under the chart data in \cref{def:c2r-minimal-representation-data}, \(U,Q\)
satisfy
\[
  \partial_\rho U+(U\cdot\nabla)U+\nabla Q
  =
  \nu\Delta U+A(\rho)(U+y\cdot\nabla U)+B(\rho)\cdot\nabla U,
  \qquad \nabla\cdot U=0,
\]
where
\[
  A(\rho(t))=\lambda(t)\lambda_t(t),\qquad
  B(\rho(t))=\lambda(t)x_c'(t).
\]
\end{lemma}

\begin{proof}
The chart data are the raw orbit data with raw time \(\rho\).  The additive
pressure function \(c(t)\) has zero spatial gradient.  The chain-rule
computation is the one in \cref{lem:c2-raw-renormalized-equation}, with
\(\tau\) replaced by \(\rho\).  Thus spatial derivatives scale as
\(\nabla_xu=\lambda^{-2}\nabla_yU\), \(\Delta_xu=\lambda^{-3}\Delta_yU\), the
convection term scales as \(\lambda^{-3}(U\cdot\nabla)U\), and
\(\rho_t=\lambda^{-2}\).  The terms generated by \(\lambda_t\) and \(x_c'\)
give \(A=\lambda\lambda_t\) and \(B=\lambda x_c'\).  Incompressibility pulls
back as \(\nabla_x\cdot u=\lambda^{-2}\nabla_y\cdot U\).
\end{proof}

\begin{lemma}[Pressure pullback]
\label{lem:c2r-pressure-pullback}
Under the chart data,
\[
  -\Delta Q=\partial_i\partial_j(U_iU_j)
\]
in distributions, modulo functions of \(\rho\).  After any admissible
repaired-gauge correction, the transformed pressure satisfies
\[
  -\Delta P=\partial_i\partial_j(V_iV_j)
\]
in distributions, modulo functions of \(\tau\).
\end{lemma}

\begin{proof}
Taking divergence of the physical Navier--Stokes equation gives
\[
  -\Delta p=\partial_i\partial_j(u_iu_j),
\]
with pressure determined up to a function of time.  Pulling this identity back
through \(p=\lambda^{-2}Q+c(t)\), \(u=\lambda^{-1}U\), and
\(y=(x-x_c)/\lambda\) gives the pressure equation for \(Q\).  If
\[
  V(Y)=\mu U(\mu Y+q),\qquad P(Y)=\mu^2Q(\mu Y+q),
\]
then
\[
  -\Delta_YP
  =\mu^4(-\Delta Q)(\mu Y+q)
  =\mu^4\partial_i\partial_j(U_iU_j)(\mu Y+q)
  =\partial_i\partial_j(V_iV_j).
\]
Adding a function of \(\rho\) or \(\tau\) to the pressure does not affect the
identity.
\end{proof}

\begin{lemma}[Forced modulation coefficients]
\label{lem:c2r-forced-modulation}
Assume the minimal data of \cref{def:c2r-minimal-representation-data}.  Then
\((V,P)\) satisfies
\[
  \partial_\tau V+(V\cdot\nabla)V+\nabla P
  =
  \nu\Delta V
  +a(\tau)(V+Y\cdot\nabla V)
  +b(\tau)\cdot\nabla V,
\]
where
\[
  a(\tau)=\mu(\tau)^2A(\rho(\tau))+\frac{\mu_\tau(\tau)}{\mu(\tau)}
\]
and
\[
  b(\tau)=\mu(\tau)B(\rho(\tau))
      +\mu(\tau)A(\rho(\tau))q(\tau)
      +\frac{q_\tau(\tau)}{\mu(\tau)}.
\]
\end{lemma}

\begin{proof}
Let
\[
  z=\mu(\tau)Y+q(\tau),\qquad
  V(Y,\tau)=\mu(\tau)U(z,\rho(\tau)),\qquad
  \rho_\tau=\mu^2.
\]
The raw equation in \((z,\rho)\) is
\[
  U_\rho+(U\cdot\nabla_z)U+\nabla_zQ
  =
  \nu\Delta_zU+A(U+z\cdot\nabla_zU)+B\cdot\nabla_zU.
\]
The identities
\[
  \nabla_YV=\mu^2\nabla_zU,\qquad
  \Delta_YV=\mu^3\Delta_zU,\qquad
  (V\cdot\nabla_Y)V=\mu^3(U\cdot\nabla_z)U,\qquad
  \nabla_YP=\mu^3\nabla_zQ
\]
together with differentiation in \(\tau\) give
\[
  V_\tau+(V\cdot\nabla_Y)V+\nabla_YP-\nu\Delta_YV
  =
  \left(\frac{\mu_\tau}{\mu}+\mu^2A\right)(V+Y\cdot\nabla_YV)
  +
  \left(\mu B+\mu Aq+\frac{q_\tau}{\mu}\right)\cdot\nabla_YV.
\]
This is the displayed equation.  Since \(\mu,q,\rho\) are absolutely continuous
and \(A,B\) are the raw chart coefficients, \(a,b\) are measurable and locally
integrable on every finite renormalized window.
\end{proof}

\begin{theorem}[Minimal representation theorem]
\label{thm:c2r-minimal-representation}
If a Type II branch satisfying the analytic data has the minimal data of
\cref{def:c2r-minimal-representation-data}, then it has a repaired-gauge
representation in the sense of \cref{def:c2-representation}.
\end{theorem}

\begin{proof}
\Cref{lem:c2r-chart-raw-orbit} gives the raw orbit.  The gauge correction
places the profile on the repaired scale and centering surface, hence gives the
repaired-gauge realization.  \Cref{lem:c2r-pressure-pullback} gives pressure
reconstruction modulo functions of time.  \Cref{lem:c2r-forced-modulation}
gives the final modulation coefficients.  These are the four inputs in
\cref{def:c2-representation}.
\end{proof}

\begin{definition}[Minimal representation failures]
\label{def:c2r-minimal-failures}
Outside the admissible repaired-gauge representation class, the ordered failures
are:
\begin{enumerate}[label=\textup{(\roman*)}]
\item no absolutely continuous concentration chart is available;
\item the chart exists but the repaired scale and centering gauges cannot be
solved admissibly and absolutely continuously.
\end{enumerate}
Pressure-pullback and modulation-coefficient failures are not independent
alternatives when the chart and the absolutely continuous gauge solve exist.
\end{definition}

\begin{corollary}[Ordered minimal representation classification]
\label{cor:c2r-ordered-classification}
Every Type II branch satisfying the analytic data has exactly one ordered
alternative: a repaired-gauge representation, failure of the raw chart, or
failure of the repaired-gauge solve.
\end{corollary}

\begin{proof}
If the chart is absent, the first failure occurs.  If the chart exists but the
admissible absolutely continuous gauge solve is absent, the second failure
occurs.  If both are present, \cref{thm:c2r-minimal-representation} gives the
repaired-gauge representation.  The alternatives are disjoint by their order
and exhaustive for the representation row.
\end{proof}

\subsection{\texorpdfstring{Critical \(L^3\) annulus}{Critical L3 annulus}}

\begin{definition}[Critical mass and positive finite annulus]
\label{def:c3-critical-mass}
For a repaired-gauge renormalized Type II branch \((V,P,a,b)\) on
\([\tau_0,\infty)\), define
\[
  N_3(\tau)=\|V(\tau)\|_{L^3(\mathbb R^3)}.
\]
The orbit has positive finite critical mass if there are constants
\(0<\eta\le M<\infty\) such that
\[
  \eta\le N_3(\tau)\le M,\qquad \tau\ge\tau_0.
\]
This condition replaces the exact normalization
\(\|V(\tau)\|_{L^3}=1\); no amplitude rescaling is performed.
\end{definition}

\begin{definition}[Critical-mass alternatives]
\label{def:c3-critical-mass-alternatives}
For a repaired-gauge orbit, the ordered critical-mass alternatives are:
\begin{enumerate}[label=\textup{(\roman*)}]
\item \(N_3\) is not a well-defined extended measurable critical norm on the
renormalized branch;
\item \(N_3(\tau)=\infty\) for some \(\tau\), or
\(\sup_{\tau\ge\tau_0}N_3(\tau)=\infty\);
\item \(\sup_{\tau\ge\tau_0}N_3(\tau)<\infty\) and
\(\liminf_{\tau\to\infty}N_3(\tau)=0\);
\item after increasing the initial time at most once, the positive
finite critical annulus of \cref{def:c3-critical-mass} holds.
\end{enumerate}
The first applicable case is the corresponding alternative.
\end{definition}

\begin{theorem}[Nonzero-mass compact Type II cost divergence]
\label{thm:c3-nonzero-mass-cost-divergence}
Let \((V,P,a,b)\) be a repaired-gauge renormalized Type II branch.  Assume:
\begin{enumerate}[label=\textup{(\roman*)}]
\item the positive finite critical annulus holds:
\[
  0<\eta\le \|V(\tau)\|_{L^3(\mathbb R^3)}\le M<\infty;
\]
\item uniform critical tightness holds:
\[
  \forall\varepsilon>0\ \exists R_\varepsilon:\quad
  \sup_{\tau\ge\tau_0}
  \int_{|y|>R_\varepsilon}|V(y,\tau)|^3\,dy<\varepsilon;
\]
\item local windowed \(H^1\) control holds:
\[
  \forall m\ge1:\quad
  \sup_{n\in\mathbb N}
  \int_{\tau_0+n}^{\tau_0+n+1}
  \|V(\tau)\|_{H^1(B_m)}^2\,d\tau<\infty.
\]
\end{enumerate}
Then finite total localized renormalization cost is impossible:
\[
  \int_{\tau_0}^{\infty}
  \tilde{\mathfrak D}_{R_0}(\tau)\,d\tau=\infty.
\]
\end{theorem}

\begin{proof}
The proof is the compact Type II cost-divergence proof with the exact normalization
\(\|V(\tau)\|_{L^3}=1\) replaced by the lower bound
\(N_3(\tau)\ge\eta\).  Assume for contradiction that the total localized
renormalization cost is finite.  The unit-window selection and compactness
argument gives times \(\sigma_j\to\infty\) such that,
after passing to a subsequence,
\[
  V(\sigma_j)\to V_\infty\quad\text{strongly in }L^3_{\rm loc},
\]
the selected sequence has vanishing local dissipation on the exhaustion used by
the compactness argument, and consequently
\[
  \nabla V_\infty=0
\]
distributionally on every fixed ball.  These conclusions use only tightness,
local windowed \(H^1\) control, and finite total cost.

The low-dissipation limit is spatially constant on each fixed ball.  The
constants are compatible as the balls increase, so \(V_\infty\equiv c\) on
\(\mathbb R^3\).  The upper critical-mass bound gives, for every \(R\),
\[
  |c|^3|B_R|
  =
  \lim_{j\to\infty}\int_{B_R}|V(y,\sigma_j)|^3\,dy
  \le M^3 .
\]
Letting \(R\to\infty\) forces \(c=0\).  Thus
\[
  V(\sigma_j)\to0\quad\text{strongly in }L^3_{\rm loc}.
\]
Choose \(R\) so that the critical \(L^3\)-tail outside \(B_R\) is
\(<\eta^3/4\) uniformly in \(\tau\).  Strong convergence on \(B_R\) gives
\[
  \int_{B_R}|V(y,\sigma_j)|^3\,dy<\eta^3/4
\]
for \(j\) large.  Therefore
\[
  \|V(\sigma_j)\|_{L^3(\mathbb R^3)}^3<\eta^3/2,
\]
contradicting \(N_3(\sigma_j)\ge\eta\).  Hence the total localized
renormalization cost is infinite.
\end{proof}

\begin{theorem}[\texorpdfstring{Critical \(L^3\)-annulus classification}{Critical L3-annulus classification}]
\label{thm:c3-l3-annulus-classification}
For a repaired-gauge Type II branch, exactly one of the ordered
alternatives in \cref{def:c3-critical-mass-alternatives} occurs.  In the fourth
case the compact Type II cost-divergence criterion remains valid with the positive finite
critical annulus in place of exact unit \(L^3\)-normalization.
\end{theorem}

\begin{proof}
If \(N_3\) is not a well-defined measurable extended norm, the first
alternative holds.  Otherwise \(N_3\) is an extended nonnegative measurable
function.  If it is infinite at some time, or has infinite supremum on the
tail, the second alternative holds.  If the supremum is finite but
\(\liminf_{\tau\to\infty}N_3(\tau)=0\), the third alternative holds.  In the
remaining case,
\[
  0<\liminf_{\tau\to\infty}N_3(\tau)
  \quad\text{and}\quad
  \sup_{\tau\ge\tau_0}N_3(\tau)<\infty.
\]
After increasing the initial time, there exist
\(0<\eta\le M<\infty\) such that
\[
  \eta\le N_3(\tau)\le M
\]
for all remaining \(\tau\).  Relabeling this tail initial time as \(\tau_0\)
gives the fourth alternative.  The cost-divergence statement is
\cref{thm:c3-nonzero-mass-cost-divergence}.
\end{proof}

\subsection{Localized cost and the Type II scale-exclusion condition}

\begin{definition}[Localized Type II cost]
\label{def:c4-localized-cost}
For a repaired-gauge Type II branch, let
\[
  \tilde{\mathfrak D}_{R_0}(\tau)
  =
  \nu\int_{\mathbb R^3}|\nabla V(y,\tau)|^2\phi_{R_0}(y)\,dy
  +
  a_+(\tau)\int_{\mathbb R^3}|V(y,\tau)|^2\phi_{R_0}(y)\,dy,
\]
where \(a_+=\max(a,0)\), \(\phi_{R_0}\) is the compact Type II cutoff, and
\(\nu>0\).  A Type II scale-exclusion cost is a measurable nonnegative function
\(\mathfrak C_{\mathrm{II}}\) that is locally integrable on finite
\(\tau\)-intervals and whose divergence on \([\tau_0,\infty)\) gives the Type
II scale-exclusion condition.
\end{definition}

\begin{definition}[Cost comparison]
\label{def:c4-cost-comparison}
A cost comparison for a repaired-gauge Type II branch consists of
\(R_0,\phi_{R_0},\mathfrak C_{\mathrm{II}},c_0,\tau_1\), with \(c_0>0\) and
\(\tau_1\ge\tau_0\), such that
\[
  \mathfrak C_{\mathrm{II}}(\tau)
  \ge c_0\,\tilde{\mathfrak D}_{R_0}(\tau)
  \qquad\text{for a.e. }\tau\ge\tau_1.
\]
The canonical Navier--Stokes choice is
\[
  \mathfrak C_{\mathrm{II}}(\tau)=\tilde{\mathfrak D}_{R_0}(\tau),
  \qquad c_0=1,\qquad \tau_1=\tau_0.
\]
\end{definition}

\begin{lemma}[Identity-cost comparison]
\label{lem:c4-identity-cost}
If \(\tilde{\mathfrak D}_{R_0}\) is measurable, nonnegative, and locally
integrable on finite \(\tau\)-intervals, then the canonical choice
\(\mathfrak C_{\mathrm{II}}=\tilde{\mathfrak D}_{R_0}\) defines a valid Type II
scale-exclusion cost and satisfies the cost comparison with \(c_0=1\) and
\(\tau_1=\tau_0\).
\end{lemma}

\begin{proof}
The assumed measurability, nonnegativity, and local finite-interval
integrability give the cost properties in \cref{def:c4-localized-cost}.  The
identity
\[
  \mathfrak C_{\mathrm{II}}=\tilde{\mathfrak D}_{R_0}
\]
gives the comparison in \cref{def:c4-cost-comparison} with \(c_0=1\) and
\(\tau_1=\tau_0\).
\end{proof}

\begin{theorem}[Cost comparison theorem]
\label{thm:c4-cost-comparison}
Let \((V,P,a,b)\) be a repaired-gauge Type II branch.  Assume:
\begin{enumerate}[label=\textup{(\roman*)}]
\item \(\tilde{\mathfrak D}_{R_0}\) is measurable, nonnegative, and locally
integrable on finite \(\tau\)-intervals;
\item
\[
  \int_{\tau_0}^{\infty}
  \tilde{\mathfrak D}_{R_0}(\tau)\,d\tau=\infty;
\]
\item either the canonical cost
\(\mathfrak C_{\mathrm{II}}=\tilde{\mathfrak D}_{R_0}\) is used, or a cost
comparison as in \cref{def:c4-cost-comparison} is available.
\end{enumerate}
Then the Type II scale-exclusion condition associated with
\(\mathfrak C_{\mathrm{II}}\) holds.
\end{theorem}

\begin{proof}
In the canonical case, \cref{lem:c4-identity-cost} supplies the cost comparison.
In the general case, the comparison is assumed.  Hence there are
\(c_0>0\) and \(\tau_1\ge\tau_0\) such that
\[
  \mathfrak C_{\mathrm{II}}(\tau)
  \ge c_0\tilde{\mathfrak D}_{R_0}(\tau)
\]
for a.e. \(\tau\ge\tau_1\).  Since
\(\tilde{\mathfrak D}_{R_0}\ge0\) and is locally integrable on finite
intervals, divergence on \([\tau_0,\infty)\) implies
\[
  \int_{\tau_1}^{\infty}
  \tilde{\mathfrak D}_{R_0}(\tau)\,d\tau=\infty.
\]
Therefore
\[
  \int_{\tau_1}^{\infty}
  \mathfrak C_{\mathrm{II}}(\tau)\,d\tau
  \ge
  c_0\int_{\tau_1}^{\infty}
  \tilde{\mathfrak D}_{R_0}(\tau)\,d\tau
  =\infty.
\]
Since \(\mathfrak C_{\mathrm{II}}\) is nonnegative and locally integrable on
finite intervals,
\[
  \int_{\tau_0}^{\infty}
  \mathfrak C_{\mathrm{II}}(\tau)\,d\tau
  =
  \int_{\tau_0}^{\tau_1}\mathfrak C_{\mathrm{II}}(\tau)\,d\tau
  +
  \int_{\tau_1}^{\infty}\mathfrak C_{\mathrm{II}}(\tau)\,d\tau
  =\infty.
\]
By definition of the Type II scale-exclusion cost, this divergence is the Type II
scale-exclusion condition.
\end{proof}

\begin{corollary}[Compact cost divergence implies scale exclusion]
\label{cor:c4-compact-cost-divergence-implies-scale-exclusion}
Assume the hypotheses of \cref{thm:c3-nonzero-mass-cost-divergence} and use the canonical
Navier--Stokes cost \(\mathfrak C_{\mathrm{II}}=\tilde{\mathfrak D}_{R_0}\).
Then the Type II scale-exclusion condition holds.  If, in addition, the
supercritical scale alternative is present and the branch remains in the
retained compact cost-divergence stratum after the ordered validation tests of
\cref{paper7:lem:no-unrecorded-cost-escape}, then the Type II exclusion
conclusion holds.
\end{corollary}

\begin{proof}
\Cref{thm:c3-nonzero-mass-cost-divergence} gives
\[
  \int_{\tau_0}^{\infty}
  \tilde{\mathfrak D}_{R_0}(\tau)\,d\tau=\infty.
\]
By \cref{lem:c4-identity-cost}, the canonical cost satisfies the cost comparison.
\Cref{thm:c4-cost-comparison} converts the infinite PDE cost into the Type II
scale-exclusion condition, hence the validated compact cost-divergence
conclusion.  Under the retained-stratum validation hypothesis in the statement,
\cref{paper7:cor:retained-stratum-cost-exclusion} applies and gives the Type
II exclusion conclusion.  Thus no separate Navier--Stokes cost-exclusion
hypothesis is used.
\end{proof}

\subsection{Two-alternative finite-cost classification}

\begin{definition}[Remaining compact cost-divergence tests]
\label{def:c5-compact-tests}
For a repaired-gauge Type II orbit, define the critical-tightness
condition by
\[
  \forall\varepsilon>0\ \exists R_\varepsilon:\quad
  \sup_{\tau\ge\tau_0}
  \int_{|y|>R_\varepsilon}|V(y,\tau)|^3\,dy<\varepsilon.
\]
Failure of critical tightness is the condition
\[
  \exists\varepsilon_0>0\ \forall R>0\ \exists\tau_R\ge\tau_0:\quad
  \int_{|y|>R}|V(y,\tau_R)|^3\,dy\ge\varepsilon_0.
\]
Define local windowed \(H^1\) control by
\[
  \forall m\ge1:\quad
  \sup_{n\in\mathbb N}
  \int_{\tau_0+n}^{\tau_0+n+1}
  \|V(\tau)\|_{H^1(B_m)}^2\,d\tau<\infty.
\]
Failure of local windowed \(H^1\) control is the condition
\[
  \exists m\ge1:\quad
  \sup_{n\in\mathbb N}
  \int_{\tau_0+n}^{\tau_0+n+1}
  \|V(\tau)\|_{H^1(B_m)}^2\,d\tau=\infty.
\]
\end{definition}

\begin{definition}[Compact Type II data]
\label{def:c5-compact-typeII-data}
A Type II branch has compact Type II data if
it has:
\begin{enumerate}[label=\textup{(\roman*)}]
\item the Type II analytic data of \cref{def:c1-typeII-analytic-data};
\item the repaired-gauge representation of \cref{def:c2-representation};
\item the positive finite critical annulus of \cref{def:c3-critical-mass};
\item the locally integrable nonnegative cost and cost comparison of
\cref{def:c4-localized-cost,def:c4-cost-comparison}.
\end{enumerate}
This condition does not include critical tightness or windowed \(H^1\) control;
those are the two remaining compact cost-divergence tests.
\end{definition}

\begin{lemma}[Compact positive case is excluded]
\label{lem:c5-compact-positive-excluded}
Assume an admissible Type II branch has compact Type II data and satisfies both
compact cost-divergence tests in \cref{def:c5-compact-tests}.  If the branch
remains in the retained compact cost-divergence stratum after the ordered
validation tests of \cref{paper7:lem:no-unrecorded-cost-escape}, then the
branch satisfies the Type II exclusion conclusion.
\end{lemma}

\begin{proof}
The compact Type II data give the supercritical scale alternative,
a repaired-gauge orbit \((V,P,a,b)\), the positive finite critical annulus, and
the cost comparison.  Together with critical tightness and local windowed
\(H^1\) control, the hypotheses of \cref{thm:c3-nonzero-mass-cost-divergence} hold.
Thus
\[
  \int_{\tau_0}^{\infty}
  \tilde{\mathfrak D}_{R_0}(\tau)\,d\tau=\infty.
\]
\Cref{thm:c4-cost-comparison} converts this divergence into the Type II
scale-exclusion condition, hence the validated compact cost-divergence
conclusion.  Since the branch remains in the retained compact cost-divergence
stratum, \cref{paper7:cor:retained-stratum-cost-exclusion} gives the Type II
exclusion conclusion.
\end{proof}

\begin{theorem}[Compact exclusion or two-alternative classification]
\label{thm:c5-exclusion-or-two-alternative}
Assume every admissible Type II branch has compact Type II data and assume
adjacent-stratum closure for the compact-cost channel in the sense of
\cref{paper7:def:compact-cost-adjacent-closure}.  Then each admissible Type II
branch has exactly one ordered alternative:
\begin{enumerate}[label=\textup{(\roman*)}]
\item the Type II exclusion conclusion, either on the retained compact
cost-divergence stratum or on a treated adjacent stratum;
\item failure of critical tightness;
\item failure of local windowed \(H^1\) control
after critical tightness has been verified.
\end{enumerate}
\end{theorem}

\begin{proof}
Fix an admissible Type II branch.  Since it has compact Type II data, only the
two compact cost-divergence tests remain.  First evaluate critical tightness.  If
tightness fails, the failure-of-tightness alternative holds.  If
tightness holds, evaluate local windowed \(H^1\) control.  If the windowed
\(H^1\) condition fails, the windowed-control alternative holds.  If it
holds, \cref{paper7:thm:classwise-retained-compact-coverage} places the branch
either in the retained compact cost-divergence stratum or in a named adjacent
exit stratum.  In the retained case,
\cref{lem:c5-compact-positive-excluded} gives the Type II exclusion conclusion.
In an adjacent exit stratum, adjacent-stratum closure gives the same Type II
exclusion conclusion.  The alternatives are exhaustive and disjoint by this
ordered evaluation.
\end{proof}

\begin{definition}[Finite-cost branch not excluded by compact cost divergence]
\label{def:c5-finite-cost-not-excluded}
An admissible Type II branch is finite-cost if
\[
  \int_{\tau_0}^{\infty}\mathfrak C_{\mathrm{II}}(\tau)\,d\tau<\infty.
\]
In the canonical Navier--Stokes cost this is equivalent to finite localized PDE
cost, because
\(\mathfrak C_{\mathrm{II}}=\tilde{\mathfrak D}_{R_0}\).  It is
not excluded by compact cost divergence if the compact Type II exclusion
conclusion does not hold for that branch.
\end{definition}

\begin{corollary}[Two-alternative finite-cost classification]
\label{cor:c5-two-alternative-finite-cost}
Under the hypotheses of \cref{thm:c5-exclusion-or-two-alternative}, every
finite-cost admissible Type II branch not excluded by the retained compact
cost-divergence argument or by treated adjacent-stratum routing is either
non-tight or fails local windowed \(H^1\) control.
\end{corollary}

\begin{proof}
Apply \cref{thm:c5-exclusion-or-two-alternative}.  Since the branch is not
excluded by retained compact cost-divergence or by treated adjacent-stratum
routing, the first alternative is unavailable.  Therefore the
ordered alternative is failure of critical tightness or failure of local
windowed \(H^1\) control.  The finite-cost assumption records that the resulting
branch remains finite-cost.
\end{proof}

\begin{remark}[Content of the classification inputs]
The preceding statements put the Type II classification inputs in PDE form.
Exhaustion failure, missing repaired-gauge representation,
ambiguous critical \(L^3\) normalization, and cost-comparison failure are all
replaced by ordered alternatives.  Once those alternatives are excluded or
assumed absent, the only finite-cost Type II alternatives not excluded by the
compact cost-divergence criterion are failure of critical tightness and failure
of local windowed \(H^1\) control.
\end{remark}

\clearpage
\section{Conclusion}

The paper gives a local branchwise exclusion mechanism for Type II
Navier--Stokes concentration in the stated class of alternatives.  Beginning
with a positive Caffarelli--Kohn--Nirenberg concentration sequence, the argument
passes to repaired-gauge variables around a retained core and keeps the analysis
inside local PDE estimates: suitable energy inequalities, pressure
reconstruction, compactness on fixed cylinders, Caccioppoli estimates,
profile decoupling, and localized scale-cost identities.

The compact single-core mechanism is the central obstruction.  A retained core
with positive local CKN density and vanishing localized dissipation has only two
compact limits after passing to selected windows.  The zero limit contradicts
the persistence of the velocity-pressure concentration by strong convergence
and pressure stability, while a nonzero spatially constant limit gives a
bounded selected-window profile and hence exits the Type II regime.  The rest
of the proof shows that the apparent escape routes from this compact
alternative are also accounted for: rough cores are reduced to compact-cylinder
control, active multibubbles and cascades are separated by the profile
analysis, and finite-cost scale collapse is incompatible with a nonvanishing
retained core.

The scale-collapse analysis is intentionally explicit about its limits.  Finite
scale-collapsing cost, and hence finite absolute scale cost in the covered
setting, are excluded within the renormalized framework.  Thick autonomous
windows produce nontrivial autonomous limits only after the stated compactness
and modulation hypotheses are available; such limits are eliminated precisely in
the parameter regimes where the corresponding stationary or generalized
self-similar class is covered by a stated Liouville or rigidity theorem, such as
the Nečas--Růžička--Šverák self-similar class.  Outside those regimes the paper
records the remaining thin-drift, modulation, compactness, or rigidity
alternative rather than concealing it inside the conclusion.

Thus, under the analytic inputs and rigidity hypotheses made explicit in the
text, every covered positive local Type II concentration sequence is eliminated.
The final local dichotomy is that a singular point in the covered class must
belong to the local Type I alternative once the Type II branch has been removed.
The contribution is not an unconditional resolution of the Navier--Stokes
regularity problem; it is a precise local Type II exclusion framework whose
unresolved cases are isolated as concrete PDE compactness, decoupling, cost, or
Liouville inputs.

\clearpage
\bibliographystyle{alpha}
\bibliography{type_II_regularity}

\end{document}
